\documentclass[11pt,a4paper,oneside,openany]{book}
\usepackage[T1]{fontenc}
\usepackage[utf8]{inputenc}
\usepackage{lmodern}
\usepackage[english]{babel}
\usepackage{microtype}
\usepackage[margin=27mm,top=25mm,bottom=26mm,headheight=17pt,headsep=10mm]{geometry}
\usepackage{amsmath,amssymb,amsthm,mathtools}
\usepackage{xcolor,graphicx,tikz}
\usepackage[most]{tcolorbox}
\usepackage{titlesec,fancyhdr,enumitem,booktabs,longtable,array,needspace}
\usepackage{xurl}
\usepackage[unicode,bookmarksnumbered,bookmarksopen,colorlinks]{hyperref}
\definecolor{Navy}{HTML}{17324D}
\definecolor{Teal}{HTML}{087E8B}
\definecolor{Gold}{HTML}{B68B3D}
\definecolor{PaleTeal}{HTML}{F1F8F8}
\definecolor{PaleBlue}{HTML}{F3F6FA}
\definecolor{Ink}{HTML}{243342}
\definecolor{Muted}{HTML}{657786}
\hypersetup{linkcolor=Navy,citecolor=Teal,urlcolor=Teal,
 pdftitle={Notes on Complex Geometry},
 pdfauthor={Xiangsen Qin},
 pdfsubject={Notes on Complex Geometry}}
\urlstyle{same}
\renewcommand{\contentsname}{Contents}
\renewcommand{\bibname}{Bibliography}
\renewcommand{\proofname}{Proof}
\renewcommand{\chaptername}{Chapter}
\renewcommand{\appendixname}{Appendix}
\setcounter{tocdepth}{1}
\setcounter{secnumdepth}{2}
\makeatletter
\renewcommand*\@pnumwidth{2.4em}
\renewcommand*\@tocrmarg{3.2em}
\renewcommand*\l@section{\@dottedtocline{1}{1.5em}{3.65em}}
\makeatother
\linespread{1.08}
\setlength{\parindent}{1.25em}
\setlength{\parskip}{3pt plus 1pt}
\setlength{\emergencystretch}{2.3em}
\clubpenalty=8000
\widowpenalty=8000
\displaywidowpenalty=3000
\allowdisplaybreaks[2]
\raggedbottom
\setlist{topsep=5pt,itemsep=3pt,parsep=2pt,leftmargin=2em}
\titleformat{\chapter}[display]
 {\sffamily\bfseries\color{Navy}\raggedright}
 {\large\color{Teal}Chapter \thechapter}
 {8pt}{\Huge}[\vspace{10pt}{\color{Gold}\titlerule[1.2pt]}]
\titlespacing*{\chapter}{0pt}{6pt}{22pt}
\titleformat{\section}{\Large\sffamily\bfseries\color{Navy}\raggedright}{\thesection}{.65em}{}
\titlespacing*{\section}{0pt}{19pt plus 3pt}{8pt}
\titleformat{\subsection}{\large\sffamily\bfseries\color{Teal}}{\thesubsection}{.65em}{}
\titleformat{\paragraph}[runin]{\sffamily\bfseries\color{Navy}}{}{0pt}{}
\titlespacing*{\paragraph}{0pt}{9pt}{.6em}
\pagestyle{fancy}
\fancyhf{}
\fancyhead[L]{\small\sffamily\color{Muted}Complex Geometry}
\fancyhead[R]{\small\sffamily\color{Muted}\nouppercase{\leftmark}}
\fancyfoot[C]{\sffamily\color{Navy}\thepage}
\renewcommand{\headrulewidth}{.3pt}
\renewcommand{\chaptermark}[1]{\markboth{Chapter \thechapter}{}}
\fancypagestyle{plain}{\fancyhf{}\fancyfoot[C]{\sffamily\color{Navy}\thepage}\renewcommand{\headrulewidth}{0pt}}
\newtheoremstyle{cgplain}{7pt}{7pt}{\normalfont}{}{\sffamily\bfseries\color{Navy}}{.}{.6em}{}
\theoremstyle{cgplain}
\newtheorem{theorem}{Theorem}[chapter]
\newtheorem{proposition}[theorem]{Proposition}
\newtheorem{lemma}[theorem]{Lemma}
\newtheorem{corollary}[theorem]{Corollary}
\newtheorem{definition}[theorem]{Definition}
\newtheorem{example}[theorem]{Example}
\newtheorem{remark}[theorem]{Remark}
\tcbset{cgstatement/.style={enhanced,breakable,sharp corners,boxrule=0pt,
 colback=PaleBlue,frame hidden,borderline west={1.4pt}{0pt}{Teal},
 left=9pt,right=9pt,top=2pt,bottom=2pt,before skip=9pt,after skip=8pt,lines before break=4}}
\tcolorboxenvironment{theorem}{cgstatement,colback=PaleTeal}
\tcolorboxenvironment{proposition}{cgstatement}
\tcolorboxenvironment{lemma}{cgstatement,colback=white,borderline west={.8pt}{0pt}{Muted}}
\tcolorboxenvironment{corollary}{cgstatement,colback=white}
\newtcolorbox{learning}{enhanced,breakable,sharp corners,boxrule=0pt,
 colback=PaleTeal,frame hidden,borderline west={2pt}{0pt}{Teal},
 left=10pt,right=10pt,top=8pt,bottom=8pt,
 title={Chapter overview},fonttitle=\sffamily\bfseries,coltitle=Navy,
 attach title to upper={\quad}}
\newcommand{\C}{\mathbb C}
\newcommand{\R}{\mathbb R}
\newcommand{\Z}{\mathbb Z}
\newcommand{\N}{\mathbb N}
\newcommand{\Q}{\mathbb Q}
\newcommand{\ii}{\mathrm i}
\DeclareRobustCommand{\Cech}{\v{C}ech}
\newcommand{\dbar}{\bar\partial}
\newcommand{\ddbar}{\partial\bar\partial}
\newcommand{\cA}{\mathcal A}
\newcommand{\cE}{\mathcal E}
\newcommand{\cF}{\mathcal F}
\newcommand{\cG}{\mathcal G}
\newcommand{\cH}{\mathcal H}
\newcommand{\cI}{\mathcal I}
\newcommand{\cJ}{\mathcal J}
\newcommand{\cO}{\mathcal O}
\DeclareMathOperator{\supp}{supp}
\DeclareMathOperator{\Hom}{Hom}
\DeclareMathOperator{\Ric}{Ric}
\DeclareMathOperator{\tr}{tr}
\numberwithin{equation}{chapter}
\pdfstringdefDisableCommands{\def\dbar{dbar}\def\C{C}\def\R{R}\def\ii{i}\def\cO{O}\def\Cech{Cech}\def\mathcal#1{#1}\def\textbf#1{#1}}
\begin{document}


\hypersetup{pageanchor=false}
\begin{titlepage}
\thispagestyle{empty}
\begin{tikzpicture}[remember picture,overlay]
 \fill[Navy] (current page.north west) rectangle ([yshift=-15mm]current page.north east);
 \fill[PaleBlue] ([xshift=-44mm]current page.north east) rectangle (current page.south east);
 \draw[Gold,line width=1.2pt] ([xshift=27mm,yshift=31mm]current page.south west) -- ([xshift=-27mm,yshift=31mm]current page.south east);
 \begin{scope}[shift={([xshift=-33mm,yshift=76mm]current page.south east)}]
  \foreach \r in {12,20,28,36,44}{\draw[Teal!18,line width=.5pt] (0,0) circle (\r mm);}
  \draw[Teal!45,line width=.8pt] (-56mm,-12mm) .. controls (-33mm,50mm) and (-5mm,-43mm) .. (36mm,17mm);
 \end{scope}
\end{tikzpicture}
\vspace*{22mm}
{\sffamily\color{Teal}\Large A TWO-SEMESTER COURSE\par}
\vspace{13mm}
{\sffamily\bfseries\color{Navy}\fontsize{36}{44}\selectfont Lectures on\par Complex Geometry\par}
\vspace{7mm}
{\sffamily\color{Navy}\par}
\vspace{17mm}
{\color{Gold}\rule{31mm}{1.5pt}\par}
\vspace{10mm}
{\sffamily\color{Navy}\LARGE Xiangsen Qin\par}
\vspace{10mm}
{\large For advanced undergraduates\par and beginning graduate students\par}
\vfill
{\sffamily\color{Navy}Lecture notes\par}
\vspace{3mm}
{\color{Muted}November 2026\par}
\vspace{16mm}
\end{titlepage}
\frontmatter
\hypersetup{pageanchor=true}
\chapter*{Preface}
\addcontentsline{toc}{chapter}{Preface}
\markboth{Preface}{}

These lecture notes are based primarily on Jean-Pierre Demailly's
\emph{Complex Analytic and Differential Geometry} \cite{Demailly2012}
and \emph{Analytic Methods in Algebraic Geometry} \cite{DemaillyAnalytic2012},
and on the lecture notes I prepared while participating in the seminar
organized by Professor Fusheng Deng at the University of the Chinese
Academy of Sciences in the spring of 2025. I am grateful to Professor Deng
for organizing the seminar. Demailly's two books are the principal sources
for the mathematical material; other sources are cited where their results
or methods enter the discussion.

The aim is to explain how local complex analysis leads to global geometry.
A holomorphic function is first studied by a Cauchy integral; a bundle is
then described by its changes of frame; and the failure of local solutions
to agree becomes a cohomology class. Curvature estimates give a way to make
such classes vanish. This leads to Hodge theory, projective embeddings,
and the construction of holomorphic sections. Singular weights and heat
kernels extend the same questions to ideals, asymptotic dimension estimates,
and problems about positivity.

\section*{Intended readers}
The classroom route is intended for advanced undergraduates, usually in
their third or fourth year, who have taken proof-based courses in analysis
and algebra and are ready for sustained arguments involving estimates and
limits. Beginning graduate students may also use it as an introduction to
analytic complex geometry. A first course in one complex variable is
essential; several-variable complex analysis, sheaf cohomology, and
algebraic geometry are developed as needed rather than assumed at the start.

The book contains more material than can be taught in full in two semesters.
The teaching plan below selects a core course and indicates material for
seminar presentations or subsequent reading. In particular, the full
analytic-space theory, technical heat-kernel localization, singular-weight
estimates on general complete manifolds, index theory, and birational
applications belong to the more advanced part of the book. They are useful
destinations for the course, but are not all prerequisites for following
the first-semester arguments.

\section*{Prerequisites}
The following preparation is needed before beginning the core course.
\begin{itemize}
\item \textbf{Linear algebra and elementary algebra.}
Vector spaces over $\R$ and $\C$, dual spaces, quotient spaces, bilinear
and Hermitian forms, eigenvalues, and the spectral theorem for Hermitian
matrices. Rings, ideals, and quotients should be familiar. Tensor and
exterior products are introduced through the bundle and form calculations.
\item \textbf{Analysis.}
Differentiation in several real variables, the inverse and implicit
function theorems, compactness arguments, and integration. The analytic
chapters require Lebesgue measure, $L^p$ spaces, Fubini's theorem, and the
monotone and dominated convergence theorems. Here $L^p$ means integrability
of the $p$th power with respect to the specified measure.
\item \textbf{One-variable complex analysis.}
Cauchy's theorem and integral formula, power and Laurent series, the
identity and maximum principles, the argument principle, and residues. The several-variable
versions are obtained from these tools in Chapter~\ref{bas:chapter}.
\item \textbf{Topology and geometric language.}
Open and closed sets, compactness, connectedness, continuity, quotient
topologies, and a countable basis. Familiarity with charts and tangent
spaces is helpful. Complex manifolds, differential forms, and vector
bundles are introduced in Chapters~\ref{bas:chapter}--\ref{bun:chapter}.
\end{itemize}

Before the analytic part of the first semester, the reader should be comfortable with Hilbert
spaces, bounded operators, orthogonal complements, and weak convergence.
The arguments introduce the closed operators, spectral decompositions,
weak derivatives, and elliptic estimates that they need. The Hahn--Banach
and Baire theorems are used as standard results of functional analysis;
a reader meeting them for the first time should study them alongside the
analytic chapters. No course in partial differential equations is required
in advance, but the regularity arguments in Chapter~\ref{ana:chapter}
deserve time and worked examples.

\section*{A route through the book}
For the first semester, begin with charts, forms, and bundles. In
Chapter~\ref{sh:chapter}, concentrate first on stalks, exact sequences,
\Cech{} cocycles, and the Dolbeault resolution: these explain how local
data give a global cohomology group. The general constructions with
flabby and soft sheaves can be studied after the basic examples. In
Chapter~\ref{alg:chapter}, first learn Weierstrass division, finite
generation, and the meaning of coherence; the longer flatness and finite
projection arguments can be reserved for a seminar.

The main geometric goals are the vanishing results of
Chapter~\ref{pos:chapter}, Hodge theory in Chapter~\ref{hod:chapter},
Kodaira embedding in Chapter~\ref{emb:chapter}, and Riemann--Roch for curves
in Chapter~\ref{prj:chapter}. The flat torus, projective space, and line
bundles on curves should be kept in view while reading these arguments.
They give concrete meanings to the abstract constructions.

The second semester begins with positive currents and singular weights.
Monomial ideals provide examples on which the integrability exponents can
be computed directly before one studies general division. The heat-kernel
chapters should be read from the Gaussian and torus models toward
localization and Morse inequalities. An undergraduate course can prove
the model calculations in class and treat the longer localization estimates
in guided reading, before using the resulting theorem in geometric
applications. The applications to Nadel vanishing, section growth, and
Stein manifolds form the main continuation. Effective bounds, index theory,
and birational geometry are suitable choices for later presentations.

Chapter~\ref{tutorial:chapter} contains sixteen problems with solutions.
They are most useful when attempted before the solution is read: write
down the objects and hypotheses, identify the earlier result that applies,
and test the conclusion on a curve or a trivial bundle. The final chapter
offers directions for further study, with first steps and estimates of
the background required. Its scores concern the suitability of a student
project, not the probability of solving an open problem.

\section*{Proofs and background results}
The text gives the proofs needed along the core route and develops the
intermediate constructions explicitly. Some later applications use major
theorems from topology, analytic spaces, or algebraic geometry whose full
proofs would require separate courses. Appendix~\ref{dep:chapter} states
these results, records the hypotheses used here, and gives references.
At each such use, the text distinguishes the background theorem from the
deduction being made. A forward reference indicates the particular result
needed, so that the reader can follow the proof in a suitable order.

The two books of Demailly remain useful companions throughout the course.
Huybrechts and Wells \cite{Huybrechts2005,Wells2008} offer complementary
introductions, and the bibliography points to more specialized treatments
as the subject broadens.

\begin{flushright}
Xiangsen Qin
\end{flushright}

\chapter*{Notation and Conventions}
\addcontentsline{toc}{chapter}{Notation and Conventions}
\markboth{Notation}{}
\label{notation:guide}
The following tables collect the conventions used in the text. They are meant for consultation: unfamiliar objects are introduced in the chapters indicated, and need not be learned from the tables before starting the course.

\section*{Spaces, indices, and estimates}
\renewcommand{\arraystretch}{1.16}
\begin{longtable}{@{}>{\raggedright\arraybackslash}p{.27\textwidth}>{\raggedright\arraybackslash}p{.67\textwidth}@{}}
\toprule Notation & Meaning and convention\\ \midrule
$\C,\R,\Z,\N,\N_0$ & Complex numbers, real numbers, integers, positive integers, and nonnegative integers, respectively; $\ii^2=-1$.\\
$X$, $n$, $E$, $r$, $L$ & $X$ is a complex manifold of complex dimension $n$ (real dimension $2n$); $E$ is a vector bundle of rank $r$; $L$ is a line bundle. Extra hypotheses such as compactness are stated locally.\\
$U\Subset V$, $\supp u$ & The closure of $U$ is a compact subset of $V$. The support of $u$ is the complement of the largest open set on which it vanishes.\\
$\alpha\in\N_0^n$ & A multi-index: $|\alpha|=\sum\alpha_j$, $\alpha!=\prod\alpha_j!$, and $z^\alpha=\prod z_j^{\alpha_j}$. Differential-form indices $I,J$ are strictly increasing index lists unless stated otherwise.\\
$L^p$, $L^p_{\mathrm{loc}}$ & For $0<p<\infty$, measurable functions or sections with integrable $p$th power of their norm, globally or on every relatively compact set. A metric and a measure are specified when they matter. $L^\infty$ means essentially bounded.\\
$C$, $C_F$; $O$, $o$ & Constants may change between lines. A subscript records an allowed dependency. $a_k=O(b_k)$ means $|a_k|\le Cb_k$ for large $k$, where $b_k>0$; $a_k=o(b_k)$ means $a_k/b_k\to0$. Other limiting parameters are specified.\\
$h(v,w)$, $\langle u,v\rangle$ & Hermitian and Hilbert inner products are conjugate-linear in the first variable. For the metric matrix $H$ and column vectors in a frame, $h(v,w)=\bar v^{\mathsf T}Hw$; $A^*$ denotes the adjoint.\\
\bottomrule
\end{longtable}

\section*{Bundles, sheaves, and cohomology}
\begin{longtable}{@{}>{\raggedright\arraybackslash}p{.27\textwidth}>{\raggedright\arraybackslash}p{.67\textwidth}@{}}
\toprule Notation & Meaning and convention\\ \midrule
$\cO_X$, $\cO_{X,x}$ & The sheaf of holomorphic functions and its local ring of germs at $x$; $\mathfrak m_x$ is the ideal of germs vanishing at $x$.\\
$E_x$, $\mathcal E_x$ & $E_x$ is the vector-space fiber. If $\mathcal E$ is the sheaf of holomorphic sections of $E$, then $\mathcal E_x$ is its stalk, an $\cO_{X,x}$-module; $E_x=\mathcal E_x/\mathfrak m_x\mathcal E_x$.\\
$\Gamma(U,\mathcal F)$ & Sections of a sheaf $\mathcal F$ over $U$. In $H^q(X,E)$, a bundle denotes its sheaf of holomorphic sections. $h^q=\dim_\C H^q$ when finite; $\chi=\sum_q(-1)^qh^q$.\\
$\cA_X^{p,q}(E)$ & Smooth $E$-valued forms with $p$ holomorphic and $q$ antiholomorphic differential factors, for $0\le p,q\le n$; other degrees give the zero space. $\Omega_X^p$ denotes holomorphic $p$-forms, and $K_X=\Omega_X^n$ is the canonical line bundle.\\
$H^{p,q}_{\dbar}(X,E)$ & The quotient of $\dbar$-closed $(p,q)$-forms by $\dbar$ of $(p,q-1)$-forms. Its sheaf interpretation is Theorem~\ref{sh:dolbeault}.\\
$H^q(X,\mathcal F)$ & Sheaf cohomology as defined in Definition~\ref{sh:cohomology-definition}. $H^j_{\rm dR}(X)$ is closed smooth $j$-forms modulo exact ones; $b_j$ is its dimension.\\
$L^k$, $A+B$, $kL$ & Tensor powers and additive notation: $A+B=A\otimes B$, $-L=L^*$, and $kL=L^{\otimes k}$. An additive line-bundle expression is not a sum of curvature forms until metrics are specified.\\
$\mathbb P(E^*)$, $\mathcal O(1)$ & Lines in the dual fibers $E_x^*$, with projection $\pi:\mathbb P(E^*)\to X$. The tautological line is $\mathcal O(-1)$, and $\mathcal O(1)$ is its dual. For holomorphic $E$, $\pi_*\mathcal O(k)=\operatorname{Sym}^kE$ for $k\ge0$ by Theorem~\ref{sh:projective-direct-image}.\\
$\mathcal I_Z$, $\cI(\varphi)$ & $\mathcal I_Z$ is a specified ideal defining an analytic subspace $Z$. The multiplier ideal $\cI(\varphi)$ consists of holomorphic germs $f$ with $|f|^2e^{-\varphi}$ locally integrable. These ideals need not be radical.\\
\bottomrule
\end{longtable}

\section*{Curvature, currents, and analytic operators}
\begin{longtable}{@{}>{\raggedright\arraybackslash}p{.27\textwidth}>{\raggedright\arraybackslash}p{.67\textwidth}@{}}
\toprule Notation & Meaning and convention\\ \midrule
$d$, $\partial$, $\dbar$ & The exterior differential and its type components, $d=\partial+\dbar$. We use $d^c=\ii(\dbar-\partial)/(2\pi)$, so $dd^c=\ii\partial\dbar/\pi$; see \eqref{pos:dc-convention}.\\
$F_h$, $\Theta_h$ & Both denote the curvature of the Chern connection, with \emph{no} factor $\ii$ included. In a holomorphic line frame $|e|_h^2=e^{-\varphi}$, $F_h=\Theta_h=\partial\dbar\varphi$. The real curvature is $\ii F_h$.\\
$c_1(L,h)$ & The normalized first Chern form $\ii F_h/(2\pi)$. Its cohomology class is $c_1(L)$; integrality is explained in Lemma~\ref{prj:integral-gluing}. Higher $c_j$ and the Chern character $\operatorname{ch}$ are defined in Theorem~\ref{bun:chern-weil}.\\
$\omega$, $dV_\omega$ & A positive real $(1,1)$-form (the Hermitian metric); $dV_\omega=\omega^n/n!$. It is K\"ahler when $d\omega=0$. We write $\Lambda_\omega$ for the pointwise adjoint of multiplication by $\omega$.\\
$[A,B]$ & The operator commutator $AB-BA$, unless a graded commutator is explicitly specified. In $[\ii F_h,\Lambda_\omega]$, curvature acts by wedge multiplication and on the bundle coefficient.\\
$[T]$, $[D]$ & $[T]$ is the cohomology class of a closed current $T$; $[D]$ integrates over a divisor with its multiplicities. For holomorphic $f$ not identically zero, $\ii\partial\dbar\log|f|^2/(2\pi)=[\operatorname{div}f]$.\\
$\dbar^*$, $\Delta''$ & The adjoint for the stated metrics and measure, and $\Delta''=\dbar\dbar^*+\dbar^*\dbar\ge0$. Closed extensions and their domains are specified before spectral or $L^2$ use. The Morse scaling is $e^{-2t\Delta''/k}$.\\
$\Delta_d$, $\Delta_{\rm eucl}$ & $\Delta_d=dd^*+d^*d$ is nonnegative. When used for potentials, $\Delta_{\rm eucl}=\sum_a\partial_{x_a}^2$ has the opposite sign. The chosen sign is stated in each local calculation.\\
$L\cdot C$, $A^n$ & Intersection numbers: $L\cdot C=\int_Cc_1(L)$ and $A^n=\int_Xc_1(A)^n$ on a compact $n$-fold. On a surface, $A\cdot B=\int_Xc_1(A)c_1(B)$.\\
\bottomrule
\end{longtable}


\chapter*{A Two-Semester Teaching Plan}
\addcontentsline{toc}{chapter}{A Two-Semester Teaching Plan}
\markboth{Teaching Plan}{}
The following plan allows sixteen weeks per semester, with three lectures
and one tutorial each week. It describes a selection from the book.
Long technical proofs listed as further reading can be discussed in a
seminar; they need not displace the examples and exercises in class.

\section*{Semester I: Local analysis and global geometry}
\footnotesize
\renewcommand{\arraystretch}{1.05}
\begin{longtable}{@{}p{.065\textwidth}>{\raggedright\arraybackslash}p{.425\textwidth}>{\raggedright\arraybackslash}p{.42\textwidth}@{}}
\toprule Week & Classroom topics & Examples and further reading\\ \midrule
1 & Holomorphic functions and Cauchy estimates. Chapter~\ref{bas:chapter}. & Power series on a polydisc; inverse function theorem.\\
2 & Complex manifolds and holomorphic maps. Chapter~\ref{bas:chapter}. & Projective charts and complex tori.\\
3 & Differential forms, types, and Stokes. Chapter~\ref{bas:forms-chapter}. & Radial homotopy and orientation signs.\\
4 & Holomorphic bundles and Hermitian metrics. Chapter~\ref{bun:chapter}. & Transition functions and the tautological bundle.\\
5 & The Chern connection and curvature. Chapter~\ref{bun:chapter}. & The curvature of $\mathcal O(1)$; symmetric powers as further reading.\\
6 & Sheaves, stalks, and exact sequences. Chapter~\ref{sh:chapter}. & Local solvability versus global sections.\\
7 & \Cech{} cocycles and the Dolbeault resolution. Chapter~\ref{sh:chapter}. & A two-chart computation on $\mathbb P^1$ (Tutorial~\ref{tut:line-cech}).\\
8 & Cohomology and a first synthesis. Chapter~\ref{sh:chapter}. & Dolbeault groups of a torus; general resolution machinery as further reading.\\
9 & Weierstrass division and coherence. Chapter~\ref{alg:chapter}. & Ideals of holomorphic germs; flatness and finite projections as further reading.\\
10 & Weak derivatives and Sobolev spaces. Chapter~\ref{ana:chapter}. & Difference quotients and integration by parts.\\
11 & Elliptic regularity and compactness. Chapter~\ref{ana:chapter}. & How a weak solution becomes smooth.\\
12 & Plurisubharmonic functions and curvature positivity. Chapter~\ref{pos:chapter}. & The complex Hessian and logarithmic potentials.\\
13 & The basic $L^2$ estimate and vanishing. Chapter~\ref{pos:chapter}. & The Gaussian weight and a line-bundle calculation.\\
14 & Hodge decomposition and duality. Chapter~\ref{hod:chapter}. & Harmonic forms on a flat torus.\\
15 & Kodaira embedding. Chapter~\ref{emb:chapter}. & Peak sections and separation of points and tangent vectors.\\
16 & Curves and Riemann--Roch. Chapter~\ref{prj:chapter}. & Degree, genus, and synthesis problems; projective bundles as further reading.\\
\bottomrule
\end{longtable}
\normalsize

\clearpage
\section*{Semester II: Singularities, estimates, and applications}
\footnotesize
\begin{longtable}{@{}p{.065\textwidth}>{\raggedright\arraybackslash}p{.425\textwidth}>{\raggedright\arraybackslash}p{.42\textwidth}@{}}
\toprule Week & Classroom topics & Examples and further reading\\ \midrule
1 & Positive currents and local potentials. Chapter~\ref{cur:chapter}. & Integration currents and measure coefficients.\\
2 & Singular metrics and multiplier ideals. Chapters~\ref{pos:chapter}, \ref{l2:chapter}. & Logarithmic poles and radial integrability.\\
3 & Complete metrics and weighted solutions. Chapter~\ref{l2:chapter}. & Cutoff functions; coherence as guided reading.\\
4 & Division and integral closure. Chapters~\ref{div:chapter}, \ref{bs:chapter}. & Monomial ideals before the general division argument.\\
5 & Extension and integrability. Chapter~\ref{int:chapter}. & One complex line and strict threshold inequalities.\\
6 & Hilbert complexes and heat operators. Chapter~\ref{heat:chapter}. & Finite-dimensional alternating sums and spectral damping.\\
7 & Gaussian kernels and model geometries. Chapter~\ref{heat:models-chapter}. & The flat torus and the constant-curvature model.\\
8 & Heat-kernel localization. Chapter~\ref{heat:local-chapter}. & Scaling and Duhamel's formula; full estimates as guided reading.\\
9 & Holomorphic Morse inequalities. Chapter~\ref{morse:chapter}. & Curvature signatures and section-growth estimates.\\
10 & Nadel vanishing. Chapter~\ref{nad:chapter}. & Multiplier ideals in an exact sequence; general approximation as guided reading.\\
11 & Generation and projective bundles. Chapters~\ref{sh:projective-chapter}, \ref{amp:chapter}. & Rank-one examples; direct-image curvature as further reading.\\
12 & Numerical positivity and big line bundles. Chapters~\ref{gp:chapter}, \ref{nad:applications}. & Counting jets and sections; a fixed twist can change growth.\\
13 & Stein manifolds and the Levi problem. Chapter~\ref{stv:chapter}. & Interpolation by a growing weight; noncompact vanishing as further reading.\\
14 & Adjoint bundles and jets. Chapter~\ref{eff:chapter}. & Selected effective bounds; the surface argument as further reading.\\
15 & Geometric examples and presentations. Chapters~\ref{exg:chapter}, \ref{pos:thin-chapter}. & Canonical bundles of curves and hypersurfaces; currents on null sets.\\
16 & Synthesis and further study. Chapters~\ref{tutorial:chapter}, \ref{research:chapter}. & Student presentations and the formulation of a manageable project.\\
\bottomrule
\end{longtable}
\normalsize

\paragraph{Topics for a subsequent seminar.}
Chapters~\ref{heat:extensions-chapter}, \ref{csg:chapter},
\ref{idx:chapter}, and \ref{bir:chapter} develop spectral extensions,
general singular estimates, index theory, and birational applications.
These can be studied after the corresponding core arguments, with
Appendix~\ref{dep:chapter} and the cited references at hand.

\tableofcontents
\mainmatter


\part{Semester I: Local Analysis and Global Geometry}


\chapter{From Holomorphic Functions to Complex Manifolds}\label{bas:chapter}
\begin{learning}
This chapter carries the familiar one-variable Cauchy formula into several variables and explains why the definition of a complex manifold must include both local coordinates and global topological conditions. Complex projective space and complex tori will serve as two recurring models throughout the book.
\end{learning}

\section{Basic Notation and Analytic Preliminaries}
We write $\ii^2=-1$, $\N_0=\{0,1,2,\ldots\}$, and use $\C^n$ with its Euclidean norm $|z|^2=\sum_j|z_j|^2$.
For an open set $U$, $C^k(U)$ means functions with continuous real derivatives through order $k$; $C^\infty(U)$ means smooth functions.
The support $\supp f$ is the closure of the set where $f$ is nonzero, and $C_c^\infty(U)$ denotes smooth functions whose support is a compact subset of $U$.
We write $V\Subset U$ when $\overline V$ is compact and contained in $U$.
Write $z_j=x_j+\ii y_j$ and define the Wirtinger derivatives
\[
 \partial_{z_j}=\tfrac12(\partial_{x_j}-\ii\partial_{y_j}),\qquad
 \partial_{\bar z_j}=\tfrac12(\partial_{x_j}+\ii\partial_{y_j}).
\]
A Hermitian inner product on a vector space is conjugate-linear in its first argument and linear in its second; on $\C^n$ the standard one is $\langle v,w\rangle=\sum_j\bar v_jw_j$.
Unless otherwise stated, manifolds are Hausdorff, second countable, and without boundary; integration on a complex manifold always uses its natural orientation.
We assume the one-variable Cauchy formula, the real inverse function theorem, Fubini's theorem, the dominated convergence theorem, and the basic properties of complete metric spaces from undergraduate courses.
The arguments below use these facts in the same form as in one-variable analysis. References \cite{Huybrechts2005,Demailly2012} provide further examples and geometric context.

\begin{definition}[Polydiscs and holomorphic functions]\label{bas:holomorphic-def}
Set $\Delta(a;r)=\{z\in\C^n:|z_j-a_j|<r_j,\ 1\le j\le n\}$.
A function on an open set is \emph{holomorphic} if, near every point, it can be expressed as a convergent power series in several variables,
$\sum_{\alpha\in\N_0^n}c_\alpha(z-a)^\alpha$. Holomorphicity for vector-valued functions is defined componentwise. A map between open subsets of complex Euclidean spaces is biholomorphic if it is bijective, holomorphic, and has a holomorphic inverse.
Here a \emph{multi-index} is $\alpha=(\alpha_1,\ldots,\alpha_n)\in\N_0^n$;
$\alpha!=\prod_j\alpha_j!$, $|\alpha|=\sum_j\alpha_j$, $(z-a)^\alpha=\prod_j(z_j-a_j)^{\alpha_j}$, and
$\partial^\alpha=\partial_{z_1}^{\alpha_1}\cdots\partial_{z_n}^{\alpha_n}$.
The unit disc is $\Delta=\{z\in\C:|z|<1\}$, and $\Delta^n$ is its $n$-fold product.
\end{definition}

\begin{theorem}[The Cauchy formula and the Cauchy--Riemann criterion]\label{bas:cauchy}
Let $f\in C^1(U)$, where $U\subset\C^n$ is open. Then $f$ is holomorphic if and only if
$\partial_{\bar z_j}f=0$ for every $j$. If $f$ is holomorphic and the closed polydisc $\overline{\Delta(a;r)}$ lies in $U$, then in its interior
\begin{equation}\label{bas:cauchy-formula}
 f(z)=\frac1{(2\pi\ii)^n}\int_{|\zeta_1-a_1|=r_1}\!\cdots\!
 \int_{|\zeta_n-a_n|=r_n}
 \frac{f(\zeta)\,d\zeta_n\cdots d\zeta_1}{\prod_{j=1}^n(\zeta_j-z_j)}.
\end{equation}
In particular, if $|f|\le M$ on the boundary torus, then
$|\partial^\alpha f(a)|\le \alpha!M\prod r_j^{-\alpha_j}$.
\end{theorem}
\begin{proof}
We first explain how the one-variable Cauchy--Riemann equation is used here.
If $f$ is $C^1$ in real coordinates and $f_{\bar z}=0$, then $f_y=\ii f_x$, so its real differentiability expansion becomes
$f(z+h)=f(z)+f_x(z)(\operatorname{Re}h+\ii\operatorname{Im}h)+o(|h|)$.
The complex difference quotient therefore has limit $f_x(z)$, and the one-variable Cauchy formula applies.
Conversely, a convergent power series can be differentiated term by term on each smaller closed disc, giving $f_{\bar z}=0$ directly.

In several variables, fixing the other coordinates reduces the hypothesis to one-variable holomorphicity in each variable.
First fix $z_2,\ldots,z_n$ and apply the one-variable formula to $z_1$:
\[
 f(z)=\frac1{2\pi\ii}\int_{|\zeta_1-a_1|=r_1}
 \frac{f(\zeta_1,z_2,\ldots,z_n)}{\zeta_1-z_1}\,d\zeta_1.
\]
Repeat with the second variable in the integrand, and continue until all variables have been replaced by integration variables.
The product of the circles is compact, $f$ is continuous, and each interior point $z_j$ has positive distance from the corresponding circle.
Thus the integrand is absolutely integrable, and Fubini's theorem justifies these finitely many iterations, yielding \eqref{bas:cauchy-formula}.

To recover the jointly convergent power series required by the definition, fix $s_j<r_j$.
For $|z_j-a_j|\le s_j$, the geometric series
\[
 \frac1{\zeta_j-z_j}
 =\sum_{m\ge0}\frac{(z_j-a_j)^m}{(\zeta_j-a_j)^{m+1}}
\]
is uniformly bounded termwise by the summable series $r_j^{-1}\sum_m(s_j/r_j)^m$.
The product of these finitely many series also converges uniformly and absolutely, so we may integrate term by term:
\[
 f(z)=\sum_{\alpha\in\N_0^n}c_\alpha(z-a)^\alpha,
 \quad c_\alpha=\frac1{(2\pi\ii)^n}
 \int\cdots\int\frac{f(\zeta)\,d\zeta_n\cdots d\zeta_1}
 {\prod_j(\zeta_j-a_j)^{\alpha_j+1}}.
\]
The product of the contour lengths is $\prod_j2\pi r_j$, hence
$|c_\alpha|\le M\prod_jr_j^{-\alpha_j}$.
Termwise differentiation on a smaller polydisc then gives
$\partial^\alpha f(a)=\alpha!c_\alpha$, proving the derivative estimate.
This argument also explains why, under the $C^1$ regularity assumption, the separate Cauchy--Riemann equations imply joint holomorphicity.
\end{proof}

\begin{corollary}[Uniform limits and the identity principle]\label{bas:identity}
A limit of holomorphic functions that is uniform on compact sets is holomorphic. A holomorphic function on a connected open set that vanishes on a nonempty open subset vanishes everywhere.
\end{corollary}
\begin{proof}
Suppose $f_\nu\to f$ uniformly on compact sets.
On any polydisc whose closure lies in the domain, the convergence is uniform on the boundary torus, and the integral kernel at an interior point is continuous and bounded. We may therefore pass to the limit in \eqref{bas:cauchy-formula}.
The resulting function $f$ satisfies the same integral formula; expansion of its kernel as a geometric series gives a convergent power series, so $f$ is holomorphic.
Differentiating the integral kernel also shows that derivatives of every order converge uniformly on each smaller polydisc.

Now let $S$ be the set where $f$ and all its complex partial derivatives vanish simultaneously.
Every derivative is continuous, and an arbitrary intersection of closed sets is closed, so $S$ is closed.
If $a\in S$, all Taylor coefficients furnished by the Cauchy formula vanish, and hence $f$ vanishes identically on a polydisc about $a$.
All derivatives vanish on that polydisc as well, so it is contained in $S$; thus $S$ is open.
The nonempty open set on which $f$ vanishes lies in $S$, so $S$ is nonempty. Connectedness forces $S=U$, and therefore $f=0$.
The argument requires vanishing on an open set, or vanishing of all Taylor coefficients at one point; vanishing of the value at a single point is insufficient.
\end{proof}

\begin{theorem}[Hartogs extension across an isolated hole]\label{bas:hartogs}
For $n\ge2$, every holomorphic function on $\Delta^n\setminus\{0\}$ extends uniquely to a holomorphic function on $\Delta^n$.
The same conclusion holds when the deleted set is $\{z_1=z_2=0\}$.
\end{theorem}
\begin{proof}
Write $z=(z_1,z')$ and fix $0<r<1$. For $|z_1|<r$ and $z'\in\Delta^{n-1}$, define
\[
 F(z_1,z')=\frac1{2\pi\ii}\int_{|\zeta|=r}\frac{f(\zeta,z')}{\zeta-z_1}\,d\zeta.
\]
The contour avoids the deleted set. On each compact subset the denominator has a positive lower bound, so we may differentiate with respect to all parameters; Theorem \ref{bas:cauchy} shows that $F$ is holomorphic.
For $z'\ne0$, the function $f(\cdot,z')$ is holomorphic on the entire unit disc, and the Cauchy formula gives $F=f$.
The two functions agree on the dense set $z'\ne0$, hence throughout the intersection of the old and new domains, and therefore glue together.
Uniqueness follows from continuity, or from Corollary \ref{bas:identity}.
When $\{z_1=z_2=0\}$ is deleted, repeat the same construction and first check $F=f$ where $z_2\ne0$; the remaining variables serve as parameters.
\end{proof}

\begin{remark}
In one variable, the isolated singularity of $1/z$ cannot be removed. Extension phenomena in higher dimensions show that one-variable intuition about singularities cannot be carried over unchanged.
Removability for general analytic sets requires the theory of analytic sets; see \cite[Chapter II]{Demailly2012}. The short proof above already supplies the coordinate-subspace version needed in this book.
\end{remark}

\begin{theorem}[Holomorphic inverse and implicit function theorems]\label{bas:holomorphic-inverse}
If $F:U\subset\C^n\to\C^n$ is holomorphic and $\det(\partial F/\partial z)(a)\ne0$, then $F$ has a holomorphic inverse near $a$.
If $G(z,w)$ is holomorphic, vanishes at a point, and has an invertible Jacobian with respect to the $r$ variables $w$ there, then near that point the solutions of $G(z,w)=0$ are uniquely of the form $w=g(z)$, with $g$ holomorphic.
\end{theorem}
\begin{proof}
The real differential of $F$ is the realification of its complex-linear differential, and its real determinant is the squared modulus of the complex determinant.
The real inverse function theorem therefore gives a local $C^1$ inverse $H$.
The chain rule gives $dH=(dF)^{-1}$. Since the inverse of an invertible complex-linear map is complex-linear, the real differential of $H$ commutes with multiplication by $\ii$ and thus satisfies all the Cauchy--Riemann equations.
Theorem \ref{bas:cauchy} now shows that $H$ is holomorphic.

For the implicit function theorem, consider $\Phi(z,w)=(z,G(z,w))$.
Its complex Jacobian is block triangular with diagonal blocks $I$ and $G_w$, and is therefore invertible.
The first group of coordinates of its local inverse must still be $z$, so write the inverse as $(z,H(z,\eta))$.
Setting $\eta=0$ gives the holomorphic function $g(z)=H(z,0)$, with $G(z,g(z))=0$.
Uniqueness follows from local injectivity of $\Phi$.
\end{proof}

\section{Gluing Coordinate Charts}
A coordinate chart is meant to identify an open part of a manifold with an open subset of $\C^n$.
The openness requirement is essential: agreement of topologies on overlaps alone does not ensure that local charts fit into one topological space.
\begin{example}[Agreement of induced topologies on intersections is insufficient]\label{bas:gluing-counterexample}
Take $X=\{a,b,c\}$. Declare the open subsets of $A_1=\{a,b\}$ to be $\varnothing,\{a\},A_1$; those of $A_2=\{b,c\}$ to be $\varnothing,\{b\},A_2$; and those of $A_3=\{c,a\}$ to be $\varnothing,\{c\},A_3$.
Every pairwise intersection is a singleton, so the induced topologies agree. Nevertheless, no topology on $X$ induces all three topologies.
Indeed, a global open set $O$ satisfying $O\cap A_1=\{a\}$ would contain $a$ but not $b$.
The topology on $A_3$ would then force $c\in O$, and the topology on $A_2$ would force $b\in O$, a contradiction.
Thus openness of the chart domains is a genuine hypothesis in the gluing proposition below.
\end{example}

\begin{proposition}[Gluing open coordinate charts]\label{bas:atlas}
Suppose a set $X$ is covered by subsets $U_\alpha$, with bijections
$\varphi_\alpha:U_\alpha\to V_\alpha\subset\C^n$, where each $V_\alpha$ is open.
Assume $\varphi_\alpha(U_\alpha\cap U_\beta)$ is open and every transition map
$\varphi_\beta\varphi_\alpha^{-1}$ is biholomorphic.
There is a unique topology for which every $U_\alpha$ is open and every $\varphi_\alpha$ is a homeomorphism.
If a countable collection of charts suffices and this topology is Hausdorff, the data define a complex manifold.
\end{proposition}
\begin{proof}
Use all sets $\varphi_\alpha^{-1}(W)$, with $W\subset V_\alpha$ open, as a proposed basis.
If $x$ belongs to two such sets, their intersection, expressed in the $\alpha$ chart, is
$W\cap\varphi_\alpha(U_\alpha\cap U_\beta)\cap
(\varphi_\beta\varphi_\alpha^{-1})^{-1}(W')$, which is open by hypothesis.
The proposed basis therefore generates a topology, and each coordinate map is a continuous open bijection.
Any topology with the required properties must contain these basis elements. Conversely, the intersection of any of its open sets with each $U_\alpha$ is a basis element of this form, proving uniqueness.
Each $V_\alpha$ has a countable basis of balls with rational centers and rational radii; the union of these bases over countably many charts is still countable.
The Hausdorff condition is an additional global requirement. Together with the local conditions above, it gives precisely the definition of a complex manifold.
\end{proof}

\begin{definition}[Holomorphic maps and tangent vectors]
A map between complex manifolds is \emph{holomorphic} if it is holomorphic in local complex coordinates.
Tangent vectors can be defined either as velocities of curves or as real derivations on germs of smooth functions satisfying the Leibniz rule.
Here two functions have the same \emph{germ} at a point if they agree on some neighborhood of that point.
In complex coordinates, the real tangent space has basis $\partial_{x_j},\partial_{y_j}$. Multiplication by $\ii$ defines a real-linear map $J$ satisfying $J^2=-I$.
The differentials of biholomorphic transition maps are complex-linear, so $J$ is independent of the coordinates.
\end{definition}

\begin{proposition}[Complex projective space]\label{bas:projective-space}
The space $\mathbb P^n$ obtained from $\C^{n+1}\setminus\{0\}$ by identifying vectors that differ by a nonzero complex scalar is a compact complex manifold.
Its standard charts are $U_j=\{[Z_0:\cdots:Z_n]:Z_j\ne0\}$, with coordinates $Z_k/Z_j$ for $k\ne j$.
\end{proposition}
\begin{proof}
Each chart is in bijection with $\C^n$. On the intersection of two charts, the coordinate change consists of reciprocating a nonzero coordinate and taking quotients of coordinates; thus it is biholomorphic, and the intersection is open.
For a column vector or matrix, a superscript $*$ denotes conjugate transpose.
To verify the global topology, consider the map
\[
 [Z]\longmapsto \frac{ZZ^*}{Z^*Z}
\]
into the space of Hermitian matrices. It is independent of the representative, and its image consists of rank-one orthogonal projections. Equal projections have the same image line, so the map is injective.
The map from the unit sphere $S^{2n+1}$ onto these projections is continuous and surjective, so its image is compact and Hausdorff.
On $U_j$, if the projection matrix is $P$, then $Z_k/Z_j=P_{kj}/P_{jj}$, proving continuity of the coordinates; the normalization formula also proves continuity of the inverse map.
Thus this compact Hausdorff topology is exactly the topology of Proposition \ref{bas:atlas}. The finitely many standard charts ensure second countability.
\end{proof}

\begin{proposition}[Complex tori]\label{bas:torus}
If $\Lambda\subset\C^n\simeq\R^{2n}$ consists of the integer linear combinations of a real basis, then $X=\C^n/\Lambda$ is a compact complex manifold.
\end{proposition}
\begin{proof}
Write $\Lambda=T\Z^{2n}$, where $T:\R^{2n}\to\R^{2n}$ is an invertible real-linear map.
Its inverse has finite operator norm $\|T^{-1}\|=\sup_{|v|=1}|T^{-1}v|$, so
$|T m|\ge\|T^{-1}\|^{-1}|m|$. The lengths of nonzero lattice vectors therefore have a uniform positive lower bound $\delta$.
A bounded set contains only finitely many lattice points, because its image under $T^{-1}$ is bounded and contains only finitely many integer points.

The quotient map $q:\C^n\to X$ is open: for any open set $W$,
$q^{-1}q(W)=\bigcup_{\lambda\in\Lambda}(W+\lambda)$ is open.
If a ball has radius less than $\delta/3$, the difference of two points in it cannot be a nonzero lattice vector.
Thus $q$ is injective on the ball and, since it is open, is a homeomorphism onto its image.
On the intersection of two charts, their local inverses differ by a lattice vector.
This vector depends continuously on the point and takes values in a discrete set, so it is constant on each connected component of the intersection.
The transition map on each such component is therefore a translation, hence holomorphic.

If $q(a)\ne q(b)$, the closed discrete set $a-b-\Lambda$ does not contain zero and meets every bounded set in finitely many points. Consequently,
$d=\inf_{\lambda\in\Lambda}|a-b-\lambda|>0$.
The quotient images of balls about $a$ and $b$ of radius less than $d/3$ are disjoint: otherwise two points in these balls would differ by a lattice vector, and the triangle inequality would give $d<2d/3$.
The quotient is therefore Hausdorff.
Images of balls with rational centers and rational radii form a countable basis, because $q$ is open.
Finally, $T[0,1]^{2n}$ is a compact closed fundamental parallelepiped containing a representative of every lattice coset.
It maps continuously onto $X$, so $X$ is compact.
\end{proof}

\begin{theorem}[Compactness and holomorphic functions]\label{bas:compact-holomorphic}
Every holomorphic function on a connected compact complex manifold is constant.
\end{theorem}
\begin{proof}
The continuous function $|f|$ attains its maximum at a point $p$. Choose a small coordinate ball centered at $p$.
Restrict $f$ to each complex line through the origin. The one-variable maximum modulus principle shows that the restriction is constant on the corresponding small disc, so $f=f(p)$ throughout the ball.
The identity principle, Corollary~\ref{bas:identity}, propagates this equality through overlapping coordinate balls: the set where $f-f(p)$ and all its local derivatives vanish is nonempty, open, and closed.
By connectedness, it is the entire manifold.
\end{proof}

\section{The Size of Analytic Zero Sets}
We introduce a brief notion of size. Cover a set $A$ by sets of diameter at most $\delta$, take the infimum of $\sum_i(\operatorname{diam}U_i)^s$ over all such covers, and then let $\delta\downarrow0$.
The result is the $s$-dimensional Hausdorff measure; its normalization is irrelevant here.
The infimum of the exponents $s$ for which this measure is zero is the \emph{Hausdorff dimension}.
A bounded $d$-dimensional cube can be covered by $O(\delta^{-d})$ small cubes of diameter $O(\delta)$.
A Lipschitz map multiplies diameters by at most a fixed constant, so the Hausdorff dimension of its image is at most $d$.
This estimate is all that is needed below to deduce dimension bounds from smooth coordinate patches.

\begin{proposition}[Dimension bounds for real and complex analytic zero sets]\label{bas:analytic-zeros}
The zero set of a nonzero real-valued real analytic function on a connected open set $U\subset\R^m$ is contained in a countable union of smooth real hypersurfaces. It therefore has Hausdorff dimension at most $m-1$ and Lebesgue measure zero.
The zero set of a nonzero holomorphic function on a connected open set $U\subset\C^n$ has Hausdorff dimension at most $2n-2$.
\end{proposition}
\begin{proof}
In the real case there can be no point where all partial derivatives vanish. At such a point the real Taylor expansion would make the function zero on a neighborhood; the common zero set of all derivatives is also closed, and the same open-and-closed argument would force the function to vanish identically.
If $f(x)=0$, let $k\ge1$ be the least order of a derivative that is nonzero at $x$.
Choose a multi-index $\alpha$ of length $k-1$ and a coordinate $j$ such that
$\partial^\alpha f(x)=0$ but $\partial_j\partial^\alpha f(x)\ne0$.
The real implicit function theorem shows that $\{\partial^\alpha f=0\}$ is a smooth hypersurface near such a point.
There are countably many multi-indices, and every regular level set admits a countable cover by coordinate patches, giving the asserted covering.
A compact coordinate patch of a smooth hypersurface is a Lipschitz image of a bounded subset of $\R^{m-1}$ and thus has Hausdorff dimension at most $m-1$. A countable union preserves this bound and consequently has zero $m$-dimensional measure.
The complex case is analogous, using complex derivatives. If $g=\partial^\alpha f$ satisfies $g(x)=0$ and $\partial_jg(x)\ne0$, the holomorphic inverse function theorem makes $g=0$ a complex hypersurface of real dimension $2n-2$.
Theorem~\ref{bas:holomorphic-inverse} supplies the required holomorphic coordinates, so the same countable-cover argument gives the complex dimension bound.
\end{proof}

\begin{remark}
The proposition does not assert that the zero set itself is smooth everywhere. For example, $z_1z_2=0$ has a crossing singularity.
The proof places each zero in a regular zero set of a suitable derivative; this is an efficient way to handle singular points.
\end{remark}

\chapter{Differential Forms and the Language of K\"ahler Geometry}\label{bas:forms-chapter}
\begin{learning}
Complex geometry retains both the real exterior derivative $d$ and the complex bidegree $(p,q)$. This chapter relates them and explains how integration, partitions of unity, and Hermitian metrics turn local calculations into constructions on the whole manifold.
\end{learning}

\section{Complexification, Types, and Two Differentials}
A smooth complex vector bundle of rank $r$ is locally a product $U\times\C^r$, with smooth invertible complex-linear transition maps on the fibers.
Write $T_xX$ for the real tangent space at $x$, $T_x^*X$ for its real dual, and $TX,T^*X$ for the corresponding bundles.
A smooth complex-valued differential $k$-form is a smooth section of $\Lambda^k(T^*X\otimes_\R\C)$, that is, an alternating complex-valued real-multilinear function of $k$ tangent vectors at each point.
We denote its space by $A^k(X)$ and put $A^k(X)=0$ for $k<0$.
The exterior product $\wedge$ is characterized in coordinates by alternating the one-form factors, so $\alpha\wedge\beta=(-1)^{ab}\beta\wedge\alpha$ for forms of degrees $a,b$.
For an increasing index set $I=(i_1<\cdots<i_k)$, write $dx_I=dx_{i_1}\wedge\cdots\wedge dx_{i_k}$; the notations $dz_I,d\bar z_I$ have the same meaning.
The exterior derivative is given in real coordinates by
$d(\sum_I a_I dx_I)=\sum_{j,I}(\partial_j a_I)dx_j\wedge dx_I$.
The chain rule makes this definition coordinate independent; equality of mixed partial derivatives and antisymmetry give $d^2=0$ and the graded Leibniz rule.
For a smooth map $F:Y\to X$, its pullback is
$(F^*\alpha)_y(v_1,\ldots,v_k)=\alpha_{F(y)}(dF_yv_1,\ldots,dF_yv_k)$.
The same coordinate formula and chain rule show $dF^*=F^*d$.

\begin{proposition}[Type decomposition]\label{bas:type-decomposition}
On a complex manifold, the complexified cotangent bundle splits as
$T^*X\otimes_\R\C=T^{*1,0}X\oplus T^{*0,1}X$, with local bases $dz_j$ and $d\bar z_j$, respectively.
Every complex-valued $k$-form therefore has a unique decomposition into $(p,q)$-forms, where $p+q=k$, and
\[
 d=\partial+\dbar,\qquad
 \partial^2=\dbar^2=0,\qquad
 \partial\dbar+\dbar\partial=0.
\]
\end{proposition}
\begin{proof}
On a real cotangent space, $J^2=-I$, so after complexification the induced map $J^*$ satisfies
$(J^*-\ii I)(J^*+\ii I)=0$.
The two roots are distinct. The projections
$\tfrac12(I-\ii J^*)$ and $\tfrac12(I+\ii J^*)$
uniquely decompose each covector into the $\ii$ and $-\ii$ eigenspaces.
From $J\partial_{x_j}=\partial_{y_j}$ and $J\partial_{y_j}=-\partial_{x_j}$, we obtain directly
$J^*dz_j=\ii dz_j$ and $J^*d\bar z_j=-\ii d\bar z_j$.
For a holomorphic change of coordinates, $dw_j$ contains no $d\bar z_k$ terms, so these subspaces are preserved by coordinate changes and define global subbundles.

Choose $p$ factors $dz$ and $q$ factors $d\bar z$, placing all $dz$ factors first.
The forms $dz_I\wedge d\bar z_J$ form a basis of the exterior algebra; distinct pairs $(p,q)$ use disjoint subsets of this basis, so the type decomposition is direct and unique.
For $u=\sum u_{IJ}dz_I\wedge d\bar z_J$, define
\[
 \partial u=\sum_{I,J,k}(\partial_{z_k}u_{IJ})
       dz_k\wedge dz_I\wedge d\bar z_J,
 \quad
 \dbar u=\sum_{I,J,k}(\partial_{\bar z_k}u_{IJ})
       d\bar z_k\wedge dz_I\wedge d\bar z_J.
\]
These terms increase the holomorphic and antiholomorphic degrees, respectively, and their sum is $du$.
Moving $d\bar z_k$ past $dz_I$ in the second expression introduces the factor $(-1)^p$; this accounts for the signs in later local solution formulas.
Finally, the three types occurring in $0=d^2u$ are
$\partial^2u$, $(\partial\dbar+\dbar\partial)u$, and $\dbar^2u$.
Uniqueness of the direct-sum decomposition forces each part to vanish.
\end{proof}

\begin{definition}[Two cohomology theories]\label{bas:cohomology-def}
Let $A^{p,q}(X)$ denote the space of smooth $(p,q)$-forms, with this space understood to be zero if either index is negative or exceeds the complex dimension.
Thus $A^k(X)=\bigoplus_{p+q=k}A^{p,q}(X)$ by Proposition~\ref{bas:type-decomposition}. Define
\[
 H^k_{\mathrm{dR}}(X,\C)=\frac{\ker(d:A^k\to A^{k+1})}{dA^{k-1}},\qquad
 H^{p,q}_{\dbar}(X)=\frac{\ker(\dbar:A^{p,q}\to A^{p,q+1})}{\dbar A^{p,q-1}}.
\]
Write $[\alpha]$ for the cohomology class of a closed form $\alpha$. Replacing complex-valued forms by real-valued forms defines $H^k_{\mathrm{dR}}(X,\R)$. For $q=0$, the denominator in the second expression is zero. The wedge product satisfies the graded Leibniz rule
$d(\alpha\wedge\beta)=d\alpha\wedge\beta+(-1)^{\deg\alpha}\alpha\wedge d\beta$;
$\partial$ and $\dbar$ satisfy the same sign rule, using total degree.
The wedge product therefore induces multiplication in these cohomology groups: changing a closed representative by an exact form changes its product with another closed representative only by an exact form.
\end{definition}

The quotient in each cohomology group identifies two closed forms when their difference has a global primitive.
In particular, a class is zero precisely when its representative is exact.
The next proposition proves local exactness on a star-shaped set; the example following it shows why global exactness is a separate question.

\begin{proposition}[The real Poincar\'e lemma and its homotopy operator]\label{bas:poincare}
Let $U\subset\R^m$ be star-shaped about the origin. Every smooth closed form of degree $k\ge1$ on $U$ is exact.
More explicitly, set
\[
 (H\alpha)_x(v_1,\ldots,v_{k-1})=
 \int_0^1t^{k-1}\alpha_{tx}(x,v_1,\ldots,v_{k-1})\,dt.
\]
Then $dH+Hd=I$ in positive degrees, while on functions $Hd f=f-f(0)$.
\end{proposition}
\begin{proof}
Write $R=\sum_jx_j\partial_{x_j}$, and let $\iota_R$ denote insertion of $R$ into the first argument of a form.
Its local flow is $\Phi_s(x)=e^s x$; define the Lie derivative by
$\mathcal L_R\alpha=\left.\frac{d}{ds}\right|_{s=0}\Phi_s^*\alpha$.
First verify Cartan's formula $\mathcal L_R=d\iota_R+\iota_Rd$.
On a function $f$, the right-hand side is $\iota_Rdf=R(f)$; on a coordinate one-form $dx_j$, it is $d(x_j)=dx_j$.
Both sides act on products of functions and forms by the same derivation rule, so equality on these generators proves the formula for all forms.

For $\alpha=\sum_{|I|=k}a_I(x)dx_I$, the radial map $F_t(x)=tx$ gives
$F_t^*\alpha=\sum_I t^ka_I(tx)dx_I$.
Differentiating term by term for $t>0$, we obtain
\[
 \frac d{dt}F_t^*\alpha
 =\sum_I\bigl(kt^{k-1}a_I(tx)+t^k\sum_jx_j\partial_ja_I(tx)\bigr)dx_I
 =t^{-1}F_t^*\mathcal L_R\alpha.
\]
Evaluating $t^{-1}F_t^*\iota_R\alpha$ on $v_1,\ldots,v_{k-1}$ gives exactly
$t^{k-1}\alpha_{tx}(x,v_1,\ldots,v_{k-1})$.
Thus the stated operator $H$ is the integral of this expression over $t\in[0,1]$.

Apply Cartan's formula and integrate over $[\varepsilon,1]$. Since pullback commutes with $d$, this gives
$F_1^*\alpha-F_\varepsilon^*\alpha=dH_\varepsilon\alpha+H_\varepsilon d\alpha$.
For any fixed compact set, the radial segments joining its points to the origin lie in a compact subset of $U$.
All coefficients and all their derivatives are bounded there, and $t^{k-1}$ is integrable near zero.
We may therefore differentiate under the integral and let $\varepsilon\downarrow0$.
For $k>0$, the coefficients of $F_\varepsilon^*\alpha$ contain the factor $\varepsilon^k$ and tend to zero; for $k=0$, the limit is the constant function $f(0)$.
This proves the homotopy identity. If $d\alpha=0$, it immediately reduces to $\alpha=dH\alpha$.
\end{proof}

\begin{example}[A radial primitive and a genuine global obstruction]\label{bas:worked-poincare}
On $\R^2$, the form $\alpha=2xy\,dx+x^2\,dy$ has primitive $x^2y$.
On $\R^2\setminus\{0\}$, the form
$\beta=(-y\,dx+x\,dy)/(x^2+y^2)$ is closed but has no global primitive.
\end{example}
\begin{proof}
Direct computation gives $d\alpha=(2x-2x)dx\wedge dy=0$.
For this one-form, the radial operator gives
\[
 H\alpha(x,y)=\int_0^1\alpha_{(tx,ty)}(x,y)\,dt
 =\int_0^1(2t^2x^2y+t^2x^2y)\,dt=x^2y.
\]
Differentiating indeed recovers $\alpha$.
For $\beta=Pdx+Qdy$, we have
$\partial_xQ=(y^2-x^2)/(x^2+y^2)^2$ and
$\partial_yP=(y^2-x^2)/(x^2+y^2)^2$, so $d\beta=0$.
But its pullback to the unit circle $(x,y)=(\cos t,\sin t)$ is $dt$, with integral $2\pi$.
If $\beta=df$, its integral around this closed curve would be the terminal value of $f$ minus its initial value, hence zero, a contradiction.
The punctured plane is not star-shaped about the origin: radial segments to the origin leave the domain.
Local primitives supplied by Proposition~\ref{bas:poincare} therefore need not glue to a global primitive.
\end{proof}

\section{Partitions of Unity, Integration, and Stokes' Theorem}
A family of subsets is locally finite if every point has a neighborhood meeting only finitely many members.
For a locally finite family of supports, sums of smooth functions and all their derivatives are ordinary finite sums on each sufficiently small neighborhood.

\begin{theorem}[Smooth partitions of unity]\label{bas:partition}
Every open cover of a smooth second countable Hausdorff manifold admits a subordinate locally finite smooth partition of unity:
$0\le\rho_i\le1$, $\sum_i\rho_i=1$, and each $\supp\rho_i$ is contained in some member of the original cover.
If $K\subset U$, with $K$ compact and $U$ open, there is a smooth function $\chi$ such that
$0\le\chi\le1$, $\chi=1$ on a neighborhood of $K$, and $\supp\chi\Subset U$.
\end{theorem}
\begin{proof}
Every point has a small coordinate ball whose closure is compactly contained in a chart. A countable basis allows us to choose countably many such balls covering the manifold.
At step $j$, include the closure of the $j$th ball in this enumeration and the closures of finitely many larger coordinate balls covering the compact set constructed at the preceding step. This gives a compact exhaustion
$K_j\subset\operatorname{int}K_{j+1}$, $\bigcup_jK_j=X$.
Cover the compact annulus $K_j\setminus\operatorname{int}K_{j-1}$ by finitely many small coordinate balls $W_{j\ell}$, and choose slightly larger balls $V_{j\ell}$ such that $\overline W_{j\ell}\subset V_{j\ell}$ and
$\overline V_{j\ell}$ is contained both in a member of the original cover and in $\operatorname{int}K_{j+1}\setminus K_{j-2}$. For the initial indices set $K_{-1}=K_0=\varnothing$.
To make these choices, first shrink balls pointwise and then take a finite subcover.
Every compact set $C$ lies in some $\operatorname{int}K_N$. For $j\ge N+2$, the corresponding $V_{j\ell}$ do not meet $C$, so $\{V_{j\ell}\}$ is locally finite.
Within each coordinate ball, the function $e^{-1/t}$ for $t>0$, extended by zero for $t\le0$, supplies a function $b_{j\ell}\in C_c^\infty(V_{j\ell})$ equal to one on $\overline W_{j\ell}$ and taking values in $[0,1]$.
The locally finite sum $b=\sum b_{j\ell}$ is smooth and everywhere positive.
Set $\rho_{j\ell}=b_{j\ell}/b$ to obtain the partition of unity. To index it by members of the original cover, sum the functions assigned to the same member.
Finally, for $K\subset U$, choose finitely many of the small balls so that their larger balls have closures in $U$.
If the associated bump functions are $b_1,\ldots,b_N$, then $\chi=1-\prod_{j=1}^N(1-b_j)$ has the required properties.
\end{proof}

\begin{proposition}[Natural orientation and Stokes' theorem]\label{bas:stokes}
Every complex manifold of dimension $n$ has a natural orientation as a real manifold of dimension $2n$.
For every compactly supported smooth $(2n-1)$-form $\eta$, we have $\int_Xd\eta=0$.
More generally, for a compact smooth domain $M\subset X$ and a smooth $(2n-1)$-form defined near $M$,
$\int_M d\eta=\int_{\partial M}\eta$, where the boundary orientation places an outward normal first.
On a compact manifold, the integral of the wedge product of two closed forms of complementary degrees depends only on their de Rham cohomology classes.
\end{proposition}
\begin{proof}
Let $A=B+\ii C$ be an invertible complex-linear matrix. Its corresponding real matrix is
$\begin{psmallmatrix}B&-C\\ C&B\end{psmallmatrix}$.
After complexification and a change of basis, this matrix is similar to $\operatorname{diag}(A,\overline A)$.
Its real determinant is therefore $\det A\det\overline A=|\det A|^2>0$.
Holomorphic coordinate changes preserve real orientation, so the orientations from complex coordinates agree.

Define integration of top-degree forms using a partition of unity subordinate to coordinate charts, as supplied by Theorem~\ref{bas:partition}.
If $\{\sigma_j\}$ is a second partition of unity, insert $1=\sum_j\sigma_j$ to split the original integral into terms involving $\rho_i\sigma_j$.
On every chart intersection, the real change-of-variables formula identifies the two coordinate integrals.
Summing again gives the integral defined by the second partition, proving independence of charts and partition of unity.

Because $\eta$ has compact support, only finitely many terms of the locally finite partition occur near its support.
Since $\sum_i\rho_i=1$ and $\sum_i d\rho_i=0$,
\[
 \sum_i d(\rho_i\eta)=\sum_i(d\rho_i\wedge\eta+\rho_i d\eta)=d\eta.
\]
In an oriented coordinate chart, write a compactly supported $(2n-1)$-form as
$\sum_j(-1)^{j-1}a_jdx_1\wedge\cdots\widehat{dx_j}\cdots\wedge dx_{2n}$.
Its differential is $(\sum_j\partial_ja_j)dx_1\wedge\cdots\wedge dx_{2n}$.
Integrating first in the $j$th variable, the fundamental theorem of calculus gives
$\int_{-\infty}^{\infty}\partial_ja_j\,dx_j=0$ because $a_j$ vanishes at both ends.
Fubini's theorem then makes the full integral zero.
Summing over charts proves Stokes' theorem.

For a domain with smooth boundary, use the real implicit function theorem to choose boundary coordinates in which the domain is a half-space, and use the same partition argument.
Integrating tangential derivatives still gives zero.
Integrating the derivative in the normal coordinate gives the boundary value by the fundamental theorem of calculus, with the sign determined by the outward-normal-first orientation.
This is precisely the integral of the restricted form on the boundary chart.
Summing these identities proves the formula for $M$.

If $d\beta=0$, then
$d(\gamma\wedge\beta)=d\gamma\wedge\beta$, since the other Leibniz term vanishes.
Thus replacing $\alpha$ by $\alpha+d\gamma$ does not change the integral of its wedge product with $\beta$.
The same argument applies to the other factor, so the integral depends only on the cohomology classes.
\end{proof}

\section{Hermitian Metrics and K\"ahler Forms}
The real tangent bundle of a complex manifold becomes a complex vector bundle through the operator $J$; equivalently one may use the $\ii$-eigenspace $T^{1,0}X$ of $J$ in $TX\otimes_\R\C$.
The holomorphic version of the vector-bundle definition and its transition functions are developed in Chapter~\ref{bun:chapter}.

\begin{definition}
A \emph{Hermitian metric} is a smoothly varying positive definite Hermitian inner product on the complex tangent bundle.
The associated real $(1,1)$-form has local expression
$\omega=\ii\sum g_{j\bar k}dz_j\wedge d\bar z_k$, where $(g_{j\bar k})$ is a positive definite Hermitian matrix.
Here we set $g_{j\bar k}=h(\partial_{z_k},\partial_{z_j})$, consistently with conjugate linearity in the first argument.
If $d\omega=0$, then $\omega$ is a \emph{K\"ahler form}, and the manifold is called \emph{K\"ahler}.
Its volume form is $dV_\omega=\omega^n/n!$, where $\omega^n$ denotes the wedge product of $n$ copies of $\omega$.
\end{definition}

\begin{proposition}[Existence and comparison of Hermitian metrics]\label{bas:metric}
Every complex manifold, and every smooth complex vector bundle over it, admits a Hermitian metric.
Two Hermitian metrics $h_1,h_2$ on a bundle over a compact manifold satisfy, pointwise on every fiber vector,
$c h_1\le h_2\le C h_1$ for constants $0<c\le C<\infty$.
\end{proposition}
\begin{proof}
Choose the standard Hermitian inner product in each local trivialization and sum them using a nonnegative partition of unity from Theorem \ref{bas:partition}.
At every point at least one weight is positive, so the sum is positive definite; local finiteness ensures smoothness.
For the second assertion, express $h_2$ relative to $h_1$ as a positive definite self-adjoint operator.
Its smallest and largest eigenvalues vary continuously: for example, the smallest is the minimum of the Rayleigh quotient on the unit sphere, and this quotient varies uniformly continuously as the matrix varies.
Compactness gives a positive lower bound for the smallest eigenvalue and a finite upper bound for the largest.
\end{proof}

A complex submanifold is a subset locally given in holomorphic coordinates by a coordinate plane, with the induced manifold structure.

\begin{proposition}[Basic operations with K\"ahler forms]\label{bas:kahler-operations}
The restriction of a K\"ahler form to a complex submanifold is K\"ahler. If $\pi_1,\pi_2$ are the two coordinate projections, the product of two K\"ahler manifolds has K\"ahler form
$\pi_1^*\omega_1+\pi_2^*\omega_2$.
On a compact complex manifold of positive dimension, a K\"ahler form cannot be exact as a real differential form.
\end{proposition}
\begin{proof}
Restriction and pullback commute with $d$, so these forms are closed.
Restriction to a submanifold's tangent space preserves positive definiteness. A tangent vector to a product has two components, at least one of which is nonzero when the vector is nonzero, so the sum of the two pullbacks is positive definite.
If $\omega=d\alpha$ on a compact manifold of complex dimension $n$, then $d\omega=0$ gives
$\omega^n=d(\alpha\wedge\omega^{n-1})$. Proposition~\ref{bas:stokes} yields $\int_X\omega^n=0$.
But $\omega^n$ is an everywhere positive volume form, so its integral is strictly positive, a contradiction.
\end{proof}

\begin{example}
The flat form $\frac{\ii}{2}\sum dz_j\wedge d\bar z_j$ is invariant under lattice translations and hence descends to a K\"ahler form on a complex torus.
The associated form of a general Hermitian metric need not be closed. A partition of unity preserves positivity, but differentiating its variable coefficients produces extra terms; it therefore does not automatically produce a K\"ahler metric.
\end{example}

\section{Local and Global Exactness}
Given a differential form, check its degree, its type, whether it is closed, and whether it is exact, in that order.
For example, $d\bar z$ is $\dbar$-exact on $\C$ because $d\bar z=\dbar\bar z$. On a complex torus, however, $\bar z$ does not descend to a single-valued function, so the same local equality need not supply a global primitive.
The distinction between local solvability and global obstructions will be expressed precisely by sheaf cohomology in the next part.


\chapter{Holomorphic Vector Bundles, Curvature, and Chern Classes}\label{bun:chapter}
\begin{learning}
Vector bundles allow vector spaces and linear maps to vary over a manifold. A connection differentiates their sections, and its curvature measures the failure of these derivatives to commute. Traces and determinants of curvature then define Chern classes, which connect local differential calculations with global topology. Projective line bundles and symmetric powers give concrete examples of these constructions.
\end{learning}

\section{Transition Functions and Sections}
\begin{definition}[Holomorphic vector bundles]
A \emph{holomorphic vector bundle} $\pi:E\to X$ of rank $r=\operatorname{rk}E$ has trivializations
$\pi^{-1}(U_i)\simeq U_i\times\C^r$ over an open cover $\{U_i\}$, with fiber coordinates related on overlaps by
$v_i=g_{ij}(x)v_j$. Here $GL_r(\C)$ is the group of invertible complex $r\times r$ matrices, and $g_{ij}:U_i\cap U_j\to GL_r(\C)$ is holomorphic and satisfies
$g_{ii}=I$ and $g_{ij}g_{jk}=g_{ik}$.
A \emph{section} $s$ satisfies $\pi\circ s=\mathrm{id}$ and is represented in trivializations by vectors of functions satisfying $s_i=g_{ij}s_j$.
A section is smooth or holomorphic when all its coefficient functions have that property.
Write $H^0(U,E)$ for the vector space of holomorphic sections over $U$; its identification with degree-zero sheaf cohomology will be explained in Definition~\ref{sh:cohomology-definition}.
The fiber $E_x=\pi^{-1}(x)$ is an $r$-dimensional complex vector space, whereas the germ of a local section remembers its values on a neighborhood, not just its value at $x$.
A rank-one bundle is a \emph{line bundle}. Write $E^*$, $E\oplus F$, $E\otimes F$, and $\det E=\Lambda^rE$ for the dual, direct sum, tensor product, and determinant bundle, respectively.
Here $\Lambda^jE$ has fiber the $j$th exterior power of $E_x$; all these operations are performed on fibers and glued by the induced transition maps.
For a line bundle, $L^k=L^{\otimes k}$, with $L^0$ the trivial line bundle and negative powers interpreted by duality.
A local frame $e=(e_1,\ldots,e_r)$ is an ordered basis of holomorphic sections on each fiber; section coordinates are column vectors, and $G^t,G^*$ mean transpose and conjugate transpose.
\end{definition}

\begin{proposition}[Constructing a bundle from a cocycle]\label{bun:gluing}
A holomorphic matrix cocycle as above determines a holomorphic vector bundle uniquely up to the natural isomorphism preserving its trivializations.
The transition functions for the dual, direct sum, tensor product, and determinant are, respectively,
$(g_{ij}^{-1})^t$, $g_{ij}\oplus h_{ij}$, $g_{ij}\otimes h_{ij}$, and $\det g_{ij}$.
\end{proposition}
\begin{proof}
Refine the trivializing cover to coordinate neighborhoods of the base. On the disjoint union $\bigsqcup_iU_i\times\C^r$, impose the relation
$(x,v)_j\sim(x,g_{ij}(x)v)_i$.
The identity $g_{ii}=I$ gives reflexivity, $g_{ij}g_{ji}=I$ gives symmetry, and
$g_{ij}g_{jk}=g_{ik}$ gives transitivity.
If two points in the same chart are equivalent, their base points agree and their fiber coordinates are unchanged by $g_{ii}=I$.
Thus each chart injects into the quotient, with image exactly $\pi^{-1}U_i$.
Equip these images with product coordinates. The overlap map
$(x,v)\mapsto(x,g_{ij}(x)v)$ is holomorphic, with inverse given by $g_{ji}$.
Proposition \ref{bas:atlas} therefore supplies the unique locally Euclidean topology in which these charts are open.

We must still verify that the total space is Hausdorff.
Points over distinct base points are separated by inverse images of disjoint open neighborhoods in the base.
For two points over the same base point, choose a trivialization about that point. Their fiber coordinates differ, so they have disjoint product neighborhoods in this open chart, which are also open in the total space.
Second countability of the base implies that every open cover has a countable subcover.
Retain countably many trivializations; the union of countable bases for their product charts is a countable basis for the total space.
Addition, scalar multiplication, and zero vectors agree on overlaps because $g_{ij}$ is linear, so the result is indeed a vector bundle.
Any other bundle with the same trivialization data is identified with it by the identity map in each chart. These maps agree on overlaps and glue uniquely to a holomorphic bundle isomorphism.

The dual transition rule follows from invariance of the pairing.
If $v_i=g_{ij}v_j$, the requirement $\lambda_i^tv_i=\lambda_j^tv_j$ for every $v_j$ gives
$\lambda_i=(g_{ij}^{-1})^t\lambda_j$.
The two summands of a direct sum transform separately, giving a block diagonal matrix.
A tensor basis transforms factor by factor, giving the Kronecker tensor matrix; the coefficient of a top exterior product is the determinant.
Each operation respects matrix multiplication, so the resulting transition functions still satisfy the cocycle condition.
\end{proof}

\begin{example}[The canonical and tangent bundles]
The Jacobian matrices of complex coordinate changes define the holomorphic tangent bundle $T^{1,0}X$; its dual is $\Omega_X^1$.
The bundle $K_X=\Lambda^n\Omega_X^1$ is the \emph{canonical bundle}, locally framed by $dz_1\wedge\cdots\wedge dz_n$.
If $w=w(z)$, then
$dw_1\wedge\cdots\wedge dw_n=\det(\partial w/\partial z)\,dz_1\wedge\cdots\wedge dz_n$.
This formula gives all transition data for $K_X$.
\end{example}

\section{Connections and the Hermitian Metric}
Differentiating the coefficient functions of a section depends on the chosen frame: a change of frame also differentiates its transition matrix.
A connection adds a matrix-valued one-form to the ordinary derivative so that the result transforms as a bundle-valued form.
The Hermitian metric will determine this extra term uniquely once its $(0,1)$ part has been fixed.
Write $A^k(X,E)$, or $A^k(E)$ when $X$ is fixed, for smooth sections of
$\Lambda^k(T^*X\otimes_\R\C)\otimes E$, and use $A^{p,q}(X,E)$ for the type-$(p,q)$ part.
The bundle $\operatorname{End}(E)=E^*\otimes E$ consists of fiberwise complex-linear endomorphisms.
In a holomorphic frame define $\dbar_E$ coefficientwise; holomorphic transition matrices show that these local operators agree, and Proposition~\ref{bas:type-decomposition} gives $\dbar_E^2=0$.
We often omit the subscript $E$ when the coefficient bundle is fixed.
For matrix-valued forms, wedge product includes matrix composition:
$(A\wedge B)_{ij}=\sum_k A_{ik}\wedge B_{kj}$.
Conjugate transpose of such a form also conjugates its scalar form coefficients; it is not a Hodge-star operator.

\begin{definition}[Connections and curvature]
A \emph{connection} on a smooth bundle is a complex-linear operator
$D:A^0(E)\to A^1(E)$ satisfying $D(fs)=df\otimes s+fDs$.
Extend it to bundle-valued forms by
$D(\alpha\otimes s)=d\alpha\otimes s+(-1)^{\deg\alpha}\alpha\wedge Ds$.
Its \emph{curvature} $F_D=D^2$ is an $\operatorname{End}(E)$-valued two-form.
A connection is compatible with a Hermitian metric if
$d\langle s,t\rangle=\langle Ds,t\rangle+\langle s,Dt\rangle$, with the pairing conjugate-linear in the first argument.
\end{definition}

\begin{proposition}[Local formulas and transformation rules]\label{bun:connection-local}
In a local frame, $D=d+A$, where $A$ is a matrix-valued one-form, and
\[
 F_D=dA+A\wedge A,\qquad DF_D=dF_D+A\wedge F_D-F_D\wedge A=0.
\]
Under the frame change $e'=eG$,
\[
 A'=G^{-1}AG+G^{-1}dG,\qquad F_D'=G^{-1}F_DG.
\]
\end{proposition}
\begin{proof}
Write the frame as a row vector $e=(e_1,\ldots,e_r)$ and a section as $s=ev$, where $v$ is a column vector.
Set $De=eA$. The Leibniz rule gives $Ds=e(dv+Av)$.
On applying $D$ again, the Leibniz sign for a one-form is negative, so
\[
 (d+A)^2v=d(Av)+A\wedge dv+A\wedge Av
 =dA\,v-A\wedge dv+A\wedge dv+A\wedge A\,v.
\]
The terms containing $dv$ cancel, leaving $F=dA+A\wedge A$.
This operator is $C^\infty$-linear, so curvature is indeed an endomorphism-valued two-form.

For an endomorphism-valued two-form, covariant differentiation is
$DF=dF+A\wedge F-F\wedge A$.
Substituting the curvature formula gives
\[
 dF=dA\wedge A-A\wedge dA,
 \qquad A\wedge F-F\wedge A=A\wedge dA-dA\wedge A.
\]
The cubic terms $A\wedge A\wedge A$ cancel, and hence $DF=0$.

If $e'=eG$, then
$D(e'v')=e\{Gdv'+(dG+AG)v'\}$.
Rewrite this as $e'\{dv'+G^{-1}(dG+AG)v'\}$ to obtain the formula for $A'$.
Likewise, $D^2(e'v')=eFGv'=e'(G^{-1}FG)v'$ gives the conjugation rule for curvature.
The connection matrix thus has an inhomogeneous term $G^{-1}dG$, whereas curvature transforms tensorially.
\end{proof}

\begin{theorem}[The Chern connection]\label{bun:chern-connection}
A Hermitian holomorphic vector bundle $(E,h)$ has a unique metric-compatible connection satisfying $D^{0,1}=\dbar$, where $D^{0,1}$ is the type-$(0,1)$ component of $D$.
In a holomorphic frame with $\langle v,w\rangle=v^*Hw$,
\begin{equation}\label{bun:chern-formula}
 A=H^{-1}\partial H,\qquad F_D=\dbar(H^{-1}\partial H).
\end{equation}
The curvature has type $(1,1)$. For a line bundle, writing $h(e,e)=e^{-\varphi}$ gives
\begin{equation}\label{bun:line-curvature}
 A=-\partial\varphi,\qquad F_D=\partial\dbar\varphi.
\end{equation}
\end{theorem}
\begin{proof}
In a holomorphic frame, $\dbar(ev)=e\dbar v$, so $D^{0,1}=\dbar$ means that the connection matrix $A$ has only a $(1,0)$ part.
For arbitrary coefficient columns $v,w$, the two sides of metric compatibility are
\begin{align*}
 d(v^*Hw)&=dv^*Hw+v^*dH\,w+v^*Hdw,\\
 \langle Dv,w\rangle+\langle v,Dw\rangle
 &=dv^*Hw+v^*(A^*H+HA)w+v^*Hdw.
\end{align*}
Since $v,w$ are arbitrary, compatibility is exactly $dH=A^*H+HA$.
The form $A^*$ has type $(0,1)$, so the $(1,0)$ part gives $\partial H=HA$.
The only possible connection matrix is therefore $A=H^{-1}\partial H$.
Conversely, taking conjugate transposes and using $H^*=H$ gives $\dbar H=A^*H$, so this $A$ satisfies the entire compatibility equation.

To check gluing, write $H'=G^*HG$. Since $G$ is holomorphic, $\partial G^*=0$, and hence
\[
 (H')^{-1}\partial H'
 =G^{-1}H^{-1}(G^*)^{-1}G^*((\partial H)G+H\partial G)
 =G^{-1}AG+G^{-1}\partial G.
\]
Because $dG=\partial G$, this is precisely the transformation rule of Proposition \ref{bun:connection-local}.
The unique local constructions therefore glue to a global connection.

Finally, differentiating $H^{-1}H=I$ by $\partial$ gives
$\partial H^{-1}=-H^{-1}(\partial H)H^{-1}$.
Since $\partial^2H=0$,
\[
 \partial A
 =-H^{-1}(\partial H)H^{-1}\wedge\partial H
 =-A\wedge A.
\]
The $(2,0)$ part of curvature vanishes; there is no $(0,2)$ part because $A^{0,1}=0$.
Thus $F=\dbar A$ has type $(1,1)$.
In rank one, $H=e^{-\varphi}$ and scalar multiplication commutes, giving $A=-\partial\varphi$ and
$F=-\dbar\partial\varphi=\partial\dbar\varphi$. The last equality uses the anticommutation in Proposition~\ref{bas:type-decomposition}.
\end{proof}

\begin{proposition}[The first Chern form and independence of the metric]\label{bun:first-chern}
The first Chern form of a line bundle $L$ is
$c_1(L,h)=\frac{\ii}{2\pi}F_D$. It is a real closed $(1,1)$-form.
If $h'=he^{-\psi}$ for a global smooth real function $\psi$, then
\[
 c_1(L,h')-c_1(L,h)=\frac{\ii}{2\pi}\partial\dbar\psi.
\]
Thus it defines a real de Rham class $c_1(L)$ independent of the metric.
Moreover, $c_1(L\otimes M)=c_1(L)+c_1(M)$ and $c_1(L^*)=-c_1(L)$.
\end{proposition}
\begin{proof}
For real $\varphi$,
$\overline{\partial\dbar\varphi}=\dbar\partial\varphi=-\partial\dbar\varphi$.
Conjugation also sends $\ii$ to $-\ii$, so
$\overline{\ii\partial\dbar\varphi}=\ii\partial\dbar\varphi$.
The Chern form is therefore real. It is closed because
$\partial^2=\dbar^2=0$ and $\partial\dbar=-\dbar\partial$.
In rank one, conjugation in the curvature transformation rule has no effect, so the local forms glue to a global form.

Under the frame change $e'=ge$, both metric coefficients are multiplied by $|g|^2$.
Their ratio $h'/h$ is therefore a global positive function. Set $\psi=-\log(h'/h)$.
The local weight changes from $\varphi$ to $\varphi+\psi$, so \eqref{bun:line-curvature} gives the difference formula.
To see that the difference is exact in the \emph{real} de Rham complex, take the global real one-form
\[
 \eta=\frac{\ii}{4\pi}(\dbar\psi-\partial\psi).
\]
Direct conjugation gives $\bar\eta=\eta$, and
$d\eta=\frac{\ii}{2\pi}\partial\dbar\psi$.
The two real closed forms thus represent the same real cohomology class.

The squared norm of a tensor product frame is the product of the two squared norms, so its weights add.
The squared norm of a dual frame is the reciprocal of the original squared norm, so its weight changes sign.
Substitution into the curvature formula gives the tensor-product and duality formulas for the first Chern class.
\end{proof}

\section{Explicit Models on Projective Space}
\begin{definition}
The fiber of $\mathcal O_{\mathbb P^n}(-1)$ over $[Z]$ is the line $\C Z\subset\C^{n+1}$. This is the \emph{tautological line bundle}.
Set $\mathcal O(1)=\mathcal O(-1)^*$ and $\mathcal O(k)=\mathcal O(1)^{\otimes k}$, interpreting negative powers by duality.
\end{definition}

For a real $C^2$ function $\varphi$, its \emph{Levi form} on a $(1,0)$ vector $v$ is
$\sum_{j,k}(\partial_{z_j}\partial_{\bar z_k}\varphi)v_j\bar v_k$.
A real $(1,1)$-form $\ii\sum a_{j\bar k}dz_j\wedge d\bar z_k$ is positive, respectively semipositive, when this coefficient quadratic form is positive on every nonzero vector, respectively nonnegative on every vector.
For a line metric, curvature positivity refers to the real form $\ii F_D$ in \eqref{bun:line-curvature}.

\begin{proposition}[The Fubini--Study form]\label{bun:fs}
The metric on $\mathcal O(1)$ induced by the standard inner product has weight
$\log(1+|z|^2)$ in the chart $Z_0\ne0$. Consequently,
\[
 \omega_{\mathrm{FS}}=\frac{\ii}{2\pi}\partial\dbar\log(1+|z|^2)
\]
is a globally defined K\"ahler form on $\mathbb P^n$ representing $c_1(\mathcal O(1))$.
Moreover, $\int_{\mathbb P^1}\omega_{\mathrm{FS}}=1$.
\end{proposition}
\begin{proof}
A local frame of the tautological bundle is $(1,z_1,\ldots,z_n)$, with squared norm $1+|z|^2$; its dual frame has the reciprocal squared norm.
By Proposition \ref{bun:first-chern}, the local Chern forms automatically glue and are closed.
For $v\in\C^n$, the Levi form is
\[
 \sum_{j,k}\partial_{z_j}\partial_{\bar z_k}\log(1+|z|^2)v_j\bar v_k
 =\frac{(1+|z|^2)|v|^2-|\langle z,v\rangle|^2}{(1+|z|^2)^2}
 \ge\frac{|v|^2}{(1+|z|^2)^2}>0
\]
when $v\ne0$, proving positive definiteness.
In the one-dimensional chart, $\omega_{\mathrm{FS}}=\frac1\pi\frac{dx\wedge dy}{(1+|z|^2)^2}$.
The omitted point at infinity has measure zero. Polar coordinates give
\[
 \int_{\mathbb P^1}\omega_{\mathrm{FS}}
 =2\int_0^\infty\frac{r}{(1+r^2)^2}\,dr=1.
\]
\end{proof}

\begin{proposition}[Sections on the projective line]\label{bun:sections-p1}
For $k\ge0$, $\dim H^0(\mathbb P^1,\mathcal O(k))=k+1$; for $k<0$, this space is zero.
In addition, $K_{\mathbb P^1}\simeq\mathcal O(-2)$.
\end{proposition}
\begin{proof}
Use the tautological frame $(1,z)$ on $U_0$ and $(w,1)$ on $U_1$.
Since $w=1/z$, the second frame is $z^{-1}$ times the first.
After dualizing, the frames of $\cO(1)$ satisfy $e_1=z e_0$.
A global section $s=f_0(z)e_0^{\otimes k}=f_1(w)e_1^{\otimes k}$ must therefore satisfy
$f_1(w)=w^kf_0(1/w)$.

The function $f_0$ is defined on the entire affine chart $\C$, so it is entire. Write
$f_0(z)=\sum_{m\ge0}a_mz^m$.
For $w\ne0$, substitution gives the convergent Laurent expansion
$w^kf_0(1/w)=\sum_{m\ge0}a_mw^{k-m}$.
Holomorphicity of $f_1$ at $w=0$, the point at infinity, is equivalent to absence of negative powers.
Laurent coefficients are uniquely determined by contour integrals, so $a_m=0$ whenever $m>k$.
If $k<0$, every $m\ge0$ satisfies $m>k$, leaving only the zero section.
If $k\ge0$, exactly $a_0,\ldots,a_k$ may be chosen freely, and every resulting polynomial is holomorphic in the other chart. The dimension is therefore $k+1$.
The corresponding homogeneous sections are $Z_0^k,Z_0^{k-1}Z_1,\ldots,Z_1^k$.

The two canonical-bundle frames are $dz,dw$, with $dw=-z^{-2}dz$.
Replace the second frame by $-dw$; the transition factor becomes $z^{-2}$, exactly the frame transition of $\cO(-2)$.
This gives the desired line-bundle isomorphism.
\end{proof}

A finite-dimensional \emph{linear system} in $L$ is a vector subspace of $H^0(X,L)$.
It is \emph{base-point-free}, or generates $L$, if its sections have no common zero.
For a basis with local coefficients $f_0,\ldots,f_N$, the associated holomorphic map is $x\mapsto[f_0(x):\cdots:f_N(x)]$; a frame change multiplies every coefficient by the same nonzero function, so this map is well defined.
An immersion has injective differential at every point; an embedding is an immersion that is also a homeomorphism onto its image.

\begin{proposition}[A base-point-free linear system gives a semipositive metric]\label{bun:generated-semipositive}
Let $m\ge1$ be an integer. If $L^m$ has finitely many holomorphic sections $s_0,\ldots,s_N$ with no common zero, then $L$ admits a smooth metric with semipositive curvature.
If the associated map $\Phi:X\to\mathbb P^N$ is an immersion, this metric is strictly positive.
\end{proposition}
\begin{proof}
In a local frame $e$ of $L$, write $s_j=f_je^{\otimes m}$.
Since there is no common zero, $S=\sum_j|f_j|^2$ is smooth and strictly positive.
Define $h(e,e)=S^{-1/m}$.
Under $e'=ge$, the coefficients satisfy $f'_j=g^{-m}f_j$, so
$(\sum|f'_j|^2)^{-1/m}=|g|^2S^{-1/m}$.
This is exactly the transformation rule for the same Hermitian metric in the new frame, proving that the metric is globally defined.
Its weight is $m^{-1}\log S$.

To compute curvature at a point, extend a $(1,0)$ tangent vector $v$ with constant coordinate coefficients and set $a=df(v)$. Write $\partial_v=\sum_jv_j\partial_{z_j}$ and $\partial_{\bar v}=\sum_j\bar v_j\partial_{\bar z_j}$ in this calculation.
Holomorphicity gives
$\partial_vS=\sum_j a_j\bar f_j$ and
$\partial_v\partial_{\bar v}S=\sum_j|a_j|^2$.
Differentiating the logarithm twice yields
\[
 \partial_v\partial_{\bar v}\log S
 =\frac{S|a|^2-|\langle f,a\rangle|^2}{S^2}\ge0,
\]
where the inequality is Cauchy--Schwarz.
Equality holds exactly when $a$ is a complex multiple of $f$.

The map $\Phi=[f_0:\cdots:f_N]$ is independent of the frame.
In a chart where $f_i\ne0$, its differential has components
$d(f_j/f_i)(v)=(a_jf_i-f_ja_i)/f_i^2$.
All these components vanish exactly when $a=(a_i/f_i)f$.
Thus the Levi form above is strictly positive on a nonzero $v$ precisely when $d\Phi(v)\ne0$.
The extra positive factor $1/m$ in curvature does not affect positivity, so an immersion gives a strictly positive metric.
\end{proof}

\begin{definition}[Initial notions of positivity]
A line bundle $L$ is \emph{semiample} if some positive tensor power is globally generated by finitely many holomorphic sections.
It is \emph{very ample} if a finite-dimensional linear system of sections of $L$ defines a holomorphic embedding of $X$ into projective space, and \emph{ample} if some positive tensor power is very ample.
It is \emph{positive} or \emph{semipositive} if it admits a smooth Hermitian metric with strictly positive or semipositive curvature, respectively.
Proposition \ref{bun:generated-semipositive} proves that semiample implies semipositive, and ample implies positive.
For compact complex manifolds, the converse assertion that positive implies ample will be proved fully in Chapter \ref{emb:chapter}; it is neither a definition nor a consequence of compactness alone.
\end{definition}

\section{Chern--Weil Theory}
Curvature transforms by conjugation, so its trace and characteristic polynomial are independent of the frame.
The Bianchi identity makes the resulting forms closed, and varying the connection changes them by exact forms.
These two facts are what turn a local matrix calculation into a cohomology class.
\begin{lemma}[The graded trace and consequences of the Bianchi identity]\label{bun:trace-closed}
For matrix-valued forms $A,B$ of degrees $a,b$,
$\tr(A\wedge B)=(-1)^{ab}\tr(B\wedge A)$.
For the curvature $F$ of any connection and every integer $q\ge0$, the form $\tr(F^q)$ is closed. Here $F^q$ means a $q$-fold wedge product with matrix composition, $F^0=I$, and $\tr$ is the sum of diagonal entries.
\end{lemma}
\begin{proof}
Combining matrix multiplication and wedge product gives
$\tr(A\wedge B)=\sum_{i,j}A_{ij}\wedge B_{ji}$.
Interchanging the scalar forms introduces $(-1)^{ab}$; exchanging the indices $i,j$ then gives the graded cyclicity formula.

For an endomorphism-valued $k$-form $C$, the induced covariant derivative is
$DC=dC+A\wedge C-(-1)^kC\wedge A$.
After taking the trace, the last two terms cancel by graded cyclicity, so $d\tr C=\tr(DC)$.
Apply the Leibniz rule to a product of curvatures. Since $\deg F=2$ is even,
\[
 D(F^q)=\sum_{j=0}^{q-1}F^j\wedge(DF)\wedge F^{q-1-j}.
\]
Every term vanishes by the Bianchi identity in Proposition~\ref{bun:connection-local}, giving $d\tr(F^q)=0$. For $q=0$ the trace is the constant rank.
Closedness thus follows from the differential identity for curvature, without requiring diagonalizability of its matrix coefficients.
\end{proof}

\begin{theorem}[Chern forms and independence of the connection]\label{bun:chern-weil}
Choose a connection compatible with a Hermitian metric, put $\Omega=\frac{\ii}{2\pi}F$, and define
\[
 \det(I+t\Omega)=\sum_{j=0}^r c_j(D)t^j,\qquad
 \operatorname{ch}(D)=\tr(e^\Omega)=\sum_{q\ge0}\frac1{q!}\tr(\Omega^q).
\]
Here $t$ is a formal variable, $c_j(D)$ has form degree $2j$, and
$\operatorname{ch}_q(D)=\tr(\Omega^q)/q!$ has form degree $2q$.
These are real closed forms. Their de Rham classes are independent of the chosen metric-compatible connection and are denoted by $c_j(E)$ and $\operatorname{ch}_q(E)$.
Write $c(E)=\sum_jc_j(E)$ for the total Chern class and $\operatorname{ch}(E)=\sum_q\operatorname{ch}_q(E)$ for the Chern character.
Products of classes are induced by wedge product as in Definition~\ref{bas:cohomology-def}; powers denote these products.
We have $c_0=1$ and $c_j=0$ for $j>r$; the top class $c_r$ need not vanish.
\end{theorem}
\begin{proof}
First check that the definitions make sense. Scalar forms of even degree commute, so the matrix entries lie in a commutative ring, and the determinant is defined by the usual sum over permutations.
Under a frame change, $\Omega'=G^{-1}\Omega G$. Determinants and traces are conjugation invariant, so the expressions agree on overlaps and define global differential forms.
Every smooth Hermitian bundle admits a compatible connection: take the coefficientwise differential in local smooth unitary frames, then form an affine combination using a nonnegative partition of unity from Theorem~\ref{bas:partition}.
Since the weights sum to one, the coefficient of $df\otimes s$ in the Leibniz rule remains one, and the metric compatibility identities also add to give compatibility.

To prove closedness, use the formal identity
\begin{equation}\label{bun:det-trace}
 \det(I+t\Omega)=\exp\left(\sum_{q\ge1}
 \frac{(-1)^{q-1}}q\,t^q\tr(\Omega^q)\right).
\end{equation}
For a diagonal matrix, the right-hand side is
$\prod_j\exp(\log(1+t\lambda_j))=\prod_j(1+t\lambda_j)$, which equals the left-hand side.
For each fixed power of $t$, the two coefficients are conjugation-invariant polynomials in the entries of an ordinary complex matrix.
Matrices with distinct eigenvalues are dense and diagonalizable, so these polynomials agree on all complex matrices.
This is an identity of polynomial coefficients and therefore remains valid after substitution in any commutative algebra over $\C$, in particular the algebra of even-degree differential forms.
By Lemma \ref{bun:trace-closed}, all power traces are closed.
Finite sums and products of closed even forms are closed, so every $c_j$ and $\operatorname{ch}_q$ is closed.
Although the formula is written as an infinite series, only finitely many terms contribute in each form degree; no analytic convergence is involved.

Next check reality. In a local unitary frame, the connection matrix evaluated on real tangent vectors is skew-Hermitian, and its curvature satisfies $F^*=-F$.
Hence $\Omega^*=\Omega$, where the star takes the conjugate transpose of the matrix and conjugates the form coefficients.
Since curvature has even degree, transposition introduces no wedge-product sign.
Conjugation therefore leaves both the characteristic-determinant coefficients and the power traces unchanged, proving that the forms are real.

Finally, compare two connections. Their difference $B=D_1-D_0$ is a global endomorphism-valued one-form, because the $df$ terms in the two Leibniz rules cancel.
Set $D_t=D_0+tB$ and write $\dot F_t=dF_t/dt$. A local calculation gives
\[
 F_t=F_0+tD_0B+t^2B\wedge B,
 \qquad\dot F_t=D_tB.
\]
When differentiating a power trace, the matrix factors need not commute, but graded cyclicity moves all terms to the same position.
Since $F_t$ has even degree, these moves introduce no sign. Thus
\[
 \frac d{dt}\tr(F_t^q)
 =\sum_{j=0}^{q-1}\tr(F_t^j\dot F_tF_t^{q-1-j})
 =q\tr((D_tB)F_t^{q-1}).
\]
Proposition~\ref{bun:connection-local} gives $D_t(F_t^{q-1})=0$, so the right-hand side is
$q\,d\tr(BF_t^{q-1})$.
Integration produces the explicit primitive
\[
 \tr(F_1^q)-\tr(F_0^q)
 =d\left(q\int_0^1\tr(BF_t^{q-1})\,dt\right).
\]
Thus all power traces have connection-independent cohomology classes.
If $\alpha'-\alpha=d\eta$ and $\beta$ is closed, then
$\alpha'\wedge\beta-\alpha\wedge\beta=d(\eta\wedge\beta)$.
Apply this successively to each factor of a polynomial. Identity \eqref{bun:det-trace} then proves independence of all Chern classes as well.
The intermediate connections need not be compatible with a single metric; the algebraic identities hold for arbitrary complex connections.
The endpoint forms are real, so taking the real part of any complex primitive proves exactness in the real complex.
Finally, the determinant of an $r\times r$ matrix has degree at most $r$ in $t$ and constant term one.
Hence $c_0=1$ and $c_j=0$ for $j>r$, with no requirement that $c_r$ vanish.
\end{proof}

\begin{remark}
The Chern forms of a general complex connection need not be pointwise real, but their classes agree with the real classes obtained from a unitary connection.
Lifting these classes to integral cohomology requires an additional topological construction.
Proposition \ref{idx:integral-chern} in Chapter \ref{idx:chapter} makes this comparison using projective bundles and the splitting principle, and explicitly distinguishes real classes from possible integral torsion classes.
At this stage the formal determinant definition alone does not imply that all relevant characteristic numbers are integers; see \cite{BottTu1982} for the topological background.
\end{remark}

\begin{proposition}[Direct sums, duals, and tensor products]\label{bun:chern-functors}
Let $\operatorname{rk}E=r$ and $\operatorname{rk}F=s$. Then
\begin{align*}
 c(E\oplus F)&=c(E)c(F),& c_j(E^*)&=(-1)^jc_j(E),\\
 \operatorname{ch}(E\oplus F)&=\operatorname{ch}(E)+\operatorname{ch}(F),&
 \operatorname{ch}(E\otimes F)&=\operatorname{ch}(E)\operatorname{ch}(F),\\
 c_1(E\otimes F)&=s\,c_1(E)+r\,c_1(F).&&
\end{align*}
In particular, $\operatorname{ch}(L)=e^{c_1(L)}$ for a line bundle.
\end{proposition}
\begin{proof}
Choose connections on the two bundles. The direct-sum connection has curvature
$\begin{psmallmatrix}F_E&0\\0&F_F\end{psmallmatrix}$
in a direct-sum frame. Its characteristic determinant is the product of the two determinants, and the trace of its exponential is the sum of the two traces, proving the direct-sum formulas.
These are identities of representative forms and hence also of cohomology classes.

The dual connection is defined by the Leibniz rule for the pairing:
$d(\lambda^ts)=(D_{E^*}\lambda)^ts+\lambda^tD_Es$.
Its connection matrix is $-A_E^t$.
Differentiating once more and accounting for the sign from interchanging one-forms gives curvature $-F_E^t$.
Curvature has even degree, so transposition does not change its determinant. Hence
$\det(I+t\frac\ii{2\pi}F_{E^*})=\det(I-t\frac\ii{2\pi}F_E)$.
Comparing coefficients of $t^j$ gives $c_j(E^*)=(-1)^jc_j(E)$.

The tensor-product connection is $D(s\otimes t)=Ds\otimes t+s\otimes Dt$.
On squaring, the mixed first-order terms cancel by the graded signs, leaving
$F_{E\otimes F}=F_E\otimes I+I\otimes F_F$.
These two even-degree matrices commute, so the binomial expansion of the exponential gives
\[
 e^{\Omega_E\otimes I+I\otimes\Omega_F}
   =e^{\Omega_E}\otimes e^{\Omega_F}.
\]
Summing diagonal entries in a tensor-product basis gives $\tr(P\otimes Q)=\tr P\,\tr Q$, proving multiplicativity of the Chern character.
Compare the degree-two form components: the right-hand side is
$s\operatorname{ch}_1(E)+r\operatorname{ch}_1(F)$.
Since $\operatorname{ch}_1=c_1$, the formula for the first Chern class follows.
For a line bundle, curvature is a scalar two-form, so the trace adds no summation and $\operatorname{ch}(L)=e^{c_1(L)}$.
\end{proof}

\begin{corollary}[Vanishing in the presence of independent sections]\label{bun:independent-sections}
If a smooth complex vector bundle of rank $r$ admits $k$ everywhere linearly independent smooth sections, then $c_j(E)=0$ for $j>r-k$.
\end{corollary}
\begin{proof}
Let the independent sections be $s_1,\ldots,s_k$.
The map $X\times\C^k\to E$, $(x,a)\mapsto\sum a_js_j(x)$, is injective on each fiber and has constant rank $k$.
Near a point where a $k\times k$ minor is nonzero, its coefficients can be solved for smoothly.
Its image is therefore a smooth subbundle $S$, and the sections give a global frame, making $S$ trivial.

Choose a Hermitian metric by Proposition~\ref{bas:metric}. In a local frame, assemble the sections as the columns of a matrix $T$.
The orthogonal projection onto $S$ is
$P=T(T^*HT)^{-1}T^*H$.
The Gram matrix in the middle is positive definite and has smooth inverse, so $P$ is smooth.
Its kernel $F$ has constant rank $r-k$, and local matrix elimination supplies smooth frames. Thus $F$ is a smooth subbundle and $E=S\oplus F$.
Use the flat connection in the global frame of the trivial bundle $S$. Its curvature is zero, so $c(S)=1$.
Proposition~\ref{bun:chern-functors} gives $c(E)=c(F)$. The characteristic determinant of a rank-$r-k$ bundle has no terms of higher degree, so $c_j(E)=0$ for $j>r-k$.
\end{proof}

\section{Chern Classes of Symmetric Powers}
For an integer $k\ge0$, $\operatorname{Sym}^kE$ is the bundle of symmetric degree-$k$ tensors: its fiber is the quotient of $E_x^{\otimes k}$ by the relations permuting factors.
Over $\C$, averaging over the $k!$ permutations identifies this quotient with the invariant subspace of $E_x^{\otimes k}$.
Thus a local basis consists of monomials $e_1^{\alpha_1}\cdots e_r^{\alpha_r}$ with $|\alpha|=k$; for $k=0$ the bundle is the trivial line bundle.

\begin{lemma}[Three counts for weak compositions]\label{bun:composition-counts}
Let $k\ge0$ and $r\ge1$, with the convention $\binom ab=0$ for nonnegative integers $a<b$.
Set
\[
 N=\binom{k+r-1}{r-1},\qquad A=\binom{k+r-1}{r},\qquad B=\binom{k+r-1}{r+1}.
\]
Summing over all $\alpha\in\N_0^r$ with $|\alpha|=k$, we have
\[
 \sum1=N,\quad \sum\alpha_i=A,\quad
 \sum\alpha_i^2=A+2B,\quad \sum\alpha_i\alpha_j=B\quad(i\ne j).
\]
Moreover, $A=kN/r$.
\end{lemma}
\begin{proof}
Place $k$ identical symbols in a row and use $r-1$ separators to divide them into $r$ groups, allowing empty groups.
The group sizes are $\alpha_1,\ldots,\alpha_r$, so the number of weak compositions is $\binom{k+r-1}{r-1}=N$.
Equivalently, if $[t^k]$ denotes extraction of the coefficient of $t^k$, then
$N=[t^k](1-t)^{-r}$.
All power series below are formal, so no numerical convergence in $t$ is needed.

Differentiating the geometric series gives
$\sum_{a\ge0}a t^a=t(1-t)^{-2}$.
Consequently,
\[
 \sum_{k\ge0}\left(\sum_{|\alpha|=k}\alpha_1\right)t^k
 =t(1-t)^{-2}(1-t)^{-(r-1)}=t(1-t)^{-r-1}.
\]
Its coefficient of $t^k$ is $[t^{k-1}](1-t)^{-r-1}=\binom{k+r-1}{r}=A$.
A second differentiation gives
$\sum_{a\ge0}a(a-1)t^a=2t^2(1-t)^{-3}$.
The corresponding generating function is therefore $2t^2(1-t)^{-r-2}$, whose coefficient of $t^k$ is $2B$.
Using $a^2=a(a-1)+a$ gives the square sum $A+2B$.

Differentiating once in each of two distinct coordinates gives generating function
$t^2(1-t)^{-4}(1-t)^{-(r-2)}=t^2(1-t)^{-r-2}$,
so the mixed sum is $B$. When $r=1$, there are no mixed terms and this calculation is unnecessary.
Permutation symmetry makes the coordinate sums equal.
Summing $\sum_i\alpha_i=k$ over all $N$ compositions gives $rA=kN$, hence $A=kN/r$.
The cases $k=0,1$ are included automatically by the convention on zero binomial coefficients.
\end{proof}

\begin{theorem}[The first two Chern classes of symmetric powers in arbitrary rank]\label{bun:symmetric-chern}
For a complex vector bundle $E$ of rank $r\ge1$ and an integer $k\ge0$, $\operatorname{rk}(\operatorname{Sym}^kE)=N$, and
\begin{align}
 c_1(\operatorname{Sym}^kE)&=A\,c_1(E),\label{bun:sym-c1}\\
 \operatorname{ch}_2(\operatorname{Sym}^kE)
 &=\tfrac12(A+2B)c_1(E)^2-(A+B)c_2(E),\label{bun:sym-ch2}\\
 c_2(\operatorname{Sym}^kE)
 &=\tfrac12(A^2-A-2B)c_1(E)^2+(A+B)c_2(E).\label{bun:sym-c2}
\end{align}
Here $N,A,B$ are as in Lemma \ref{bun:composition-counts}.
\end{theorem}
\begin{proof}
Symmetric tensors of degree $k$ have monomial basis
$e_1^{\alpha_1}\cdots e_r^{\alpha_r}$ with $|\alpha|=k$, so the rank is the number $N$ of weak compositions.
First compute for a diagonal infinitesimal matrix $\operatorname{diag}(x_1,\ldots,x_r)$.
Its action is the sum of its actions on the factors of each monomial; hence the eigenvalue is
$w_\alpha=\sum_i\alpha_ix_i$, rather than a product of the $x_i$.
Lemma~\ref{bun:composition-counts} gives
\[
 \sum_\alpha w_\alpha=\sum_i\left(\sum_\alpha\alpha_i\right)x_i
 =A\sum_i x_i.
\]
Expand the square, grouping equal and distinct indices:
\begin{align*}
 \frac12\sum_\alpha w_\alpha^2
 &=\frac12\sum_i\left(\sum_\alpha\alpha_i^2\right)x_i^2
   +\sum_{i<j}\left(\sum_\alpha\alpha_i\alpha_j\right)x_ix_j\\
 &=\tfrac12(A+2B)\sum_i x_i^2+B\sum_{i<j}x_ix_j.
\end{align*}
Substitute $c_1=\sum_i x_i$ and $c_2=\sum_{i<j}x_ix_j$.
Since $\sum_i x_i^2=c_1^2-2c_2$, the expression becomes
$\tfrac12(A+2B)c_1^2-(A+B)c_2$.

We must explain why this diagonal calculation applies even when the vector bundle does not split.
The entries of the infinitesimal symmetric-power representation depend linearly on the original matrix entries, and its trace and the trace of its square are conjugation-invariant polynomials.
Complex matrices with distinct eigenvalues are dense and diagonalizable, so the identities just proved hold for every complex matrix.
They are therefore polynomial identities and remain valid in any commutative algebra over $\C$, in particular in even-degree differential forms.
On $E^{\otimes k}$, let the connection act factor by factor.
On squaring, the cross-factor first-derivative terms cancel by the wedge-product sign, leaving the sum of the actions of $F_E$ on the individual factors.
This operation commutes with permutations of tensor factors, so it preserves the symmetric subbundle.
Its restriction is exactly the infinitesimal symmetric-power representation applied to $F_E$.
Substitute $\frac\ii{2\pi}F_E$ into the polynomial identities to obtain the formulas for $c_1$ and $\operatorname{ch}_2$.

Finally, expand \eqref{bun:det-trace} to second order:
$c_1=\tr\Omega$ and
$c_2=\tfrac12((\tr\Omega)^2-\tr\Omega^2)$.
Thus $\operatorname{ch}_2=\tfrac12(c_1^2-2c_2)$.
Solving for $c_2=\tfrac12c_1^2-\operatorname{ch}_2$ and substituting the two expressions just computed gives \eqref{bun:sym-c2}.
\end{proof}

\begin{example}[Checking the formula directly]
For $r=3,k=2$, we have $N=6,A=4,B=1$, so
$c_1(\operatorname{Sym}^2E)=4c_1(E)$.
Indeed, the six monomials are $e_1^2,e_2^2,e_3^2,e_1e_2,e_1e_3,e_2e_3$.
The sum of the exponents of each fixed $e_i$ is $2+1+1=4$, which gives the coefficient of $c_1(E)$ directly.
When $r=2$, formula \eqref{bun:sym-c2} becomes
\[
 c_2(\operatorname{Sym}^kE)
 =\frac{k(k-1)(k+1)(3k+2)}{24}c_1(E)^2
 +\frac{k(k+1)(k+2)}6c_2(E).
\]
For $k=1$, the coefficient of $c_1(E)^2$ is zero and that of $c_2(E)$ is one, as required by $\operatorname{Sym}^1E=E$.
\end{example}

\begin{example}[Checking a symmetric power term by term in a split model]\label{bun:worked-symmetric}
On $\mathbb P^2$, set $h=c_1(\cO(1))$ and
$E=\cO(1)\oplus\cO(2)$. Then
\[
 c_1(\operatorname{Sym}^2E)=9h,\qquad
 c_2(\operatorname{Sym}^2E)=26h^2,\qquad
 \operatorname{ch}_2(\operatorname{Sym}^2E)=\tfrac{29}{2}h^2.
\]
\end{example}
\begin{proof}
Proposition~\ref{bun:chern-functors} gives $c(E)=(1+h)(1+2h)=1+3h+2h^2$.
In local line-bundle frames $e_1,e_2$, the three basis elements of $\operatorname{Sym}^2E$ are
$e_1^2,e_1e_2,e_2^2$. Their transition functions are those of
$\cO(2),\cO(3),\cO(4)$, respectively, so
$\operatorname{Sym}^2E=\cO(2)\oplus\cO(3)\oplus\cO(4)$.
Hence
\[
 c=(1+2h)(1+3h)(1+4h)
  =1+9h+(6+8+12)h^2,
\]
where $h^3=0$ because it is a degree-six class on a manifold of real dimension four.
The Chern character is $e^{2h}+e^{3h}+e^{4h}$, whose coefficient of $h^2$ is
$(2^2+3^2+4^2)/2=29/2$.

Now check the general formulas: for $r=2,k=2$, we have $N=3,A=3,B=1$.
Thus $c_1=3(3h)=9h$ and
$c_2=\tfrac12(9-3-2)(3h)^2+(3+1)(2h^2)=26h^2$.
The two independent calculations agree. They also show that the second Chern character may have a denominator; its components should not be confused with integral Chern classes.
\end{proof}

\begin{proposition}[A degree bound for higher Chern characters]\label{bun:symmetric-degree}
For fixed $r,q$, the expression $\operatorname{ch}_q(\operatorname{Sym}^kE)$ is a polynomial in the integer $k\ge0$, with coefficients in $H^{2q}_{\mathrm{dR}}(X,\C)$ and degree at most $r+q-1$.
\end{proposition}
\begin{proof}
The diagonal-matrix calculation gives
$q!\operatorname{ch}_q=\sum_{|\alpha|=k}(\sum\alpha_ix_i)^q$.
After expansion, it suffices to study sums of monomials $\prod\alpha_i^{m_i}$ of total degree $q$.
For an integer variable, $a^m$ is a linear combination of falling factorials
$(a)_\ell=a(a-1)\cdots(a-\ell+1)$, $0\le\ell\le m$.
This follows by induction on degree, since the leading term of $(a)_m$ is $a^m$.
Differentiating the geometric series $\ell$ times gives
$\sum_{a\ge0}(a)_\ell t^a=\ell!t^\ell(1-t)^{-\ell-1}$.
The generating function of each resulting term is therefore a constant times $t^d(1-t)^{-r-d}$, where $d\le q$.
Its coefficient is $\binom{k+r-1}{r+d-1}$, a polynomial of degree $r+d-1$ that automatically vanishes for $k<d$.
This gives the bound $r+q-1$. As in Theorem \ref{bun:symmetric-chern}, the polynomial identities extend from diagonal matrices to curvature, completing the proof.
\end{proof}

\begin{remark}
The same counting question arises for tensor constructions associated with other representations.
Schur functors and their Chern characters require more systematic representation theory; see \cite{Fulton1998,BottTu1982} for further reading.
In particular, the fact that each expanded monomial has degree $q$ does not make the total degree in $k$ equal to $q$: the number of summands also grows with $k$.
\end{remark}


\chapter{Sheaves, Cohomology, and the Dolbeault Theorem}\label{sh:chapter}

This chapter addresses a recurring question in complex geometry: why can an equation have local solutions but no global solution?
Sheaves record local objects and their restrictions; cohomology records the obstructions to gluing local solutions.
The Dolbeault theorem will identify sheaf cohomology with closed bundle-valued forms modulo forms having a global primitive.
This identification, Theorem~\ref{sh:dolbeault}, requires neither compactness nor the K\"ahler condition.
We define cohomology by the Godement resolution and prove the comparison theorem needed to replace that resolution by differential forms.
The prerequisites are differential forms, vector bundles, Stokes' theorem, and smooth partitions of unity from the preceding chapters.
For comparison with standard terminology, see \cite[Chapters II and IV]{Demailly2012}
and \cite[Sheaves on Spaces; Cohomology of Sheaves]{Stacks}.
On a first reading, the main points are stalkwise exactness, the obstruction illustrated by local logarithms, and the passage from a fine resolution to Dolbeault cohomology.
The general pullback formulas can be read when they are needed for projective bundles.
In the topological extension arguments, it is enough at first to keep a second countable manifold in mind; its covers can be reduced to countable ones, so the successive gluing constructions have a concrete form.

\section{From Local Functions to Sheaves}

\begin{definition}[Presheaves and sheaves]\label{sh:def-sheaf}
Let $X$ be a topological space. A \emph{presheaf of abelian groups} $\cF$ assigns an abelian group $\cF(U)$ to every open set $U$ and a restriction homomorphism
$r_{UV}:\cF(U)\to\cF(V)$ to every inclusion $V\subset U$, satisfying
$r_{UU}=\operatorname{id}$ and $r_{VW}r_{UV}=r_{UW}$.
We usually write $r_{UV}(s)=s|_V$.
It is a \emph{sheaf} if, for every open cover $U=\bigcup_iU_i$, the following hold:
\begin{enumerate}
\item If $s,t\in\cF(U)$ agree on every $U_i$, then $s=t$.
\item If $s_i\in\cF(U_i)$ satisfy $s_i|_{U_i\cap U_j}=s_j|_{U_i\cap U_j}$,
there is $s\in\cF(U)$ with $s|_{U_i}=s_i$.
\end{enumerate}
The first condition makes the section in the second condition unique.
A sheaf morphism is a family of homomorphisms on open sets that commutes with all restriction maps.
We also write $\Gamma(U,\cF)=\cF(U)$.
\end{definition}

\begin{example}\label{sh:example-basic}
Continuous functions, smooth functions, holomorphic functions, and smooth differential forms all form sheaves.
On a complex manifold we denote the holomorphic function sheaf by $\cO_X$, the smooth function sheaf by $\cA_X^0$, and the sheaves of smooth forms by $\cA_X^k$ and $\cA_X^{p,q}$.
Thus $\cA_X^{p,q}(U)=A^{p,q}(U)$ in the notation of Definition~\ref{bas:cohomology-def}.
For instance, glue compatible smooth functions pointwise to obtain $f$.
Every point has a neighborhood where $f$ equals one of the given smooth functions, so $f$ is smooth, proving existence; pointwise equality proves uniqueness.
Continuous, smooth, and holomorphic sections of vector bundles have the same property, since regularity can be checked in local frames.
\end{example}

\begin{example}[A constant presheaf need not be a sheaf]\label{sh:constant-presheaf}
Assigning the same group $A$ to every nonempty open set, with identity restriction maps and value zero on the empty set, generally gives only a presheaf.
If $U=U_1\sqcup U_2$, choose different constants on the two nonempty open sets.
They are compatible on the empty intersection but cannot arise from one constant on $U$.
The \emph{constant sheaf} $\underline A$ is the sheaf of locally constant $A$-valued functions, with $A$ regarded as discrete.
\end{example}

\begin{definition}[Stalks and germs]\label{sh:def-stalk}
Fix $x\in X$. Consider pairs $(U,s)$, where $U$ is an open neighborhood of $x$ and $s\in\cF(U)$.
Declare $(U,s)\sim(V,t)$ if $s,t$ agree on some common open neighborhood of $x$.
The equivalence classes form the \emph{stalk} $\cF_x$ of $\cF$ at $x$.
The class of $s$ is denoted $s_x$ and called its \emph{germ} at $x$.
This is exactly the direct limit $\varinjlim_{U\ni x}\cF(U)$ of the restriction system: representatives become equal once restricted far enough to a common smaller neighborhood.
The support of a sheaf means $\supp\cF=\{x:\cF_x\ne0\}$; for a general sheaf this set need not be closed.
\end{definition}

For holomorphic functions, a germ contains more information than a value.
At $0\in\C$, the functions $z$ and $0$ have the same value but different germs, since they do not agree on any neighborhood of zero.
Passing to a stalk allows the neighborhood to shrink; it does not discard the function's local behavior.

\begin{lemma}[Germs detect sections]\label{sh:germs-detect}
For a sheaf $\cF$ and sections $s,t\in\cF(U)$, we have $s=t$ if and only if
$s_x=t_x$ for every $x\in U$.
\end{lemma}
\begin{proof}
One implication follows from the definition of restriction.
Conversely, equality of germs gives, for each $x$, an open neighborhood $V_x\subset U$ on which $s|_{V_x}=t|_{V_x}$.
These neighborhoods cover $U$, so the uniqueness axiom gives $s=t$.
\end{proof}

\begin{proposition}[An explicit construction of sheafification]\label{sh:sheafification}
For a presheaf of abelian groups $\mathcal P$, let $\mathcal P^+(U)$ consist of families of germs
$(a_x)_{x\in U}$, with $a_x\in\mathcal P_x$, that satisfy the following condition: every $x$ has a neighborhood $V\subset U$ and a section $s\in\mathcal P(V)$ such that $a_y=s_y$ for all $y\in V$.
Pointwise addition and restriction make $\mathcal P^+$ a sheaf.
The natural map $\eta:\mathcal P\to\mathcal P^+$, $s\mapsto(s_x)$, induces an isomorphism on every stalk.
If $\cF$ is a sheaf, every presheaf morphism $a:\mathcal P\to\cF$ factors uniquely as
$a=\widetilde a\eta$, where $\widetilde a:\mathcal P^+\to\cF$ is a sheaf morphism.
\end{proposition}
\begin{proof}
Compatible families of germs glue pointwise on an open cover, and the glued family is still locally represented by a presheaf section. Thus the sheaf axioms hold.
If $\sigma\in\mathcal P^+(U)$ is represented near $x$ by $s\in\mathcal P(V)$, then $\sigma_x$ comes from $s_x$, proving surjectivity of $\eta_x$.
If the image of $s_x$ is zero, then $\eta(s)$ vanishes on a neighborhood $W$ of $x$.
In particular its component at $x$ is $s_x=0$, proving injectivity.

To construct $\widetilde a$, choose local representatives $s_i\in\mathcal P(U_i)$ of $\sigma$ and take $a(s_i)\in\cF(U_i)$.
At each point of an overlap the germs of $s_i,s_j$ agree, so the germs of $a(s_i),a(s_j)$ agree.
Lemma \ref{sh:germs-detect} makes these sections equal on the overlap.
They therefore glue uniquely to $\widetilde a(\sigma)$.
Changing local representatives changes no germ and hence no resulting section.
Addition and restriction are preserved, as is clear on local representatives.
Any factorization must take these values on each $U_i$, proving uniqueness.
\end{proof}

\begin{proposition}[Isomorphisms can be checked on stalks]\label{sh:stalk-isomorphism}
A sheaf morphism $a:\cF\to\cG$ is an isomorphism if and only if every
$a_x:\cF_x\to\cG_x$ is an isomorphism.
\end{proposition}
\begin{proof}
Necessity follows by taking germs of the inverse morphism.
Conversely, if $a(s)=0$, then every $s_x=0$, so Lemma~\ref{sh:germs-detect} gives $s=0$; thus the map on each open set is injective.
Given $t\in\cG(U)$, use surjectivity of $a_x$ at each $x$ to choose a preimage germ represented by $s_x'\in\cF(V_x)$.
Shrink $V_x$ until $a(s_x')=t|_{V_x}$, which is possible because the germs agree at $x$.
On overlaps, $s_x'$ and $s_y'$ have the same image, so the injectivity just proved makes them equal.
They glue to $s\in\cF(U)$ with $a(s)=t$.
These inverse maps on open sets automatically commute with restrictions.
\end{proof}

\section{The Local Meaning of Exactness}

\begin{definition}[Kernels, images, quotients, and exact sequences]\label{sh:def-exact}
For a sheaf morphism $a:\cF\to\cG$, define its kernel by
$(\ker a)(U)=\ker(\cF(U)\to\cG(U))$.
Its image and cokernel sheaves are the sheafifications of the presheaves
$U\mapsto a(\cF(U))$ and $U\mapsto\cG(U)/a(\cF(U))$, respectively.
A sequence of sheaf morphisms is \emph{exact} if at each position the image sheaf equals the kernel sheaf of the next map.
In particular, exactness of $\cG\to\cH\to0$ means that \emph{sections lift locally}; it does not assert global lifting on every open set.
\end{definition}

\begin{proposition}[The stalk criterion for exactness]\label{sh:stalk-exact}
A sequence of sheaves of abelian groups is exact if and only if its sequence of stalks is exact at every point.
For a short exact sequence
\begin{equation}\label{sh:ses}
 0\longrightarrow\cF\xrightarrow{a}\cG\xrightarrow{b}\cH\longrightarrow0,
\end{equation}
and every open set $U$, the sequence
\[
0\longrightarrow\cF(U)\xrightarrow{a}\cG(U)\xrightarrow{b}\cH(U)
\]
is left exact.
\end{proposition}
\begin{proof}
Kernels commute with passage to stalks: if $a(s)_x=0$, shrink the neighborhood representing $s$ until $a(s)=0$, so that the germ comes from the kernel sheaf.
Images are handled by local representatives.
For cokernels, Proposition \ref{sh:sheafification} says that sheafification does not change stalks; restrict the finitely many representatives and their equations to a common neighborhood.
Thus stalks of kernels, images, and cokernels are the corresponding kernels, images, and cokernels of group homomorphisms.
Proposition \ref{sh:stalk-isomorphism} now gives the criterion.

Suppose $t\in\cG(U)$ satisfies $b(t)=0$.
Stalkwise exactness gives local sections $s_i$ with $t=a(s_i)$.
On overlaps, $a(s_i-s_j)=0$; stalkwise injectivity and Lemma \ref{sh:germs-detect} imply $s_i=s_j$.
The sections glue to $s\in\cF(U)$ with $a(s)=t$.
The same stalkwise injectivity proves injectivity of $\cF(U)\to\cG(U)$.
\end{proof}

\begin{example}[Local logarithms need not give a global logarithm]\label{sh:exp-obstruction}
Write $\C^*=\C\setminus\{0\}$ and let $\cO_X^*$ be the sheaf of nowhere-zero holomorphic functions, with multiplication as its abelian-group law. On $X=\C^*$ there is a short exact sequence
\[
0\longrightarrow\underline\Z\longrightarrow\cO_X
 \xrightarrow{\ f\mapsto e^{2\pi\ii f}\ }\cO_X^*
 \longrightarrow1.
\]
The kernel consists of locally constant integer-valued functions.
A nowhere-zero holomorphic function has a holomorphic logarithm on sufficiently small discs: divide by its value at the center, shrink the disc until the quotient takes values in $|w-1|<1$, and use the power series for $\log(1+z)$, together with a chosen logarithm of the constant value.
This proves surjectivity as a sheaf morphism.
But the section $g(z)=z$ has no global logarithm.
If $e^{2\pi\ii f(z)}=z$, differentiation gives $2\pi\ii f'(z)=1/z$.
Integration around the unit circle gives zero on the left and $2\pi\ii$ on the right, a contradiction.
Thus global sections do not generally preserve surjectivity at the end of a short exact sequence.
\end{example}

\begin{definition}[Direct images, inverse images, and sections on subspaces]\label{sh:inverse-image}
Let $f:X\to Y$ be continuous. The direct image of a sheaf $\cF$ on $X$ is
$(f_*\cF)(V)=\cF(f^{-1}V)$.
The inverse image $f^{-1}\cG$ of a sheaf $\cG$ on $Y$ is the sheafification of the presheaf
\[
 U\longmapsto\varinjlim_{\substack{V\text{ open in }Y\\ f(U)\subset V}}
 \cG(V).
\]
For the inclusion $i:S\hookrightarrow X$ of any subspace, define
$\cF|_S=i^{-1}\cF$ and $\cF(S)=\Gamma(S,i^{-1}\cF)$.
Here $i^{-1}$ acts on a \emph{sheaf}; the expression is not $i^{-1}(S)$.
A section in $\cF(S)$ is a family of germs $s(x)\in\cF_x$ locally represented near each point of $S$ by a section on an ambient open subset of $X$.
\end{definition}

\begin{proposition}[Exactness of inverse and direct images]\label{sh:inverse-exact}
There is a natural isomorphism $(f^{-1}\cG)_x\simeq\cG_{f(x)}$.
Consequently, $f^{-1}$ preserves exact sequences of sheaves of abelian groups, whereas $f_*$ is left exact and need not preserve surjectivity.
\end{proposition}
\begin{proof}
An inverse-image presheaf section is represented by a section on some $V\supset f(U)$.
Send it to its germ at $f(x)$ to obtain the proposed map.
Every neighborhood $V$ of $f(x)$ contains $f(f^{-1}V)$, so a section on $V$ represents an inverse-image section on $U=f^{-1}V$; this proves surjectivity.
If two representatives have the same germ at $f(x)$, they agree on a smaller neighborhood $W$.
Restriction to $f^{-1}W$ shows that their inverse-image germs agree, proving injectivity.
Proposition~\ref{sh:sheafification} says that sheafification does not change stalks, so the assertion holds for $f^{-1}\cG$.
Exactness follows from Proposition \ref{sh:stalk-exact}.
Left exactness of the direct image follows by applying the left exactness assertion of that proposition on $f^{-1}V$.
Finally, for $f:\C^*\to\{\mathrm{pt}\}$, Example \ref{sh:exp-obstruction} shows that the direct image need not preserve surjectivity.
\end{proof}

\begin{proposition}[Inverse-image sections as families of germs]\label{sh:inverse-germ-family}
Elements of $\Gamma(U,f^{-1}\cG)$ are equivalent to families assigning
$a_x\in\cG_{f(x)}$ to each $x\in U$, with this condition: every $x$ has an open neighborhood $W\subset U$, an open neighborhood $V$ of $f(x)$, and $s\in\cG(V)$ such that
$f(W)\subset V$ and $a_y=s_{f(y)}$ for all $y\in W$.
The family must be indexed by $x\in U$, rather than merely by points of $f(U)$.
\end{proposition}
\begin{proof}
The inverse image is a sheafification, which Proposition \ref{sh:sheafification} describes by locally representable families of germs.
Proposition \ref{sh:inverse-exact} identifies the stalk of its defining presheaf at $x$ with $\cG_{f(x)}$.
A local representative is therefore precisely a section $s$ on some $V\supset f(W)$, as stated.
Conversely, a family satisfying this condition is represented on each $W$ by an inverse-image presheaf section, hence defines a unique section by the construction of sheafification.
The two constructions are inverse.

The following example explains why the indexing set matters.
Let $X$ consist of two discrete points, let $Y$ be a point, let $f$ be constant, and let $\cG(Y)=A$.
The inverse image has section group $A\times A$ on $X$, since the local sections at the two points can be chosen independently.
Indexing only by the single image point would give just one copy of $A$, omitting sections with different values on the two components.
\end{proof}

\begin{proposition}[Gluing over a locally finite closed cover]\label{sh:closed-gluing}
Let $\{S_i\}$ be a locally finite closed cover of $X$.
If $s_i\in\cF(S_i)$ agree on every $S_i\cap S_j$, there is a unique
$s\in\cF(X)$ with $s|_{S_i}=s_i$.
In particular, this holds for every finite closed cover.
Equivalently, the sequence
\[
0\to\cF(X)\to\prod_i\cF(S_i)
 \xrightarrow{\ (s_i)\mapsto(s_j|_{S_i\cap S_j}-s_i|_{S_i\cap S_j})\ }
 \prod_{i,j}\cF(S_i\cap S_j)
\]
is exact.
\end{proposition}
\begin{proof}
Compatibility first defines a family of germs $s(x)$. Fix $x$.
Local finiteness gives a neighborhood meeting only finitely many $S_i$.
Remove from it all these closed sets that do not contain $x$, leaving an open neighborhood $U$.
Every $S_i$ meeting $U$ now contains $x$.
For each such $i$, the section $s_i$ is represented near $x$ by an ambient section $t_i$.
Agreement of the germs on intersections gives $(t_i)_x=(t_j)_x$.
Only finitely many pairs occur, so shrink $U$ simultaneously until all $t_i$ agree there and represent the respective $s_i$ on $U\cap S_i$.
Their common section on $U$ represents the entire germ family under consideration.
The family is therefore locally represented by sections and glues globally by the sheaf axiom, or equivalently by Proposition~\ref{sh:sheafification}.
Uniqueness follows because every germ has been prescribed.
\end{proof}

\begin{remark}\label{sh:closed-caution}
Local finiteness is essential.
For example, on the compact space $X=\{0\}\cup\{1/n:n\ge1\}$, take the closed cover by all singletons.
For the sheaf of locally constant real functions, prescribe the zero germ at $0$ and the germ of the constant function $1$ at each $1/n$.
Distinct singletons have empty intersection, so compatibility imposes no condition.
These germs cannot glue to a function locally constant near $0$.
Thus the gluing property for locally finite closed covers does not extend to arbitrary closed covers.
\end{remark}

\section{Module Sheaves, Locally Free Sheaves, and Pullbacks}

\begin{definition}\label{sh:ringed-space}
A \emph{ringed space} $(X,\cO_X)$ is a topological space equipped with a sheaf of commutative rings with identity.
For an open subset $U$, write $\cO_U=\cO_X|_U$.
Modules are unital, and $\Hom_{\cO_X}(\cF,\cG)$ denotes the group of $\cO_X$-linear sheaf morphisms.
An $\cO_X$-module sheaf has restriction homomorphisms compatible with local scalar multiplication.
A module sheaf $\cF$ is \emph{locally free of finite rank} if every point has a neighborhood $U$ and a finite integer $r$ such that
$\cF|_U\simeq\cO_U^{\oplus r}$.
The case $r=1$ is that of an invertible sheaf.
On a complex manifold, the sheaf of holomorphic sections of a holomorphic vector bundle $E$ is denoted $\cO(E)$.
A holomorphic local frame gives $\cO(E)|_U\simeq\cO_U^{\oplus r}$.
\end{definition}

For modules over a commutative ring $R$, the tensor product $M\otimes_RN$ is generated by symbols $m\otimes n$, subject to additivity in each variable and the balancing relations $(rm)\otimes n=m\otimes(rn)$.
Consequently, an $R$-balanced bilinear map on $M\times N$ factors uniquely through $M\otimes_RN$; this is its universal property.

\begin{proposition}[Local properties of tensor products]\label{sh:tensor}
The tensor product of module sheaves is the sheafification of
$U\mapsto\cF(U)\otimes_{\cO_X(U)}\cG(U)$, and satisfies
\[
 (\cF\otimes_{\cO_X}\cG)_x
 \simeq\cF_x\otimes_{\cO_{X,x}}\cG_x.
\]
Tensor product is right exact. If one factor is locally free, tensoring with it preserves every exact sequence.
\end{proposition}
\begin{proof}
Send a local tensor $s\otimes t$ to $s_x\otimes t_x$.
Every tensor of germs is a finite sum; all germs involved can be represented on a common neighborhood, proving surjectivity.
Tensor relations likewise involve only finitely many terms and scalar germs.
If a relation holds at the stalk, restrict all representatives and equations to a smaller common neighborhood to obtain the same relation among local tensors.
This proves injectivity.

For modules over a ring $R$, the universal property of tensor product directly gives
\[
 M\otimes_R(N/N')
 \simeq (M\otimes_RN)/\operatorname{im}(M\otimes_RN').
\]
This proves right exactness for modules, and the stalk criterion gives the sheaf assertion.
If $\cF|_U=\cO_U^{\oplus r}$, then
$(\cF\otimes\cG)|_U=\cG|_U^{\oplus r}$.
Finite direct sums preserve exactness componentwise, so tensoring with a locally free sheaf is exact at every stalk.
\end{proof}

An $R$-module $M$ is \emph{flat} if tensoring with it preserves injections of $R$-modules.
Right exactness of tensor product, proved in Proposition~\ref{sh:tensor}, makes this equivalent to preservation of short exact sequences.

\begin{proposition}[Flatness can be checked on stalks]\label{sh:flat-stalks}
The tensor functor $-\otimes\cF$ is exact if and only if each
$\cF_x$ is flat over $\cO_{X,x}$.
In particular, locally free sheaves are flat.
\end{proposition}
\begin{proof}
If all stalks are flat, take stalks of any sheaf injection, tensor, and apply Proposition \ref{sh:tensor}.
The resulting maps on stalks are injective, so Proposition~\ref{sh:stalk-exact} makes the tensor-product sheaf morphism injective.

Conversely, fix $x$, set $A=\cO_{X,x}$, and take any injection $M'\to M$ of $A$-modules.
Construct the skyscraper module sheaf at $x$: its value on an open set containing $x$ is $M$, and its value on an open set not containing $x$ is zero.
Local functions act through their germs at $x$.
The same construction for $M'$ gives a sheaf injection whose stalk at $x$ is the original module injection.
Tensoring with $\cF$ preserves injectivity, and taking the stalk at $x$ gives an injection
$M'\otimes_A\cF_x\to M\otimes_A\cF_x$.
As $M',M$ were arbitrary, $\cF_x$ is flat.
Tensoring with a locally free module is a direct sum, which preserves injections componentwise, proving the last assertion.
\end{proof}

\begin{definition}[Pullback of module sheaves]\label{sh:module-pullback}
A morphism of ringed spaces $f:(X,\cO_X)\to(Y,\cO_Y)$ includes a ring-sheaf morphism
$f^\sharp:f^{-1}\cO_Y\to\cO_X$.
For a module sheaf on $Y$, define
\[
 f^*\cG=\cO_X\otimes_{f^{-1}\cO_Y}f^{-1}\cG,
 \qquad
 (f^*\cG)_x
 =\cO_{X,x}\otimes_{\cO_{Y,f(x)}}\cG_{f(x)}.
\]
The stalk formula follows from Propositions \ref{sh:inverse-exact} and \ref{sh:tensor}.
Inverse image $f^{-1}$ and module pullback $f^*$ are different constructions; $f^*$ is generally only right exact.
\end{definition}

\begin{example}[Pullback can destroy injectivity]\label{sh:pullback-not-exact}
Let $i:\{0\}\hookrightarrow\C$, with structure ring $\C$ on the point.
Multiplication by the coordinate gives an injection $\cO_\C\xrightarrow{z}\cO_\C$:
if $zf=0$, then $f$ vanishes where $z\ne0$ and hence also at $0$ by continuity.
After pullback to the point, however, the map becomes $\C\xrightarrow{0}\C$, which is not injective.
A general projection formula therefore cannot be justified by claiming that $f^*$ is always exact.
\end{example}

\begin{proposition}[Pullback--direct-image adjunction]\label{sh:adjunction}
Let $f:(X,\cO_X)\to(Y,\cO_Y)$ be a morphism of ringed spaces, let $\cG$ be a module sheaf on $Y$, and let $\cH$ be one on $X$.
There is a natural bijection
\[
 \Hom_{\cO_X}(f^*\cG,\cH)
 \simeq \Hom_{\cO_Y}(\cG,f_*\cH).
\]
If $\mathcal{H}om$ denotes the internal Hom sheaf, defined by
$\mathcal{H}om(\cF,\cG)(U)=\Hom(\cF|_U,\cG|_U)$, there is also a sheaf isomorphism
\[
 f_*\mathcal{H}om_{\cO_X}(f^*\cG,\cH)
 \simeq\mathcal{H}om_{\cO_Y}(\cG,f_*\cH).
\]
\end{proposition}
\begin{proof}
First establish the adjunction between $f^{-1}$ and $f_*$ for sheaves of abelian groups.
Given $a:\cG\to f_*\cH$, take a representative $s\in\cG(V)$ of an inverse-image presheaf section on $U$, where $f(U)\subset V$.
Apply $a$ to obtain $a(s)\in\cH(f^{-1}V)$, then restrict to $U$.
Shrinking $V$ does not change this section, so the construction passes to the direct limit and defines a presheaf morphism.
Proposition~\ref{sh:sheafification} gives $\widetilde a:f^{-1}\cG\to\cH$.
Conversely, $s\in\cG(V)$ has a natural inverse-image section on $f^{-1}V$; apply $\widetilde a$ to it to recover $a(s)$.
Both composites are the identity on these local representatives, so the constructions are inverse.

If $a$ is $\cO_Y$-linear, then $\widetilde a$ is $f^{-1}\cO_Y$-linear.
The universal property of tensor product extends it uniquely to the $\cO_X$-linear morphism
\[
 \cO_X\otimes_{f^{-1}\cO_Y}f^{-1}\cG\longrightarrow\cH,
 \qquad u\otimes s\longmapsto u\,\widetilde a(s).
\]
Conversely, restriction to $1\otimes s$ recovers $\widetilde a$.
This proves the first formula.
Apply it on every open set $V\subset Y$ to the restricted map $f^{-1}V\to V$.
The constructions commute with restriction, so the resulting bijections form the sheaf isomorphism in the second formula.
\end{proof}

\begin{example}[Direct image cannot be moved inside Hom without hypotheses]\label{sh:hom-counterexample}
Omitting $f^*$ from the preceding proposition and writing
$f_*\mathcal{H}om(\cF,\cG)=\mathcal{H}om(f_*\cF,f_*\cG)$ is generally incorrect.
Let $X$ consist of two discrete points, let $Y$ be one point, give every point structure ring $\C$, and take $\cF=\cG=\cO_X$.
The left side consists of pointwise linear endomorphisms and equals $\C^2$.
The right side is $\Hom_\C(\C^2,\C^2)$, of dimension four.
The natural map produces only diagonal matrices and cannot mix the components belonging to different points.
\end{example}

\begin{proposition}[Pullback preserves tensor products and the unit]\label{sh:pullback-tensor}
For a morphism of ringed spaces $f:X\to Y$, there are natural isomorphisms
\[
 f^*(\cF\otimes_{\cO_Y}\cG)
 \simeq f^*\cF\otimes_{\cO_X}f^*\cG,
 \qquad f^*\cO_Y\simeq\cO_X.
\]
\end{proposition}
\begin{proof}
Fix $x\in X$ and write $A=\cO_{Y,f(x)}$, $B=\cO_{X,x}$,
$M=\cF_{f(x)}$, and $N=\cG_{f(x)}$.
The two directions of the first isomorphism on stalks are
\begin{align*}
 B\otimes_A(M\otimes_AN)&\longrightarrow
       (B\otimes_AM)\otimes_B(B\otimes_AN),\\
 b\otimes(m\otimes n)&\longmapsto(b\otimes m)\otimes(1\otimes n),\\[2pt]
 (B\otimes_AM)\otimes_B(B\otimes_AN)&\longrightarrow
       B\otimes_A(M\otimes_AN),\\
 (b\otimes m)\otimes(c\otimes n)&\longmapsto bc\otimes(m\otimes n).
\end{align*}
Moving an $A$-scalar from $M$ or $N$ to $B$, or moving a $B$-scalar across the outer tensor product, leaves the formulas unchanged.
Thus all tensor relations are respected.
The first map followed by the second is the identity.
The other composite sends $(b\otimes m)\otimes(c\otimes n)$ to
$(bc\otimes m)\otimes(1\otimes n)$, equal to the original element by the $B$-balancing relation.
The same formulas apply to local sections and commute with restriction, defining sheaf morphisms.
Proposition~\ref{sh:stalk-isomorphism} turns their inverseness on stalks into the asserted sheaf isomorphism.

For the second formula, the maps are $b\otimes a\mapsto bf^\sharp(a)$ and $b\mapsto b\otimes1$.
The balancing identity $b\otimes a=bf^\sharp(a)\otimes1$ makes them inverse.
No flatness assumption is used.
\end{proof}

\begin{proposition}[Evaluation and the general projection morphism]\label{sh:projection-morphism}
For a module sheaf $\cF$ on $X$ and a module sheaf $\cG$ on $Y$, there are natural morphisms
\[
 f^*((f_*\cF)\otimes\cG)\longrightarrow\cF\otimes f^*\cG,
 \qquad
 (f_*\cF)\otimes\cG\longrightarrow f_*(\cF\otimes f^*\cG).
\]
\end{proposition}
\begin{proof}
First construct evaluation $\varepsilon:f^*f_*\cF\to\cF$.
A local representative of $f^{-1}f_*\cF$ on $U\subset X$ is a section of the direct image on some $V\supset f(U)$, namely $t\in\cF(f^{-1}V)$.
Restrict it to $U$.
This is compatible with shrinking $V$ and therefore sheafifies to $f^{-1}f_*\cF\to\cF$.
Then $b\otimes t\mapsto b\,t|_U$ respects the balancing relations and gives $\varepsilon$.
Under Proposition \ref{sh:pullback-tensor}, the first morphism is $\varepsilon\otimes\operatorname{id}_{f^*\cG}$.

The second follows by adjunction, Proposition \ref{sh:adjunction}, from the first.
It can also be written directly on $V\subset Y$ as
\[
 s\otimes t\longmapsto s\otimes(1\otimes f^{-1}t),
 \quad s\in\cF(f^{-1}V),\quad t\in\cG(V).
\]
For $a\in\cO_Y(V)$, the relation $(as)\otimes t=s\otimes(at)$ on the left is preserved on the right by moving $f^\sharp(a)$ between tensor factors.
The map is therefore well defined and sheafifies.
Its compatibility with restrictions proves naturality.
\end{proof}

\begin{proposition}[The projection formula for a locally free factor of finite rank]\label{sh:projection-formula}
Let $f:(X,\cO_X)\to(Y,\cO_Y)$ be a morphism of ringed spaces,
let $\cF$ be an $\cO_X$-module sheaf, and let $\cG$ be a locally free $\cO_Y$-module sheaf of finite rank.
Then there is a natural isomorphism
\[
 (f_*\cF)\otimes_{\cO_Y}\cG
 \simeq f_*(\cF\otimes_{\cO_X}f^*\cG).
\]
\end{proposition}
\begin{proof}
Take the second natural morphism of Proposition \ref{sh:projection-morphism}.
It suffices to check that it is an isomorphism on an open cover of $Y$.
Choose a frame $e_1,\ldots,e_r$ for $\cG|_V\simeq\cO_V^r$.
Every section of the left side over $V$ is uniquely $\sum_{j=1}^rs_j\otimes e_j$, with $s_j\in(f_*\cF)(V)$.
Proposition \ref{sh:pullback-tensor} identifies $(f^*\cG)|_{f^{-1}V}$ with $\cO_{f^{-1}V}^r$, framed by $1\otimes e_j$.
Every section of the right side is therefore uniquely $\sum_js_j\otimes(1\otimes e_j)$.
The projection morphism sends the first expression to the second, and its inverse recovers the coefficients $s_j$.

On $V$, the map is thus the identity on $r$ copies of $(f_*\cF)|_V$.
Under a frame change both sides transform by the same matrix, with coefficients pulled back by $f^\sharp$.
The local isomorphisms agree with the coordinate-independent natural morphism and glue globally.
Finite rank ensures that only finite direct sums occur; a general direct image need not commute with arbitrary infinite direct sums.
\end{proof}

\begin{remark}[Coherence and finitely generated relations]\label{sh:coherence-scope}
A module sheaf is locally finitely generated if near each point finitely many sections generate every stalk, equivalently if there is a local surjection from a finite free sheaf. It is coherent if it is locally finitely generated and the kernel of every morphism from a finite free sheaf to it is again locally finitely generated.
Chapter~\ref{alg:chapter} proves coherence of the structure sheaf from Weierstrass division (Theorem~\ref{alg:oka}), then establishes finite presentations and the kernel, cokernel, and pullback properties of coherent sheaves.
The Dolbeault theorem applies to holomorphic vector bundles without requiring this coherence theorem.
\end{remark}

\section{Closed Neighborhoods, Partitions of Unity, and Extension}

The next arguments turn compatible local sections near a closed set into a single section on a neighborhood.
A space is \emph{paracompact} if every open cover admits a locally finite open refinement.
A family of sets is locally finite if each point has a neighborhood meeting only finitely many members.
A Hausdorff space is \emph{regular} if a point and a disjoint closed set have disjoint open neighborhoods, and \emph{normal} if any two disjoint closed sets have disjoint open neighborhoods.
Smooth manifolds satisfy the assumptions used here, but the proofs begin with Hausdorff paracompact spaces so that the precise requirements of soft-sheaf theory remain visible.

\begin{lemma}[Closed refinements and partitions of unity on paracompact spaces]\label{sh:paracompact-topology}
A Hausdorff paracompact space is normal.
Every open cover has a locally finite closed refinement $\{F_j\}$ with
$F_j\subset U_{i(j)}$ and $\bigcup_j\operatorname{int}F_j=X$.
Every locally finite open cover $\{U_i\}$ admits a subordinate continuous partition of unity
$\rho_i\ge0$, $\supp\rho_i\subset U_i$, $\sum_i\rho_i=1$.
Closed subspaces are also paracompact.
\end{lemma}
\begin{proof}
First prove regularity. Given a closed set $A$ and $x\notin A$, for each $a\in A$ choose disjoint open neighborhoods $V_a\ni a$, $W_a\ni x$, using the Hausdorff property.
Take a locally finite open refinement $\{T_j\}$ of $\{X\setminus A\}\cup\{V_a\}$.
Retain only members meeting $A$; each lies in some $V_a$, so $x\notin\overline{T_j}$.
The closures of a locally finite family remain locally finite, and their union is closed.
Thus $\bigcup_jT_j$ is an open neighborhood of $A$ whose closure does not contain $x$.
Its closure's complement separates $x$ from $A$.

For disjoint closed sets $A,B$, regularity supplies, for each $a\in A$, an open neighborhood $V_a$ with $\overline{V_a}\cap B=\varnothing$.
Refine $\{X\setminus A\}\cup\{V_a\}$ again and retain the members meeting $A$.
Their union $V$ contains $A$, while $\overline V$ misses $B$.
The sets $V$ and $X\setminus\overline V$ are the required disjoint open neighborhoods, proving normality.
The fact used here, that a locally finite union of closed sets is closed, is checked pointwise: near a point outside the union, only finitely many closed sets need be avoided; all other members already miss the neighborhood.

For an arbitrary open cover, choose at each point $V_x$ with
$x\in V_x$ and $\overline{V_x}\subset U_{i(x)}$.
Take a locally finite open refinement $T_j\subset V_{x(j)}$ and put $F_j=\overline{T_j}$.
These form the claimed closed refinement.
If the original cover $\{U_i\}$ is locally finite, combine all $F_j$ assigned to the same $U_i$ into a set $F_i$.
Every subunion of the locally finite closed family is closed, so $F_i$ is closed, lies in $U_i$, and the $F_i$ cover $X$.
Normality gives an open set $V_i$ with $F_i\subset V_i\subset\overline{V_i}\subset U_i$.

For completeness, we construct the continuous functions separating closed sets.
Given closed $F\subset V$ with $V$ open, recursively choose open sets $O_t$ for dyadic numbers $0<t<1$, with
$F\subset O_t$ and $\overline{O_s}\subset O_t$ for $s<t$, and set $O_1=V$.
At each step we insert an open set between an already chosen closed set and an open set, as normality permits.
Define $g(x)=\inf\{t:x\in O_t\}$, taking the infimum of the empty set to be $1$.
Then $g=0$ on $F$ and $g=1$ on $X\setminus V$.
Moreover, $\{g<t\}=\bigcup_{s<t}O_s$ and
$\{g>t\}=\bigcup_{s>t}(X\setminus\overline{O_s})$ for $0\le t<1$, with the endpoint cases immediate; hence $g$ is continuous.
Apply this to $(F_i,V_i)$ and set $h_i=1-g_i$.
Then $h_i=1$ on $F_i$ and $\supp h_i\subset\overline{V_i}\subset U_i$.
Local finiteness makes $h=\sum_i h_i$ continuous, and the covering property gives $h\ge1$.
Set $\rho_i=h_i/h$ to obtain the partition of unity.

Finally, for a closed subspace $S$, write an open cover of $S$ as $\{S\cap U_i\}$.
Refine the open cover $\{U_i\}\cup\{X\setminus S\}$ of $X$, then intersect its members with $S$.
The result remains locally finite and refines the original cover, proving paracompactness of $S$.
\end{proof}

\begin{proposition}[Why manifolds are paracompact]\label{sh:second-countable-paracompact}
Every second countable, locally compact Hausdorff space is paracompact.
In particular, all real and complex manifolds considered in these lectures are Hausdorff and paracompact.
\end{proposition}
\begin{proof}
A locally compact Hausdorff space is regular: first choose a compact neighborhood $K$ of $x$, then separate $x$ from the closed set to be avoided inside the compact Hausdorff space $K$.
In a compact Hausdorff space, such separation follows by taking finitely many intersections and unions of Hausdorff separating neighborhoods.
Thus every point has arbitrarily small relatively compact open neighborhoods.
Second countability gives a countable subcover of these neighborhoods: for every element of a countable basis contained in some cover member, choose one such member.
We can therefore construct recursively compact sets
\[
 K_1\subset\operatorname{int}K_2\subset K_2
 \subset\operatorname{int}K_3\subset\cdots,\qquad\bigcup_jK_j=X.
\]
At each step, cover the preceding compact set by finitely many relatively compact open sets and add the closures of the first $j$ members of the countable cover.
Their finite union gives the next compact set.

Now fix any open cover, and set $K_j=\varnothing$ for $j\le0$.
Each compact layer $K_j\setminus\operatorname{int}K_{j-1}$ admits a finite open cover whose members each lie in a given cover member and in
$\operatorname{int}K_{j+1}\setminus K_{j-2}$.
This is possible because the layer misses $K_{j-2}$ and lies in $\operatorname{int}K_{j+1}$.
The union of these finite families covers $X$.
For $x\in\operatorname{int}K_N$, a neighborhood contained in $\operatorname{int}K_N$ can meet only families with $j\le N+2$.
Together these have finitely many members.
We have obtained a locally finite open refinement, proving paracompactness.
\end{proof}

\begin{proposition}[Sections over closed sets extend to open neighborhoods]\label{sh:closed-neighborhood-extension}
Let $X$ be Hausdorff and paracompact, let $S\subset X$ be closed, and let $\cF$ be any sheaf of abelian groups.
Every $s\in\cF(S)$ extends to a section $t\in\cF(W)$ on some open neighborhood $W\supset S$.
Consequently, the natural map
$\varinjlim_{W\supset S}\cF(W)\to\cF(S)$ is an isomorphism.
\end{proposition}
\begin{proof}
By the definition of the inverse-image sheaf, choose ambient open sets $U_i$ covering $S$ and sections $s_i\in\cF(U_i)$ representing $s$ on $S\cap U_i$.
Add $U_0=X\setminus S$ and use Lemma~\ref{sh:paracompact-topology} to take a locally finite closed refinement $\{F_j\}$ of the resulting cover, assigning each $F_j$ to a containing set $U_{i(j)}$.

For $i(j),i(k)\ne0$, put
\[
 E_{jk}=\{x\in U_{i(j)}\cap U_{i(k)}:
              (s_{i(j)})_x=(s_{i(k)})_x\}.
\]
Equality of germs persists on a neighborhood, so $E_{jk}$ is open.
Compatibility on $S$ gives $S\cap U_{i(j)}\cap U_{i(k)}\subset E_{jk}$.
The bad set $B_{jk}=(F_j\cap F_k)\setminus E_{jk}$ is closed and disjoint from $S$.
These bad sets are locally finite: a neighborhood meeting only finitely many $F_j$ meets sets from only finitely many pairs $(j,k)$.
Hence
\[
 B=\bigcup_{i(j)=0}F_j\ \cup\ \bigcup_{j,k}B_{jk}
\]
is closed and disjoint from $S$. Set $W=X\setminus B$.
In the locally finite closed cover $\{F_j\cap W\}$ of $W$, members assigned to $i(j)=0$ have disappeared.
On the remaining members the sections $s_{i(j)}$ have equal germs on pairwise intersections and thus agree there.
Proposition \ref{sh:closed-gluing} glues them to $t\in\cF(W)$.
Its germs along $S$ equal those of $s$, giving the desired extension.

If two neighborhood extensions agree on $S$, they agree on some open neighborhood of each $x\in S$.
The union of these neighborhoods is still a neighborhood of $S$, so the extensions have the same direct-limit class.
This proves injectivity of the natural map.
\end{proof}

\begin{definition}[Support, fine sheaves, soft sheaves, and flabby sheaves]\label{sh:fine-soft-flabby}
The support of a sheaf endomorphism $\phi:\cF\to\cF$ is
$\overline{\{x:\phi_x\ne0\}}$.
A sheaf $\cF$ is \emph{fine} if for every locally finite open cover $\{U_i\}$ there are endomorphisms $\phi_i$ with
$\supp\phi_i\subset U_i$ and $\sum_i\phi_i=\operatorname{id}$; the sum is locally finite.
It is \emph{soft} if every section over a closed subset $S$ extends to $X$.
It is \emph{flabby}, or \emph{flasque}, if every section over an open subset extends to $X$.
Flabbiness is equivalent to surjectivity of every restriction $\cF(V)\to\cF(U)$ for open $U\subset V$:
one direction follows by extending to $X$ and restricting to $V$, and the other by taking $V=X$.
\end{definition}

\begin{lemma}[The precise condition for extension by zero]\label{sh:extension-zero}
If $\phi$ is an endomorphism of $\cF$ with $\supp\phi\subset U$, then for every
$s\in\cF(U)$, the section $\phi(s)$ has a unique global extension vanishing on
$X\setminus\supp\phi$.
\end{lemma}
\begin{proof}
On the open set $U\setminus\supp\phi$, all germs of $\phi(s)$ vanish, so the section is zero.
It is therefore compatible with the zero section on $X\setminus\supp\phi$.
The original section is defined on $U$, and $U$ together with $X\setminus\supp\phi$ covers $X$. The sheaf axiom gives the unique gluing on these two open sets.
The section $s$ itself need not admit extension by zero; the support condition makes the gluing possible.
\end{proof}

\begin{proposition}[Modules over continuous or smooth functions are fine]\label{sh:smooth-fine}
On a Hausdorff paracompact space, every module sheaf over the sheaf of continuous functions is fine.
On a smooth manifold, every module sheaf over the smooth function sheaf $\cA_X^0$ is fine.
In particular, sheaves of smooth sections of vector bundles and of bundle-valued differential forms are fine.
\end{proposition}
\begin{proof}
In the continuous case use the partition of unity from Lemma \ref{sh:paracompact-topology}; in the smooth case use Theorem \ref{bas:partition}.
In either case define $\phi_i(s)=\rho_i s$.
The module axioms and compatibility with restrictions make this a sheaf endomorphism, and
$\sum\rho_i=1$ gives $\sum\phi_i=\operatorname{id}$.
If $x\notin\supp\rho_i$, then $\rho_i$ vanishes on a neighborhood of $x$, so $\phi_i$ vanishes there as well.
Thus $\supp\phi_i\subset\supp\rho_i\subset U_i$.
Only containment of supports is asserted: on the zero module sheaf, for example, every multiplication endomorphism has empty support.
\end{proof}

\begin{proposition}[Fine and flabby sheaves are soft]\label{sh:fine-soft}
On a Hausdorff paracompact space, every fine sheaf and every flabby sheaf is soft.
\end{proposition}
\begin{proof}
Given $s$ over a closed set $S$, Proposition \ref{sh:closed-neighborhood-extension} first extends it to $t$ on an open neighborhood $W$.
For a flabby sheaf, extend once more to $X$.
For a fine sheaf, choose a decomposition of the identity $\phi_W+\phi_0$ subordinate to the two-set cover $\{W,X\setminus S\}$.
Since $\supp\phi_0\subset X\setminus S$, its stalk map is zero on $S$, so $(\phi_W)_x=\operatorname{id}$ there.
Extend $\phi_W(t)$ by zero to $X$ using Lemma~\ref{sh:extension-zero}.
Its germs on $S$ agree with $s$, giving the desired global extension.
\end{proof}

\begin{lemma}[A soft kernel permits global lifting]\label{sh:soft-lifting}
Let $X$ be Hausdorff and paracompact, and suppose
$0\to\cF\xrightarrow a\cG\xrightarrow b\cH\to0$ is exact.
If $\cF$ is soft, then $\cG(X)\to\cH(X)$ is surjective.
\end{lemma}
\begin{proof}
Fix $h\in\cH(X)$. Sheaf surjectivity gives local lifts $g_i\in\cG(U_i)$.
Take a subordinate locally finite closed cover $\{F_j\}_{j\in J}$ and restrict the assigned lift to each $F_j$, calling it $t_j$.
For every $J'\subset J$, the union $S_{J'}=\bigcup_{j\in J'}F_j$ is closed.
Consider pairs $(J',t)$, where $t\in\cG(S_{J'})$ and $b(t)=h|_{S_{J'}}$, ordered by inclusion of indices and extension of sections.

Every chain has an upper bound: take the union of its index sets.
On each $F_j$ restrict a section from a chain member containing $j$.
This is independent of the choice because chain members are comparable, and the restrictions agree on all $F_j\cap F_k$.
Proposition \ref{sh:closed-gluing} glues them over the union.
Zorn's lemma therefore gives a maximal pair $(J',t)$.
If $j\notin J'$, then on the closed set $S_{J'}\cap F_j$, the difference $t_j-t$ lies in $\ker b$.
Proposition~\ref{sh:inverse-exact} preserves exactness after restriction to this closed set, and Proposition~\ref{sh:stalk-exact} identifies the kernel on global sections. Hence the difference is uniquely $a(r)$ for some
$r\in\cF(S_{J'}\cap F_j)$.
Softness extends $r$ to $\widetilde r\in\cF(X)$.
Replace $t_j$ on $F_j$ by $t_j-a(\widetilde r)|_{F_j}$; it now agrees with $t$ on the intersection.
Gluing over the two closed sets extends $t$ to $S_{J'}\cup F_j$, contradicting maximality.
Thus $J'=J$ and $S_{J'}=X$, giving a global lift.

On a second countable manifold the closed cover may be chosen countable, and the same corrections can be made successively along $F_1,F_2,\ldots$.
Each step retains the previous section; locally finite closed gluing completes the construction.
Zorn's lemma is the version of this argument for an arbitrary index set.
\end{proof}

\begin{proposition}[Softness in short exact sequences]\label{sh:soft-exact}
On a Hausdorff paracompact space, if the kernel $\cF$ of a short exact sequence is soft, then the middle sheaf $\cG$ is soft if and only if the quotient $\cH$ is soft.
\end{proposition}
\begin{proof}
Let $S\subset X$ be closed. It remains Hausdorff and paracompact.
The restriction $\cF|_S$ is soft: a closed subset of $S$ is closed in $X$, so its section extends first to $X$ and then restricts to $S$.
The restricted short exact sequence remains exact by Proposition~\ref{sh:inverse-exact}, and Lemma~\ref{sh:soft-lifting} gives surjectivity of $\cG(S)\to\cH(S)$.

If $\cG$ is soft, lift $h\in\cH(S)$ to $g\in\cG(S)$, extend $g$ to $X$, and take its image to extend $h$.
Conversely, suppose $\cH$ is soft and take $g\in\cG(S)$.
Extend $b(g)$ to $\widetilde h\in\cH(X)$, and use the soft-kernel lifting lemma to choose a global lift $\widetilde g$.
The difference $g-\widetilde g|_S$ has zero image, so it comes from some $r\in\cF(S)$.
Extend $r$ to $\widetilde r\in\cF(X)$.
Then $\widetilde g+a(\widetilde r)$ extends $g$ globally.
\end{proof}

\begin{proposition}[Closed discrete support]\label{sh:discrete-support}
If $X$ is Hausdorff and the support of $\cF$ is contained in a closed discrete set $S$, then for every open set $U$,
\[
 \cF(U)\simeq\prod_{x\in S\cap U}\cF_x.
\]
Consequently, $\cF$ is flabby.
\end{proposition}
\begin{proof}
The germ map is injective by Lemma~\ref{sh:germs-detect}, because all stalks outside $S$ vanish.
Given a family of germs on the right, choose for each $x\in S\cap U$ a neighborhood $V_x\subset U$ meeting $S$ only at $x$, and represent the germ on a smaller neighborhood $W_x$.
Every germ of this representative on $W_x\setminus\{x\}$ vanishes, so it is the zero section there.
For $x\ne y$, the intersection $W_x\cap W_y$ misses $S$, and both representatives vanish there.
They therefore glue uniquely, together with the zero section on $U\setminus S$, over an open cover of $U$.
This proves surjectivity.
Restriction forgets some coordinates; inserting zeros in the omitted coordinates gives a lift, proving flabbiness.
\end{proof}

For coherent sheaves, a zero-dimensional analytic support is locally discrete; Corollary~\ref{sh:zero-dimensional-coherent} will deduce flabbiness from the proposition once analytic dimension has been developed.

\section{The Godement Resolution and Cohomology}
A resolution embeds a sheaf into a sequence of sheaves with better extension properties.
The construction below starts by allowing arbitrary families of germs, which always extend by assigning zero at additional points.
Its successive quotients record the failure of those families to come from actual sections.
A cochain complex $C^\bullet$ is a sequence of groups $C^q$ with homomorphisms $d:C^q\to C^{q+1}$ satisfying $d^2=0$.
Its cohomology is $H^q(C^\bullet)=\ker(d:C^q\to C^{q+1})/\operatorname{im}(d:C^{q-1}\to C^q)$; elements of the kernel are cocycles, and elements of the image are boundaries.
A complex morphism is a family of maps commuting with the differentials.
A resolution of a sheaf is an exact augmented complex starting with that sheaf, as in \eqref{sh:godement-resolution} below.


\begin{definition}[Unrestricted families of germs]\label{sh:godement-definition}
For any sheaf of abelian groups $\cF$, define
\[
 G^0(\cF)(U)=\prod_{x\in U}\cF_x.
\]
These families of germs need \emph{not} be locally represented by one section.
Restriction simply forgets coordinates, so $G^0(\cF)$ is a flabby sheaf.
The map $j_\cF:\cF\to G^0(\cF)$ sends a section to its family of germs and is injective by Lemma~\ref{sh:germs-detect}.
Let $Q(\cF)=G^0(\cF)/\cF$ be the quotient sheaf, and recursively define
\[
 Q^0(\cF)=\cF,\quad Q^{r+1}(\cF)=Q(Q^r(\cF)),
 \quad G^r(\cF)=G^0(Q^r(\cF)).
\]
The differential $d:G^r(\cF)\to G^{r+1}(\cF)$ is the quotient map to
$Q^{r+1}(\cF)$ followed by its natural injection.
\end{definition}

\begin{proposition}[Exactness and functoriality of the Godement resolution]\label{sh:godement-exact}
The construction gives a natural exact sequence
\begin{equation}\label{sh:godement-resolution}
0\longrightarrow\cF\longrightarrow G^0(\cF)
 \xrightarrow dG^1(\cF)\xrightarrow dG^2(\cF)\longrightarrow\cdots.
\end{equation}
The functors $G^0$, $Q$, and every $G^r$ are additive and exact.
For every open set $U$, the functor $\cF\mapsto\Gamma(U,G^r(\cF))$ is also exact.
These constructions commute with restriction to open subspaces.
\end{proposition}
\begin{proof}
By definition there are short exact sequences
\[
0\to Q^r(\cF)\to G^r(\cF)\to Q^{r+1}(\cF)\to0.
\]
Splicing consecutive sequences gives \eqref{sh:godement-resolution}, with $d^2=0$.
A sheaf morphism acts stalkwise on families of germs and then on quotient sheaves, giving functoriality and additivity.

Given $0\to\cF'\to\cF\to\cF''\to0$, applying $G^0$ on any $U$ gives the product of the exact stalk sequences.
Choosing a preimage in each coordinate proves surjectivity at the end, so the sequence is exact on every open set.
Exactness of $Q$ follows by an element chase.
At a fixed stalk, consider the natural short exact sequences
$0\to\cF_x\to G^0(\cF)_x\to Q(\cF)_x\to0$
and the corresponding sequences with primes and double primes.
If an element of $G^0(\cF')_x$ comes from $\cF_x$ after mapping into $G^0(\cF)_x$, that element of $\cF_x$ has zero image in $\cF''_x$ and therefore comes from $\cF'_x$.
This proves injectivity of $Q(\cF')_x\to Q(\cF)_x$.
If a class in $Q(\cF)_x$ maps to zero, its representative has image in $G^0(\cF'')_x$ coming from $\cF''_x$.
Lift that element to $\cF_x$ and subtract it from the representative.
The new representative maps to zero in $G^0(\cF'')_x$, hence comes from $G^0(\cF')_x$.
This proves middle exactness.
Surjectivity at the end follows from surjectivity of $G^0(\cF)_x\to G^0(\cF'')_x$.
Thus $Q$ is exact, and induction gives exactness of $Q^r$ and $G^r$.

Finally,
$\Gamma(U,G^r(\cF))=\prod_{x\in U}(Q^r(\cF))_x$,
so the same product argument proves exactness of this functor.
On an open subspace the stalks and local restrictions are unchanged, and quotient sheaves also commute with open restriction.
All the constructions therefore commute with open restriction.
\end{proof}

\begin{definition}[Sheaf cohomology]\label{sh:cohomology-definition}
Define
\[
 H^q(X,\cF)=H^q(\Gamma(X,G^\bullet(\cF)))
 =\frac{\ker(\Gamma(X,G^q(\cF))\to\Gamma(X,G^{q+1}(\cF)))}
 {\operatorname{im}(\Gamma(X,G^{q-1}(\cF))\to\Gamma(X,G^q(\cF)))}.
\]
For $q=0$, the denominator is zero.
Proposition~\ref{sh:stalk-exact} gives $H^0(X,\cF)=\Gamma(X,\cF)$.
If $H^q(X,\cF)=0$ for every $q>0$, then $\cF$ is \emph{acyclic for global sections}, or simply \emph{acyclic}.
This terminology does not require $H^0$ to vanish.
\end{definition}

\begin{proposition}[The long exact cohomology sequence]\label{sh:long-exact}
The short exact sequence \eqref{sh:ses} naturally induces a long exact sequence
\[
0\to H^0(X,\cF)\to H^0(X,\cG)\to H^0(X,\cH)
 \xrightarrow\delta H^1(X,\cF)\to\cdots.
\]
Its general segment is
\[
H^q(X,\cF)\to H^q(X,\cG)\to H^q(X,\cH)
 \xrightarrow\delta H^{q+1}(X,\cF).
\]
\end{proposition}
\begin{proof}
Put $A^q=\Gamma(X,G^q(\cF))$, $B^q=\Gamma(X,G^q(\cG))$, and $C^q=\Gamma(X,G^q(\cH))$.
Proposition~\ref{sh:godement-exact} gives the degreewise short exact sequence
$0\to A^\bullet\xrightarrow iB^\bullet\xrightarrow\pi C^\bullet\to0$; here $i,\pi$ are the induced maps, not coordinate indices.
For a closed $c\in C^q$, choose $b\in B^q$ with $\pi b=c$.
Then $\pi(db)=dc=0$, so there is a unique $a\in A^{q+1}$ with $ia=db$.
Since $i(da)=d^2b=0$ and $i$ is injective, $a$ is closed.
Define $\delta[c]=[a]$.
Changing the lift adds $ia'$ to $b$ and therefore adds $da'$ to $a$.
Replacing $c$ by $c+dc'$ and lifting $c'$ allows us to replace $b$ by $b+db'$, leaving $a$ unchanged.
The connecting homomorphism is thus well defined.

Check exactness at each position.
If a closed $b$ has image $\pi b=dc'$, lift $c'$ to $b'$; then $b-db'$ comes from a closed element of $A^q$.
This proves exactness at $H^q(B)$.
If $\delta[c]=0$, then the element above satisfies $a=da'$, so $b-ia'$ is closed and still lifts $c$.
This proves exactness at $H^q(C)$.
If a closed $a\in A^{q+1}$ satisfies $ia=db$, then $\pi b$ is closed and
$\delta[\pi b]=[a]$.
This proves exactness at $H^{q+1}(A)$.
Consecutive composites vanish directly by construction, and left exactness supplies the initial degree-zero segment.
\end{proof}

\begin{lemma}[A flabby kernel permits lifting of sections]\label{sh:flabby-lifting}
If $\cF$ in \eqref{sh:ses} is flabby, then $\cG(U)\to\cH(U)$ is surjective for every open set $U$.
If $\cG$ is also flabby, then $\cH$ is flabby.
\end{lemma}
\begin{proof}
Given $s\in\cH(U)$, consider all local lifts $(V,t)$, where $V\subset U$ is open,
$t\in\cG(V)$, and $b(t)=s|_V$, ordered by extension.
Sections in a chain agree and hence glue over the union of their open domains, giving an upper bound.
Zorn's lemma yields a maximal local lift $(V,t)$.
If $x\in U\setminus V$, sheaf surjectivity gives a lift $t'\in\cG(W)$ near $x$.
On $V\cap W$, the difference $t'-t$ comes uniquely from $r\in\cF(V\cap W)$.
Extend $r$ to $W$ and subtract its image from $t'$.
The resulting lift glues with $t$ on $V\cup W$, contradicting maximality.
Thus $V=U$.
On a second countable space one may instead choose a countable lifting cover, glue its members successively by this correction, and take the union of the increasing open domains, avoiding maximal elements altogether.

If $\cG$ is flabby, lift $s\in\cH(U)$ to $\cG(U)$, extend to $\cG(X)$, and project to $\cH(X)$.
The result extends $s$ globally.
\end{proof}

A sheaf $\cI$ is \emph{injective} if every morphism $\cA\to\cI$ extends across every injection $\cA\hookrightarrow\mathcal B$.
This extension property concerns the source of the injection; it is separate from exactness of global sections.

\begin{proposition}[Direct summands of flabby sheaves and injective sheaves]\label{sh:flabby-summand}
Every direct summand of a flabby sheaf is flabby.
Every injective sheaf of abelian groups is flabby.
\end{proposition}
\begin{proof}
Suppose $\mathcal{B}=\cF\oplus\mathcal{K}$ is flabby.
For $s\in\cF(U)$, view $(s,0)$ as a section of $\mathcal{B}(U)$, extend it to $\mathcal{B}(X)$, and project to $\cF(X)$.
Projection commutes with restriction, so the result extends $s$.

Apply the defining extension property to $\cI\hookrightarrow G^0(\cI)$ and the identity of $\cI$.
It gives a retraction $r:G^0(\cI)\to\cI$ with $r\circ j=\operatorname{id}$.
Consequently, $G^0(\cI)=j(\cI)\oplus\ker r$: each local section $b$ is uniquely
$j(r(b))+(b-j(r(b)))$.
Since $G^0(\cI)$ is flabby, the first paragraph proves that $\cI$ is flabby.
The splitting here concerns an injection whose \emph{source is injective}; arbitrary short exact sequences need not split.
\end{proof}

\begin{proposition}[Short exact sequences with flabby kernel]\label{sh:flabby-exact}
If the kernel $\cF$ in $0\to\cF\to\cG\to\cH\to0$ is flabby, then $\cG$ is flabby if and only if $\cH$ is flabby.
\end{proposition}
\begin{proof}
If $\cG$ is flabby, lift $h\in\cH(U)$ to $g\in\cG(U)$ by Lemma \ref{sh:flabby-lifting}, extend $g$ to $X$, and take its image.
Conversely, suppose $\cH$ is flabby and take $g\in\cG(U)$.
Extend its image to $h\in\cH(X)$, then use the same lemma to choose a global lift $g'\in\cG(X)$.
The difference $g-g'|_U$ comes from $r\in\cF(U)$.
Extend $r$ to $\widetilde r\in\cF(X)$; then $g'+\widetilde r$, with the second summand included through the kernel map, extends $g$.
\end{proof}

\begin{theorem}[Acyclicity of flabby, soft, and fine sheaves]\label{sh:acyclic}
On any topological space, every flabby sheaf is acyclic.
On a Hausdorff paracompact space, every soft sheaf, and hence every fine sheaf, is acyclic.
\end{theorem}
\begin{proof}
Suppose $\cF$ is flabby. Since $G^0(\cF)$ is flabby, Proposition \ref{sh:flabby-exact} makes its quotient $Q(\cF)$ flabby.
Lemma \ref{sh:flabby-lifting} also gives surjectivity of
\[
 \Gamma(X,G^0(\cF))\longrightarrow\Gamma(X,Q(\cF)).
\]
Repeating the argument shows that every $Q^r(\cF)$ is flabby and every map
$\Gamma(X,G^{r-1}(\cF))\to\Gamma(X,Q^r(\cF))$ is surjective.

If $u\in\Gamma(X,G^r(\cF))$ satisfies $du=0$, exactness of the resolution and left exactness of sections show that $u$ comes uniquely from
$v\in\Gamma(X,Q^r(\cF))$.
For $r>0$, the preceding surjection lifts $v$ to
$w\in\Gamma(X,G^{r-1}(\cF))$.
The definition of the differential gives $dw=u$.
Thus every positive-degree cocycle is a boundary.

Now let $X$ be Hausdorff and paracompact and $\cF$ soft.
Every $G^r$ is flabby and hence soft by Proposition \ref{sh:fine-soft}.
In $0\to\cF\to G^0(\cF)\to Q(\cF)\to0$, Proposition \ref{sh:soft-exact} makes $Q(\cF)$ soft, and Lemma \ref{sh:soft-lifting} gives the same surjection on global sections.
Induction makes every $Q^r$ soft and all quotient maps surjective on global sections.
The cocycle-lifting argument of the preceding paragraph then proves acyclicity.
Finally, Proposition~\ref{sh:fine-soft} makes fine sheaves soft, so they are acyclic.
\end{proof}

\begin{remark}[Relation to injective resolutions]\label{sh:injective}
Proposition \ref{sh:flabby-summand} proves that injective sheaves are flabby, hence acyclic by Theorem \ref{sh:acyclic}.
Theorem~\ref{sh:resolution-theorem} will show that any injective resolution computes the same groups as the Godement definition; its proof uses the acyclicity just proved.
The existence of enough injective objects is therefore not needed as a prerequisite for the Dolbeault theorem.
\end{remark}

\section{\Cech{} Cohomology: Writing Down Gluing Obstructions}

\begin{definition}[The \Cech{} complex]\label{sh:cech-definition}
For an open cover $\mathfrak U=\{U_i\}_{i\in I}$, write
$U_{i_0\dots i_q}=U_{i_0}\cap\cdots\cap U_{i_q}$.
The group $\check C^q(\mathfrak U,\cF)$ consists of alternating families of sections
$c_{i_0\dots i_q}\in\cF(U_{i_0\dots i_q})$:
interchanging two indices changes the sign, and components with repeated indices are zero.
Define
\[
 (\delta c)_{i_0\dots i_{q+1}}
 =\sum_{a=0}^{q+1}(-1)^a
 c_{i_0\dots\widehat{i_a}\dots i_{q+1}}|_{U_{i_0\dots i_{q+1}}}.
\]
In $\delta^2$, the term obtained by omitting positions $a,b$ occurs twice, with signs
$(-1)^{a+b-1}$ and $(-1)^{a+b}$, so $\delta^2=0$.
Write its cohomology as $\check H^q(\mathfrak U,\cF)$.
The sheaf axioms directly give $\check H^0(\mathfrak U,\cF)=\Gamma(X,\cF)$.
\end{definition}

\begin{proposition}[The concrete meaning of the first obstruction]\label{sh:cech-obstruction}
In \eqref{sh:ses}, let $s\in\cH(X)$ and choose lifts $t_i\in\cG(U_i)$ on an open cover.
Write uniquely $t_j-t_i=a(c_{ij})$.
Then $c=(c_{ij})$ is a \Cech{} one-cocycle with values in $\cF$.
Its class is independent of the lifts and vanishes if and only if $s$ has a global lift.
\end{proposition}
\begin{proof}
On a triple intersection,
$a(c_{jk}-c_{ik}+c_{ij})=(t_k-t_j)-(t_k-t_i)+(t_j-t_i)=0$.
Injectivity of $a$ gives $\delta c=0$.
Replacing $t_i$ by $t_i+a(v_i)$ changes the cocycle by $v_j-v_i=\delta v$, preserving its class.
If $c=\delta v$, then $t_i-a(v_i)$ agree on overlaps and glue to a global lift.
Conversely, a global lift $t$ gives $t_i-t|_{U_i}=a(v_i)$, so
$c_{ij}=v_j-v_i$ and the class is zero.
\end{proof}

\begin{proposition}[An explicit \Cech{} contraction for fine sheaves]\label{sh:cech-contraction}
If $\cF$ is fine and $\mathfrak U$ is a locally finite open cover, then
$\check H^q(\mathfrak U,\cF)=0$ for $q>0$.
\end{proposition}
\begin{proof}
Choose a subordinate decomposition of the identity $\{\phi_j\}$.
For $q\ge1$, define
\[
 (hc)_{i_0\dots i_{q-1}}
 =\sum_j\phi_j(c_{j i_0\dots i_{q-1}}),
\]
where each summand is first extended by zero from $U_j\cap U_{i_0\dots i_{q-1}}$ to
$U_{i_0\dots i_{q-1}}$.
Lemma \ref{sh:extension-zero} justifies the extension, and the sum is locally finite.
For a fixed component $(i_0,\dots,i_q)$,
\[
\begin{aligned}
 (h\delta c)_{i_0\dots i_q}
 &=\sum_j\phi_j\left(c_{i_0\dots i_q}
   +\sum_{a=0}^q(-1)^{a+1}c_{j i_0\dots\widehat{i_a}\dots i_q}\right)\\
 &=c_{i_0\dots i_q}-(\delta hc)_{i_0\dots i_q}.
\end{aligned}
\]
Thus $h\delta+\delta h=\operatorname{id}$.
If $\delta c=0$ and $q>0$, then $c=\delta(hc)$, proving the assertion.
The calculation identifies exactly what a partition of unity removes: the \Cech{} obstruction to gluing.
\end{proof}

\begin{remark}[A fixed cover and sheaf cohomology]\label{sh:cech-caution}
The group $\check H^q(\mathfrak U,\cF)$ for a fixed cover cannot replace
$H^q(X,\cF)$ without hypotheses.
For example, the one-member cover $\{X\}$ has zero \Cech{} cohomology in every positive degree, while sheaf cohomology on $X$ may be nonzero.
One can usually take a limit over refinements of covers.
We will not need that general comparison theory; instead, we prove a comparison theorem for a fixed cover under explicit assumptions.
\end{remark}

\section{Why Resolutions Compute the Same Cohomology}
The comparison argument uses two indices: one for the resolution being studied and one for its Godement resolution.
The elementary elimination lemma below shows that either order of computation gives the same answer.
In each total degree only finitely many components occur, so the proof consists of a finite sequence of subtractions of boundaries.

\begin{lemma}[Elimination in a double complex]\label{sh:double-complex}
Let $B^{a,b}$ be a first-quadrant double complex of abelian groups, meaning $B^{a,b}=0$ unless $a,b\ge0$, with horizontal and vertical differentials
$d_h,d_v$ satisfying $d_h^2=d_v^2=0$ and $d_hd_v=d_vd_h$.
Its total complex is
\[
 \operatorname{Tot}^N B=\bigoplus_{a+b=N}B^{a,b},
 \qquad D|_{B^{a,b}}=d_h+(-1)^a d_v.
\]
If each horizontal row has a compatible exact augmentation
$0\to V^b\to B^{0,b}\to B^{1,b}\to\cdots$,
then $V^\bullet\to\operatorname{Tot}B$ induces an isomorphism in cohomology.
The analogous assertion holds for compatible exact augmentations of all vertical columns.
\end{lemma}
\begin{proof}
The mixed terms have opposite signs, so $D^2=0$.
Consider horizontal rows first. Write a total-degree-$N$ cocycle as
$c=c^{0,N}+\cdots+c^{N,0}$.
If its rightmost nonzero component has $a=r>0$, the $(r+1,N-r)$ component of $Dc=0$ gives
$d_hc^{r,N-r}=0$.
Row exactness provides $c^{r,N-r}=d_hb^{r-1,N-r}$.
Subtract $Db^{r-1,N-r}$ from $c$.
This removes the rightmost component and changes only the component immediately to its left.
After finitely many repetitions, the class is represented by $c'$ supported in $B^{0,N}$.
Since $d_hc'=0$, it comes from $V^N$, and $Dc'=0$ makes that element closed in $V^\bullet$.
This proves surjectivity.

If a cocycle from $V^N$ equals $Db$, apply the same elimination to the rightmost nonzero component of $b$.
Because $Db$ vanishes at every position with $a>0$, this component is horizontally closed.
Subtracting a $D$-boundary from $b$ removes it without changing $Db$.
Eventually $b$ is supported in $B^{0,N-1}$ and is horizontally closed, hence comes from $V^{N-1}$.
The original cocycle is already a boundary in $V^\bullet$, proving injectivity.
For $N=0$ there are no components to eliminate; the assertion is exactness of the augmentation.

For columns, interchange $a,b$.
The resulting total differential is $d_v+(-1)^b d_h$.
Multiplication by $(-1)^{ab}$ on $B^{a,b}$ intertwines this differential with the original one.
The row argument therefore applies to columns as well.
\end{proof}

\begin{proposition}[The natural comparison map for an arbitrary resolution]\label{sh:general-resolution-map}
For any exact resolution $0\to\cF\to\mathcal L^0\to\mathcal L^1\to\cdots$, even without acyclicity assumptions on the $\mathcal L^q$, there is a natural homomorphism
\[
 H^q(\Gamma(X,\mathcal L^\bullet))\longrightarrow H^q(X,\cF).
\]
In degree zero it is the identity on $\Gamma(X,\cF)$.
\end{proposition}
\begin{proof}
Set $B^{a,b}=\Gamma(X,G^a(\mathcal L^b))$, with the Godement differential horizontal and the given resolution differential vertical.
Use total differential $D=d_h+(-1)^a d_v$.
The natural inclusions $\mathcal L^b\to G^0(\mathcal L^b)$ map
$\Gamma(X,\mathcal L^\bullet)$ into column zero of the total complex.
This is a complex morphism: its horizontal differential vanishes, and at $a=0$ the vertical sign is positive.

On the other hand, $\Gamma(X,G^a(-))$ is exact, so each vertical column has exact augmentation
$0\to\Gamma(X,G^a(\cF))\to B^{a,0}\to B^{a,1}\to\cdots$.
Lemma~\ref{sh:double-complex} gives a natural isomorphism
$H^q(X,\cF)\to H^q(\operatorname{Tot}B)$.
Compose the map in cohomology induced by the first complex morphism with the inverse of this isomorphism.
This is the desired homomorphism.
The entire construction is functorial in resolutions and their morphisms, hence natural.
In degree zero it simply embeds the same global section in two natural ways, so it is the identity.
\end{proof}

\begin{theorem}[Comparison with an acyclic resolution]\label{sh:resolution-theorem}
Suppose
\[
0\to\cF\to\mathcal L^0\xrightarrow{d_{\mathcal L}}\mathcal L^1
 \xrightarrow{d_{\mathcal L}}\cdots
\]
is an exact sequence of sheaves of abelian groups, with each $\mathcal L^b$ acyclic.
Then there is a natural isomorphism
\[
 H^q(X,\cF)\simeq H^q(\Gamma(X,\mathcal L^\bullet)).
\]
In particular, fine resolutions on smooth manifolds compute sheaf cohomology.
\end{theorem}
\begin{proof}
Take $B^{a,b}=\Gamma(X,G^a(\mathcal L^b))$, with the Godement differential horizontal and the differential induced by $d_{\mathcal L}$ vertical.
Functoriality makes them commute.
For fixed $b$, acyclicity of $\mathcal L^b$ gives exactness of
\[
0\to\Gamma(X,\mathcal L^b)\to B^{0,b}\to B^{1,b}\to\cdots.
\]
Lemma~\ref{sh:double-complex} identifies total cohomology with the cohomology of the global-section complex in the statement.
For fixed $a$, the exact functor $\Gamma(X,G^a(-))$ of Proposition \ref{sh:godement-exact} turns the given resolution into the exact sequence
\[
0\to\Gamma(X,G^a(\cF))\to B^{a,0}\to B^{a,1}\to\cdots.
\]
A second application of Lemma~\ref{sh:double-complex} identifies total cohomology with
\[H^q(\Gamma(X,G^\bullet(\cF)))=H^q(X,\cF).\]
Both augmentations are natural, so the induced isomorphisms and their inverses are natural.
The assertion for fine resolutions follows from Theorem \ref{sh:acyclic}.
\end{proof}

\begin{theorem}[\Cech{} comparison for covers with acyclic intersections]\label{sh:leray}
Suppose an open cover $\mathfrak U$ satisfies, on every nonempty finite intersection $U_{i_0\dots i_a}$,
\[
 H^b(U_{i_0\dots i_a},\cF|_{U_{i_0\dots i_a}})=0\quad(b>0).
\]
Then there is a natural isomorphism
$\check H^q(\mathfrak U,\cF)\simeq H^q(X,\cF)$.
\end{theorem}
\begin{proof}
Set $B^{a,b}=\check C^a(\mathfrak U,G^b(\cF))$, with the \Cech{} differential horizontal and the Godement differential vertical.
First prove exactness of each horizontal row augmented by $\Gamma(X,G^b(\cF))$.
For each point $x$, choose an index $i(x)$ whose cover member contains $x$.
Since sections of $G^b(\cF)=G^0(Q^b(\cF))$ are arbitrary families of germs, define pointwise
\[
 (hc)_{i_0\dots i_{a-1}}(x)
 =c_{i(x),i_0,\dots,i_{a-1}}(x).
\]
This is a section of $G^b(\cF)$; the choice $i(x)$ need be neither continuous nor locally constant.
The termwise calculation of Proposition \ref{sh:cech-contraction} gives
$h\delta+\delta h=\operatorname{id}$.
The degree-zero kernel is the space of global sections by the sheaf axiom.
Lemma~\ref{sh:double-complex} applied to these rows identifies total cohomology with $H^q(X,\cF)$.

For fixed $a$, the column is a product of Godement complexes on finite intersections, using one ordering of the distinct indices.
The Godement construction commutes with open restriction, so the hypothesis makes each such complex exact in positive degrees, with degree-zero kernel
$\cF(U_{i_0\dots i_a})$.
Products remain exact by choosing preimages coordinatewise.
Thus each column has exact augmentation
$0\to\check C^a(\mathfrak U,\cF)\to B^{a,0}\to B^{a,1}\to\cdots$.
The column version of Lemma~\ref{sh:double-complex} identifies total cohomology with $\check H^q(\mathfrak U,\cF)$, completing the proof.
\end{proof}

\begin{remark}\label{sh:no-circularity}
The intersection hypothesis in this comparison theorem requires a separate proof.
Local solvability of the $\dbar$ equation near each point gives exactness on stalks; by itself it does not imply vanishing of higher cohomology on a fixed open set.
The Dolbeault theorem below uses the fine-resolution comparison theorem \ref{sh:resolution-theorem}, so no coordinate cover needs first to be shown to satisfy the special hypothesis above.
\end{remark}

\section{A One-Variable Integral Operator and the Local Dolbeault Lemma}

\begin{lemma}[The one-variable Cauchy operator]\label{sh:cauchy-operator}
For $f\in C_c^\infty(\C)$, write $dA=dx\,dy$ for plane area measure and define
\[
 (Tf)(z)=\frac1\pi\int_{\C}\frac{f(\zeta)}{z-\zeta}\,dA(\zeta).
\]
Then $Tf$ is smooth and $\partial(Tf)/\partial\bar z=f$.
If $f$ also depends smoothly on other parameters, with support in the integration variable locally uniformly contained in a compact set, then $T$ commutes with every partial derivative in those parameters.
\end{lemma}
\begin{proof}
The kernel $|z-\zeta|^{-1}$ is locally integrable in the plane because
$\int_{|w|<\epsilon}|w|^{-1}dA(w)=2\pi\epsilon$.
Substitute $\zeta=z+w$ to obtain
\[
 (Tf)(z)=-\frac1\pi\int_{\C}\frac{f(z+w)}{w}\,dA(w).
\]
As $z$ ranges over a fixed compact set, all nonzero integrands are supported in one bounded region.
Derivatives of any order act only on the smooth factor $f(z+w)$ and remain bounded by a constant times $|w|^{-1}$.
Repeated differentiation under the integral is therefore valid, proving smoothness and
$\partial_{\bar z}Tf=T(\partial_{\bar z}f)$.
With parameters present, the same locally uniform bounds justify commuting their derivatives with $T$.

Fix $z$. Take a large disc containing both $\supp f$ and $z$, and remove the disc of radius $\epsilon$ centered at $z$.
On the punctured region,
\[
 \partial_{\bar\zeta}\!\left(\frac{f(\zeta)}{\zeta-z}\right)
 =\frac{\partial_{\bar\zeta}f(\zeta)}{\zeta-z}.
\]
Since $d(g\,d\zeta)=2\ii\,\partial_{\bar\zeta}g\,dA$, the boundary formula in Proposition~\ref{bas:stokes} applies.
The outer boundary integral vanishes and the inner boundary is oriented clockwise, giving
\[
 \int_{\C}\frac{\partial_{\bar\zeta}f(\zeta)}{\zeta-z}\,dA(\zeta)
 =\lim_{\epsilon\to0}\frac{-1}{2\ii}
   \int_{|\zeta-z|=\epsilon}^{\rm counterclockwise}
      \frac{f(\zeta)}{\zeta-z}\,d\zeta
 =-\pi f(z).
\]
The last limit follows by writing $\zeta=z+\epsilon e^{\ii t}$ and using continuity of $f$.
The omitted small-disc integral tends to zero by local integrability of the kernel.
The negative sign in the definition of $T$ with denominator $\zeta-z$ now gives $\partial_{\bar z}Tf=f$.
\end{proof}

\begin{theorem}[The local Dolbeault--Grothendieck lemma]\label{sh:dolbeault-lemma}
Let $P=D_1\times\cdots\times D_n\subset\C^n$ be an open polydisc, and let
$P'=D_1'\times\cdots\times D_n'$ with $\overline{D_j'}\subset D_j$.
If $\alpha$ is a smooth $(p,q)$-form on $P$, $q>0$, and $\dbar\alpha=0$, then there is a smooth $(p,q-1)$-form $u$ on $P'$ satisfying $\dbar u=\alpha|_{P'}$.
\end{theorem}
\begin{proof}
First take $p=0$.
We eliminate $d\bar z_n,d\bar z_{n-1},\dots,d\bar z_1$ successively.
At step $k$, suppose the closed form $\beta$ is defined on
$D_1\times\cdots\times D_k\times D_{k+1}'\times\cdots\times D_n'$
and contains none of $d\bar z_{k+1},\dots,d\bar z_n$.
Write uniquely
\[
 \beta=d\bar z_k\wedge g+h,
\]
where $g,h$ contain only $d\bar z_1,\dots,d\bar z_{k-1}$.
Comparing coefficients containing $d\bar z_j$, $j>k$, in $\dbar\beta=0$ shows that every coefficient of $g,h$ has zero derivative with respect to $\bar z_j$.

Choose $\chi_k\in C_c^\infty(D_k)$ equal to one near $\overline{D_k'}$.
Extend each coefficient of $\chi_k g$ by zero to $\C$ in variable $k$, then apply the operator $T_k$ of Lemma \ref{sh:cauchy-operator}.
Call the resulting $(0,q-1)$-form $v_k=T_k(\chi_k g)$.
On the product with its $k$th disc now shrunk, it satisfies
\[
 \partial_{\bar z_k}v_k=g,\qquad
 \partial_{\bar z_j}v_k=0\quad(j>k).
\]
The second set of equations uses commutation with parameter derivatives; the cutoff depends only on $z_k$.
The part of $\dbar v_k$ containing $d\bar z_k$ is exactly $d\bar z_k\wedge g$, and no higher-index $d\bar z_j$ occurs.
Thus $\beta-\dbar v_k$ remains closed and contains none of $d\bar z_k,\dots,d\bar z_n$.
This completes one elimination step while shrinking only the $k$th disc.

Start at $k=n$ and repeat finitely many times.
The remaining form has positive antiholomorphic degree but contains no $d\bar z_j$, so it is zero.
Restrict all $v_k$ to the common domain $P'$ and add them to obtain $u$.
Every domain and shrinkage has been specified; the construction makes no claim of global solvability on an arbitrary open set.

For general $p$, write $\alpha=\sum_{|I|=p}dz_I\wedge\alpha_I$.
Since
$\dbar(dz_I\wedge\alpha_I)=(-1)^p dz_I\wedge\dbar\alpha_I$,
closedness is equivalent to closedness of each $\alpha_I$.
Choose $\dbar v_I=\alpha_I$ by the case just proved, and set
$u=(-1)^p\sum_I dz_I\wedge v_I$.
The same sign formula gives $\dbar u=\alpha$.
\end{proof}

\begin{remark}
The primitive $u$ has antiholomorphic degree $q-1$, not $q$.
For $q=0$ the conclusion is different: the kernel of $\dbar$ consists of $(p,0)$-forms with holomorphic coefficients, rather than the image of forms of a lower degree.
The proof also shows that if the original coefficients are holomorphic in additional complex parameters and the cutoffs are independent of them, the constructed primitive remains holomorphic in those parameters.
\end{remark}

\begin{corollary}[A primitive on a larger domain and the bundle-valued version]\label{sh:bundle-local-dolbeault}
In Theorem \ref{sh:dolbeault-lemma}, $u$ may be chosen smooth on all of $P$, provided the equality $\dbar u=\alpha$ is required only on the prescribed $P'\Subset P$.
The same assertion holds for closed forms with values in a holomorphic vector bundle within a coordinate neighborhood carrying a holomorphic trivialization.
\end{corollary}
\begin{proof}
Choose intermediate polydiscs
$\overline{P'}\subset P_1\Subset P_2\Subset P$.
The local lemma gives a primitive $v$ on $P_2$.
Choose a smooth cutoff $\chi$ supported in $P_2$ and equal to one on $P_1$ by Theorem~\ref{bas:partition}.
Extend $\chi v$ by zero to $P$ to obtain a smooth form $u$.
On $P'$, we have $\dbar\chi=0$ and $\chi=1$, so $\dbar u=\dbar v=\alpha$.
The same equation need not hold on $P\setminus P_1$.

For a holomorphic frame $e_1,\ldots,e_r$ of $E$, write $\alpha=\sum\alpha_j e_j$.
Since $\dbar_Ee_j=0$, closedness is equivalent to $\dbar\alpha_j=0$ for each component.
Apply the integral construction to the finitely many components on the same polydisc and set $u=\sum u_j e_j$.
Then $\dbar_Eu=\alpha$ on $P'$.
\end{proof}

\section{The Dolbeault Theorem and Its Scope}

\begin{definition}\label{sh:dolbeault-complex}
Let $E$ be a holomorphic vector bundle of finite rank over a complex manifold $X$.
Write $\cA^{p,q}(E)$ for the sheaf of smooth $E$-valued $(p,q)$-forms and $A^{p,q}(X,E)$ for its global sections.
In a local holomorphic frame, apply $\dbar$ to each coefficient to define
$\dbar_E:\cA^{p,q}(E)\to\cA^{p,q+1}(E)$.
Transition matrices are holomorphic, so $\dbar_E$ is compatible with frame changes and satisfies $\dbar_E^2=0$.
Define Dolbeault cohomology by
\[
 H_{\dbar}^{p,q}(X,E)
 =\frac{\ker(\dbar_E:A^{p,q}(X,E)\to A^{p,q+1}(X,E))}
 {\operatorname{im}(\dbar_E:A^{p,q-1}(X,E)\to A^{p,q}(X,E))}.
\]
For $q=0$, the denominator is zero.
When $E$ is the trivial line bundle, omit it from the notation.
\end{definition}

\begin{proposition}[The fine Dolbeault resolution]\label{sh:dolbeault-resolution}
Let $\Omega_X^p$ be the sheaf of holomorphic $p$-forms.
There is an exact sequence
\[
0\to\Omega_X^p\otimes_{\cO_X}\cO(E)
 \to\cA^{p,0}(E)\xrightarrow{\dbar_E}\cA^{p,1}(E)
 \xrightarrow{\dbar_E}\cdots\xrightarrow{\dbar_E}\cA^{p,n}(E)\to0,
\]
and every $\cA^{p,q}(E)$ is fine.
\end{proposition}
\begin{proof}
In local holomorphic coordinates and a holomorphic frame, closedness in degree $q=0$ means that each coefficient satisfies
$\partial f/\partial\bar z_j=0$, hence is holomorphic by Theorem \ref{bas:cauchy}.
This identifies the degree-zero kernel with $\Omega_X^p\otimes\cO(E)$.
For $q>0$, apply Theorem \ref{sh:dolbeault-lemma} componentwise to obtain a primitive on a smaller neighborhood of every point.
The complex is therefore exact on stalks and hence, by Proposition \ref{sh:stalk-exact}, exact as a sequence of sheaves.
Finally, each form sheaf is a module over smooth functions, so it is fine by Proposition \ref{sh:smooth-fine}.
\end{proof}

\begin{theorem}[The Dolbeault isomorphism]\label{sh:dolbeault}
For every Hausdorff second countable complex manifold $X$ and every holomorphic vector bundle $E$ of finite rank, there is a natural isomorphism
\[
 \boxed{\quad
 H^q(X,\Omega_X^p\otimes_{\cO_X}\cO(E))
 \simeq H_{\dbar}^{p,q}(X,E).
 \quad}
\]
It is compatible with maps induced by holomorphic vector-bundle morphisms.
\end{theorem}
\begin{proof}
Proposition \ref{sh:dolbeault-resolution} gives a fine resolution of the sheaf in question.
Theorem \ref{sh:acyclic} proves that its fine sheaves are acyclic.
Theorem \ref{sh:resolution-theorem} therefore identifies sheaf cohomology with the cohomology of the global-section complex, which is by definition $H_{\dbar}^{p,q}(X,E)$.
A holomorphic bundle morphism $E\to F$ has a holomorphic matrix in local holomorphic frames.
It commutes with $\dbar$ and thus defines a morphism of fine resolutions.
Naturality of the comparison theorem gives the claimed compatibility.
\end{proof}

\begin{corollary}[An exact criterion for global solvability]\label{sh:global-obstruction}
For $q>0$, a global $\dbar_E$-closed $(p,q)$-form $\alpha$ has a global smooth primitive if and only if its corresponding class in
$H^q(X,\Omega_X^p\otimes\cO(E))$ vanishes.
In particular, vanishing of this group implies global solvability for every such equation.
\end{corollary}
\begin{proof}
By definition, $[\alpha]=0$ in the Dolbeault quotient means that some global
$u\in A^{p,q-1}(X,E)$ satisfies $\alpha=\dbar_Eu$.
Theorem~\ref{sh:dolbeault} identifies the zero classes, proving equivalence.
\end{proof}

\begin{remark}[Fineness does not make every global form complex exact]\label{sh:fine-not-exact}
Acyclicity of each individual sheaf $\cA^{p,q}(E)$ does not imply that the cohomology of the complex $A^{p,\bullet}(X,E)$ vanishes.
The former concerns the \emph{Godement resolution of one sheaf}; the latter concerns \emph{a collection of different sheaves connected by $\dbar$}.
The comparison theorem says that the latter computes cohomology of the initial holomorphic sheaf, which is usually not fine.
The next example exhibits a nonzero obstruction explicitly.
\end{remark}

\section{A Complete Example: The Square Complex Torus}

\begin{proposition}\label{sh:torus}
Let $T=\C/(\Z+\ii\Z)$, the complex torus of Proposition~\ref{bas:torus}. On its covering space use $z=x+\ii y$; functions on $T$ have doubly periodic coefficients, and $d\bar z$ descends to $T$.
Then
\[
 H^0(T,\cO_T)=\C,
 \qquad H^1(T,\cO_T)=\C[d\bar z].
\]
More explicitly, $a(x,y)d\bar z$ is $\dbar$-exact if and only if
$\int_0^1\int_0^1a(x,y)\,dx\,dy=0$.
\end{proposition}
\begin{proof}
In complex dimension one, every $(0,1)$-form is automatically closed.
If $a=\tfrac12(f_x+\ii f_y)$ with $f$ doubly periodic, integration cancels the boundary terms, so the mean of $a$ is zero.
This proves necessity and shows that $d\bar z$ is not globally $\dbar$-exact.

For sufficiency, define its Fourier coefficients by
$a_{mn}=\int_0^1\int_0^1a(x,y)e^{-2\pi\ii(mx+ny)}dx\,dy$ for $(m,n)\in\Z^2$.
Consider the Fourier series
\[
 a(x,y)=\sum_{(m,n)\in\Z^2}a_{mn}e^{2\pi\ii(mx+ny)}.
\]
Repeated integration by parts shows that for every positive integer $N$,
$|a_{mn}|\le C_N(1+|m|+|n|)^{-N}$.
For example, apply $(1-\partial_x^2-\partial_y^2)^N$ to $a$ and take Fourier coefficients.
Boundedness of the resulting integral gives the estimate.
We can also verify directly that this absolutely convergent Fourier series equals $a$, without assuming a separate theorem on pointwise Fourier expansion.
Set
\[
 K_N(t)=\frac1N\left|\sum_{j=0}^{N-1}e^{2\pi\ii jt}\right|^2.
\]
This kernel is nonnegative, has integral one over a period, and satisfies
$K_N\le C_\epsilon/N$ where the distance from the integers is at least $\epsilon>0$; the last estimate follows from the finite geometric-series formula.
Uniform continuity, after splitting the convolution integral into parts near and far from the origin, shows that
$a*(K_N\otimes K_N)$ converges uniformly to $a$.
Here $K_N\otimes K_N(s,t)=K_N(s)K_N(t)$ and the periodic convolution is
$\int_0^1\int_0^1a(x-s,y-t)K_N(s)K_N(t)ds\,dt$.
Thus a continuous function with all Fourier coefficients zero must vanish.
The sum of the absolutely convergent series has the same coefficients as $a$, so the two are equal.
For a nonzero index, the multiplier of $\partial_{\bar z}$ is $\pi(\ii m-n)$, with
$|\pi(\ii m-n)|=\pi\sqrt{m^2+n^2}$.
If $a_{00}=0$, set
\[
 f(x,y)=\sum_{(m,n)\ne(0,0)}
 \frac{a_{mn}}{\pi(\ii m-n)}e^{2\pi\ii(mx+ny)}.
\]
Growth of the denominator and arbitrarily rapid Fourier decay make this series and all its derivatives uniformly convergent.
Indeed, group indices by $k\le |m|+|n|<k+1$; each group has $O(k+1)$ terms, and a sufficiently large decay exponent permits comparison with a convergent numerical series.
Thus $f$ is smooth and doubly periodic, and termwise differentiation gives $\partial_{\bar z}f=a$.
Every class consequently has a unique representative that is a constant multiple of $d\bar z$.
If $\dbar f=0$, the same multiplier calculation eliminates every nonzero Fourier coefficient, making $f$ constant.
Finally, Theorem \ref{sh:dolbeault} identifies these two computations with the asserted sheaf cohomology groups.
\end{proof}

\section{Two Routes Beyond This Chapter}

The Dolbeault quotient measures the obstruction to finding a global primitive for a closed form.
For the square torus, this obstruction is its constant Fourier coefficient; on a general manifold, it is a sheaf cohomology class.

Two complementary routes follow.
Chapter \ref{alg:chapter} proves Oka coherence, analytic flatness, and local finite free resolutions from Weierstrass division, explaining why finitely many relations hold uniformly on a neighborhood.
After Hodge theory and curvature vanishing, Chapter \ref{sh:projective-chapter} uses explicit contractions on compact fibers to compute direct images of projective bundles, allowing arbitrary coherent sheaves as coefficients.
Both routes use the sheaf theory already established here and retain the distinction between local solvability and cohomology vanishing on an arbitrary open set.


\chapter{Convergent Power Series, Coherent Sheaves, and Local Algebra}\label{alg:chapter}

\begin{learning}
Near a point, a holomorphic function is represented by a convergent power series.
Weierstrass division reduces suitable equations to polynomial equations in one variable, making finite-dimensional algebra available in a local analytic problem.
Coherence strengthens finite generation at one point to finite generation on a neighborhood.
The later sections use these ideas to study resolutions, finite projections, and flat families.
\end{learning}

The one-variable tools required here are the Cauchy formula, the argument principle, and elementary polynomial theory.
For local analytic geometry see \cite[Chapters I and II]{Demailly2012}; for commutative algebra see \cite[Commutative Algebra]{Stacks}.
For a first course, the principal results are Weierstrass division and Oka's coherence theorem.
The sections on resolutions, completion, and finite projections support later applications to singularities and direct images; they can be studied when those applications arise.

\section{Convergent Power Series and Division with Parameters}
All rings in this chapter are commutative with identity, and all modules are unital.
For elements $a_1,\ldots,a_s$ of a ring $A$, $(a_1,\ldots,a_s)$ is their generated ideal, consisting of the finite sums $\sum b_ja_j$ with $b_j\in A$.
An $A$-module is \emph{finite} if it is finitely generated; a finite ring extension means a ring homomorphism for which the target is finite as a module over the source.
A \emph{unit} has a multiplicative inverse; an \emph{integral domain} has no zero divisors and has $1\ne0$.
An ideal is \emph{maximal} if it is maximal among proper ideals, and a ring is \emph{local} if it has exactly one maximal ideal.
A ring is \emph{Noetherian} if every ideal is finitely generated.
Equivalently, increasing chains of ideals stabilize: their union is an ideal generated by finitely many members of the chain, while a non-finitely-generated ideal would give a strictly increasing chain by adjoining generators one at a time.


Write
\[
 R_n=\C\{z_1,\ldots,z_n\},\qquad
 \mathfrak m_n=(z_1,\ldots,z_n).
\]
Elements of $R_n$ are power series convergent on some neighborhood of the origin; equality means equality of the germs they define.
Taylor expansion from Theorem~\ref{bas:cauchy} identifies $R_n$ with the stalk $\cO_{\C^n,0}$ of holomorphic germs.
Write $z'=(z_1,\ldots,z_{n-1})$ and $w=z_n$, and put $R_0=\C$.
A function $f(z',w)$ is \emph{regular of order $d$ in $w$} if $f(0,w)$ has vanishing order $d\ge0$ at $w=0$; order zero means that $f(0,0)\ne0$.

\begin{lemma}[Units, the maximal ideal, and regular coordinates]\label{alg:local-ring}
The ring $R_n$ is a local integral domain with unique maximal ideal $\mathfrak m_n$.
An element $f$ is a unit if and only if $f(0)\ne0$.
After an invertible complex-linear coordinate change, every nonzero germ is regular in the last variable.
\end{lemma}
\begin{proof}
If $f(0)\ne0$, continuity gives a neighborhood where $|f-f(0)|<|f(0)|/2$.
The geometric series
\[
 \frac1f=\frac1{f(0)}\sum_{k\ge0}
 \left(-\frac{f-f(0)}{f(0)}\right)^k
\]
converges uniformly on smaller compact subsets and gives a holomorphic inverse by Corollary~\ref{bas:identity}.
If $f(0)=0$, evaluating $fg=1$ at the origin gives a contradiction.
Every series with zero constant term is $\sum_jz_jg_j$: group monomials by their first occurring variable, or take
$g_j(z)=\int_0^1(\partial f/\partial z_j)(tz)\,dt$.
Thus the nonunits form exactly the ideal $\mathfrak m_n$, whose quotient is $\C$, and every proper ideal is contained in it.
A product of nonzero holomorphic germs cannot vanish identically.
Indeed, if $fg=0$ and $f$ is nonzero at some nearby point, then $g$ vanishes near that point, and Corollary~\ref{bas:identity} gives $g=0$.
This proves that $R_n$ is an integral domain.

Finally, write $f=f_d+f_{d+1}+\cdots$, with $f_d$ the first nonzero homogeneous polynomial.
A nonzero polynomial cannot vanish at every $v\in\C^n$, so choose $v$ with $f_d(v)\ne0$ and complete $v$ to a linear coordinate basis.
In the new coordinates, $f(0,w)=f_d(v)w^d+O(w^{d+1})$.
\end{proof}

\begin{lemma}[Analytic division by a monic polynomial]\label{alg:polynomial-division}
Let
\[
 W(z',w)=w^d+a_1(z')w^{d-1}+\cdots+a_d(z'),\qquad a_j(0)=0.
\]
Every $g\in R_n$ has a unique expression
\[
 g=Wq+r,\qquad q\in R_n,\quad
 r=\sum_{j=0}^{d-1}r_j(z')w^j,\quad r_j\in R_{n-1}.
\]
The quotient and remainder coefficients are also holomorphic in any additional holomorphic parameters.
\end{lemma}
\begin{proof}
Shrink to a product neighborhood where $g,W$ are defined.
Choose a small circle $|\zeta|=\rho$.
At $z'=0$, the polynomial $W=\zeta^d$ is nonzero on this circle.
Shrink the base further so that $W(z',\zeta)\ne0$ on the circle and all $d$ roots remain inside it.
For $|w|<\rho$, define
\begin{align*}
 q(z',w)&=\frac1{2\pi\ii}\int_{|\zeta|=\rho}
       \frac{g(z',\zeta)}{W(z',\zeta)(\zeta-w)}\,d\zeta,\\
 r(z',w)&=\frac1{2\pi\ii}\int_{|\zeta|=\rho}
       \frac{g(z',\zeta)}{W(z',\zeta)}
       \frac{W(z',\zeta)-W(z',w)}{\zeta-w}\,d\zeta.
\end{align*}
The contour is independent of the parameters, and the denominators have uniform positive lower bounds on compact parameter subsets.
Differentiation in the parameters is therefore permitted, proving holomorphicity.
The difference quotient is a polynomial of degree at most $d-1$ in $w$, giving the required form of $r$.
Adding the formulas yields
\[
 W(z',w)q(z',w)+r(z',w)
 =\frac1{2\pi\ii}\int_{|\zeta|=\rho}
                  \frac{g(z',\zeta)}{\zeta-w}\,d\zeta=g(z',w).
\]

For uniqueness, the two proposed expressions may initially be defined only on smaller neighborhoods.
Subtract them to obtain $Wq=-r$, then choose a smaller circle and base neighborhood on which the identity holds and all roots of $W$ lie inside the circle.
For fixed $z'$, if $\alpha$ is a root of multiplicity $e$ of $W(z',\cdot)$, then
$r(z',\cdot)=-Wq$ has a zero of order at least $e$ at $\alpha$.
There are $d$ zeros counted with multiplicity, whereas $\deg_wr<d$, so $r(z',\cdot)=0$.
This holds for all sufficiently small $z'$, giving $r=0$, and Lemma~\ref{alg:local-ring} gives $q=0$.
Additional parameters occur only in the numerator and coefficients of the integral; the same argument proves the final assertion.
\end{proof}

\begin{theorem}[Weierstrass preparation]\label{alg:weierstrass-preparation}
If $f\in R_n$ is regular of order $d$ in $w=z_n$, there are a unique unit $u\in R_n$ and a unique monic polynomial
\[
 W=w^d+a_1(z')w^{d-1}+\cdots+a_d(z'),\qquad a_j(0)=0,
\]
such that $f=uW$.
The polynomial $W$ is the \emph{Weierstrass polynomial} of $f$.
\end{theorem}
\begin{proof}
Choose $\rho>0$ such that $f(0,w)$ has no zero for $0<|w|\le\rho$.
After shrinking the base, $f(z',w)$ is nonzero on $|w|=\rho$.
By the argument principle, the number of zeros inside, counted with multiplicity, is a continuous integer-valued function of $z'$.
It is therefore constant and equals its value $d$ at $z'=0$.
For $1\le k\le d$, set
\[
 s_k(z')=\frac1{2\pi\ii}\int_{|\zeta|=\rho}
       \zeta^k\frac{\partial f/\partial\zeta}{f}(z',\zeta)\,d\zeta.
\]
The residue theorem gives $s_k=\sum_{\nu=1}^d\alpha_\nu^k$, listing roots with multiplicity.
No continuous or holomorphic labeling of individual roots is needed; the integral itself proves that $s_k$ is holomorphic.
Starting with $e_0=1$, define recursively by Newton's identities
\[
 k e_k=\sum_{j=1}^k(-1)^{j-1}e_{k-j}s_j.
\]
To verify the recursion, write $E(t)=\prod_{\nu=1}^d(1+\alpha_\nu t)=\sum_ke_kt^k$.
Formal logarithmic differentiation gives
$E'(t)=E(t)\sum_{j\ge1}(-1)^{j-1}s_jt^{j-1}$.
Comparing the coefficients of $t^{k-1}$ gives the displayed identities, so the recursively defined $e_k$ are precisely the elementary symmetric functions of the roots.
Thus $W=w^d-e_1w^{d-1}+\cdots+(-1)^de_d$ has roots $\alpha_\nu$.
At $z'=0$ all roots are zero, so every nonleading coefficient vanishes at the origin.

Apply Lemma \ref{alg:polynomial-division} to write $f=Wq+r$.
For fixed $z'$, the functions $f,W$ have the same root multiplicities, so $r=f-Wq$ vanishes to at least those multiplicities.
As in the uniqueness argument, $\deg_wr<d$ forces $r=0$.
At $z'=0$, we have $q(0,w)=f(0,w)/w^d$, so $q(0,0)\ne0$.
Taking $u=q$ proves existence.

If also $f=\widetilde u\widetilde W$, both units are nonzero on a sufficiently small neighborhood.
For each small $z'$, the roots and multiplicities of $W$ and $\widetilde W$ agree with those of $f$.
Since both are monic of degree $d$, they are equal, and then $u=\widetilde u$.
\end{proof}

\begin{theorem}[Weierstrass division]\label{alg:weierstrass-division}
Under the preceding hypotheses, every $g\in R_n$ has a unique expression
\[
 g=fq+\sum_{j=0}^{d-1}r_j(z')w^j,\qquad q\in R_n,\quad r_j\in R_{n-1}.
\]
In particular, $R_n/(f)$ is free over $R_{n-1}$ with basis $1,w,\ldots,w^{d-1}$.
\end{theorem}
\begin{proof}
Write $f=uW$.
Divide by $W$ to obtain $g=WQ+r$, then set $q=u^{-1}Q$.
Uniqueness follows from uniqueness of division by $W$.
In the quotient, every class has a representative of degree less than $d$ in $w$, proving generation by the stated classes.
If such a polynomial lies in $(f)=(W)$, uniqueness of division makes it zero, proving linear independence.
\end{proof}

\begin{example}[A division that can be checked directly]
For $W(z,w)=w^2-z$,
$w^5=(w^2-z)(w^3+zw)+z^2w$.
Thus $R_2/(w^2-z)$ has basis $1,w$ over $\C\{z\}$.
This quotient is also isomorphic to $\C\{w\}$, but its $\C\{z\}$-module structure is determined by $z\mapsto w^2$.
Being free of rank two over $\C\{z\}$ and having one local parameter as a complex curve describe different information.
\end{example}

\section{Finite Generation, Nakayama's Lemma, and Closedness}

A module is Noetherian if every submodule is finitely generated, equivalently if increasing chains of submodules stabilize; the same union argument used for ideals proves the equivalence.

\begin{lemma}[Finite modules over a Noetherian ring]\label{alg:noetherian-modules}
If every ideal of a ring $A$ is finitely generated, then every submodule of $A^r$ is finitely generated.
Consequently, every submodule of a finitely generated $A$-module is finitely generated.
\end{lemma}
\begin{proof}
Induct on $r$. The case $r=1$ is the hypothesis.
For $N\subset A^r$, let $J$ be its image under projection onto the last coordinate.
Choose $v_1,\ldots,v_s\in N$ whose last coordinates generate $J$.
Subtracting a linear combination of the $v_j$ from any $v\in N$ makes its last coordinate zero.
The remainder lies in $N\cap(A^{r-1}\times0)$, which has finitely many generators by induction.
The two sets of generators together generate $N$.
If $M$ has $r$ generators, choose a surjection $A^r\to M$.
The inverse image of a submodule $L\subset M$ is finitely generated, and the images of its generators generate $L$.
\end{proof}

An analytic set is locally the common zero set of finitely many holomorphic functions.
The possibly nonreduced analytic subspace defined by an ideal sheaf $\cI=(f_1,\ldots,f_s)$ has structure sheaf $(\cO_U/\cI)|_{Z(f_1,\ldots,f_s)}$.
Thus its local ring $\cO_{Z,x}$ at a point is a quotient of a convergent power-series ring, and $\mathfrak m_x$ denotes its maximal ideal of germs vanishing at $x$.
A complex analytic space is a ringed space locally of this form; later we identify the reduced structure using the ideal of all vanishing functions.

\begin{theorem}[Noetherianity of analytic local rings]\label{alg:noetherian}
The ring $R_n=\C\{z_1,\ldots,z_n\}$ is Noetherian.
Its quotients, in particular local rings of complex analytic spaces, are Noetherian as well.
\end{theorem}
\begin{proof}
Induct on $n$. For $n=0$ the ring is the field $\C$.
Let $I\subset R_n$ be nonzero and choose $0\ne f\in I$.
By Lemma \ref{alg:local-ring}, make a linear coordinate change so that $f$ is regular in $w$.
Theorem~\ref{alg:weierstrass-preparation} gives $f=uW$; since $u$ is a unit, $W\in I$.
Divide all elements of $I$ by $W$. Their remainders form
\[
 N=I\cap\left\{\sum_{j<d}a_j(z')w^j\right\}\subset R_{n-1}^d.
\]
A remainder stays in $I$ because $g-Wq\in I$, and these remainders are closed under addition and multiplication by $R_{n-1}$.
The induction hypothesis and Lemma \ref{alg:noetherian-modules} give generators $r_1,\ldots,r_s$.
Every remainder is an $R_{n-1}$-linear combination of them, so
$I=(W,r_1,\ldots,r_s)$.
The zero ideal is finitely generated as well.
Finally, ideals of $A/J$ are images of ideals of $A$ containing $J$, so quotients inherit Noetherianity.
\end{proof}

In any local ring $(A,\mathfrak m)$, an element outside $\mathfrak m$ is a unit: otherwise its principal ideal would be proper and would lie in a maximal ideal, necessarily $\mathfrak m$.

\begin{lemma}[Nakayama's lemma]\label{alg:nakayama}
Let $(A,\mathfrak m)$ be a local ring, let $M$ be a finitely generated $A$-module, and let $I\subset\mathfrak m$.
If $M=IM$, then $M=0$.
Equivalently, if the images of $m_1,\ldots,m_r$ generate $M/\mathfrak mM$, then these elements generate $M$.
\end{lemma}
\begin{proof}
Choose generators $v_1,\ldots,v_s$.
The identity $M=IM$ gives $v_i=\sum_j a_{ij}v_j$, with $a_{ij}\in I$.
Put $B=(a_{ij})$; the ring remains denoted by $A$. In matrix form, $(\operatorname{Id}-B)v=0$.
Left multiplication by the adjugate gives $\det(\operatorname{Id}-B)v_i=0$ for every $i$.
The determinant is congruent to $1$ modulo $\mathfrak m$, hence lies outside the maximal ideal and is a unit.
All $v_i$ therefore vanish.
For the equivalent statement, set $N=\sum_iAm_i$.
The hypothesis says $M/N=\mathfrak m(M/N)$, and $M/N$ is finitely generated.
The assertion just proved gives $M=N$.
Conversely, if $M=IM$ with $I\subset\mathfrak m$, then $M/\mathfrak mM=0$; apply the second statement to the empty generating set.
\end{proof}

For a holomorphic vector bundle, reducing its module of section germs modulo the maximal ideal gives the fiber at the chosen point.
Nakayama's lemma says that a basis of this quotient lifts to a generating set of the module of germs, provided that module is finite.
It does not yet give generators on one fixed neighborhood; that is the additional issue addressed by coherence.

\begin{lemma}[Hilbert's basis theorem]\label{alg:hilbert-basis}
If $A$ is Noetherian, then $A[t]$ is Noetherian.
\end{lemma}
\begin{proof}
Let $J\subset A[t]$ be an ideal; the case $J=0$ is immediate.
The leading coefficients of nonzero polynomials in $J$, together with zero, form an ideal $L$.
For addition, first multiply by powers of $t$ to align degrees; the sum of leading coefficients is either the new leading coefficient or zero.
Multiplication by elements of $A$ gives the same conclusion.
Choose $f_1,\ldots,f_s\in J$ whose leading coefficients generate $L$, and put $d=\max\deg f_i$.
Elements of $J$ of degree less than $d$ form a submodule of $A^d$, so Lemma~\ref{alg:noetherian-modules} gives finitely many generators $g_1,\ldots,g_t$.
For $f\in J$ with $\deg f=e\ge d$, express its leading coefficient as a combination of those of the $f_i$.
Subtract the corresponding $\sum_i a_it^{e-\deg f_i}f_i$ to lower the degree strictly.
After finitely many repetitions the remainder has degree less than $d$ and is generated by the $g_j$.
Thus $J=(f_1,\ldots,f_s,g_1,\ldots,g_t)$.
\end{proof}

A graded ring or module is a direct sum of pieces indexed by nonnegative degrees, with multiplication adding degrees.
An element in one degree is \emph{homogeneous}; below $t$ is a formal variable recording that degree, and $I^k$ is the ideal generated by products of $k$ elements of $I$, with $I^0=A$.

\begin{theorem}[The Artin--Rees lemma]\label{alg:artin-rees}
Let $A$ be Noetherian, $I\subset A$ an ideal, and $N\subset M$ a submodule of a finitely generated module.
There is an integer $c\ge0$ such that for all $k\ge c$,
\[
 I^kM\cap N=I^{k-c}(I^cM\cap N).
\]
\end{theorem}
\begin{proof}
Write $I=(a_1,\ldots,a_s)$.
The graded ring
\[
 \mathcal R=\bigoplus_{k\ge0}I^kt^k=A[a_1t,\ldots,a_st]
\]
is a quotient of $A[T_1,\ldots,T_s]$, hence Noetherian by Lemma~\ref{alg:hilbert-basis}.
The module $\mathcal M=\bigoplus_{k\ge0}I^kMt^k$ is generated over $\mathcal R$ by any finite generating set of $M$ placed in degree zero.
It is therefore a finite $\mathcal R$-module.
Its submodule
$\mathcal N=\bigoplus_{k\ge0}(I^kM\cap N)t^k$
is finitely generated.
Generators may be chosen homogeneous: replace each member of a finite generating set by all its homogeneous components.
Let their degrees be $d_1,\ldots,d_l$ and take $c=\max d_i$; if $\mathcal N=0$, take $c=0$.

The degree-$k$ part for $k\ge c$ is generated by contributions contained in
$I^{k-d_i}(I^{d_i}M\cap N)$.
Since
$I^{k-d_i}=I^{k-c}I^{c-d_i}$ and
$I^{c-d_i}(I^{d_i}M\cap N)\subset I^cM\cap N$,
the left side lies in the right side of the claimed equality.
Conversely, $I^{k-c}(I^cM\cap N)$ lies both in $I^kM$ and in $N$.
This proves equality.
\end{proof}

\begin{theorem}[Krull intersection and closedness of submodules]\label{alg:krull-intersection}
If $(A,\mathfrak m)$ is a Noetherian local ring and $M$ is finitely generated, then
\[
 \bigcap_{k\ge0}\mathfrak m^kM=0.
\]
\end{theorem}
\begin{proof}
Set $K=\bigcap_k\mathfrak m^kM$.
Lemma~\ref{alg:noetherian-modules} makes $K$ finitely generated.
Apply Theorem~\ref{alg:artin-rees} to $K\subset M$ and take $k=c+1$.
Since $K\subset\mathfrak m^kM$ and $K\subset\mathfrak m^cM$,
\[
 K=\mathfrak m^{c+1}M\cap K
   =\mathfrak m(\mathfrak m^cM\cap K)=\mathfrak mK.
\]
Lemma~\ref{alg:nakayama} gives $K=0$.
\end{proof}

\begin{corollary}[Approximation modulo powers of the maximal ideal detects containment]\label{alg:krull-closure}
Under the preceding hypotheses, every submodule $N\subset M$ satisfies
\[
 \bigcap_{k\ge0}(N+\mathfrak m^kM)=N.
\]
In particular, if ideals $I,J\subset A$ satisfy
$J\subset I+\mathfrak m^k$ for every $k$, then $J\subset I$.
\end{corollary}
\begin{proof}
The module $M/N$ is finitely generated, and
$(N+\mathfrak m^kM)/N=\mathfrak m^k(M/N)$.
Theorem~\ref{alg:krull-intersection} applied to $M/N$ shows that the image of the intersection is zero, so the original intersection lies in $N$.
The reverse containment is immediate.
Take $M=A$ and $N=I$ for the last assertion.
\end{proof}

\section{From Finite Generation at a Point to Coherence on a Neighborhood}
For functions $f_1,\ldots,f_r$, a relation is a tuple of holomorphic functions $(a_1,\ldots,a_r)$ satisfying $\sum_j a_jf_j=0$.
For instance, $(-w,z)$ is a relation between the coordinate functions $(z,w)$.
Coherence asks for finitely many such tuples which generate every relation near each point of one common neighborhood.

\begin{definition}[Finite presentation and coherence]
On a ringed space $(X,\cO_X)$, a module sheaf $\cF$ is \emph{locally finitely generated} if near every point there is a surjection $\cO_X^r\to\cF$.
It is \emph{locally finitely presented} if near every point there is an exact sequence
$\cO_X^s\to\cO_X^r\to\cF\to0$.
It is \emph{coherent} if it is locally finitely generated and the kernel of every local morphism $\cO_X^r\to\cF$ is locally finitely generated.
The word ``every'' ensures independence from a chosen generating set.
\end{definition}

\begin{lemma}[A polynomial model for a finite projection]\label{alg:finite-polynomial-model}
Let $W(z',w)$ be a Weierstrass polynomial of degree $d$.
Shrink the base and the disc so that all roots lie inside the disc.
Then the direct image to the base of the structure sheaf of the hypersurface $Z=(W=0)$ is locally represented uniquely by
$1,w,\ldots,w^{d-1}$.
Writing $B$ for the base and $\cO_{a'}=\cO_{B,a'}$, the stalk at $a'$ of this rank-$d$ algebra is isomorphic to the finite product
\[
 \prod_{b:\,W(a',b)=0}\cO_{Z,(a',b)}.
\]
\end{lemma}
\begin{proof}
First explain the product decomposition.
Choose disjoint small discs about the distinct roots $b_1,\ldots,b_l$ of $W(a',\cdot)$.
After shrinking the base near $a'$, no root lies on their boundaries.
Apply the contour integrals and Newton identities from Theorem~\ref{alg:weierstrass-preparation} separately to each boundary.
They give monic factors $W_\nu$ whose roots are exactly those in the corresponding disc, with $W=\prod_\nu W_\nu$.
At the base point, two factors have disjoint roots and are coprime.
The coefficient matrix for finding a polynomial linear combination of the two factors equal to $1$ is invertible there (its determinant is their resultant); solving this finite system holomorphically nearby gives
$A_{\nu\mu}W_\nu+B_{\nu\mu}W_\mu=1$.
The factors are pairwise coprime.

The Chinese remainder theorem can be verified directly here.
The coprimality identities construct polynomial classes $e_\nu$ equal to $1$ modulo $W_\nu$ and $0$ modulo the other factors.
The inverse of $P\mapsto(P\bmod W_\nu)_\nu$ is
$(P_\nu)_\nu\mapsto\sum_\nu e_\nu P_\nu$.
Hence
\[
 \cO_{a'}[w]/(W)\simeq\prod_\nu\cO_{a'}[w]/(W_\nu).
\]
At $(a',b_\nu)$, all other factors are units.
Division in the local coordinate $w-b_\nu$ gives
$\cO_{Z,(a',b_\nu)}\simeq\cO_{a'}[w]/(W_\nu)$.
This proves the asserted stalk isomorphism.

On the base, consider the free module sheaf $\cO^d$ and the natural map
$(r_0,\ldots,r_{d-1})\mapsto[\sum r_jw^j]$.
The stalk calculation shows that it is an isomorphism at every point and therefore a sheaf isomorphism by Proposition \ref{sh:stalk-isomorphism}.
This also handles functions initially defined only in separate small neighborhoods of the roots; such functions cannot simply be assumed to extend to the whole vertical disc.
\end{proof}

\begin{theorem}[Oka's coherence theorem]\label{alg:oka}
The structure sheaf $\cO_X$ of a complex manifold is coherent.
More explicitly, for every finite holomorphic matrix $A$, the relation sheaf
$\ker(A:\cO_X^r\to\cO_X^s)$ has finitely many generators near each point.
\end{theorem}
\begin{proof}
The problem is local, so work in $\C^n$ and induct on $n$.
For $n=0$ this is finite-dimensional linear algebra.
Assume that in dimension $n-1$ the relation sheaf of every finite matrix is locally finitely generated.
First prove the assertion in dimension $n$ for a single row $(f_1,\ldots,f_r)$.
If all germs $f_i$ at the point are zero, they vanish on a common neighborhood and the relation sheaf is $\cO^r$.
Otherwise reorder so that $f_1\ne0$.
Choose regular coordinates and write $f_1=uW$, with $u$ a unit.
Multiplying the first unknown by $u$ is an automorphism of the free module, so it suffices to treat $(W,f_2,\ldots,f_r)$.

Let $B$ be a small base neighborhood and $D$ a vertical disc.
By Lemma \ref{alg:finite-polynomial-model}, functions modulo $W$ form the algebra represented by $\cO_B^d$.
Multiplication by $f_j$ followed by taking the remainder is a $d\times d$ matrix $M_j$ with holomorphic coefficients on $B$.
By the dimension induction hypothesis, after shrinking $B$ the kernel $K$ of
\[
 (\cO_B^d)^{r-1}\longrightarrow\cO_B^d,
 \qquad (a_2,\ldots,a_r)\longmapsto\sum_{j=2}^rM_ja_j
\]
is generated by finitely many sections $k^1,\ldots,k^N$.
Represent each $k^\nu_j$ by a polynomial $a_j^\nu(z',w)$ of degree less than $d$.
Their relation in the quotient algebra says that the remainder of $\sum_{j\ge2}f_ja_j^\nu$ upon division by $W$ vanishes.
Thus this sum is $Wq^\nu$, with $q^\nu$ holomorphic on $B\times D$.
Consequently,
\[
 v^\nu=(-q^\nu,a_2^\nu,\ldots,a_r^\nu)
\]
is an actual relation.
For each $2\le j\le r$, also define $e^j$ with first component $-f_j$, $j$th component $W$, and all other components zero.
These too are relations.

We verify that these finitely many relations generate at \emph{every} nearby point.
Where $W(x)\ne0$, any relation $(g_1,\ldots,g_r)$ equals
$\sum_{j\ge2}(g_j/W)e^j$.
If $W(x)=0$, write $x=(a',b)$.
Place the classes modulo $W$ of $(g_2,\ldots,g_r)$ in the factor corresponding to the fiber point $b$ and put zero in the factors at the other distinct fiber points.
Lemma~\ref{alg:finite-polynomial-model} identifies this tuple with a germ of the quotient algebra on the base.
Since $\sum f_jg_j=-Wg_1$, the germ lies in $K_{a'}$ and is therefore
$\sum_\nu c_\nu k^\nu$ with $c_\nu\in\cO_{B,a'}$.
Subtract $\sum_\nu c_\nu v^\nu$ from the relation.
Every component indexed by $j\ge2$ is now divisible by $W$ in $\cO_{X,x}$; write it as $Wh_j$.
Subtract $\sum_{j\ge2}h_je^j$.
Only a first component $h_1$ remains, and it satisfies $Wh_1=0$.
Lemma~\ref{alg:local-ring} says that the local ring is an integral domain; since $W\ne0$, we have $h_1=0$.
The single-row relation sheaf is thus finitely generated on a common neighborhood.

Now induct on the number of rows $s$.
The first-row kernel has finitely many generators, giving a surjection
$P:\cO^t\to\ker A_1$.
Compose the remaining $s-1$ rows with $P$ to obtain a matrix $B$.
By row induction, $\ker B$ is locally finitely generated.
Images under $P$ of these generators generate $\ker A$:
a germ $v\in\ker A$ lifts through $P$ to $u$, and the remaining rows give $Bu=0$.
Express $u$ in the chosen generators and apply $P$ to express $v$ in their images.
This proves the assertion for every finite matrix.
The case $s=1$ gives coherence of the structure sheaf.
\end{proof}

\begin{proposition}[Basic operations on coherent sheaves]\label{alg:coherent-category}
If a sheaf of rings $\cO$ is coherent, its finitely presented module sheaves are exactly its coherent module sheaves.
Kernels, images, and cokernels of morphisms between coherent module sheaves are coherent, as are extensions and finite direct sums.
\end{proposition}
\begin{proof}
Work first on a small neighborhood.
The sheaf $\cO^r$ is coherent: take kernels row by row for a matrix into $\cO^r$, using the lifting of generators in the last paragraph of the proof of Theorem~\ref{alg:oka}.
This uses only coherence of $\cO$.
If $\cF$ is coherent, choose a finite generating surjection $p:\cO^r\to\cF$.
Its kernel is finitely generated, giving a finite presentation.
Conversely, suppose $\cO^s\xrightarrow{B}\cO^r\to\cF\to0$ is exact.
Given $a:\cO^t\to\cF$, shrink and lift its finitely many basis sections to a matrix
$A:\cO^t\to\cO^r$.
The equation $a(v)=0$ means that locally $Av=Bw$ for some $w$.
Thus $\ker a$ is the projection onto the first factor of the kernel of $(A,-B)$.
That matrix kernel is finitely generated, so its image is finitely generated.
This proves coherence of $\cF$.

Let $a:\cF\to\cG$ be a morphism between finitely presented sheaves.
Choose $\cO^r\twoheadrightarrow\cF$ and lift the composite into $\cG$ to its free presentation.
The same matrix-relation construction gives finitely many generators of
$\ker(\cO^r\to\cG)$, which surjects onto $\ker a$.
This also presents the source kernel as the image of a matrix kernel between finite free sheaves; applying the same construction to its relations yields a finite presentation of $\ker a$.
More directly, any map $\cO^t\to\ker a$ can be viewed as a map into $\cF$, and its kernel is finitely generated by coherence of $\cF$.
Together with the finite generation already obtained, this proves that $\ker a$ is coherent.
The image $\operatorname{im}a$ is finitely generated, and the kernel of a free-sheaf map into it is the kernel of the composite into $\cG$.
Thus the image is coherent as well.
To present the cokernel, add the images of finitely many generators of $\cF$ to the finite relation list of $\cG$.
The cokernel is consequently finitely presented and coherent.

Finally, suppose $0\to\cF\to\cH\to\cG\to0$ has coherent endpoints.
Lift a finite generating set of $\cG$ and adjoin generators of $\cF$ to generate $\cH$ locally.
For $u:\cO^t\to\cH$, the kernel $K$ of the composite into $\cG$ is coherent.
The sheaf $\ker u$ is the kernel of $K\to\cF$, hence coherent by the kernel assertion just proved.
Thus $\cH$ is coherent.
Finite direct sums are split extensions.
\end{proof}

\begin{corollary}[Pullback of coherent sheaves]\label{alg:coherent-pullback}
If $f:X\to Y$ is a holomorphic map of complex manifolds or complex analytic spaces and $\cF$ is coherent on $Y$, then $f^*\cF$ is coherent on $X$.
\end{corollary}
\begin{proof}
Choose a local finite presentation $\cO_Y^s\to\cO_Y^r\to\cF\to0$.
Propositions~\ref{sh:inverse-exact} and \ref{sh:tensor}, with Definition~\ref{sh:module-pullback}, give
$\cO_X^s\to\cO_X^r\to f^*\cF\to0$.
The pullback is finitely presented, hence coherent by Proposition \ref{alg:coherent-category}.
For a complex analytic space, the structure sheaf is locally $\cO_U/\cI$, with $U\subset\C^n$ and $\cI$ generated by finitely many local defining equations.
Theorem~\ref{alg:oka} and Proposition~\ref{alg:coherent-category} make this quotient, and its finitely presented modules, coherent.
Only right exactness of pullback is required; no flatness of $f$ is assumed.
\end{proof}

\begin{proposition}[Support and the openness of injectivity and surjectivity]\label{alg:coherent-support}
The support $\supp\cF=\{x:\cF_x\ne0\}$ of a coherent sheaf is a closed analytic subset.
For a morphism $a:\cF\to\cG$ of coherent sheaves, both the set where $a_x$ is injective and the set where it is surjective are open.
\end{proposition}
\begin{proof}
Choose a local presentation $\cO^s\xrightarrow{B}\cO^r\to\cF\to0$.
Lemma~\ref{alg:nakayama} gives
\[
 \cF_x=0\iff\cF_x/\mathfrak m_x\cF_x=0
 \iff B(x):\C^s\to\C^r\text{ is surjective}.
\]
The last condition means that some $r\times r$ minor is nonzero.
The support is therefore the common zero set of all $r$th-order minors.
If $s<r$, there are no such minors and the support is the whole neighborhood; if $r=0$, it is empty.
These local analytic closed sets define a globally closed set because every point outside the support has a neighborhood on which the sheaf vanishes.

Both $\ker a$ and $\operatorname{coker}a$ are coherent by Proposition~\ref{alg:coherent-category}.
Injectivity of $a_x$ is equivalent to $(\ker a)_x=0$, and surjectivity to $(\operatorname{coker}a)_x=0$.
The two sets in question are thus complements of closed supports.
\end{proof}

\section{Koszul Complexes and Finite Free Resolutions}
Choosing generators gives a surjection from a free module, and choosing generators of its kernel gives a presentation.
Repeating this procedure records relations among relations.
A free resolution organizes the resulting sequence, while the Koszul complex provides an explicit model for the coordinate functions.
A chain complex $C_\bullet$ has differentials $d:C_j\to C_{j-1}$ with $d^2=0$, and its homology is
$H_j(C_\bullet)=\ker(d:C_j\to C_{j-1})/\operatorname{im}(d:C_{j+1}\to C_j)$.
A free resolution of an $R$-module $M$ is an exact sequence of free $R$-modules ending in a surjection to $M$.
These algebraic differentials are module homomorphisms, distinct from exterior differentiation of forms.


\begin{lemma}[The Koszul complex of the coordinates]\label{alg:koszul}
Let $R=R_n$ and $\mathfrak m=\mathfrak m_n$, let $e_1,\ldots,e_n$ be the standard basis of $R^n$, set $K_j=\bigwedge^jR^n$, and define
\[
 d(e_{i_1}\wedge\cdots\wedge e_{i_j})
 =\sum_{\nu=1}^j(-1)^{\nu-1}z_{i_\nu}
 e_{i_1}\wedge\cdots\widehat e_{i_\nu}\cdots\wedge e_{i_j}.
\]
Then $d^2=0$, and $K_\bullet\to R/\mathfrak m=\C$ is a free resolution of length $n$.
\end{lemma}
\begin{proof}
Every term of $d^2$ deleting two distinct basis vectors appears twice with opposite signs and cancels.
Prove exactness by induction on the number of coordinates used.
For one coordinate the complex is
$0\to R\xrightarrow{z_1}R\to R/(z_1)\to0$.
Multiplication by $z_1$ is injective because $R$ is a domain.
On adjoining $z_j$, the new degree-$q$ term is the sum of the old terms $K_q\oplus K_{q-1}e_j$, and its differential is
$(u,v)\mapsto(du+(-1)^{q-1}z_jv,dv)$.
By induction the old homology is concentrated in degree zero and equals
$R/(z_1,\ldots,z_{j-1})=\C\{z_j,\ldots,z_n\}$, on which multiplication by $z_j$ is injective by Lemma~\ref{alg:local-ring}.
Here is the elimination explicitly. If $(u,v)$ is a cycle in degree $q>1$, then $dv=0$ and old exactness gives $v=db$.
If $q=1$, the equation $du+z_jv=0$ says $z_j[v]=0$ in old degree-zero homology; injectivity again gives $v=db$.
In both cases subtract the new boundary of $(0,b)$, whose second component is $db=v$.
The remaining second component is zero, and the first component is an old positive-degree cycle, hence an old boundary.
It too is a new boundary. All positive-degree homology therefore vanishes.
In degree zero the new boundaries add $z_j$ to the old coordinate ideal, so the quotient is $R/(z_1,\ldots,z_j)$.
Set $j=n$ to conclude.
\end{proof}

\begin{theorem}[Finite local free resolutions on analytic manifolds]\label{alg:regular-resolution}
Every finite $R_n$-module has a finite free resolution of length at most $n$.
Consequently, on an $n$-dimensional complex manifold, every coherent sheaf has, near each point, an exact sequence
\[
 0\longrightarrow\cO^{r_n}\longrightarrow\cdots
 \longrightarrow\cO^{r_1}\longrightarrow\cO^{r_0}
 \longrightarrow\cF\longrightarrow0.
\]
\end{theorem}
\begin{proof}
First work over $R=R_n$ with maximal ideal $\mathfrak m=\mathfrak m_n$.
Choose a basis of $M/\mathfrak mM$ and lift it to $M$.
Lemma~\ref{alg:nakayama} gives a surjection $F_0\to M$ with $F_0$ finite free.
Its kernel is finitely generated by Theorem~\ref{alg:noetherian} and Lemma~\ref{alg:noetherian-modules}.
Repeat on each kernel to obtain a possibly infinite free resolution
$\cdots\to F_2\to F_1\to F_0\to M\to0$.
At every step use a minimal generating set modulo $\mathfrak m$.
Then $d(F_i)\subset\mathfrak mF_{i-1}$: a relation with a unit coefficient would express one generator in terms of the others, contradicting minimality.
All differentials of $F_\bullet\otimes_R\C$ are therefore zero.

Consider the double complex $F_p\otimes_RK_q$, with $K$ the Koszul resolution of Lemma~\ref{alg:koszul}.
On bidegree $(p,q)$ the total differential is $d_F\otimes1+(-1)^p1\otimes d_K$, and total degree is $p+q$.
For fixed $p$, freeness of $F_p$ makes homology in the $q$ direction vanish except at $q=0$, where it is $F_p\otimes\C$.
For fixed $q$, freeness of $K_q$ makes homology in the $p$ direction vanish except at $p=0$, where it is $M\otimes K_q$.
Both calculations give the same total homology by the chain-complex version of Lemma~\ref{sh:double-complex}, as can also be seen directly by elimination.
For a total cycle in a fixed degree, express an outermost component outside the surviving homology as a boundary and subtract the corresponding total boundary.
This removes one component and introduces only a neighboring component in the other direction.
There are only finitely many components in a fixed total degree, so elimination terminates after finitely many steps.
The same operation on boundaries proves injectivity of the resulting identification.
Thus
\[
 F_i/\mathfrak mF_i
 =H_i(F_\bullet\otimes\C)
 \simeq H_i(M\otimes K_\bullet)=0\qquad(i>n),
\]
since $K_i=0$ for $i>n$.
Lemma~\ref{alg:nakayama} gives $F_i=0$ for $i>n$, so the resolution ends at step $n$.

Finally, represent all the finitely many matrix germs involved on a common neighborhood of $x$, together with the final map to $\cF$.
Their composites vanish as germs and hence vanish identically after shrinking.
The homology sheaves of this finite complex are coherent by Proposition \ref{alg:coherent-category}, and their stalks at $x$ vanish.
Proposition \ref{alg:coherent-support} permits a further shrinking on which these finitely many homology sheaves vanish simultaneously.
This gives the asserted exact sheaf resolution.
\end{proof}

\begin{remark}[How finite generation must be used]
To prove coherence of an analytic ideal sheaf, it is insufficient to observe that each stalk is an ideal in a Noetherian ring.
That observation gives finitely many generators separately at each point, without giving a uniform finite generating set on a neighborhood.
One must construct finitely many holomorphic functions on a neighborhood and prove that they generate the required ideal at every nearby point.
Once local finite generation is established, Oka's theorem supplies finite generation of the relation sheaf.
Krull closedness, Corollary \ref{alg:krull-closure}, often upgrades containments valid to arbitrarily high order into actual containments.
\end{remark}

\section{Completion and Flatness of Analytic Products}
Flatness describes when relations remain exact after a change of base.
For example, the projection from a product should introduce new variables without creating a new relation among functions of the base alone.
Proving this for convergent power series requires some additional algebra, developed in this section.

\begin{definition}[Flat maps]\label{alg:flat-map}
A ring homomorphism $A\to B$ is \emph{flat} if $B$ is flat as an $A$-module.
A morphism of ringed spaces $f:X\to Y$ is flat at $x$ if
$\cO_{Y,f(x)}\to\cO_{X,x}$ is flat, and is flat if it is flat at every point.
By Definition~\ref{sh:module-pullback} and Proposition~\ref{sh:flat-stalks}, pullback along a flat map preserves exact sequences.
A flat ring map is \emph{faithfully flat} if $M\otimes_A B=0$ implies $M=0$ for every $A$-module $M$.
\end{definition}

A proper ideal $\mathfrak p$ is \emph{prime} if $ab\in\mathfrak p$ implies $a\in\mathfrak p$ or $b\in\mathfrak p$; $\operatorname{Spec}A$ is the set of prime ideals of $A$.
For a multiplicatively closed set $S\subset A$, the localization $S^{-1}A$ consists of fractions $a/s$, with $a/s=b/t$ when some $u\in S$ satisfies $u(ta-sb)=0$; module localization is defined by the same rule.
Write $A_{\mathfrak p}=(A\setminus\mathfrak p)^{-1}A$, $A[1/h]$ when powers of $h$ are inverted, and $\operatorname{Frac}(A)$ when $A$ is a domain and all nonzero elements are inverted.
The residue field is $\kappa(\mathfrak p)=A_{\mathfrak p}/\mathfrak pA_{\mathfrak p}$.
For a ring map $A\to B$, $\mathfrak q\cap A$ denotes the inverse image of $\mathfrak q$, even if the map is not injective.
Finally, $\operatorname{Ann}_A M=\{a:aM=0\}$ and $\operatorname{Ann}_A(v)=\{a:av=0\}$ are annihilators.
A module is faithful if its annihilator is zero. A module over a domain is torsion-free if multiplication by every nonzero ring element is injective.

\begin{lemma}[The ideal criterion and the localization criterion]\label{alg:flat-criteria}
Let $A$ be Noetherian.
An $A$-module $M$ is flat if and only if the natural map $I\otimes_AM\to M$ is injective for every ideal $I\subset A$.
For any ring homomorphism $A\to B$, the ring $B$ is flat over $A$ if and only if, for each
$\mathfrak q\in\operatorname{Spec}B$ and $\mathfrak p=\mathfrak q\cap A$, the ring $B_{\mathfrak q}$ is flat over $A_{\mathfrak p}$.
\end{lemma}
\begin{proof}
Necessity of the first condition follows by tensoring the injection $I\hookrightarrow A$.
For sufficiency, first prove injectivity after tensoring a submodule $N\subset A^r$, by induction on $r$.
For $r=1$ this is the hypothesis.
Project $N$ onto the last coordinate, with image ideal $I$ and kernel $N'\subset A^{r-1}$.
If $v\in N\otimes M$ maps to zero in $A^r\otimes M$, its image in $I\otimes M$ is zero by injectivity of $I\otimes M\to M$.
Right exactness lifts $v$ from some $v'\in N'\otimes M$.
The image of $v'$ in $A^{r-1}\otimes M$ is zero, so induction gives $v'=0$ and hence $v=0$.
A tensor element in an arbitrary free module involves only finitely many coordinates.
Intersecting the submodule with the finite free direct summand spanned by these coordinates extends the result to arbitrary free modules.

Now take an injection $L\subset P$, choose a free surjection $E\to P$, let $K$ be its kernel, and let $N$ be the inverse image of $L$.
If $\ell\in L\otimes M$ maps to zero in $P\otimes M$, lift it by right exactness to
$n\in N\otimes M$.
Its image in $E\otimes M$ comes from $K\otimes M$.
Subtract this preimage from $n$; the difference maps to zero in $E\otimes M$ and is zero by the injectivity already proved for $N\subset E$.
Since $K$ maps to zero in $L$, we obtain $\ell=0$.
This proves flatness.

Localization is exact: a fraction vanishes if some allowed denominator annihilates its numerator, and finitely many such conditions can be handled with the product of the denominators.
Thus localization rings are flat, and composites of flat homomorphisms are flat.
If $B$ is $A$-flat, localization at $\mathfrak q$ remains $A$-flat and, since the elements outside $\mathfrak p$ act invertibly, is $A_{\mathfrak p}$-flat. This proves necessity in the second assertion.
Conversely, if an injection fails to remain injective after tensoring with $B$, take a nonzero element $v$ in its kernel and a maximal ideal $\mathfrak q\subset B$ containing its annihilator.
Then $v/1\ne0$ after localization at $\mathfrak q$.
But the localized map is obtained by localizing first at $A_{\mathfrak p}$ and tensoring with $B_{\mathfrak q}$, so the hypothesis makes it injective, a contradiction.
\end{proof}

The $\mathfrak m$-adic topology on a module $M$ has neighborhoods $\mathfrak m^kM$ of zero.
The inverse limit $\varprojlim M/\mathfrak m^kM$ consists of compatible sequences of residue classes, one at each order.
The formal power-series ring $\C[[z_1,\ldots,z_n]]$ allows arbitrary coefficient series, without an analytic convergence requirement.
The associated graded ring is
$\operatorname{gr}_{\mathfrak m}A=\bigoplus_{j\ge0}\mathfrak m^j/\mathfrak m^{j+1}$;
the initial form of a nonzero element is its first nonzero homogeneous class.

\begin{proposition}[Completion is exact and faithfully flat]\label{alg:completion-flat}
Let $(A,\mathfrak m)$ be Noetherian local, and write
$\widehat M=\varprojlim M/\mathfrak m^kM$.
Completion preserves short exact sequences of finite modules, and
$\widehat M\simeq M\otimes_A\widehat A$.
The map $A\to\widehat A$ is faithfully flat, and $\widehat A$ is Noetherian.
For an analytic local ring $A=\C\{z\}/I$,
$\widehat A=\C[[z]]/(I)$.
\end{proposition}
\begin{proof}
For $0\to N\to M\to P\to0$, at level $k$ there is an exact sequence
\[
0\to N/(N\cap\mathfrak m^kM)\to M/\mathfrak m^kM
\to P/\mathfrak m^kP\to0.
\]
Transition maps on the left are surjective.
Given a compatible sequence on the right, choose a lift at the first level.
After choosing a lift at level $k+1$, its image differs from the chosen level-$k$ lift by an element on the left.
Lift that difference and subtract it to make the lifts compatible.
This proves surjectivity at the end after taking inverse limits; middle exactness and injectivity are checked levelwise.
Theorem~\ref{alg:artin-rees} shows that the filtrations $N\cap\mathfrak m^kM$ and $\mathfrak m^kN$ are cofinal: every term of either contains a sufficiently far term of the other.
They therefore define the same compatible approximations and the same inverse limit $\widehat N$.
Completion is therefore exact.

Take a finite presentation $A^s\to A^r\to M\to0$ and apply both tensor product with $\widehat A$ and completion.
In each case the final module is the cokernel of the same matrix
$\widehat A^s\to\widehat A^r$.
The natural map $M\otimes\widehat A\to\widehat M$ is thus an isomorphism.
Exact completion of $I\hookrightarrow A$ gives
$I\otimes\widehat A\hookrightarrow\widehat A$ for every ideal $I$.
Lemma~\ref{alg:flat-criteria} proves flatness.
For any nonzero module $M$, choose a nonzero cyclic submodule $A/J\subset M$.
Since $J$ is proper, $J\subset\mathfrak m$, and
$\widehat A/J\widehat A\twoheadrightarrow\widehat A/\mathfrak m\widehat A=A/\mathfrak m\ne0$.
Flatness embeds this nonzero module into $M\otimes\widehat A$, proving faithful flatness.

Noetherianity of the completion also follows without an additional hypothesis.
The associated graded ring $\operatorname{gr}_{\mathfrak m}A$ is a finitely generated algebra over $A/\mathfrak m$ and hence Noetherian by Lemma~\ref{alg:hilbert-basis}.
It is also the associated graded ring of $\widehat A$.
For an ideal $J\subset\widehat A$, choose $f_1,\ldots,f_s\in J$ whose initial forms generate its initial ideal.
For any $f\in J$, use those initial forms to cancel the lowest-degree terms, leaving a remainder of strictly higher order.
Repeated cancellation produces coefficients of the $f_i$ that converge in the $\mathfrak m$-adic topology.
Completeness supplies limits $a_i\in\widehat A$, and separatedness gives $f=\sum_i a_if_i$.
Thus every ideal is finitely generated.
Finally, successive Taylor truncations identify $\widehat{\C\{z\}}$ with $\C[[z]]$.
Apply exact completion to $0\to I\to\C\{z\}\to A\to0$ to obtain the quotient formula.
\end{proof}

\begin{theorem}[Analytic product projections are flat]\label{alg:analytic-product-flat}
For any complex analytic spaces $X,Y$, the projection $X\times Y\to Y$ is flat.
Consequently, locally trivial holomorphic maps are flat, and pullbacks along vector-bundle and projective-bundle projections preserve exact sequences.
\end{theorem}
\begin{proof}
Locally write the ring of the target $Y$ as $A$, the ring of the other factor $X$ as $T$, and the ring of their product as $B$:
\[
 A=\C\{z\}/I,\quad T=\C\{w\}/J,\quad
 B=\C\{z,w\}/(I(z),J(w)).
\]
We must prove that $B$ is flat over $A$.
Begin with an explicit linear-algebra calculation after completion.
In $\widehat T=\C[[w]]/J$, let $\mathfrak m_T$ denote the ideal generated by the images of the $w$ coordinates, and choose representatives $e_{d,1},\ldots,e_{d,r_d}$ of order $d$ whose classes form a $\C$-basis of
$\mathfrak m_T^d/\mathfrak m_T^{d+1}$.
Successive subtraction of initial terms gives a unique convergent formal sum
$\sum_{d,j}c_{d,j}e_{d,j}$.
Theorem~\ref{alg:krull-intersection} gives uniqueness and completeness gives existence.
Choose lifts of order at least $d$ for these representatives in $\C[[w]]$.
The decomposition then defines a continuous, filtration-preserving $\C$-linear section
$\widehat T\to\C[[w]]$.

Apply this decomposition to each coefficient in $z$ to obtain
\[
 \C[[z,w]]/J(w)\simeq\prod_{d,j}\C[[z]]
\]
as $\C[[z]]$-modules.
There is no requirement of a common analytic radius of convergence for the coefficients, since these are formal power series.
As $d$ increases, the $w$-order increases, so every such coefficient family defines a formal series.
Because $J$ has finitely many generators, having every $z$-coefficient in $J$ is equivalent to membership of the entire series in the extended ideal $J(w)$:
express each coefficient using those generators and collect by $z$-degree.
Likewise, finite generation of $I$ gives
$I\cdot\prod_{d,j}\C[[z]]=\prod_{d,j}I$.
Consequently,
\[
 \widehat B\simeq\prod_{d,j}\widehat A
\]
as $\widehat A$-modules.

Over a Noetherian ring, an arbitrary product of copies of the ring is flat.
Indeed, every ideal $K$ is finitely presented, and its finite presentation gives
$K\otimes_A\prod A\simeq\prod K$:
finite matrix kernels and images are computed coordinatewise in a product.
The natural map to $\prod A$ is therefore injective, so Lemma \ref{alg:flat-criteria} proves flatness.
Thus $\widehat B$ is flat over $\widehat A$, and hence over $A$ by Proposition~\ref{alg:completion-flat}.
The same proposition makes $B\to\widehat B$ faithfully flat.
If an injection developed a nonzero kernel after tensoring with $B$, tensoring further with $\widehat B$ would retain that kernel by exactness and faithfulness, contradicting $A$-flatness of $\widehat B$.
Therefore $B$ is flat over $A$.
A locally trivial map is locally such a projection; the stalk criterion and Proposition \ref{sh:flat-stalks} give the remaining assertions.
\end{proof}

For the following optional algebraic counterpart, equip $\operatorname{Spec}A$ with basic open sets $D(a)=\{\mathfrak p:a\notin\mathfrak p\}$ and the structure sheaf whose sections on $D(a)$ are $A[1/a]$.
This ringed space is an \emph{affine scheme}; a scheme is locally an affine scheme, and a scheme morphism is a morphism of the resulting locally ringed spaces (the stalk maps send the source maximal ideal into the target maximal ideal).
The stalk at $\mathfrak p$ is $A_{\mathfrak p}$, since passing to neighborhoods inverts precisely the elements outside $\mathfrak p$.
The next proof needs only this description of affine stalks.

\begin{proposition}[Local criteria for flatness of scheme morphisms]\label{alg:scheme-flatness}
For a scheme morphism $f:X\to Y$, the following are equivalent: every stalk-ring map is flat; the map $\cO_Y(V)\to\cO_X(U)$ is flat for every pair of affine open sets with $f(U)\subset V$; and such affine covers can be chosen with all corresponding ring maps flat.
Flatness may also be checked after choosing an arbitrary open cover and checking within its members.
Every open restriction of a flat morphism is flat, and every open immersion is flat.
\end{proposition}
\begin{proof}
We need only the facts that points of $\operatorname{Spec}A$ are prime ideals and the stalk at $\mathfrak p$ is $A_{\mathfrak p}$.
For a pair of affine opens, the stalk criterion is exactly the second assertion of Lemma \ref{alg:flat-criteria}.
Thus flatness of all stalk maps gives flatness for every affine pair, which implies the existence of a cover by such pairs.
Conversely, flatness on a covering gives flatness at every stalk.
The formulation with arbitrary open covers holds because restriction to an open neighborhood does not change the stalk at the point.
For an open immersion the corresponding stalk map is the identity, hence flat.
\end{proof}

\section{Finite Projections, Analytic Dimension, and Uniform Neighborhoods}
The model $w^2=z$ gives a finite map from a curve to the $z$-disc, with two sheets away from the branch point.
Finite projections generalize this picture to arbitrary analytic germs.
They reduce questions about dimension and singularities to finite modules over a power-series ring with fewer variables.
The radical of an ideal is $\sqrt I=\{a:a^k\in I\text{ for some }k\ge1\}$.
It equals the intersection of the prime ideals containing $I$: if $a\notin\sqrt I$, the nonzero ring $(A/I)[1/a]$ has a maximal ideal, whose inverse image is a prime containing $I$ but not $a$; the opposite containment follows from primeness.
A ring is \emph{reduced} if it has no nonzero nilpotent elements, equivalently its nilradical $\sqrt{(0)}$ is zero.
A finite-dimensional local $\C$-algebra has nilpotent maximal ideal: its descending powers stabilize as vector spaces, and Lemma~\ref{alg:nakayama} makes the stable power zero.
Such rings are examples of \emph{Artinian} rings, in which descending ideal chains stabilize.
A module has \emph{finite length} if it has a finite filtration with simple quotients; over the analytic local rings here the simple quotient is $\C$, so its length equals its complex dimension when finite.
For a homomorphism of analytic local rings, $A\{w\}$ denotes the convergent power-series algebra over $A$: if $A=\C\{y\}/I$, it means $\C\{y,w\}/I\C\{y,w\}$.


\begin{lemma}[The analytic finiteness criterion]\label{alg:analytic-finiteness}
Let $A\to B$ be a homomorphism of analytic local rings, with
$B=A\{w_1,\ldots,w_s\}/J$.
If $B/\mathfrak m_AB$ is finite-dimensional over $\C$, then $B$ is a finite $A$-module.
The corresponding analytic map has sufficiently small representatives that are finite, meaning proper with finite fibers; proper means that inverse images of compact sets are compact.
The direct image of a coherent sheaf under a finite map is coherent, and its stalk at a base point is the direct sum of the stalks over that finite fiber.
\end{lemma}
\begin{proof}
First prove finite generation, taking care that the elimination actually converges.
By the nilpotence observation above, choose $N$ large enough that every degree-$N$ monomial vanishes in the finite-dimensional local fiber algebra $B/\mathfrak m_AB$.
Lift the relations to $A\{w\}$. For every $|\alpha|=N$,
\[
 w^\alpha=\sum_{j=1}^m y_jh_{\alpha j}(y,w)
                 +\sum_{l=1}^t c_{\alpha l}(y,w)F_l(y,w),
\]
where the $y_j$ generate $\mathfrak m_A$ and the $F_l$ are finitely many defining relations.
Work in the convergent power-series ring before taking the quotient, including the defining relations of $A$ among the $F_l$.
Every power series admits a uniquely prescribed linear decomposition
\[
 g=R(g)+\sum_{|\alpha|=N}w^\alpha Q_\alpha(g),
 \qquad \deg_wR(g)<N.
\]
To prescribe it, assign each monomial of $w$-degree at least $N$, by a fixed rule, to a degree-$N$ monomial dividing it; place the lower-degree terms in $R$.

For the absolute coefficient norm
$\|g\|_{\epsilon,r}=\sum|g_{\beta\gamma}|\epsilon^{|\beta|}r^{|\gamma|}$,
we have $\|R\|\le1$ and $\|Q_\alpha\|\le r^{-N}$. The series of finite absolute coefficient norm form a Banach space, being a weighted $\ell^1$ space; thus the geometric operator series below converges in that norm.
First fix a sufficiently small $w$-radius $r$ on which the finitely many relations converge.
Then shrink the $y$-radius $\epsilon$ until
\[
 T(g)=\sum_{\alpha,j}y_jh_{\alpha j}Q_\alpha(g)
 \quad\text{satisfies}\quad \|T\|<\tfrac12.
\]
The series $\sum_{k\ge0}RT^kg$ therefore converges and is a combination of finitely many low-degree $w$-monomials, with convergent coefficients depending only on $y$.
The relations show that $g-Rg-Tg$ is a combination of the $F_l$, whose coefficients are bounded linear operators applied to $g$.
Summing after substituting $T^kg$, these coefficient series also converge geometrically.
Thus $g-\sum_{k\ge0}RT^kg$ belongs to the analytic ideal $(F_l)$ itself, not merely to its formal completion.
The $w$-monomials of degree less than $N$ generate $B$.
For any given germ, choose a smaller $r$ on which it converges and then choose $\epsilon$; the finite generating set remains the same.

For the geometric assertion, finiteness makes each $w_j$ satisfy a monic equation.
The central fiber algebra is Artinian local, so its maximal ideal is nilpotent.
A power of $w_j$ therefore lies in $\mathfrak m_AB$; the determinant trick applied to multiplication by this power on a finite generating set produces a monic equation that specializes at $y=0$ to $w_j^d=0$.
Fix a small vertical polydisc and shrink the base until all roots remain in its interior.
Over any compact subset of the base, the inverse image is closed and bounded, with no escape through the vertical boundary, hence compact.
The map is proper.
Each coordinate in a fiber belongs to the finite set of roots of its monic equation, so every fiber is finite.

Coherence of the direct image requires control of relations on the same neighborhood.
Let $P_j(y,w_j)$ be the monic coordinate equations just obtained, each involving only the base variables and one vertical variable.
Successive Weierstrass division represents the algebra modulo these equations as a finite free module on the base, with basis
$\prod_jw_j^{\alpha_j}$, $0\le\alpha_j<\deg P_j$.
The finitely many defining equations $F_l$ act as finite matrices in this algebra.
Their ideal is generated as a base module by the finitely many products of basis monomials with the $F_l$.
The actual direct image is therefore the cokernel of a finite holomorphic matrix and is coherent by Theorem~\ref{alg:oka} and Proposition~\ref{alg:coherent-category}.
For a general coherent sheaf, reduce the matrix coefficients of its finite presentation by the same division process to obtain a finite presentation on the base.
The local quotients at nearby fiber points agree with this finite algebra, as one verifies by applying the root-cluster decomposition of Lemma \ref{alg:finite-polynomial-model} one variable at a time.
Thus the construction computes the direct-image sheaf on a neighborhood, not only a finite module at the center.
For a given base point, choose disjoint neighborhoods of its finitely many fiber points.
Properness ensures that after shrinking the base its entire inverse image lies in their union.
Sections on these disjoint neighborhoods are independent, so the direct-image stalk is the finite direct sum of the source stalks.
All the arguments apply to coherent sheaves with finitely many module generators, as well as to the structure sheaf.
\end{proof}

An element is integral over $A$ if it satisfies a monic polynomial with coefficients in $A$.
A finite ring extension is integral: multiplication by an element on a finite generating list gives a matrix relation, and its adjugate produces a monic equation, as in the determinant argument of Lemma~\ref{alg:nakayama}.
For any ring, its Krull dimension is the supremum of lengths of strict chains of prime ideals, with a chain of $d+1$ primes having length $d$.

\begin{lemma}[Dimension under finite integral extension]\label{alg:integral-dimension}
For a finite ring extension $A\hookrightarrow B$, we have $\dim A=\dim B$, where dimension is the supremum of lengths of strict chains of prime ideals.
The ring $R_n$ has dimension $n$, and every nonzero ideal $I\subset R_n$ satisfies $\dim R_n/I\le n-1$.
\end{lemma}
\begin{proof}
We establish three properties of prime ideals.
For $\mathfrak p\subset A$, the finite module $B_{\mathfrak p}$ is nonzero and
$\mathfrak pB_{\mathfrak p}\ne B_{\mathfrak p}$.
Otherwise Lemma~\ref{alg:nakayama} would make it zero, contradicting the injection of $A_{\mathfrak p}$.
A maximal ideal of the nonzero finite-dimensional algebra
$B_{\mathfrak p}/\mathfrak pB_{\mathfrak p}$ pulls back to a prime of $B$ lying over $\mathfrak p$.
This is \emph{lying over}.
Apply it to the quotient extension $A/\mathfrak p\hookrightarrow B/\mathfrak q$ to lift a chain in $A$ step by step to $B$; this is \emph{going up}.

If $\mathfrak q_1\subset\mathfrak q_2$ contract to the same $\mathfrak p$, first quotient by $\mathfrak q_1$, then invert the nonzero elements of $A/\mathfrak p$.
The resulting domain is integral over a field and is itself a field:
a monic equation for a nonzero element, after removing any zero constant-term factor, expresses its inverse.
There can therefore be no surviving nonzero proper prime.
Localization is injective on the domain, so $\mathfrak q_1=\mathfrak q_2$.
Strict chains in $B$ consequently remain strict after contraction.
Going up gives $\dim B\ge\dim A$, and this \emph{incomparability} gives the reverse inequality.

Induct on $n$ for the dimension of $R_n$.
The chain
$(0)\subset(z_1)\subset(z_1,z_2)\subset\cdots\subset\mathfrak m_n$
gives the lower bound $n$.
If $I=R_n$, the quotient is zero and there is nothing to prove; otherwise take a nonzero element of the proper ideal $I$, choose regular coordinates, and let $W$ be its Weierstrass polynomial.
Theorem~\ref{alg:weierstrass-division} makes $R_n/(W)$ finite free of positive rank, hence faithful, over $R_{n-1}$; the first part of this lemma gives dimension $n-1$.
Its quotient $R_n/I$ has dimension at most $n-1$.
Any prime chain in $R_n$, after removing a possible initial zero ideal, descends to a chain in a quotient by a nonzero ideal and therefore has total length at most $n$.
This completes the induction.
\end{proof}

For an analytic space $Z$, define $\dim_xZ=\dim\cO_{Z,x}$ and $\dim Z=\sup_x\dim_xZ$, with $\dim\varnothing=-\infty$.
A germ $(Z,x)$ records the analytic space on an unspecified sufficiently small neighborhood of $x$.
The fiber $X_y$ of a holomorphic map $f:X\to Y$ has local ring
$\cO_{X,x}/\mathfrak m_y\cO_{X,x}$ at a point $x$ over $y$; here $\mathfrak m_y$ is the maximal ideal of $\cO_{Y,y}$.

\begin{theorem}[Finite projections of local analytic spaces]\label{alg:finite-projection}
After a linear coordinate change, every nonempty analytic space germ $(Z,0)$ embedded in $(\C^N,0)$, including a nonreduced space, admits a finite projection representative to $\C^d$, where $d=\dim\cO_{Z,0}$.
Its local ring is a finite faithful module over $\C\{t_1,\ldots,t_d\}$.
If the fiber of $f:(X,x)\to(Y,y)$ has dimension $d$ at $x$, one can choose $d$ linear functions $t$ such that
$(f,t):(X,x)\to(Y\times\C^d,(y,0))$ is finite.
\end{theorem}
\begin{proof}
For the algebraic assertion, induct on the ambient dimension $N$.
If the defining ideal is zero, use the identity projection.
Otherwise choose a nonzero defining equation, make it regular in the last coordinate, and apply Weierstrass division.
This makes $\cO_{Z,0}$ a finite $R_{N-1}$-module.
If $J$ is the kernel of $R_{N-1}\to\cO_{Z,0}$, the extension of $R_{N-1}/J$ is finite and faithful.
Continue by induction to obtain a finite faithful extension of a convergent power-series ring $R_d$.
At each stage change only the remaining coordinates linearly, so the final base coordinates are linear functions of the original ambient coordinates.
Lemma \ref{alg:integral-dimension} identifies $d$ with the dimension of the local ring.

The central fiber of a finite extension is a finite-dimensional local algebra.
Lemma \ref{alg:analytic-finiteness} therefore gives a proper finite representative of the projection.
Its image is the whole small base.
Indeed, the finite direct-image module has a finite presentation, and its support is defined by maximal minors.
Each such minor annihilates the whole module, as follows by applying the adjugate identity to the selected square block.
Faithfulness makes the annihilator zero, so all those minors vanish identically and the support is the whole base.
Equivalently, the zeroth Fitting ideal (the ideal of these maximal presentation minors) is contained in the annihilator and its zero set is the support; the containment is precisely the same columnwise adjugate calculation.

Finally, apply the first assertion to the fiber ring
$\cO_{X,x}/\mathfrak m_y\cO_{X,x}$.
Choose $t_1,\ldots,t_d$ so that further quotienting by these functions gives a finite-dimensional complex vector space.
The finiteness criterion applied to
$\cO_{Y\times\C^d,(y,0)}\to\cO_{X,x}$
gives the required relative finite map.
\end{proof}

For a monic polynomial with roots $\lambda_1,\ldots,\lambda_r$ in a splitting field, its discriminant is $\prod_{\mu<\nu}(\lambda_\mu-\lambda_\nu)^2$.
This symmetric expression is a polynomial in its coefficients and is nonzero exactly when the roots are distinct.

\begin{lemma}[Finite reduced algebras away from a discriminant]\label{alg:generic-graphs}
Let $D=\C\{t_1,\ldots,t_d\}$ inject into the reduced local algebra $B$ of a finite analytic projection, with $B$ finite over $D$.
There are a nonzero $h\in D$ and a monic polynomial $P\in D[1/h][T]$ with invertible discriminant such that
\[
 B[1/h]\simeq D[1/h][T]/(P).
\]
On sufficiently small analytic representatives, the projection over $h\ne0$ is therefore locally a finite union of disjoint holomorphic graphs.
\end{lemma}
\begin{proof}
Put $K=\operatorname{Frac}(D)$ and $B_K=B\otimes_DK$.
This is a nonzero finite-dimensional reduced commutative $K$-algebra.
Each prime quotient is a finite-dimensional domain over $K$ and hence a field: multiplication by a nonzero element is injective and therefore surjective.
There are finitely many maximal ideals, since the Chinese remainder map for any finite list of distinct maximal ideals is surjective and their number is bounded by $\dim_K B_K$.
Their intersection is the nilradical, which is zero.
The same Chinese remainder construction as in Lemma~\ref{alg:finite-polynomial-model} writes $B_K$ as a finite product of finite field extensions of $K$.

These extensions are separable: in characteristic zero an irreducible polynomial has nonzero derivative and is coprime to that derivative, so its roots are distinct.
We spell out the single-generator step.
Choose finitely many algebra generators and a field in which their minimal polynomials split. In a product of fields a minimal polynomial is the least common multiple of the minimal polynomials in the factors, so it too has distinct roots.
The commuting multiplication matrices of the generators are diagonalizable because those polynomials have distinct roots, so they are simultaneously diagonalizable: diagonalize one and restrict the others to its invariant eigenspaces, then continue by induction.
Their regular representation, meaning multiplication on the algebra itself, is faithful and identifies the extended algebra with a product of $r=\dim_K B_K$ copies of the splitting field.
Indeed, if there were fewer than $r$ distinct joint eigenspaces, all multiplication matrices would be determined by fewer than $r$ eigenvalues, contradicting the algebra's dimension $r$.
Since $K$ is infinite, choose a $K$-linear combination $\theta$ of the generators whose values in these $r$ factors are distinct, avoiding the finitely many linear equations that equate two values.
Writing $\theta_1,\ldots,\theta_r$ for these values, the Vandermonde determinant $\prod_{\mu<\nu}(\theta_\nu-\theta_\mu)\ne0$ shows that $1,\theta,\ldots,\theta^{r-1}$ form a basis after extension of scalars and hence already over $K$.
Thus $B_K=K[T]/(P)$ for the monic minimal polynomial $P$ of $\theta$, and the discriminant (the squared product of differences of its roots) is nonzero.

Clear the denominators of $\theta$ and of the coefficients of $P$. The element $P(\theta)$ becomes zero over $K$, so also include a nonzero annihilator of this element in $h$; evaluation then defines a homomorphism on the polynomial quotient.
The resulting map from $D[1/h][T]/(P)$ to $B[1/h]$ is generically an isomorphism between finite modules over a Noetherian ring.
Its kernel and cokernel are finite and become zero over $K$; each of their finitely many generators is killed by some nonzero element of $D$.
Include the product of those elements, and the numerator of the discriminant, in $h$.
This gives the displayed isomorphism and the invertible discriminant.
This step also removes any $D$-torsion that vanished on passing from $B$ to $B_K$.

The finitely many coefficient identities hold on fixed small analytic representatives.
Over $h\ne0$, the roots of $P$ are simple, so Theorem~\ref{bas:holomorphic-inverse} writes each root locally as a holomorphic function of the base variables.
All other algebra generators are polynomial expressions in this root with holomorphic coefficients there.
These expressions give the claimed graphs, and the distinct roots make them disjoint.
\end{proof}

\begin{proposition}[The local analytic Nullstellensatz and reduced structures]\label{alg:analytic-nullstellensatz}
For an ideal $I\subset\C\{z_1,\ldots,z_N\}$, the function germs vanishing on its zero-set germ are exactly
$\sqrt I=\{g:g^k\in I\text{ for some }k\}$.
The vanishing ideal sheaf of an analytic set is coherent, and its quotient defines the reduced structure sheaf.
\end{proposition}
\begin{proof}
The radical of an ideal in a Noetherian ring is a finite intersection of prime ideals.
To prove this finiteness, suppose some radical ideal fails to be such an intersection, and choose a maximal one $J$ using the ascending chain condition.
It is not prime, so choose $a,b\notin J$ with $ab\in J$.
Expansion of powers directly verifies
$J=\sqrt{J+(a)}\cap\sqrt{J+(b)}$.
Both ideals on the right strictly contain $J$ and thus have finite prime decompositions by maximality, a contradiction.
Removing redundant factors leaves finitely many minimal primes.

First suppose $I=\mathfrak p$ is prime.
Finite projection makes the domain $B=R_N/\mathfrak p$ a finite torsion-free module over $D=\C\{t_1,\ldots,t_d\}$.
Choose a basis over the fraction field and clear denominators to obtain an injection $B\hookrightarrow D^r$.
Lemma~\ref{alg:analytic-finiteness} makes the direct image coherent, and Proposition~\ref{alg:coherent-support} extends this injection to a fixed smaller neighborhood.
Over the fraction field of $D$, the algebra $B$ is a finite field extension, separable in characteristic zero.
Lemma~\ref{alg:generic-graphs}, which uses only finite projection and finite-dimensional algebra, gives after inverting a nonzero $h$ an algebra defined by a polynomial with distinct roots.
Geometrically it is locally a union of holomorphic graphs of simple roots.
If $g$ vanishes on the original zero set, it vanishes on all these graphs and is therefore zero in $B[1/h]$.
Torsion-freeness gives $g=0$ in $B$, or $g\in\mathfrak p$.
The same argument shows that a representative of this prime germ can be chosen reduced:
a nilpotent section vanishes on the graph locus, and the injection into $\cO^r$ excludes any nonzero section supported only on $h=0$.

In general, write $\sqrt I=\bigcap_{j=1}^s\mathfrak p_j$.
Its zero set is the union of the zero sets of these primes.
Indeed, at a point outside every such zero set, choose $f_j\in\mathfrak p_j$ nonzero at the point; their product belongs to the intersection and remains nonzero there.
A function $g$ vanishing on $Z(I)$ therefore lies in every $\mathfrak p_j$, hence in $\sqrt I$.
The converse follows immediately from $g^k\in I$.

Finally, choose finite generators of $\sqrt I$ at the center and let them define a coherent ideal sheaf $\cJ$ nearby.
The central injection
$R_N/\sqrt I\hookrightarrow\bigoplus_jR_N/\mathfrak p_j$
remains injective after shrinking because its kernel is coherent and has zero central stalk.
The representatives on the right can already be chosen reduced, so $\cO/\cJ$ is reduced on the whole small neighborhood.
Applying the just-proved Nullstellensatz at each nearby point identifies $\cJ$ with all germs vanishing on the zero set.
Thus $\cJ$ is its vanishing ideal sheaf, proving coherence and existence of the reduced structure.
\end{proof}

\begin{corollary}[Coherent sheaves with zero-dimensional support]\label{sh:zero-dimensional-coherent}
If a coherent sheaf $\cF$ on a compact complex manifold has zero-dimensional support, then $\cF$ is flabby.
For any module sheaf $\cG$, the tensor product $\cF\otimes\cG$ is also flabby.
\end{corollary}
\begin{proof}
Proposition \ref{alg:coherent-support} ensures that the support is a closed analytic set.
A zero-dimensional analytic set has isolated points locally by Theorem~\ref{alg:finite-projection}, which gives a finite projection to a point.
One can also see discreteness directly from its local ring: let $(A,\mathfrak m)$ be the local ring of the support. This is a zero-dimensional Noetherian local ring, so its only prime ideal is $\mathfrak m$ and $\mathfrak m=\sqrt{(0)}$.
Finite generation then implies $\mathfrak m^N=0$ for some $N$.
The local defining ideal therefore contains a power of each coordinate, and its zero set in a sufficiently small neighborhood consists only of the center.
To justify the ring-theoretic step, a radical ideal in a Noetherian ring has finitely many generators, each with some power in the original ideal.
Expanding a sufficiently high ideal power gives a uniform exponent $N$.
A zero-dimensional local ring has only one prime ideal, its maximal ideal.
The support is thus closed and discrete, and compactness makes it finite.
The stalk of a tensor product is the tensor product of stalks, so
$\supp(\cF\otimes\cG)\subset\supp\cF$.
Apply Proposition~\ref{sh:discrete-support} to both sheaves; its proof uses only closed discreteness, independently of the analytic arguments here.
\end{proof}

An analytic space is reduced when its local rings are reduced.
An analytic set is irreducible if it is not a union of two proper analytic subsets; its local irreducible components correspond to the minimal prime ideals, by Proposition~\ref{alg:analytic-nullstellensatz}.
It is pure $d$-dimensional when every local irreducible component has dimension $d$.
A point is regular if a neighborhood is a complex submanifold of the ambient manifold; the singular locus consists of the remaining points.
An unramified covering here means a finite covering whose local inverse branches are biholomorphic.
A smooth local stratum is a locally closed smooth submanifold in a chosen coordinate neighborhood.

\begin{proposition}[Density of the regular locus and dimension drop of the exceptional set]\label{alg:analytic-stratification}
A reduced pure $d$-dimensional analytic set admits locally a finite projection to a $d$-dimensional polydisc.
There is a nonzero holomorphic function $h$ on the base such that over $h\ne0$ the projection is an unramified finite covering, locally a union of finitely many holomorphic graphs.
The exceptional subset has dimension at most $d-1$.
The regular locus is therefore dense, and the singular locus is a closed analytic subset of strictly smaller dimension.
On every relatively compact neighborhood, the analytic set can be covered by finitely many smooth local strata.
\end{proposition}
\begin{proof}
Work in an ambient coordinate neighborhood $Z\subset\C^N$. Use Theorem~\ref{alg:finite-projection} and write $D=\C\{t_1,\ldots,t_d\}$ and $B=\cO_{Z,0}$.
Lemma~\ref{alg:generic-graphs} gives a nonzero $h\in D$ and holomorphic graph descriptions over $h\ne0$.
These points are smooth, and the projection is locally biholomorphic.

Purity means that every component corresponding to a minimal prime of $B$ has dimension $d$.
Such a prime contracts to zero in $D$: otherwise Lemma \ref{alg:integral-dimension} would bound that component's dimension by $d-1$.
Thus $h$ lies in no minimal prime, and $h\ne0$ is dense in every component.
The quotient $B/(h)$ is finite over $D/(h)$, so its dimension is at most $d-1$.

Choose finitely many generators of the reduced defining ideal locally.
If their Jacobian has rank $N-d$, Theorem~\ref{bas:holomorphic-inverse} gives a smooth $d$-dimensional zero set containing $Z$.
If $Z$ were a proper analytic subset of it, a nonzero equation and Lemma \ref{alg:integral-dimension} would bound its dimension by $d-1$.
Thus the two agree locally.
Conversely, at a smooth point of dimension $d$, the $N-d$ transverse coordinates generate the ideal of all vanishing functions; this follows by Taylor integration in the transverse variables.
For a reduced analytic structure this is the defining ideal by the already proved Proposition~\ref{alg:analytic-nullstellensatz}.
The differentials of the chosen generators therefore span the conormal space, and the Jacobian rank is $N-d$.
The singular locus is consequently defined by the vanishing of all $(N-d)$th-order minors and is closed analytic.
Repeat Theorem~\ref{alg:finite-projection} on the singular locus; dimension drops by at least one each time.
Noetherianity gives finitely many irreducible components of each germ: repeated decomposition into two proper analytic subgerms would otherwise produce a strictly increasing chain of defining ideals.
A relatively compact region admits finitely many such local descriptions, yielding finitely many smooth local strata.
Strata may be separately indexed in different coordinate patches; no Whitney conditions are asserted here.
\end{proof}

\begin{lemma}[A uniform criterion for vanishing germs of coherent subsheaves]\label{alg:uniform-germ}
For a coherent sheaf $\cF$ near $x$, one can choose fixed neighborhoods $V\subset U$ such that every coherent subsheaf $\mathcal K\subset\cF$ defined on $U$ satisfies:
if $\mathcal K_x=0$, then $\mathcal K|_V=0$.
The neighborhood $V$ depends only on $\cF$ and $x$, not on $\mathcal K$.
\end{lemma}
\begin{proof}
First construct a finite prime filtration.
In a nonzero Noetherian module $M$, choose a nonzero $v$ whose annihilator is maximal among annihilators of nonzero elements.
If $abv=0$ but $bv\ne0$, then
$\operatorname{Ann}(bv)\supset\operatorname{Ann}(v)$.
Maximality makes the two equal, giving $av=0$.
Thus $\operatorname{Ann}(v)$ is prime and $Av\simeq A/\mathfrak p$.
Repeat on the quotient; Noetherianity makes the ascending submodule chain terminate.
Hence $M=\cF_x$ has a finite filtration whose quotients are $\cO_{X,x}/\mathfrak p_j$.
Represent the finitely many generators and relations on a common neighborhood and use the closed-support properties of coherent kernels and cokernels to shrink until the filtration is exact.
This gives a finite filtration of $\cF$ by coherent subsheaves, with quotients $\mathcal Q_j$ representing prime analytic germs.

For one such prime factor, Theorem~\ref{alg:finite-projection} makes its central local ring $B$ finite over
$D=\C\{t_1,\ldots,t_d\}$ and torsion-free, since $B$ is a domain.
Choose a basis of $B\otimes_D\operatorname{Frac}(D)$ and clear denominators in a finite generating set to obtain a $D$-linear injection $B\hookrightarrow D^r$.
The finite direct image is coherent, so these finitely many linear formulas define a morphism on a base neighborhood.
Its kernel vanishes at the center and hence vanishes after shrinking.
Writing $\pi$ for this finite projection, on a fixed representative $\pi_*\mathcal Q_j$ embeds in a free sheaf $\cO^r$.

If a coherent subsheaf $\mathcal K'\subset\mathcal Q_j$ has $\mathcal K'_x=0$, it first vanishes on some neighborhood that may depend on $\mathcal K'$.
The central fiber is just $x$, so properness makes $\pi_*\mathcal K'$ vanish near the base center.
It is also a coherent subsheaf of $\cO^r$.
On a connected base, the support of a nonzero coherent subsheaf of a free sheaf is both closed and open.
Closedness follows from coherence.
At a point with nonzero stalk, choose a nonzero holomorphic vector-function representative in the subsheaf; on a small connected neighborhood Corollary~\ref{bas:identity} makes its germ nonzero at every point, proving openness of the support.
Vanishing near the center therefore forces vanishing on the whole fixed connected base.
Since finite direct-image stalks are finite direct sums of source stalks, $\mathcal K'$ vanishes on the whole fixed projection representative.
This supplies a neighborhood $V_j$ for each prime factor, independent of the subsheaf.

For an arbitrary $\mathcal K\subset\cF$, intersect it with the filtration terms.
The resulting quotients are coherent subsheaves of the corresponding $\mathcal Q_j$, with zero stalk at $x$.
Each vanishes on $V_j$.
On the finite intersection $V=\bigcap_jV_j$, successive filtration terms agree, and therefore $\mathcal K|_V=0$.
\end{proof}

\begin{corollary}[Arbitrary families holomorphic on one common open set]\label{alg:holomorphic-family}
If all $\{f_\lambda\}_{\lambda\in\Lambda}$ are holomorphic on the same open set $U$, their generated ideal sheaf $\sum_\lambda\cO_Uf_\lambda$ is coherent.
\end{corollary}
\begin{proof}
Fix $x\in U$.
Noetherianity provides finitely many family members $f_{\lambda_1},\ldots,f_{\lambda_s}$ whose germs generate the ideal generated by all the germs $f_\lambda$.
Let $\cI$ be the coherent ideal sheaf generated by these finitely many functions, and put $\mathcal Q=\cO_U/\cI$.
Each $f_\lambda$ generates a cyclic coherent subsheaf $\mathcal K_\lambda$ of $\mathcal Q$ with zero stalk at $x$.
Lemma \ref{alg:uniform-germ} supplies one neighborhood $V$, \emph{the same for every $\lambda$}, on which all $\mathcal K_\lambda$ vanish.
All $f_\lambda$ thus belong to $\cI$ on $V$, and the original ideal sheaf equals $\cI$ there.
It is locally finitely generated, so Oka's theorem makes it coherent.
\end{proof}

\begin{corollary}[Strong Noetherianity on compact spaces]\label{alg:strong-noetherian}
On a compact complex analytic space, every increasing chain of coherent subsheaves of a coherent sheaf stabilizes after finitely many steps.
In particular, every increasing chain of coherent ideal sheaves stabilizes.
\end{corollary}
\begin{proof}
Let $\cF_1\subset\cF_2\subset\cdots\subset\cF$.
Each stalk $\cF_x$ is a Noetherian module, so there is $N_x$ with
$(\cF_j)_x=(\cF_{N_x})_x$ for all $j\ge N_x$.
In the coherent quotient $\cF/\cF_{N_x}$, all the subsheaves
$\cF_j/\cF_{N_x}$ have zero stalk at $x$.
The uniform vanishing criterion gives a common neighborhood $V_x$ on which they all vanish.
Compactness gives a finite cover $V_{x_1},\ldots,V_{x_l}$.
Set $N=\max_iN_{x_i}$.
For $j\ge N$, the sheaves $\cF_j$ and $\cF_N$ agree on every cover member and hence globally.
\end{proof}

For the following volume comparison, set
$\omega_{\C^m}=\frac{\ii}{2}\sum_{j=1}^m dz_j\wedge d\bar z_j$ and $dV_{\C^m}=\omega_{\C^m}^m/m!$.
The complex Grassmannian $\operatorname{Gr}_{\C}(d,N)$ is the space of complex $d$-planes in $\C^N$.
It identifies with the closed bounded set of Hermitian projections of rank $d$, so it is compact, by the same matrix description used in Proposition~\ref{bas:projective-space}.

\begin{lemma}[Perturbing finite projections and controlling volume]\label{alg:finite-projections-dominate}
If a linear projection $P_0:\C^N\to\C^d$ is finite on the analytic germ $(Z,0)$, then every sufficiently small linear perturbation is also finite on the germ.
Among these projections one can choose finitely many $P_1,\ldots,P_s$ and a constant $c>0$ such that, on every complex $d$-dimensional tangent plane,
\[
 \sum_{j=1}^s P_j^*\frac{\omega_{\C^d}^d}{d!}
 \ \ge\ c\,\frac{\omega_{\C^N}^d}{d!}.
\]
Thus bounds on the numbers of sheets of these finitely many projections control the local volume of the regular locus.
\end{lemma}
\begin{proof}
In coordinates $(y,w)$ adapted to $P_0$, Lemma~\ref{alg:analytic-finiteness} gives, for a fixed elimination degree $\ell$, relations
$w^\alpha=\sum y_jh_{\alpha j}+\sum c_{\alpha l}F_l$.
For a nearby projection $u=P(y,w)$, solve for
$y=A_Pu+B_Pw$, with $A_P$ bounded and $B_P\to0$.
In a fixed absolute coefficient norm, the additional terms in the elimination operator have norm at most
$C\|B_P\|r^{1-\ell}$, where $\ell$ is the monomial degree used in that criterion, not the ambient dimension $N$.
First choose $P$ close enough to $P_0$, then shrink the $u$-radius.
The total operator norm remains less than $1/2$.
The same geometric-series argument proves finite generation over the $u$ variables, so the perturbed projection remains finite.
This establishes stability of the convergent construction, beyond mere set-theoretic isolation of one fiber.

Fix a complex $d$-plane $L$.
The expression $\det(P|_L)$ is a polynomial in the entries of $P$, and is not the zero polynomial, since some $P$ restricts to an isomorphism on $L$.
It cannot vanish throughout the open neighborhood of $P_0$.
Choose $P_L$ there for which $P_L^*dV$ is positive on $L$.
This positivity persists in a Grassmannian neighborhood of $L$.
Compactness of the complex Grassmannian gives finitely many such neighborhoods covering it, and corresponding $P_1,\ldots,P_s$.
Compute $\det(P_j|_L)$ using an orthonormal basis of $L$; its squared modulus is independent of that choice. The sum $\sum_j|\det(P_j|_L)|^2$ is continuous and positive on the compact Grassmannian, so its minimum is some $c>0$.
This is the asserted form inequality.
Integrating over the regular locus gives the volume-control assertion; the sheet-counting integral is detailed in the section on integration currents of analytic sets.
\end{proof}

\begin{lemma}[Numbers of parameters and dimension bounds]\label{alg:parameter-dimension}
If $r$ elements of an analytic local ring $(B,\mathfrak m)$ generate an ideal with radical $\mathfrak m$, then $\dim B\le r$.
Consequently, a subspace defined by $r$ equations has dimension at least $\dim B-r$.
\end{lemma}
\begin{proof}
First derive explicit growth estimates from Theorem~\ref{alg:finite-projection}.
If $\dim B=d$, choose a finite faithful subring
$D=\C\{t_1,\ldots,t_d\}\subset B$ and put $\mathfrak n=(t_1,\ldots,t_d)$.
If $B$ has $s$ generators over $D$, then
\[
 \dim_\C B/\mathfrak m^k
 \le\dim_\C B/\mathfrak n^kB
 \le s\binom{k+d-1}{d}.
\]
The finite central fiber gives $a$ with $\mathfrak m^a\subset\mathfrak nB$.
Theorem~\ref{alg:artin-rees} applied to the inclusion of finite $D$-modules $D\subset B$ gives $c$ with
$D\cap\mathfrak n^lB\subset\mathfrak n^{l-c}D$.
Thus, for large $l$,
\[
 \dim_\C B/\mathfrak m^{al}
 \ge \dim_\C D/(D\cap\mathfrak m^{al})
 \ge \binom{l-c+d-1}{d}.
\]
Monotonicity in $k$ extends this lower bound to all sufficiently large $k$.
Hence $\dim_\C B/\mathfrak m^k$ is bounded above and below by positive constants times $k^d$.

If $I=(a_1,\ldots,a_r)$ has radical $\mathfrak m$, finite generation of $\mathfrak m$ gives $\mathfrak m^a\subset I$ for some $a$, so $l_0=\dim_\C(B/I)<\infty$.
For $r=0$, the maximal ideal is nilpotent and $B$ is zero-dimensional, giving the assertion directly.
For $r>0$, the quotient $I^j/I^{j+1}$ is generated over $B/I$ by the degree-$j$ monomials in $r$ variables, so its complex dimension is at most
$l_0\binom{j+r-1}{r-1}$.
Summing the layers gives
\[
 \dim_\C B/I^k\le l_0\binom{k+r-1}{r}.
\]
Since $I^k\subset\mathfrak m^k$, the preceding $k^d$ lower bound cannot exceed this $k^r$ upper bound.
Therefore $d\le r$.
If $J$ has $r$ generators, put $e=\dim B/J$.
Choose the $e$ finite-projection parameters of $B/J$ and lift them to $B$.
Together with the $r$ generators of $J$, they generate an ideal with radical $\mathfrak m$.
Thus $\dim B\le e+r$, proving the last assertion.
\end{proof}

\begin{theorem}[The dimension formula for flat maps]\label{alg:flat-dimension}
If $f:X\to Y$ is a flat map of complex analytic spaces, then for every $x$ and $y=f(x)$,
\[
 \dim_xX=\dim_yY+\dim_xX_y.
\]
\end{theorem}
\begin{proof}
Put $A=\cO_{Y,y}$, $B=\cO_{X,x}$, $r=\dim A$, and
$d=\dim B/\mathfrak m_AB$.
Theorem~\ref{alg:finite-projection} gives $r$ parameters $a_i$ in $A$ generating an ideal with radical $\mathfrak m_A$.
Choose $d$ parameters in the fiber ring and lift them to $b_j\in B$.
Then $(a_1,\ldots,a_r,b_1,\ldots,b_d)$ has radical $\mathfrak m_B$.
Lemma \ref{alg:parameter-dimension} gives $\dim B\le r+d$.

For the reverse inequality, prove the required going-down property of flatness.
A flat local homomorphism $A\to B$ is faithfully flat:
a nonzero module contains a nonzero cyclic submodule $A/I$, while
$IB\subset\mathfrak m_B$, so $B/IB\ne0$.
Flatness embeds this nonzero quotient into the tensor product of the original module.
If $\mathfrak p'\subset\mathfrak p=\mathfrak q\cap A$, then
$A_{\mathfrak p}\to B_{\mathfrak q}$ is also flat and local, hence faithful.
Therefore
$B_{\mathfrak q}\otimes_{A_{\mathfrak p}}\kappa(\mathfrak p')\ne0$.
Choose a prime in this nonzero ring and pull it back to obtain
$\mathfrak q'\subset\mathfrak q$ contracting exactly to $\mathfrak p'$.

Take a prime chain of length $d$ in the fiber ring and lift it to
$\mathfrak q_0\subsetneq\cdots\subsetneq\mathfrak q_d=\mathfrak m_B$,
all contracting to $\mathfrak m_A$.
Choose a chain of length $r$ in $A$ ending at $\mathfrak m_A$.
Starting from $\mathfrak q_0$, apply going down successively to append a chain below it, producing a chain of length $r+d$ in $B$.
Hence $\dim B\ge r+d$, proving equality.
\end{proof}

\begin{theorem}[Flat analytic maps are open]\label{alg:flat-open}
A flat holomorphic map of complex analytic spaces is an open map.
\end{theorem}
\begin{proof}
Fix $x\in X$ and any open neighborhood $U$ of $x$.
We must show that $f(U)$ contains a neighborhood of $y=f(x)$.
Let $d=\dim_xX_y$.
The relative assertion of Theorem~\ref{alg:finite-projection} gives a representative inside $U$ for which
$g=(f,t)$ maps finitely to $Y\times D^d$, with the central fiber projecting finitely onto all of $D^d$.
Set $\mathcal M=g_*\cO_X$, a coherent module sheaf.
On the central fiber, $\mathcal M/\mathfrak m_y\mathcal M$ is a finite module sheaf on $D^d$ supported on the whole polydisc.
Outside the zero set of some nonzero holomorphic function $h$, it is free of positive rank:
choose a basis over the fraction field, clear the finitely many denominators, and use closedness of kernel and cokernel supports.
We may thus choose $t'$ arbitrarily close to zero where this fiber module is free and nonzero.

Put $C=\cO_{Y\times D^d,(y,t')}$, $M=\mathcal M_{(y,t')}$, and $A=\cO_{Y,y}$.
A finite direct-image stalk is the finite direct sum of the source stalks, each of which is $A$-flat by hypothesis.
Hence $M$ is $A$-flat.
Lift a basis of $M/\mathfrak m_AM$ over $C/\mathfrak m_AC$ to a map
$\phi:C^r\to M$.
Its cokernel is finite and vanishes modulo $\mathfrak m_A$, so Lemma~\ref{alg:nakayama} makes $\phi$ surjective.
Let $K$ be its kernel.
Since the quotient $M$ is $A$-flat, tensoring with $A/\mathfrak m_A$ leaves an exact sequence
\[
0\to K/\mathfrak m_AK\to(C/\mathfrak m_AC)^r
 \xrightarrow{\bar\phi}M/\mathfrak m_AM\to0.
\]
We justify left exactness here.
Generally, if $0\to K\to N\to M\to0$ has flat quotient $M$, tensoring with any module $Q$ preserves injectivity at the left.
Choose a free presentation $0\to R\to F\to Q\to0$ and lift an element in the proposed kernel to $K\otimes F$.
Its image in $N\otimes F$ comes from some $w\in N\otimes R$.
The image of $w$ in $M\otimes R$ maps to zero in $M\otimes F$, so is zero by flatness.
Right exactness therefore lifts $w$ from $K\otimes R$.
Subtract this lift; the adjusted element of $K\otimes F$ maps to zero in $N\otimes F$.
Freeness of $F$ makes $K\otimes F\to N\otimes F$ injective, so the original element is already zero in $K\otimes Q$.
Apply this chase with $Q=A/\mathfrak m_A$ to obtain the displayed exact sequence.
Since $\bar\phi$ is an isomorphism, $K/\mathfrak m_AK=0$.
The module $K$ is finite, so Lemma~\ref{alg:nakayama} again gives $K=0$.

Thus $M\simeq C^r$ with $r>0$.
Coherence extends this isomorphism to a neighborhood of $(y,t')$, on which $\mathcal M$ has full support.
The image of a finite map is the support of the direct image of its structure sheaf, so the image of $g$ contains that neighborhood.
Projecting to $Y$ shows that $f(U)$ contains an open neighborhood of $y$.
Since $x,U$ were arbitrary, $f$ is open.
\end{proof}

\begin{lemma}[Freeness for finite maps between smooth germs]\label{alg:finite-smooth-flat}
For a finite holomorphic map $(\C^n,0)\to(\C^n,0)$, the source local ring is a finite free module over the target local ring.
\end{lemma}
\begin{proof}
Write $B=\C\{x_1,\ldots,x_n\}$ for the source ring and $g_1,\ldots,g_n$ for the images of the target coordinates.
Finiteness says that $(g_1,\ldots,g_n)$ has radical equal to the source maximal ideal.
We first prove positive-degree exactness of the Koszul complex of $g$; counting equations alone would not prove this.

Consider the coordinate localization complex
\[
 C^q(B)=\bigoplus_{|I|=q}B[1/x_i:i\in I],\qquad 0\le q\le n,
\]
with $B$ in the term $I=\varnothing$ and alternating restriction differential.
Its cohomology can occur only in degree $n$.
To see this directly, each localized element has a convergent Laurent expansion whose negative exponents have a uniform finite lower bound.
Group terms according to the set $S$ of coordinates with negative exponents; there are only $2^n$ groups.
If $S\ne\{1,\ldots,n\}$, choose $j\notin S$ and contract that group by the alternating operator inserting or deleting $j$ from the localization index set.
A termwise calculation gives $dh+hd=\operatorname{id}$.
Deleting $j$ is permitted because there are no negative powers of $x_j$ in the group.
The projections preserve convergence and finite pole orders, so the contraction acts on the actual localization rings.
Only the group with $S$ equal to all coordinates occurs solely in degree $n$ and cannot be contracted.

Form the double complex of $C^\bullet(B)$ with the Koszul chain complex $K_\bullet(g;B)$, using total degree $q-p$.
First take cohomology in the localization direction.
Every Koszul term is free, so only the row $q=n$ survives.
Since $p\le n$, total cohomology is zero in negative degrees.
In the other direction, every Koszul homology module is annihilated by all $g_i$, because exterior multiplication by $e_i$ satisfies
$dh_i+h_id=g_i\operatorname{id}$.
These finite modules therefore have finite length and vanish upon localization at any $x_j$.
The second computation leaves only $q=0$, with total-degree-$-p$ cohomology equal to $H_p(K(g;B))$.
Both computations can be implemented by the finite elimination argument of Lemma \ref{sh:double-complex}; the double complex is bounded, so no infinite interchange is involved.
Vanishing in negative total degree gives $H_p(K(g;B))=0$ for $p>0$.

Choose a minimal free resolution $F_\bullet\to B$ over $A=\C\{y_1,\ldots,y_n\}$. By Theorem~\ref{alg:regular-resolution}, its differentials vanish modulo $\mathfrak m_A$. The Koszul resolution of the target coordinates computes the same tensor homology, so
$F_1/\mathfrak m_AF_1\simeq H_1(B\otimes_AK(y;A))=H_1(K(g;B))=0$.
Lemma~\ref{alg:nakayama} gives $F_1=0$, so $F_0\to B$ is an isomorphism and $B$ is finite free.
\end{proof}

\begin{theorem}[Flatness, openness, and equidimensional fibers between manifolds]\label{alg:manifold-flatness}
Let $f:X^n\to Y^m$ be a holomorphic map between nonempty connected complex manifolds.
The following are equivalent: $f$ is flat; $f$ is open; every nonempty fiber is pure of dimension $n-m$.
\end{theorem}
\begin{proof}
By Theorems~\ref{alg:flat-open} and~\ref{alg:flat-dimension}, a flat map is open and its fibers have local dimension $n-m$. Computing at a regular point of each component outside the other components gives pure dimension $n-m$.

Now suppose $f$ is open. Fix $x,y$ and set $d=\dim_xX_y$.
The fiber is defined by the $m$ target-coordinate equations, so Lemma~\ref{alg:parameter-dimension} gives $d\ge n-m$.
Choose a relative finite projection $g=(f,t)$ by Theorem~\ref{alg:finite-projection}.
Its image $Z\subset Y\times\C^d$ is an $n$-dimensional analytic set satisfying
$Z\cap(\{y\}\times D^d)=\{y\}\times D^d$.
If $d>n-m$, then $Z$ is a proper analytic subset of the higher-dimensional manifold $Y\times\C^d$.
In coordinates with $y=0$, choose a nonzero holomorphic function $H(y,t)$ vanishing on $Z$.
Since the whole central vertical polydisc lies in $Z$, $H(0,t)=0$.
Expand $H=\sum_\alpha y^\alpha H_\alpha(t)$ and let $k$ be the least nonzero homogeneous degree.
Choose a constant vector $v$ for which
$P(t)=\sum_{|\alpha|=k}v^\alpha H_\alpha(t)$ is not identically zero, then choose $t_0$ with $P(t_0)\ne0$.
On a small compact neighborhood of $t_0$, the convergence
$s^{-k}H(sv,t)\to P(t)$ is uniform.
For sufficiently small $s\ne0$, no point of $Z$ in that neighborhood projects to $sv$.
Thus $Z\to Y$ is not open at $(0,t_0)$.
But $g$ is onto its image, and the projection of any relatively open $V\subset Z$ is
$f(g^{-1}V)$, which must be open because $f$ is open.
This contradiction gives $d=n-m$ at every point, hence pure-dimensional fibers.

Conversely, assume pure fiber dimension $d=n-m$ and choose a relative finite map $g:(X,x)\to(Y\times\C^d,(y,0))$. Both germs are smooth of dimension $n$, so Lemma~\ref{alg:finite-smooth-flat} makes its ring map finite free, hence flat. The intermediate ring is flat over $\cO_{Y,y}$ by Theorem~\ref{alg:analytic-product-flat}; composition therefore makes $f$ flat at every point.
\end{proof}


\chapter{Weak Derivatives and Elliptic Regularity}\label{ana:chapter}
\begin{learning}
An integral estimate controls the size of a solution but need not give pointwise derivatives. Elliptic regularity supplies those derivatives: a weak solution becomes smooth when the equation has smooth coefficients and a smooth right-hand side. We build this argument from convolution, integration by parts, and difference quotients. Theorem~\ref{ana:elliptic-regularity} is the main result needed for curvature vanishing; the final section gives the uniform estimates used in heat-kernel theory.
\end{learning}

\section{Weak derivatives and Sobolev spaces}
Throughout this chapter, $U\subset\R^d$ is open, $dx$ is Lebesgue measure, and $\|u\|_2=(\int_U|u|^2dx)^{1/2}$.
In real coordinates, $\partial^\alpha=\partial_{x_1}^{\alpha_1}\cdots\partial_{x_d}^{\alpha_d}$; this is the real-coordinate version of the multi-index convention in Definition~\ref{bas:holomorphic-def}.
A distribution on $U$ is a continuous linear functional on $C_c^\infty(U)$, with continuity meaning convergence on test functions having a common compact support and convergence of all their derivatives.
A locally integrable function defines a distribution by integration against a test function.

\begin{definition}
The weak derivative $\partial_j u=v$ of a locally integrable function $u$ is defined by
$\int u\,\partial_j\phi=-\int v\phi$ for every compactly supported smooth test function.
For $U\subset\R^d$, $H^m(U)$ is the space of functions whose weak derivatives $\partial^\alpha u$, $|\alpha|\le m$, all belong to $L^2$, with norm
$\|u\|_{H^m}^2=\sum_{|\alpha|\le m}\|\partial^\alpha u\|_2^2$.
Membership in $H^m_{\mathrm{loc}}(U)$ means that this property holds after multiplication by any $\chi\in C_c^\infty(U)$ and extension by zero to $\R^d$.
The space $H_0^1(U)$ is the closure of $C_c^\infty(U)$ in the $H^1$ norm.
The space $H^{-1}(\R^d)$ is the continuous dual of $H^1(\R^d)$, interpreted through distributional pairing.
A distribution lies in $H^{-1}_{\mathrm{loc}}(U)$ when every such cutoff product belongs to $H^{-1}(\R^d)$.
For complex functions, $\langle u,v\rangle_{L^2}=\int\bar u v$ is conjugate-linear in its first argument; the dual pairing is used with the same convention.
On a manifold or a bundle, Sobolev spaces are defined in coordinate charts and local frames with cutoffs. On a compact manifold a finite atlas and any smooth positive metric and volume form give equivalent norms, by the chain rule and Proposition~\ref{bas:metric}.
For functions with values in a finite-dimensional vector space, these definitions apply componentwise.
\end{definition}

For example, $u(x)=|x|$ on $(-1,1)$ has weak derivative $\operatorname{sgn}(x)$, although it has no classical derivative at zero. Indeed, split the integral $\int |x|\phi'(x)\,dx$ at zero and integrate by parts on the two intervals. The boundary terms at zero vanish because $u(0)=0$, leaving
$\int |x|\phi'=-\int\operatorname{sgn}(x)\phi$.
Thus $u\in H^1(-1,1)$. Weak derivatives let us keep integration by parts even when pointwise differentiability is unavailable.

\begin{lemma}[Smooth approximation and difference quotients]\label{ana:mollifiers}
Choose a nonnegative $\rho\in C_c^\infty(\R^d)$ with $\int\rho=1$, put $\rho_\varepsilon(x)=\varepsilon^{-d}\rho(x/\varepsilon)$, and set
$J_\varepsilon u=\rho_\varepsilon*u$, where $(f*u)(x)=\int f(y)u(x-y)dy$.
Then $J_\varepsilon u\to u$ in $L^2$; for $u\in H^m$, convergence also holds in $H^m$.
For the $j$th standard unit vector $e_j$ and $h\ne0$, write $D_h^ju(x)=(u(x+he_j)-u(x))/h$. If $u\in H^1$, then
$\|D_h^ju\|_2\le\|\partial_j u\|_2$. Conversely, if $u\in L^2$ and these difference quotients are uniformly bounded as $h\to0$, then $\partial_j u\in L^2$.
\end{lemma}
\begin{proof}
Continuous functions of compact support are dense in $L^2$: first truncate the size and support of a function, then approximate by simple functions and approximate measurable sets in measure by open and closed sets, and finally approximate their indicators by continuous cutoffs.
Translation is continuous in $L^2$ for continuous functions of compact support. Since translations preserve the norm, density extends this assertion to every $L^2$ function.
Minkowski's inequality gives
$\|J_\varepsilon u-u\|_2\le\int\rho_\varepsilon(y)\|u(\cdot-y)-u\|_2dy\to0$.
Convolution commutes with weak derivatives, so convergence in $H^m$ follows by applying this argument to each derivative.
For a smooth function, $D_h^ju(x)=\int_0^1\partial_j u(x+the_j)dt$; another application of Minkowski's inequality proves the difference-quotient bound. Smooth approximation extends it to $H^1$ functions.
For the converse, choose a weakly convergent subsequence $D_{h_\ell}^ju\rightharpoonup v$.
A change of variables gives, for every test function $\phi$,
$\int D_h^ju\,\phi=-\int uD_{-h}^j\phi$. Letting $h\to0$ identifies $v$ with $\partial_j u$.
The weakly convergent Hilbert-space subsequence used here can be constructed directly: choose an orthonormal basis of the separable closed subspace generated by the bounded sequence, take a diagonal subsequence on its coordinates, and use Bessel's inequality and boundedness to identify the weak limit.
\end{proof}

\begin{lemma}[Sobolev embedding via the Fourier transform]\label{ana:sobolev-embedding}
If $u\in H^m(\R^d)$ and $m>k+d/2$, then $u$ has a $C^k$ representative satisfying
$\|u\|_{C^k}\le C_{m,k,d}\|u\|_{H^m}$, where $\|u\|_{C^k}=\max_{|\alpha|\le k}\sup_x|\partial^\alpha u(x)|$.
\end{lemma}
\begin{proof}
Use the convention $\widehat u(\xi)=(2\pi)^{-d/2}\int e^{-\ii x\cdot\xi}u(x)dx$, initially for integrable $u$; $x\cdot\xi=\sum_jx_j\xi_j$.
For a compactly supported smooth function, integration by parts shows that its Fourier transform decreases rapidly.
Parseval's identity follows from a short regularization argument: insert $e^{-\varepsilon|\xi|^2}$ into
$\int\overline{\widehat u}\widehat v$, interchange the integrals, and evaluate the product of one-dimensional Gaussian integrals.
The result is the pairing of $u$ with a Gaussian convolution of $v$; the Gaussian approximate identity converges to $v$ as $\varepsilon\downarrow0$.
Taking limits gives $\int\overline{\widehat u}\widehat v=\int\bar uv$.
The same Gaussian argument proves Fourier inversion, and density in $L^2$ then extends Parseval's identity.
Thus the $H^m$ norm is equivalent to $\|(1+|\xi|^2)^{m/2}\widehat u\|_2$.
For $|\alpha|\le k$, Cauchy--Schwarz gives
\[
 \int|\xi^\alpha\widehat u(\xi)|d\xi
 \le\left(\int(1+|\xi|^2)^{k-m}d\xi\right)^{1/2}
 \|(1+|\xi|^2)^{m/2}\widehat u\|_2.
\]
The first integral is finite because $m-k>d/2$.
The inversion integral therefore converges absolutely and uniformly, may be differentiated up to order $k$, and defines continuous derivatives satisfying the stated bound.
For general $u$, smooth approximation from Lemma~\ref{ana:mollifiers} and the same estimate give the conclusion by passage to the limit.
\end{proof}

\section{Gaining the first weak derivative}
For operators $S,T$, write $[S,T]=ST-TS$ for their commutator.
For a differential operator $P=\sum_{|\alpha|\le m}A_\alpha(x)\partial^\alpha$, its principal symbol is
$\sigma_P(x,\xi)=\sum_{|\alpha|=m}A_\alpha(x)(\ii\xi)^\alpha$ for a real covector $\xi$.
Thus symbols multiply under composition, and the principal symbol of $-\sum a^{ij}\partial_i\partial_j$ is $\sum a^{ij}\xi_i\xi_j$.
A formal adjoint $P^*$ is characterized by $\langle Pu,v\rangle=\langle u,P^*v\rangle$ for compactly supported smooth sections, using the chosen metric and volume form; this is an integration-by-parts identity, not yet a choice of domain for an unbounded operator.

In this section we consider functions with values in a finite-dimensional Hermitian vector space and operators
\begin{equation}\label{ana:operator}
 Pu=-\sum_{i,j=1}^d\partial_i(a^{ij}\partial_j u)+\sum_jb^j\partial_j u+cu.
\end{equation}
Here $a^{ij}=a^{ji}$ are real smooth functions, and there is a $\lambda>0$ such that
$\sum a^{ij}(x)\xi_i\xi_j\ge\lambda|\xi|^2$. The coefficients $b^j,c$ may be smooth complex matrices.
In local arguments, all coefficients and their required derivatives are bounded on a fixed compact set.
Repeated coordinate indices in later formulas in this chapter are summed from $1$ to $d$, and $\nabla u=(\partial_1u,\ldots,\partial_du)$ denotes the coordinate gradient.
The model is $P=-\sum_j\partial_j^2$: integration by parts turns $\langle Pu,u\rangle$ into $\|\nabla u\|_2^2$. Uniform ellipticity preserves this control when the principal coefficients vary; the lower-order terms are absorbed by a small part of the gradient term.

\begin{lemma}[Energy estimate]\label{ana:energy}
For smooth $w$ supported in a fixed compact set, there is a constant $C$, independent of $w$, such that
\[
 \|w\|_{H^1}\le C\bigl(\|Pw\|_{H^{-1}}+\|w\|_2\bigr).
\]
\end{lemma}
\begin{proof}
Pair $Pw$ with $w$ and integrate the principal part by parts. Its real part is at least $\lambda\|\nabla w\|_2^2$.
The absolute value of the lower-order terms is at most $C\|\nabla w\|_2\|w\|_2+C\|w\|_2^2$.
Apply $ab\le\varepsilon a^2+(4\varepsilon)^{-1}b^2$ with $\varepsilon$ sufficiently small, absorbing half the gradient term into the left-hand side.
This gives
$\|w\|_{H^1}^2\le C(|\langle Pw,w\rangle|+\|w\|_2^2)$.
The dual norm yields $|\langle Pw,w\rangle|\le\|Pw\|_{H^{-1}}\|w\|_{H^1}$.
Apply the same quadratic inequality to absorb this term, then take square roots.
\end{proof}

\begin{lemma}[The order of the convolution commutator]\label{ana:commutator}
If $a$ is smooth with bounded first derivatives, then $[a,J_\varepsilon]\partial_j$ has $L^2$ operator norm bounded independently of $\varepsilon$.
Consequently, for \eqref{ana:operator},
$\|[P,J_\varepsilon]v\|_{H^{-1}}\le C\|v\|_2$.
\end{lemma}
\begin{proof}
For smooth $v$, integrate
$\int(a(x)-a(y))\rho_\varepsilon(x-y)\partial_jv(y)dy$ by parts in $y_j$. The resulting kernel is
\[
 (\partial_j a)(y)\rho_\varepsilon(x-y)
 +(a(x)-a(y))\partial_j\rho_\varepsilon(x-y).
\]
Its absolute value is bounded by
$\|\nabla a\|_\infty\bigl(\rho_\varepsilon(z)+|z|\,|\partial_j\rho_\varepsilon(z)|\bigr)$.
This dominating convolution kernel has $L^1$ norm independent of $\varepsilon$, so Young's inequality gives the required $L^2$ bound.
Density extends the formula to distributional derivatives of $L^2$ functions.
The commutator of the principal part is
$-\partial_i([a^{ij},J_\varepsilon]\partial_jv)$, a first derivative of a function with a uniform $L^2$ bound, and hence uniformly bounded in $H^{-1}$.
For the first-order terms use the bound just proved; for the zeroth-order term use $2\|c\|_\infty\|v\|_2$. The conclusion follows.
\end{proof}

\begin{proposition}[From $L^2$ to $H^1$]\label{ana:first-gain}
If $u\in L^2_{\mathrm{loc}}(U)$ and $Pu\in H^{-1}_{\mathrm{loc}}(U)$, then $u\in H^1_{\mathrm{loc}}(U)$.
\end{proposition}
\begin{proof}
Choose two concentric balls with closures contained in $U$, take $\chi$ supported in the smaller ball, and set $v=\chi u$, extended by zero to $\R^d$.
Choose a second smooth function $0\le\chi_2\le1$, equal to one near $\supp\chi$ and supported in the larger ball.
Extend the principal coefficients by $\chi_2 a+(1-\chi_2)I$, preserving uniform positive definiteness, and extend the remaining coefficients using smooth cutoffs.
A first-order term maps $L^2$ to $H^{-1}$ because $a\partial_j u=\partial_j(au)-(\partial_j a)u$.
Since $[P,\chi]$ is first order and $u\in L^2_{\mathrm{loc}}$, we obtain
$Pv=\chi Pu+[P,\chi]u\in H^{-1}(\R^d)$.
Set $v_\varepsilon=J_\varepsilon v$. For small $\varepsilon$, its support remains in a fixed compact set, and
\[
 Pv_\varepsilon=J_\varepsilon(Pv)+[P,J_\varepsilon]v.
\]
Convolution is bounded on $H^{-1}$: apply Young's inequality in $H^1$ to test functions and pass to the dual.
Lemma \ref{ana:commutator} therefore bounds the right-hand side uniformly in $H^{-1}$.
Lemma \ref{ana:energy} gives a uniform $H^1$ bound for $v_\varepsilon$.
Take a weakly convergent subsequence. Its $L^2$ limit is already identified as $v$ by Lemma \ref{ana:mollifiers}.
Passing to the limit also in the test-function pairings for weak derivatives shows that $v\in H^1$.
Taking $\chi=1$ near any prescribed smaller compact set proves the local assertion.
\end{proof}

\section{Difference quotients give the second weak derivative}
\begin{proposition}[Interior $H^2$ regularity]\label{ana:second-gain}
If $u\in H^1_{\mathrm{loc}}(U)$ and $Pu=f\in L^2_{\mathrm{loc}}(U)$, then $u\in H^2_{\mathrm{loc}}(U)$.
For $V\Subset W\Subset U$, there is a constant $C$ such that
$\|u\|_{H^2(V)}\le C(\|f\|_{L^2(W)}+\|u\|_{H^1(W)})$.
\end{proposition}
\begin{proof}
Move the lower-order terms to the right to obtain
$-\partial_i(a^{ij}\partial_ju)=g$, where
$\|g\|_{L^2(W)}\le C(\|f\|_{L^2(W)}+\|u\|_{H^1(W)})$.
Choose $\eta\in C_c^\infty(W)$ with $\eta=1$ on $V$ by Theorem~\ref{bas:partition}, and take $|h|$ small enough that all relevant translates remain in $W$.
Fix a coordinate direction and write $w=D_hu$. In the weak equation use the test function
$-D_{-h}(\eta^2w)$; it belongs to $H_0^1(W)$ and is admissible by smooth approximation.
Discrete integration by parts, $\int aD_{-h}b=-\int(D_ha)b$, and the identity
$D_h(a\partial_ju)=a(x+h)\partial_jw+(D_ha)\partial_ju$
give
\[
 \int a^{ij}(x+h)\partial_jw\,\overline{\partial_i(\eta^2w)}
 =-\int g\,\overline{D_{-h}(\eta^2w)}
 -\int(D_ha^{ij})\partial_ju\,\overline{\partial_i(\eta^2w)}.
\]
Taking real parts, the principal term on the left is at least $\lambda\|\eta\nabla w\|_2^2$; the terms containing $\nabla\eta$ are bounded by
$C\|w\|_2\|\eta\nabla w\|_2+C\|w\|_2^2$.
Lemma~\ref{ana:mollifiers} gives $\|D_{-h}(\eta^2w)\|_2\le C(\|\eta\nabla w\|_2+\|w\|_2)$,
and $\|D_ha^{ij}\|_\infty\le\|\partial a^{ij}\|_\infty$.
Apply Cauchy--Schwarz and the quadratic inequality with a small parameter to each term, absorbing every product linear in $\|\eta\nabla w\|_2$. The result is
\[
 \|\eta\nabla D_hu\|_2^2
 \le C\bigl(\|g\|_{L^2(W)}^2+\|u\|_{H^1(W)}^2+\|D_hu\|_{L^2(W')}^2\bigr),
\]
where $W'$ is a fixed, slightly smaller region than $W$ containing all the required translates.
The last term is bounded by $\|\nabla u\|_{L^2(W)}^2$.
Thus the difference quotients of each $\partial_ju$ in every coordinate direction are uniformly bounded on $V$.
Take a weakly convergent subsequence in $L^2(V)$ and integrate by parts discretely against an arbitrary compactly supported test function on $V$; sufficiently small translates still lie in $W$.
As $h\to0$, the argument of Lemma \ref{ana:mollifiers} identifies the limit with the corresponding weak derivative.
All second weak derivatives therefore belong to $L^2(V)$ and satisfy the preceding estimate.
\end{proof}

\begin{theorem}[Local elliptic regularity]\label{ana:elliptic-regularity}
For a uniformly elliptic operator \eqref{ana:operator} with smooth coefficients, if
$u\in L^2_{\mathrm{loc}}$ and $Pu$ is smooth, then $u$ is smooth.
The same conclusion holds for second-order operators on manifolds whose principal symbol is a positive definite real quadratic form times the identity on the fiber.
\end{theorem}
\begin{proof}
Proposition \ref{ana:first-gain} first gives $u\in H^1_{\mathrm{loc}}$, and Proposition \ref{ana:second-gain} then gives $u\in H^2_{\mathrm{loc}}$.
Suppose that $u\in H^{m+1}_{\mathrm{loc}}$ is already known, and take $|\alpha|=m$.
After writing the operator in nondivergence form, distributional differentiation gives
\[
 P(\partial^\alpha u)=\partial^\alpha(Pu)+[P,\partial^\alpha]u.
\]
Every commutator term has at least one derivative on a coefficient, so at most $m+1$ derivatives fall on $u$.
The right-hand side therefore belongs to $L^2_{\mathrm{loc}}$, while $\partial^\alpha u\in H^1_{\mathrm{loc}}$.
Proposition \ref{ana:second-gain} places it in $H^2_{\mathrm{loc}}$, giving $u\in H^{m+2}_{\mathrm{loc}}$.
Induction yields Sobolev regularity of every order. After multiplying by a local cutoff, apply Lemma \ref{ana:sobolev-embedding}: for each $k$, choosing $m>k+d/2$ gives $u\in C^k$.
These representatives agree because they are continuous and equal almost everywhere, so $u$ is smooth.
In coordinates and local frames, the manifold case reduces to the operators above. Rewriting the second-order terms in divergence form changes only first-order coefficients, and the local conclusions patch together.
\end{proof}

\begin{example}[Why the Dolbeault Laplacian satisfies the hypotheses]
Choose Hermitian metrics on the manifold and on the coefficient bundle, and let $\dbar^*$ be the formal adjoint just defined.
For a real covector $\xi$, let $\xi^{0,1}$ be its component in the type decomposition of Proposition~\ref{bas:type-decomposition}, and let $\varepsilon(\xi^{0,1})$ mean exterior multiplication by that covector.
With our symbol convention, $\sigma_{\dbar}(\xi)=\ii\varepsilon(\xi^{0,1})$ and $\sigma_{\dbar^*}(\xi)=-\ii\varepsilon(\xi^{0,1})^*$.
On an orthonormal wedge basis, insertion and exterior multiplication satisfy
$\varepsilon\varepsilon^*+\varepsilon^*\varepsilon=|\xi^{0,1}|^2I$: each basis form either contains the covector direction or does not, and precisely one summand contributes.
Consequently the Dolbeault Laplacian $\Delta''=\dbar\dbar^*+\dbar^*\dbar$ has
$\sigma_{\Delta''}(\xi)=|\xi^{0,1}|^2I$.
For a nonzero real cotangent vector this number is strictly positive, so $\Delta''$ has positive definite scalar principal symbol.
This explains why the minimal-norm weak solutions in Chapter \ref{pos:chapter} become smooth, without first invoking the full Hodge theory.
\end{example}

\section{A compactness tool: Rellich's lemma}
\begin{theorem}[Compact embedding on a compact manifold]\label{ana:rellich}
On a compact smooth manifold without boundary, for a smooth vector bundle of finite rank, the inclusion $H^1\hookrightarrow L^2$ is compact.
\end{theorem}
\begin{proof}
Using finitely many coordinate charts and Theorem~\ref{bas:partition}, split an $H^1$-bounded sequence into finitely many components supported compactly inside coordinate cubes.
Extend each component by zero to a larger cube, then periodically, to obtain an $H^1$-bounded sequence on a fixed real torus. The product Fej\'er-kernel argument in Proposition~\ref{sh:torus}, in $d$ real coordinates, makes the trigonometric functions a complete orthonormal system; their finite sums therefore give Parseval's identity by $L^2$ approximation.
In the Fourier basis $e^{2\pi\ii k\cdot x}$, Parseval's identity and the derivative formula give
$\sum_k(1+|k|^2)|c_k|^2\le C$.
For every member of the sequence, the squared $L^2$ norm of the Fourier tail $|k|>R$ is therefore at most $C/R^2$.
Here $k\in\Z^d$ is the frequency index and $c_k$ is the corresponding Fourier coefficient. For fixed $R$ only finitely many bounded coefficients remain. Choose a subsequence on which they all converge, and diagonalize over integer $R$ to obtain a subsequence converging at every low frequency.
Given an error tolerance, first choose $R$ so that both tails are small, then make the low-frequency difference small. The subsequence is thus Cauchy in $L^2$.
Diagonalize once more over the finitely many charts and bundle components, then sum to obtain an $L^2$-convergent subsequence of the original sequence.
\end{proof}

\begin{remark}
The result proved here is an interior statement on a compact manifold. Problems with boundary also require specified boundary conditions; this chapter does not by itself prove regularity for arbitrary boundary problems.
For more general elliptic systems, boundary regularity, and versions with distributional data, see \cite{Evans2010}. The vanishing and embedding arguments below use the interior estimates with smooth coefficients established here.
\end{remark}

\begingroup
\setlength{\parskip}{2pt plus 1pt}
\setlength{\abovedisplayskip}{6pt plus 2pt minus 2pt}
\setlength{\belowdisplayskip}{6pt plus 2pt minus 2pt}
\setlength{\abovedisplayshortskip}{0pt plus 2pt}
\setlength{\belowdisplayshortskip}{4.5pt plus 1pt minus 1pt}
\section{Quantitative interior estimates and first-order operators}
The heat-kernel arguments below require regularity for whole families of operators. Smoothness of each weak solution is therefore insufficient: we must identify the data on which the estimate constants depend.

\begin{lemma}[Interpolation and interior absorption]\label{ana:interpolation}
For a compactly supported smooth function, an integer $m\ge1$, and $0<\varepsilon\le1$,
\[
 \|u\|_{H^{m-1}}\le\varepsilon\|u\|_{H^m}
       +C_m\varepsilon^{-(m-1)}\|u\|_2.
\]
Suppose $M>0$, $R>r_0$, and a nonnegative bounded function $F$ on $[r_0,R]$ satisfies
$F(r)\le\theta F(s)+A(s-r)^{-M}$ for $r<s$, with a sufficiently small $\theta>0$. Then
$F(r_0)\le C_{M,R-r_0}A$.
\end{lemma}
\begin{proof}
Split Fourier space into $|\xi|\le 2/\varepsilon$ and its complement. On the first region, $(1+|\xi|^2)^{(m-1)/2}\le C_m\varepsilon^{-(m-1)}$; on the second it is at most $C\varepsilon(1+|\xi|^2)^{m/2}$. Estimate the two integrals separately, take square roots, and adjust $\varepsilon$ to obtain the first inequality.
For the second assertion, put $r_j=r_0+(R-r_0)(1-2^{-j})$. Iterating the inequality $N$ times gives
\[
 F(r_0)\le\theta^N F(r_N)+A(R-r_0)^{-M}2^M
                    \sum_{j=0}^{N-1}(\theta2^M)^j.
\]
Fix $\theta<2^{-M-1}$. Since $F$ is bounded, the first term tends to zero, while the geometric sum is at most $2$. This is the explicit mechanism behind absorbing an interior higher-order norm.
\end{proof}

\begin{theorem}[Uniform interior elliptic estimates]\label{ana:interior-estimate}
Let $P$ be as in \eqref{ana:operator}, let $U\Subset V$, and let $s\ge0$ be an integer. Then
\begin{equation}\label{ana:quantitative-estimate}
 \|u\|_{H^{s+2}(U)}\le C\bigl(\|Pu\|_{H^s(V)}+\|u\|_{L^2(V)}\bigr).
\end{equation}
The estimate holds first for smooth $u$, and also for $u\in L^2(V)$ whenever the right-hand side is finite.
The constant depends only on $U\Subset V$, the ellipticity lower bound, the matrix size, and common bounds for finitely many derivatives of the coefficients on $V$. It can be chosen uniformly for families in which these data are uniformly controlled.
\end{theorem}
\begin{proof}
First work on concentric balls. Keeping track of the cutoff derivatives in Proposition \ref{ana:second-gain} gives
\[
 \|u\|_{H^2(B_r)}\le C(\rho-r)^{-M}
       \bigl(\|Pu\|_{L^2(B_\rho)}+\|u\|_{H^1(B_\rho)}\bigr),
\]
where $M$ is a fixed integer whose precise value is immaterial here. Choose a cutoff supported on a slightly larger ball and equal to one on $B_\rho$, and apply Lemma \ref{ana:interpolation} to its product with $u$. Differentiating the product produces only lower-order terms. Taking the interpolation parameter sufficiently small relative to $(\rho-r)^M$ gives
\[
 \|u\|_{H^2(B_r)}\le\theta\|u\|_{H^2(B_t)}
  +C_\theta(t-r)^{-M'}\bigl(\|Pu\|_{L^2(B_t)}+\|u\|_{L^2(B_t)}\bigr)
\]
with $\rho=(r+t)/2$ between $r$ and $t$. Lemma~\ref{ana:interpolation} eliminates the first term, proving the case of Sobolev order $s=0$.

Proceed by induction to higher orders. For $|\alpha|=s$, the proof of Theorem~\ref{ana:elliptic-regularity} gives
\[
 P\partial^\alpha u=\partial^\alpha Pu+[P,\partial^\alpha]u,
 \qquad
 \|[P,\partial^\alpha]u\|_2\le C_s\|u\|_{H^{s+1}}.
\]
Apply the $H^2$ estimate to each $\partial^\alpha u$ and sum. The interior $H^{s+2}$ norm is bounded by $\|Pu\|_{H^s}$ and the $H^{s+1}$ norm on a slightly larger region. Apply interpolation to a cutoff product and the same interior absorption lemma to eliminate $H^{s+1}$, obtaining \eqref{ana:quantitative-estimate}.
For general $U\Subset V$, use a finite covering by such balls and sum.

For weak solutions, Propositions \ref{ana:first-gain} and \ref{ana:second-gain} first supply the first two weak derivatives. Differentiate weakly and repeat the difference-quotient argument. At every stage the bound uses only the stated right-hand side, with no assumption of additional derivatives. Equivalently, one may first convolve in each difference-quotient test identity and then pass to a weak limit. Weak lower semicontinuity preserves the estimate for the resulting weak derivatives.
The calculation uses only finitely many coefficient derivatives, fixed cutoff bounds, and the ellipticity constant. This also proves the final uniformity assertion.
\end{proof}

\begin{proposition}[Uniform interior estimates for Dirac-type operators]\label{ana:dirac-estimate}
Let $D=\sum_j A^j(x)\partial_j+B(x)$ be a formally self-adjoint first-order operator with smooth coefficients.
Assume its symbol satisfies the Clifford relation
$\sigma_D(x,\xi)\sigma_D(x,\eta)+\sigma_D(x,\eta)\sigma_D(x,\xi)=2g_x^{-1}(\xi,\eta)I$,
where $g_x^{-1}$ is a positive definite real inner product on cotangent vectors.
In particular, $\sigma_{D^2}(x,\xi)=|\xi|_{g^{-1}}^2I$, a positive definite real quadratic form times the identity matrix. For every integer $s\ge0$,
\[
 \|u\|_{H^{s+1}(U)}\le C_s\bigl(\|Du\|_{H^s(V)}+\|u\|_{L^2(V)}\bigr),
 \qquad U\Subset V.
\]
The constants have the same uniformity as in Theorem \ref{ana:interior-estimate}.
\end{proposition}
\begin{proof}
For compactly supported smooth $v$, formal self-adjointness gives $\|Dv\|_2^2=\langle D^2v,v\rangle$. Integrate the second-order principal terms by parts. Positive definiteness yields
\[
 \lambda\|\nabla v\|_2^2
 \le\|Dv\|_2^2+C\|\nabla v\|_2\|v\|_2+C\|v\|_2^2.
\]
Absorb the middle term with the quadratic inequality with a small parameter to get
$\|v\|_{H^1}\le C(\|Dv\|_2+\|v\|_2)$.
Set $v=\chi u$ and use $D(\chi u)=\chi Du+\sum_jA^j(\partial_j\chi)u$ to obtain the interior estimate for $s=0$. The commutator $[D,\chi]=\sum_jA^j(\partial_j\chi)$ has order zero and is bounded on each fixed compact set.

For $|\alpha|=s$,
$D\partial^\alpha u=\partial^\alpha Du+[D,\partial^\alpha]u$.
The commutator has order at most $s$, so the preceding $H^1$ estimate gives
\[
 \|u\|_{H^{s+1}(U)}\le C\bigl(\|Du\|_{H^s(V_1)}+\|u\|_{H^s(V_1)}\bigr)
\]
for $U\Subset V_1\Subset V$. Apply the induction hypothesis to $V_1\Subset V$ to bound the last $H^s$ norm by $\|Du\|_{H^{s-1}(V)}+\|u\|_2$, which is already controlled by the required right-hand side. This completes the induction.
For weak solutions, the first-order Friedrichs convolution commutator has kernels of the form $(A(x)-A(y))\nabla\rho_\varepsilon(x-y)$ and $\nabla A\,\rho_\varepsilon$, dominated by kernels with uniformly bounded $L^1$ norms. Convolution, cutoffs, and weak compactness give the same conclusion.
\end{proof}
\endgroup


\chapter{Plurisubharmonic Functions, Positive Currents, and Curvature Vanishing}
\label{pos:chapter}

The complex Hessian gives an analogue of convexity adapted to holomorphic functions: its nonnegativity means subharmonicity on each complex line. This condition also makes sense for singular functions through their associated currents. For a line bundle, the same Hessian describes curvature, and a positive lower bound yields an estimate for solving $\dbar u=f$. We develop these connections first in explicit local models, then on compact K\"ahler manifolds, where solvability gives cohomology vanishing.

Throughout this chapter we use the conventions
\begin{equation}\label{pos:dc-convention}
 d^c=\frac{\ii}{2\pi}(\dbar-\partial),\qquad
 dd^c=\frac{\ii}{\pi}\partial\dbar.
\end{equation}
Write $\Theta_h=F_h$ for the Chern curvature of Theorem~\ref{bun:chern-connection}; the symbol does not include the factor $\ii$.
In a holomorphic frame $e$, a line bundle metric is written $|e|_h^2=e^{-\varphi}$, so that
$\ii\Theta_h=\ii\partial\dbar\varphi$ by \eqref{bun:line-curvature}.
The integrability condition defining multiplier ideals will be $|f|^2e^{-\varphi}\in L^1_{\mathrm{loc}}$.
To compare with the convention $e^{-2\psi}$ used in some references, set $\varphi=2\psi$;
the two weight conventions must not be mixed within a proof.
For general background, see \cite{Demailly2012,Hormander1965}.

\section{From the Mean Inequality in One Variable to the Complex Hessian}

\begin{definition}\label{pos:psh-definition}
An extended-real-valued function on an open set $\Omega\subset\C^n$,
$u:\Omega\to[-\infty,+\infty)$, is called plurisubharmonic, abbreviated psh,
if it is upper semicontinuous, locally integrable, and satisfies
the following inequality on every closed disk $\{a+\zeta v:|\zeta|\leq r\}\subset\Omega$ in a complex affine line:
\[
 u(a)\leq\frac1{2\pi}\int_0^{2\pi}u(a+re^{\ii t}v)\,dt.
\]
Its restriction to an individual line is allowed to be identically $-\infty$.
For the discussion of monotone limits, we also allow the identically $-\infty$ function as a separate degenerate case.
Every psh function whose distributional derivatives are taken in this chapter belongs to $L^1_{\mathrm{loc}}$. We write $\operatorname{PSH}(\Omega)$ for the nondegenerate psh functions on $\Omega$.
\end{definition}

Local integrability is included in the definition to make the hypotheses of each distributional calculation explicit.
The usual definition requires only subharmonicity on complex lines; after the identically $-\infty$ case is excluded on a connected domain,
local integrability follows from the complete argument in Lemma \ref{cur:psh-integrable}.

\begin{lemma}[The smooth criterion in one variable]\label{pos:subharmonic-smooth}
Let $v$ be a real $C^2$ function on a plane domain, and put $\Delta=\partial_x^2+\partial_y^2$. This local scalar Laplacian has the opposite sign from the nonnegative spectral Laplacians used later.
It satisfies the circular mean inequality on every closed disk if and only if $\Delta v\geq0$.
\end{lemma}
\begin{proof}
For a fixed center $a$, put
$M(r)=(2\pi)^{-1}\int_0^{2\pi}v(a+re^{\ii t})\,dt$.
Differentiation for $r>0$, followed by the divergence theorem, gives
\[
 M'(r)=\frac1{2\pi r}\int_{B(a,r)}\Delta v\,dx\,dy.
\]
If $\Delta v\geq0$, then $M$ is nondecreasing, and
$\lim_{r\downarrow0}M(r)=v(a)$, which proves the mean inequality.
Conversely, the second-order Taylor expansion gives
\[
 M(r)=v(a)+\frac{r^2}{4}\Delta v(a)+o(r^2).
\]
The circular mean of the linear term is zero; the means of $x^2$ and $y^2$ are both $r^2/2$,
and that of $xy$ is zero. Since $M(r)\geq v(a)$, division by $r^2$ and passage to the limit $r\to0$
give $\Delta v(a)\geq0$.
\end{proof}

\begin{theorem}[The Levi criterion and smooth pullbacks]\label{pos:levi}
Let $u\in C^2(\Omega,\R)$. The following conditions are equivalent:
\begin{enumerate}
\item $u$ is psh;
\item the Hermitian matrix $\bigl(u_{j\bar k}\bigr)$ is positive semidefinite, where
$u_{j\bar k}=\partial^2u/\partial z_j\partial\bar z_k$;
\item $u\circ F$ is subharmonic for every holomorphic map $F:\Delta\to\Omega$, where $\Delta=\{\zeta\in\C:|\zeta|<1\}$.
\end{enumerate}
In particular, holomorphic pullbacks of smooth psh functions are psh.
\end{theorem}
\begin{proof}
Along the complex line $F(\zeta)=a+\zeta v$, we have
\[
 \Delta_\zeta(u\circ F)(0)
 =4\sum_{j,k}u_{j\bar k}(a)v_j\bar v_k.
\]
Lemma~\ref{pos:subharmonic-smooth} proves the equivalence of the first two conditions.
For a general holomorphic $F$, the chain rule gives
\[
 \frac{\partial^2(u\circ F)}{\partial\zeta\partial\bar\zeta}
 =\sum_{j,k}u_{j\bar k}(F(\zeta))F_j'(\zeta)
          \overline{F_k'(\zeta)}.
\]
No term involving $F''$ occurs, because
$\partial_{\bar\zeta}F_j=0$ and
$\partial_\zeta\overline{F_k}=0$.
The second condition makes the right-hand side nonnegative, proving the third; applied to affine maps, the third
implies the first. For an arbitrary holomorphic map $G$, apply this calculation
to the composition of $G$ with each affine line to obtain the final assertion.
\end{proof}

\begin{proposition}[Basic stability properties]\label{pos:psh-stability}
Plurisubharmonic functions have the following stability properties.
\begin{enumerate}
\item A finite nonnegative linear combination of psh functions is psh.
\item $\max(u,v)$ is psh.
\item If $u_\nu\downarrow u$ and $u\in L^1_{\mathrm{loc}}$, then $u$ is psh.
\item If $u$ is smooth and psh, and $\chi\in C^2(\R)$ satisfies
$\chi'\geq0$ and $\chi''\geq0$, then $\chi\circ u$ is psh.
\end{enumerate}
\end{proposition}
\begin{proof}
For the first assertion, upper semicontinuity, local integrability, and the mean inequality are preserved under the indicated sums (terms with zero coefficient are omitted).
For the second, a finite maximum of upper semicontinuous functions is upper semicontinuous.
The two functions have a common upper bound on every compact set, and the negative part of their maximum is bounded by the negative part of $u$;
hence the maximum is locally integrable. If $u(a)\geq v(a)$ at the center, then
\[
 \max(u,v)(a)=u(a)
 \leq\frac1{2\pi}\int u(a+re^{\ii t}w)\,dt
 \leq\frac1{2\pi}\int\max(u,v)(a+re^{\ii t}w)\,dt.
\]
The other case follows by interchanging $u$ and $v$.

For the third assertion, $u=\inf_\nu u_\nu$ is upper semicontinuous.
Near any fixed circle, $u_1$ has a finite upper bound $M$.
The monotone convergence theorem applied to the nonnegative increasing sequence $M-u_\nu$ permits passage to the limit
on both sides of $u_\nu(a)\leq(2\pi)^{-1}\int u_\nu$.
The limiting right-hand side may equal $-\infty$; the assumption $u\in L^1_{\mathrm{loc}}$
ensures that the limit satisfies the nondegenerate definition used here.

For the fourth assertion, direct differentiation gives
\[
 (\chi\circ u)_{j\bar k}
 =\chi'(u)u_{j\bar k}+\chi''(u)u_j u_{\bar k}.
\]
The first term is positive semidefinite, and the second is a nonnegative multiple of a positive semidefinite matrix of rank at most one.
Theorem~\ref{pos:levi} proves the claim.
\end{proof}

\begin{example}[Logarithmic weights]\label{pos:log-norm}
Let $g_1,\ldots,g_r$ be holomorphic functions and $\epsilon>0$. Then
\[
 u_\epsilon=\log\bigl(\epsilon+|g_1|^2+\cdots+|g_r|^2\bigr)
\]
is smooth and psh. If $\log\sum|g_j|^2$ is locally integrable, it is also psh.
\end{example}
\begin{proof}
Put $G=\epsilon+\sum|g_j|^2$. For a tangent vector $v$, write
$a_j=dg_j(v)$. Then
\[
 \sum_{k,l}(u_\epsilon)_{k\bar l}v_k\bar v_l
 =\frac{G\sum_j|a_j|^2-|\sum_j a_j\bar g_j|^2}{G^2}\geq0.
\]
The final inequality is Cauchy--Schwarz.
Let $\epsilon\downarrow0$ and apply the third assertion of Proposition \ref{pos:psh-stability}.
For example, if $g=z_1^{m_1}\cdots z_s^{m_s}$, local integrability follows coordinate by coordinate from
$\int_0^\delta r|\log r|\,dr<\infty$ and Fubini's theorem.
\end{proof}



\section{Positivity Cones and Currents: Degrees, Signs, and Measures}

Let $V$ be a complex vector space of dimension $n$. Its complex orientation declares
$\ii dz_1\wedge d\bar z_1\wedge\cdots\wedge\ii dz_n\wedge d\bar z_n$
to be a positive top-degree form.

\begin{definition}\label{pos:positive-forms}
A real $(p,p)$ form $\alpha$ is called strongly positive if it is a finite sum
\[
 \alpha=\sum_\nu c_\nu
 (\ii\theta_{\nu,1}\wedge\bar\theta_{\nu,1})\wedge\cdots\wedge
 (\ii\theta_{\nu,p}\wedge\bar\theta_{\nu,p}),\qquad c_\nu\geq0,
\]
where every $\theta_{\nu,j}$ is a $(1,0)$ covector. It is called weakly positive if, for all $(1,0)$ covectors
$\eta_1,\ldots,\eta_{n-p}$, the top-degree form
\[
 \alpha\wedge\ii\eta_1\wedge\bar\eta_1\wedge\cdots\wedge
 \ii\eta_{n-p}\wedge\bar\eta_{n-p}
\]
is nonnegative. In this chapter, the unqualified term ``positive current'' means positivity dual to strongly positive test forms.
\end{definition}

\begin{proposition}\label{pos:positive-algebra}
Every strongly positive form is weakly positive. The wedge product of strongly positive forms is strongly positive.
The wedge product of a weakly positive form with a strongly positive form is weakly positive, provided its bidegree does not exceed $(n,n)$.
In bidegree $(1,1)$, strong and weak positivity coincide and are equivalent to positive semidefiniteness of the Hermitian coefficient matrix.
\end{proposition}
\begin{proof}
For each generator of the strongly positive cone,
\[
 \bigwedge_{j=1}^{p}(\ii\theta_j\wedge\bar\theta_j)
 =\ii^{p^2}(\theta_1\wedge\cdots\wedge\theta_p)
                 \wedge\overline{(\theta_1\wedge\cdots\wedge\theta_p)}.
\]
After wedging with $n-p$ test factors, the coefficient of the resulting top-degree form is
the squared absolute value of the corresponding determinant. It is therefore nonnegative; summing proves the first assertion.
The wedge product of two strongly positive generators is again a generator of the same type, proving the second assertion.
If $\alpha$ is weakly positive and $\beta$ strongly positive, expand $\beta$ into generators
and append test factors to reach bidegree $(n,n)$; weak positivity gives nonnegativity term by term.

Finally, write a real $(1,1)$ form as
$\alpha=\ii\sum a_{j\bar k}dz_j\wedge d\bar z_k$.
The spectral theorem puts it, in a unitary basis, in the form
$\ii\sum\lambda_j dz_j\wedge d\bar z_j$.
If all $\lambda_j\geq0$, this is a strongly positive decomposition. If some $\lambda_j<0$,
wedging with the strongly positive factors in the other $n-1$ coordinate directions produces a negative top-degree form,
contradicting weak positivity.
\end{proof}

\begin{remark}
In intermediate degrees, the possibility of arranging the coefficients in a matrix does not identify the different notions of positivity.
For $p>1$, strong positivity requires the $p$-forms in its generating expressions to be decomposable.
For instance, $dz_1\wedge dz_2+dz_3\wedge dz_4$ has nonzero wedge square,
so it cannot be the wedge product of two one-forms.
The $L^2$ curvature calculations below mainly use $(1,1)$ positivity and the separately defined notion of Nakano positivity,
so these cones need not be conflated.
\end{remark}

\begin{definition}\label{pos:currents}
On an oriented manifold of real dimension $m=2n$, a current $T$ of degree $k$ is a continuous linear functional
on the space of compactly supported smooth $(m-k)$-forms.
A smooth $k$-form $\alpha$ defines the current
$\langle T_\alpha,\eta\rangle=\int_X\alpha\wedge\eta$.
Its derivative is defined by
\[
 \langle dT,\eta\rangle=(-1)^{k+1}\langle T,d\eta\rangle.
\]
Continuity means continuity under convergence of every derivative on a fixed compact support. A $(p,q)$ current on a complex manifold acts on $(n-p,n-q)$ test forms.
A $(p,p)$ current is positive if it takes a nonnegative value on every compactly supported
strongly positive $(n-p,n-p)$ test form.
Weak convergence $T_\nu\to T$ means convergence of the value on every fixed test form. A current has \emph{order zero} if, on each fixed compact support, its value is bounded by a constant times the supremum norm of the test-form coefficients.
\end{definition}

\begin{proposition}\label{pos:current-calculus}
These definitions have the following properties.
\begin{enumerate}
\item $dT_\alpha=T_{d\alpha}$ and $d^2=0$.
\item If $Y\subset X$ is a closed embedded complex submanifold without boundary, of complex codimension $p$, then
$[Y](\eta)=\int_Y\eta$ is a closed positive $(p,p)$ current.
\item Weak limits of positive currents are positive, and weak limits of closed currents are closed.
\end{enumerate}
\end{proposition}
\begin{proof}
Compact support allows Stokes' theorem (Proposition~\ref{bas:stokes}) to give
$0=\int d(\alpha\wedge\eta)
=\int d\alpha\wedge\eta+(-1)^k\int\alpha\wedge d\eta$,
which is precisely the derivative formula in the first assertion. Differentiating again gives $d^2\eta=0$.

In the second assertion, the closed embedding ensures that a compactly supported test form has compact support when restricted to $Y$,
so the integral is finite and depends continuously on the test form.
A strongly positive test form restricts to a nonnegative top-degree form on $Y$, hence has nonnegative integral.
The definition of $d[Y]$ gives, up to a constant sign,
$\int_Yd\eta=\int_{\partial Y}\eta=0$.

The third assertion follows by taking limits against each test form.
For example, if $T_\nu\to T$ and $dT_\nu=0$, then
$\langle dT,\eta\rangle=(-1)^{k+1}\lim_\nu\langle T_\nu,d\eta\rangle=0$.
\end{proof}

\begin{proposition}[Local smoothing and positive Hessian currents]\label{pos:regularization}
Let $u\in\operatorname{PSH}(\Omega)\cap L^1_{\mathrm{loc}}(\Omega)$.
On every $\Omega'\Subset\Omega$, convolution with nonnegative smooth kernels produces
smooth psh functions $u_\epsilon$ such that
$u_\epsilon\to u$ in $L^1_{\mathrm{loc}}(\Omega')$.
Consequently, $\ii\partial\dbar u$ is a positive $(1,1)$ current.
\end{proposition}
\begin{proof}
Choose a nonnegative $\rho\in C_c^\infty(\C^n)$ supported in the unit ball with $\int\rho=1$,
and put $\rho_\epsilon(z)=\epsilon^{-2n}\rho(z/\epsilon)$.
When $\epsilon$ is smaller than the distance to the boundary, $u_\epsilon=u*\rho_\epsilon$ is defined and smooth.
Take a closed disk in a complex line contained in $\Omega'$.
For each translation vector $w$, multiply the mean inequality for $u$ on the translated disk by
$\rho_\epsilon(w)$ and integrate. This gives the mean inequality for $u_\epsilon$.
The interchange of integrals is justified by applying Tonelli's theorem to a local upper bound minus $u$, and then using local integrability.
Thus $u_\epsilon$ is psh.

The $L^1$ approximation property of convolution follows from continuity of translations: on a slightly larger compact set,
\[
 \|u*\rho_\epsilon-u\|_{L^1}
 \leq\int\rho_\epsilon(w)\|u(\,\cdot-w)-u\|_{L^1}\,dw\longrightarrow0.
\]
Translation continuity is first checked for continuous compactly supported functions and then follows from their density in $L^1$.
Distributional differentiation is continuous, so
$\ii\partial\dbar u_\epsilon\to\ii\partial\dbar u$.
Theorem~\ref{pos:levi} and Proposition~\ref{pos:positive-algebra} make the smooth currents on the left positive. Taking limits against strongly positive test forms, as in Proposition~\ref{pos:current-calculus}, proves positivity of the limiting current.
\end{proof}

\begin{lemma}[Nonnegative distributions are measures]\label{pos:positive-distribution}
If a distribution $R$ satisfies $R(f)\geq0$ for every nonnegative $f\in C_c^\infty$,
then $R$ has order zero and is represented by a locally finite positive Radon measure.
\end{lemma}
\begin{proof}
Fix a compact set $K$ and choose $\chi\in C_c^\infty$ with $\chi\geq0$ and
$\chi=1$ near $K$. For a real smooth function $f$ supported in $K$,
$-\|f\|_\infty\chi\leq f\leq\|f\|_\infty\chi$, whence
\[
 |R(f)|\leq R(\chi)\|f\|_\infty.
\]
Applying this separately to the real and imaginary parts gives the corresponding order-zero bound for complex functions.
Uniform approximation therefore extends $R$ uniquely to a positive functional on compactly supported continuous functions.
The Riesz representation theorem represents this functional by a Radon measure.
Local finiteness also follows directly from $R(\chi)<\infty$.
Here we use the Riesz representation theorem from real analysis; no finiteness of the measure on the entire noncompact manifold is assumed.
\end{proof}

\begin{theorem}[Positive currents have no higher-order distributional coefficients]\label{pos:order-zero}
Every local coefficient of a positive $(p,p)$ current is a complex Radon measure. Thus the current has order zero.
It is also real.
\end{theorem}
\begin{proof}
In a coordinate chart, take any constant-coefficient strongly positive $(n-p,n-p)$ form $\gamma$.
The functional $f\mapsto T(f\gamma)$ is a nonnegative distribution and hence a measure by Lemma~\ref{pos:positive-distribution}.
We show that the complex span of such $\gamma$ contains all constant-coefficient
$(n-p,n-p)$ forms. For any one-forms $a,b$, the identity
\begin{align*}
4a\wedge\bar b={}&(a+b)\wedge\overline{a+b}
 -(a-b)\wedge\overline{a-b}\\
&+\ii(a+\ii b)\wedge\overline{a+\ii b}
 -\ii(a-\ii b)\wedge\overline{a-\ii b}
\end{align*}
is verified by expansion.
Reorder each basis form $dz_I\wedge d\bar z_J$ into a wedge product of pairs of one-forms,
then apply this identity to each pair. The result is a finite complex linear combination of strongly positive forms.
Every coordinate coefficient of $T$ is therefore a finite linear combination of measures and is itself a measure.

Every strongly positive form is real. The spanning property just proved also shows that their real span
contains all real $(n-p,n-p)$ forms: average a complex linear representation of a real form with its conjugate.
The value of $T$ on a real function times a strongly positive form is real,
so its value on every real test form is real; that is, $T=\bar T$.
\end{proof}

\begin{remark}
Measure coefficients allow a current to be multiplied by a bounded Borel function; arbitrary distributions do not have this property.
Even for a current of order zero, preservation of closedness under extension by zero across a closed set needs a separate proof.
Theorem~\ref{cur:el-mir} proves El Mir extension by controlling the energy of cutoff functions near the excluded set. Finite mass defines the extension; the energy estimate proves that it is closed.
\end{remark}

\section{The Poincar\'e--Lelong Formula and Local Potentials}

\begin{lemma}\label{pos:one-variable-lelong}
On $\C$, with convention \eqref{pos:dc-convention},
$dd^c\log|z|=\delta_0$, where $\delta_0$ is the Dirac mass at the origin.
\end{lemma}
\begin{proof}
The function $\log|z|$ is locally integrable because
$\int_0^\delta r|\log r|\,dr<\infty$.
In one variable, $dd^cv=(2\pi)^{-1}\Delta v\,dx\wedge dy$,
so it suffices to prove $\Delta\log|z|=2\pi\delta_0$.
Write $dA=dx\,dy$ for planar Lebesgue measure. Take $\chi\in C_c^\infty(\C)$ and integrate by parts twice on a large disk
with $B(0,\epsilon)$ removed. The outer boundary contributes nothing because $\chi$ is compactly supported;
the outward normal on the inner boundary is $-\partial_r$, so
\[
 \int_{|z|>\epsilon}\log|z|\Delta\chi\,dA
 =\int_0^{2\pi}
 \bigl[\chi(\epsilon e^{\ii t})
       -\epsilon\log\epsilon\,\partial_r\chi(\epsilon e^{\ii t})\bigr]dt.
\]
The right-hand side tends to $2\pi\chi(0)$. By local integrability, the left-hand side tends to
$\langle\Delta\log|z|,\chi\rangle$, proving the assertion.
\end{proof}

\begin{theorem}[The Poincar\'e--Lelong formula for normal crossings]
\label{pos:poincare-lelong}
Suppose a holomorphic function on a coordinate polydisk has the form
\[
 f(z)=a(z)z_1^{m_1}\cdots z_s^{m_s},\qquad
 a(z)\ne0,\quad m_j\in\N.
\]
Then $\log|f|$ is locally integrable and
\[
 dd^c\log|f|=\sum_{j=1}^s m_j[\{z_j=0\}].
\]
If the zero divisor $D$ of a holomorphic line bundle section $s$ has this form locally,
then every smooth Hermitian metric $h$ satisfies the global current identity
\begin{equation}\label{pos:section-lelong}
 dd^c\log|s|_h=[D]-\frac{\ii}{2\pi}\Theta_h.
\end{equation}
\end{theorem}
\begin{proof}
Local integrability follows from $\log|f|=\log|a|+\sum m_j\log|z_j|$.
After shrinking the chart, the nowhere-vanishing holomorphic function $a$ has a holomorphic logarithm,
so $\partial\dbar\log|a|=0$.
The derivatives of $\log|z_j|$ occur only in the $j$th variable.
Expand an $(n-1,n-1)$ test form in coordinates. Only terms not containing
$dz_j$ or $d\bar z_j$ pair nontrivially with $dd^c\log|z_j|$.
Apply Lemma \ref{pos:one-variable-lelong} in the $j$th variable,
then Fubini's theorem in the remaining variables, to obtain
$dd^c\log|z_j|=[\{z_j=0\}]$.
Summation proves the local formula.

Locally write $s=fe$ and $|e|_h^2=e^{-\varphi}$, so that
$\log|s|_h=\log|f|-\varphi/2$.
The local formula together with
$dd^c(\varphi/2)=\ii\Theta_h/(2\pi)$ gives
\eqref{pos:section-lelong}. Both sides are independent of the frame and therefore patch globally.
\end{proof}

\begin{remark}
The zero set of a general holomorphic function need not have normal crossings. The general Poincar\'e--Lelong formula remains valid,
but its proof also requires multiplicities of analytic hypersurfaces, locally finite volume, and extension arguments near the singular locus.
The present theorem proves only the stated case; one cannot enlarge its scope by simply declaring singular points negligible.
The general version is proved in Theorem \ref{cur:poincare-lelong};
finite mass and closedness of analytic integration currents are established in Theorem \ref{cur:analytic-integration}.
\end{remark}

\begin{proposition}[Local potentials for smooth closed real $(1,1)$ forms]
\label{pos:smooth-local-potential}
Let $\alpha$ be a smooth closed real $(1,1)$ form. Near every point there is a smooth real function $v$
such that $\alpha=\ii\partial\dbar v$. If $\alpha\geq0$, then $v$ is psh.
\end{proposition}
\begin{proof}
On a small ball, Proposition~\ref{bas:poincare} gives a real one-form $\beta$ with $d\beta=\alpha$.
Write $\beta=\beta^{1,0}+\beta^{0,1}$.
Comparison of the $(0,2)$ components gives $\dbar\beta^{0,1}=0$.
Theorem~\ref{sh:dolbeault-lemma} supplies a smooth complex function $w$ such that
$\beta^{0,1}=\dbar w$; reality then gives
$\beta^{1,0}=\partial\bar w$. Consequently,
\[
 \alpha=d\beta
 =\partial\dbar w+\dbar\partial\bar w
 =\partial\dbar(w-\bar w)
 =\ii\partial\dbar(2\operatorname{Im}w).
\]
Take $v=2\operatorname{Im}w$. If $\alpha\geq0$, its coefficient matrix is
the Levi matrix of $v$, so $v$ is psh.
The imaginary part is essential here: the real part would not satisfy the preceding identity.
\end{proof}

General closed positive $(1,1)$ currents also admit local psh potentials.
Extending the argument above to currents requires distributional versions of the Poincar\'e and Dolbeault lemmas,
together with the regularization fact that a distribution with nonnegative complex Hessian has a unique psh representative.
Theorems~\ref{cur:distribution-psh} and~\ref{cur:local-potential} establish these distributional facts and yield the local potentials.

\section{Singular Metrics and Explicit Multiplier Ideals}

\begin{definition}\label{pos:multiplier-definition}
Let $\varphi\in L^1_{\mathrm{loc}}(U)$ be psh, or the sum of a psh function and a smooth real function.
A function locally of the latter form is called \emph{quasi-plurisubharmonic}, or quasi-psh. Its multiplier ideal at $x$ is defined by
\[
 \cI(\varphi)_x
 =\{f\in\cO_{X,x}: |f|^2e^{-\varphi}\text{ is integrable on some neighborhood of }x\}.
\]
This is a \emph{local} condition: the integration domain may have to shrink with the representative and the point.
For a singular line bundle metric $h$, write $|e|_h^2=e^{-\varphi}$ locally
and use this definition to obtain $\cI(h)$.
\end{definition}

\begin{proposition}\label{pos:multiplier-invariance}
The definition gives an ideal sheaf of $\cO_X$, and $\cI(h)$ is independent of the holomorphic frame
and the smooth positive volume form. If $\varphi-\psi$ is locally bounded, then
$\cI(\varphi)=\cI(\psi)$.
\end{proposition}
\begin{proof}
If $f$ and $g$ have integrable weighted squares,
$|f+g|^2\leq2|f|^2+2|g|^2$ proves the same for their sum.
If $a$ is holomorphic, it is bounded after shrinking to a relatively compact neighborhood,
and hence $|af|^2e^{-\varphi}\leq C|f|^2e^{-\varphi}$.
Thus every stalk is an ideal. Integrability is checked locally: cover a compact subset by finitely many such neighborhoods
and add the finite integrals. The local definitions therefore glue to a sheaf.

Under a frame change $e'=be$, where $b$ is nowhere-vanishing and holomorphic,
the new weight is $\varphi' =\varphi-\log|b|^2$.
Both $|b|$ and $|b|^{-1}$ are bounded on a sufficiently small relatively compact neighborhood,
so weighted integrability is unchanged. The ratio of two smooth positive volume forms is likewise locally bounded above and below by positive constants.
Finally, if $|\varphi-\psi|\leq C$, then
$e^{-C}e^{-\psi}\leq e^{-\varphi}\leq e^Ce^{-\psi}$,
which proves the equivalence of the two integrability conditions.
\end{proof}

\begin{theorem}[Multiplier ideals of normal crossing weights]\label{pos:snc-ideal}
Suppose, near the origin, that
\[
 \varphi(z)=\sum_{j=1}^s a_j\log|z_j|^2+b(z),\qquad
 a_j\geq0,
\]
where $b$ is locally bounded. Write $\lfloor a\rfloor$ for the greatest integer at most $a$. Then
\[
 \cI(\varphi)_0
 =\bigl(z_1^{\lfloor a_1\rfloor}\cdots
        z_s^{\lfloor a_s\rfloor}\bigr)\subset\cO_{\C^n,0}.
\]
\end{theorem}
\begin{proof}
Proposition~\ref{pos:multiplier-invariance} allows us to omit $b$.
Integrability of a monomial $z^\alpha$ separates into the one-variable conditions
\[
 \int_0^\delta r^{2\alpha_j-2a_j+1}\,dr<\infty
 \quad\Longleftrightarrow\quad
 \alpha_j>a_j-1
 \quad\Longleftrightarrow\quad
 \alpha_j\geq\lfloor a_j\rfloor.
\]
If $f$ is divisible by the displayed monomial, the quotient is bounded on a small closed polydisk,
and these integrals immediately prove weighted integrability of $f$.

Conversely, expand $f=\sum_{\alpha\in\N_0^n}c_\alpha z^\alpha$.
On the torus with fixed radii $(r_1,\ldots,r_n)$, Fourier orthogonality gives
\[
 \frac1{(2\pi)^n}\int_{[0,2\pi]^n}
 |f(r_1e^{\ii t_1},\ldots,r_ne^{\ii t_n})|^2\,dt
 =\sum_\alpha |c_\alpha|^2r_1^{2\alpha_1}\cdots r_n^{2\alpha_n}.
\]
One may first use the uniformly convergent power series on a strictly smaller polydisk and then pass to the limit.
Multiply by the nonnegative radial weight and integrate; Tonelli's theorem permits interchanging the sum and integral.
If some $c_\alpha\ne0$ and $\alpha_j\leq a_j-1$ for some $j$,
the corresponding monomial integral diverges, and therefore so does the full integral.
Every nonzero coefficient thus satisfies $\alpha_j\geq\lfloor a_j\rfloor$ for each $j$.
Factoring these common powers from the convergent power series proves the required divisibility.
\end{proof}

\begin{proposition}[Radial weights and vanishing orders]\label{pos:radial-ideal}
Let $c\geq0$ and, near the origin in $\C^n$, set
$\varphi=c\log|z|^2$. Writing the maximal ideal as
$\mathfrak m_0=(z_1,\ldots,z_n)$, we have
\[
 \cI(\varphi)_0
 =\mathfrak m_0^{\max\{0,\lfloor c\rfloor-n+1\}}.
\]
In particular, when $c=n+s$ with $s\in\N_0$, weighted integrability forces
all Taylor coefficients of $f$ of degree at most $s$ to vanish.
\end{proposition}
\begin{proof}
For $f\ne0$ of vanishing order $k$, write its homogeneous expansion
$f=P_k+P_{k+1}+\cdots$, with $P_k\ne0$.
Let $d\sigma$ denote the surface measure on the unit sphere $S^{2n-1}$. Homogeneous polynomials of different degrees are orthogonal on this sphere:
replace the spherical variable $\zeta$ by $e^{\ii t}\zeta$ and integrate in $t$
to eliminate the mixed terms of different degrees.
Thus, for small $r$,
\[
 \int_{|\zeta|=1}|f(r\zeta)|^2d\sigma(\zeta)
 =\sum_{\ell\geq k}r^{2\ell}\|P_\ell\|_{L^2(S^{2n-1})}^2.
\]
The nonzero leading term gives a lower bound $C_1r^{2k}$; uniform convergence of the Taylor expansion on a small ball
gives an upper bound $C_2r^{2k}$. The required integrability is therefore equivalent to
\[
 \int_0^\delta r^{2k-2c+2n-1}\,dr<\infty
 \quad\Longleftrightarrow\quad k>c-n.
\]
For an integer $k\geq0$, this says precisely
$k\geq\max\{0,\lfloor c\rfloor-n+1\}$.
A holomorphic germ has vanishing order at least $m$ if and only if it belongs to $\mathfrak m_0^m$,
as follows directly from its convergent power series. Taking $c=n+s$ proves the last assertion.
\end{proof}

\begin{remark}
The strict inequality at the threshold cannot be dropped: when $k=c-n$, the radial integral is
$\int_0^\delta dr/r$, which diverges.
This is why coefficients of logarithmic weights must be recorded precisely in interpolation arguments.
Multiplier ideals of arbitrary psh weights are coherent, but coherence does not follow automatically from the ideal property established here;
it requires local $L^2$ solvability and a finite-generation argument.
\end{remark}

\section{From Estimates to Existence: A Hilbert Complex Lemma}

For a matrix equation $Tu=f$, every vector in $\ker T^*$ is orthogonal to the range of $T$. In infinite-dimensional spaces, that orthogonality alone need not put $f$ in the range, because the range may not be closed. The estimate below supplies the missing norm control. The second operator $S$ encodes the compatibility condition: in the Dolbeault application, both operators are $\dbar$ in successive degrees and $Sf=0$ means that the datum is closed.

In what follows, Hilbert inner products are conjugate-linear in the first variable and linear in the second.
For a densely defined operator $T:H_0\supset\operatorname{Dom}T\to H_1$,
the adjoint satisfies
$\langle Tu,v\rangle=\langle u,T^*v\rangle$.
The notation ``$ST=0$'' includes the domain requirement
$T(\operatorname{Dom}T)\subset\ker S\subset\operatorname{Dom}S$.
These requirements cannot be omitted in arguments involving unbounded operators.

\begin{lemma}[Existence in a Hilbert complex]\label{pos:hilbert-existence}
Let $H_0,H_1,H_2$ be Hilbert spaces, and let
$T:H_0\to H_1$ and $S:H_1\to H_2$ be closed densely defined operators satisfying $ST=0$.
Let $f\in\ker S$. If there is a constant $B\geq0$ such that
\begin{equation}\label{pos:datum-estimate}
 |\langle f,v\rangle|^2
 \leq B\bigl(\|T^*v\|^2+\|Sv\|^2\bigr),
 \qquad v\in\operatorname{Dom}T^*\cap\operatorname{Dom}S,
\end{equation}
then there is $u\in\operatorname{Dom}T$ satisfying
$Tu=f$, $\|u\|^2\leq B$.
Exactly one solution is orthogonal to $\ker T$; it has the smallest norm among all solutions.
\end{lemma}
\begin{proof}
Since $S$ is closed, $K=\ker S$ is a closed subspace. Let $P:H_1\to K$ be the orthogonal projection.
For $v\in\operatorname{Dom}T^*$, write $v=v_0+v_1$, where
$v_0=Pv\in K$,$v_1\perp K$.
Since $\operatorname{Ran}T\subset K$, every $x\in\operatorname{Dom}T$ satisfies
$\langle Tx,v_1\rangle=0$.
By the definition of the adjoint, this means that
$v_1\in\operatorname{Dom}T^*$ and $T^*v_1=0$.
Consequently, $v_0\in\operatorname{Dom}T^*\cap\operatorname{Dom}S$,
$T^*v_0=T^*v$, $Sv_0=0$.
Moreover, $f\in K$ implies $\langle f,v\rangle=\langle f,v_0\rangle$.
Applying \eqref{pos:datum-estimate} to $v_0$ gives
\begin{equation}\label{pos:adjoint-functional-bound}
 |\langle f,v\rangle|\leq\sqrt B\,\|T^*v\|,
 \qquad v\in\operatorname{Dom}T^*.
\end{equation}

Define a linear functional on $\operatorname{Ran}T^*$ by
$\ell(T^*v)=\langle f,v\rangle$.
If $T^*v=T^*w$, apply \eqref{pos:adjoint-functional-bound} to
$v-w$ to obtain $\langle f,v-w\rangle=0$. Thus $\ell$ is well defined.
The same estimate gives $|\ell(a)|\leq\sqrt B\|a\|$.
Extend $\ell$ continuously to $\overline{\operatorname{Ran}T^*}$
and set it equal to zero on the orthogonal complement. This gives a linear functional on $H_0$ of norm at most $\sqrt B$.
The Hilbert-space representation principle gives $u\in H_0$ such that
\[
 \ell(a)=\langle u,a\rangle,\qquad\|u\|\leq\sqrt B.
\]
For completeness, if $\ell\ne0$, project a vector outside its closed kernel onto $(\ker\ell)^\perp$ and normalize the result to a vector $z$ with $\ell(z)=1$. Since $a-\ell(a)z\in\ker\ell$, one has $\ell(a)=\langle z,a\rangle/\|z\|^2$; take $u=z/\|z\|^2$. The case $\ell=0$ gives $u=0$. The orthogonal projection used here exists by the parallelogram identity applied to a sequence minimizing distance to a closed subspace. This proves the representation directly, with $\|u\|=\|\ell\|$.
Thus, for every $v\in\operatorname{Dom}T^*$,
$\langle u,T^*v\rangle=\langle f,v\rangle$.
To verify directly that this implies $Tu=f$, let $G$ be the closed graph of $T$.
The definition of the adjoint gives
\[
 G^\perp=\{(-T^*v,v):v\in\operatorname{Dom}T^*\}.
\]
Indeed, $(a,b)\perp G$ is equivalent to
$\langle Tx,b\rangle=-\langle x,a\rangle$ for every $x\in\operatorname{Dom}T$,
which says that $b\in\operatorname{Dom}T^*$ and $T^*b=-a$.
The identity just obtained implies $(u,f)\perp G^\perp$, hence
$(u,f)\in G^{\perp\perp}=G$. This is exactly $u\in\operatorname{Dom}T$ and $Tu=f$.

The extension above vanishes on $(\operatorname{Ran}T^*)^\perp=\ker T$,
so $u\perp\ker T$.
Any other solution $u'$ satisfies $u'-u\in\ker T$, whence
$\|u'\|^2=\|u\|^2+\|u'-u\|^2$, proving minimality and uniqueness.
\end{proof}

\begin{corollary}\label{pos:coercive-existence}
Under the hypotheses of Lemma~\ref{pos:hilbert-existence}, suppose there is a constant $c>0$ such that
\[
 \|T^*v\|^2+\|Sv\|^2\geq c\|v\|^2
\]
on the common domain. Then every $f\in\ker S$ has a solution $Tu=f$ with
$\|u\|^2\leq c^{-1}\|f\|^2$. In particular, $\operatorname{Ran}T=\ker S$ is closed.
\end{corollary}
\begin{proof}
Cauchy--Schwarz gives
$|\langle f,v\rangle|^2\leq\|f\|^2\|v\|^2
\leq c^{-1}\|f\|^2(\|T^*v\|^2+\|Sv\|^2)$.
Apply Lemma~\ref{pos:hilbert-existence} and use the inclusion $\operatorname{Ran}T\subset\ker S$ to obtain all the assertions.
\end{proof}

\begin{proposition}[Interpreting inverse curvature in the degenerate case]\label{pos:inverse-curvature}
Fix a smooth positive volume form. Let $A(x)$ be a measurable positive semidefinite operator on a finite-rank Hermitian vector bundle, and let $f\in L^2$ for this metric and volume form.
Use the usual inverse on the positive eigenspaces of $A$, and declare the inverse energy to be $+\infty$
whenever $f(x)$ has a nonzero component in $\ker A(x)$.
If
\[
 B=\int\langle f,A^{-1}f\rangle<\infty,
\]
then, for every $v$ for which the right-hand side is defined,
\[
 |\langle f,v\rangle_{L^2}|^2
 \leq B\int\langle v,Av\rangle.
\]
Consequently, for the Hilbert complex $T,S$ of Lemma~\ref{pos:hilbert-existence}, if the common operator domain also satisfies
$\int\langle v,Av\rangle\leq\|T^*v\|^2+\|Sv\|^2$
and $Sf=0$, Lemma~\ref{pos:hilbert-existence} gives a solution with squared norm at most $B$.
\end{proposition}
\begin{proof}
Finite inverse energy implies that $f$ is orthogonal to $\ker A$ almost everywhere.
On the positive eigenspaces at each point, write
$\langle f,v\rangle=\langle A^{-1/2}f,A^{1/2}v\rangle$.
Cauchy--Schwarz, first pointwise and then after integration, proves the estimate.
Combining it with the assumed basic estimate gives precisely
\eqref{pos:datum-estimate}, so Lemma~\ref{pos:hilbert-existence} applies.
Thus the inverse of a positive semidefinite $A$ cannot be treated as an ordinary matrix inverse defined everywhere,
nor can inverse energy in a zero-eigenvalue direction simply be declared zero.
\end{proof}

\section{Extending Test-Form Estimates to Operator Domains}

We use convolution to approximate forms having only weak derivatives by smooth forms.
The approximation must control both the forms themselves and all operators appearing in the estimate. For a differential operator $P$ and multiplication by a scalar function $a$, write $[P,a]u=P(au)-aPu$.

\begin{lemma}[The convolution commutator for first-order operators]\label{pos:friedrichs}
On an open subset of $\R^m$, let
\[
 P=\sum_{j=1}^m a_j(x)\partial_{x_j}+b(x)
\]
be a first-order linear operator with smooth matrix coefficients. If $u\in L^2$ has compact support and $Pu\in L^2$ distributionally,
then convolution kernels with sufficiently small support satisfy
\[
 u*\rho_\epsilon\to u,\qquad
 P(u*\rho_\epsilon)\to Pu
 \quad\text{in }L^2.
\]
The same convolution sequence works simultaneously for finitely many first-order operators.
\end{lemma}
\begin{proof}
Extend the coefficients and functions from a fixed compact neighborhood of the support so that all convolutions take place in Euclidean space.
Distributional integration by parts gives
\begin{align*}
R_\epsilon u(x)
&:=P(u*\rho_\epsilon)(x)-(Pu)*\rho_\epsilon(x)\\
&=\sum_j\int[a_j(x)-a_j(y)]
       \partial_{x_j}\rho_\epsilon(x-y)u(y)\,dy\\
&\quad+\sum_j\int\rho_\epsilon(x-y)(\partial_j a_j)(y)u(y)\,dy\\
&\quad+\int[b(x)-b(y)]\rho_\epsilon(x-y)u(y)\,dy.
\end{align*}
The coefficient-derivative term arises when integration by parts differentiates $a_j(y)$.
Since $a_j$ is Lipschitz on this compact neighborhood, the absolute value of the first kernel is bounded by
$C|x-y||\nabla\rho_\epsilon(x-y)|$.
Its $L^1$ norm is independent of $\epsilon$: scaling gives
\[
 \int |z|\,|\nabla\rho_\epsilon(z)|\,dz
 =\int|w|\,|\nabla\rho(w)|\,dw.
\]
The other two terms likewise have uniformly bounded $L^1$ kernels.
Young's inequality therefore gives
$\|R_\epsilon u\|_2\leq C\|u\|_2$, with $C$ independent of $\epsilon$.

If $u$ is smooth with compact support, then
$P(u*\rho_\epsilon)\to Pu$ and $(Pu)*\rho_\epsilon\to Pu$ directly,
so $R_\epsilon u\to0$.
For arbitrary $u\in L^2$, choose a smooth compactly supported $w$ with $\|u-w\|_2$ as small as desired,
and use
\[
 \|R_\epsilon u\|_2
 \leq C\|u-w\|_2+\|R_\epsilon w\|_2.
\]
First let $\epsilon\to0$, then let $w\to u$, to obtain $R_\epsilon u\to0$.
Together with $(Pu)*\rho_\epsilon\to Pu$ and $u*\rho_\epsilon\to u$, this proves the assertion.
For finitely many operators, apply the same estimate to each one.
\end{proof}

\begin{proposition}[A common core on a compact manifold]\label{pos:compact-core}
Let $X$ be a compact manifold without boundary, with smooth vector bundles and metrics.
For finitely many first-order operators $P_1,\ldots,P_N$ with smooth coefficients, define their common maximal domain by
the distributional conditions $u\in L^2$ and $P_j u\in L^2$ for every $j$.
Smooth sections are dense in this domain for the norm
\[
 \|u\|_2+\sum_j\|P_ju\|_2.
\]
Moreover, the Hilbert adjoint of a maximal operator is the maximal realization of its formal adjoint.
\end{proposition}
\begin{proof}
Choose finitely many coordinate charts, local frames, and a subordinate smooth partition of unity $\chi_\nu$.
Each $\chi_\nu u$ is supported inside a coordinate chart, and
\[
 P_j(\chi_\nu u)=\chi_\nu P_ju+[P_j,\chi_\nu]u.
\]
The commutator is a bounded smooth operator of order zero, so the right-hand side belongs to $L^2$.
Apply Lemma \ref{pos:friedrichs} in each chart, then add the smooth approximations.
Because the cover is finite, the sum of the errors still tends to zero, and
$\sum_\nu\chi_\nu u=u$. This proves the common-core assertion.

Let $P^\dagger$ denote the formal adjoint defined by integration by parts.
If $v\in\operatorname{Dom}P^*$, testing the adjoint identity on smooth compactly supported sections
gives $P^\dagger v=P^*v\in L^2$ distributionally.
Conversely, if the distribution $P^\dagger v$ belongs to $L^2$, then every smooth $u$ satisfies
$\langle Pu,v\rangle=\langle u,P^\dagger v\rangle$.
The graph-norm density just proved extends the identity to all
$u\in\operatorname{Dom}P$, proving $v\in\operatorname{Dom}P^*$.
\end{proof}

\begin{remark}
On a domain with boundary, convolution cannot simply cross the boundary, and adjoint domains may impose boundary conditions.
This is why the proposition explicitly assumes that there is no boundary.
On a noncompact manifold one also needs cutoffs with expanding supports and controlled commutators.
We first carry out this step completely on $\C^n$, then return to compact manifolds.
\end{remark}

\section{Weighted Solvability on Euclidean Space: A Complete Model}

On $\C^n$, let $dV$ be Lebesgue measure and declare the forms $d\bar z_J$ to be an orthonormal basis. For a smooth real weight $\varphi$,
define
\[
 \|u\|_\varphi^2=\int_{\C^n}\sum_{|J|=q}'|u_J|^2e^{-\varphi}\,dV.
\]
The prime means that the multi-index is strictly increasing. For $|K|=q-1$, let $u_{jK}$ denote
the coefficient extended by antisymmetry; if $j$ occurs in $K$, put $u_{jK}=0$.

\begin{proposition}[The weighted basic identity]\label{pos:euclidean-identity}
For $1\le q\le n$, every smooth compactly supported $(0,q)$ form $u$ satisfies
\begin{align}
 \|\dbar u\|_\varphi^2+\|\dbar_\varphi^*u\|_\varphi^2
 &=\sum_{|J|=q}'\sum_{j=1}^n
          \|\partial_{\bar z_j}u_J\|_\varphi^2\notag\\
 &\quad+\int\sum_{|K|=q-1}'\sum_{j,k}
      \varphi_{j\bar k}u_{jK}\bar u_{kK}e^{-\varphi}\,dV.
 \label{pos:euclidean-bkn}
\end{align}
If $\bigl(\varphi_{j\bar k}\bigr)\geq cI$, the right-hand side is at least
$qc\|u\|_\varphi^2$.
\end{proposition}
\begin{proof}
Let $a_j$ be exterior multiplication by $d\bar z_j$ and $b_j=a_j^*$ the corresponding contraction.
They satisfy
$b_ja_k+a_kb_j=\delta_{jk}$, together with the anticommutation relations for operators of the same type.
Put $\varphi_j=\partial_{z_j}\varphi$ and $\delta_j=\partial_{z_j}-\varphi_j$.
Weighted integration by parts gives
$\partial_{\bar z_j}^*=-\delta_j$, so
\[
 \dbar=\sum_ja_j\partial_{\bar z_j},\qquad
 \dbar_\varphi^*=-\sum_jb_j\delta_j.
\]
Multiplying out and using the anticommutation relations and
$[\partial_{\bar z_k},\delta_j]=-\varphi_{j\bar k}$ gives
\[
 \dbar\dbar_\varphi^*+\dbar_\varphi^*\dbar
 =-\sum_j\delta_j\partial_{\bar z_j}
     +\sum_{j,k}\varphi_{j\bar k}a_kb_j.
\]
Apply this identity to $u$ and take its $L^2$ inner product with $u$.
Another integration by parts in the first term gives
$\sum_{J,j}'\|\partial_{\bar z_j}u_J\|_\varphi^2$.
In the second term, use $a_k^*=b_k$ to obtain
$\sum_{K,j,k}'\varphi_{j\bar k}u_{jK}\bar u_{kK}$, which proves the identity.
If the Hessian is bounded below by $cI$, this quadratic form is at least
$c\sum_{K,j}'|u_{jK}|^2$.
Every increasing $q$-index $J$ occurs exactly $q$ times in the latter sum,
so $\sum_{K,j}'|u_{jK}|^2=q\sum_J'|u_J|^2$.
\end{proof}

\begin{lemma}[A weighted common core on the whole space]\label{pos:euclidean-core}
In these weighted $L^2$ spaces, give $\dbar$ its maximal distributional realization.
Smooth compactly supported forms are dense in
$\operatorname{Dom}\dbar\cap\operatorname{Dom}\dbar_\varphi^*$
for the common graph norm.
\end{lemma}
\begin{proof}
First identify the adjoint domain. Testing an element of the Hilbert adjoint domain against smooth compactly supported forms puts it in the maximal distributional domain of the formal adjoint.
For the converse, apply the cutoff and smoothing argument below to the maximal domain of $\dbar$ alone and pass to the limit in integration by parts.
For this domain identification, the cutoff and smoothing argument is applied to $\dbar$ alone; simultaneous approximation is then obtained by using it for both differential expressions.
Choose a fixed smooth cutoff $\chi$ equal to one on the unit ball,
zero outside the ball of radius two, and satisfying $0\leq\chi\leq1$. Put $\chi_R(z)=\chi(z/R)$.
Then $|d\chi_R|\leq C/R$. For $u$ in the common domain,
\[
 \dbar(\chi_Ru)=\chi_R\dbar u+(\dbar\chi_R)\wedge u,
\]
\[
 \dbar_\varphi^*(\chi_Ru)
 =\chi_R\dbar_\varphi^*u-\sum_j(\partial_{z_j}\chi_R)b_ju.
\]
Both commutators have weighted $L^2$ norm at most $C_qR^{-1}\|u\|_\varphi$.
The tails outside the supports tend to zero by $L^2$ integrability, so $\chi_Ru\to u$ in the common graph norm.

For fixed $R$, $e^{-\varphi}$ and $e^{\varphi}$ are bounded on a slightly larger compact set.
The weighted and unweighted $L^2$ norms are equivalent there.
Apply Lemma \ref{pos:friedrichs} simultaneously to $\dbar$ and its formal adjoint
on $\chi_Ru$ to obtain smooth compactly supported approximations.
The identification of the Hilbert adjoint with the maximal formal adjoint follows by integration by parts against smooth test forms,
followed by the same cutoff and smoothing approximation on the maximal domain, exactly as in
Proposition \ref{pos:compact-core}.
Finally let the smoothing scale tend to zero and choose a diagonal sequence with $R\to\infty$.
\end{proof}

\begin{theorem}[Solvability with a strictly psh weight on the whole space]\label{pos:euclidean-solving}
Let $\varphi\in C^\infty(\C^n,\R)$ satisfy
$\bigl(\varphi_{j\bar k}\bigr)\geq cI$ for some $c>0$.
If $1\leq q\leq n$ and $f$ is a weighted $L^2$ $(0,q)$ form
with $\dbar f=0$ in the distributional sense, then there is a weighted $L^2$
$(0,q-1)$ form $u$ such that
\[
 \dbar u=f,\qquad
 \|u\|_\varphi^2\leq\frac1{qc}\|f\|_\varphi^2.
\]
There is a unique minimal-norm solution, characterized by orthogonality to all closed weighted $L^2$ $(0,q-1)$ forms.
\end{theorem}
\begin{proof}
On the weighted $L^2$ spaces in adjacent degrees, set $T=\dbar$ and $S=\dbar$,
both with maximal distributional domains. If $q=n$, take $S=0$ on the entire top-degree space.
These operators are densely defined because smooth compactly supported forms are dense.
They are closed: if $u_\nu\to u$ and $\dbar u_\nu\to v$ in weighted $L^2$,
then convergence also holds in ordinary $L^2$ on compact sets, and hence in distributions;
continuity of distributional differentiation gives $\dbar u=v$.
The distributional identity $\dbar^2=0$ also proves $ST=0$, including the domain requirement.

Proposition \ref{pos:euclidean-identity} gives, for smooth compactly supported forms, the estimate
$\|T^*v\|^2+\|Sv\|^2\geq qc\|v\|^2$.
Lemma \ref{pos:euclidean-core} extends it to the common domain.
Since $f\in\ker S$, Corollary \ref{pos:coercive-existence} gives the solution and its norm bound.
Minimality and uniqueness follow from the last paragraph of Lemma \ref{pos:hilbert-existence}.
\end{proof}

\begin{example}[Sharpness of the estimate]\label{pos:gaussian-example}
Take $\varphi(z)=|z|^2$, $q=1$, and $f=d\bar z_1$.
The function $u=\bar z_1$ satisfies $\dbar u=f$, and
\[
 \|f\|_\varphi^2=\pi^n,\qquad
 \|u\|_\varphi^2=\pi^n.
\]
Thus the constant cannot be decreased uniformly when $q=c=1$.
\end{example}
\begin{proof}
The equation follows by differentiation of the coefficients.
Fubini's theorem and the one-variable polar-coordinate integrals
$\int_\C e^{-|z|^2}dA=\pi$,
$\int_\C|z|^2e^{-|z|^2}dA=\pi$ give the two norms.
To show that no solution has smaller norm, again write $T=\dbar$ and take $v=d\bar z_1$.
Both $v$ and its formal adjoint image $T^*v=\bar z_1$ belong to weighted $L^2$.
Lemma~\ref{pos:euclidean-core} identifies the formal and Hilbert adjoints, so this is an admissible test form.
For any solution $w$,
\[
 \pi^n=\langle f,v\rangle_\varphi
       =\langle w,T^*v\rangle_\varphi
       \leq\|w\|_\varphi\,\pi^{n/2}.
\]
It follows that $\|w\|_\varphi^2\geq\pi^n$, proving sharpness.
\end{proof}

\section{K\"ahler Commutators and the Bochner--Kodaira--Nakano Identity}

Let $(X,\omega)$ be K\"ahler and $(E,h)$ a Hermitian holomorphic vector bundle.
Use the Chern connection of Theorem~\ref{bun:chern-connection}, written $D=D'+D''$, with $D''=\dbar$, and curvature $F=D^2$. Norms and formal adjoints in this section use $h$, $\omega$, and $dV_\omega=\omega^n/n!$.
Regard $F$ as the operator given by exterior multiplication combined with its action on the $E$ component.
Define $L u=\omega\wedge u$, $\Lambda=L^*$, and
\[
 \Delta'=D'D'^*+D'^*D',\qquad
 \Delta''=D''D''^*+D''^*D''.
\]
We first work with differential expressions on smooth forms; global operator domains will be specified when estimates are used.

\begin{lemma}[First-order normalizations]\label{pos:normal-frames}
Near every $x\in X$ there are holomorphic coordinates and a holomorphic bundle frame for which, at $x$,
$\omega=\ii\sum_j dz_j\wedge d\bar z_j$, and all first derivatives
of the base metric and the bundle metric matrix vanish.
\end{lemma}
\begin{proof}
First use a complex linear coordinate change to make $g_{j\bar k}(x)$ the identity, and place $x$ at the origin.
The K\"ahler condition $d\omega=0$ gives
$\partial_\ell g_{j\bar k}(0)=\partial_jg_{\ell\bar k}(0)$.
Put
$\Gamma^k_{j\ell}=\partial_jg_{\ell\bar k}(0)$; these coefficients are symmetric in $j,\ell$.
Make the local holomorphic coordinate change
\[
 z_k=w_k-\frac12\sum_{j,\ell}\Gamma^k_{j\ell}w_jw_\ell.
\]
Its differential at the origin is the identity, so it is locally invertible.
In the first-order metric transformation formula, the original $\partial g$ terms cancel the quadratic coordinate-change terms,
giving $\partial g(0)=0$ in the new coordinates; conjugation also gives $\bar\partial g(0)=0$.

For the bundle metric, first make a constant frame change so that $H(0)=I$, and write
\[
 H(z)=I+\sum_j(A_jz_j+A_j^*\bar z_j)+O(|z|^2).
\]
Change the frame by right multiplication by the holomorphic matrix $B(z)=I-\sum_jA_jz_j$.
The new metric is $B^*HB=I+O(|z|^2)$,
so all its first derivatives vanish as well.
\end{proof}

\begin{proposition}[The K\"ahler commutation identities]\label{pos:kahler-identities}
With these conventions,
\[
 [\Lambda,D']=\ii D''^*,\qquad
 [\Lambda,D'']=-\ii D'^*.
\]
Here $[A,B]=AB-BA$. Since $\Lambda$ has even degree,
this agrees with the graded commutator in these identities.
\end{proposition}
\begin{proof}
These are identities of first-order differential operators, so they can be checked pointwise in the coordinates and frame of Lemma~\ref{pos:normal-frames}.
At the chosen point, let $e_j,\bar e_j$ denote exterior multiplication by $dz_j,d\bar z_j$,
and let $\iota_j,\bar\iota_j$ be their Hermitian adjoints.
Exterior multiplication and contraction, of the same or different types, obey the exterior-algebra anticommutation rules;
for example, $\iota_je_k+e_k\iota_j=\delta_{jk}$ and
$\bar\iota_je_k+e_k\bar\iota_j=0$.
At this point,
\[
 L=\ii\sum_je_j\bar e_j,\qquad
 \Lambda=-\ii\sum_j\bar\iota_j\iota_j.
\]
Applying the anticommutation rules term by term gives
\[
 [\Lambda,e_k]=-\ii\bar\iota_k,\qquad
 [\Lambda,\bar e_k]=\ii\iota_k.
\]
Write $\nabla_j=D_{\partial/\partial z_j}$ and $\nabla_{\bar j}=D_{\partial/\partial\bar z_j}$ for covariant derivatives. The normalizations eliminate additional first-order coefficient terms from the metrics, volume, and connection at this point. Thus
\[
 D'=\sum_je_j\nabla_j,\quad D''=\sum_j\bar e_j\nabla_{\bar j},\quad
 D'^*=-\sum_j\iota_j\nabla_{\bar j},\quad
 D''^*=-\sum_j\bar\iota_j\nabla_j.
\]
Multiply the two exterior-algebra commutators by the corresponding covariant derivatives and sum
to obtain the desired identities. Since the point was arbitrary, they hold globally.
\end{proof}

\begin{theorem}[The Bochner--Kodaira--Nakano identity]\label{pos:bkn}
On smooth $E$-valued forms on a K\"ahler manifold,
\begin{equation}\label{pos:bkn-formula}
 \Delta''=\Delta'+[\ii F,\Lambda].
\end{equation}
If $X$ is compact without boundary, or if the form has compact support, then
\begin{align}
 \|D''u\|^2+\|D''^*u\|^2
 &=\|D'u\|^2+\|D'^*u\|^2
   +\int_X\langle u,[\ii F,\Lambda]u\rangle\,dV_\omega\notag\\
 &\geq\int_X\langle u,[\ii F,\Lambda]u\rangle\,dV_\omega.
 \label{pos:basic-estimate}
\end{align}
\end{theorem}
\begin{proof}
Proposition~\ref{pos:kahler-identities} gives $D''^*=-\ii[\Lambda,D']$ and
$D'^*=\ii[\Lambda,D'']$.
Substitution into the two Laplacians and complete expansion give
\begin{align*}
 \Delta''&=-\ii\bigl(D''\Lambda D'-D''D'\Lambda
                      +\Lambda D'D''-D'\Lambda D''\bigr),\\
 \Delta'&=\ii\bigl(D'\Lambda D''-D'D''\Lambda
                      +\Lambda D''D'-D''\Lambda D'\bigr).
\end{align*}
On subtracting, the terms containing $D''\Lambda D'$ and $D'\Lambda D''$ cancel, leaving
\[
 \Delta''-\Delta'
 =\ii(D'D''+D''D')\Lambda-\ii\Lambda(D'D''+D''D').
\]
Theorem~\ref{bun:chern-connection} shows that Chern curvature has type $(1,1)$, so $D'^2=D''^2=0$ and
$D'D''+D''D'=D^2=F$.
This proves \eqref{pos:bkn-formula}.
Integration by parts, using compact support or compactness without boundary, gives
$\langle u,\Delta'u\rangle=\|D'u\|^2+\|D'^*u\|^2$,
and the analogous identity for $\Delta''$. This proves the equality; dropping the two nonnegative squared norms gives the inequality.
\end{proof}

\begin{proposition}[All eigenvalues of the line bundle curvature operator]\label{pos:curvature-eigenvalues}
Suppose $E$ is a line bundle, and choose a unitary coframe at a point so that
\[
 \omega=\ii\sum_jdz_j\wedge d\bar z_j,\qquad
 \ii F=\ii\sum_j\lambda_jdz_j\wedge d\bar z_j.
\]
Then the eigenvalue of $[\ii F,\Lambda]$ on the $(p,q)$ basis form
$dz_I\wedge d\bar z_J$ is
\begin{equation}\label{pos:line-eigenvalue}
 \sum_{j\in I}\lambda_j+\sum_{j\in J}\lambda_j-\sum_{j=1}^n\lambda_j.
\end{equation}
In particular, on $(n,q)$ forms it is $\sum_{j\in J}\lambda_j$;
if $\omega=\ii F>0$, the operator on $(p,q)$ forms is
$(p+q-n)\operatorname{Id}$.
\end{proposition}
\begin{proof}
Decompose the exterior algebra by coordinate directions. The $j$th direction has four basis states:
$1,dz_j,d\bar z_j,dz_j\wedge d\bar z_j$.
The commutator of exterior multiplication by $\ii dz_j\wedge d\bar z_j$ with its adjoint
acts on these four states by $-1,0,0,1$, respectively.
This follows by checking contraction followed by multiplication and multiplication followed by contraction on each state.
Even-degree operators belonging to different directions commute. Thus the $j$th contribution is
$\lambda_j(\mathbf1_{j\in I}+\mathbf1_{j\in J}-1)$, where $\mathbf1_{j\in I}$ is one when $j\in I$ and zero otherwise.
Summing gives \eqref{pos:line-eigenvalue}.
If $I=\{1,\ldots,n\}$, the first and last sums cancel;
if all $\lambda_j=1$, the result is $p+q-n$.
\end{proof}

\begin{remark}
Even when line bundle curvature is positive in every direction, with an arbitrary fixed background metric
its operator on general $(p,q)$ forms cannot simply be described as
the sum of a number of smallest eigenvalues.
For example, take $n=3$, $p=q=2$, $I=J=\{1,2\}$, and
$(\lambda_1,\lambda_2,\lambda_3)=(1,1,100)$.
Formula \eqref{pos:line-eigenvalue} gives $-98$.
This explains the role of choosing $\omega=\ii F$ in AKN vanishing.
For a general vector bundle, curvature is not even a scalar $(1,1)$ form with only $n$ eigenvalues;
the bundle indices introduced below must be retained.
\end{remark}

\begin{definition}\label{pos:nakano-definition}
In unitary tangent and bundle frames, write
\[
 F e_\alpha=\sum_{j,k,\beta}R_{j\bar k\alpha\bar\beta}
             dz_j\wedge d\bar z_k\otimes e_\beta.
\]
Thus the first bundle index is the input index and the second is the output index: $R_{j\bar k\alpha\bar\beta}=h(e_\beta,F_{j\bar k}e_\alpha)$. This specifies all conjugations for our inner-product convention.
The bundle $(E,h)$ is Nakano positive if, for every nonzero array
$w=(w^{j\alpha})$,
\[
 \sum_{j,k,\alpha,\beta}
 R_{j\bar k\alpha\bar\beta}w^{j\alpha}\bar w^{k\beta}>0.
\]
If the inequality is required only for nonzero rank-one arrays $w^{j\alpha}=\xi_jv^\alpha$,
the bundle is Griffiths positive. These definitions are invariant under unitary frame changes because the arrays transform
as tensors in $T^{1,0}X\otimes E$, and the quadratic form is the pairing with the same tensor.
\end{definition}

\begin{proposition}[Nakano positivity in top holomorphic degree]
\label{pos:nakano-operator}
For $1\le q\le n$, write an $E$-valued $(n,q)$ form as $u=dz_1\wedge\cdots\wedge dz_n\wedge\sum_{J,\alpha}u_J^\alpha d\bar z_J\otimes e_\alpha$ in unitary frames. Extend the coefficients $u_{jK}^\alpha$ antisymmetrically as in Proposition~\ref{pos:euclidean-identity}. Then
\begin{equation}\label{pos:nakano-quadratic}
 \langle u,[\ii F,\Lambda]u\rangle
 =\sum_{|K|=q-1}'\sum_{j,k,\alpha,\beta}
 R_{j\bar k\alpha\bar\beta}
 u_{jK}^{\alpha}\bar u_{kK}^{\beta}.
\end{equation}
Consequently, if the Nakano quadratic form is bounded below by $c|w|^2$, then
$[\ii F,\Lambda]\geq qc\operatorname{Id}$ on $(n,q)$ forms.
\end{proposition}
\begin{proof}
In a normal frame, let $C_{j\bar k}$ be the curvature coefficient matrices,
so that $C_{j\bar k}e_\alpha=\sum_\beta R_{j\bar k\alpha\bar\beta}e_\beta$.
Since the holomorphic degree is already $n$, exterior multiplication by any $dz_j$ gives zero.
Expand $\ii F=\ii\sum C_{j\bar k}e_j\bar e_k$ and
$\Lambda=-\ii\sum\bar\iota_\ell\iota_\ell$.
In top holomorphic degree, the term applying $F$ first vanishes;
if $\iota_\ell$ first removes a holomorphic factor, only $j=\ell$ can contribute when a factor is inserted again.
Commuting through the remaining factors with the anticommutation rules leaves exactly
\[
 [\ii F,\Lambda]u=\sum_{j,k}C_{j\bar k}\bar e_k\bar\iota_j u.
\]
Pair with $u$ and transfer $\bar e_k$ to its adjoint $\bar\iota_k$
to obtain \eqref{pos:nakano-quadratic}.
For each fixed $K$, apply the Nakano lower bound to the array $(u_{jK}^\alpha)$.
Then use $\sum_{K,j,\alpha}'|u_{jK}^\alpha|^2=q|u|^2$ to conclude.
\end{proof}

\section{From Weak Solutions to Smooth Solutions and Compact Vanishing}

This section needs only Theorem~\ref{ana:elliptic-regularity} from Chapter~\ref{ana:chapter}:
\emph{If a smooth second-order operator has positive definite scalar principal symbol,
then $u\in L^2_{\mathrm{loc}}$ and smoothness of $Pu$ imply smoothness of $u$.}
This is an independent local analytic input. Existence is established directly by Lemma~\ref{pos:hilbert-existence},
which constructs a preimage before any use of a global Hodge decomposition.

\begin{theorem}[Curvature solvability on a compact K\"ahler manifold]\label{pos:compact-solving}
Let $(X,\omega)$ be compact K\"ahler and $(E,h)$ a smooth Hermitian holomorphic vector bundle.
Fix $p,q$ with $q\geq1$. If, on $E$-valued $(p,q)$ forms,
\[
 A=[\ii F,\Lambda]\geq c\operatorname{Id},\qquad c>0,
\]
then every smooth $\dbar$-closed $E$-valued $(p,q)$ form $f$
has a smooth $E$-valued $(p,q-1)$ form $u$ satisfying
\[
 \dbar u=f,\qquad \|u\|^2\leq c^{-1}\|f\|^2.
\]
In particular, $H^{p,q}_{\dbar}(X,E)=0$.
\end{theorem}
\begin{proof}
Take the maximal distributional realizations $T=D''$ and $S=D''$ in adjacent $L^2$ degrees.
Closedness, dense definition, and $ST=0$ follow exactly as in
Theorem \ref{pos:euclidean-solving}, working locally in coordinates and frames.
The estimate \eqref{pos:basic-estimate} gives, on smooth forms,
\[
 \|T^*v\|^2+\|Sv\|^2\geq c\|v\|^2.
\]
Proposition \ref{pos:compact-core} extends this to the common operator domain.
Lemma~\ref{pos:hilbert-existence} therefore gives a minimal-norm $L^2$ solution $u$ satisfying
$Tu=f$, $u\perp\ker T$, and the asserted bound.

We now prove smoothness. For every smooth $(p,q-2)$ form $a$,
$D''a\in\ker T$, so $\langle u,D''a\rangle=0$.
This is exactly the distributional identity $D''^*u=0$.
If $q=1$, negative antiholomorphic degrees do not occur, and this identity is automatic.
Thus
\[
 \Delta''u=D''D''^*u+D''^*D''u=D''^*f
\]
holds distributionally, with a smooth right-hand side.
In unitary coordinates at any point, the principal part of $\Delta''$ is
$-\sum_j\partial_{z_j}\partial_{\bar z_j}$ times the identity matrix,
a positive multiple of the nonnegative real Laplacian principal part; its scalar principal symbol remains positive definite under coordinate changes.
Since $u\in L^2$, Theorem~\ref{ana:elliptic-regularity} makes $u$ smooth.
Every smooth closed form therefore has a smooth preimage, proving Dolbeault cohomology vanishing.
\end{proof}

\begin{corollary}[Kodaira vanishing]\label{pos:kodaira}
Let $L$ be a holomorphic line bundle on a compact complex manifold $X$ with a smooth positive-curvature metric.
Then
\[
 H^q(X,K_X\otimes L)=0,\qquad q\geq1.
\]
\end{corollary}
\begin{proof}
The curvature form $\omega=\ii F_L$ is real, closed, and positive, so it is itself a K\"ahler metric.
By Proposition \ref{pos:curvature-eigenvalues}, on $L$-valued $(n,q)$ forms,
$[\ii F_L,\Lambda]=q\operatorname{Id}$.
Theorem \ref{pos:compact-solving} gives $H^{n,q}_{\dbar}(X,L)=0$.
Theorem~\ref{sh:dolbeault} identifies this group with
$H^q(X,\Omega_X^n\otimes L)=H^q(X,K_X\otimes L)$.
\end{proof}

\begin{corollary}[Akizuki--Kodaira--Nakano vanishing]\label{pos:akn}
Under the assumptions of Corollary~\ref{pos:kodaira},
\[
 H^q(X,\Omega_X^p\otimes L)=0\qquad\text{if }p+q>n.
\]
\end{corollary}
\begin{proof}
Again choose $\omega=\ii F_L$.
On $(p,q)$ forms, the curvature operator equals $(p+q-n)\operatorname{Id}$,
which has a strictly positive lower bound when $p+q>n$.
In this range $q\geq1$, since $p\leq n$.
Apply Theorem~\ref{pos:compact-solving} and the Dolbeault isomorphism, Theorem~\ref{sh:dolbeault}.
\end{proof}

\begin{corollary}[Nakano vanishing]\label{pos:nakano-vanishing}
Let $(E,h)$ be a Nakano positive bundle on a compact K\"ahler manifold $(X,\omega)$.
Then $H^q(X,K_X\otimes E)=0$ for all $q\geq1$.
\end{corollary}
\begin{proof}
The Nakano quadratic form is a continuous strictly positive function on the unit sphere bundle of $T^{1,0}X\otimes E$.
The base and sphere fibers are compact, so it has a positive minimum $c>0$.
Proposition \ref{pos:nakano-operator} gives the operator bound on $(n,q)$ forms
$[\ii F,\Lambda]\geq qc\operatorname{Id}$.
Theorems~\ref{pos:compact-solving} and~\ref{sh:dolbeault} prove the conclusion.
\end{proof}

For an effective divisor $D$ locally given by a nonzero holomorphic equation $f_D$, the sheaf $\cO_X(D)$ consists locally of meromorphic multiples $a/f_D$ with $a$ holomorphic. It is a line bundle, dual to the ideal line bundle $\cO_X(-D)=(f_D)$. Local equations differ by units, so these local definitions glue.

\begin{corollary}[Singular-metric vanishing for an integral divisor]
\label{pos:integral-divisor-nadel}
Let $X$ be compact, let $A$ be a smooth positive line bundle, and let $D$ be an effective integral divisor with normal crossing support.
Give $L=A\otimes\cO(D)$ the singular metric $h$ with local weight
$\varphi_A+\log|f_D|^2$, where $f_D$ is a local equation for $D$.
Then
\[
 H^q(X,K_X\otimes L\otimes\cI(h))=0,\qquad q\geq1.
\]
\end{corollary}
\begin{proof}
Theorem~\ref{pos:snc-ideal} has no rounding error for integral coefficients, and gives
$\cI(h)=\cO(-D)$.
Consequently,
$K_X\otimes L\otimes\cI(h)
 =K_X\otimes A\otimes\cO(D)\otimes\cO(-D)
 \simeq K_X\otimes A$.
The higher cohomology of the right-hand side vanishes by Corollary~\ref{pos:kodaira}.
This proves the stated singular case without using the general Nadel theorem.
\end{proof}

\begin{corollary}[Quantitative solvability for high powers of a positive line bundle]\label{pos:high-power-solving}
Let $(X,\omega)$ be compact K\"ahler and $(L,h)$ a positive line bundle.
Choose constants $a>0$, $b\geq0$ such that
\[
 \ii F_L\geq a\omega,\qquad -b\omega\leq\ii F_{K_X}\leq b\omega,
\]
where the metric on $K_X$ is induced by $\omega$.
For every integer $k\geq\max\{1,2b/a\}$ and every smooth $\dbar$-closed
$L^k$-valued $(0,1)$ form $f$, there is a smooth $L^k$-valued function $u$ such that
\begin{equation}\label{pos:high-power-estimate}
 \dbar u=f,\qquad
 \|u\|_{h^k,\omega}^2\leq\frac{2}{ak}\|f\|_{h^k,\omega}^2.
\end{equation}
\end{corollary}
\begin{proof}
Compactness bounds the relative curvature eigenvalues uniformly, so such constants $a,b$ exist.
Put $E_k=L^k\otimes K_X^{-1}$ with its product metric. The induced-connection calculation of Proposition~\ref{bun:chern-functors} gives
\[
 \ii F_{E_k}=k\ii F_L-\ii F_{K_X}
 \geq(ka-b)\omega\geq\frac{ka}{2}\omega.
\]
Thus on $E_k$-valued $(n,1)$ forms the curvature operator is at least $(ka/2)\operatorname{Id}$.

To specify the identification precisely, for a local nowhere-zero holomorphic $n$-form $\tau$, define
\[
 I_q(v)=(-1)^{nq}\tau\wedge v\otimes\tau^{-1}.
\]
The two frame-change factors cancel, so this is a globally defined map
from $L^k$-valued $(0,q)$ forms to $E_k$-valued $(n,q)$ forms.
Because $K_X$ carries the induced metric, $I_q$ is a pointwise isometry.
The inserted sign absorbs the factor $(-1)^n$ from the Leibniz rule, giving
$\dbar I_q=I_{q+1}\dbar$.

Apply Theorem \ref{pos:compact-solving} to the closed form $I_1(f)$
to obtain a smooth $U$ satisfying
\[
 \dbar U=I_1(f),\qquad \|U\|^2\leq\frac{2}{ak}\|I_1(f)\|^2.
\]
Set $u=I_0^{-1}U$. The isometry and compatibility with the complexes give
\eqref{pos:high-power-estimate}.
\end{proof}

\section{From Smooth Estimates to the General Theory}

The Hilbert complex argument and the BKN calculation provide the common foundation for the general theory developed next.
Chapter \ref{cur:chapter} proves local potentials for positive $(1,1)$ currents, El Mir extension, and analytic integration currents;
Chapter \ref{l2:chapter} develops complete metrics, singular-weight solvability on pseudoconvex domains, and coherence of multiplier ideals;
Chapter \ref{div:chapter} proves general Skoda division and the Brian\c{c}on--Skoda theorem in detail;
Chapter \ref{int:chapter} establishes Skoda integrability;
and Chapter \ref{nad:chapter} uses controlled global regularization to prove general Nadel vanishing.
We finish with two immediately accessible bridges that recur throughout those proofs.

\begin{lemma}[Integral dependence implies a growth bound]\label{pos:integral-growth}
Let $I=(g_1,\ldots,g_r)\subset\cO_{\C^n,0}$, and suppose a germ $f$ satisfies
\[
 f^d+a_1f^{d-1}+\cdots+a_d=0,\qquad a_j\in I^j.
\]
Such a monic relation says that $f$ is \emph{integral over} $I$; $I^j$ is the ordinary product ideal. Then $|f|\leq C|g|$ on a neighborhood of the origin, where
$|g|=(\sum|g_j|^2)^{1/2}$.
\end{lemma}
\begin{proof}
Write each $a_j$ as a finite sum of products of $j$ generators with holomorphic coefficients.
Bound these coefficients on a small relatively compact neighborhood to obtain
$|a_j|\leq C_j|g|^j$.
At a fixed point, put $M=\max_j|a_j|^{1/j}$.
If $|f|>2M$, divide the integral equation by $f^d$ to obtain
\[
 1\leq\sum_{j=1}^d\frac{|a_j|}{|f|^j}
 \leq\sum_{j=1}^d\left(\frac M{|f|}\right)^j
 <\sum_{j=1}^\infty2^{-j}=1,
\]
a contradiction. Hence $|f|\leq2M\leq2\max_j C_j^{1/j}|g|$.
If $M=0$, the equation gives $f=0$ directly, so the conclusion still holds.
\end{proof}

The converse, that a growth bound implies integral dependence, does not follow by reversing this root bound;
Theorem \ref{div:integral-closure} proves it using Skoda division, Noetherianity, and the determinant trick.

\begin{proposition}[From vanishing to interpolation]\label{pos:jet-consequence}
Let $B$ be a holomorphic line bundle, let $\mathcal B=\cO_X(B)$ be its sheaf of holomorphic sections, and let $\cJ\subset\cO_X$ be a coherent ideal sheaf.
Suppose $\cO_X/\cJ$ has finite support
$\{x_1,\ldots,x_N\}$ and $H^1(X,B\otimes\cJ)=0$.
Then the restriction map
\[
 H^0(X,B)\longrightarrow
 \bigoplus_{j=1}^N \mathcal B_{x_j}\otimes_{\cO_{X,x_j}}
                  (\cO_X/\cJ)_{x_j}
\]
is surjective. If $\cJ_{x_j}\subset\mathfrak m_{x_j}^{s_j+1}$,
it in particular realizes arbitrary prescribed jets of order $s_j$ at these points. Here $\mathfrak m_x$ is the maximal ideal of germs vanishing at $x$, and a jet of order $s$ is a class in $\mathcal B_x/\mathfrak m_x^{s+1}\mathcal B_x$; for $s=0$ this quotient is the fiber $B_x$.
\end{proposition}
\begin{proof}
A line bundle is locally free, so Proposition~\ref{sh:tensor} shows that tensoring preserves short exact sequences and gives
\[
 0\longrightarrow B\otimes\cJ\longrightarrow B
 \longrightarrow B\otimes(\cO_X/\cJ)\longrightarrow0.
\]
In the long exact sequence of Proposition~\ref{sh:long-exact}, $H^1(X,B\otimes\cJ)=0$ makes
$H^0(X,B)\to H^0(X,B\otimes(\cO_X/\cJ))$ surjective.
Proposition~\ref{sh:discrete-support} identifies global sections of a sheaf with finite support with the direct sum of its stalks at the support points:
choose disjoint small neighborhoods and represent each germ on a sufficiently small neighborhood;
the sheaf vanishes away from the support point there, so the representative glues to the zero section outside.
This proves the stated description of the target.

If $\cJ_{x_j}\subset\mathfrak m_{x_j}^{s_j+1}$, there is a quotient map
$\cO_{X,x_j}/\cJ_{x_j}\twoheadrightarrow
 \cO_{X,x_j}/\mathfrak m_{x_j}^{s_j+1}$.
Lift each prescribed jet to the first quotient, then use the surjectivity just proved to lift to a global section.
\end{proof}

\begin{remark}
In an analytic interpolation construction, first cut off a local section to obtain $\chi s$, then solve
$\dbar u=\dbar(\chi s)$.
The correction $u$ is smooth or $L^2$ and is generally not holomorphic;
the holomorphic section is $\chi s-u$.
If $u$ were assumed from the outset to lie in $H^0(X,B)$, the left-hand side would vanish,
and in general could not equal the specified right-hand side.
\end{remark}

Singular positivity has different consequences from smooth positivity. In particular, strictly positive singular curvature need not produce a smooth positive metric, and it does not by itself imply vanishing for high powers with one fixed multiplier ideal. Examples~\ref{nad:regularization-counterexamples} and~\ref{nad:fixed-ideal-counterexample} illustrate these two distinctions. For a measurable vector bundle metric, the curvature condition also requires a suitable definition of positivity, discussed at the end of Chapter~\ref{l2:chapter}.


\chapter{Hodge Theory, Duality, and Smooth Representatives of Currents}\label{hod:chapter}
\begin{learning}
Dolbeault cohomology identifies two closed forms when their difference is exact. A Hermitian metric selects a distinguished representative of each class: the harmonic form, orthogonal to all exact forms. On a compact manifold, elliptic estimates make this representative smooth and the space of such representatives finite-dimensional. We construct the harmonic projection and Green operator, then use them to prove duality and the K\"ahler decomposition.
\end{learning}

\section{Constructing the self-adjoint operator from its energy form}
Let $X$ be a compact Hermitian complex manifold and $E$ a Hermitian holomorphic vector bundle, and fix a bidegree $(p,q)$. Here $L^2=L^2_{p,q}(X,E)$ and $H^s=H^s_{p,q}(X,E)$ are the metric square-integrable and Sobolev spaces of these forms, as in Chapter~\ref{ana:chapter}. Inner products are conjugate-linear in the first variable and use $dV_\omega=\omega^n/n!$. The formal adjoint $\dbar^*$ has the maximal realization described in Proposition~\ref{pos:compact-core}. Initially on smooth forms, write
\[
 Q(u,v)=\langle\dbar u,\dbar v\rangle+
                 \langle\dbar^*u,\dbar^*v\rangle,
 \qquad Q(u)=Q(u,u).
\]
The local anticommutation relation between exterior multiplication and contraction shows that
$D=\sqrt2(\dbar+\dbar^*)$, acting on the direct sum of all antiholomorphic degrees, has square with positive definite scalar principal symbol; its restriction in the fixed bidegree is $D^2=2\Delta''$.
\begin{lemma}[The energy domain and compact inverse]\label{hod:compact-resolvent}
The domain of the closed quadratic form $Q$ is $H^1$. For every $f\in L^2$, there is a unique $u\in H^1$ such that
\begin{equation}\label{hod:weak-resolvent}
 Q(v,u)+\langle v,u\rangle=\langle v,f\rangle
                      \quad(v\in H^1).
\end{equation}
The map $K:f\mapsto u$ is an injective, self-adjoint, positive compact operator on $L^2$, and $u\in H^2$.
\end{lemma}
\begin{proof}
For smooth $u$, the localized first-order energy estimate and a finite partition of unity give
\[
 C^{-1}\|u\|_{H^1}^2\le Q(u)+\|u\|_2^2\le C\|u\|_{H^1}^2.
\]
The upper bound uses only that $\dbar$ and $\dbar^*$ are first-order operators. The lower bound follows from Proposition \ref{ana:dirac-estimate}; orthogonality of distinct degrees gives
$\|Du\|^2=2Q(u)$. Smooth forms are dense in $H^1$, so completion in the energy norm gives exactly $H^1$, with an equivalent norm.

Set $a(v,u)=Q(v,u)+\langle v,u\rangle$. This is a Hilbert inner product on $H^1$, and $v\mapsto\langle v,f\rangle$ is a continuous conjugate-linear functional. The Riesz representation theorem gives a unique $u$. Here the representation theorem follows from the projection theorem: if the functional is nonzero, its kernel is a closed hyperplane; choose a unit vector in the orthogonal complement and use the value of the functional there to determine the representing vector. The projection theorem itself follows from the parallelogram identity applied to a sequence minimizing distance. Taking $v=u$ yields
$\|u\|_{H^1}\le C\|f\|_2$.
Thus $K:L^2\to H^1\hookrightarrow L^2$ is compact, since the last inclusion is compact by Theorem \ref{ana:rellich}.

For $u=Kf,w=Kg$, the equation gives
$\langle f,Kg\rangle=a(Kf,Kg)=\langle Kf,g\rangle$, so $K$ is self-adjoint. Moreover,
$\langle f,Kf\rangle=a(Kf,Kf)\ge0$.
If $Kf=0$, then \eqref{hod:weak-resolvent} makes $f$ orthogonal to every smooth form, hence $f=0$.
Finally, the weak equation is $\Delta''u=f-u$. Its right-hand side belongs to $L^2$; summing Theorem \ref{ana:interior-estimate} over finitely many charts gives $u\in H^2$.
\end{proof}

Define the self-adjoint realization of $\Delta''$ by $\Delta''=K^{-1}-I$ on $\operatorname{Ran}K$. Its domain is exactly $H^2$: Lemma~\ref{hod:compact-resolvent} gives $\operatorname{Ran}K\subset H^2$, while for $u\in H^2$ the weak integration-by-parts identity gives $K(\Delta''u+u)=u$. In particular, the differential expression and this operator realization agree on smooth forms. The spectral construction below also verifies self-adjointness directly.

\begin{lemma}[Spectral decomposition of a positive compact operator]\label{hod:compact-spectral}
If $K$ is an injective, positive, self-adjoint compact operator on a separable Hilbert space, then it has a complete orthonormal system of eigenvectors $e_j$, with eigenvalues $\mu_j>0$. In infinite dimension, $\mu_j\to0$.
\end{lemma}
\begin{proof}
For a positive operator, $\|K\|=\sup_{\|v\|=1}\langle v,Kv\rangle$. This identity follows from polarization and the Cauchy--Schwarz inequality
$|\langle x,Ky\rangle|^2\le\langle x,Kx\rangle\langle y,Ky\rangle$ for the positive form; the latter is the discriminant condition for the nonnegative quadratic expression $\langle x+ty,K(x+ty)\rangle\ge0$.
Choose unit vectors $v_j$ such that the quadratic form tends to $\mu=\|K\|$. Compactness gives a convergent subsequence of $Kv_j$, and
\[
 \|(K-\mu)v_j\|^2
 =\|Kv_j\|^2-2\mu\langle v_j,Kv_j\rangle+\mu^2\longrightarrow0,
\]
since $\|Kv_j\|^2\le\mu^2$ and the left-hand side is nonnegative. If $\mu>0$, then $v_j=\mu^{-1}(Kv_j-(K-\mu)v_j)$ also converges, and its limit is a unit eigenvector.

Its orthogonal complement is invariant under $K$, so repeat the argument. Each nonzero eigenspace is finite-dimensional: otherwise an orthonormal sequence in it would have images under $K$ with no convergent subsequence. If the resulting eigenvalues failed to tend to zero, the same compactness argument would give a contradiction. Finally, if $K$ were nonzero on the orthogonal complement of all eigenvectors, the first paragraph would produce another eigenvector there. Hence $K$ vanishes on that complement. Injectivity makes the complement zero, proving completeness.
\end{proof}

\begin{theorem}[Hodge decomposition on a compact Hermitian manifold]\label{hod:hodge}
The operator $\Delta''=\dbar\dbar^*+\dbar^*\dbar$ has an orthonormal basis of smooth eigenforms, with eigenvalues
$0\le\lambda_1\le\lambda_2\le\cdots\to\infty$, each of finite multiplicity.
Write $\mathcal H^{p,q}(X,E)=\ker\Delta''$, and let $H$ be the orthogonal projection onto it. There is a Green operator $G$ satisfying
\begin{equation}\label{hod:green-equation}
 \Delta''G=G\Delta''=I-H,\qquad HG=GH=0.
\end{equation}
For every smooth form $u$, there is an orthogonal decomposition
\begin{equation}\label{hod:decomposition}
 u=Hu+\dbar\dbar^*Gu+\dbar^*\dbar Gu.
\end{equation}
In particular, Dolbeault cohomology is naturally isomorphic to $\mathcal H^{p,q}(X,E)$ and is therefore finite-dimensional.
\end{theorem}
\begin{proof}
Apply Lemma~\ref{hod:compact-spectral} to $K$ from Lemma~\ref{hod:compact-resolvent}. The identity $Ke_j=\mu_je_j$ and the weak equation imply
$\Delta''e_j=(\mu_j^{-1}-1)e_j$. Nonnegativity gives $\lambda_j=\mu_j^{-1}-1\ge0$. Theorem~\ref{ana:elliptic-regularity} first gives $H^2$ and then, by iteration, every $H^m$; Sobolev embedding (Lemma~\ref{ana:sobolev-embedding}) makes the eigenvectors smooth. Since $\mu_j\to0$, we have $\lambda_j\to\infty$.
The zero eigenspace is finite-dimensional; otherwise $K$ would have an infinite-dimensional eigenspace with eigenvalue $1$, contradicting compactness.

Writing $u=\sum_j u_je_j$, define
\[
 Hu=\sum_{\lambda_j=0}u_je_j,
 \qquad Gu=\sum_{\lambda_j>0}\lambda_j^{-1}u_je_j.
\]
The positive eigenvalues are bounded away from zero, so $G$ is bounded on $L^2$. Verify \eqref{hod:green-equation} first on finite sums, then pass to limits using $L^2$ convergence and closedness of the operator. This gives the identities distributionally, with $G\Delta''$ interpreted on the domain of $\Delta''$.
For smooth $u$, the right-hand side is smooth, and regularity makes $Gu$ smooth. Expanding the Green identity yields \eqref{hod:decomposition}.
The two image terms are orthogonal because
$\langle\dbar a,\dbar^*b\rangle=\langle\dbar^2a,b\rangle=0$; both are orthogonal to harmonic forms, since a harmonic form satisfies $\dbar h=\dbar^*h=0$.

On smooth forms, $\Delta''\dbar=\dbar\Delta''$, simply by expansion and $\dbar^2=0$. Thus $\dbar$ preserves each positive eigenspace and commutes with $G$; elliptic estimates, or smooth convergence of the spectral series, extend this to every smooth form. If $\dbar u=0$, then $\dbar Gu=G\dbar u=0$, so
$u=Hu+\dbar(\dbar^*Gu)$. Every class therefore has a harmonic representative.
If a harmonic form satisfies $h=\dbar v$, then
$\|h\|^2=\langle h,\dbar v\rangle=\langle\dbar^*h,v\rangle=0$, proving uniqueness.
\end{proof}

\begin{proposition}[The Green operator on currents]\label{hod:green}
The operator $G$ maps smooth forms continuously to smooth forms and extends by transposition to all currents; \eqref{hod:green-equation} and \eqref{hod:decomposition} remain valid for currents.
If a current $T$ satisfies $\Delta''T\in C^\infty$, then $T$ is smooth.
\end{proposition}
\begin{proof}
Theorem~\ref{ana:interior-estimate}, summed over a finite cover, and $L^2$ boundedness successively give
\[
 \|Gu\|_{H^{s+2}}\le C_s\bigl(\|(I-H)u\|_{H^s}+\|Gu\|_2\bigr)
                    \le C'_s\|u\|_{H^s}.
\]
Since $H$ is projection onto finitely many smooth basis vectors, it is continuous on every Sobolev space. Together with Lemma~\ref{ana:sobolev-embedding}, the estimate proves continuity of $G:C^\infty\to C^\infty$.
Use the metric and volume form to identify a current with a generalized section acting conjugate-linearly on smooth test sections of the same bidegree and bundle. In this pairing a smooth section $u$ is represented by $T_u(v)=\langle v,u\rangle_{L^2}$. Define $(GT)(v)=T(Gv)$; then $T_u(Gv)=\langle Gv,u\rangle=\langle v,Gu\rangle$, so this extension agrees exactly with the original operator. For a general differential operator $P$, the corresponding distributional action is $(PT)(v)=T(P^*v)$. These formulas specify the conjugations and the test spaces before transposing the identities.
Define $HT$ in the same way; it is a finite linear combination of smooth harmonic forms.
Apply the established operator identities to smooth test forms and transpose to obtain the Green identities on currents.
If $\Delta''T=f$ is smooth, then $T=HT+Gf$, whose two terms are smooth.
The transposition argument applies to arbitrary currents; no $L^2$ hypothesis on $T$ is needed.
\end{proof}

\section{Serre duality and the K\"ahler decomposition}
\begin{theorem}[The analytic form of Serre duality]\label{hod:serre-duality}
Write $H^{p,q}(X,E)=H^{p,q}_{\dbar}(X,E)$. In a product $u\wedge v$, the symbol $\tr$ denotes contraction $E\otimes E^*\to\C$, followed by the wedge product of the form factors. Integration defines a nondegenerate pairing
\[
 H^{p,q}(X,E)\times H^{n-p,n-q}(X,E^*)\longrightarrow\C,
 \qquad ([u],[v])\longmapsto\int_X\tr(u\wedge v).
\]
In particular, $H^q(X,E)$ and $H^{n-q}(X,K_X\otimes E^*)$ are dual spaces.
\end{theorem}
\begin{proof}
First check independence of representatives. Replacing $u$ by $u+\dbar a$ when $v$ is closed changes the integrand, by the Leibniz rule, by a $\dbar$-exact form; its integral vanishes by Proposition~\ref{bas:stokes}. The other variable is treated in the same way.
Pointwise, the metric defines an antilinear isomorphism
$S:\Lambda^{p,q}T^*X\otimes E\to\Lambda^{n-p,n-q}T^*X\otimes E^*$, characterized by
$\tr(u\wedge Su)=|u|^2dV$. In a unitary coframe, $S$ sends each exterior basis element to the complementary exterior product with the sign that makes this identity hold. Thus $S$ is invertible and isometric.
Integration by parts for compactly supported forms gives
$S\dbar=\pm\dbar^*S$ and $S\dbar^*=\pm\dbar S$, where the operators on the right act on $E^*$ with its dual metric and the signs depend only on total degree. To see this, apply the Leibniz rule and Stokes' theorem to
$\int\tr(\dbar a\wedge b)$; this yields the asserted adjoint relation between $\dbar$ and $\dbar^*$ in complementary degrees. The signs need not be evaluated, since they do not affect kernels.
Consequently $S$ gives a bijection between harmonic forms and harmonic forms of complementary degree with values in the dual bundle.
For $[u]\ne0$, choose its nonzero harmonic representative and set $v=Su$. The integral is $\|u\|^2>0$, so the pairing of the two finite-dimensional spaces is nondegenerate.
Finally, setting $(p,q)=(0,q)$, using $\Omega_X^n=K_X$, and applying Theorem~\ref{sh:dolbeault} gives the sheaf-cohomology formulation.
\end{proof}

\begin{proposition}[The three Laplacians in the K\"ahler case]\label{hod:kahler-laplacians}
For ordinary complex-valued forms on a compact K\"ahler manifold, put $\Delta_d=dd^*+d^*d$ and $\Delta_\partial=\partial\partial^*+\partial^*\partial$. Then
\[
 \Delta_d=2\Delta_\partial=2\Delta'',
 \qquad
 H^k_{\mathrm{dR}}(X,\C)=\bigoplus_{p+q=k}H^{p,q}(X).
\]
Complex conjugation gives $h^{p,q}=h^{q,p}$; in particular, $b_1=2h^{0,1}$.
\end{proposition}
\begin{proof}
Apply the K\"ahler commutation relations (Proposition \ref{pos:kahler-identities}) to the trivial line bundle. Its curvature is zero, so Theorem~\ref{pos:bkn} gives $\Delta_\partial=\Delta''$. The mixed terms also vanish. For example, substitute $\dbar^*=-\ii[\Lambda,\partial]$ into
$\partial\dbar^*+\dbar^*\partial$ and expand; the terms cancel in pairs because $\partial^2=0$. Take adjoints for the other mixed term. Expanding the Laplacian of $d=\partial+\dbar$ proves the first identity.
Repeat Lemma~\ref{hod:compact-resolvent} and Theorem~\ref{hod:hodge} for $d+d^*$ to obtain a unique $d$-harmonic representative of every de Rham class. Its squared principal symbol is again positive definite and scalar, so all analytic steps are identical.
The first identity says that a form is $d$-harmonic precisely when each bidegree component is $\dbar$-harmonic, proving the direct-sum decomposition.
Complex conjugation exchanges $\partial$ and $\dbar$ and commutes with $\Delta_d$, hence gives an antilinear bijection between the harmonic spaces of bidegrees $(p,q)$ and $(q,p)$. Taking $k=1$ proves the last identity.
\end{proof}

\begin{theorem}[The $\partial\bar\partial$ lemma, including currents]\label{hod:ddbar-lemma}
Let $X$ be compact K\"ahler, and let $T$ be a $d$-closed current of pure bidegree $(p,q)$. The following are equivalent: $T$ is $d$-exact, $\partial$-exact, $\dbar$-exact, or $\partial\dbar$-exact.
If $T$ is smooth, the potential in the last assertion may also be chosen smooth. If $p=0$ or $q=0$, any such exact current is zero.
\end{theorem}
\begin{proof}
The last condition implies the other three by the type identities of Proposition~\ref{bas:type-decomposition}, in particular $\partial\dbar=-\dbar\partial$.
Conversely, any of the first three conditions implies $HT=0$: pair an exact expression with a harmonic test form and integrate by parts. The result is zero because K\"ahler harmonic forms are annihilated by $d^*,\partial^*,\dbar^*$.
Since $T$ has pure bidegree and $dT=0$, we also have $\partial T=\dbar T=0$. The Green operator commutes with $\partial,\dbar$ (first on smooth eigenspaces, then by transposition on currents), so
$\partial GT=\dbar GT=0$.
Now calculate in succession:
\[
 T=\Delta''GT=\dbar\dbar^*GT,
 \qquad
 \dbar^*GT=-\ii[\Lambda,\partial]GT=\ii\partial\Lambda GT.
\]
Therefore
\[
 T=-\ii\partial\dbar\Lambda GT.
\]
The potential can be taken to be $-\ii\Lambda GT$; when $T$ is smooth, Proposition \ref{hod:green} makes it smooth.
If $p=0$ or $q=0$, one component of the bidegree of $\Lambda GT$ is negative, so it vanishes and the formula gives $T=0$.
For currents, the compositions with $G$ and the commutation identities are understood through the transposition construction in Proposition~\ref{hod:green}.
\end{proof}

\section{Smooth comparison on noncompact manifolds}
The comparison of smooth forms and currents also holds without compactness or the K\"ahler condition. Here the local Dolbeault lemma replaces the global Green operator. A finite double-complex argument will compare the two kinds of representatives.

\begin{lemma}[A bounded double complex with exact rows and columns]\label{hod:double-exact}
Let $Q^{p,q}$ be a double complex nonzero in only finitely many bidegrees. Write its differentials as $a:Q^{p,q}\to Q^{p+1,q}$ and $b:Q^{p,q}\to Q^{p,q+1}$, with $a^2=b^2=ab+ba=0$. If every row and column is exact, then
\[
 \ker(ab)=\operatorname{im}a+\operatorname{im}b,
 \qquad \ker a\cap\ker b=\operatorname{im}(ab).
\]
\end{lemma}
\begin{proof}
Suppose first that $abx=0$. Row exactness gives $bx=au_1$. Applying $b$ yields $abu_1=0$, so $bu_1=au_2$. Continue in this way, decreasing $p$ by one and increasing $q$ by one at every step. Boundedness eventually gives $bu_m=0$. Column exactness gives $u_m=bv_m$. Substituting in the previous identity yields
$b(u_{m-1}+av_m)=0$, whence $u_{m-1}+av_m=bv_{m-1}$.
Back-substitute until $u_1=bv_1-av_2$. Then $bx=abv_1=-bav_1$, so $b(x+av_1)=0$. Column exactness once more gives $x=-av_1+bw$, proving the first identity.
If $ax=bx=0$, row exactness first gives $x=au$. Thus $abu=0$, and the first identity gives $u=av+bw$, hence $x=abw$. The reverse inclusion follows immediately from anticommutation.
\end{proof}

\begin{lemma}[Dolbeault comparison for smooth forms and currents]\label{hod:current-dolbeault}
On every complex manifold, the inclusion of smooth forms into currents induces an isomorphism on global cohomology in each $\dbar$ complex and each $\partial$ complex.
\end{lemma}
\begin{proof}
First verify local exactness. The scalar distribution kernel $K(z)=1/(\pi z)$ satisfies $\partial_{\bar z}K=\delta_0$, or equivalently $\dbar K=\delta_0\,d\bar z$. This is the distributional kernel identity behind Lemma~\ref{sh:cauchy-operator}: apply Stokes' theorem outside the disk of radius $\varepsilon$ and shrink its boundary circle to zero; the boundary integral is the test-function value at zero.
For a distribution compactly supported in one coordinate direction, convolution by $K$ therefore gives a right inverse for $\dbar$ in that direction. For a general local distribution, first multiply by a cutoff equal to one on a smaller polydisk. The extra terms are supported outside the smaller coordinate disk, so their convolutions are smooth in the variable under consideration. Eliminate the components containing $d\bar z_j$ one coordinate at a time, by the same finite coordinatewise elimination as in Theorem~\ref{sh:dolbeault-lemma}. Convolution and differentiation remain valid distributionally, yielding local primitives for closed currents of positive degree.
A scalar $\dbar$-closed distribution, hence a current of antiholomorphic degree zero, is a holomorphic function: its smooth convolutions are holomorphic; on a smaller polydisk, the Cauchy integral formula expresses each derivative as the pairing of the smoothed function with a fixed smooth test kernel. Distributional convergence makes these pairings locally uniformly bounded and convergent, so the original distribution is represented by a holomorphic function. With a fixed $(p,0)$ basis, apply this componentwise.

Smooth forms and currents thus provide two exact resolutions of the sheaf of holomorphic $p$-forms. Both sheaves admit multiplication by smooth partitions of unity and are therefore fine by Proposition~\ref{sh:smooth-fine}. The resolution-comparison theorem, Theorem~\ref{sh:resolution-theorem}, shows that the global complexes of both fine resolutions compute the same sheaf cohomology, and that inclusion induces the identity. This proves the $\dbar$ comparison; complex conjugation gives the $\partial$ comparison.
\end{proof}

Write $\mathcal D'^{p,q}(X)$ for the space of $(p,q)$ currents and $\mathcal A^{p,q}(X)$ for smooth $(p,q)$ forms. For either of these double complexes $V^{\bullet,\bullet}$, define
\[
 \begin{aligned}
 H_{\rm BC}^{p,q}(V)&=\frac{\ker\partial\cap\ker\dbar\ \text{ in }V^{p,q}}{\partial\dbar V^{p-1,q-1}},\\
 H_{\rm A}^{p,q}(V)&=\frac{\ker(\partial\dbar:V^{p,q}\to V^{p+1,q+1})}{\partial V^{p-1,q}+\dbar V^{p,q-1}}.
 \end{aligned}
\]
These are Bott--Chern and Aeppli cohomology; spaces in negative degrees are zero.

\begin{theorem}[Smooth representatives in Bott--Chern and Aeppli cohomology]\label{hod:bc-aeppli}
The following hold on every complex manifold $X$.
If $T$ is a current with smooth $\partial\dbar T$, then
\[
 T=\omega+\partial\alpha+\dbar\beta,
 \qquad \omega\text{ smooth}.
\]
If a smooth form has an expression $T=\partial u+\dbar v$ with currents $u,v$, it also has such an expression with smooth $u,v$.
The inclusion of smooth forms into currents induces isomorphisms $H_{\rm BC}^{p,q}(\mathcal A)\simeq H_{\rm BC}^{p,q}(\mathcal D')$ and $H_{\rm A}^{p,q}(\mathcal A)\simeq H_{\rm A}^{p,q}(\mathcal D')$.
If $X$ is compact K\"ahler and $\partial\dbar T=0$, the form $\omega$ in the first expression can be chosen harmonic.
\end{theorem}
\begin{proof}
Set $Q^{p,q}=\mathcal D'^{p,q}(X)/\mathcal A^{p,q}(X)$, the quotient of currents by smooth forms. Lemma \ref{hod:current-dolbeault} and the long exact sequence of a short exact sequence of complexes show that every row and column of $Q$ is exact.
If $\partial\dbar T$ is smooth, its image $[T]\in Q$ satisfies $\partial\dbar[T]=0$. Lemma \ref{hod:double-exact} gives $[T]=\partial[\alpha]+\dbar[\beta]$. Choose any current representatives and move terms to the other side; the difference is precisely a smooth form $\omega$, proving the first assertion.

If $T=\partial u+\dbar v$ is smooth, then $\partial[u]+\dbar[v]=0$ in $Q$, hence $\partial\dbar[v]=0$. The same lemma gives $[v]=\partial[a]+\dbar[b]$. Substitution yields
$\partial([u]-\dbar[a])=0$, and row exactness gives $[u]-\dbar[a]=\partial[c]$.
Choose current representatives and put $u'=u-\dbar a-\partial c$, $v'=v-\partial a-\dbar b$. Both are smooth, and
\[
 \partial u'+\dbar v'
 =T-\partial\dbar a-\dbar\partial a=T.
\]
This proves injectivity for Aeppli comparison; the first paragraph proves surjectivity.

For Bott--Chern surjectivity, suppose $T$ is both $\partial$- and $\dbar$-closed. The second identity of the lemma gives $[T]=\partial\dbar[S]$, so $T-\partial\dbar S$ is a smooth closed representative.
For injectivity, suppose the smooth form $T$ is $\partial\dbar S$. In $Q$, $\partial\dbar[S]=0$, and the first identity gives a smooth form $S-\partial a-\dbar b=\gamma$. Thus $T=\partial\dbar\gamma$.

Finally, suppose $X$ is compact K\"ahler. Apply Theorem~\ref{hod:hodge} to the smooth Aeppli representative $\omega$ obtained above. Alternatively, apply directly to currents the identity
$T=HT+\Delta''GT$. The term $\dbar\dbar^*GT$ is already $\dbar$-exact. The remaining term $\dbar^*\dbar GT$ is $\partial$-exact by
$\dbar^*=-\ii[\Lambda,\partial]$ and $\partial\dbar GT=G\partial\dbar T=0$.
Explicitly, $\dbar^*\dbar GT=\ii\partial\Lambda\dbar GT$.
Hence $T=HT+\partial(\ii\Lambda\dbar GT)+\dbar(\dbar^*GT)$. If $T$ is smooth, both potentials are smooth.
In the first assertion, the smooth representative is obtained modulo $\partial$-exact and $\dbar$-exact currents. This equivalence relation is the one used in Aeppli cohomology, rather than equivalence modulo $d$-exact currents.
\end{proof}


\chapter{From Positive Curvature to Projective Embeddings}\label{emb:chapter}
\begin{learning}
Holomorphic functions on a compact connected complex manifold are constant by Theorem~\ref{bas:compact-holomorphic}, so projective coordinates must come from sections of a line bundle. Positive curvature lets us construct these sections: cut off local coordinate functions, solve away their $\dbar$ error, and use the resulting section ratios as local coordinates. Compactness then gives one tensor power whose sections separate all points and tangent vectors.
\end{learning}

Throughout this chapter, $d\lambda$ denotes Euclidean Lebesgue measure in a coordinate chart; global norms use the stated bundle metric and a fixed smooth positive volume form. For a multi-index $\alpha$, $\partial^\alpha=\partial_{z_1}^{\alpha_1}\cdots\partial_{z_n}^{\alpha_n}$. A scalar radius in $\Delta(a;r)$ means that all polydisk radii equal $r$, in the notation of Definition~\ref{bas:holomorphic-def}.

\section{Local norm control for holomorphic functions}
\begin{lemma}[Interior Cauchy--$L^2$ estimate]\label{emb:cauchy-l2}
Let $K\Subset U\subset\C^n$, with $K$ compact. For every multi-index $\alpha$, there is a constant $C$ such that every holomorphic function on $U$ satisfies
\[
 \sup_K|\partial^\alpha f|\le C\|f\|_{L^2(U)}.
\]
\end{lemma}
\begin{proof}
First take a polydisk centered at $a$, with all radii equal to $r$, and write
$f=\sum_\beta c_\beta(z-a)^\beta$.
Integrate the uniformly convergent series on smaller polydisks and then let their radii increase to $r$.
Distinct monomials are orthogonal in each angular variable, since
$\int_0^{2\pi}e^{\ii(m-l)\theta}\,d\theta=0$ for $m\ne l$.
Thus
\[
 \int_{\Delta(a;r)}|f|^2
 =\sum_\beta|c_\beta|^2\prod_{j=1}^n
             \frac{\pi r^{2\beta_j+2}}{\beta_j+1}
 \ge |c_0|^2\pi^nr^{2n}.
\]
All terms are nonnegative, justifying passage to the limit, and $c_0=f(a)$.
This proves the mean-square pointwise estimate.

Since $K\Subset U$, fix $r>0$ so that the doubled polydisk
$\overline{\Delta(a;2r)}$ lies in $U$ for every $a\in K$.
If $z\in\overline{\Delta(a;r)}$, its polydisk of radius $r$ still lies in
$\Delta(a;2r)$, whence
$|f(z)|\le\pi^{-n/2}r^{-n}\|f\|_{L^2(U)}$.
Apply the derivative estimate of Theorem~\ref{bas:cauchy} at $a$ to obtain
\[
 |\partial^\alpha f(a)|
 \le\alpha!\,\pi^{-n/2}r^{-n-|\alpha|}\|f\|_{L^2(U)}.
\]
The constant is independent of $a$ and $f$; taking the supremum over $a\in K$ proves the assertion.
If the right-hand side is infinite, the inequality still holds, so boundedness of $U$ is unnecessary.
\end{proof}

\begin{proposition}[Finite-dimensional section spaces on compact manifolds]\label{emb:finite-sections}
If $X$ is compact and $E$ is a holomorphic vector bundle of finite rank, then $H^0(X,E)$ is finite-dimensional.
\end{proposition}
\begin{proof}
Choose finitely many holomorphic trivializations $U_i$ and relatively compact subsets $V_i\Subset U_i$ covering $X$.
Fix a smooth Hermitian metric and volume form.
On a slightly larger compact neighborhood of $\overline{V_i}$, the metric matrix has its smallest eigenvalue bounded positively below and its largest eigenvalue bounded above; the volume density has corresponding positive lower and finite upper bounds.
Thus sections with global $L^2$ norm at most one have uniformly bounded Euclidean $L^2$ norms for their local coefficients.
Lemma~\ref{emb:cauchy-l2} gives uniform bounds for these coefficients and their first derivatives, hence uniform boundedness and equicontinuity on $\overline{V_i}$.

To extract a uniformly convergent subsequence, choose a countable dense subset of each $\overline{V_i}$.
At each point, the finitely many coefficient values lie in bounded subsets of $\C$.
Repeated subsequence extraction and a diagonal choice make all these values converge simultaneously.
Given $\epsilon>0$, equicontinuity gives a common $\delta$.
Compactness provides a finite covering of $\overline{V_i}$ by $\delta$-balls centered at dense-set points.
For sufficiently large subsequence indices, the differences at these finitely many centers are less than $\epsilon$.
Applying equicontinuity twice bounds the difference on the whole compact set by $3\epsilon$.
The subsequence therefore converges uniformly on each $\overline{V_i}$.
Since there are only finitely many charts, this can be arranged simultaneously.

Corollary~\ref{bas:identity} makes the limiting coefficients holomorphic.
The transition relations $s_i=g_{ij}s_j$ are continuous matrix identities and survive the limit, so the limits patch to a global holomorphic section.
Uniform convergence on finitely many charts, finite volume, and metric comparison also give global $L^2$ convergence.
Every sequence in the unit ball therefore has an $L^2$-convergent subsequence.
If the section space were infinite-dimensional, choose successively sections outside the span of their predecessors, subtract their orthogonal projections, and normalize. This produces infinitely many orthonormal sections.
Any two have distance $\sqrt2$, contradicting the subsequence compactness just proved.
\end{proof}

\section{Normalizing the local weight}
\begin{lemma}[The local quadratic model of a positive line bundle]\label{emb:normal-weight}
Suppose $L$ has a smooth metric $h$ with strictly positive curvature, and let $p\in X$.
There are complex coordinates $z$ centered at $p$ and a local holomorphic frame $e$ such that $|e|_h^2=e^{-\varphi}$ and
\[
 \varphi(z)=\sum_{j,k}a_{j\bar k}z_j\bar z_k+O(|z|^3),\qquad
 c|z|^2\le\varphi(z)\le C|z|^2
\]
on a small ball, where $(a_{j\bar k})$ is positive definite and $c,C>0$.
\end{lemma}
\begin{proof}
In arbitrary centered coordinates, group the second-order Taylor expansion of the real smooth weight by type:
\[
 \varphi(z)=c+2\operatorname{Re}\!\left(\sum_jb_jz_j+
              \frac12\sum_{j,k}c_{jk}z_jz_k\right)
       +\sum_{j,k}a_{j\bar k}z_j\bar z_k+R(z),
 \qquad |R(z)|\le C_0|z|^3.
\]
Reality of the weight makes the mixed matrix Hermitian; its entries are
$a_{j\bar k}=\partial_j\partial_{\bar k}\varphi(0)$.
Define the holomorphic polynomial
$g(z)=c/2+\sum_jb_jz_j+\tfrac12\sum_{j,k}c_{jk}z_jz_k$
and replace the frame by $e'=e\exp g$.
The exponential never vanishes, so this remains a holomorphic frame; its norm is
$|e'|^2=e^{-\varphi+2\operatorname{Re}g}$.
The new weight therefore has no constant, linear, or pure quadratic terms, leaving only the mixed quadratic part and the remainder.

The Levi matrix of the curvature is $(a_{j\bar k})$. Strict positivity gives a smallest eigenvalue
$\lambda>0$ and a largest eigenvalue $\lambda_{\max}$.
Choose $r$ with $C_0r\le\lambda/2$.
For $|z|<r$,
\[
 \frac\lambda2|z|^2
 \le\lambda|z|^2-C_0|z|^3
 \le\varphi(z)
 \le\lambda_{\max}|z|^2+C_0|z|^3
 \le(\lambda_{\max}+\lambda/2)|z|^2.
\]
Take $c=\lambda/2$ and $C=\lambda_{\max}+\lambda/2$.
In particular, the new frame has norm one at the center, a fact used below when comparing point values of sections of high powers.
\end{proof}

\section{Peak sections: cut off and solve}
From now on, $X$ is compact and $L$ is positive. The form $\omega=\ii F_L$ is itself K\"ahler.
Corollary \ref{pos:high-power-solving} therefore applies to $L^k$: there are $k_0,C_0$ such that for every $k\ge k_0$ and every smooth closed $L^k$-valued $(0,1)$-form $f$, there is a smooth solution
\begin{equation}\label{emb:solver}
 \dbar u=f,\qquad \|u\|_2^2\le\frac{C_0}{k}\|f\|_2^2.
\end{equation}
All norms use $h^k$ and a fixed volume form; in particular, the exponential estimates below concern weighted sections, not unweighted coefficients.

\begin{proposition}[Approximating local coordinates by global sections]\label{emb:peak}
Given $p\in X$, there are a coordinate ball $B$, a holomorphic frame $e$, and constants $C,\gamma>0$ such that, for every sufficiently large integer $k$, global holomorphic sections
$s_{0,k},s_{1,k},\ldots,s_{n,k}\in H^0(X,L^k)$ have on $B$ the expressions
\[
 s_{0,k}=e^{\otimes k}(1-v_{0,k}),\qquad
 s_{j,k}=e^{\otimes k}(z_j-v_{j,k}),
\]
with $\|v_{j,k}\|_{C^1(B)}\le Ce^{-\gamma k}$.
\end{proposition}
\begin{proof}
Choose the normalized frame of Lemma~\ref{emb:normal-weight} on a ball $B_r=B(0,r)$, and a fixed cutoff
$\chi\in C_c^\infty(B_r)$ equal to $1$ near $\overline{B_{r/2}}$.
Set $a_0=1$, $a_j=z_j$, and extend
$\sigma_{j,k}=\chi a_je^{\otimes k}$ by zero to $X$.
Since the support of $\chi$ is compactly contained in the frame domain, this extension vanishes near its boundary and is smooth.
The functions $a_j$ and the frame are holomorphic, so
\[
 f_{j,k}=\dbar\sigma_{j,k}=(\dbar\chi)a_je^{\otimes k},
 \qquad\dbar f_{j,k}=\dbar^2\sigma_{j,k}=0.
\]
The error is supported in a fixed compact annulus where $|z|\ge r/2$,
and hence $\varphi\ge cr^2/4=:\mu>0$.
On this annulus, $|\dbar\chi|$, $|a_j|$, and the volume density have bounds independent of $k$;
the squared norm of the $k$th power of the frame is exactly $e^{-k\varphi}$.
Therefore
\begin{equation}\label{emb:error-size}
 \|f_{j,k}\|_2^2\le C_1e^{-\mu k}.
\end{equation}

Solve with \eqref{emb:solver} to obtain smooth $u_{j,k}$ satisfying
$\dbar u_{j,k}=f_{j,k}$ and
$\|u_{j,k}\|_2^2\le C_2k^{-1}e^{-\mu k}$.
Then $s_{j,k}=\sigma_{j,k}-u_{j,k}$ is globally holomorphic.
On $B_{r/2}$, where the error vanishes, write $u_{j,k}=v_{j,k}e^{\otimes k}$.
Holomorphicity of the frame makes $\dbar u_{j,k}=0$ equivalent to $\dbar v_{j,k}=0$,
so the holomorphic-function estimates apply to $v_{j,k}$.

Choose a fixed $\rho<r/2$ with $\sup_{B_\rho}\varphi\le\mu/2$.
The fixed volume density has a positive lower bound, and
$|u_{j,k}|_{h^k}^2=|v_{j,k}|^2e^{-k\varphi}$, so
\[
 \int_{B_\rho}|v_{j,k}|^2d\lambda
 \le C_3 e^{k\mu/2}\|u_{j,k}\|_2^2
 \le C_4k^{-1}e^{-\mu k/2}.
\]
Applying Lemma \ref{emb:cauchy-l2} to the zeroth and first derivatives gives
$\|v_{j,k}\|_{C^1(B_{\rho/2})}
 \le C_5k^{-1/2}e^{-\mu k/4}\le C_5e^{-\mu k/4}$.
There are finitely many $j$, so choose common constants and set $B=B_{\rho/2}$, $\gamma=\mu/4$.
The correction $u_{j,k}$ is holomorphic only off the error support; it is the difference $s_{j,k}$ that is globally holomorphic.
\end{proof}

For a globally generated power $L^k$, choose a basis $s_0,\ldots,s_N$ of the finite-dimensional space $H^0(X,L^k)$ from Proposition~\ref{emb:finite-sections}. In a local frame write $s_j=f_je^{\otimes k}$ and define the \emph{complete linear-system map}
\[
 \Phi_k:X\longrightarrow\mathbb P(H^0(X,L^k)^*),\qquad x\longmapsto[f_0(x):\cdots:f_N(x)].
\]
The common frame-change factor cancels in projective coordinates, so this map is well defined. A different basis changes only the projective coordinates. A power is \emph{base-point free} when its global sections have no common zero.

\begin{corollary}[Base-point freeness and local separation]\label{emb:local-separation}
For all sufficiently large $k$, the bundle $L^k$ is generated by global sections.
For each $p$, there are also a neighborhood $U$ and such an integer $k$ for which the complete linear-system map of $L^k$ is injective and immersive on $U$.
\end{corollary}
\begin{proof}
On the ball used for the peak-section construction, the coefficient of $s_{0,k}$ is $1-v_{0,k}$.
For $k$ sufficiently large that $\sup|v_{0,k}|<1/2$, this coefficient is nowhere zero.
Choose such a fixed ball and threshold at every center; compactness gives a finite subcover.
Taking the largest of the finitely many thresholds ensures that, for every larger $k$, each point admits a nonzero global section of $L^k$.
This is base-point freeness, equivalently global generation for a line bundle.

On a fixed ball, use $s_{0,k}$ as the projective-coordinate denominator:
\[
 F_k(z)_j=\frac{z_j-v_{j,k}}{1-v_{0,k}}
 =z_j+\frac{z_jv_{0,k}-v_{j,k}}{1-v_{0,k}}.
\]
The denominator has modulus at least $1/2$. Differentiating the quotient, every error term contains some $v_{j,k}$ or a first derivative of it, while $z$ stays bounded on the fixed ball. Thus $F_k(z)=z+E_k(z)$ with $\|E_k\|_{C^1}\to0$.
Shrink to a convex ball and increase $k$ until the real derivative satisfies $\|DE_k\|\le1/2$.
Integration along the segment from $w$ to $z$ gives
$E_k(z)-E_k(w)=\int_0^1DE_k(w+t(z-w))(z-w)\,dt$,
hence $|F_k(z)-F_k(w)|\ge\tfrac12|z-w|$.
Likewise $|dF_k(v)|\ge\tfrac12|v|$, so $F_k$ is injective and immersive.

By Proposition \ref{emb:finite-sections}, the section space has a finite basis.
The displayed $n+1$ sections are linear combinations of that basis, so their ratios can be recovered from the projective coordinates of the complete linear system.
If the complete system identified two points, all these ratios would agree;
if its differential annihilated a tangent vector, all their differentials would vanish on it.
Both contradict the properties of $F_k$, so the complete system has the same local separation properties.
\end{proof}

\begin{lemma}[A pointwise estimate away from the error support]\label{emb:point-estimate}
Let $u_k$ be an $L^2$ section of $L^k$ on $X$, holomorphic on a fixed neighborhood of $q$. For sufficiently large $k$,
$|u_k(q)|_{h^k}^2\le Ck^n\|u_k\|_{L^2(X)}^2$, with $C$ allowed to depend on the neighborhood and on $q$.
\end{lemma}
\begin{proof}
In coordinates near $q$, remove the constant and real linear terms of the weight by a holomorphic change of frame.
Strict positivity is unnecessary here: Taylor's theorem alone gives
$\varphi(0)=0$ and $|\varphi(z)|\le C_1|z|^2$.
Write $u_k=v_ke^{\otimes k}$, with $v_k$ holomorphic on this neighborhood.
For fixed small $\varepsilon$ and sufficiently large $k$, the polydisk of radius
$\varepsilon k^{-1/2}$ lies in the neighborhood.
There $k|\varphi|\le C_1n\varepsilon^2$, so weighted and unweighted norms are comparable by fixed constants.

The mean-square estimate proved in Lemma~\ref{emb:cauchy-l2} gives
\[
 |v_k(0)|^2\le
 \frac{k^n}{\pi^n\varepsilon^{2n}}
 \int_{\Delta(0;\varepsilon k^{-1/2})}|v_k|^2\,d\lambda
 \le Ck^n\|u_k\|_{L^2(X)}^2.
\]
The last step also compares the fixed background volume with $d\lambda$; its constant is independent of $k$.
Since $\varphi(0)=0$, the frame has $|e(q)|=1$, and $|v_k(0)|=|u_k(q)|_{h^k}$.
The factor $k^n$ is the reciprocal volume scale of a ball of radius $k^{-1/2}$ in complex dimension $n$.
\end{proof}

\begin{proposition}[Separating any pair of points by a high power]\label{emb:pair-separation}
If $p\ne q$, then for all sufficiently large $k$, the complete linear-system map of $L^k$ sends $p$ and $q$ to distinct points.
\end{proposition}
\begin{proof}
Choose disjoint coordinate neighborhoods of $p$ and $q$, and apply Proposition~\ref{emb:peak} in normalized frames on each to obtain
$s_{p,k}=\sigma_{p,k}-u_{p,k}$ and $s_{q,k}=\sigma_{q,k}-u_{q,k}$.
At the respective centers the frames have norm one, and the section coefficients $1-v_{0,k}$ tend to one.
On a fixed neighborhood of $q$, $\sigma_{p,k}=0$ and $\dbar u_{p,k}=0$,
so $s_{p,k}=-u_{p,k}$ is holomorphic there.
Lemma~\ref{emb:point-estimate} and the correction bound following \eqref{emb:error-size} give
\[
 |s_{p,k}(q)|_{h^k}^2\le Ck^n\|u_{p,k}\|_2^2
 \le C'k^{n-1}e^{-\mu k}\longrightarrow0.
\]
Similarly, $s_{q,k}(p)\to0$.

Identify the two fibers with $\C$ using the unit frames at the centers. The evaluation matrix satisfies
\[
 M_k=\begin{pmatrix}s_{p,k}(p)&s_{q,k}(p)\\
                         s_{p,k}(q)&s_{q,k}(q)\end{pmatrix}
 \longrightarrow\begin{pmatrix}1&0\\0&1\end{pmatrix}.
\]
Thus $\det M_k\ne0$ for large $k$, and evaluation onto $L_p^k\oplus L_q^k$ is surjective.
For example, subtracting a suitable multiple of $s_{q,k}$ from $s_{p,k}$ gives a section vanishing at $q$ but not at $p$; exchanging the points gives the converse.
Increase $k$ beyond the global base-point-free threshold, so that the complete linear system is defined everywhere.
If the projective images of the two points coincided, the two evaluations of every section would have a common proportionality factor.
The sections just constructed contradict this, proving separation.
\end{proof}

\begin{lemma}[Combining finitely many powers]\label{emb:veronese}
If $m,K$ are positive integers, $L^m$ is base-point free, and $m\mid K$, then $L^K$ is base-point free, and its complete linear system separates every pair of points and every tangent vector already separated by that of $L^m$.
\end{lemma}
\begin{proof}
Set $\ell=K/m$ and choose a basis $s_0,\ldots,s_N$ of $H^0(X,L^m)$.
For each multi-index $|\alpha|=\ell$, the product $s^\alpha$ is a section of $L^K$.
At every point some $s_i\ne0$, hence $s_i^\ell\ne0$; these products have no common zero.
Their projective map is $\Phi_m$ followed by the degree-$\ell$ Veronese map, defined by $[Z_0:\cdots:Z_N]\mapsto[Z^\alpha]_{|\alpha|=\ell}$.
Even if the products satisfy linear relations, they still form a base-point-free generating family.

In the chart $s_i\ne0$, the product coordinates recover
\[
 \frac{s_j}{s_i}
 =\frac{s_js_i^{\ell-1}}{s_i^\ell}.
\]
Equality of the product-map values therefore forces equality of all affine coordinates of the original map.
Likewise, vanishing of its differential on a vector forces vanishing of the differentials of all original coordinate ratios.
It thus preserves the existing separation of points and tangent vectors.

The span of the product sections is a subspace of $H^0(X,L^K)$.
Extend a basis of that subspace to a basis of the full space; the preceding ratios remain quotients of coordinates of the complete projective map.
Adding sections therefore preserves separation and cannot introduce new base points.
\end{proof}

\begin{theorem}[Kodaira embedding: positive and ample line bundles]\label{emb:kodaira}
A holomorphic line bundle $L$ on a compact complex manifold is positive if and only if it is ample.
In particular, a compact complex manifold carrying a positive line bundle admits a holomorphic embedding into some complex projective space.
\end{theorem}
\begin{proof}
Proposition \ref{bun:generated-semipositive} already proves that ample implies positive.
Conversely, use Corollary \ref{emb:local-separation} to choose finitely many open sets $U_1,\ldots,U_a$ covering $X$
and base-point-free powers $m_1,\ldots,m_a$ such that $\Phi_{m_i}$ is injective and immersive on $U_i$.
The set $W=\bigcup_i U_i\times U_i$ is an open neighborhood of the diagonal, so
$C=(X\times X)\setminus W$ is compact and contains no diagonal point.
For every $(p,q)\in C$, Proposition \ref{emb:pair-separation} gives a base-point-free power whose map separates $p$ and $q$.
Separation persists on a neighborhood of $(p,q)$ because projective space is Hausdorff and the map is continuous.
Compactness gives a finite subcover by these neighborhoods, involving only finitely many additional powers.
Let $K$ be a common multiple of all these powers. Lemma \ref{emb:veronese} makes $\Phi_K$ injective and immersive on every $U_i$ and separates every pair in $C$.
Thus $\Phi_K$ is globally injective and immersive everywhere.
We finish the passage from an injective immersion to an embedding.
Closed subsets of a compact space are compact, and their continuous images in Hausdorff projective space are closed.
Hence the continuous injection $\Phi_K$ is a homeomorphism onto its image.
At each point choose projective affine coordinates in which the differential has a nonzero $n\times n$ minor.
By Theorem~\ref{bas:holomorphic-inverse}, those $n$ components are local coordinates,
and the other components are holomorphic functions of them, so the image is locally a holomorphic graph.
The homeomorphism property rules out additional local branches coming from distant points.
Therefore $\Phi_K$ is a holomorphic embedding, and the linear system of these $L^K$ sections makes $L^K$ very ample.
\end{proof}

\section{Scope and effective questions}
\begin{corollary}[A uniform positive lower bound relative to any background metric]\label{emb:uniform-positive}
Fix any Hermitian form $\omega_0$ on a compact complex manifold.
A line bundle $L$ is ample if and only if it has a smooth metric $h$ and a constant $\varepsilon>0$ with $\ii F_{L,h}\ge\varepsilon\omega_0$.
\end{corollary}
\begin{proof}
The inequality gives strictly positive curvature, so Theorem \ref{emb:kodaira} makes $L$ ample.
Conversely, Theorem~\ref{emb:kodaira} gives a positive-curvature metric, whose smallest curvature eigenvalue relative to $\omega_0$ is a continuous positive function.
Its minimum on the compact space supplies $\varepsilon>0$.
\end{proof}

\begin{remark}[Existence versus an explicit power]
This chapter produces a usable power, without asserting that it depends only on dimension.
Coordinate radii, curvature lower bounds, metric derivative bounds, and global geometry enter the peak-section constants.
Bounds for a very ample power require further information about vanishing orders and zero-dimensional subschemes.
Chapter~\ref{eff:chapter} uses multiplier ideals and Nadel vanishing to obtain effective jet-separation bounds.
For these methods, see \cite{Demailly2012,DemaillyAnalytic2012,Lazarsfeld2004}.
\end{remark}

\begin{remark}[Ampleness does not mean triviality of the canonical bundle]
Proposition \ref{bun:sections-p1} gives $K_{\mathbb P^1}=\mathcal O(-2)$, while $\mathbb P^1$ is itself projective.
Projectivity therefore does not imply $c_1(K_X)=0$. By contrast, the first Chern class of a positive line bundle has strictly positive integral over every projective line lying in $X$.
\end{remark}


\begin{example}[An explicit conic from the abstract embedding argument]\label{emb:worked-veronese}
The three sections $Z_0^2,Z_0Z_1,Z_1^2$ of $\cO(2)$ on $\mathbb P^1$ define an embedding
\[
 \nu_2([Z_0:Z_1])=[Z_0^2:Z_0Z_1:Z_1^2],
\]
whose image is exactly the smooth conic $X_0X_2=X_1^2\subset\mathbb P^2$.
\end{example}
\begin{proof}
Simultaneous vanishing of the three coordinates would force $Z_0=Z_1=0$, which is impossible, so the map is everywhere defined.
In the chart $Z_0\ne0$, put $z=Z_1/Z_0$; the map is $z\mapsto[1:z:z^2]$.
On the image chart $X_0\ne0$, recover the source coordinate as $z=X_1/X_0$.
Its derivative is one, so the map is injective and immersive in this chart.
In the chart $Z_1\ne0$, recover $w=Z_0/Z_1=X_1/X_2$ in the same way.
These two charts cover the source, and the inverse formulas satisfy $w=1/z$ on their overlap. Thus the map is globally injective and immersive.
Compactness makes it a homeomorphism onto its image, hence a holomorphic embedding.

Image points satisfy $X_0X_2=X_1^2$. Conversely, a projective point satisfying this equation with $X_0\ne0$
has the form $[1:z:z^2]$. If $X_0=0$, the equation forces $X_1=0$,
leaving only $[0:0:1]=\nu_2([0:1])$.
The image is therefore precisely this conic.
The gradient $(X_2,-2X_1,X_0)$ of its defining equation cannot vanish at a projective point,
so the curve is smooth.
The coordinate ratios give local inverses to the projective map, while the equation identifies its image globally.
\end{proof}


\chapter{Integral Classes, Projectivity, and Compact Riemann Surfaces}\label{prj:chapter}
\begin{learning}
The curvature of a line bundle represents an integral cohomology class. Conversely, an integral real $(1,1)$ form determines a holomorphic line bundle with that curvature. This links Kodaira's embedding theorem to topology. On a compact curve, the exact sequence at one point gives Riemann--Roch, and its numerical consequences show that every compact Riemann surface embeds in $\mathbb P^3$.
\end{learning}

\section{Constructing line bundles from integral degree-two classes}
A \emph{good cover} is an open cover for which every nonempty finite intersection is contractible. We use the standard topological comparison stated in Appendix~\ref{dep:good-cover}: on a manifold, constant-sheaf cohomology agrees with singular cohomology; a good cover computes it by its \Cech{} complex, and the real de Rham comparison identifies the class of a closed form with a real constant cocycle. This comparison is a topological input, not a consequence of fine-sheaf acyclicity for the constant integer sheaf; see \cite{BottTu1982}.
The class of a real closed two-form $\alpha$ is called integral if it lies in the image of $H^2(X,\Z)\to H^2(X,\R)$. The concrete gluing argument realizing an integral class is supplied next.

\begin{lemma}[Integrality and gluing connections]\label{prj:integral-gluing}
A real closed two-form $\alpha$ is the first Chern form of a unitary connection on a smooth Hermitian line bundle
if and only if $[\alpha]$ is integral.
If $X$ is complex and $\alpha$ has type $(1,1)$, the line bundle admits a holomorphic structure for which the given connection is the Chern connection.
\end{lemma}
\begin{proof}
Choose a locally finite good covering $\{U_i\}$, meaning that all nonempty finite intersections are contractible.
On a manifold, sufficiently small strongly convex metric balls, chosen locally finitely, give such a covering:
the unique short minimizing geodesic between two points lies in every strongly convex ball containing them, hence also in their intersection;
geodesic contraction toward a fixed point contracts the intersection.
Proposition~\ref{bas:poincare} gives real one-forms $a_i$ on $U_i$ with $da_i=2\pi\alpha$.
On $U_i\cap U_j$, the form $a_i-a_j$ is closed, so write it as $df_{ij}$ with $f_{ij}$ real;
choose $f_{ji}=-f_{ij}$ and $f_{ii}=0$.
On triple intersections,
\[
 c_{ijk}=\frac{f_{ij}+f_{jk}+f_{ki}}{2\pi}
\]
has zero differential and is therefore constant. The alternating sum on quadruple intersections cancels term by term, so $\delta c=0$.

The \Cech{}--de Rham double complex identifies $[\alpha]$ with the real constant cocycle $[c]$:
first take local primitives $a_i/(2\pi)$, then primitives $f_{ij}/(2\pi)$ on overlaps,
and finally take their alternating differences to obtain $c$.
These are exactly the two elimination steps in Lemma \ref{sh:double-complex}.
The constant integer sheaf has no higher cohomology on the contractible intersections of a good covering,
so an integral class is represented on the same covering by an integer two-cocycle $n_{ijk}$.
Integrality is consequently equivalent to the existence of a real constant one-cochain $b_{ij}$ with
$c-n=\delta b$.
Replacing $f_{ij}$ by $f_{ij}-2\pi b_{ij}$ leaves $df_{ij}=a_i-a_j$ unchanged,
but makes the triple sums exactly $2\pi n_{ijk}$.
This adjustment by constants ensures that the connection realizes the prescribed form $\alpha$, not merely some form in the same cohomology class.

Assuming integrality, set $g_{ij}=e^{\ii f_{ij}}$.
The triple sums are integral multiples of $2\pi$, so $g_{ij}g_{jk}g_{ki}=1$.
Glue unit frames by $e_j=e_i g_{ij}$ using Proposition~\ref{bun:gluing} to obtain a Hermitian line bundle.
Take local connection forms $A_i=-\ii a_i$. Since
\[
 A_i+g_{ij}^{-1}dg_{ij}
 =-\ii a_i+\ii\,df_{ij}=-\ii a_j=A_j,
\]
they glue to a unitary connection. Its curvature is
$F=dA_i=-\ii\,da_i=-2\pi\ii\alpha$, giving $\ii F/(2\pi)=\alpha$.
Conversely, for a unitary connection, the transition functions of local unit frames take values in the unit circle.
Choose real arguments $f_{ij}$ on contractible overlaps.
The transition cocycle makes the triple sums integral multiples of $2\pi$, while the connection transformation law gives
$a_i-a_j=df_{ij}$, where $A_i=-\ii a_i$.
The same double-complex computation sends the first Chern class to this integer cocycle, so the class is integral.

Finally, let $\alpha$ have type $(1,1)$. The $(0,2)$ curvature component vanishes, hence
$(D^{0,1})^2=0$.
In a local smooth frame write $D^{0,1}=\dbar+A^{0,1}$; then
$\dbar A^{0,1}=0$.
Theorem~\ref{sh:dolbeault-lemma} gives $A^{0,1}=\dbar u$.
Change the frame to $e'=e e^{-u}$. Direct calculation gives
$D^{0,1}e'=e^{-u}(A^{0,1}-\dbar u)e=0$.
The transition functions between two such frames are annihilated by $\dbar$, hence holomorphic.
These frames define a holomorphic structure. The original connection preserves the metric, and its $(0,1)$ part is the $\dbar$ of this structure,
so the uniqueness assertion in Theorem~\ref{bun:chern-connection} identifies it with the Chern connection.
\end{proof}

\begin{theorem}[Kodaira's integral K\"ahler class criterion]\label{prj:hodge-projective}
For a compact complex manifold $X$, the following are equivalent: $X$ admits a holomorphic embedding in projective space; $X$ carries a positive holomorphic line bundle; $X$ carries a K\"ahler form with integral cohomology class.
If $X$ is compact K\"ahler and $H^{2,0}(X)=0$, then $X$ is projective.
\end{theorem}
\begin{proof}
For an embedding $i:X\hookrightarrow\mathbb P^N$, the Fubini--Study calculation in Proposition~\ref{bun:fs} makes the curvature of $i^*\mathcal O(1)$ strictly positive and its first Chern class is integral by Lemma \ref{prj:integral-gluing}. The first Chern form of a positive line bundle is an integral K\"ahler form. Conversely, the same lemma realizes an integral K\"ahler form as the positive first Chern form of a holomorphic line bundle, and Theorem \ref{emb:kodaira} gives an embedding. This proves the three equivalences.

Now suppose $H^{2,0}=0$. Proposition~\ref{hod:kahler-laplacians} gives $H^{0,2}=0$ and shows that every real degree-two cohomology class has type $(1,1)$.
Fix a K\"ahler form $\omega$. Rational classes are dense in the finite-dimensional space $H^2(X,\R)$: after choosing a basis for integer cohomology modulo torsion, this is density of $\Q^b$ in $\R^b$. Choose rational classes $a_j\to[\omega]$, and represent their differences from $[\omega]$ by smooth real harmonic $(1,1)$-forms $\eta_j$. The harmonic space is finite-dimensional, so convergence of coefficients gives $\eta_j\to0$ in every $C^k$ norm.
Positive definiteness is open in the $C^0$ topology, so $\omega+\eta_j$ remains K\"ahler for large $j$. Multiplication by a positive integer clears the denominators of $a_j$, giving an integral K\"ahler class. The first part applies.
\end{proof}

\begin{proposition}[Reducing the dimension of a projective embedding]\label{prj:dimension-bound}
A compact complex manifold of dimension $n$ that embeds in some projective space also embeds in $\mathbb P^{2n+1}$.
\end{proposition}
\begin{proof}
Suppose $\iota:X\hookrightarrow\mathbb P^N$ is an embedding and $N>2n+1$.
We must estimate the set of all points on secant lines and tangent spaces, including the position parameters on the lines themselves.
Locally choose nonzero vector representatives $v(p),v(q)\in\C^{N+1}$ for $\iota(p),\iota(q)$.
For $p\ne q$, these vectors are linearly independent, so
\[
 (p,q,[a:b])\longmapsto[av(p)+bv(q)]
\]
is everywhere defined on the corresponding off-diagonal parameter domain and covers every point of these secants.
The parameter space has complex dimension $n+n+1=2n+1$; the final dimension records position along the secant line.

In a coordinate chart take a vector lift $v(z)$ of the embedding.
The vectors $v,\partial_{z_1}v,\ldots,\partial_{z_n}v$ are linearly independent because $\iota$ is immersive.
All points of the projective tangent spaces are parametrized by
\[
 (z,[a_0:\cdots:a_n])\longmapsto
 [a_0v(z)+\textstyle\sum_{j=1}^n a_j\partial_{z_j}v(z)]
\]
with complex parameter dimension $n+n=2n$.
These tangent spaces include $X$ itself, so a projection center avoiding them lies outside $X$.

Second countability provides countably many such parameter charts. Exhaust each by countably many compact coordinate blocks;
the smooth map to the target is Lipschitz on each block.
A real $d$-dimensional parameter block can be covered by $O(\varepsilon^{-d})$ cubes of side $\varepsilon$,
whose images are covered by the same number of balls of radius $C\varepsilon$.
For $d<2N$, the resulting bound on $2N$-dimensional outer measure is
$C'\varepsilon^{2N-d}\to0$.
For secant points $d\le4n+2$, and for tangent-space points $d\le4n$.
Since $2N>4n+2$, both countable unions have measure zero in projective space.
Choose a point $a$ outside both sets.

Project from $a=[a_0]$ to $\mathbb P(\C^{N+1}/\C a_0)\simeq\mathbb P^{N-1}$.
This projection is defined on $X$.
If distinct points had the same image, their vector representatives would be proportional modulo $\C a_0$,
putting $a_0$ in their span; this says that $a$ lies on their secant line, a contradiction.
If the differential at $p=[v]$ annihilated a nonzero tangent vector represented by $\dot v$,
then $\dot v$ modulo $\C a_0$ would lie in the line of $v$, namely
$\dot v\in\operatorname{span}(v,a_0)$.
The tangent vector is nonzero, so $\dot v\notin\C v$, hence
$a_0\in\operatorname{span}(v,\dot v)$. This puts $a$ in the projective tangent space at $p$, again a contradiction.
The restricted projection is therefore an injective holomorphic immersion. The compactness and local-graph argument at the end of Theorem~\ref{emb:kodaira} makes it an embedding.
Each projection lowers the ambient dimension by one; repeat until $N=2n+1$.
If the original $N\le2n+1$, simply compose with a linear embedding.
\end{proof}

\section{Divisors, the exact sequence at a point, and Riemann--Roch}
Let $C$ be a connected compact Riemann surface and set $g=h^1(C,\mathcal O_C)$, its genus. For a line bundle $L$, write $\mathcal L=\cO_C(L)$ for its section sheaf, $H^q(L)=H^q(C,\mathcal L)$, $h^q(L)=\dim_{\C}H^q(L)$, and $\chi(L)=h^0(L)-h^1(L)$. The fiber is $L_p=\mathcal L_p/\mathfrak m_p\mathcal L_p$, while $\mathcal L_p$ is the stalk of holomorphic germs.
A divisor is a finite sum $D=\sum_p m_pp$, $m_p\in\Z$. A nonzero meromorphic function has the local expression $f=z^m a(z)$ with $a(0)\ne0$ holomorphic; define $\operatorname{ord}_p(f)=m$, allowing negative orders at poles. Sections of $\mathcal O(D)$ are local meromorphic functions satisfying $\operatorname{ord}_p(f)+m_p\ge0$.
Multiplication locally by $z^{m_p}$ trivializes this sheaf, so it is a line bundle. Define $\deg D=\sum_p m_p$, and write $L(D)=L\otimes\mathcal O(D)$.
The twist $L(p)$ allows a section of $L$ to have a simple pole at $p$; the twist $L(-p)$ requires it to vanish there. Evaluation at $p$ therefore compares $L(-p)$ with $L$, as the next exact sequence records.

\begin{lemma}[The length of a point and change in Euler characteristic]\label{cur:point-sequence}
For a line bundle $L$ and a point $p\in C$, there is a short exact sequence
\[
 0\longrightarrow L(-p)\longrightarrow L\longrightarrow \left.L\right|_p\longrightarrow0,
\]
whose final term is the skyscraper sheaf at $p$ with value the one-dimensional fiber $L_p=\mathcal L_p/\mathfrak m_p\mathcal L_p$; its global section space is one-dimensional. All line-bundle cohomology groups are finite-dimensional, and
$\chi(L)-\chi(L(-p))=1$.
Every line bundle admits a nonzero meromorphic section.
\end{lemma}
\begin{proof}
In a local coordinate at $p$, the map on holomorphic coefficients is $f\mapsto f(p)$, with kernel $z\mathcal O$ and quotient $\C$. At every other point the quotient sheaf vanishes. This proves exactness on stalks, hence exactness of sheaves by Proposition~\ref{sh:stalk-exact}.
Finite-dimensionality follows from Theorems~\ref{sh:dolbeault} and~\ref{hod:hodge} for compact Hermitian manifolds; on a curve there is no Dolbeault cohomology above degree one. Proposition~\ref{sh:discrete-support} and Theorem~\ref{sh:acyclic} give zero higher cohomology for the skyscraper sheaf: a section over a closed set requires only the value at that point. Proposition~\ref{sh:long-exact} therefore gives the finite-dimensional sequence
\[
 0\to H^0(L(-p))\to H^0(L)\to \left.L\right|_p\to H^1(L(-p))\to H^1(L)\to0.
\]
Taking alternating dimensions gives the asserted difference.

Finally, fix a point $p$ and, for a positive integer $m$, consider
$0\to L\to L(mp)\to Q_m\to0$.
Near $p$, choose a coordinate $z$ and a holomorphic frame $e$ of $L$. The quotient
$z^{-m}\mathcal O e/\mathcal O e$ has complex basis given by the classes of $z^{-1}e,\ldots,z^{-m}e$.
Since $Q_m$ vanishes at all other points, $h^0(Q_m)=m$.
The kernel of $H^0(Q_m)\to H^1(L)$ in the long exact sequence has dimension at least $m-h^1(L)$.
Choose $m>h^1(L)$ and a nonzero element of this kernel.
Exactness lifts it to a global section of $L(mp)$, necessarily nonzero since its principal part is nonzero.
Viewed as a section of $L$, it has a pole of order at most $m$ at $p$ and is holomorphic elsewhere, giving the required meromorphic section.
The only inputs are the proved finite-dimensionality and the exact sequence at a point, so the argument applies to every compact Riemann surface before projectivity is established.
\end{proof}

\begin{proposition}[Three expressions for degree]\label{cur:degree}
The zero-pole divisor $D_s$ of a nonzero meromorphic section $s$ satisfies $L\simeq\mathcal O(D_s)$, and
\[
 \deg L:=\int_Cc_1(L)=\deg D_s.
\]
The total number of zeros of a nonzero holomorphic section, counted with multiplicity, equals $\deg L$; a line bundle of negative degree has no nonzero holomorphic section.
\end{proposition}
\begin{proof}
Zeros and poles of a nonzero meromorphic section are isolated. Compactness makes them finite in number, so $D_s$ is a finite divisor.
In a local frame $e$, write $s=f e$, with $f=z^m u$ and $u$ holomorphic and nowhere zero. Multiplication by $s$ maps local sections of $\mathcal O(D_s)$ invertibly to holomorphic sections of $L$. The transition relations make these local isomorphisms glue.
Choose any smooth Hermitian metric. Apply Lemma~\ref{pos:one-variable-lelong} to each signed order $m$, including negative orders at poles, and use the metric calculation in \eqref{pos:section-lelong}. This gives
$\frac{\ii}{2\pi}\partial\dbar\log|s|_h^2=[D_s]-c_1(L,h)$.
Pair with the constant function $1$. The left-hand side vanishes because a derivative of a current acts as zero on a constant; the first term on the right is the sum of zero and pole multiplicities. This proves the integral identity.
A holomorphic section has no negative local multiplicities, so the total is nonnegative. A nonzero holomorphic section of a negative-degree line bundle would therefore be a contradiction.
\end{proof}

\begin{theorem}[Riemann--Roch for curves]\label{cur:rr}
Every line bundle $L$ on a compact Riemann surface satisfies
\[
 h^0(L)-h^0(K_C\otimes L^{-1})=\deg L+1-g.
\]
Moreover, $\deg K_C=2g-2$.
\end{theorem}
\begin{proof}
Take a nonzero meromorphic section by Lemma \ref{cur:point-sequence} and apply Proposition \ref{cur:degree}; every line bundle is thus isomorphic to some $\mathcal O(D)$. Increase or decrease the divisor coefficients one point at a time. Lemma~\ref{cur:point-sequence} changes $\chi$ by one at each unit step, so
$\chi(L)=\deg D+\chi(\mathcal O_C)$.
Theorem~\ref{bas:compact-holomorphic} makes holomorphic functions on a compact connected Riemann surface constant, hence $h^0(\mathcal O_C)=1$; the definition of $g$ gives $\chi(\mathcal O_C)=1-g$. With the degree formula this yields
$h^0(L)-h^1(L)=\deg L+1-g$.
Serre duality, Theorem \ref{hod:serre-duality}, gives $h^1(L)=h^0(K_C\otimes L^{-1})$, proving the formula.
Now take $L=K_C$. Duality gives $h^0(K_C)=h^1(\mathcal O_C)=g$ and $h^1(K_C)=h^0(\mathcal O_C)=1$. Substitution in the Euler-characteristic formula gives
$g-1=\deg K_C+1-g$, equivalently $\deg K_C=2g-2$.
This is an independent proof for curves, without deducing the result from higher-dimensional Hirzebruch--Riemann--Roch.
\end{proof}

\begin{theorem}[Generation, very ampleness, and ampleness on curves]\label{cur:ample}
If $\deg L\ge2g$, then $L$ is globally generated; if $\deg L\ge2g+1$, then $L$ is very ample.
Consequently $L$ is ample if and only if $\deg L>0$. Every compact Riemann surface is projective and embeds in $\mathbb P^3$.
\end{theorem}
\begin{proof}
For $\deg L\ge2g$, Theorem~\ref{cur:rr} makes the degree of $K_C\otimes L^{-1}(p)$ at most $-1$, so Proposition~\ref{cur:degree} excludes nonzero holomorphic sections. Theorem~\ref{hod:serre-duality} gives $H^1(L(-p))=0$. Lemma~\ref{cur:point-sequence} then makes evaluation $H^0(L)\to L_p$ surjective. Since $p$ was arbitrary, there are no base points.
If $\deg L\ge2g+1$, then for any points $p,q$, allowing equality, the same argument gives
$H^1(L(-p-q))=0$. For $p\ne q$, the quotient sheaf is $\left.L\right|_p\oplus\left.L\right|_q$; surjectivity of evaluation prescribes the two values independently, so the linear system separates the points.
For $p=q$, the local quotient is $\mathcal O/(z^2)$, and surjectivity prescribes constant and linear terms independently. Use a section nonzero at $p$ as denominator and one with zero constant term and nonzero linear term as numerator. Their ratio has nonzero derivative at $p$, so the map is immersive.
The complete linear-system construction before Corollary~\ref{emb:local-separation} and the final argument of Theorem~\ref{emb:kodaira} now give a holomorphic embedding.

If $\deg L>0$, choose $m$ with $m\deg L\ge2g+1$. The preceding argument makes $L^m$ very ample, so $L$ is ample. Conversely, Theorem~\ref{emb:kodaira} gives positive curvature, whose first Chern form has strictly positive integral over the curve.
For any point $p$, the bundle $\mathcal O(p)$ has degree one and is therefore ample. Thus every compact complex curve is projective; Proposition \ref{prj:dimension-bound}, with $n=1$, gives an embedding in $\mathbb P^3$.
\end{proof}

\begin{proposition}[Total residue and the obstruction to a unique simple pole]\label{cur:residues}
The sum of the residues of a meromorphic one-form on a compact Riemann surface is zero. Such a form therefore cannot have exactly one pole if that pole is simple.
If $g\ge1$, there is also no meromorphic function on the compact surface with exactly one pole of order one.
\end{proposition}
\begin{proof}
Remove disjoint small disks around all poles. On the remaining surface with boundary the meromorphic one-form is holomorphic, hence closed as a complex-valued real one-form. Proposition~\ref{bas:stokes} makes the sum of its boundary integrals zero. The boundary orientation is opposite to the positive disk orientation, and each circle integral is $2\pi\ii$ times the residue, so the total residue is zero. At a simple pole the Laurent coefficient $a_{-1}$ is nonzero and is precisely the residue, ruling out a unique simple pole.
If a meromorphic function $f$ had only a simple pole at $p$, use $g=h^0(K_C)\ge1$ to choose a nonzero holomorphic one-form $\eta$. Write its vanishing order at $p$ as $k\ge0$, so $\eta=z^k u(z)dz$ with $u(0)\ne0$. Since $f=z^{-1}v(z)$ with $v(0)\ne0$, the form $f^{k+1}\eta=z^{-1}v(z)^{k+1}u(z)dz$ has a simple pole at $p$ and is holomorphic elsewhere. This contradicts the preceding paragraph.
\end{proof}

\section{Worked computations}
\begin{example}[The projective line]
Proposition \ref{bun:sections-p1} gives $K=\mathcal O(-2)$ on $\mathbb P^1$, with no nonzero holomorphic sections. Serre duality yields $g=h^1(\mathcal O)=h^0(K)=0$. Riemann--Roch gives
$h^0(\mathcal O(d))=d+1$ for $d\ge0$, and $h^1(\mathcal O(d))=0$ for $d\ge-1$; for $d\le-2$, one has $h^0=0$ and $h^1=-d-1$.
Indeed, for nonnegative $d$, the sections are homogeneous binary polynomials of degree $d$, with $d+1$ monomials; negative-degree bundles have no sections. The identity $h^1(\mathcal O(d))=h^0(\mathcal O(-d-2))$ gives the remaining cases.
\end{example}
\begin{example}[An elliptic curve]
For $C=\C/(\Z+\tau\Z)$ with $\operatorname{Im}\tau>0$, the form $dz$ descends to a nowhere-vanishing holomorphic one-form, so $K_C\simeq\mathcal O_C$. Compact connectedness and Serre duality give $g=h^0(K_C)=h^0(\mathcal O_C)=1$. A line bundle of positive degree $d$ has $h^0(L)=d$, because $L^{-1}$ has negative degree and no sections. Theorem~\ref{cur:ample} and $h^0(L)=3$ show that degree three guarantees an embedding in $\mathbb P^2$. Degree one gives only one section, which must have a zero, so the bundle is not yet globally generated. The degree thresholds on curves thus have substantive content.
\end{example}


\chapter{Cohomology of Projective Bundles and Relative Computations}\label{sh:projective-chapter}

A projective bundle is locally a product with projective space. Its fiberwise homogeneous polynomials become symmetric powers of the original vector bundle, and this correspondence extends to cohomology. We first compute on projective space, then use Hodge operators to solve along the fibers while preserving dependence on the base point. Local free resolutions extend the resulting formulas from bundles to coherent sheaves.
The analytic prerequisites are Hodge decomposition (Theorem~\ref{hod:hodge}) and Kodaira vanishing (Corollary~\ref{pos:kodaira}); the algebraic prerequisite is the finite local resolution of Theorem~\ref{alg:regular-resolution}.

\section{Line bundles and differential forms on projective space}

\begin{lemma}[The Euler sequence and the canonical bundle of projective space]\label{sh:projective-euler}
Let $V$ be a complex vector space of dimension $r$, let $P=\mathbb P(V^*)$ parametrize lines in $V^*$,
and let $L=\cO_P(1)$ be the dual of the negative tautological line bundle. Then
\[
 0\longrightarrow\Omega_P^1\otimes L
 \longrightarrow V\otimes\cO_P\longrightarrow L\longrightarrow0
\]
is exact, and $K_P\simeq L^{-r}\otimes\det V$.
Choosing a volume element of $V$ gives the usual $K_{\mathbb P^{r-1}}\simeq\cO(-r)$.
\end{lemma}
\begin{proof}
At $[\ell]\in P$, the fiber is $L_{[\ell]}=\ell^*$.
The map $V\to\ell^*$ restricts linear functions in $V=(V^*)^*$ to $\ell$;
it is surjective, with kernel the functions annihilating $\ell$.
An infinitesimal variation of a line is represented by a graph, so
$T_{P,[\ell]}=\Hom(\ell,V^*/\ell)$.
Dualizing gives
$T^*_{P,[\ell]}\otimes\ell^*=(V^*/\ell)^*$, 
which is exactly that kernel.
In projective coordinate charts these identifications are given by holomorphic matrices, hence patch to the stated exact sequence.

Take determinants of the rank-$(r-1)$ bundle $\Omega_P^1$.
In a short exact sequence, the determinant of the middle term is the tensor product of the determinants of the two outer terms;
this follows directly by lifting local frames and writing the resulting block upper-triangular transition matrix.
Consequently
\[
 \det V\otimes\cO_P
 =\det(\Omega_P^1\otimes L)\otimes L
 =K_P\otimes L^{r-1}\otimes L.
\]
Rearranging gives the formula.
\end{proof}

\begin{proposition}[The required line-bundle cohomology on projective space]\label{sh:projective-line-cohomology}
On $P=\mathbb P(V^*)$, for $r\ge2$,
\[
 H^0(P,\cO(k))=\operatorname{Sym}^kV\quad(k\ge0),
 \qquad H^0(P,\cO(k))=0\quad(k<0),
\]
and
\[
 H^q(P,\cO(k))=0\qquad(q>0,\ k\ge-r+1).
\]
For $r=1$, $P$ is a point and all corresponding positive-degree groups vanish.
\end{proposition}
\begin{proof}
First compute global sections.
In the affine chart $Z_0\ne0$, put $z_j=Z_j/Z_0$; a section is represented by an entire function $f(z)$.
Its Fubini--Study norm is
$|f(z)|(1+|z|^2)^{-k/2}$.
Compactness of $P$ bounds this norm by a common constant $C$, hence
$|f(z)|\le C(1+|z|^2)^{k/2}$.
For $k\ge0$, apply Theorem~\ref{bas:cauchy} on a polydisk of radius $R\ge1$ to any multi-index $\alpha$:
\[
 |\partial^\alpha f(0)|\le C_\alpha R^{k-|\alpha|}.
\]
If $|\alpha|>k$, letting $R\to\infty$ makes this derivative zero.
Thus $f$ is a polynomial of degree at most $k$. Its homogenization
$Z_0^kf(Z_1/Z_0,\ldots,Z_{r-1}/Z_0)$ is a degree-$k$ homogeneous polynomial on $V^*$,
that is, an element of $\operatorname{Sym}^kV$.
Conversely, every such homogeneous polynomial defines a section of $L^k$ by its scaling law, and the constructions are inverse.
For $k<0$, the same bound makes $f$ bounded, hence constant by the Cauchy estimate;
the growth bound as $|z|\to\infty$ then forces the constant to be zero.

For higher cohomology use the already proved Kodaira vanishing theorem.
The Fubini--Study curvature of $\cO(1)$ is positive by Proposition \ref{bun:fs}.
Lemma~\ref{sh:projective-euler} gives
$\cO(k)=K_P\otimes\cO(k+r)\otimes(\det V)^*$.
The constant line factor does not affect vanishing. When $k\ge-r+1$, we have $k+r\ge1$,
so $\cO(k+r)$ has positive curvature.
Corollary \ref{pos:kodaira} gives vanishing for every $q>0$.
\end{proof}

\begin{proposition}[Projective differential forms with one positive twist]\label{sh:projective-differentials}
For $1\le j\le r-1$ and every $q\ge0$,
\[
 H^q(\mathbb P(V^*),\Omega_P^j\otimes\cO(1))=0.
\]
\end{proposition}
\begin{proof}
Take the Euler surjection of Lemma~\ref{sh:projective-euler}. Its exterior-power Koszul sequence, with the contraction convention of Lemma~\ref{alg:koszul}, is
\[
0\to\Omega_P^j\otimes L
\to\bigwedge^jV\otimes L^{1-j}
\to\bigwedge^{j-1}V\otimes L^{2-j}\to\cdots
\to V\otimes\cO\to L\to0.
\]
Its exactness can be checked directly. On a small open set where $L$ is trivial,
choose a vector section $v$ mapping to $1$.
The free sheaf splits as its kernel plus $\cO v$; the Koszul differential is contraction by this surjection,
and exterior multiplication by $v$ gives a contracting homotopy with $dh+hd=\operatorname{id}$.
The leftmost kernel is the $j$th exterior power of the kernel sheaf, twisted by $L^{1-j}$,
namely $\bigwedge^j(\Omega_P^1\otimes L)\otimes L^{1-j}=\Omega_P^j\otimes L$.
This proves exactness of the entire sequence.

The degrees of all line bundles on the right range from $1-j$ to $1$, hence are at least $-r+1$.
Proposition \ref{sh:projective-line-cohomology} therefore makes them acyclic in positive degrees.
The resolution comparison theorem \ref{sh:resolution-theorem} computes the cohomology of the leftmost term from their global section complex.
Negative-degree line bundles have no sections. The last two nonzero terms are
$V\to H^0(P,L)=V$, the map viewing a linear function as a section, which is the identity.
The global section complex is therefore exact in every degree, and all the asserted groups vanish.
\end{proof}

\section{Fiberwise contractions with parameters}

For a product $U\times\mathbb P^{r-1}$, a section of $\mathcal O(1)$ is linear in the homogeneous fiber coordinates, with coefficients that are holomorphic functions on $U$. This is the degree-zero case of the formulas below. In higher degrees we must also solve the $\dbar$ equation along each fiber and show that its solutions depend smoothly on the base parameters.

\begin{definition}[Higher direct images]
For a continuous map $\pi:Y\to X$ and a sheaf $\mathcal{B}$, define
$R^q\pi_*\mathcal{B}$ as the $q$th cohomology sheaf of $\pi_*G^\bullet(\mathcal{B})$, where $G^\bullet(\mathcal B)$ is the Godement resolution \eqref{sh:godement-resolution}.
Because filtered direct limits preserve finite exact sequences of abelian groups,
\[
 (R^q\pi_*\mathcal{B})_x
 =\varinjlim_{U\ni x}H^q(\pi^{-1}U,\mathcal{B}).
\]
For $q=0$ this is the ordinary direct image; positive degrees measure cohomology that persists as neighborhoods in the base shrink.
\end{definition}

\begin{lemma}[A parameter computation for finite-dimensional fiber cohomology]\label{sh:fiber-hodge-parameter}
Let $M$ be compact Hermitian and $B\to M$ a holomorphic vector bundle with
$H^b(M,\cO(B))=0$ for every $b>0$.
Here $\cO(B)$ is the sheaf of holomorphic sections of $B$. Set $V=H^0(M,\cO(B))$, and let $\operatorname{pr}_M:U\times M\to M$ be the second projection. For the first projection $p:U\times M\to U$, where $U$ is a complex manifold,
\[
 p_*\cO(\operatorname{pr}_M^*B)=\cO_U\otimes V,
 \qquad R^qp_*\cO(\operatorname{pr}_M^*B)=0\quad(q>0).
\]
If all fiber cohomology groups vanish, including $b=0$, then all direct images, including degree zero, vanish.
\end{lemma}
\begin{proof}
Take a coordinate polydisk $D\subset U$ in the base.
Decompose the smooth Dolbeault complex on its product with $M$ according to antiholomorphic base degree $a$ and fiber degree $b$,
with differential $D_{\rm tot}=\dbar_U+(-1)^a\dbar_M$.
Fiberwise Hodge theory supplies continuous operators on smooth forms,
$H$ and $h=\dbar_M^*G$, satisfying
\[
 \dbar_Mh+h\dbar_M=I-H.
\]
Continuity follows from the elliptic estimates in Theorem \ref{hod:hodge} and Proposition \ref{hod:green}.
Thus, on forms depending smoothly on base parameters, $h,H$ preserve smooth parameter dependence and commute with parameter derivatives.
Apply continuity to difference quotients and pass to the limit, successively in each derivative order, to verify this assertion.

In bidegree $(a,b)$, put $K=(-1)^ah$.
The two mixed terms involving $\dbar_U$ have opposite signs and cancel; the fiber terms give the Hodge identity above. Hence
\[
 D_{\rm tot}K+KD_{\rm tot}=I-H.
\]
The product Dolbeault complex is therefore homotopy equivalent to the base Dolbeault complex with values in fiberwise harmonic forms.
The hypothesis confines these harmonic forms to $b=0$, where they are exactly the finite-dimensional holomorphic section space $V$.
The resulting complex is $\mathcal A^{0,\bullet}(D)\otimes V$, the smooth base Dolbeault complex.

At a fixed base point, Theorem~\ref{sh:dolbeault-lemma} solves every positive-degree closed element on a smaller base polydisk, component by component in
$V$. Passing to the direct limit over neighborhoods $D$ therefore makes all positive-degree cohomology vanish.
Degree-zero closed elements are finite sums of a fixed basis of $V$ with holomorphic coefficients, so the direct image is $\cO_U\otimes V$.
Theorem~\ref{sh:dolbeault} computes sheaf cohomology on the product, so this is exactly the stalk computation of higher direct images.
If all fiber cohomology vanishes, then $H=0$, and the homotopy identity contracts the whole complex, proving the final assertion.
\end{proof}

\begin{lemma}[Long exact sequences and acyclic direct images]\label{sh:direct-image-comparison}
A short exact sequence of sheaves induces a long exact sequence of $R^q\pi_*$.
If $R^q\pi_*\mathcal{B}=0$ for $q>0$, then naturally
\[
 H^l(Y,\mathcal{B})\simeq H^l(X,\pi_*\mathcal{B})\qquad(l\ge0).
\]
\end{lemma}
\begin{proof}
The functor $\Gamma(\pi^{-1}U,G^a(-))$ is exact by Proposition \ref{sh:godement-exact}.
Applying $\pi_*G^\bullet$ termwise to a short exact sequence therefore gives a termwise exact sequence of complexes.
The same construction by lifting, differentiating, and taking kernels as in Proposition \ref{sh:long-exact}
gives the long exact sequence on open sets and on stalks. The connecting maps commute with restriction, so they assemble into sheaf morphisms.

The direct image of a flabby sheaf is flabby: restriction on base open sets is restriction between their inverse images, and remains surjective.
Thus every $\pi_*G^a(\mathcal{B})$ is flabby.
Vanishing of higher direct images says precisely that
$0\to\pi_*\mathcal{B}\to\pi_*G^0(\mathcal{B})\to\pi_*G^1(\mathcal{B})\to\cdots$
is a flabby resolution.
Its global section complex is termwise identical to $\Gamma(Y,G^\bullet(\mathcal{B}))$.
Theorem~\ref{sh:resolution-theorem} on both spaces gives the asserted natural isomorphism.
\end{proof}

\section{Projective-bundle formulas with coherent coefficients}

\begin{theorem}[All direct images of nonnegative twists on a projective bundle]\label{sh:projective-direct-image}
Let $E$ be a rank-$r$ holomorphic vector bundle on a complex manifold $X$,
$\pi:Y=\mathbb P(E^*)\to X$, $L=\cO_Y(1)$.
For every coherent sheaf $\cF$ and integer $k\ge0$, there is a natural isomorphism and vanishing
\[
 \pi_*(\pi^*\cF\otimes L^k)=\cF\otimes\operatorname{Sym}^kE,
 \qquad
 R^q\pi_*(\pi^*\cF\otimes L^k)=0\quad(q>0).
\]
Consequently, for every $l\ge0$,
\[
 H^l(Y,\pi^*\cF\otimes L^k)
 \simeq H^l(X,\cF\otimes\operatorname{Sym}^kE).
\]
\end{theorem}
\begin{proof}
The direct-image assertions are local on the base. Choose $U$ with $E|_U\simeq U\times V$,
so $Y|_U=U\times\mathbb P(V^*)$.
For $\cF=\cO_U$, Proposition \ref{sh:projective-line-cohomology} and
the parameter lemma \ref{sh:fiber-hodge-parameter} give the result immediately.
For finite free $\cF$, apply it componentwise.

For general coherent $\cF$, take near each base point a finite free resolution from Theorem \ref{alg:regular-resolution},
$0\to F_m\to\cdots\to F_0\to\cF\to0$.
The projective-bundle projection is flat by Theorem \ref{alg:analytic-product-flat}.
Pullback therefore preserves exactness by Proposition~\ref{sh:flat-stalks}, as does tensoring with the locally free $L^k$ by Proposition~\ref{sh:tensor}.
Write the resulting terms as $B_i=\pi^*F_i\otimes L^k$.
Their positive-degree direct images vanish, and their degree-zero images are
$F_i\otimes\operatorname{Sym}^kV$.

Starting at the left end, split the resolution into short exact sequences.
If $0\to A\to B\to C\to0$ has vanishing positive-degree direct images for $A$ and $B$,
Lemma~\ref{sh:direct-image-comparison} gives $R^q\pi_*C=0$ for $q>0$,
and exactness of $0\to\pi_*A\to\pi_*B\to\pi_*C\to0$.
Apply this first to the cokernel of $B_m\to B_{m-1}$, then proceed to the right.
This proves that the final term $\pi^*\cF\otimes L^k$ has no higher direct images,
and that its degree-zero direct image is the cokernel of $F_1\otimes\operatorname{Sym}^kV\to
F_0\otimes\operatorname{Sym}^kV$.
Right exactness of tensor product identifies this cokernel with $\cF\otimes\operatorname{Sym}^kV$.
The resolution is taken on a neighborhood; a general global free resolution need not exist.

Finally, check that the isomorphisms are natural and glue across changes of frame.
The $k$th symmetric power of the tautological evaluation $\pi^*E\to L$ gives
$\pi^*\operatorname{Sym}^kE\to L^k$.
Tensor with $\pi^*\cF$ and use Proposition~\ref{sh:adjunction} to obtain the natural morphism
$\cF\otimes\operatorname{Sym}^kE\to\pi_*(\pi^*\cF\otimes L^k)$.
In the chosen trivialization, this views a homogeneous polynomial as a section on each fiber,
so it is the local isomorphism already proved. Naturality makes these local identifications agree, yielding a global isomorphism.
The global cohomology formula follows from Lemma \ref{sh:direct-image-comparison}.
\end{proof}

For the submersion $\pi:Y\to X$, define $\Omega^1_{Y/X}=\Omega_Y^1/\pi^*\Omega_X^1$ and $\Omega^j_{Y/X}=\bigwedge^j\Omega^1_{Y/X}$. These locally free sheaves describe holomorphic differentials in the fiber directions; on a product chart they are the pullbacks of the cotangent exterior powers of the fiber.

\begin{corollary}[Vanishing for relative differential forms]\label{sh:projective-relative-forms}
Under the hypotheses of Theorem~\ref{sh:projective-direct-image}, for $1\le j\le r-1$ and every $q\ge0$,
\[
 R^q\pi_*(\pi^*\cF\otimes\Omega_{Y/X}^j\otimes L)=0.
\]
For $j=0$, the degree-zero direct image is $\cF\otimes E$, and all higher images vanish.
\end{corollary}
\begin{proof}
In a trivialization of $E$, the relative differential-form bundle is the pullback of the fiber bundle
$\Omega_{\mathbb P(V^*)}^j$.
Proposition \ref{sh:projective-differentials} makes all fiber cohomology vanish after twisting by $L$.
Lemma~\ref{sh:fiber-hodge-parameter} therefore makes all direct images vanish for free coefficients.
For general $\cF$, use the finite local free resolution of Theorem~\ref{alg:regular-resolution}, exactly as in the proof of Theorem~\ref{sh:projective-direct-image}.
Flat pullback and tensoring with the relative vector bundle preserve exactness.
The long exact sequences, moving from left to right, now force even the degree-zero direct images to vanish at every step, proving the assertions.
The case $j=0$ is Theorem~\ref{sh:projective-direct-image} with $k=1$.
\end{proof}


\part{Semester II: Singular Weights and Heat Kernels}


\chapter{Measure Structure, Local Potentials, and Extension of Positive Currents}
\label{cur:chapter}

Positive currents allow curvature and integration over analytic sets to be treated in the same language. Their coefficients are measures, so positivity can be examined almost everywhere through matrix and exterior-algebra calculations. Closedness is subtler: extending a measure across a singular set can create a boundary term. The cutoff estimates in this chapter show when that term vanishes and lead to the general Poincar\'e--Lelong formula.
We use the convention for $dd^c$ in \eqref{pos:dc-convention}.

\section{Positivity Cones in Exterior Algebra}

A nonzero decomposable vector $v_1\wedge\cdots\wedge v_p$ determines the $p$-dimensional plane spanned by its factors. Changing the basis of that plane multiplies the vector by the determinant of the basis change. The generators of the strongly positive cone in Definition~\ref{pos:positive-forms} arise from such planes; this is why decomposability matters for positivity in higher bidegrees.

\begin{lemma}[A kernel criterion for decomposable multivectors]\label{cur:decomposable}
Let $V$ be a complex vector space of dimension $N$ and $0\ne a\in\Lambda^pV$. Put
$L_a(v)=v\wedge a$. Then $\dim\ker L_a\le p$, with equality if and only if
$a$ is the wedge product of $p$ vectors. The decomposable multivectors, together with zero, form a closed set.
\end{lemma}
\begin{proof}
Choose a basis $e_1,\ldots,e_k$ of $\ker L_a$ and extend it to a basis of $V$.
Expand $a=\sum_{|I|=p}a_Ie_I$. The identity $e_j\wedge a=0$ implies
$a_I=0$ whenever $I$ does not contain $j$: the corresponding $(p+1)$-fold basis wedges are distinct
and cannot cancel. Every nonzero term therefore contains $1,\ldots,k$, so $k\le p$.
If $k=p$, the only possible nonzero term is $c e_1\wedge\cdots\wedge e_p$.
Conversely, if $a=e_1\wedge\cdots\wedge e_p\ne0$, expanding an arbitrary vector shows that
$v\wedge a=0$ is equivalent to $v\in\operatorname{span}(e_1,\ldots,e_p)$.

Every matrix coefficient of $L_a$ is linear in the $a_I$.
The condition $\dim\ker L_a\ge p$ is equivalent to $\operatorname{rank}L_a\le N-p$,
or to the vanishing of all minors of size $N-p+1$. These are closed polynomial conditions.
Their common zero set consists exactly of the nonzero decomposable vectors and zero, and is therefore closed.
\end{proof}

\begin{example}[The difference between dimensions three and four]\label{cur:small-dimension}
If $\dim V=3$, every $a\in\Lambda^2V$ is decomposable:
the range of $L_a:V\to\Lambda^3V$ has dimension at most one, so its kernel has dimension at least two;
in the nonzero case, Lemma \ref{cur:decomposable} applies.
If $e_1,e_2,e_3,e_4$ are linearly independent, however,
$a=e_1\wedge e_2+e_3\wedge e_4$ is not decomposable,
because $a\wedge a=2e_1\wedge e_2\wedge e_3\wedge e_4\ne0$,
whereas the wedge square of a decomposable bivector is zero.
\end{example}

\begin{lemma}[When a compactly generated cone is closed]\label{cur:compact-cone}
Let $K$ be compact in a finite-dimensional real vector space, and let $\operatorname{conv}K$ be its convex hull, the set of finite convex combinations of points of $K$. Assume $0\notin\operatorname{conv}K$.
Then $\{t y:t\ge0,\ y\in\operatorname{conv}K\}$ is a closed convex cone.
The conclusion can fail if the assumption $0\notin\operatorname{conv}K$ is removed.
\end{lemma}
\begin{proof}
First, $\operatorname{conv}K$ is compact. In dimension $d$, every convex combination reduces to
at most $d+1$ terms: if there are more, the points are affinely dependent, so there are real numbers $b_j$, not all zero,
with $\sum b_j=0$ and $\sum b_jx_j=0$. Replace each weight $a_j$ by $a_j-tb_j$,
taking the largest $t\ge0$ that leaves every weight nonnegative. At least one weight then vanishes.
Iteration gives the required number of terms. The convex hull is consequently a continuous image of the compact set
$K^{d+1}\times\{a_j\ge0,\sum a_j=1\}$.

Compactness and exclusion of the origin give $\delta>0$ such that $|y|\ge\delta$
throughout the convex hull. If $t_jy_j\to x$, the inequality $t_j\delta\le|t_jy_j|$
bounds $t_j$. Passing to a subsequence with $t_j\to t$ and $y_j\to y$ gives $x=ty$.
The cone is therefore closed. For a counterexample without the extra assumption, take $K=\{(s,s^2):0\le s\le1\}$: its cone contains
$(1,s)$ for $s>0$ but not $(1,0)$, so the cone is not closed.
\end{proof}

\begin{proposition}[Closedness and bipolarity of the strongly positive cone]\label{cur:cones}
On a complex vector space $V$ of dimension $n$, use the positivity definitions of Definition~\ref{pos:positive-forms}. The strongly positive cone of real $(p,p)$ forms is closed and convex. The weakly positive cone is dual to the strongly positive cone in complementary degree,
and the strongly positive cone is dual to the complementary weakly positive cone. Both consist of real forms.
\end{proposition}
\begin{proof}
Choose a Hermitian inner product on $\Lambda^{p,0}V^*$.
The form $\ii^{p^2}a\wedge\bar a$ corresponds to the positive semidefinite rank-one matrix $aa^*$.
For decomposable $a$ with $|a|=1$, these matrices form a compact set $K$:
Lemma \ref{cur:decomposable} gives closedness, and the unit sphere gives boundedness.
Every such matrix has trace one, as does every matrix in its convex hull. The origin is therefore excluded.
Lemma \ref{cur:compact-cone} shows that the generated convex cone is closed; by definition this is the strongly positive cone.

For a cone $C$ in a Hausdorff locally convex real space, define its dual cone by $C^*=\{\ell:\ell\text{ is continuous and real linear},\ \ell(x)\ge0\text{ for every }x\in C\}$. Pairing the space with its continuous dual, every closed convex cone satisfies $C^{**}=C$.
Indeed, $C\subset C^{**}$ is immediate. If $x\notin C$, Hahn--Banach separation
gives a continuous real linear functional $\ell$ and a real number $a$ such that
$\ell(x)<a\le\ell(y)$ for all $y\in C$. Since $0\in C$, we have $a\le0$.
If $\ell(y)<0$ for some $y\in C$, large positive multiples of $y$ contradict the separating inequality.
Thus $\ell\ge0$ on $C$ while $\ell(x)<0$, so $x\notin C^{**}$.
Apply this to the nondegenerate wedge pairing to obtain the two duality statements.
For an arbitrary convex cone, a continuous functional is nonnegative on the cone exactly when it is nonnegative on its closure,
giving the general formula $C^{**}=\overline C$.

Strongly positive forms are visibly real. If weak positivity is defined solely by nonnegative pairing with all strongly positive test forms,
one must additionally show that the imaginary part vanishes. The polarization identity
\[
4a\wedge\bar b=(a+b)\wedge\overline{a+b}-(a-b)\wedge\overline{a-b}
 +\ii(a+\ii b)\wedge\overline{a+\ii b}
 -\ii(a-\ii b)\wedge\overline{a-\ii b}
\]
applied to individual one-form factors shows, after taking real parts, that strongly positive forms span all real $(p,p)$ forms over $\R$.
The imaginary part pairs to zero with this spanning set, so nondegeneracy forces it to vanish.
\end{proof}

\begin{corollary}[A basis of strongly positive forms]\label{cur:positive-basis}
The real vector space of real $(p,p)$ forms has a real basis consisting of decomposable strongly positive forms;
the same forms are a complex basis of all complex $(p,p)$ forms.
\end{corollary}
\begin{proof}
The factorwise polarization at the end of Proposition \ref{cur:cones} shows that decomposable strongly positive forms span
all real $(p,p)$ forms. Successively choose elements outside the span of those already selected;
finite dimensionality produces a basis after finitely many steps.
Every complex $(p,p)$ form has a unique expression $u+\ii v$ with $u,v$ real.
Thus the real basis spans over $\C$, and separating real and imaginary parts shows that every complex linear relation is trivial.
\end{proof}

A $(p,p)$ current is called \emph{weakly positive} when it is positive in the sense of Definition~\ref{pos:currents}; it is \emph{strongly positive} when it is nonnegative on every weakly positive complementary test form. For a smooth complementary even-degree form $a$, define $(T\wedge a)(\chi)=T(\chi a)$ on scalar test functions $\chi$.

\begin{lemma}[Wedge products and positivity tests]\label{cur:wedge-test}
Wedge products of strongly positive forms are strongly positive; wedge products of weakly and strongly positive forms are weakly positive.
A current $T$ is weakly positive (respectively strongly positive) if and only if its wedge product with every smooth strongly positive
(respectively weakly positive) form of complementary degree is a positive measure. It suffices to use compactly supported test forms.
\end{lemma}
\begin{proof}
Expand strongly positive forms into nonnegative sums of decomposable positive generators.
Their wedge product is again such a sum, proving the first assertion.
If $u$ is weakly positive and $v$ strongly positive, wedge also with any complementary strongly positive form $w$.
Then $v\wedge w$ is strongly positive, so $u\wedge v\wedge w\ge0$, proving the second assertion.
For currents, fix a smooth form $a$ in the appropriate dual positivity cone. If $T$ is positive and $\chi$ is a nonnegative compactly supported function,
$\chi a$ has the same positivity, so $(T\wedge a)(\chi)=T(\chi a)\ge0$.
Lemma~\ref{pos:positive-distribution} makes a nonnegative distribution a positive measure, proving necessity.
Conversely, for a compactly supported positive test form $a$, choose a nonnegative cutoff $\chi$ equal to one near its support.
Then $T(a)=(T\wedge a)(\chi)\ge0$. A test form without compact support can be multiplied, near any fixed compact set,
by a cutoff equal to one there, so the two testing procedures are equivalent.
\end{proof}

\begin{proposition}[Two notions of positivity for rank-one forms]\label{cur:rank-one}
For every $a\in\Lambda^{p,0}V^*$, the form
$\ii^{p^2}a\wedge\bar a$ is weakly positive; it is strongly positive exactly when $a$ is decomposable, with zero included.
For $2\le p\le n-2$, the two cones are distinct.
\end{proposition}
\begin{proof}
Wedge the form with a complementary decomposable positive form. The result is the squared modulus of a complex top-degree coefficient
times a positive volume form, proving weak positivity. The case $a=0$ is immediate, so assume $a\ne0$ below.
If it is strongly positive, write $aa^*=\sum_{j=1}^N b_jb_j^*$ with every $b_j$ decomposable.
For $v\perp a$, evaluating the quadratic forms at $v$ gives
$0=\sum_j|\langle b_j,v\rangle|^2$. Hence every $b_j$ lies in the line $\C a$.
At least one $b_j$ is nonzero, so $a$ is decomposable. The converse is the definition of strong positivity.

Take
\[
a=(dz_1\wedge dz_2+dz_3\wedge dz_4)
       \wedge dz_5\wedge\cdots\wedge dz_{p+2}.
\]
If $v=\sum c_jdz_j$ satisfies $v\wedge a=0$, comparison of the basis terms containing, respectively,
$\{1,3,4\}$, $\{2,3,4\}$, $\{1,2,3\}$, and $\{1,2,4\}$
gives $c_1=\cdots=c_4=0$; terms containing an index $j>p+2$ then give $c_j=0$.
Thus $\ker L_a=\operatorname{span}(dz_5,\ldots,dz_{p+2})$, of dimension $p-2$.
Lemma \ref{cur:decomposable} proves that $a$ is not decomposable, giving the desired example.
\end{proof}

\section{Weak Topologies, Measure Coefficients, and Almost-Everywhere Positivity}

\begin{proposition}[Currents and distributional coefficients]\label{cur:distribution-coefficients}
Every degree-$k$ current in a coordinate chart has a unique expression $\sum_{|I|=k}T_I dx_I$, with distributional coefficients $T_I$.
Weak convergence of currents is equivalent to distributional convergence of their finitely many coordinate coefficients.
\end{proposition}
\begin{proof}
Let $I^c$ be the ordered complementary index, and define $s_I$ by $dx_I\wedge dx_{I^c}=s_I dV$.
Set $T_I(f)=s_I T(fdx_{I^c})$. The map $f\mapsto fdx_{I^c}$ is continuous in the test-function topology,
so $T_I$ is a distribution. Expansion of a test form and termwise pairing give the representation and its uniqueness.
Since there are only finitely many coefficients, coefficientwise convergence and convergence against every test form imply one another.
The same calculation applies to the type decomposition of complex forms.
\end{proof}

Weak convergence is defined by evaluation on test functions. The topology on the entire space of continuous functionals can nevertheless fail to be metrizable, even when sequences are useful in particular compactness arguments. The next proposition explains this distinction; it is not needed for the subsequent coefficient-measure calculations.

\begin{proposition}[Countability in the weak-star topology]\label{cur:weak-star}
Let $V$ be a Hausdorff locally convex real or complex topological vector space. Its continuous dual $V'$ carries the weak-star topology of pointwise convergence on $V$, generated by seminorms $\ell\mapsto|\ell(v)|$. A Hamel basis is an algebraic basis, so its finite linear combinations span the space.
If the weak-star topology on its continuous dual $V'$ has a countable neighborhood basis at zero, then $V$ has at most countable Hamel dimension.
Consequently, the weak-star topology on the entire dual of an infinite-dimensional Banach space is not first countable.
The last conclusion fails if ``Banach'' is replaced merely by ``normed.''
\end{proposition}
\begin{proof}
Let $U_j$ be a countable neighborhood basis. Each $U_j$ contains a basic neighborhood of the form
$W_j=\{\ell:|\ell(v_{j1})|,\ldots,|\ell(v_{jN_j})|<\epsilon_j\}$.
For each fixed $v\in V$, some $W_j$ is contained in $\{\ell:|\ell(v)|<1\}$.
If $v\notin F_j=\operatorname{span}(v_{j1},\ldots,v_{jN_j})$,
closedness of finite-dimensional subspaces and Hahn--Banach give a continuous functional with
$\ell|_{F_j}=0$ and $\ell(v)=1$. Every multiple $t\ell$ belongs to $W_j$, a contradiction.
Thus $V=\bigcup_jF_j$, so $V$ has a countable Hamel spanning set.

An infinite-dimensional Banach space cannot have a countable Hamel basis: otherwise it would be the union of finite-dimensional closed subspaces
$F_1\subset F_2\subset\cdots$, each with empty interior, contradicting Baire's theorem.
Completeness is precisely the hypothesis needed for this application of Baire.
For a counterexample in a normed space, take the finitely supported sequences $c_{00}$ with the $\ell^2$ norm.
Its continuous dual is $\ell^2$, and weak-star convergence is coordinatewise convergence,
because each element of $c_{00}$ reads only finitely many coordinates. This topology is induced by the metric
$\sum_{j\ge1}2^{-j}\min(1,|x_j-y_j|)$, and is first countable.
\end{proof}

\begin{lemma}[A finite-measure proof of Radon--Nikodym]\label{cur:radon-nikodym}
If a positive $\sigma$-finite measure $\nu$ is absolutely continuous with respect to a positive $\sigma$-finite measure $\mu$,
there is $f\ge0$ with $d\nu=f\,d\mu$, uniquely determined outside a $\mu$-null set.
A complex measure whose total variation is $\sigma$-finite and absolutely continuous with respect to $\mu$ has a corresponding complex-valued density.
\end{lemma}
\begin{proof}
First suppose both measures are finite, and put $\lambda=\mu+\nu$.
The functional $F(v)=\int v\,d\nu$ on $L^2(\lambda)$ is bounded, since
$|F(v)|\le\nu(X)^{1/2}\|v\|_{L^2(\lambda)}$.
The Hilbert-space representation argument in the proof of Lemma~\ref{pos:hilbert-existence}, with the conjugation absorbed into the density, gives $g\in L^2(\lambda)$ satisfying
$\nu(B)=\int_Bg\,d\lambda$ for every measurable $B$, by taking $v=\mathbf1_B$.
Examine the imaginary part of $g$ and the sets where its real part is negative or exceeds one. The inequalities
$0\le\nu(B)\le\lambda(B)$ show that $g$ is real and $0\le g\le1$ almost everywhere.
Hence $d\mu=(1-g)d\lambda$.
The set $\{g=1\}$ is $\mu$-null, thus $\nu$-null by absolute continuity, and consequently $\lambda$-null.
Put $f=g/(1-g)$, choosing any value on this null set. For nonnegative simple functions $v$,
$\int vf\,d\mu=\int vg\,d\lambda=\int v\,d\nu$.
Monotone convergence extends this identity to all nonnegative measurable functions and proves existence.
If two densities differed, integration over a set where a real or imaginary part of their difference has fixed sign and modulus greater than $1/m$ would contradict equality of the measures.
This also proves uniqueness.

In the general case, choose a countable measurable cover on which both measures are finite, disjointize it, and apply the finite case on each piece;
uniqueness gives compatible densities. For a complex measure, take the Jordan decompositions of its real and imaginary parts,
apply the positive-measure result to the four resulting measures, and subtract.
\end{proof}

\begin{lemma}[Testing nonnegativity of a density]\label{cur:nonnegative-density}
Let $\mu$ be an inner regular positive measure and $f$ a real locally integrable function.
If $\int f\eta\,d\mu\ge0$ for every nonnegative bounded compactly supported measurable function $\eta$,
then $f\ge0$ almost everywhere.
\end{lemma}
\begin{proof}
If $\{f<0\}$ had positive measure, some $E_m=\{f\le-1/m\}$ would have positive measure.
Inner regularity supplies a compact $K\subset E_m$ with $\mu(K)>0$.
Local integrability gives $\int_K|f|<\infty$. Taking $\eta=\mathbf1_K$
would give $\int_Kf\le-\mu(K)/m<0$, a contradiction.
\end{proof}

\begin{theorem}[A density criterion for positive currents]\label{cur:density}
Let $T$ be a real $(p,p)$ current of order zero. Locally choose a positive measure $\mu$ dominating the total variations of all coefficients,
and write $T=t\mu$. Then $T$ is weakly positive if and only if $t(x)$ lies in the weakly positive cone almost everywhere;
the analogous criterion holds for strong positivity. Each positivity cone of currents is convex and weakly closed.
\end{theorem}
\begin{proof}
Apply Lemma \ref{cur:radon-nikodym} coefficient by coefficient to obtain the density.
If $t(x)$ belongs to the required cone almost everywhere, its pairing with each test form in the dual cone integrates to a nonnegative value.

For the converse, the exceptional null set may initially depend on the test form, and this must be addressed.
In a coordinate chart, choose a countable dense set $\{a_j\}$ in the constant-coefficient dual positivity cone of complementary degree.
For every nonnegative $\chi\in C_c^\infty$, positivity gives
$\int\chi\,\langle t,a_j\rangle\,d\mu\ge0$.
This implies $\langle t,a_j\rangle\ge0$ almost everywhere. Otherwise, a set
$\{\langle t,a_j\rangle<-1/m\}$ would have positive measure. Radon regularity gives
a compact subset of positive measure, and continuous cutoffs followed by smooth approximation would produce a negative integral.
Remove the union of these countably many exceptional null sets. All the inequalities then hold simultaneously.
Continuity of the finite-dimensional pairing extends them to the whole dual cone, and Proposition \ref{cur:cones} gives the conclusion by bipolarity.
For strong positivity, interchange the two cones.

Finally, finite nonnegative linear combinations of positive currents are positive. If $T_j\to T$, then for each fixed positive test form
$a$, $T(a)=\lim_jT_j(a)\ge0$, proving weak closedness. Weak and strong positivity remain distinct throughout this argument.
\end{proof}

\begin{proposition}[Control of coefficient measures]\label{cur:coefficient-bound}
Let $T$ be a positive $(r,r)$ current. Use increasing $r$-indices $I,J$ and write $\lambda_I=\prod_{j\in I}\lambda_j$. Locally write
$T=\ii^{r^2}\sum_{|I|=|J|=r}T_{I\bar J}dz_I\wedge d\bar z_J$.
Then $T_{I\bar I}\ge0$, $T_{I\bar J}=\overline{T_{J\bar I}}$, and
\[
 \lambda_I\lambda_J|T_{I\bar J}|
 \le 2^{n-r}\!\sum_{I\cap J\subset K\subset I\cup J,\ |K|=r}
           \lambda_K^2T_{K\bar K},\qquad \lambda_j\ge0.
\]
In particular, the trace measure $\sum_I T_{I\bar I}$ dominates all coefficient measures up to a dimensional constant.
For $r=1$, positivity is equivalent to nonnegativity, for every constant vector $a$, of the measure
$\sum_{j,k}\bar a_j a_kT_{j\bar k}$. In this degree, weak and strong positivity agree.
\end{proposition}
\begin{proof}
Reality and measure coefficients were established in Theorem \ref{pos:order-zero}.
Pairing with the complementary strongly positive basis form gives $T_{I\bar I}\ge0$; reality gives conjugate symmetry.
To control off-diagonal terms, pair the one-form factors of $dz_{I^c}\wedge d\bar z_{J^c}$,
and apply the four-term polarization identity of Proposition \ref{cur:cones} to each pair.
After taking total variations, each coefficient has modulus $4^{-(n-r)}$, and every resulting term is a positive measure.
Summing over the four phases $1,-1,\ii,-\ii$ cancels mixed terms, since
\[
 \frac14\sum_{\zeta^4=1}
    \ii(a+\zeta b)\wedge\overline{a+\zeta b}
 =\ii a\wedge\bar a+\ii b\wedge\bar b.
\]
Summing factor by factor leaves only diagonal terms. The expansion has at most $2^{n-r}$ terms; repeated-index wedges vanish.
Every surviving term is some $T_{K\bar K}$ with $K$ containing $I\cap J$.
This gives the unweighted bound, with multiplicities bounded by $2^{n-r}$.
On the measure densities, make the linear cotangent basis change
$dz_j=\lambda_jdw_j$, first with all $\lambda_j>0$.
The coefficients are multiplied by $\lambda_I\lambda_J$ and $\lambda_K^2$, respectively;
all use the same dominating measure, so this is purely linear algebra.
If surplus diagonal terms were retained in the expansion, let the $\lambda_j$ outside $I\cup J$ tend to zero
to remove the terms whose indices are not contained in that union.
Finally pass to the limit to allow any $\lambda_j=0$, proving the displayed inequality.

When $r=1$, the density criterion reduces positivity to positive semidefiniteness of a Hermitian matrix.
Testing nonnegativity of the measures for a countable dense set of vectors $a$ ensures that the density matrix is positive semidefinite almost everywhere.
Diagonalization expresses such a matrix as a sum of nonnegative rank-one matrices, so the two notions of positivity coincide.
\end{proof}

\begin{example}[A derivative of a locally integrable function need not be a measure]\label{cur:derivative-not-measure}
On $\R$, set $f(x)=-2|x|^{-1/2}$. It is locally integrable, but its distributional derivative $f'$ does not have order zero.
Indeed, for $x\ne0$, the derivative equals $\operatorname{sgn}(x)|x|^{-3/2}$.
If it were a locally finite measure, its total variation on $[\epsilon,a]$ would be at least
$\int_\epsilon^a x^{-3/2}dx=2\epsilon^{-1/2}-2a^{-1/2}$.
Letting $\epsilon\downarrow0$ contradicts finite variation on compact sets.
We are not treating this nonintegrable function as a distribution on the whole line: the global derivative is defined by
$f'(\chi)=-\int f\chi'$, and the displayed function represents it only away from the origin.
\end{example}

\section{From a Distributional Complex Hessian to a Psh Function}

\begin{lemma}[Deriving local integrability]\label{cur:psh-integrable}
Let $u$ be upper semicontinuous, with values in $[-\infty,\infty)$, and suppose its restriction to every complex line
satisfies the subharmonic mean inequality. On a connected domain, if $u\not\equiv-\infty$, then $u\in L^1_{\rm loc}$.
\end{lemma}
\begin{proof}
Choose $a$ with $u(a)>-\infty$, and a polydisk $P$ centered at $a$ whose closure lies in the domain.
Upper semicontinuity and compactness give a finite upper bound $M$. Apply the circular mean inequality in each variable
and integrate over the radii to obtain
$u(a)\le |P|^{-1}\int_Pu$.
For full rigor, apply Tonelli throughout to the nonnegative function $M-u$. The successive mean inequalities
give $\int_P(M-u)\le |P|(M-u(a))<\infty$, proving integrability on $P$.

Let $G$ be the set of points near which $u$ is integrable. It is nonempty and open.
For $x\in\overline G$, choose a somewhat larger coordinate ball $B$ around $x$.
Choose $a\in G\cap B$ sufficiently close to $x$ with $u(a)>-\infty$; such points exist
because $u$ is locally integrable on $G$. A polydisk centered at $a$ can be chosen
to contain a neighborhood of $x$ while its closure remains in $B$.
The first paragraph proves integrability on that polydisk, so $x\in G$.
Thus $G$ is also closed, and connectedness makes it the whole domain.
\end{proof}

In this section write $\Delta_0=\sum_{a=1}^{2n}\partial_{x_a}^2=4\sum_{j=1}^n\partial_{z_j}\partial_{\bar z_j}$ for the scalar Euclidean Laplacian.

\begin{lemma}[Harmonic distributions are smooth]\label{cur:weyl}
A distribution $w$ on a Euclidean domain satisfying $\Delta_0w=0$ is represented by a smooth harmonic function.
\end{lemma}
\begin{proof}
Work on successively smaller balls. Radial convolution gives smooth harmonic functions $w_\epsilon=w*\rho_\epsilon$. The mean-value calculation of Lemma~\ref{pos:subharmonic-smooth}, with spheres in place of circles, shows that a harmonic function equals its convolution with any fixed nonnegative radial smooth kernel $\eta$ of integral one, provided the kernel support stays inside the domain. Hence $w_\epsilon=w_\epsilon*\eta$ on a fixed smaller ball.
As $\epsilon\downarrow0$, the right-hand side converges in every smooth norm to $w*\eta$: every derivative of $\rho_\epsilon*\eta$ converges uniformly to the corresponding derivative of $\eta$, on one fixed compact support, and distributions are continuous on such test functions. The left-hand side converges to $w$ distributionally. Thus $w=w*\eta$ there, which is smooth; its Laplacian remains zero. Such smaller balls cover the domain.
\end{proof}

\begin{theorem}[The distributional psh criterion]\label{cur:distribution-psh}
A real distribution $u$ on a complex manifold satisfies $\ii\partial\dbar u\ge0$ if and only if it has a psh representative.
This representative is unique in its distribution class. For an $L^1$ function, the conclusion is equality almost everywhere with a psh function;
arbitrary changes on a null set need not preserve plurisubharmonicity.
\end{theorem}
\begin{proof}
For the forward implication, take the trace and use Lemma~\ref{pos:positive-distribution}: $\Delta_0u$ is a positive measure. In real dimension $m=2n$, a fundamental solution of $\Delta_0$ is
\[
 \Gamma_m(x)=
 \begin{cases}
 |x|^{2-m}/((2-m)\sigma_{m-1}),&m>2,\\
 (2\pi)^{-1}\log|x|,&m=2,
 \end{cases}
\]
where $\sigma_{m-1}$ is the area of the unit sphere in $\R^m$. It is harmonic off the origin, and its outward normal derivative has integral one on every centered sphere; integration by parts therefore gives $\Delta_0\Gamma_m=\delta_0$, as in Lemma~\ref{pos:one-variable-lelong}. Convolve $\Gamma_m$ with a compactly supported cutoff of $\Delta_0u$ that equals one on a smaller ball. The resulting $v$ satisfies $\Delta_0v=\Delta_0u$ there. Local integrability of $\Gamma_m$ and Fubini's theorem give $v\in L^1_{\rm loc}$. By Lemma~\ref{cur:weyl}, $u-v$ is smooth. Thus $u$ has an $L^1_{\rm loc}$ representative.

Choose nonnegative radial smooth mollifiers $\rho_\epsilon$ and set $u_\epsilon=u*\rho_\epsilon$.
Its complex Hessian is the convolution of a positive current, so $u_\epsilon$ is smooth and psh.
Since $\Delta_0u\ge0$, spherical means are nondecreasing in the radius; using a fixed radial kernel under scaling makes
$u_\epsilon$ decrease as $\epsilon\downarrow0$.
To justify this first for distributions, apply the spherical-mean derivative calculation in the proof of Lemma \ref{pos:subharmonic-smooth}, in real dimension $2n$, to $u*\rho_\delta$,
and then let $\delta\downarrow0$ with the other smoothing scale fixed.
The decreasing limit $u^*$ is upper semicontinuous and satisfies the line mean inequalities. Since
$u_\epsilon\to u$ in $L^1_{\rm loc}$, we have $u^*=u$ almost everywhere.
The limit is not identically $-\infty$, and is the required representative.

If two psh functions agree almost everywhere, all their ball averages agree. As the balls shrink, their averages tend to their respective pointwise values,
so the functions agree everywhere. The reverse implication is Proposition~\ref{pos:regularization}, together with the weak closedness in Proposition~\ref{pos:current-calculus}.
\end{proof}

A smooth real function is \emph{pluriharmonic} if its complex Hessian $\partial\dbar u$ vanishes.

\begin{theorem}[Local potentials for general closed positive $(1,1)$ currents]\label{cur:local-potential}
Every closed positive $(1,1)$ current $T$ can be written near each point as
$T=\ii\partial\dbar u$ with $u$ psh.
Two real local potentials defining the same current differ by a smooth pluriharmonic function.
\end{theorem}
\begin{proof}
In a coordinate ball, write $T=\ii\sum T_{j\bar k}dz_j\wedge d\bar z_k$ and abbreviate $\partial_j=\partial_{z_j}$, $\bar\partial_j=\partial_{\bar z_j}$.
Closedness gives the coefficient identities
$\partial_lT_{j\bar k}=\partial_jT_{l\bar k}$ and
$\bar\partial_lT_{j\bar k}=\bar\partial_kT_{j\bar l}$.
Consequently, distributionally,
\[
 \Delta_0 T_{j\bar k}
   =4\partial_j\bar\partial_k\sum_lT_{l\bar l}.
\]
As in the fundamental-solution construction in Theorem~\ref{cur:distribution-psh}, choose a real locally integrable $v$ with
$\Delta_0 v=4\sum_lT_{l\bar l}$.
Every coefficient of $A=T-\ii\partial\dbar v$ is then a harmonic distribution, so Lemma~\ref{cur:weyl} makes $A$ a smooth closed real $(1,1)$ form.
Proposition \ref{pos:smooth-local-potential} gives a smooth real $w$ such that
$A=\ii\partial\dbar w$. Hence $T=\ii\partial\dbar(v+w)$.
Theorem \ref{cur:distribution-psh} replaces $v+w$ by its unique psh representative.

If $\partial\dbar(u_1-u_2)=0$, taking the trace gives $\Delta_0(u_1-u_2)=0$,
so the difference is smooth; its vanishing complex Hessian then makes it pluriharmonic.
On a simply connected small ball, $\partial(u_1-u_2)$ is a closed holomorphic one-form. Integration expresses the difference as the real part of a holomorphic function.
\end{proof}

\section{Extension by Zero across Complete Pluripolar Sets}

\begin{definition}
A closed set $E\subset X$ is locally complete pluripolar if, near every point, there is a psh function $v$, not identically $-\infty$,
with $E=\{v=-\infty\}$. Proper analytic subsets have this property: choose finitely many local defining functions $f_j$;
then $\log\sum|f_j|^2$ is the desired function.
By Lemma \ref{cur:psh-integrable}, the $-\infty$ locus of a nondegenerate psh function is Lebesgue-null.
\end{definition}

On a coordinate chart with $\omega_0=\ii\sum_jdz_j\wedge d\bar z_j$, the \emph{trace mass} of a positive $(p,p)$ current is the measure $\|T\|=T\wedge\omega_0^{n-p}$. Proposition~\ref{cur:coefficient-bound} makes local finiteness of this measure equivalent to local finiteness of all coefficient variations. This is the meaning of finite mass below. For the integration current of a complex submanifold, trace mass is a fixed positive multiple of its induced volume; the hypothesis therefore generalizes locally finite volume.

\begin{theorem}[El Mir's extension-by-zero theorem]\label{cur:el-mir}
Let $E$ be closed and locally complete pluripolar, and let $T$ be a closed positive $(p,p)$ current on $X\setminus E$.
Assume that $T$ has finite mass near $E$ in every relatively compact coordinate chart.
Extending all coefficient measures to assign zero mass to $E$ produces a closed positive current $\widetilde T$.
\end{theorem}
\begin{proof}
Positivity and order zero follow directly from extension of the coefficient measures and Theorem \ref{cur:density}.
It remains to prove $d\widetilde T=0$ locally. If $p=n$, this is automatic in top degree;
assume $p<n$. On a ball $B$, write $E=\{v=-\infty\}$. By shrinking the ball and adding a constant,
assume $v<-2$. Fix $K\Subset B$ and choose $\eta\in C_c^\infty(B)$ with
$0\le\eta\le1$ and $\eta=1$ near $K$. Put
$\omega_0=\ii\sum dz_j\wedge d\bar z_j$.

\emph{Step 1: cutoffs with uniformly bounded energy.}
On a slightly larger ball, convolve $v$ to obtain smooth psh functions $v_\epsilon\downarrow v$.
For a fixed large integer $j$, choose a smooth convex nondecreasing function $M$ satisfying
$M(t)=-1$ for $t\le-1$, $M(0)=0$, and strict increase on $(-1,0)$.
For instance, integrate $e^{-1/(t+1)}$ on $t>-1$, extend constantly to the left, and normalize affinely.
Thus $-1\le M(t)\le0$ for the negative arguments used here. Set $b_{j,\epsilon}=M(v_\epsilon/j)$.
It is smooth and psh, takes values in $[-1,0]$, and equals $-1$ on a neighborhood of $E\cap\operatorname{supp}\eta$
once $\epsilon$ is sufficiently small for the fixed $j$, by compactness and upper semicontinuity.
Therefore $\eta(b_{j,\epsilon}+1)^2$ has support disjoint from $E$.
Closedness of $T$ and two integrations by parts give
\begin{align*}
2\int\eta\,T\wedge\ii\partial b_{j,\epsilon}\wedge\dbar b_{j,\epsilon}
                         \wedge\omega_0^{n-p-1}
&\le \int\eta\,T\wedge\ii\partial\dbar(b_{j,\epsilon}+1)^2
                         \wedge\omega_0^{n-p-1}\\
&=\int(b_{j,\epsilon}+1)^2T\wedge\ii\partial\dbar\eta
                         \wedge\omega_0^{n-p-1}\le C.
\end{align*}
The first inequality uses $b+1\ge0$ and $\ii\partial\dbar b\ge0$.
The last constant depends on the second derivatives of $\eta$ and the finite mass of $T$ on a fixed compact set,
but not on $j,\epsilon$. The integrations by parts are legitimate precisely because $b+1$ vanishes near $E$.

\emph{Step 2: vanishing of pairings with cutoff derivatives.}
Choose a smooth $a:[-1,0]\to[0,1]$ equal to zero on $[-1,-3/4]$
and one on $[-1/2,0]$. Set $\chi_{j,\epsilon}=a(b_{j,\epsilon})$.
Put $c=-M^{-1}(-1/2)>0$. Its derivatives are supported where $v_\epsilon<-cj$, and
$\ii\partial\chi_{j,\epsilon}\wedge\dbar\chi_{j,\epsilon}
\le C_a\ii\partial b_{j,\epsilon}\wedge\dbar b_{j,\epsilon}$.
After fixing $n-p-1$ further decomposable positive factors, wedging them with $T$
produces a positive $(n-1,n-1)$ current. Its pairing in the remaining one-form variables is Hermitian positive semidefinite
and satisfies Cauchy--Schwarz, first pointwise for its density matrix and then after integration.
Thus every fixed test form $\alpha$ of degree $2n-2p-1$ supported in $K$ satisfies
\[
 |T(d\chi_{j,\epsilon}\wedge\alpha)|
 \le C_\alpha
 \left(\int_KT\wedge\ii\partial b_{j,\epsilon}\wedge\dbar b_{j,\epsilon}
                 \wedge\omega_0^{n-p-1}\right)^{1/2}
 \left(\|T\|\bigl(K\cap\{v_\epsilon<-cj\}\bigr)\right)^{1/2}.
\]
For a test form that is not a single decomposable term, apply this estimate to its finitely many coordinate basis terms and sum;
the spanning property of strongly positive forms supplies one fixed constant.
The first factor is uniformly bounded. As $\epsilon\downarrow0$, the second factor is bounded by
$\|T\|(K\cap\{v<-cj\})^{1/2}$. Then let $j\to\infty$:
these sets decrease to the empty set in $X\setminus E$, so finite mass makes their measures tend to zero.

\emph{Step 3: testing closedness.}
The form $\chi_{j,\epsilon}\alpha$ has support away from $E$, so on $X\setminus E$,
$T(d(\chi_{j,\epsilon}\alpha))=0$.
After expansion, the term containing $d\chi$ tends to zero by Step 2. In the successive limits $\epsilon\downarrow0$, $j\to\infty$, $\chi$ tends to one on $X\setminus E$,
and bounded convergence makes the other term tend to $\widetilde T(d\alpha)$.
Thus $\widetilde T(d\alpha)=0$ for every test form, which is $d\widetilde T=0$.
\end{proof}

\begin{lemma}[The dimension principle for closed currents of order zero]\label{cur:dimension-principle}
If a closed $(p,p)$ current $S$ of order zero is supported in an analytic set of complex codimension strictly greater than $p$,
then $S=0$.
\end{lemma}
\begin{proof}
First suppose the support lies in the smooth submanifold $Y=\{z_1=\cdots=z_c=0\}$.
Since the coefficients of $S$ are measures supported on $Y$, $z_jS=\bar z_jS=0$.
Differentiate these identities and use $dS=0$ to obtain
$dz_j\wedge S=d\bar z_j\wedge S=0$ ($1\le j\le c$).
Expansion in the basis wedges shows that every nonzero term must contain $dz_1,\ldots,dz_c$
and $d\bar z_1,\ldots,d\bar z_c$, which is impossible when $c>p$.

For an analytic set, apply this conclusion near each point of its regular locus.
The support of $S$ is then contained in the singular locus, which is analytic and has smaller dimension.
Repeat the argument. Dimension is a nonnegative integer, so finitely many repetitions leave the empty set, and $S=0$.
The density and dimension-drop facts used here are Proposition~\ref{alg:analytic-stratification}.
\end{proof}

\begin{theorem}[The general Poincar\'e--Lelong formula]\label{cur:poincare-lelong}
Let $f$ be holomorphic and not identically zero on any connected component. Then $\log|f|\in L^1_{\rm loc}$ and
\[
 dd^c\log|f|=[\operatorname{div}(f)].
\]
The right-hand side integrates over each irreducible hypersurface with its vanishing multiplicity and is extended by zero across the singular locus.
Consequently, \eqref{pos:section-lelong} holds for every holomorphic line bundle section nonzero on each connected component, without a normal crossing assumption.
\end{theorem}
\begin{proof}
Example~\ref{pos:log-norm} makes $\tfrac12\log(|f|^2+\epsilon)$ smooth and psh. Its decreasing limit $\log|f|$ is upper semicontinuous, and the line mean inequalities pass to the limit by monotone convergence, exactly as in Proposition~\ref{pos:psh-stability}. This limit is finite wherever $f\ne0$; Lemma~\ref{cur:psh-integrable} then supplies local integrability, so the nondegenerate psh definition applies.
Thus $T=dd^c\log|f|$ is a closed positive $(1,1)$ current.
At a smooth point of the reduced zero hypersurface, write locally $f=u z_1^m$ with $u$ nowhere zero.
The one-variable calculation in Theorem \ref{pos:poincare-lelong} gives $T=m[\{z_1=0\}]$.

Let $Z$ be the union of the singular locus of the reduced zero hypersurface and intersections of distinct irreducible branches.
Its complex codimension is at least two. The preceding calculation shows that, on $X\setminus Z$,
$T=[\operatorname{div}(f)]$. The mass on the right is bounded by the mass of $T$,
so it is finite near $Z$. Theorem \ref{cur:el-mir} gives its closed positive extension by zero $R$.
The difference $T-R$ is a closed $(1,1)$ current of order zero supported in $Z$.
Lemma \ref{cur:dimension-principle} forces it to vanish, proving the formula.
The global formula for line bundle sections again follows from $\log|s|_h=\log|f|-\varphi/2$.
\end{proof}

\begin{theorem}[Lelong's integration current of an analytic set]\label{cur:analytic-integration}
Let $A\subset X$ be a pure $d$-dimensional analytic subset, and write $A_{\rm reg}$ and $A_{\rm sing}$ for its regular and singular loci. Its regular locus has locally finite volume,
and $[A](\alpha)=\int_{A_{\rm reg}}\alpha$ defines a closed positive $(n-d,n-d)$ current.
\end{theorem}
\begin{proof}
First prove locally finite volume. By Lemma \ref{alg:finite-projections-dominate},
a coordinate neighborhood admits finitely many linear projections $P_1,\ldots,P_N$ to $\C^d$
that restrict locally to finite maps on the analytic set. On its regular locus they satisfy
$\sum_j(P_j|_{A_{\rm reg}})^*dV_{\C^d}\ge c\,dV_{A_{\rm reg}}$ with a fixed $c>0$, where $dV_{A_{\rm reg}}$ is the induced real $2d$-dimensional volume form.
After shrinking, the number of sheets of each finite map has a fixed upper bound $m_j$.
Away from the discriminant it is a finite unramified covering. Its critical locus is an analytic set of smaller dimension,
hence has zero $2d$-dimensional volume in the regular locus.
For a nonnegative compactly supported cutoff $\chi$, the real change-of-variables formula gives
\[
 \int_{A_{\rm reg}}\chi\,P_j^*dV_{\C^d}
 =\int_{\C^d}\sum_{x\in P_j^{-1}(y)\cap A_{\rm reg}}\chi(x)\,dV_y
 \le m_j\|\chi\|_\infty\operatorname{Vol}(P_j(\operatorname{supp}\chi)).
\]
Summing over the finitely many projections controls the ambient volume of $A_{\rm reg}$ and proves local finiteness.

In $X\setminus A_{\rm sing}$, the regular locus is a closed complex submanifold,
so Proposition \ref{pos:current-calculus} makes its integration current closed and positive.
The singular locus is analytic and therefore locally complete pluripolar. The volume bound just proved is exactly
the finite-mass hypothesis of Theorem \ref{cur:el-mir}.
That theorem shows that extension by zero remains closed and positive.
The extension is the current $[A]$ defined above; no extra mass is assigned to the singular locus.
\end{proof}

\begin{proposition}[Local mass growth of closed positive currents]\label{cur:mass-growth}
Let $T$ be a closed positive $(p,p)$ current in a ball, set $k=n-p$, and put $\beta=dd^c|z|^2$.
Then
$r^{-2k}\int_{|z-a|<r}T\wedge\beta^k$
is nondecreasing in $r$, for balls contained in the domain. In particular, for centers in a fixed interior compact set and sufficiently small radii,
\[
 \int_{|z-a|<r}T\wedge\beta^k\le C r^{2k}.
\]
For a $(1,1)$ current, the total variations of its coefficients therefore satisfy a local $O(r^{2n-2})$ mass bound.
\end{proposition}
\begin{proof}
If $k=0$, this is simply monotonicity of the mass of a positive measure on increasing balls.
Assume $k\ge1$, first with $T$ smooth. Translate so that $a=0$, and put $\alpha=dd^c\log|z|^2$.
The form $\alpha$ is semipositive on the punctured ball. Stokes' theorem gives
\[
 \int_{|z|=r}T\wedge d^c\log|z|^2\wedge\alpha^{k-1}
 =r^{-2k}\int_{|z|<r}T\wedge\beta^k.
\]
To check the powers of $r$, restrict to the sphere, where $d|z|^2=0$. Thus
$\alpha^{k-1}=r^{-2(k-1)}\beta^{k-1}$ and
$d^c\log|z|^2=r^{-2}d^c|z|^2$.
Applying Stokes again on the annulus between two spheres identifies the difference of the normalized masses with
$\int_{r<|z|<R}T\wedge\alpha^k\ge0$.
The already treated case $k=0$ requires no punctured-ball calculation.

In the general case, convolve $T$ on a slightly larger ball, preserving closedness and positivity.
Pass to the weak limit at radii whose spheres carry no coefficient-measure mass to obtain the same inequality.
Approximate any radius from below by such radii and use monotone convergence of measures to extend it to all radii.
Put $\mu=T\wedge\beta^k$. For a fixed larger uniform radius $R$, monotonicity of normalized mass gives
$\mu(B(a,r))\le(r/R)^{2k}\mu(B(a,R))$.
When $a$ ranges over an interior compact set, all the larger balls lie in one fixed compact set, so the constant is uniform.
Finally apply Proposition \ref{cur:coefficient-bound} to transfer the trace-measure estimate to each coefficient's total variation.
\end{proof}


\chapter{Complete Metrics, Singular Weights, and Local Multiplier Ideals}
\label{l2:chapter}

The Bochner argument on a compact manifold has an easily overlooked advantage: integration by parts produces no boundary at infinity.
We first recover this advantage using complete metrics, then show that the auxiliary metric can be changed without losing control of the relevant quantities.
After these two steps, singular weights on pseudoconvex domains can be treated by convolution and weak limits.
Finally, solving an equation with a pole in the weight turns the analytically defined multiplier ideal into a locally finitely generated holomorphic ideal.
These steps justify integration by parts, preserve the estimates, and establish coherence of the local object.

We use $F_{E,h}=\Theta(E,h)$ for Chern curvature, as in Theorem~\ref{bun:chern-connection}, and $\Lambda_\omega$ for the adjoint of wedging with $\omega$. The spaces $L^2_{p,q}(X,E)$ consist of measurable $E$-valued $(p,q)$ forms with finite squared norm for $h$ and $dV_\omega=\omega^n/n!$. For a real function $a$, $(\partial a)\mathbin{\lrcorner}$ denotes the pointwise adjoint of the operator $v\mapsto\dbar a\wedge v$.
An \emph{exhaustion} is a continuous function bounded below whose sublevel sets are relatively compact. A domain in $\C^n$ is called \emph{pseudoconvex} here when it admits a continuous psh exhaustion. A continuous function is \emph{strictly psh} if locally $u-\delta|z|^2$ is psh for some $\delta>0$.

\section{Cutoff Functions from Completeness}

\begin{lemma}[Cutoffs and a common core on a complete manifold]\label{l2:complete-core}
Let $g$ be a complete Riemannian metric. There are compactly supported Lipschitz functions $\chi_j$ such that
$0\le\chi_j\le1$, $\chi_j\to1$ on every compact set, and $|d\chi_j|_g\le C/j$.
On a complete Hermitian manifold with a smooth Hermitian vector bundle, smooth compactly supported forms
are a common graph-norm core for $\dbar$ and its Hilbert adjoint.
\end{lemma}
\begin{proof}
First, closed bounded balls for a complete Riemannian distance are compact. This can be proved without a global geodesic existence theorem.
Riemannian distance is a length distance, and each point has a small compact closed ball.
Suppose the supremum of radii of compact closed balls centered at $x_0$ were finite, equal to $R$. For every small $\epsilon>0$,
$\overline B(x_0,R-\epsilon)$ is covered by finitely many $\epsilon$-balls.
For $x\in\overline B(x_0,R)$ outside that smaller ball, follow a curve of length less than $d(x_0,x)+\epsilon$;
a point where the curve first reaches distance $R-\epsilon$ from $x_0$ lies within $3\epsilon$ of $x$.
Thus $\overline B(x_0,R)$ is totally bounded. It is closed and complete, hence compact.
Cover it by finitely many neighborhoods with compact closure. A sufficiently small enlargement of its radius is still contained in their union, contradicting the definition of $R$.
Therefore all closed balls are compact.

On a connected component, the distance function $r(x)=d(x_0,x)$ is $1$-Lipschitz.
Choose a smooth real function $a$ equal to one on $[0,1]$, zero on $[2,\infty)$,
and satisfying $0\le a\le1$. Set $\chi_j=a(r/j)$. The almost-everywhere chain rule gives the gradient bound. If the manifold is disconnected, second countability gives at most countably many components; apply this construction on the first $j$ components and put $\chi_j=0$ on the others. Each compact set meets only finitely many components, so the same conclusion holds.

For $v\in\operatorname{Dom}\dbar\cap\operatorname{Dom}\dbar^*$, the distributional product rule gives
\[
 \dbar(\chi_jv)=\chi_j\dbar v+\dbar\chi_j\wedge v,
 \qquad
 \dbar^*(\chi_jv)=\chi_j\dbar^*v-(\partial\chi_j)\mathbin{\lrcorner}v.
\]
Exterior multiplication and contraction have norms bounded by a fixed constant times $|d\chi_j||v|$.
It follows that $\chi_jv\to v$ and both operator images converge in $L^2$.
For fixed $j$, the support lies in a compact set. A partition of unity and convolution in finitely many coordinate charts,
together with the first-order commutator estimate of Lemma \ref{pos:friedrichs}, approximate the form and both derivatives simultaneously.
The Lipschitz cutoffs may also be smoothed in these charts while retaining uniform bounds on their weak first derivatives.
First let the convolution radius tend to zero, then let $j\to\infty$, to obtain the common core.
The domain identification precedes the simultaneous approximation: apply the cutoff and convolution argument to $\dbar$ alone, then extend integration by parts from that core. This identifies its Hilbert adjoint with the maximal distributional formal adjoint and justifies the product rule used above.
\end{proof}

\begin{proposition}[A complete K\"ahler metric from a smooth psh exhaustion]
\label{l2:complete-metric}
If a K\"ahler manifold $(X,\omega)$ has a smooth psh exhaustion $\rho\ge1$, then
\[
 \widehat\omega=\omega+\ii\partial\dbar\rho^2
 =\omega+2\rho\,\ii\partial\dbar\rho
       +2\ii\partial\rho\wedge\dbar\rho
\]
is a complete K\"ahler metric. Every sublevel set $X_c=\{\rho<c\}$ also admits a complete K\"ahler metric.
\end{proposition}
\begin{proof}
The displayed form is closed and dominates $\omega$, hence is K\"ahler.
Its last term controls the real differential $d\rho$, so there is a fixed constant $C$ such that every piecewise smooth curve satisfies
$|\rho(\gamma(b))-\rho(\gamma(a))|\le C\operatorname{length}_{\widehat\omega}(\gamma)$.
Thus the $\rho$-values of a $\widehat\omega$-Cauchy sequence are bounded.
The exhaustion property confines the sequence to a compact sublevel set. Smooth metrics are uniformly equivalent near a compact set,
so the sequence converges, proving completeness.
On $X_c$, use the positive psh exhaustion $\rho_c=1+1/(c-\rho)$ and apply the same formula.
\end{proof}

\begin{lemma}[Smoothing continuous strictly psh functions on Euclidean domains]\label{l2:richberg-domain}
Let $u$ be continuous and strictly psh on a Euclidean domain, and let $\epsilon$ be a positive continuous function.
There is a smooth strictly psh $v$ with $u<v<u+\epsilon$.
Consequently, a domain with a continuous psh exhaustion also has a smooth strictly psh exhaustion.
\end{lemma}
\begin{proof}
Set $u_0=u$. Choose a locally finite family of pairs of relatively compact coordinate balls $B_j'\Subset B_j$, with the inner balls covering the domain.
Strictness allows the balls to be chosen with constants $\delta_j>0$ such that
$u-\delta_j|z|^2$ is psh on $B_j$. Convolutions of continuous functions converge uniformly on compact sets,
so sufficiently small convolution scales give smooth psh approximations on slightly smaller neighborhoods,
with arbitrarily small error in value and at least half the original strict lower bound.

We describe the gluing. Choose a smooth real $b_j$ positive on $B_j'$
and negative near the inside of $\partial B_j$. Multiply it by a sufficiently small constant that the negative part of its complex Hessian
is less than half the strict lower bound still available at this step.
Convolve the current continuous strictly psh function $u_{j-1}$ and add this perturbation to obtain $w_j$,
with $w_j>u_{j-1}$ near $B_j'$ and $w_j<u_{j-1}$ near $\partial B_j$.
Replace the maximum of the two functions by a regularized maximum.
An explicit choice is
\[
 M_\tau(a,b)=\int_{\R^2}\max(a+s,b+t)\,\vartheta_\tau(s)\vartheta_\tau(t)\,dsdt,
\]
where the kernels are nonnegative, even, supported in $[-\tau,\tau]$, and have integral one. This function is convex and nondecreasing in both variables;
it equals the larger argument when $|a-b|>2\tau$.
Define $u_j=M_{\tau_j}(u_{j-1},w_j)$ on $B_j$ and $u_j=u_{j-1}$ outside $B_j$, choosing $\tau_j$ smaller than the value gaps near the inner ball and boundary. Thus the replacement leaves the original function unchanged near the boundary, is smooth near the inner ball,
and remains strictly psh. Arrange that its change in value is less than $2^{-j-2}\epsilon$
and its curvature loss less than $2^{-j-2}$ times a preassigned local positive lower bound. Shrinking the convolution scale and perturbation achieves both.

Local finiteness ensures that only finitely many replacements occur near each point.
Every point belongs to an inner ball, and the function becomes smooth near it at that step. Subsequent replacements use two functions
that are both smooth there, so smoothness persists. The total error is less than $\epsilon/2$,
and the total curvature loss is less than half the allocated positive lower bound. Each point receives a strict increase at its inner-ball step;
subsequent regularized maxima never decrease it. Hence $u<v<u+\epsilon$.
If $u_0$ is only a continuous psh exhaustion, first set $u=u_0+|z|^2$ and take $\epsilon=1$.
The resulting $v$ is still an exhaustion.
\end{proof}

\section{Estimates on Complete Manifolds and Removal of the Auxiliary Metric}

Completeness supplies cutoffs for integration by parts, whereas the final norm may be defined by a different metric. These are separate roles. In bidegree $(n,q)$, an explicit cancellation between the form norm and the volume element will let us use a complete metric temporarily and then remove it from the estimate.

\begin{theorem}[The basic $L^2$ theorem on a complete K\"ahler manifold]\label{l2:complete-solving}
Let $(X,\omega)$ be complete K\"ahler and $(E,h)$ a smooth Hermitian holomorphic vector bundle,
with $A=[\ii\Theta(E,h),\Lambda_\omega]\ge0$ on $E$-valued $(p,q)$ forms, where $q\ge1$.
If $f\in L^2_{p,q}(X,E)$, $\dbar f=0$, and
\[
 I(f)=\int_X\langle A^{-1}f,f\rangle\,dV_\omega<\infty,
\]
then there is $u\in L^2_{p,q-1}(X,E)$ satisfying $\dbar u=f$ and $\|u\|^2\le I(f)$.
The generalized inverse is understood in the sense of Proposition \ref{pos:inverse-curvature}.
If $f$ is smooth, $u$ can be chosen smooth.
\end{theorem}
\begin{proof}
For a smooth compactly supported $v$, Theorem~\ref{pos:bkn} gives
\[
 \|\dbar^*v\|^2+\|\dbar v\|^2\ge\int_X\langle Av,v\rangle dV_\omega.
\]
Lemma \ref{l2:complete-core} extends this inequality to the common domain, using lower semicontinuity for the nonnegative curvature integral.
Apply Cauchy--Schwarz pointwise on the positive eigenspaces of $A$. The component of $f$ in its zero eigenspace vanishes because
$I(f)<\infty$. Thus
\[
 |\langle f,v\rangle_{L^2}|^2\le I(f)\int_X\langle Av,v\rangle dV_\omega
 \le I(f)(\|\dbar^*v\|^2+\|\dbar v\|^2).
\]
Lemma \ref{pos:hilbert-existence} gives the solution and estimate.

Choose the solution of least norm. It is orthogonal to $\ker\dbar$, in particular to every
$\dbar\beta$ with $\beta$ smooth and compactly supported. Hence $\dbar^*u=0$ distributionally.
Thus $\Delta''u=\dbar^*f$. If $f$ is smooth, the right-hand side is smooth and Theorem~\ref{ana:elliptic-regularity} makes $u$ smooth.
This smoothness uses a Laplacian with smooth coefficients; the last step cannot simply be repeated for singular metrics later.
\end{proof}

\begin{definition}[$m$-semipositivity]\label{l2:m-positive}
Let $n=\dim_{\C}X$ and $r=\operatorname{rank}E$. In unitary tangent and bundle frames define the coefficients by
$F_{E,h}e_a=\sum_{j,k,b}c_{j\bar k a\bar b}dz_j\wedge d\bar z_k\otimes e_b$, with input index $a$ and output index $b$, as in Definition~\ref{pos:nakano-definition}.
The bundle is called $m$-semipositive if the Hermitian curvature form
$\sum c_{j\bar k a\bar b}\xi_{ja}\bar\xi_{kb}$
is nonnegative on every complex $n\times r$ matrix $\xi$ of rank at most $m$.
Rank one gives Griffiths semipositivity; for $m\ge\min(n,r)$ this is Nakano semipositivity.
Strict positivity means strict inequality for every nonzero matrix of the indicated rank.
\end{definition}

\begin{lemma}[Metric monotonicity in bidegree $(n,q)$]\label{l2:metric-monotonicity}
Let $\omega'\ge\omega$ be Hermitian metrics, and suppose the curvature of $E$ is $m$-semipositive,
with $m\ge\min(n-q+1,r)$. Put $A_\omega=[\ii\Theta(E),\Lambda_\omega]$.
Then $A_\omega\ge0$, and
\[
 |f|_{\omega'}^2dV_{\omega'}\le|f|_\omega^2dV_\omega,
 \qquad
 \langle A_{\omega'}^{-1}f,f\rangle_{\omega'}dV_{\omega'}
 \le\langle A_\omega^{-1}f,f\rangle_\omega dV_\omega
\]
for $(n,q)$ forms. The first inequality holds for all $(n,k)$ forms as well.
\end{lemma}
\begin{proof}
At the point under consideration, make $\omega$ the identity and simultaneously write
$\omega'=\ii\sum_j\lambda_jdz_j\wedge d\bar z_j$, with $\lambda_j\ge1$.
Use increasing indices $J,K$, put $\lambda_J=\prod_{j\in J}\lambda_j$ (with $\lambda_\varnothing=1$), and extend coefficients antisymmetrically as in Proposition~\ref{pos:nakano-operator}. If $f=\sum_{|J|=q}f_J dz_1\wedge\cdots\wedge dz_n\wedge d\bar z_J$,
the metric factor in top holomorphic degree cancels the volume factor, leaving
$|f|_{\omega'}^2dV_{\omega'}=\sum_J|f_J|^2\lambda_J^{-1}dV_\omega$.
This proves the first inequality.

Write the test-form coefficients as $v_J=\lambda_Jw_J$. Write $Q_\omega(v)=\langle v,A_\omega v\rangle_\omega dV_\omega$ for its quadratic-form density. Proposition~\ref{pos:nakano-operator} gives
\[
 Q_{\omega'}(v)=
 \sum_{|K|=q-1}\lambda_K
 \sum_{j,k,a,b}c_{j\bar k a\bar b}w_{jK,a}\overline{w_{kK,b}}\,dV_\omega.
\]
Antisymmetry gives $w_{jK,a}=0$ for $j\in K$, so each inner matrix has at most
$n-q+1$ nonzero rows and rank at most $\min(n-q+1,r)$.
Every inner quadratic form is nonnegative. This proves $A_\omega\ge0$ and shows that replacing each $\lambda_K$ by one
can only decrease the form. On the other hand, the pairing times volume is
$\langle f,v\rangle_{\omega'}dV_{\omega'}=\sum_J\langle f_J,w_J\rangle dV_\omega$,
which no longer involves $\lambda$.
The generalized-inverse variational formula
\[
 \langle A^{-1}f,f\rangle
   =\sup_v\{2\operatorname{Re}\langle f,v\rangle-\langle Av,v\rangle\}
\]
follows by completing the square in each eigendirection, including the case of value $+\infty$.
Substitution of the two quadratic forms into this formula gives the second inequality.
\end{proof}

\begin{theorem}[The metric used in the estimate need not be complete]\label{l2:incomplete-metric}
Suppose $X$ admits some complete K\"ahler metric, and let $\omega$ be any K\"ahler metric.
Let $E$ be smooth and $m$-semipositive, with $m\ge\min(n-q+1,r)$.
Then Theorem \ref{l2:complete-solving} remains valid in bidegree $(n,q)$, with every norm in the estimate computed using $\omega$.
\end{theorem}
\begin{proof}
Set $\omega_\delta=\omega+\delta\widehat\omega$, where $\widehat\omega$ is complete.
This is complete: its distance dominates $\sqrt\delta$ times the complete distance, and local equivalence of the smooth metrics gives convergence in $\omega_\delta$ of each Cauchy sequence.
By Lemma~\ref{l2:metric-monotonicity}, neither the $L^2$ norm nor the inverse-curvature energy of $f$ increases.
Theorem \ref{l2:complete-solving} gives $\dbar u_\delta=f$ and
$\int|u_\delta|_{\omega_\delta}^2dV_{\omega_\delta}\le I(f)$.

Fix a compact set $K$. On $K$, $\omega_\delta\le(1+C_K\delta)\omega$.
The explicit formula for the $(n,q-1)$ norm gives
\[
 \int_K|u_\delta|_\omega^2dV_\omega
 \le(1+C_K\delta)^{q-1}I(f).
\]
Successively take weakly convergent subsequences on a countable compact exhaustion, then a diagonal subsequence, to obtain $u\in L^2_{\rm loc}$.
Distributional differentiation is continuous under this weak convergence, so $\dbar u=f$.
Weak lower semicontinuity of the Hilbert norm first gives the estimate on each $K$;
exhausting $X$ and applying monotone convergence gives $\|u\|_\omega^2\le I(f)$ globally.
For smooth $f$, choose the minimal-norm solution in the original metric and use the elliptic argument in Theorem~\ref{l2:complete-solving} to obtain a smooth solution.
\end{proof}

\section{General Singular Weights on Pseudoconvex Domains}

\begin{theorem}[H\"ormander's theorem for singular psh weights on a domain]\label{l2:domain-singular}
Let $\Omega\subset\C^n$ admit a continuous psh exhaustion, and let $\varphi$ be psh.
Let $B$ be a continuous positive semidefinite Hermitian matrix field satisfying
$\ii\partial\dbar\varphi\ge\ii\sum B_{j\bar k}dz_j\wedge d\bar z_k$ as currents.
Use Euclidean Lebesgue measure $dV$ and the orthonormal basis $d\bar z_J$ of Proposition~\ref{pos:euclidean-identity}. For $1\le q\le n$, let $B_q$ be the induced curvature operator on $(0,q)$ forms from that identity; in an eigenframe its eigenvalues are
$\lambda_{j_1}+\cdots+\lambda_{j_q}$.
If a $(0,q)$ form $f$ satisfies $\dbar f=0$ and
\[
 I=\int_\Omega\langle B_q^{-1}f,f\rangle e^{-\varphi}dV<\infty,
\]
then there is a $(0,q-1)$ form $u$ satisfying $\dbar u=f$ and
$\int_\Omega|u|^2e^{-\varphi}dV\le I$.
The datum $f$ may initially be only locally integrable; finite energy ensures the required local weighted $L^2$ integrability.
\end{theorem}
\begin{proof}
Use Lemma \ref{l2:richberg-domain} to choose a smooth strictly psh exhaustion $\rho$, and put
$\Omega_j=\{\rho<c_j\}\Subset\Omega$, where $c_j\to\infty$.
On Euclidean space, exterior multiplication by the fixed form $dz_1\wedge\cdots\wedge dz_n$
identifies $(0,q)$ forms isometrically with $(n,q)$ forms valued in the trivial line bundle, commuting with $\dbar$ up to a fixed sign.

First assume $B$ is positive definite. On a neighborhood of $\overline\Omega_j$, convolve to obtain
$\varphi_\epsilon=\varphi*\rho_\epsilon\searrow\varphi$.
Convolution preserves the current inequality, and $B*\rho_\epsilon\to B$ uniformly.
Thus, for every small $\delta>0$ and sufficiently small $\epsilon$,
$\ii\partial\dbar\varphi_\epsilon\ge\ii\sum (B-\delta I)_{j\bar k}dz_j\wedge d\bar z_k$ on $\Omega_j$.
Choose $\delta$ smaller than half the least eigenvalue on this compact set.
Theorem \ref{l2:incomplete-metric}, using the complete auxiliary metrics from Proposition \ref{l2:complete-metric}, gives a solution $u_{j,\epsilon,\delta}$ on $\Omega_j$,
whose squared weighted norm is bounded by the energy of $f$ computed with $(B-\delta I)_q^{-1}$.
Since $e^{-\varphi_\epsilon}\le e^{-\varphi}$, this is bounded by the corresponding integral with the original weight.

Fix $j,\delta$, and first let $\epsilon\downarrow0$.
The weights $\varphi_\epsilon$ have a common finite upper bound on compact sets, so the solutions are locally bounded in ordinary $L^2$.
Take a diagonal weak limit $u_{j,\delta}$. To retain the singular-weight estimate rigorously, fix an earlier scale $\epsilon_0$;
for all $\epsilon<\epsilon_0$,
$e^{-\varphi_{\epsilon_0}}\le e^{-\varphi_\epsilon}$.
Use weak lower semicontinuity with this fixed smooth weight, then let $\epsilon_0\downarrow0$.
Monotone convergence gives the estimate with the singular weight. Next let $\delta\downarrow0$;
because $(B-\delta I)_q^{-1}\le2B_q^{-1}$ and converges pointwise, dominated convergence restores constant one.
Finally let $j\to\infty$ and take local weak limits in the same way to obtain a solution on all of $\Omega$.

If $B$ is only semipositive, replace the weight by $\varphi+t|z|^2$ and the lower bound by $B+tI$.
The variational formula for the generalized inverse gives $(B+tI)_q^{-1}\le B_q^{-1}$.
The preceding construction therefore gives solutions with estimates bounded by $I$.
As $t\downarrow0$, take weak limits on compact sets. Fix an earlier $t_0$, use lower semicontinuity, and then use monotone convergence
to recover the original-weight estimate. Every limit preserves $\dbar u=f$.
Local weighted $L^2$ integrability follows from the upper bound for $B$ on compact sets:
$|f|^2\le C_K\langle B_q^{-1}f,f\rangle$ on the finite-energy subspace.
\end{proof}

\begin{lemma}[$L^2$ removability across analytic sets]\label{l2:removable}
Let $Z$ be a proper analytic subset. If $f$ is holomorphic on $\Omega\setminus Z$ and
$f\in L^2_{\rm loc}(\Omega)$, then $f$ extends uniquely holomorphically to $\Omega$.
\end{lemma}
\begin{proof}
Locally choose a nonzero holomorphic function $g$ with $Z\subset\{g=0\}$.
A generic linear coordinate change makes $g$ regular in each coordinate direction,
meaning that its lowest nonzero homogeneous part does not vanish on any coordinate axis direction.
Theorem~\ref{alg:weierstrass-preparation} shows that, for almost every fixed choice of the other coordinates, a small disk in that direction
meets $\{g=0\}$ in finitely many points.
Fubini also gives square integrability of $f$ on almost every such disk.
In a one-variable Laurent expansion, any negative-degree term makes the radial integral of the circular squared mean diverge;
therefore all these isolated singularities are removable.
Integration by parts in each coordinate direction gives $\partial f/\partial\bar z_j=0$ distributionally.
Hence $\Delta_0 f=4\sum_j\partial_{z_j}\partial_{\bar z_j}f=0$. Lemma~\ref{cur:weyl} makes $f$ smooth,
and Theorem~\ref{bas:cauchy} makes it holomorphic. Equality on the dense complement proves uniqueness.
\end{proof}

\section{Why Multiplier Ideals Are Coherent}

\begin{lemma}[The ideal sheaf generated by a family of holomorphic functions]\label{l2:holomorphic-family}
The ideal sheaf generated by any family of holomorphic functions on a fixed open set is coherent.
\end{lemma}
\begin{proof}
This is a direct application of Corollary~\ref{alg:holomorphic-family} from local algebra.
Its common-neighborhood argument deserves emphasis: Noetherianity first selects finitely many functions generating the ideal
at the stalk at $x$. Apply the uniform vanishing lemma \ref{alg:uniform-germ}
to the fixed coherent quotient $\cO/(f_1,\ldots,f_N)$.
The image of each other generator generates a coherent subsheaf with zero stalk at $x$;
the lemma makes all these subsheaves vanish on one common neighborhood.
The finitely many selected functions therefore generate the whole ideal there, and Theorem~\ref{alg:oka} supplies the coherent relation sheaf.
Shrinking a neighborhood separately for each function would not establish this common-neighborhood conclusion.
\end{proof}

\begin{theorem}[Coherence of multiplier ideals]\label{l2:multiplier-coherent}
For every psh function $\varphi$, the multiplier ideal $\cI(\varphi)$ is coherent.
The same holds when $\varphi$ is locally the sum of a psh function and a smooth function.
\end{theorem}
\begin{proof}
This is local, so work in a larger coordinate ball $B$.
Put
\[
 \mathcal H=\{g\in\cO(B):\int_B|g|^2e^{-\varphi}dV<\infty\},
\]
and let $\cJ$ be the ideal sheaf generated by all $g\in\mathcal H$.
Lemma \ref{l2:holomorphic-family} makes $\cJ$ coherent.
Clearly $\cJ\subset\cI(\varphi)$. We prove the reverse inclusion at every stalk.

Fix $x\in B$, $f\in\cI(\varphi)_x$, and an integer $k\ge1$.
Represent $f$ on a small ball about $x$ and choose a smooth compactly supported cutoff $\chi$ equal to one on a smaller ball.
The form $v=\dbar(\chi f)$ is supported on an annulus away from $x$, and
$|v|^2e^{-\varphi}$ is integrable.
On the whole ball $B$, use the weight
\[
 \Phi(z)=\varphi(z)+(n+k)\log|z-x|^2+|z|^2
\]
in Theorem \ref{l2:domain-singular}. Its curvature is at least $\omega_0=\ii\sum_jdz_j\wedge d\bar z_j$,
and the support of $v$ avoids the pole, so its inverse-curvature energy is finite.
We obtain $\dbar u=v$ and $\int_B|u|^2e^{-\Phi}<\infty$.
Thus $g=\chi f-u$ is holomorphic on $B$, by the distributional Cauchy--Riemann argument of Lemma~\ref{l2:removable}.
Since $B$ is bounded, $e^{-\Phi}$ is bounded below by a positive constant times $e^{-\varphi}$,
so $u$ is square integrable with the original weight, and $g\in\mathcal H$.

Near $x$, $u$ is holomorphic. The psh weight is bounded above on a small closed ball, so
$\int|u|^2|z-x|^{-2(n+k)}<\infty$.
Use the radial calculation of Proposition~\ref{pos:radial-ideal}: expand its Taylor series and take squared spherical means. If the lowest nonzero homogeneous term has degree $d$,
its radial integral contains $\int_0^\delta r^{2d-2(n+k)+2n-1}dr$;
integrability requires $d>k$. Hence $u\in\mathfrak m_x^{k+1}$, also trivially when $u$ is zero.
We obtain $f\in\cJ_x+\mathfrak m_x^{k+1}$ for every $k$.
The Krull closure statement \ref{alg:krull-closure} gives
$f\in\bigcap_k(\cJ_x+\mathfrak m_x^k)=\cJ_x$.
The two sheaves therefore coincide, proving coherence.
Proposition~\ref{pos:multiplier-invariance} shows that adding a smooth function does not change the multiplier ideal on a small ball because that function is bounded there. This proves the last assertion.
\end{proof}

\begin{proposition}[The local weighted Dolbeault resolution]\label{l2:local-resolution}
Let $L$ be a holomorphic line bundle whose singular metric $h=e^{-\varphi}$ has local weights equal to a psh function plus a smooth function.
Let $\mathcal L^{n,q}_{(2)}(L,h)$ be the sheaf of forms locally square integrable with respect to $h$
whose distributional $\dbar$ is also locally square integrable with respect to $h$.
Then
\[
 0\longrightarrow\cO(K_X\otimes L)\otimes\cI(h)
 \longrightarrow\mathcal L^{n,0}_{(2)}(L,h)
 \xrightarrow{\dbar}\mathcal L^{n,1}_{(2)}(L,h)
 \xrightarrow{\dbar}\cdots
\]
is an exact fine resolution.
\end{proposition}
\begin{proof}
The degree-zero kernel consists of holomorphic sections by the distributional Cauchy--Riemann argument in Lemma~\ref{l2:removable};
the weighted integrability condition is precisely the multiplier-ideal condition.
In positive degree, shrink to a coordinate ball and add $C|z|^2$ to make the curvature strictly positive.
The new and old weights differ by a bounded function there, so they define the same local $L^2$ spaces.
Theorem \ref{l2:domain-singular} gives a local preimage of every closed form.
If $\dbar u=f$ and both $u,f$ are locally square integrable, $u$ belongs exactly to the graph-domain sheaf defined above.
Finally, multiplication by a smooth function $a$ gives
$\dbar(au)=\dbar a\wedge u+a\dbar u$, which remains integrable on compact sets.
The partitions of unity from Theorem~\ref{bas:partition} act on these sheaves; Proposition~\ref{sh:smooth-fine} therefore proves that they are fine.
\end{proof}

\Needspace{11\baselineskip}
\begin{remark}[Singular metrics on general vector bundles need a separate definition]
For a line bundle, singular curvature is the distributional complex Hessian of a locally integrable scalar weight. A measurable matrix metric on a vector bundle does not have this scalar description: its curvature requires an appropriate Nakano-type $L^2$ positivity condition or a controlled approximation class. The $n$ eigenvalues of a scalar complex Hessian cannot describe its bundle indices.
The domain theorem here uses scalar singular weights; the smooth vector bundle statements consistently use the curvature operator.
Moreover, a minimal-norm solution for a singular weight need not be smooth even if $f$ is smooth,
because its weighted adjoint is no longer a differential operator with smooth coefficients.
\end{remark}


\chapter{General Skoda Division and the Brian\c{c}on--Skoda Theorem}
\label{div:chapter}

Division by one holomorphic function asks whether a quotient extends across its zero set. For several generators $g_1,\ldots,g_r$, the corresponding problem is to solve $f=\sum g_i h_i$ with every $h_i$ holomorphic. Away from their common zero set, $s_i=f\bar g_i/|g|^2$ already solves the algebraic equation. We will correct this smooth solution within the kernel of the map $h\mapsto\sum g_i h_i$, preserving the equation while making $h$ holomorphic.
A Cauchy--Schwarz inequality with a rank factor provides the curvature estimate. The resulting division theorem leads to an integral-closure inclusion in the local ring. Chapter~\ref{bs:chapter} gives a complementary proof for monomial ideals using convex geometry.

\section{A Matrix Inequality and Curvature of the Kernel Bundle}

\begin{lemma}[Skoda's matrix inequality]\label{div:skoda-inequality}
Let $B=(b_{j\alpha})$ be an $n\times s$ complex matrix of rank $\ell$, let $b_\alpha$ denote its columns, and put $Q=BB^*=(Q_{j\bar k})$. For every $n\times s$ matrix $\Xi=(\xi_{j\alpha})$,
\[
 \left|\sum_{j,\alpha}b_{j\alpha}\xi_{j\alpha}\right|^2
 \le\ell\sum_{j,k,\alpha}Q_{j\bar k}\xi_{j\alpha}\bar\xi_{k\alpha}.
\]
The contraction on the left is bilinear, without a conjugate on $\xi$; the right side uses exactly the displayed order of conjugations. Let $Q^+$ denote the Moore--Penrose inverse, which inverts the positive eigenvalues and is zero on $\ker Q$. Then
$\sum_\alpha\langle Q^+b_\alpha,b_\alpha\rangle=\ell$.
This use of $Q^+$ is compatible with Proposition~\ref{pos:inverse-curvature}, because every $b_\alpha$ lies in $\operatorname{Ran}Q$; it does not assign finite inverse energy to a nonzero kernel component.
\end{lemma}
\begin{proof}
Put $D=B^T\Xi$. The left-hand side is $|\tr D|^2$, and the sum on the right, before the factor $\ell$, equals $\sum_{\alpha,\beta}|D_{\beta\alpha}|^2$. Under unitary changes $B'=UBV$ and $\Xi'=\bar U\Xi\bar V$, one has $D'=V^TD\bar V$. Both its trace and the sum of squared moduli of its entries are unchanged. A singular value decomposition therefore reduces to $b_{11}=\sigma_1,\ldots,b_{\ell\ell}=\sigma_\ell>0$, with all other entries zero.
The left-hand side becomes $|\sum_{j=1}^\ell\sigma_j\xi_{jj}|^2$, which Cauchy--Schwarz bounds by $\ell\sum_{j=1}^\ell\sigma_j^2|\xi_{jj}|^2$. Adding the remaining nonnegative terms gives the stated right-hand side.
Finally,
\[
 \sum_\alpha\langle Q^+b_\alpha,b_\alpha\rangle
 =\tr(Q^+BB^*)=\tr(Q^+Q)=\ell,
\]
because $Q^+Q$ is the orthogonal projection onto the $\ell$-dimensional range of $Q$. If $\ell=0$, both assertions are immediate.
\end{proof}

\begin{lemma}[Negative kernel-bundle curvature and compensation by a logarithmic norm]\label{div:kernel-curvature}
Let $g=(g_1,\ldots,g_r)$ be holomorphic on a domain $\Omega\subset\C^n$ and have no common zero. Write $\underline\C^r=\Omega\times\C^r$ for the trivial bundle, let $g$ act by $v\mapsto\sum_i g_i v_i$, and give
$F=\ker(g:\underline\C^r\to\underline\C)$ the induced metric. Put
$Q=\ii\partial\dbar\log|g|^2$,$m=\min(n,r-1)$.
After multiplying the metric on $F$ by the scalar weight $|g|^{-2(m+\epsilon)}e^{-\varphi}$,
its Nakano curvature is bounded below by $\epsilon Q\otimes\operatorname{Id}_F$, where
$\epsilon>0$ and $\varphi$ is smooth and psh.
\end{lemma}
\begin{proof}
Fix a point and use a constant unitary change to arrange $g=(a,0,\ldots,0)$ there, with $a>0$. Let $e_1,\ldots,e_r$ denote the resulting standard unitary frame of the trivial bundle.
Near that point, the kernel bundle has holomorphic frame
$e_\alpha-(g_\alpha/g_1)e_1$ for $2\le\alpha\le r$.
At the point, its metric matrix is the identity and its first derivatives vanish. The Chern formula \eqref{bun:chern-formula}, with input/output coefficient indices as in Definition~\ref{pos:nakano-definition}, gives
\[
 \Theta(F)_{j\bar k\alpha\bar\beta}
   =-\frac{\partial_{z_j}g_\alpha\overline{\partial_{z_k}g_\beta}}{a^2}.
\]
Equivalently, the connection matrix for the orthogonal splitting shows that kernel-bundle curvature is trivial-bundle curvature minus the square of the second fundamental form.
Put $b_{j\alpha}=\partial_{z_j}g_\alpha/a$ for $2\le\alpha\le r$, and let $B=(b_{j\alpha})$. Write $Q=\ii\sum_{j,k}Q_{j\bar k}dz_j\wedge d\bar z_k$ and denote its coefficient matrix by $\mathsf Q=(Q_{j\bar k})$. Direct differentiation of $\log\sum|g_i|^2$ gives
$Q_{j\bar k}=\sum_{\alpha=2}^rb_{j\alpha}\bar b_{k\alpha}$, so $\mathsf Q=BB^*$. The terms involving $\partial g_1$ cancel against the subtraction term from differentiating the logarithm.
In the convention of Definition~\ref{pos:nakano-definition}, the kernel bundle's Nakano quadratic form is therefore $-|\sum b_{j\alpha}\xi_{j\alpha}|^2$.
Lemma \ref{div:skoda-inequality} bounds it below by
$-m\sum Q_{j\bar k}\xi_{j\alpha}\bar\xi_{k\alpha}$.
Multiplication by the scalar weight adds to the curvature
$(m+\epsilon)Q+\ii\partial\dbar\varphi$, tensored with the identity.
The last Hessian is semipositive, and subtraction of the first two contributions leaves $\epsilon Q$, as required.
\end{proof}

\begin{lemma}[A complete auxiliary metric on the complement of the zero set]\label{div:complement-complete}
Let $\Omega_0\Subset\Omega\subset\C^n$ be domains, with $\Omega_0$ admitting a smooth psh exhaustion.
Let $g$ be a holomorphic tuple on $\Omega$, not identically zero, and put $Z=\{g=0\}$.
Then $\Omega_0\setminus Z$ admits a complete K\"ahler metric,
which can be chosen to dominate $Q=\ii\partial\dbar\log|g|^2$.
\end{lemma}
\begin{proof}
First give $\Omega_0$ a complete metric $\omega_c$ using Proposition \ref{l2:complete-metric}.
Since $g$ is bounded on $\overline\Omega_0$, subtract a constant to arrange
$v=\log|g|^2-C<-1$.
Outside the zero set, direct differentiation gives
\[
 \ii\partial\dbar\bigl(-\sqrt{-v}\bigr)
 =\frac{\ii\partial\dbar v}{2\sqrt{-v}}
   +\frac{\ii\partial v\wedge\dbar v}{4(-v)^{3/2}}
 \ge4\ii\partial(-v)^{1/4}\wedge\dbar(-v)^{1/4}.
\]
Thus $\omega_c+Q+\ii\partial\dbar(-\sqrt{-v})$ is K\"ahler.
A curve approaching $\partial\Omega_0$ has infinite length because of $\omega_c$;
one approaching $Z$ has infinite length because $(-v)^{1/4}\to\infty$ and the last gradient bound controls its variation.
Equivalently, a Cauchy sequence can approach neither the boundary nor the zero set,
so it stays in a compact subset and converges. The new metric is complete and clearly dominates $Q$.
\end{proof}

\section{General Division with a Precise Estimate}

\begin{theorem}[Skoda division]\label{div:skoda}
Let $\Omega\subset\C^n$ be a pseudoconvex domain, let $\varphi$ be psh,
and let $g_1,\ldots,g_r\in\cO(\Omega)$, not all identically zero.
Put $|g|^2=\sum_i|g_i|^2$ and $m=\min(n,r-1)$.
If $f\in\cO(\Omega)$ satisfies
\[
 I=\int_\Omega |f|^2|g|^{-2(m+1+\epsilon)}e^{-\varphi}dV<\infty,
 \qquad\epsilon>0,
\]
then there are $h_1,\ldots,h_r\in\cO(\Omega)$ such that
\[
 f=\sum_i g_i h_i,
 \qquad
 \int_\Omega\sum_i|h_i|^2|g|^{-2(m+\epsilon)}e^{-\varphi}dV
       \le\left(1+\frac m\epsilon\right)I.
\]
\end{theorem}
\begin{proof}
If $r=1$, set $h_1=f/g_1$ on $\{g_1\ne0\}$.
The required estimate is exactly the hypothesis. The weight has a positive lower bound on compact sets, so $h_1$ is locally square integrable.
Lemma \ref{l2:removable} extends it across the zero set, proving the conclusion with constant one.
Assume henceforth $r\ge2$.

\emph{Step 1: a smooth weight on a relatively compact exhaustion domain.}
Fix $\Omega_0\Subset\Omega$ with a smooth psh exhaustion, and temporarily assume that $\varphi$ is smooth near $\overline\Omega_0$. All integrals in this step are over $\Omega_0\setminus Z$, where $Z=\{g=0\}$, and their input integral is bounded by the original $I$. Such subdomains are supplied by Lemma~\ref{l2:richberg-domain}. Off $Z$, take the smooth minimal-norm lifting
$s=f\bar g/|g|^2$, which satisfies $g\cdot s=f$.
It need not be holomorphic; its defect is $\eta=\dbar s$.
Since $g$ is holomorphic, $g\cdot\eta=\dbar(g\cdot s)=0$,
so $\eta$ is a closed $(0,1)$ form valued in the kernel bundle $F$.
Equip $F$ with the weight of Lemma \ref{div:kernel-curvature}.

At the point and in the frame used in that lemma, $s=(f/a,0,\ldots,0)$ and
$\eta_{\bar j,\alpha}=(f/a)\overline{b_{j\alpha}}$.
Compute first in the Euclidean metric and factor out the common scalar bundle weight. Let $A$ be the curvature operator of the weighted kernel bundle on $(n,1)$ forms, as in Proposition~\ref{pos:nakano-operator}. The scalar operator associated to $Q$ acts on $(0,1)$ coefficient columns by $\mathsf Q^T$, because its exterior-algebra expression is $\sum_{j,k}Q_{j\bar k}\bar e_k\bar\iota_j$. Lemma~\ref{div:kernel-curvature} bounds $A$ below by $\epsilon$ times this operator, tensored with $\operatorname{Id}_F$. Since the columns of $\eta$ are $(f/a)\bar b_\alpha$, the variational formula in Lemma~\ref{l2:metric-monotonicity} gives
\[
 \langle A^{-1}\eta,\eta\rangle
 \le\frac1\epsilon\frac{|f|^2}{|g|^2}
       \sum_\alpha\langle \mathsf Q^+b_\alpha,b_\alpha\rangle
 \le\frac m\epsilon\frac{|f|^2}{|g|^2}.
\]
Here every $b_\alpha$ lies in $\operatorname{Ran}\mathsf Q$, so the inverse energy is finite even when $\mathsf Q$ is singular; the last bound is Lemma~\ref{div:skoda-inequality}. The exterior operators $\bar e_k,\bar\iota_j$ are the multiplication and contraction operators of Proposition~\ref{pos:kahler-identities}.
In bidegree $(n,1)$, inverse-curvature energy times volume is independent of the auxiliary K\"ahler metric;
this also follows directly from the empty-index case $K=\varnothing$ in Lemma \ref{l2:metric-monotonicity}.
Choose the complete metric in Lemma \ref{div:complement-complete}, dominating $Q$.
In this metric the squared covector norm of the columns $\bar b_\alpha$ has sum at most $\operatorname{rank}\mathsf Q\le m$, since $\omega\ge Q$. Thus the ordinary weighted $L^2$ norm of $\eta$ is bounded by the integral of $m|f|^2/|g|^2$ with the stated weight,
so every hypothesis of Theorem \ref{l2:complete-solving} is satisfied.
It gives an $F$-valued form $u$ with $\dbar u=\eta$ and
\[
 \int|u|^2|g|^{-2(m+\epsilon)}e^{-\varphi}
       \le\frac m\epsilon I.
\]
The solution is an $(n,0)$ form, whose squared norm times volume is also independent of the auxiliary metric, so it may again be regarded as an $F$-valued function.

Set $h=s-u$. It is holomorphic and $g\cdot h=f$.
The section $s$ is pointwise orthogonal to $F$, whereas $u$ is $F$-valued. Hence
$|h|^2=|s|^2+|u|^2=|f|^2/|g|^2+|u|^2$.
Integration gives the precise constant $1+m/\epsilon$ by orthogonality.
The weight is locally bounded below by a positive constant, so $h$ belongs to ordinary local $L^2$.
Lemma \ref{l2:removable}, applied componentwise, extends $h$ holomorphically across $Z$.

\emph{Step 2: removal of smoothness and relative compactness.}
Choose relatively compact pseudoconvex exhaustion domains $\Omega_j$.
On a neighborhood of $\overline\Omega_j$, convolve $\varphi$ to obtain
$\varphi_\delta\searrow\varphi$. Step 1 gives $h_{j,\delta}$,
with right-hand integral at most $I$ because $e^{-\varphi_\delta}\le e^{-\varphi}$.
These holomorphic functions have a uniform ordinary $L^2$ bound on each compact set.
Lemma~\ref{emb:cauchy-l2} gives uniform bounds for values and derivatives on smaller compact sets,
so a subsequence converges locally uniformly. The limit is holomorphic and satisfies the same division identity.
Fatou's lemma on the complement of the null set $Z$ gives the estimate with the original singular weight.
A further diagonal subsequence as $j\to\infty$ produces $h$ on all of $\Omega$.
Neither limiting process increases the constant, completing the proof.
\end{proof}

\section{Growth Ideals and Reduction of Generators}

\begin{definition}
Let $I=(g_1,\ldots,g_r)\ne0$ be an ideal of $\cO_{\C^n,0}$.
For a real number $k\ge0$, define the following sets of holomorphic germs $f\in\cO_{\C^n,0}$; $L^1_{\rm loc,0}$ means integrable on some neighborhood of the origin:
\[
 \begin{aligned}
 I^{(k)}&=\{f:|f|\le C|g|^k\text{ on some neighborhood of the origin}\},\\
 J^{(k)}&=\{f:|f|^2|g|^{-2(k+\epsilon)}\in L^1_{\rm loc,0}
                      \text{ for some }\epsilon>0\}.
 \end{aligned}
\]
Parenthesized superscripts do not denote ordinary powers; $I^k$, for integral $k\ge0$, always denotes the usual product ideal.
These sets are ideals: multiplication by a holomorphic function adds only a locally bounded factor; finite sums are handled by the triangle inequality,
taking the smallest admissible positive $\epsilon$ in the integrability condition.
Two finite generating systems are holomorphic matrix combinations of each other.
After shrinking, the matrices are bounded, so the two norms are comparable and the definitions are independent of the generators.
\end{definition}

\begin{lemma}[Integrability of small negative powers]\label{div:small-power}
If $0\ne g\in\cO_{\C^n,0}$ has vanishing order $d\ge1$, then
$|g|^{-2c}$ is locally integrable for $0\le c<1/d$. If $g(0)\ne0$, it is locally integrable for every finite $c\ge0$.
Consequently, $I^{(k)}\subset J^{(k)}$; for integers $k\ge1$, also $I^k\subset I^{(k)}$; and
$I^{(a)}I^{(b)}\subset I^{(a+b)}$, $I^{(a)}J^{(b)}\subset J^{(a+b)}$.
\end{lemma}
\begin{proof}
Choose the last coordinate direction so that the lowest nonzero homogeneous term does not vanish on it.
Weierstrass preparation, Theorem \ref{alg:weierstrass-preparation}, gives
$g=uP$, where $u$ is nowhere zero and
$P(z',w)=\prod_{j=1}^d(w-\alpha_j(z'))$ is monic in $w$; the roots need not be chosen holomorphically.
After shrinking the polydisk, all roots lie in a fixed bounded disk.
The arithmetic--geometric mean inequality gives
\[
 |P|^{-2c}=\prod_{j=1}^d|w-\alpha_j|^{-2c}
 \le\frac1d\sum_{j=1}^d|w-\alpha_j|^{-2dc}.
\]
If $dc<1$, the integral in $w$ of the right-hand side has a finite bound independent of the root locations,
because translation bounds it by the integral of $|w|^{-2dc}$ over a fixed disk.
Integrate in $z'$ and use a positive lower bound for $|u|$ to prove integrability.
When $g(0)\ne0$, a positive lower bound for $|g|$ proves the assertion directly.

If $|f|\le C|g|^k$, then
$|f|^2|g|^{-2(k+\epsilon)}\le C^2|g|^{-2\epsilon}
\le C^2|g_j|^{-2\epsilon}$ for any nonzero generator $g_j$.
Choose $\epsilon>0$ sufficiently small.
The ordinary power $I^k$ is generated by products of $k$ generators, each bounded in modulus by $|g|^k$,
so $I^k\subset I^{(k)}$. The two product inclusions follow, respectively, by multiplying growth bounds and
by inserting a growth bound into the weighted integral. For finite sums, take the smallest admissible positive $\epsilon$.
\end{proof}

\begin{lemma}[Norm reduction to at most $n$ generators]\label{div:reduction}
If $r>n$, there are $n$ constant-coefficient linear combinations
$h_\alpha=\sum_i a_{\alpha i}g_i\in I$ such that
$c|g|\le|h|\le C|g|$ on a neighborhood of the origin.
Thus every $I^{(k)}$ and $J^{(k)}$ defined using $g$ agrees with the corresponding object defined using $h$.
\end{lemma}
\begin{proof}
On a small connected neighborhood $U$, consider
$U\setminus Z\to\mathbf P^{r-1}$,$z\mapsto[g_1(z):\cdots:g_r(z)]$.
The closure of its graph, $\Gamma\subset U\times\mathbf P^{r-1}$, is an irreducible analytic set of dimension $n$.
Analyticity can be checked in projective coordinate charts: write the incidence equations
$g_i\xi_j-g_j\xi_i=0$ and retain the irreducible component dominating $U$, which contains the graph over $U\setminus Z$.
Away from $Z$, this component is exactly the graph. Density of the regular locus and the identity principle identify its closure with that component.
Components contained entirely in $Z\times\mathbf P^{r-1}$ do not belong to $\Gamma$.
Compactness of the projective factor makes $\pi:\Gamma\to U$ proper.

The fiber $F=\pi^{-1}(0)$ is a proper analytic subset of the irreducible $n$-dimensional graph $\Gamma$. The dimension-drop argument of Proposition~\ref{alg:analytic-stratification} therefore gives $\dim F\le n-1$.
Choose projective hyperplanes $H_1,\ldots,H_n$ successively so that
$F\cap H_1\cap\cdots\cap H_n=\varnothing$.
At each step, choose a hyperplane that contains none of the finitely many current irreducible components.
The zero set of a nonidentically zero linear form has smaller dimension on each positive-dimensional component.
Once only finitely many points remain, choose a hyperplane avoiding them.
Each choice avoids finitely many proper linear subspaces of the space of linear forms, so it is possible.
Write $\ell_\alpha$ for an equation of $H_\alpha$, and set $h_\alpha=\ell_\alpha(g)$.

The function
$a([\xi])=(\sum_\alpha|\ell_\alpha(\xi)|^2)^{1/2}/|\xi|$
has a positive lower bound on the compact set $F$. After shrinking $U$, properness gives a positive lower bound
$c$ on $\pi^{-1}(U)$ as well. Otherwise there would be $z_j\to0$ and
$\xi_j\in\pi^{-1}(z_j)$ with $a(\xi_j)\to0$; projective compactness would give
a limiting point in $F$, a contradiction.
Along the graph, substitute $\xi=g(z)$ to obtain $|h(z)|\ge c|g(z)|$.
The opposite inequality is the norm bound for a fixed linear map. Continuity extends both inequalities over the zero set.
Norm comparability makes all growth and integrability conditions equivalent.
\end{proof}

\section{The General Brian\c{c}on--Skoda Theorem and Integral Closure}

\begin{theorem}[The general Brian\c{c}on--Skoda theorem]\label{div:briancon-skoda}
Let $0\ne I\subset\cO_{\C^n,0}$ be generated by $r$ elements, and put
$p=\min(n-1,r-1)$. Then, for real $k\ge p$ in the equality and integers $k\ge1$ in the inclusion,
\[
 J^{(k+1)}=I J^{(k)}=I^{(1)}J^{(k)}\quad(k\ge p),
 \qquad J^{(k+p)}\subset I^k\quad(k\ge1).
\]
In particular, $I^{(k+p)}\subset I^k$ for every integer $k\ge1$.
\end{theorem}
\begin{proof}
First suppose $r\le n$, so $p=r-1$.
Take $f\in J^{(k+1)}$ and choose $\epsilon>0$ making its defining integral finite.
On a sufficiently small ball, apply Theorem \ref{div:skoda} with the additional weight
$\varphi=(k-p)\log|g|^2$. It is psh by Example~\ref{pos:log-norm} and Lemma~\ref{cur:psh-integrable}, since $k\ge p$; when $k=p$, take it to be zero.
The input exponent in Skoda's theorem is
$p+1+\epsilon+(k-p)=k+1+\epsilon$, exactly the known integrability exponent.
The output is $f=\sum_i g_i h_i$ with
$\int|h_i|^2|g|^{-2(k+\epsilon)}<\infty$.
Thus $h_i\in J^{(k)}$, proving $J^{(k+1)}\subset I J^{(k)}$.
The reverse inclusion and the second equality follow from
\[
 I J^{(k)}\subset I^{(1)}J^{(k)}\subset J^{(k+1)},
\]
where the last inclusion is the multiplication property in Lemma \ref{div:small-power}.

If $r>n$, replace $g$ by the $n$ elements $h$ of Lemma \ref{div:reduction}.
Their ideal $I_0$ is contained in $I$, while all corresponding $J^{(k)}$ are unchanged.
The case already proved gives $J^{(k+1)}=I_0J^{(k)}\subset I J^{(k)}$;
the reverse inclusions hold as before, giving the assertion with $p=n-1$.

For an integer $k\ge1$, apply the equality successively at exponents $k+p,k+p-1,\ldots,p+1$ to obtain
$J^{(k+p)}=I^kJ^{(p)}\subset I^k$.
Finally use $I^{(k+p)}\subset J^{(k+p)}$.
\end{proof}

\begin{lemma}[A polynomial root bound]\label{div:root-bound}
Let $d\ge1$ be an integer and $x,a_1,\ldots,a_d\in\C$. If $x^d+a_1x^{d-1}+\cdots+a_d=0$, then
$|x|\le2\max_{1\le j\le d}|a_j|^{1/j}$.
For a nonzero leading coefficient, first divide the equation by it.
\end{lemma}
\begin{proof}
If the right-hand side vanishes, the equation is $x^d=0$. Otherwise suppose
$|x|>2\max_j|a_j|^{1/j}$. Then
\[
 \left|\sum_{j=1}^d a_jx^{d-j}\right|
 \le |x|^d\sum_{j=1}^d\frac{|a_j|}{|x|^j}
 <|x|^d\sum_{j=1}^d2^{-j}<|x|^d,
\]
contradicting the equation. The root bound follows.
\end{proof}

\begin{theorem}[Growth characterizes integral closure]\label{div:integral-closure}
The ideal $I^{(1)}$ is exactly the integral closure $\overline I$:
a germ $f$ belongs to $I^{(1)}$ if and only if it satisfies a monic relation
$f^d+a_1f^{d-1}+\cdots+a_d=0$ with $a_j\in I^j$.
Consequently, the usual Brian\c{c}on--Skoda inclusion is
\[
 \boxed{\ \overline{I^{k+p}}\subset I^k,
       \qquad p=\min(n-1,r-1),\ k\ge1.\ }
\]
\end{theorem}
\begin{proof}
Given an integral relation, each $a_j$ is a finite holomorphic linear combination of products of $j$ generators,
so $|a_j|\le C_j|g|^j$ on a small neighborhood.
Lemma \ref{div:root-bound} gives $|f|\le C|g|$, that is, $f\in I^{(1)}$.

For the converse, use a finite module to encode the division relations.
Put $p=\min(n-1,r-1)$ as in Theorem~\ref{div:briancon-skoda}, and let $M=J^{(p)}$. Theorem~\ref{alg:noetherian} makes this ideal finitely generated.
It is nonzero: a sufficiently high power of a nonzero generator $g_j$ belongs to $M$.
For $f\in I^{(1)}$, Theorem \ref{div:briancon-skoda} gives
$fM\subset J^{(p+1)}=IM$.
Choose generators $w_1,\ldots,w_N$ and write
$fw_\alpha=\sum_\beta a_{\alpha\beta}w_\beta$ with $a_{\alpha\beta}\in I$.
Multiply $(f\operatorname{Id}-A)w=0$ by the adjugate matrix to obtain
$\det(f\operatorname{Id}-A)w_\alpha=0$ for every $\alpha$.
The local holomorphic ring is an integral domain, and some $w_\alpha\ne0$,
so $\det(f\operatorname{Id}-A)=0$.
Expanding the determinant gives a monic relation. The coefficient of $f^{N-j}$ is a sum of products of $j$ entries $a_{\alpha\beta}$,
and belongs to $I^j$. This is the required integral relation.

Finally, standard generators of $I^s$ are products of $s$ generators $g_j$.
Their sum of squared moduli differs from $|g|^{2s}$ only by fixed positive combinatorial coefficients, so their norms are comparable.
Apply the result just proved to $I^s$ to obtain $\overline{I^s}=I^{(s)}$.
Take $s=k+p$ and apply Theorem~\ref{div:briancon-skoda} to obtain the boxed inclusion.
\end{proof}

\begin{example}[Why the dimension loss cannot simply be removed]
In $\C^2$, take $I=(x^2,y^2)$.
Although $xy\notin I$, we have $|xy|\le(|x|^2+|y|^2)/2\le C(|x|^4+|y|^4)^{1/2}$,
so $xy\in\overline I$. It also satisfies the integral relation
$(xy)^2-x^2y^2=0$, whose constant term belongs to $I^2$.
Thus $\overline I\subset I$ is false in general.
Here $p=1$, and the theorem gives the correct inclusion $\overline{I^{k+1}}\subset I^k$.
\end{example}


\chapter{Brian\c{c}on--Skoda for Monomial Ideals}\label{bs:chapter}
\begin{learning}
For a monomial ideal, divisibility can be read directly from exponent vectors. Integral dependence imposes convex inequalities on those vectors, and a rounding argument turns the inequalities back into divisibility by products of generators. This gives a proof of the Brian\c{c}on--Skoda inclusion using convergent power series and finite-dimensional convex geometry; the analytic division theorem of Chapter~\ref{div:chapter} is not required for this proof. For further algebraic background, see \cite[Chapters 1 and 13]{HunekeSwanson2006}.
\end{learning}

\section{Integral dependence and weighted vanishing orders}
Write $R=\cO_{\C^n,0}=\C\{z_1,\ldots,z_n\}$ for the ring of convergent power series at the origin. Let
\[
 I=(z^{a_1},\ldots,z^{a_m}),\qquad
 a_j\in\N_0^n\setminus\{0\},\qquad z^a=\prod_{i=1}^nz_i^{a_i},
\]
where $m\ge1$. For exponent vectors, $w\cdot b=\sum_iw_i b_i$ is the ordinary real dot product. Redundant generators are allowed. The ideal $I^q$ is generated by products of $q$ generators, and hence is again a finitely generated monomial ideal. All vector inequalities, such as $b\ge a$, are coordinatewise.

\begin{definition}[Integral dependence over an ideal]\label{bs:integral-definition}
An element $f\in R$ is \emph{integral over} an ideal $J\subset R$ if it satisfies a monic equation
\[
 f^d+c_1f^{d-1}+\cdots+c_d=0,
 \qquad d\ge1,\qquad c_j\in J^j\quad(1\le j\le d).
\]
The set of such elements is denoted by $\overline J$ and called the integral closure of $J$. We use this equation as the definition, without presupposing other structural results about integral closure. In particular, $f\in\overline{I^t}$ requires $c_j\in(I^t)^j=I^{tj}$; membership only in $I^t$ would not suffice.
\end{definition}

\begin{lemma}[Vanishing order for a strictly positive weight]\label{bs:weighted-order}
Fix $w=(w_1,\ldots,w_n)$ with every $w_i>0$. For nonzero $f=\sum_b f_bz^b$, define
\[
 \nu_w(f)=\min_{f_b\ne0}w\cdot b,
 \qquad \nu_w(0)=+\infty.
\]
The minimum exists, and
\[
 \nu_w(fg)=\nu_w(f)+\nu_w(g),\qquad
 \nu_w(f+g)\ge\min\{\nu_w(f),\nu_w(g)\}.
\]
The second inequality is an equality when the two orders on the right differ. If $h(w)=\min_jw\cdot a_j$, then $c\in I^s$ implies $\nu_w(c)\ge sh(w)$.
\end{lemma}
\begin{proof}
Because $w_i>0$, the set $\{b\in\N_0^n:w\cdot b\le M\}$ is finite for every $M$. Choose any exponent in the support to give an upper bound. A minimum then exists and only finitely many terms attain it. Their sum, denoted by $\operatorname{in}_w(f)$, is a nonzero polynomial.
The smallest possible weight in $fg$ is $\nu_w(f)+\nu_w(g)$, and the terms of this weight form exactly
$\operatorname{in}_w(f)\operatorname{in}_w(g)$.
A product of nonzero polynomials is nonzero: order their finitely many exponent vectors lexicographically; the product of the largest monomials is the unique largest product and cannot cancel. This proves multiplicativity.
Addition produces no term below the smaller of the two orders. If these orders differ, the lowest-weight terms cannot cancel.
Finally, express $c\in I^s$ as a finite sum of products of $s$ generators with analytic coefficients. Each coefficient has nonnegative order and each product has order at least $sh(w)$. The two established properties give the claim.
\end{proof}

\begin{proposition}[Integral dependence gives every weighted lower bound]\label{bs:integral-weight}
If $t\ge1$ is an integer and $f\in\overline{I^t}$, then for every strictly positive weight $w$,
\[
 \nu_w(f)\ge th(w).
\]
\end{proposition}
\begin{proof}
There is nothing to prove when $f=0$. Otherwise suppose $v=\nu_w(f)<th(w)$ and take an integral-dependence equation. Its leading term $f^d$ has order $dv$, whereas for $j\ge1$,
\[
 \nu_w(c_jf^{d-j})\ge tjh(w)+(d-j)v>dv.
\]
Only $f^d$ contributes nonzero terms of the lowest weight. Their sum cannot vanish, contradicting the equation.
\end{proof}

\section{The exponent polyhedron and its boundary}
For $I=(z_1^2,z_2^3)$, a monomial $z_1^u z_2^v$ belongs to $I$ precisely when $u\ge2$ or $v\ge3$. Its exponent set is therefore the union of the integer regions above $(2,0)$ and $(0,3)$ in the coordinatewise order. The following convex set fills the region between these two corners; in general it records all convex combinations of the generator exponents and all larger vectors.

The Newton polyhedron is
\[
 P=\operatorname{conv}\{a_1,\ldots,a_m\}+\R_+^n,
 \qquad \R_+^n=\{c:c_i\ge0\}.
\]
Here $\operatorname{conv}$ consists of nonnegative linear combinations whose coefficients sum to one, and the plus sign is the Minkowski sum of sets. For $t>0$, write $tP=\{tx:x\in P\}$, and put $\mathbf1=(1,\ldots,1)$. Passing to $P$ temporarily fills the integral gaps between monomial exponents.

\begin{lemma}[Separation and short convex combinations]\label{bs:polyhedron}
The set $P$ is closed and convex, and has the following properties.
\begin{enumerate}
\item A point $x$ belongs to $P$ if and only if
$w\cdot x\ge\min_jw\cdot a_j$ for every $w\in\R_+^n$.
\item With $\ell=\min(n,m)$, every $x\in P$ admits a representation
\[
 x=\sum_{j=1}^s\lambda_j a_{i_j}+c,
 \qquad s\le\ell,\quad \lambda_j\ge0,\quad
 \sum_{j=1}^s\lambda_j=1,\quad c\in\R_+^n.
\]
\end{enumerate}
\end{lemma}
\begin{proof}
The convex hull of finitely many points is compact. Its sum with a closed set is closed: for a convergent sequence of sums, pass to a convergent subsequence of the compact components and subtract to obtain the limit of the nonnegative components. Convexity is immediate, and $P+\R_+^n=P$.

For the first assertion, membership plainly implies the inequalities. If $x\notin P$, choose a nearest point $y\in P$. It exists because the distance minimization may be restricted to a sufficiently large closed ball, whose intersection with $P$ is compact. Put $w=y-x\ne0$.
For $z\in P$, the segment $y+u(z-y)$ lies in $P$ for $0\le u\le1$. The right derivative at $u=0$ of its squared distance to $x$ is nonnegative, giving $w\cdot(z-y)\ge0$.
Set $z=y+ve_i$ with $v>0$. This proves $w_i\ge0$ for every $i$. Thus
\[
 w\cdot x=w\cdot y-|w|^2
 <\inf_{z\in P}w\cdot z
 =\min_jw\cdot a_j.
\]
The last equality uses $w\ge0$ and the definition of $P$. This separates $x$ from $P$.

If $m\le n$, the second assertion follows directly from the definition of the convex hull. Suppose $m>n$, write ${\bf1}=(1,\ldots,1)$, and set
\[
 \tau=\max\{u\ge0:x-u{\bf1}\in P\},\qquad p=x-\tau{\bf1}.
\]
The admissible set is nonempty and closed. Since $P\subset\R_+^n$, it is bounded above by $\min_i x_i$, so the maximum exists. The point $p$ is on the boundary; otherwise it could be moved a little farther in the direction $-{\bf1}$.
We construct a nonzero supporting normal, without requiring strict positivity. Take $x_\epsilon=p-\epsilon{\bf1}\notin P$, choose a nearest point $y_\epsilon\in P$, and normalize $y_\epsilon-x_\epsilon$ to a unit vector $w_\epsilon$. The preceding argument gives $w_\epsilon\ge0$ and
$w_\epsilon\cdot(z-y_\epsilon)\ge0$ for all $z\in P$.
Since $p\in P$, $|y_\epsilon-p|\le2\epsilon\sqrt n$. Along a convergent subsequence of the unit vectors we obtain
\[
 |w|=1,\quad w\ge0,\quad w\cdot(z-p)\ge0\quad(z\in P).
\]
Write $p=\sum_j\lambda_j a_j+c$ with $c\ge0$ and $\sum_j\lambda_j=1$. The supporting value is $h=\min_jw\cdot a_j=w\cdot p$, so
\[
 0=\sum_j\lambda_j(w\cdot a_j-h)+w\cdot c.
\]
Every summand is nonnegative. Hence every vertex with $\lambda_j>0$ lies in the hyperplane $w\cdot a_j=h$, of affine dimension at most $n-1$.

If more than $n$ coefficients are positive, the corresponding vertices are affinely dependent. Choose numbers $\eta_j$, not all zero, with
$\sum_j\eta_j=0$ and $\sum_j\eta_ja_j=0$.
There are both positive and negative $\eta_j$. Put
$v=\min_{\eta_j>0}\lambda_j/\eta_j>0$ and replace $\lambda_j$ by $\lambda_j-v\eta_j$.
The new coefficients remain nonnegative, still sum to one, and represent the same point, while at least one positive coefficient becomes zero. Repetition leaves at most $n$ terms. Finally add back $\tau{\bf1}$ and absorb it into the nonnegative remainder. This is the required representation of $x$.
\end{proof}

\begin{corollary}[Every support exponent is constrained]\label{bs:support}
Let $t\ge1$ be an integer. If $f=\sum_bf_bz^b\in\overline{I^t}$ and $f_b\ne0$, then $b\in tP$.
\end{corollary}
\begin{proof}
For every strictly positive $w$, Proposition~\ref{bs:integral-weight} gives
$w\cdot b\ge\nu_w(f)\ge th(w)$.
For a merely nonnegative weight, apply this to $w+\epsilon{\bf1}$ and let $\epsilon\downarrow0$. A minimum of finitely many continuous functions is continuous, so the inequality persists. Apply the first part of Lemma~\ref{bs:polyhedron} to $b/t$ to obtain $b/t\in P$.
\end{proof}

\section{Rounding and analytic division}
\begin{lemma}[Termwise monomial divisibility implies analytic ideal membership]\label{bs:analytic-division}
Let $J=(z^{\beta_1},\ldots,z^{\beta_N})\subset R$.
A convergent series $f=\sum_bf_bz^b$ belongs to $J$ if and only if every exponent with $f_b\ne0$ satisfies $b\ge\beta_j$ for some $j$.
\end{lemma}
\begin{proof}
If $f=\sum_jz^{\beta_j}g_j$, each term in the expansions satisfies this condition. Cancellation only removes exponents and cannot create new ones, proving necessity.
Conversely, assign each support exponent to the \emph{first} index $j$ with $\beta_j\le b$. This partitions the support into disjoint sets $S_j$. Define
\[
 g_j(z)=\sum_{b\in S_j}f_bz^{b-\beta_j}.
\]
We must check convergence. Choose a sufficiently small polyradius $\rho=(\rho_1,\ldots,\rho_n)$ with all $\rho_i>0$ such that $\sum_b|f_b|\rho^b<\infty$. Then
\[
 \sum_{b\in S_j}|f_b|\rho^{b-\beta_j}
 \le\rho^{-\beta_j}\sum_b|f_b|\rho^b<\infty.
\]
Thus each $g_j$ converges absolutely on this polydisc and defines an analytic germ. Finite regrouping of an absolutely convergent series gives $f=\sum_jz^{\beta_j}g_j\in J$. This convergence check distinguishes analytic division from merely formal division.
\end{proof}

\begin{theorem}[The monomial Brian\c{c}on--Skoda inclusion]\label{bs:theorem}
Suppose $I\subset\cO_{\C^n,0}$ is generated by $m\ge1$ nonunit monomials and $\ell=\min(n,m)$. Then, for every integer $q\ge1$,
\[
 \boxed{\quad\overline{I^{q+\ell-1}}\subset I^q.\quad}
\]
\end{theorem}
\begin{proof}
Put $t=q+\ell-1$ and let $f\in\overline{I^t}$. By Corollary~\ref{bs:support}, every exponent $b$ in its support satisfies $b/t\in P$. The second part of Lemma~\ref{bs:polyhedron} gives
\[
 b\ge\sum_{j=1}^s\Lambda_ja_{i_j},\qquad
 \Lambda_j\ge0,\quad\sum_j\Lambda_j=t,\quad s\le\ell.
\]
Set $r_j=\lfloor\Lambda_j\rfloor$, the greatest integer at most $\Lambda_j$. Since $r_j>\Lambda_j-1$,
\[
 \sum_jr_j>t-s\ge t-\ell=q-1.
\]
The left-hand side is an integer, so it is at least $q$. All coordinates of $a_{i_j}$ are nonnegative, hence $b\ge\sum_jr_ja_{i_j}$.
Therefore $z^b$ is divisible by a product of at least $q$ generators and belongs to $I^q$. Equivalently, choose exactly $q$ of those factors to obtain a monomial generator of $I^q$ dividing $z^b$.
There are finitely many such generators. Applying Lemma~\ref{bs:analytic-division} to the full convergent series proves $f\in I^q$.
\end{proof}

\section{Examples and the loss in the exponent}
\begin{example}[No loss in one variable]\label{bs:one-variable}
When $n=1$, let $d$ be the smallest generator exponent. Then $I=(z^d)$ and $\ell=1$, so the theorem gives $\overline{I^q}\subset I^q$. The reverse inclusion follows from the monic equation $T-f=0$ for $f\in I^q$, and therefore $\overline{I^q}=I^q$.
Alternatively, if $f$ is integral over $(z^{dq})$, Proposition~\ref{bs:integral-weight} forces its lowest degree to be at least $dq$. Its convergent series is consequently divisible by $z^{dq}$.
The same loss-free conclusion holds for a principal ideal generated by one monomial in any dimension.
\end{example}

\begin{example}[Two generators]\label{bs:two-generators}
Take $I=(z_1^a,z_2^b)$ with $a,b\ge1$. Then
\[
 P=\{(u,v)\in\R_+^2:u/a+v/b\ge1\}.
\]
Indeed, a convex combination of the two exponent vectors is $(a\lambda,b(1-\lambda))$. Choosing $0\le\lambda\le1$ so that this vector is at most $(u,v)$ is equivalent to $u/a+v/b\ge1$.
If $f\in\overline{I^{q+1}}$, every support exponent satisfies $u/a+v/b\ge q+1$. Therefore
\[
 \left\lfloor\frac ua\right\rfloor+
 \left\lfloor\frac vb\right\rfloor
 >\frac ua+\frac vb-2\ge q-1.
\]
This integer is at least $q$, so the monomial is divisible by one of
$z_1^{ai}z_2^{b(q-i)}$ with $0\le i\le q$.
For $I=(z_1^2,z_2^2)$, the element $z_1z_2$ does not belong to $I$ but satisfies
$T^2-z_1^2z_2^2=0$, whose constant term belongs to $I^2$. Thus $z_1z_2\in\overline I$. Even for monomial ideals, the theorem cannot in general be strengthened to $\overline{I^q}\subset I^q$.
\end{example}

\begin{example}[The exponent cannot uniformly be lowered by one]\label{bs:sharpness}
In $\C^r$, with $r\ge2$, take
$I=(z_1^r,\ldots,z_r^r)$ and $g=(z_1\cdots z_r)^{r-1}$.
Every variable has exponent less than $r$, so $g\notin I$. On the other hand,
\[
 g^r=\prod_{j=1}^r(z_j^r)^{r-1}\in I^{r(r-1)}=(I^{r-1})^r.
\]
The equation $T^r-g^r=0$ shows that $g$ is integral over $I^{r-1}$, so
$g\in\overline{I^{r-1}}\setminus I$.
With $q=1$ and $n=m=r$, the exponent $r$ in the theorem therefore cannot uniformly be replaced by $r-1$.
\end{example}

\begin{remark}[Relation with the general theorem]\label{bs:scope}
The loss $\ell-1$ comes from explicit counting. A convex combination on the supporting face needs at most $\ell$ terms, rounding down loses strictly less than one per term, and the final sum is an integer.
This argument does not turn arbitrary holomorphic generators into monomials. Cancellation between terms and the relation between integral closure and local size estimates require additional tools. Those tools were established in Chapter~\ref{div:chapter}: Theorem~\ref{div:skoda} proves general Skoda division, Theorem~\ref{div:briancon-skoda} deduces general Brian\c{c}on--Skoda, and Theorem~\ref{div:integral-closure} identifies integral closure through growth estimates.
The present convex-geometric proof provides a computable model and explains the exponent loss. The two proofs can be read side by side. For further algebraic context see \cite{Skoda1972,Demailly2012,HunekeSwanson2006}.
\end{remark}


\chapter{From a Complex Line to Skoda Integrability}
\label{int:chapter}

A multiplier ideal records which holomorphic functions have integrable weighted squares.
To relate it to the geometric strength of a singularity, we need a complementary statement:
when the singularity is sufficiently mild, the constant function is itself integrable.
Skoda's integrability theorem answers this question through the Lelong number. Its proof combines one-variable logarithmic potentials, a local $L^2$ extension theorem, and restriction to a generic complex line.
The extension argument follows Chen's auxiliary-weight method \cite{Chen2011}. It extends the constant function from a line to a neighborhood of the point; the resulting holomorphic function stays away from zero on a smaller neighborhood and transfers the weighted integrability to the constant function.

Write $\Delta_R=\{w\in\C:|w|<R\}$. On a product $P'\times\Delta_R$, $dV'$ and $dV=dV'\,dA(w)$ denote Euclidean Lebesgue measures, and $z=(z',w)$ denotes the full coordinate tuple.

\section{A Local \texorpdfstring{$L^2$}{L2} Extension Theorem}

The construction again starts by cutting off a holomorphic function and solving away the $\dbar$ error. This time the correction must vanish on the slice $w=0$, so that the prescribed values survive. A factor $|w|^{-2}$ in its squared norm forces this vanishing: a holomorphic function with nonzero value on the slice would give the divergent radial integral $\int_0^\delta dr/r$.

\begin{theorem}[$L^2$ extension on a product polydisk]\label{int:extension}
Let $P=P'\times\Delta_R\subset\C^{n-1}\times\C$ be a bounded polydisk,
with $R<e^{-1/2}$ and coordinates $(z',w)$.
Let $\varphi$ be psh on $P$, and suppose $f\in\cO(P')$ satisfies
$\int_{P'}|f|^2e^{-\varphi(z',0)}dV'<\infty$.
Then there is $F\in\cO(P)$ with $F(z',0)=f(z')$ and
\[
 \int_P\frac{|F|^2e^{-\varphi}}{|w|^2\log^2|w|^2}\,dV
 \le C\int_{P'}|f|^2e^{-\varphi(z',0)}dV',
\]
where $C$ is an absolute constant independent of $f$ and $\varphi$.
Removing the additional positive weight on the left also gives an estimate with a constant depending only on $R$.
\end{theorem}
\begin{proof}
First assume that $f$ and a smooth psh weight $\varphi$ are defined near the corresponding closed polydisks.
These two temporary assumptions will be removed at the end.
Choose a smooth $0\le\chi\le1$ with $\chi(t)=1$ for $t\le1/2$ and $\chi(t)=0$ for $t\ge1$.
Put $t=|w|^2$, and choose $\epsilon>0$ so small that $t+\epsilon^2<e^{-1}$ throughout $P$.
Define
\[
 \rho=\log(t+\epsilon^2),\quad s=-\rho>1,\quad
 \eta=s+\log s,\quad\psi=-\log\eta,\quad
 \Phi=\varphi+\log|w|^2.
\]
For fixed $\epsilon$, $\psi$ is smooth and bounded, whereas $\Phi$ has a logarithmic pole along $w=0$.

\emph{Step 1: exact positivity of the auxiliary weight.}
Since $\partial\eta=-(1+1/s)\partial\rho$, differentiation gives
\begin{equation}\label{int:aux-hessian}
 \partial\dbar\psi
 =\frac{1+1/s}{\eta}\partial\dbar\rho
   +\frac{\partial\rho\wedge\dbar\rho}{\eta s^2}
   +\frac{\partial\eta\wedge\dbar\eta}{\eta^2}.
\end{equation}
After multiplication by $\ii$, all three terms are semipositive, so $\psi$ is psh.
The last two terms together dominate
$\bigl(1+\eta/(s+1)^2\bigr)\partial\psi\wedge\dbar\psi$.
Let $B$ be the Levi matrix of the smooth weight $\varphi+\psi$, and let $B_1$ be its induced operator on $(0,1)$ forms from Theorem~\ref{l2:domain-singular}. The notation $|a|_{B^{-1}}^2$ means $\langle B_1^{-1}a,a\rangle$, with the infinite-energy convention of Proposition~\ref{pos:inverse-curvature}. The variational formula in Lemma~\ref{l2:metric-monotonicity} gives
\begin{equation}\label{int:gradient-global}
 |\dbar\psi|_{B^{-1}}^2
 \le\frac{1}{1+\eta/(s+1)^2},\qquad
 1-|\dbar\psi|_{B^{-1}}^2\ge\frac1{6\eta}.
\end{equation}
The second inequality uses $s\le\eta\le2s$ and
$(s+1)^2+\eta\le6s^2\le6\eta^2$.
On the other hand,
$\rho_{w\bar w}=\epsilon^2/(t+\epsilon^2)^2$, so
\begin{equation}\label{int:gradient-annulus}
 B\ge\frac{\epsilon^2}{\eta(t+\epsilon^2)^2}\,dw\otimes d\bar w,
 \qquad
 |\dbar\psi|_{B^{-1}}^2\le\frac4\eta
       \quad\text{when }\epsilon^2/2\le t\le\epsilon^2.
\end{equation}
For the final bound, substitute
$|\psi_{\bar w}|=(1+1/s)|w|/[\eta(t+\epsilon^2)]$ into the preceding inequality.

\emph{Step 2: an appropriate minimal-norm solution.}
Extend $f$ independently of $w$ and put
$v=f\,\dbar\chi(t/\epsilon^2)$.
This form is closed and supported in the thin annulus above.
Apply Theorem \ref{l2:domain-singular} with weight $\Phi+|z|^2$ to obtain initially
a solution of $\dbar u=v$ that is square integrable with weight $e^{-\Phi}$;
the additional $|z|^2$ is bounded on the polydisk, so the two norms are equivalent.
Among these solutions choose one minimizing the $L^2(e^{-\Phi})$ norm. The solution set is closed: weighted convergence implies distributional convergence on compact sets, and $\dbar$ is distributionally continuous. Orthogonal projection onto this closed affine space, as in Lemma~\ref{pos:hilbert-existence}, gives a unique minimizer. It is orthogonal to all holomorphic functions in that space.
Because $\psi$ is bounded, $ue^\psi$ is orthogonal to all holomorphic functions in
$L^2(e^{-\Phi-\psi})$ as well: the exponential factors in the pairings cancel.
Thus $ue^\psi$ is the minimal-norm solution of
$\dbar U=e^\psi(v+u\dbar\psi)$.

Apply Theorem~\ref{l2:domain-singular} again to this equation with weight $\Phi+\psi$ and continuous curvature lower bound $B$; the additional current $\ii\partial\dbar\log|w|^2$ is nonnegative by Theorem~\ref{cur:poincare-lelong}.
The parameterized Cauchy--Schwarz inequality then gives, for every $a>0$,
\begin{align*}
 \int|u|^2e^{\psi-\Phi}
 &\le\int|v+u\dbar\psi|_{B^{-1}}^2e^{\psi-\Phi}\\
 &\le(1+a^{-1})\int|v|_{B^{-1}}^2e^{\psi-\Phi}
    +\int|u|^2|\dbar\psi|_{B^{-1}}^2e^{\psi-\Phi}\\
 &\hspace{1cm}+a\int_{\operatorname{supp}v}
                |u|^2|\dbar\psi|_{B^{-1}}^2e^{\psi-\Phi}.
\end{align*}
Where $v=0$, no parameterized inequality was used, so the last error term is confined to the thin annulus.
Take $a=1/48$. By \eqref{int:gradient-global} and \eqref{int:gradient-annulus},
moving the two terms involving $u$ to the left leaves a coefficient at least $1/(12\eta)$.
Since $e^\psi=1/\eta$, we obtain
\begin{equation}\label{int:absorbed}
 \int\frac{|u|^2e^{-\Phi}}{\eta^2}
 \le588\int |v|_{B^{-1}}^2e^{\psi-\Phi}.
\end{equation}
All integrals are finite at this stage: the initial $u$ belongs to $L^2(e^{-\Phi})$, $\psi$ is bounded,
and the inverse-curvature norm of the gradient is bounded by one. The absorption therefore involves only finite quantities.

\emph{Step 3: uniform control of the thin-annulus integral.}
Write $v=v_{\bar w}d\bar w$. On $\operatorname{supp}v$,
$|v_{\bar w}|\le\|\chi'\|_\infty |f||w|/\epsilon^2$.
Using the first inequality in \eqref{int:gradient-annulus} and canceling the factors $\eta$ and $t$ gives
\[
 |v|_{B^{-1}}^2e^{\psi-\Phi}
 \le \|\chi'\|_\infty^2|f|^2
       \frac{(t+\epsilon^2)^2}{\epsilon^6}e^{-\varphi}
 \le\frac{4\|\chi'\|_\infty^2}{\epsilon^2}|f|^2e^{-\varphi}.
\]
The annulus has area $O(\epsilon^2)$. Integrate first in $w$ and use uniform continuity of the smooth weight on the closed polydisk to obtain
the bound
\[
 \limsup_{\epsilon\downarrow0}
 \int |v|_{B^{-1}}^2e^{\psi-\Phi}
 \le C_0\int_{P'}|f(z')|^2e^{-\varphi(z',0)}dV'.
\]
The constant depends only on the fixed cutoff, not on the size of the derivatives of the weight.

\emph{Step 4: preserving the extension value and passing to the limit.}
For $|w|<\epsilon/\sqrt2$, $v=0$, so $u$ is holomorphic by the distributional Cauchy--Riemann argument in Lemma~\ref{l2:removable}.
The condition $\int|u|^2|w|^{-2}e^{-\varphi}<\infty$ forces $u(z',0)=0$:
a nonzero value would, by continuity, give a positive lower bound on a small product neighborhood,
contradicting divergence of $\int_0^\delta dr/r$.
Thus $F_\epsilon=\chi(t/\epsilon^2)f-u$ is holomorphic throughout $P$ and
$F_\epsilon(z',0)=f(z')$.

Since $\eta\le2s\le2(-\log t)$, \eqref{int:absorbed} controls
the squared integral of $u$ with weight $e^{-\varphi}/(t\log^2t)$.
The same weighted integral of the cutoff term tends to zero: its integral in $w$ is bounded by a constant times
$\int_0^{\epsilon^2}dt/(t\log^2t)=1/|\log\epsilon^2|$.
Consequently, $F_\epsilon$ has a uniform bound in the target norm.
Lemma~\ref{emb:cauchy-l2} gives uniform local bounds for values and derivatives; diagonal subsequence extraction gives local uniform convergence. Corollary~\ref{bas:identity} makes the limit $F$ holomorphic, and it satisfies
$F(z',0)=f(z')$, and Fatou's lemma gives the estimate.

Finally, exhaust $P$ by smaller product polydisks, and near each closed polydisk convolve
$\varphi$ to obtain $\varphi_\delta\searrow\varphi$. The preceding proof applies to each smooth weight.
The right-hand side has a common bound because $e^{-\varphi_\delta(z',0)}\le e^{-\varphi(z',0)}$.
First let $\delta\downarrow0$, then exhaust $P$. Two normal-family subsequence arguments and Fatou's lemma
give the conclusion for general psh weights and general $f$.
\end{proof}

\section{Lelong Numbers and Generic Complex Lines}

\begin{definition}[The normalization of Lelong numbers used here]\label{int:lelong-definition}
For a psh weight $\varphi$ near the origin that is not identically $-\infty$, define
\[
 \nu(\varphi,0)=\lim_{r\downarrow0}
      \frac{\sup_{|z|\le r}\varphi(z)}{\log r^2}.
\]
The next lemma proves existence of the limit. At another point $x$, use centered holomorphic coordinates and write $\nu(\varphi,x)$; the same lemma proves coordinate independence. Thus $\nu(a\log|z|^2,0)=a$ for $a\ge0$, consistently with the metric convention $h=e^{-\varphi}$.
In references using $h=e^{-2u}$, our $\nu(\varphi,0)$ is the usual $\nu(u,0)$.
\end{definition}

Line directions form $\mathbb P^{n-1}$. In the phrase ``almost every line,'' use the probability measure obtained by projecting normalized surface measure on the unit sphere to its complex lines; equivalently, this is the normalized Fubini--Study measure.

\begin{lemma}[Existence of the Lelong number and equality on generic lines]\label{int:generic-line}
The displayed limit exists, is finite and nonnegative, and is independent of centered holomorphic coordinates.
For almost every complex line $\ell$ through the origin, the restriction $\varphi|_\ell$ is not identically $-\infty$ and
$\nu(\varphi|_\ell,0)=\nu(\varphi,0)$.
Moreover, near the origin,
$\varphi(z)\le\nu(\varphi,0)\log|z|^2+C$.
\end{lemma}
\begin{proof}
Put $u=\varphi/2$, add a constant so that $u\le0$ on a ball, and write
$M(r)=\sup_{|z|\le r}u(z)$.
Apply the subharmonic maximum principle on an annulus in each complex line.
Interpolating the two boundary-circle upper bounds by $a\log|\zeta|+b$ shows that $M(e^t)$ is convex in $t$.
It is also nondecreasing and finite for every positive radius.
Its asymptotic slope as $t\to-\infty$ exists and lies between zero and any fixed finite secant slope.
Hence $\nu=\lim M(r)/\log r$ is finite and nonnegative. If $z$ and $w$ are two centered holomorphic coordinate systems, Theorem~\ref{bas:holomorphic-inverse} gives constants $c,C>0$ with $c|z|\le|w|\le C|z|$ near the center. Their ball suprema are therefore bounded between $M(cr)$ and $M(Cr)$. Since $\log(cr)/\log r$ and $\log(Cr)/\log r$ tend to one, the limit is unchanged.
Convexity also makes $M(e^t)-\nu t$ nondecreasing, so
$M(r)\le\nu\log r+C$, proving the final inequality after multiplying by two.

To compare with lines, let $A(r)$ be the spherical average of $u$ on $|z|=r$.
Equivalently, average first over the circle in each complex line and then over line directions.
Each circular mean is convex in $\log r$, so $A(e^t)$ is convex.
Local integrability and Fubini give finite averages for almost every radius and nondegenerate restrictions for almost every line. Monotonicity of circular means then gives finiteness at every positive radius by comparison with a smaller radius having finite average.
Put $\nu_A=\lim A(r)/\log r$.
The inequality $A(r)\le M(r)$ first gives $\nu_A\ge\nu$.

For the reverse inequality, use the Poisson inequality on a real ball. For smooth subharmonic $u$, its harmonic Poisson extension $h$ has the same boundary values and dominates $u$: the function $u-h+\epsilon(|z|^2-s^2)$ is strictly subharmonic, is zero on the sphere of radius $s$, and cannot have a positive interior maximum by the second-derivative test. Let $\epsilon\downarrow0$. For $|z|\le r<s$, the normalized Poisson kernel is at least
$c(r/s)=(1-(r/s)^2)/(1+r/s)^{2n}$.
Because the boundary values satisfy $u\le0$, we obtain $u(z)\le c(r/s)A(s)$, and hence
$M(r)\le c(r/s)A(s)$.
Take $s=r^a$, with $0<a<1$. Then $r/s\to0$; dividing by the negative quantity $\log r$ and letting $r\downarrow0$
gives $\nu\ge a\nu_A$. Letting $a\uparrow1$ proves $\nu_A=\nu$.
For nonsmooth $u$, obtain the Poisson inequality by convolution followed by a decreasing limit; boundary continuity is unnecessary.

Fix a small radius $R$, and denote the circular mean on a line by $m_\ell(r)$.
For almost every line, the nonnegative secant slopes
\[
 b_\ell(r)=\frac{m_\ell(R)-m_\ell(r)}{\log R-\log r}
\]
decrease to $\nu_\ell$ as $r\downarrow0$.
For fixed $r_0<R$ and $r<r_0$, $b_\ell(r)\le b_\ell(r_0)$, and this dominating function is integrable over directions,
with integral $[A(R)-A(r_0)]/[\log R-\log r_0]$.
Dominated convergence shows that the average of $\nu_\ell$ is $\nu_A=\nu$.
On every nondegenerate line, $m_\ell(r)\le M(r)$ implies $\nu_\ell\ge\nu$.
A nonnegative function with zero integral vanishes almost everywhere, so $\nu_\ell=\nu$ for almost every line.
\end{proof}

\begin{lemma}[Exponential integrability in one variable]\label{int:one-variable}
Let $u$ be subharmonic near the origin and suppose
$\lim_{r\downarrow0}\sup_{|z|\le r}u(z)/\log r<1$.
Then $e^{-2u}$ is locally integrable.
\end{lemma}
\begin{proof}
Put $\mu=dd^cu$, a positive measure.
The circular mean $m(r)$ satisfies
$m(R)-m(r)=\int_r^R\mu(\Delta_t)\,dt/t$.
For smooth $u$, this follows from Stokes' theorem; in general, convolve and pass to the limit using monotonicity and weak convergence of measures.
It follows that $\lim m(r)/\log r=\mu(\{0\})$.
The one-variable case of Lemma~\ref{int:generic-line} identifies this with the number in the hypothesis, so $\mu(\{0\})<1$.
Choose $R$ small enough that $a=\mu(\Delta_R)<1$.

Define the logarithmic potential $p(z)=\int_{\Delta_R}\log|z-\zeta|\,d\mu(\zeta)$.
Fubini gives local integrability, and $dd^cp=\mu$ on $\Delta_R$.
The kernel identity of Lemma~\ref{pos:one-variable-lelong} gives $dd^c(u-p)=0$ distributionally. Lemma~\ref{cur:weyl} makes $u-p$ a smooth harmonic function, which has a finite lower bound on a smaller closed disk.
If $a=0$, $u$ is harmonic there and the assertion is immediate.
If $a>0$, apply Jensen's inequality for the exponential to the probability measure $d\mu/a$:
\[
 e^{-2p(z)}
 =\exp\left(\int\log|z-\zeta|^{-2a}\frac{d\mu(\zeta)}a\right)
 \le\int|z-\zeta|^{-2a}\frac{d\mu(\zeta)}a.
\]
Since $a<1$, integration of the right-hand side over $z$ in a smaller disk has a finite bound independent of $\zeta$.
Tonelli proves integrability of $e^{-2p}$. Multiplying by the bounded exponential of the harmonic part proves integrability of $e^{-2u}$.
\end{proof}

\section{Skoda Integrability and Vanishing Orders of Jets}

\begin{theorem}[Skoda's integrability theorem]\label{int:skoda-integrability}
If $\varphi$ is psh and $\nu(\varphi,x)<1$, then $e^{-\varphi}$ is integrable near $x$.
Consequently, $\cI(\varphi)_x=\cO_{X,x}$.
\end{theorem}
\begin{proof}
Choose coordinates with $x=0$.
Lemma \ref{int:generic-line} gives a complex line $\ell$ such that
$\nu(\varphi|_\ell,0)=\nu(\varphi,0)<1$.
Apply Lemma \ref{int:one-variable} to $u=\varphi|_\ell/2$.
The constant function one is then square integrable with weight $e^{-\varphi}$ on a small disk in this line.

Choose linear coordinates making the line the first coordinate axis.
In a sufficiently small product polydisk, Theorem \ref{int:extension} extends the constant one first to
the plane of the first two coordinates, then to the first three coordinates, and successively to $\C^n$.
At every step the ordinary weighted integral is finite, because the extension theorem's extra weight has a positive lower bound.
We obtain a holomorphic $F$ satisfying $F(0)=1$ and
$\int|F|^2e^{-\varphi}<\infty$.
By continuity, $|F|\ge1/2$ on a smaller neighborhood of the origin,
so $e^{-\varphi}\le4|F|^2e^{-\varphi}$ is integrable there.
Every holomorphic germ is bounded on a sufficiently small closed ball, so integrability of the constant is equivalent to the multiplier ideal being the whole local ring.
\end{proof}

\begin{proposition}[Large Lelong numbers force high-order vanishing]\label{int:jet-order}
If $\nu(\varphi,x)\ge n+s$, with an integer $s\ge0$, then
$\cI(\varphi)_x\subset\mathfrak m_x^{s+1}$.
\end{proposition}
\begin{proof}
The growth bound of Lemma \ref{int:generic-line} gives
$\varphi(z)\le(n+s)\log|z-x|^2+C$ on a sufficiently small ball.
Hence $e^{-\varphi}\ge C'|z-x|^{-2(n+s)}$.
As in Proposition~\ref{pos:radial-ideal}, if the lowest nonzero homogeneous term of a nonzero holomorphic $f$ has degree $d$,
its squared spherical mean is a positive constant times $r^{2d}$ plus $O(r^{2d+1})$.
Thus its weighted radial integral is bounded below by
$c\int_0^\delta r^{2d-2(n+s)+2n-1}dr
=c\int_0^\delta r^{2d-2s-1}dr$.
This diverges for $d\le s$, so membership in the multiplier ideal forces $d\ge s+1$.
\end{proof}

\Needspace{8\baselineskip}
\begin{remark}
The two bounds answer different questions: $\nu<1$ makes the multiplier ideal trivial, whereas $\nu\ge n+s$ forces the jet of order $s$ to vanish.
For $1\le\nu<n$, the Lelong number alone does not determine the ideal.
For instance, $\varphi=a\log|z_1|^2$ and $\varphi=a\log|z|^2$ have the same Lelong number,
but their integrability thresholds are governed, respectively, by Theorem~\ref{pos:snc-ideal} and Proposition~\ref{pos:radial-ideal}, namely one-variable transverse integration and $2n$-dimensional radial integration.
\end{remark}


\chapter{Hilbert Complexes, Heat Semigroups, and Spectral Comparison}
\label{heat:chapter}

Holomorphic Morse inequalities relate cohomology dimensions to curvature integrals.
The proof combines three ideas: positive-eigenvalue complexes are exact;
at scale $k^{-1/2}$ the heat kernel approaches a constant-magnetic-field model; the model's long-time limit selects a specified curvature index.
This chapter establishes the first step and the required operator tools. Chapter~\ref{heat:models-chapter} computes the exact models,
Chapter~\ref{heat:local-chapter} proves the general local scaling theorem,
and Chapter~\ref{morse:chapter} completes the proof on compact complex Hermitian manifolds.
At each step we distinguish operator norms, heat traces, and pointwise kernel norms; these comparisons are not interchangeable.

As throughout these notes, inner products are conjugate-linear in the first variable and linear in the second.
The finite-dimensional part requires only linear algebra; the infinite-dimensional part uses the Sobolev spaces,
compact embeddings, and elliptic regularity established in Chapter~\ref{ana:chapter};
for the compact Dolbeault realization, see Lemma~\ref{hod:compact-resolvent}.
The parameter $k$ later denotes a positive tensor-power integer; $t>0$ always denotes heat time.
For a nonnegative Hermitian matrix, the heat operator multiplies a vector of eigenvalue $\lambda$ by $e^{-t\lambda}$.
As $t$ increases, the zero eigenspace stays fixed and every positive eigenspace decays.
The first section shows how this elementary observation detects the cohomology of a complex.

\section{Finite-dimensional complexes: the algebra behind Morse inequalities}
\label{heat:finite-section}

\begin{definition}[Finite-dimensional complexes and adjoints]\label{heat:complex-definition}
A finite-dimensional Hilbert complex is
\[
0\longrightarrow C^0\xrightarrow{d_0}C^1
\xrightarrow{d_1}\cdots\xrightarrow{d_{n-1}}C^n
\longrightarrow0,
\qquad d_{q+1}d_q=0,
\]
where each $C^q$ is a finite-dimensional complex inner-product space.
Set $d_{-1}=d_n=0$, and write
\[
Z^q=\ker d_q,\quad B^q=\operatorname{im}d_{q-1},\quad
H^q=Z^q/B^q,
\qquad
\Delta_q=d_{q-1}d_{q-1}^*+d_q^*d_q.
\]
Here $d_q^*$ denotes the adjoint for the given inner products.
\end{definition}

\begin{proposition}[Finite-dimensional Hodge decomposition]
\label{heat:finite-hodge}
There is an orthogonal direct-sum decomposition
\[
C^q=B^q\oplus\ker\Delta_q\oplus\operatorname{im}d_q^*,
\qquad H^q\simeq\ker\Delta_q.
\]
The latter isomorphism sends a harmonic vector to its cohomology class.
\end{proposition}
\begin{proof}
By the definition of the adjoint,
$(\operatorname{im}d_q^*)^\perp=\ker d_q=Z^q$.
Finite-dimensional subspaces are closed, so
$C^q=Z^q\oplus\operatorname{im}d_q^*$.
Since also $B^q\subset Z^q$,
\[
Z^q=B^q\oplus\bigl(Z^q\cap(B^q)^\perp\bigr).
\]
For $u\in C^q$, direct calculation gives
\[
\langle\Delta_qu,u\rangle
=\|d_{q-1}^*u\|^2+\|d_qu\|^2.
\]
Thus $\ker\Delta_q=\ker d_q\cap\ker d_{q-1}^*
=Z^q\cap(B^q)^\perp$.
This proves the direct-sum decomposition.
Every closed vector has a unique decomposition $d_{q-1}v+h$ with $h\in\ker\Delta_q$,
so every cohomology class has exactly one harmonic representative.
\end{proof}

\begin{theorem}[Finite-dimensional Morse inequalities]
\label{heat:finite-morse}
Set $c_q=\dim C^q$ and $h_q=\dim H^q$.
For $0\leq q\leq n$,
\begin{align}
 h_q&\leq c_q,\label{heat:weak-algebra}\\
 \sum_{j=0}^q(-1)^{q-j}h_j
 &\leq\sum_{j=0}^q(-1)^{q-j}c_j.\label{heat:strong-algebra}
\end{align}
For $q=n$, the second inequality is an equality.
\end{theorem}
\begin{proof}
Write $b_q=\dim B^q$, with $b_0=b_{n+1}=0$.
The rank-nullity theorem gives
\[
c_q=\dim Z^q+b_{q+1}=h_q+b_q+b_{q+1}.
\]
This already proves $h_q\leq c_q$.
Taking alternating sums cancels the intermediate $b_j$ in pairs and yields the sharper identity
\begin{equation}
\sum_{j=0}^q(-1)^{q-j}(c_j-h_j)=b_{q+1}\geq0.
\label{heat:boundary-remainder}
\end{equation}
For $q=n$ the right-hand side is zero, giving the Euler--Poincar\'e identity.
The remainder in the strong inequality is precisely the dimension of the boundary space in the next degree.
\end{proof}

\begin{proposition}[The heat-trace formulation]
\label{heat:finite-heat-morse}
For every $t>0$,
\begin{equation}
\sum_{j=0}^q(-1)^{q-j}h_j
\leq\sum_{j=0}^q(-1)^{q-j}\tr(e^{-t\Delta_j}).
\label{heat:finite-trace-morse}
\end{equation}
There is equality for $q=n$; as $t\to+\infty$, each
$\tr(e^{-t\Delta_j})$ tends to $h_j$.
\end{proposition}
\begin{proof}
First, $d_q\Delta_q=\Delta_{q+1}d_q$: on expansion,
all terms containing $d_qd_{q-1}$ or $d_{q+1}d_q$ vanish.
For fixed $\lambda\geq0$, the $\lambda$-eigenspaces therefore form a subcomplex
$C_\lambda^\bullet$.
If $\lambda>0$ and $u\in C_\lambda^q$ satisfies $d_qu=0$, then
\[
u=\lambda^{-1}\Delta_qu
=d_{q-1}(\lambda^{-1}d_{q-1}^*u),
\]
so this eigenspace subcomplex is exact.
Applying Theorem~\ref{heat:finite-morse} to it gives
\[
\sum_{j=0}^q(-1)^{q-j}\dim C_\lambda^j\geq0.
\]
Multiply by $e^{-t\lambda}>0$ and sum over $\lambda>0$.
By Proposition~\ref{heat:finite-hodge}, the $\lambda=0$ contribution is precisely the left-hand side,
proving \eqref{heat:finite-trace-morse}.
The full alternating sum vanishes at every positive eigenvalue, giving equality in top degree.
Finally, finite-dimensional spectral decomposition gives
$\tr(e^{-t\Delta_j})=h_j+\sum_{\lambda>0}e^{-t\lambda}\dim C_\lambda^j$;
let $t\to\infty$.
\end{proof}

\section{From closed quadratic forms to discrete spectra}
\label{heat:spectral-section}

Spaces of differential forms on a manifold are infinite-dimensional.
To retain the preceding finite-dimensional argument, we must show that eigenspaces are finite-dimensional
and that the relevant sums converge. We treat these questions separately.

\begin{definition}[Closed quadratic forms]\label{heat:closed-form-definition}
Let $H$ be a separable Hilbert space and $V\subset H$ a dense linear subspace.
A Hermitian form $q(u,v)$ is a lower-bounded closed form if there is $C\geq0$ with
$q(u,u)\geq-C\|u\|^2$, and
\[
b(u,v)=q(u,v)+(C+1)\langle u,v\rangle
\]
makes $V$ a Hilbert space.
Write $q(u)=q(u,u)$ and $\|u\|_V^2=b(u,u)$.
The form has compact embedding if $V\hookrightarrow H$ is compact.
\end{definition}

Here compactness means that every bounded sequence in $V$ has a convergent subsequence in $H$;
it does not mean that the unit ball of $V$ is compact in $V$.
Sobolev spaces are typical form domains for Dirichlet energies,
but placing a differential operator in this framework requires the corresponding compact embedding theorem.

\begin{lemma}[Orthogonal projection and Riesz representation]
\label{heat:riesz}
If $H$ is a Hilbert space and $M\subset H$ a closed linear subspace,
each $x\in H$ has a unique decomposition $x=m+w$ with $m\in M$ and $w\perp M$.
Every continuous linear functional $\ell$ has a unique expression
$\ell(v)=\langle z,v\rangle$, with $\|z\|=\|\ell\|$.
\end{lemma}
\begin{proof}
Choose $m_j\in M$ such that $\|x-m_j\|$ tends to the distance $d$ from $x$ to $M$.
The parallelogram identity gives
\[
\|m_j-m_k\|^2
=2\|x-m_j\|^2+2\|x-m_k\|^2
-4\left\|x-\frac{m_j+m_k}{2}\right\|^2,
\]
whose right-hand side tends to zero. Thus $m_j$ converges to some $m\in M$.
Compare the distances from $x$ to $m$ and to $m+ta$, for $a\in M$ and $t\in\C$.
Taking real and purely imaginary $t$ gives $x-m\perp M$.
Subtracting two such decompositions produces a vector in both $M$ and $M^\perp$, hence zero.

If $\ell=0$, take $z=0$.
Otherwise $M=\ker\ell$ is a closed subspace of codimension one.
Choose $0\ne w\in M^\perp$. Then
$z=\overline{\ell(w)}w/\|w\|^2$ satisfies
$\langle z,w\rangle=\ell(w)$, and both sides vanish on $M$.
The identity therefore holds on all of $H=M\oplus\C w$.
Cauchy--Schwarz gives $\|\ell\|\leq\|z\|$;
taking $v=z/\|z\|$ gives the reverse inequality.
Uniqueness follows because $\langle z-z',v\rangle=0$ for every $v$.
\end{proof}

\begin{lemma}[Spectral decomposition of a compact positive operator]
\label{heat:compact-spectral}
Let $H$ be a separable Hilbert space and $S:H\to H$ a compact, self-adjoint, positive, injective bounded operator.
Then $H$ has an orthonormal basis $(e_j)$ satisfying
$Se_j=\mu_je_j$, where $\mu_j>0$.
If $H$ is infinite-dimensional, the eigenvalues may be ordered as $\mu_1\geq\mu_2\geq\cdots\to0$.
Every positive eigenvalue has finite multiplicity.
\end{lemma}
\begin{proof}
The positive form $[u,v]=\langle Su,v\rangle$ satisfies Cauchy--Schwarz:
regard $[u+tv,u+tv]\geq0$ as a quadratic expression in $t\in\C$,
and choose $t$ so that the mixed term is negative, obtaining
$|[u,v]|^2\leq[u,u][v,v]$.
Consequently, if $M=\sup_{\|u\|=1}\langle Su,u\rangle$, then
$|\langle Su,v\rangle|\leq M\|u\|\|v\|$, so $\|S\|=M$.
For $S\ne0$ we have $M>0$.

Choose unit vectors $u_j$ with $\langle Su_j,u_j\rangle\to M$.
A bounded sequence in a separable Hilbert space has a weakly convergent subsequence: extract convergent subsequences successively in the coordinates of an orthonormal basis,
then diagonalize and control the tails by Cauchy--Schwarz.
This proves the required weak compactness directly.
We may therefore suppose $u_j\rightharpoonup u$, with $\|u\|\leq1$.
Compactness gives $Su_j\to Su$, after subsequence extraction and identification by uniqueness of the weak limit.
Thus $\langle Su,u\rangle=M$, and $M\|u\|^2\geq M$ forces $\|u\|=1$.
For any $v\perp u$, consider the maximum at $t=0$ of
$\langle S(u+tv),u+tv\rangle/\|u+tv\|^2$.
Its derivatives in the real and imaginary directions vanish, giving $\langle Su,v\rangle=0$.
Hence $Su=Mu$.

Self-adjointness makes $u^\perp$ invariant under $S$, so repeat the construction there.
The resulting unit eigenvectors are orthogonal.
If infinitely many had eigenvalues at least $\epsilon>0$, the distances between their images would be
at least $\sqrt2\epsilon$, contradicting compactness. This proves both finite multiplicity and $\mu_j\to0$.
Let $W$ be the orthogonal complement of the closed span of the chosen eigenvectors.
If $S|_W\ne0$, the same maximization gives a positive eigenvalue $\nu$ on $W$.
Because the construction chooses the largest remaining eigenvalue at each step, every $\mu_j\geq\nu$, contradicting convergence to zero.
Thus $S|_W=0$, and injectivity implies $W=0$.
\end{proof}

\begin{theorem}[Spectrum and the min--max principle for compactly embedded closed forms]
\label{heat:minmax}
Let $q$ be a lower-bounded closed form as above, with $V\hookrightarrow H$ compact.
There are an orthonormal basis $(e_j)$ and real numbers
$\lambda_1\leq\lambda_2\leq\cdots$ such that
\[
q(e_j,v)=\lambda_j\langle e_j,v\rangle\qquad(v\in V).
\]
In infinite dimension, $\lambda_j\to+\infty$.
For $u=\sum_ju_je_j$,
\[
u\in V\quad\Longleftrightarrow\quad
\sum_j(\lambda_j+C+1)|u_j|^2<\infty,
\qquad q(u)=\sum_j\lambda_j|u_j|^2.
\]
Moreover,
\begin{align}
\lambda_p
 &=\min_{\substack{E\subset V\\\dim E=p}}
       \max_{\substack{u\in E\\\|u\|=1}}q(u),\label{heat:minmax-formula}\\
\lambda_p
 &=\max_{\substack{F\subset H\\\dim F=p-1}}
       \inf_{\substack{u\in V\cap F^\perp\\\|u\|=1}}q(u).
\label{heat:maxmin-formula}
\end{align}
\end{theorem}
\begin{proof}
For $f\in H$, the functional $v\mapsto\langle f,v\rangle$ is continuous on $(V,b)$,
since $\|v\|\leq\|v\|_V$.
Lemma~\ref{heat:riesz} gives a unique $Sf\in V$ with
\[
b(Sf,v)=\langle f,v\rangle\quad(v\in V),
\qquad\|Sf\|_V\leq\|f\|.
\]
As an operator into $H$, $S$ is compact by compactness of the embedding.
Substitute $v=Sg$ and exchange $f,g$ to get
$\langle Sf,g\rangle=\langle f,Sg\rangle$.
Also,
$\langle Sf,f\rangle=b(Sf,Sf)\geq0$.
If $Sf=0$, then $f$ is orthogonal to the dense subspace $V$, so $f=0$.
Lemma~\ref{heat:compact-spectral} therefore applies.

Write $Se_j=\mu_je_j$.
Since $Se_j\in V$ and $\mu_j>0$, we have $e_j\in V$ and
\[
b(e_j,v)=\mu_j^{-1}\langle e_j,v\rangle.
\]
Set $\lambda_j=\mu_j^{-1}-(C+1)$ to obtain the eigenvalue equation.
These vectors are orthogonal for $b$.
If $v\in V$ is $b$-orthogonal to all of them, then
$\langle v,e_j\rangle=0$ for every $j$, hence $v=0$.
They therefore form a complete orthogonal system also in the form norm of $V$.
Parseval's identity gives
$b(u,u)=\sum_j\mu_j^{-1}|u_j|^2$.
Conversely, if this sum is finite, its partial sums are Cauchy in $V$;
their limits in $V$ and $H$ coincide, so $u\in V$.
This proves the characterizations of the form domain and the quadratic form.

Take $E_0=\operatorname{span}\{e_1,\ldots,e_p\}$. Then
$q(u)\leq\lambda_p\|u\|^2$ for $u\in E_0$.
Every $p$-dimensional subspace $E\subset V$ contains a unit vector
$u\perp e_1,\ldots,e_{p-1}$, since there are only $p-1$ linear constraints.
The spectral expansion gives $q(u)\geq\lambda_p$, proving
\eqref{heat:minmax-formula}.
Taking $F_0=\operatorname{span}\{e_1,\ldots,e_{p-1}\}$ gives one direction of
\eqref{heat:maxmin-formula}.
For every $(p-1)$-dimensional $F\subset H$, the intersection $E_0\cap F^\perp$ contains a unit vector
with $q$-value at most $\lambda_p$, giving the other direction.
\end{proof}

\paragraph{Operator domains, coefficients, and trace conventions.}\label{heat:operator-notation}
For an operator $A$, $D(A)=\operatorname{Dom}A$ denotes its domain, endowed with the graph norm
$\|u\|_{D(A)}=(\|u\|^2+\|Au\|^2)^{1/2}$ when needed.
In the expansion above, $u_j=\langle e_j,u\rangle$.
For a bounded operator $T$, $\|T\|$ is its operator norm and $T^*$ its Hilbert adjoint.

The associated self-adjoint operator is defined by
$Au=\sum_j\lambda_ju_je_j$, with domain
$\{u:\sum_j\lambda_j^2|u_j|^2<\infty\}$.
For a bounded function $f$ on the eigenvalue set, define
$f(A)u=\sum_j f(\lambda_j)u_je_j$.
Parseval's identity gives $\|f(A)\|\leq\sup_j|f(\lambda_j)|$.
If every $f(\lambda_j)\geq0$, its positive square root is obtained by replacing each eigenvalue by its nonnegative square root.
Define the heat operator of $A$ spectrally by
\[
e^{-tA}u=\sum_je^{-t\lambda_j}u_je_j,\qquad t\geq0.
\]
The spectrum is bounded below, so this is bounded; its semigroup property follows term by term.
For a bounded positive diagonal operator $Te_j=a_je_j$, set
$\operatorname{tr}T=\sum_j a_j$, allowing $+\infty$.
This sum is independent of the basis: for any other orthonormal basis $(f_\ell)$, Parseval and Tonelli give
\[
 \sum_\ell\langle f_\ell,Tf_\ell\rangle
 =\sum_{\ell,j}a_j|\langle e_j,f_\ell\rangle|^2
 =\sum_j a_j.
\]
Thus $\operatorname{tr}(e^{-tA})=\sum_j e^{-t\lambda_j}$, possibly $+\infty$ before summability is established.
All these diagonal constructions follow directly from Theorem~\ref{heat:minmax}.

\begin{corollary}[Form comparison implies spectral comparison]
\label{heat:form-comparison}
Let $q_A,q_B$ satisfy Theorem~\ref{heat:minmax},
with form domains $V_B\subset V_A$ and $q_A(u)\leq q_B(u)$ for every $u\in V_B$.
Then $\lambda_j(A)\leq\lambda_j(B)$.
If the corresponding heat traces are finite, then
\[
\tr(e^{-tB})\leq\tr(e^{-tA})\qquad(t>0).
\]
In particular, if the forms have the same domain and $0<\epsilon<1$, $c\geq0$ satisfy
\[
(1-\epsilon)q_B(u)-c\|u\|^2
\leq q_A(u)\leq(1+\epsilon)q_B(u)+c\|u\|^2,
\]
then, whenever the relevant series converge,
\begin{equation}
 e^{-ct}\tr(e^{-(1+\epsilon)tB})
 \leq\tr(e^{-tA})
 \leq e^{ct}\tr(e^{-(1-\epsilon)tB}).
\label{heat:trace-sandwich}
\end{equation}
\end{corollary}
\begin{proof}
Restricting the subspaces in the min--max formula for $q_A$ to those in $V_B$ increases the minimum.
Applying $q_A\leq q_B$ on those subspaces gives the eigenvalue comparison.
Since $x\mapsto e^{-tx}$ decreases, termwise
$e^{-t\lambda_j(B)}\leq e^{-t\lambda_j(A)}$.
Sum finitely many terms and then pass to the limit.
Likewise, the two-sided form comparison gives
\[
(1-\epsilon)\lambda_j(B)-c
\leq\lambda_j(A)\leq(1+\epsilon)\lambda_j(B)+c,
\]
Exponentiating termwise and summing proves \eqref{heat:trace-sandwich}.
\end{proof}

\begin{example}[Discrete spectrum alone does not imply finite heat trace]
\label{heat:not-trace-class}
On the standard basis of $\ell^2(\N)$, define $Ae_j=\log(j+1)e_j$.
Its form domain is
$V=\{u:\sum_j(1+\log(j+1))|u_j|^2<\infty\}$.
The form is closed, and $V\hookrightarrow\ell^2$ is compact: on the unit ball of $V$,
the squared norm of the tail after index $N$ is at most $(1+\log(N+1))^{-1}$,
while the finite-dimensional part is compact.
Nevertheless,
\[
\tr(e^{-A})=\sum_{j=1}^\infty\frac1{j+1}=+\infty.
\]
Thus discreteness of the spectrum does not justify skipping summability of the heat trace.
For elliptic operators on compact manifolds, summability follows from elliptic estimates or heat-kernel construction;
it is a separate analytic result.
\end{example}

\begin{remark}[Closability of connection energy]
For a smooth unitary connection, if $u_j\to0$ in $L^2$ and $\nabla u_j\to v$ in $L^2$,
then for every compactly supported smooth bundle-valued one-form $a$,
$\langle v,a\rangle=\lim_j\langle u_j,\nabla^*a\rangle=0$, hence $v=0$.
The form $\|\nabla u\|_2^2$ is therefore closable. After adding a bounded Hermitian potential and a sufficiently large constant,
its form-norm completion gives the closed form used below. The argument also applies to Euclidean magnetic connections with linearly growing coefficients.
\end{remark}

\section{Self-adjoint realizations and spectral calculus on noncompact models}
\label{heat:general-calculus-section}

Heat operators on Euclidean space generally lack discrete spectra. For example, the plane waves of $-\Delta$ do not belong to
$L^2(\R^d)$. The compact spectral theorem above therefore does not apply directly to the magnetic models below.
This section supplies the infinite-dimensional version actually needed.
The spectrum $\operatorname{spec}P$ consists of those $z\in\C$ for which $P-zI$ has no everywhere-defined bounded inverse on $H$; such an inverse, when it exists, is the resolvent.
A self-adjoint operator has compact resolvent if one (hence every) resolvent is compact.
A Dirichlet realization means the operator associated with the closure of the energy form on compactly supported smooth sections of the domain.

\begin{proposition}[The self-adjoint realization of a closed form]
\label{heat:friedrichs}
Let $q$ be a densely defined, lower-bounded, closed Hermitian form with domain $V\subset H$.
There is a unique self-adjoint operator $P$ satisfying
\[
 \operatorname{Dom}P=\{u\in V:\ q(v,u)=\langle v,f\rangle
                         \text{ for some }f\in H\text{ and every }v\in V\},
 \qquad Pu=f.
\]
If $q(u)\geq-C\|u\|^2$, then $P\geq-C$.
\end{proposition}
\begin{proof}
Set $b=q+(C+1)\langle\cdot,\cdot\rangle$.
For $f\in H$, Lemma~\ref{heat:riesz} gives a unique $Sf\in V$ with
$b(v,Sf)=\langle v,f\rangle$. Thus $S:H\to H$ is bounded, positive, and self-adjoint,
with $\|S\|\leq1$. If $Sf=0$, then $f$ is orthogonal to the dense subspace $V$, so $f=0$.
Also $(\operatorname{Ran}S)^\perp=\ker S=0$, so $\operatorname{Ran}S$ is dense.
On this dense domain define $P=S^{-1}-(C+1)$.
Its domain is exactly the displayed set, and the corresponding $f$ is unique.

It remains to prove self-adjointness.
If $w\in\operatorname{Dom}P^*$ and $P^*w=z$, substitute $u=Sa$ in the adjoint identity
to obtain
\[
 \langle a-(C+1)Sa,w\rangle=\langle Sa,z\rangle\qquad(a\in H).
\]
Since $S=S^*$, it follows that $w=S(z+(C+1)w)$.
Thus $w\in\operatorname{Ran}S=\operatorname{Dom}P$ and $Pw=z$.
This proves $P^*=P$. The lower bound follows from $\langle u,Pu\rangle=q(u)$.
For any other self-adjoint operator associated with the same form, the inverse of $P+C+1$ satisfies the defining identity for $S$,
so its inverse, and therefore the operator itself, is the same.
\end{proof}

\begin{lemma}[Polynomial control of weak convergence]
\label{heat:bernstein}
For every $f\in C([0,1])$, the Bernstein polynomials
\[
 B_Nf(x)=\sum_{j=0}^N f(j/N)\binom Njx^j(1-x)^{N-j}
\]
converge uniformly to $f$.
Consequently, if finite positive measures $\nu_k$ on $[0,1]$ have uniformly bounded total mass,
and the integrals of every monomial converge to those for a finite positive measure $\nu$, then $\nu_k$ converges weakly to $\nu$, meaning convergence of integrals against every continuous function on $[0,1]$.
\end{lemma}
\begin{proof}
Interpret the binomial coefficients as binomial-distribution weights with parameter $x$.
The binomial theorem and its first two derivatives show that $j/N$ has mean $x$ and variance $x(1-x)/N\leq1/(4N)$.
For $\eta>0$, split the sum into $|j/N-x|<\eta$ and its complement.
The first part has error at most the modulus of continuity $\omega_f(\eta)=\sup_{|x-y|\leq\eta}|f(x)-f(y)|$.
The variance bound makes the second part's weight at most $1/(4N\eta^2)$, and its error at most
$2\|f\|_\infty/(4N\eta^2)$.
Let $N\to\infty$, then $\eta\downarrow0$, to obtain uniform convergence.

For the measure assertion, uniformly approximate $f$ within $\epsilon$ by a polynomial $p$.
Convergence of finitely many moments gives $\int p\,d\nu_k\to\int p\,d\nu$.
Replacing $p$ by $f$ costs at most
$\epsilon(\sup_k\nu_k([0,1])+\nu([0,1]))$.
Let $\epsilon\downarrow0$ to conclude weak convergence.
\end{proof}

\begin{lemma}[Borel functional calculus]
\label{heat:borel-calculus}
For a lower-bounded self-adjoint $P$, each bounded Borel function $f$ on its spectrum defines a bounded operator $f(P)$,
satisfying the multiplication, conjugation, and norm rules
\[
 (fg)(P)=f(P)g(P),\quad \overline f(P)=f(P)^*,\quad
 \|f(P)\|\leq\sup|f|.
\]
Uniformly bounded pointwise convergence $f_j\to f$ implies $f_j(P)u\to f(P)u$.
In particular, $e^{-tP}$ is a strongly continuous self-adjoint semigroup; for $P\geq0$,
$\cos(s\sqrt P)$ is defined and has norm at most $1$.
\end{lemma}
\begin{proof}
Choose $c$ so that $A=P+c\geq1$, and put $S=A^{-1}$.
This is an injective bounded self-adjoint operator with $0\leq S\leq1$.
For a polynomial $p$,
\begin{equation}
 \|p(S)\|\leq\max_{0\leq x\leq1}|p(x)|.
 \label{heat:polynomial-calculus}
\end{equation}
Here is a justification of this norm estimate.
The spectrum of a bounded self-adjoint $T$ is real: the inequality
$\|(T-z)u\|\geq|\operatorname{Im}z|\|u\|$ holds, and the orthogonal complement of the range is
$\ker(T-\overline z)=0$. The range is closed and dense, so $T-z$ is invertible.
For $0\leq S\leq1$, the analogous real-part estimate confines its spectrum to $[0,1]$.
Factoring a polynomial into linear factors proves
$\operatorname{spec}p(S)=p(\operatorname{spec}S)$.
Finally, self-adjoint $T$ satisfies $\|T^{2^j}\|=\|T\|^{2^j}$,
while a Cauchy integral of the resolvent around a circle enclosing the spectrum gives
$\|T^m\|\leq C_r r^m$, for $r>\sup|\operatorname{spec}T|$.
Take $m=2^j$ and let $j\to\infty$ to get
$\|T\|=\sup|\operatorname{spec}T|$.
For a polynomial with complex coefficients, apply this to $p(S)^*p(S)$, obtaining
\eqref{heat:polynomial-calculus}.
The Cauchy integral identity follows first for radius greater than $\|T\|$ by termwise integration of the Neumann series,
then by deforming the contour through the holomorphic resolvent. No spectral measure has yet been used.

Lemma~\ref{heat:bernstein} extends $p\mapsto p(S)$ uniquely to $C([0,1])$.
For $f\geq0$, polynomial approximation of $\sqrt f$ shows $f(S)\geq0$.
Thus for each $v\in H$, the positive functional
$f\mapsto\langle v,f(S)v\rangle$ is represented by a finite positive Borel measure $\mu_v$.
This is the Riesz representation theorem for finite measures: define the measure of an open set $U$ first as
\[
 \sup\{\ell(f):0\leq f\leq1,\ f\in C([0,1]),\ \supp f\subset U\},
\]
then define outer measure by approximation from open sets. Disjoint compact sets can be separated by continuous cutoffs, making this outer measure
countably additive on Borel sets. Subdivide the range of a continuous function into intervals of length $1/N$,
use inner and outer approximation, and let $N\to\infty$ to obtain $\ell(f)=\int f\,d\mu_v$.
The same argument gives uniqueness of the measure.

For fixed $v$, associate $p(S)v$ with $p(x)$. The identity
\[
 \|p(S)v\|^2=\int|p(x)|^2\,d\mu_v(x)
\]
identifies the closed cyclic subspace $\overline{\{p(S)v:p\text{ polynomial}\}}$ isometrically with $L^2(\mu_v)$.
Density of continuous functions in $L^2$ of a finite Borel measure follows from the preceding inner and outer approximations.
On this space, $S$ is multiplication by $x$.
Choose cyclic vectors successively in orthogonal complements, decomposing the separable space $H$ into at most countably many such spaces.
Injectivity of $S$ means that none of these measures has an atom at $x=0$.
Thus $P$ is multiplication by $x^{-1}-c$, and for a bounded Borel function $f$,
define $f(P)$ as multiplication by $f(x^{-1}-c)$.
All algebraic rules hold pointwise, and strong convergence is the dominated convergence theorem.
Taking $f(\lambda)=e^{-t\lambda}$ and $f(\lambda)=\cos(s\sqrt\lambda)$
gives the final assertions.
\end{proof}

\paragraph{General trace and Hilbert--Schmidt norms.}
For a bounded operator $T$, Lemma~\ref{heat:borel-calculus} defines
$|T|=(T^*T)^{1/2}$, the positive square root of $T^*T$.
It is trace class if $\|T\|_1=\sum_j\langle e_j,|T|e_j\rangle<\infty$;
then $\operatorname{tr}T=\sum_j\langle e_j,Te_j\rangle$.
It is Hilbert--Schmidt if $\|T\|_{\mathrm{HS}}^2=\sum_j\|Te_j\|^2<\infty$.
These definitions are independent of the orthonormal basis: for a positive operator this follows by applying Parseval and Tonelli to its square root;
the general trace follows by polarization and absolute convergence.
For the positive diagonal operators above, this agrees with the eigenvalue sum already defined.

\paragraph{Kernel notation.}\label{heat:kernel-notation}
For a Hermitian bundle $E\to M$ and a chosen volume measure $dV_g$, an integral kernel $H(t;x,y)$ is a linear map $E_y\to E_x$ satisfying
$(e^{-tP}u)(x)=\int_M H(t;x,y)u(y)\,dV_g(y)$.
Kernel norms below are fiber operator norms unless marked $\mathrm{HS}$, and $\operatorname{tr}_{E_x}H(t;x,x)$ is a finite-dimensional fiber trace.
In local real coordinates $x=(x_1,\ldots,x_d)$, $\partial^\alpha$ denotes a derivative with multi-index $\alpha\in\N_0^d$, $|\alpha|=\sum\alpha_i$.

\Needspace{10\baselineskip}
\begin{proposition}[Smooth kernels of heat operators]
\label{heat:smoothing-kernel}
Let $P$ be a smooth, lower-bounded self-adjoint second-order strongly elliptic operator with scalar principal symbol on a complete manifold,
or its Dirichlet realization on a smooth bounded domain.
For $t>0$, $e^{-tP}$ has a smooth interior kernel $H(t;x,y)$.
On each compact interior set $K$, all derivatives are bounded for $t_0\leq t\leq T$, where $0<t_0<T<\infty$; if $P\geq0$, one may take all $t\geq t_0$.
The bounds are uniform when the spectral lower bound, local coefficient bounds, and ellipticity lower bound are uniform.
\end{proposition}
\begin{proof}
Lemma~\ref{heat:borel-calculus}, with $P\geq-C$ and a fixed $c>C+1$, gives
$\|(P+c)^j e^{-tP}\|\leq\sup_{\lambda\geq-C}(\lambda+c)^j e^{-t\lambda}<\infty$.
Repeated application of Theorem~\ref{ana:interior-estimate} on nested relatively compact coordinate balls gives
\[
 \|e^{-tP}f\|_{H^{2j}(K)}\leq C_{j,K,t_0,T}\|f\|_{L^2}.
\]
Write $d=\dim_{\R}M$ and let $m$ be the desired derivative order. Choose $2j>m+d/2$. Sobolev embedding, Lemma~\ref{ana:sobolev-embedding},
makes $f\mapsto\partial_x^\alpha(e^{-tP}f)(x)$ a bounded $L^2$ functional.
Lemma~\ref{heat:riesz} gives the rows of the kernel and their $x$ derivatives of every order.
Factor the semigroup as $e^{-tP/2}e^{-tP/2}$ and use self-adjointness to express the kernel as an inner product of two $L^2$ rows.
Differentiating each row proves joint smoothness in $x,y$.
For time derivatives use $\partial_t e^{-tP}=-Pe^{-tP}$.
All constants come from the spectral lower bound and the finitely many local elliptic estimates used, giving the stated uniformity.
\end{proof}


\section{Duhamel's formula}
\label{heat:duhamel-section}

\begin{theorem}[Duhamel's formula for bounded operators]
\label{heat:duhamel}
Let $A,B$ be bounded operators on a Banach space. For $t\geq0$,
\begin{align}
 e^{-tA}-e^{-tB}
 &=\int_0^t e^{-(t-s)A}(B-A)e^{-sB}\,ds\label{heat:duhamel-time}\\
 &=t\int_0^1 e^{-(1-u)tA}(B-A)e^{-utB}\,du.
\label{heat:duhamel-unit}
\end{align}
The integral converges in operator norm.
If $A,B$ are self-adjoint on a Hilbert space and
$A\geq aI$, $B\geq bI$, then
\begin{equation}
\|e^{-tA}-e^{-tB}\|
\leq t\|A-B\|e^{-t\min(a,b)}.
\label{heat:duhamel-norm}
\end{equation}
\end{theorem}
\begin{proof}
The exponential series of a bounded operator converges uniformly on bounded time intervals,
as does its derivative series, so termwise differentiation is valid.
Fix $t$ and set $F(s)=e^{-(t-s)A}e^{-sB}$.
Commuting each operator only with its own exponential gives
\[
F'(s)=e^{-(t-s)A}(A-B)e^{-sB}.
\]
The first identity follows from $F(0)=e^{-tA}$, $F(t)=e^{-tB}$, and the fundamental theorem of calculus.
Substitute $s=tu$, $ds=t\,du$, to obtain the second identity; the factor $t$ is essential.
In the self-adjoint case, Lemma~\ref{heat:borel-calculus} gives
$\|e^{-rA}\|\leq e^{-ra}$ and $\|e^{-rB}\|\leq e^{-rb}$.
Take norms under the integral and use
$e^{-(t-s)a-sb}\leq e^{-t\min(a,b)}$ to obtain the last estimate.
\end{proof}

\begin{proposition}[Unbounded operators with bounded difference]
\label{heat:bounded-difference}
Let $A,B$ be lower-bounded self-adjoint operators with common domain $D$,
and suppose $B-A$ on $D$ extends to a bounded operator $V$ on $H$.
Then \eqref{heat:duhamel-time} holds in the strong operator sense,
and \eqref{heat:duhamel-norm} holds with $\|A-B\|$ replaced by $\|V\|$.
\end{proposition}
\begin{proof}
For $u\in D$, Lemma~\ref{heat:borel-calculus} gives $e^{-sB}u\in D$,
and $s\mapsto e^{-sB}u$ is continuous in the graph norm of $D$ and differentiable in $H$.
Since $A=B-V$, the graph norms of $A$ and $B$ are equivalent. Choose lower bounds $A\geq aI$ and $B\geq bI$ for the estimate below.
We may therefore differentiate $F(s)u=e^{-(t-s)A}e^{-sB}u$, obtaining
\[
\frac d{ds}F(s)u=-e^{-(t-s)A}Ve^{-sB}u.
\]
Integrate first over $[\delta,t-\delta]$, then let $\delta\downarrow0$.
Strong continuity of the semigroups supplies the endpoints; boundedness of $V$ and the lower-bound estimates ensure integrability.
This proves the formula for $u\in D$.
The norm of the right-hand side is at most
$t\|V\|e^{-t\min(a,b)}\|u\|$, a bound independent of membership in $D$.
Density of $D$ extends the formula and estimate to all of $H$.
\end{proof}

\begin{remark}
The difference of two second-order differential operators usually still contains first- or second-order derivatives
and is not bounded from $L^2$ to $L^2$.
Proposition~\ref{heat:bounded-difference} then cannot be applied directly.
One must specify a common form domain, check operator domains, or first prove kernel bounds integrable at the time endpoints.
In particular, the factor $s^{-1}$ produced by applying $\nabla^2$ to a Gaussian kernel
is not integrable by itself at $s=0$; the spatial weights or cancellation that resolve this must be explicit.
\end{remark}

\begin{proposition}[Strong solutions with a source term]
\label{heat:inhomogeneous}
Let $A\geq0$ be self-adjoint, let $u_0\in D(A)$,
and let $f:[0,T]\to D(A)$ be continuous in the graph norm.
Then
\begin{equation}
u(t)=e^{-tA}u_0+\int_0^t e^{-(t-s)A}f(s)\,ds
\label{heat:variation-constants}
\end{equation}
belongs to $C([0,T],D(A))\cap C^1([0,T],H)$
and is the unique solution in this class of $u'+Au=f$, $u(0)=u_0$.
\end{proposition}
\begin{proof}
Because $Ae^{-rA}v=e^{-rA}Av$ for $v\in D(A)$,
the integral is defined in the graph norm, with
\[
A\int_0^t e^{-(t-s)A}f(s)\,ds
=\int_0^t e^{-(t-s)A}Af(s)\,ds.
\]
Graph-norm continuity and strong continuity of the semigroup make these integrals continuous in $t$.
Take difference quotients in \eqref{heat:variation-constants},
separating the new short-time integral from the difference on the original integration interval.
Their limits are respectively $f(t)$ and $-\int_0^t e^{-(t-s)A}Af(s)\,ds$.
Thus $u'=f-Au$, and the initial condition is immediate.
For the difference $w$ of two solutions,
\[
\frac d{dt}\|w(t)\|^2=-2\langle Aw(t),w(t)\rangle\leq0.
\]
Since $w(0)=0$, this implies $w\equiv0$.
\end{proof}

\begin{proposition}[A heat-trace perturbation estimate]
\label{heat:trace-perturbation}
Under the hypotheses of Proposition~\ref{heat:bounded-difference}, assume further that
$A,B\geq0$ and both heat operators are trace class at every positive time.
Then
\begin{equation}
\|e^{-tA}-e^{-tB}\|_1
\leq\frac t2\|V\|
\bigl(\tr(e^{-tA/2})+\tr(e^{-tB/2})\bigr).
\label{heat:trace-norm-estimate}
\end{equation}
The same right-hand side therefore bounds the absolute difference of the two heat traces.
Here $\|\cdot\|_1$ denotes the trace norm.
\end{proposition}
\begin{proof}
If $P\geq0$ is trace class and $R$ is bounded, then
$\|PR\|_1,\|RP\|_1\leq\tr(P)\|R\|$.
Indeed, in the spectral expansion $P=\sum_jp_j e_j\otimes e_j^*$, where $(e_j\otimes e_j^*)u=e_j\langle e_j,u\rangle$,
multiplication of each rank-one term by $R$ gives trace norm at most $p_j\|R\|$;
apply the triangle inequality and $\sum_jp_j<\infty$.
On the part $0\leq s\leq t/2$ of the Duhamel integral,
\[
\|e^{-(t-s)A}Ve^{-sB}\|_1
\leq\|V\|\tr(e^{-(t-s)A})
\leq\|V\|\tr(e^{-tA/2}).
\]
On $t/2\leq s\leq t$, use $e^{-sB}$ instead as the trace-class factor,
giving the bound $\|V\|\tr(e^{-tB/2})$.
To verify existence of the integral in trace norm, write for $s\in[\delta,t-\delta]$
\[
 e^{-(t-s)A}Ve^{-sB}
 =e^{-(t-s-\delta/2)A}\bigl(e^{-\delta A/2}V\bigr)e^{-sB}.
\]
The middle factor is a fixed trace-class operator, and the outer semigroups are norm-continuous at positive times,
so the integrand is trace-norm continuous. The two uniform bounds above make the omitted endpoint integrals tend to zero in trace norm.
The trace-norm integral therefore exists and agrees with the strong Duhamel integral; integrating proves the estimate.
Finally, use $|\tr T|\leq\|T\|_1$.
\end{proof}

\section{Why heat-trace comparisons do not imply pointwise comparisons}
\label{heat:order-section}

\begin{proposition}[Exponentiation does not preserve self-adjoint operator order]
\label{heat:nonmonotone}
There exist positive semidefinite $2\times2$ Hermitian matrices $A\leq B$
for which $e^{-B}\leq e^{-A}$ fails.
\end{proposition}
\begin{proof}
Take
\[
A=\begin{pmatrix}0&0\\0&2\end{pmatrix},\qquad
P=\begin{pmatrix}1&1\\1&1\end{pmatrix},\qquad B_\epsilon=A+\epsilon P.
\]
Since $P=vv^*$ with $v=(1,1)^\mathsf T$, we have $P\geq0$ and $B_\epsilon\geq A$.
Theorem~\ref{heat:duhamel} gives the operator-norm limit
\[
\lim_{\epsilon\downarrow0}
\frac{e^{-A}-e^{-B_\epsilon}}{\epsilon}
=\int_0^1 e^{-(1-s)A}Pe^{-sA}\,ds
=\begin{pmatrix}1&c\\c&e^{-2}\end{pmatrix}=:D,
\quad c=\frac{1-e^{-2}}2.
\]
Since $\sinh1>1$, by its power series or integration of $\cosh t>1$,
we have $c=e^{-1}\sinh1>e^{-1}$.
For $w=(-c,1)^\mathsf T$,
$w^*Dw=e^{-2}-c^2<0$.
The norm convergence implies that, for sufficiently small $\epsilon>0$,
$w^*(e^{-A}-e^{-B_\epsilon})w<0$.
This disproves the proposed operator-order inequality.
\end{proof}

\begin{remark}
Corollary~\ref{heat:form-comparison} remains valid because it compares eigenvalues term by term.
The error would be to bypass eigenvalues and exponentiate a matrix or operator inequality directly.
Commuting self-adjoint matrices can be simultaneously unitarily diagonalized, and scalar monotonicity then does give operator comparison;
Proposition~\ref{heat:nonmonotone} shows why commutativity cannot simply be dropped.
For vector-bundle heat kernels, one also cannot substitute $\delta_x$ into an operator inequality as though it were an ordinary $L^2$ vector.
A point mass is not in $L^2$, and distributional convergence alone does not imply convergence of the associated quadratic-form values.
\end{remark}

\begin{proposition}[Integrating the heat-kernel diagonal gives the heat trace]
\label{heat:kernel-trace}
Let $X$ carry a finite Borel measure $\mu$ with full support,
and let $E\to X$ be a Hermitian vector bundle of finite rank.
Suppose a self-adjoint semigroup $T_t$ has a continuous kernel $K_{t/2}(x,y)$ at time $t/2$,
and $X$ is compact. Then $T_t=T_{t/2}T_{t/2}^*$ is trace class,
its continuous kernel satisfies
\[
K_t(x,x)=\int_X K_{t/2}(x,z)K_{t/2}(x,z)^*\,d\mu(z),
\]
and
\begin{equation}
\tr(T_t)=\int_X\tr_{E_x}K_t(x,x)\,d\mu(x).
\label{heat:trace-integral}
\end{equation}
\end{proposition}
\begin{proof}
Compactness and continuity make the squared kernel norm integrable, so $T_{t/2}$ is Hilbert--Schmidt.
Choose an orthonormal basis $(e_j)$ of $L^2(X,E)$. Parseval and Tonelli give
\[
\sum_j\|T_{t/2}e_j\|^2
=\int_X\int_X\|K_{t/2}(x,z)\|_{\mathrm{HS}}^2
\,d\mu(z)\,d\mu(x)<\infty.
\]
Thus $T_{t/2}T_{t/2}^*$ is trace class, with trace equal to this squared Hilbert--Schmidt norm.
The kernel composition formula follows by applying the iterated integral to continuous sections and using Fubini;
self-adjointness gives $K_{t/2}(z,x)=K_{t/2}(x,z)^*$.
Taking the fiber trace on the diagonal and integrating gives exactly the squared norm above.
\end{proof}

Formula \eqref{heat:trace-integral} distinguishes three objects:
$K_t(x,y)$ is a kernel, $\tr_{E_x}K_t(x,x)$ a local density,
and $\tr(T_t)$ a number.
The local density generally depends on $x$; the heat trace integrates it over all of $X$.


\section{Partitions of unity, the IMS identity, and spectral localization}

\begin{proposition}[The IMS identity]
\label{heat:ims}
On a Riemannian manifold $X$ with Hermitian bundle $E$, let $\nabla$ be unitary and $P_k=k^{-1}\nabla^*\nabla+Q$, with $k>0$ and $Q=Q^*$ lower bounded,
and let $\chi_1,\ldots,\chi_N$ be real smooth functions with $\sum_\nu\chi_\nu^2=1$.
Write $q_k$ for its closed quadratic form, and take $u$ in its form domain. Require the finitely many $d\chi_\nu$ to be bounded, so multiplication by each $\chi_\nu$ preserves that domain. Then
\[
 \sum_\nu q_k(\chi_\nu u)
 =q_k(u)+k^{-1}\int\sum_\nu|d\chi_\nu|^2|u|^2.
\]
Assume compact resolvent for the realizations in the following counting statements. If each $\chi_\nu$ is supported inside an open set $\Omega_\nu$, then
\begin{equation}
 N_X(\lambda)\leq\sum_\nu N_{\Omega_\nu}^{D}
       \left(\lambda+k^{-1}\Big\|\sum_\nu|d\chi_\nu|^2\Big\|_\infty\right).
 \label{heat:ims-counting}
\end{equation}
Here $N_X(\lambda)$ counts global eigenvalues at most $\lambda$, with multiplicity, and $N_{\Omega_\nu}^{D}$ is the corresponding count for the Dirichlet realization on $\Omega_\nu$.
For pairwise disjoint $\Omega_\nu$, one also has
$\sum_\nu N_{\Omega_\nu}^{D}(\lambda)\leq N_X(\lambda)$.
\end{proposition}
\begin{proof}
Expand $\nabla(\chi_\nu u)=d\chi_\nu\otimes u+\chi_\nu\nabla u$.
The sum of mixed terms contains $\sum_\nu\chi_\nu d\chi_\nu
=\tfrac12d(\sum_\nu\chi_\nu^2)=0$, and the potential terms sum to the original potential term.
This proves the identity first for compactly supported smooth sections, then by density in the form norm.

Let $E$ be the global spectral space for eigenvalues at most $\lambda$, and let $\delta$ be the error in the formula.
Consider the map from $E$ to the direct sum of the local spectral spaces with eigenvalues at most $\lambda+\delta$,
$u\mapsto(\Pi_\nu\chi_\nu u)_\nu$, where $\Pi_\nu$ is the orthogonal projection onto that local spectral space.
If it had a nonzero kernel vector, each nonzero $\chi_\nu u$ would contain only spectral components above $\lambda+\delta$, giving
\[
 \sum_\nu q_k(\chi_\nu u)>(\lambda+\delta)
                 \sum_\nu\|\chi_\nu u\|^2
 =(\lambda+\delta)\|u\|^2.
\]
The IMS identity and $u\in E$ give the opposite weak inequality, a contradiction.
The map is therefore injective, and dimension comparison proves \eqref{heat:ims-counting}.
For disjoint open sets, extend the sections in each Dirichlet spectral space by zero.
Their direct sum lies in the global form domain and has Rayleigh quotient $q_k(u)/\|u\|^2$ at most $\lambda$. Theorem~\ref{heat:minmax} gives the lower bound.
\end{proof}

\begin{remark}
IMS controls the dimensions of spectral spaces; finite propagation, developed in Chapter~\ref{heat:local-chapter}, controls the interior heat kernel.
These conclusions concern different objects: the spectral comparison \eqref{heat:ims-counting} does not give a pointwise order relation between two matrix-valued heat kernels.
\end{remark}


\chapter{Gaussian Kernels, Radial Integration, and Constant Magnetic Fields}
\label{heat:models-chapter}

The Euclidean Gaussian, the constant magnetic field, and the flat torus provide explicit heat kernels.
We compute them here and use them as local comparison models in the next chapter. The magnetic-field sign, volume normalization, and factor $2$ in heat time
will together determine the final factor $(2\pi)^{-n}$ in the holomorphic Morse inequalities.
Here $m$ denotes Euclidean real dimension, $k>0$ the scaling parameter, $t>0$ heat time, and Euclidean integration is with respect to Lebesgue measure.
Kernel and operator norms follow the conventions in Sections~\ref{heat:spectral-section} and \ref{heat:general-calculus-section}.

\section{Euclidean Gaussian kernels and integrable errors}
\label{heat:gaussian-section}

\begin{lemma}[Gaussian integrals and convolution]
\label{heat:gaussian-convolution}
Let $m\geq1$, $k>0$, $t>0$, and set
\[
G_{t,k}(z)=\left(\frac{k}{4\pi t}\right)^{m/2}
 \exp\left(-\frac{k|z|^2}{4t}\right).
\]
Then $\int_{\R^m}G_{t,k}=1$, and for $s,t>0$,
\begin{equation}
\int_{\R^m}G_{t,k}(x-z)G_{s,k}(z-y)\,dz
=G_{t+s,k}(x-y).
\label{heat:gaussian-semigroup}
\end{equation}
\end{lemma}
\begin{proof}
The square of the one-dimensional integral $I=\int_\R e^{-u^2}\,du$ is
\[
I^2=\int_{\R^2}e^{-(x^2+y^2)}\,dx\,dy
=2\pi\int_0^\infty re^{-r^2}\,dr=\pi.
\]
Since $I>0$, we obtain $I=\sqrt\pi$.
Integration in each coordinate followed by rescaling gives the normalization.
For convolution, complete the square:
\[
\frac{|x-z|^2}{t}+\frac{|z-y|^2}{s}
=\frac{|x-y|^2}{t+s}
 +\frac{t+s}{ts}
 \left|z-\frac{sx+ty}{t+s}\right|^2.
\]
Substitute this identity into the integral. The second term gives an explicit Gaussian integral;
its factor $(4\pi ts/(k(t+s)))^{m/2}$ combines with the prefactors
to give exactly $(k/(4\pi(t+s)))^{m/2}$.
\end{proof}

\begin{proposition}[The Euclidean heat equation]
\label{heat:euclidean-kernel}
With $G_{t,k}$ from Lemma~\ref{heat:gaussian-convolution}, for $f\in C_c^\infty(\R^m)$ define
\[
u(t,x)=\int_{\R^m}G_{t,k}(x-y)f(y)\,dy.
\]
Here $\Delta=\sum_{j=1}^m\partial_{x_j}^2$, so $-\Delta$ is nonnegative. Then $u_t=k^{-1}\Delta u$, and $u(t,\cdot)\to f$ in $L^2$
and uniformly on compact sets.
Among functions in $C([0,T],L^2)\cap C^1((0,T],L^2)$ satisfying
$u(t)\in H^2(\R^m)$ at every positive time, this equation and initial value have a unique solution.
Here $H^2$ consists of functions whose distributional derivatives through order two are square-integrable.
\end{proposition}
\begin{proof}
Direct differentiation gives
\[
\partial_tG_{t,k}
=\left(-\frac m{2t}+\frac{k|z|^2}{4t^2}\right)G_{t,k}
=\frac1k\Delta_zG_{t,k}.
\]
At positive time, Gaussian decay permits every required derivative under the integral,
so $u$ satisfies the equation.
Write
$u(t,x)-f(x)=\int G_{t,k}(z)(f(x-z)-f(x))\,dz$.
Given $\epsilon>0$, uniform continuity of $f$ supplies $\delta>0$ such that
the absolute difference is at most $\epsilon$ for $|z|<\delta$.
This part of the integral is at most $\epsilon$; the rest is bounded by
$2\|f\|_\infty\int_{|z|\geq\delta}G_{t,k}(z)\,dz$,
which tends to zero after the substitution $z=\sqrt{t/k}\,w$.
The same translation representation proves $L^2$ convergence:
\[
\|G_{t,k}*f-f\|_2
\leq\int G_{t,k}(z)\|f(\cdot-z)-f\|_2\,dz.
\]
For small $z$, use continuity of translations in $L^2$; for large $z$, use $2\|f\|_2$ and the Gaussian tail.
If $w$ is the difference of two solutions, integration by parts gives
\[
\frac d{dt}\|w(t)\|_2^2=-\frac2k\|\nabla w(t)\|_2^2\leq0.
\]
Integrate from $0<\epsilon<t$ and let $\epsilon\downarrow0$.
Continuity of the initial value in $L^2$ gives $w(t)=0$.
Integration by parts in $H^2$ follows by smooth compactly supported approximation and passage to the limit in Sobolev norm.
\end{proof}

\begin{proposition}[Gaussian bounds for first and second derivatives]
\label{heat:gaussian-derivatives}
Set $\theta=k|z|^2/t$, and use the Hilbert--Schmidt norm $|\nabla^2G|^2=\sum_{i,j}|\partial_i\partial_jG|^2$ for the Hessian.
Then
\begin{align*}
|\nabla G_{t,k}(z)|
&\leq(4\pi)^{-m/2}e^{-1/2}
       \left(\frac kt\right)^{(m+1)/2}e^{-\theta/8},\\
|\nabla^2G_{t,k}(z)|
&\leq(4\pi)^{-m/2}\left(\frac{\sqrt m}{2}+\frac2e\right)
       \left(\frac kt\right)^{(m+2)/2}e^{-\theta/8}.
\end{align*}
For every real $a>-m$, one also has
\begin{equation}
\int_{\R^m}|z|^aG_{t,k}(z)\,dz
=2^a\frac{\Gamma((m+a)/2)}{\Gamma(m/2)}
\left(\frac tk\right)^{a/2},
\label{heat:gaussian-moment}
\end{equation}
where $\Gamma(s)=\int_0^\infty r^{s-1}e^{-r}\,dr$.
\end{proposition}
\begin{proof}
The exact derivatives are
\[
\partial_iG=-\frac{kz_i}{2t}G,\qquad
\partial_i\partial_jG
=\left(\frac{k^2z_iz_j}{4t^2}-\frac{k\delta_{ij}}{2t}\right)G.
\]
Hence
\[
|\nabla G|=\frac12\left(\frac kt\right)^{1/2}\theta^{1/2}G,
\qquad
|\nabla^2G|\leq\frac kt\left(\frac\theta4+\frac{\sqrt m}{2}\right)G.
\]
Split $e^{-\theta/4}$ into two factors $e^{-\theta/8}$ and use
\[
\sup_{\theta\geq0}\theta^{1/2}e^{-\theta/8}=2e^{-1/2},
\qquad
\sup_{\theta\geq0}\theta e^{-\theta/8}=8/e,
\]
to obtain the derivative bounds.
For the moment formula, use polar coordinates and $r=2\sqrt{t/k}\sqrt s$.
The sphere area $2\pi^{m/2}/\Gamma(m/2)$ follows by comparing the same Gaussian integral in Cartesian
and polar coordinates. Substitution gives \eqref{heat:gaussian-moment}.
The condition $a>-m$ is exactly integrability at the origin.
\end{proof}

\begin{lemma}[An integrable Duhamel convolution]
\label{heat:integrable-error}
Suppose scalar or matrix kernels $H,E$ satisfy, for $0<s<t$,
\[
\|H(t-s,x,z)\|\leq C_0G_{c(t-s),k}(x-z),
\qquad
\|E(s,z,y)\|\leq C_1k^{-1/2}s^{-1/2}G_{cs,k}(z-y),
\]
where $c,C_0,C_1>0$ are independent of $s,x,y,z$.
Then the double integral converges absolutely, and
\begin{equation}
\left\|\int_0^t\int_{\R^m}H(t-s,x,z)E(s,z,y)\,dz\,ds\right\|
\leq2C_0C_1k^{-1/2}\sqrt t\,G_{ct,k}(x-y).
\label{heat:integrable-duhamel-bound}
\end{equation}
\end{lemma}
\begin{proof}
Use submultiplicativity of the matrix norm and then
\eqref{heat:gaussian-semigroup}; the spatial integral is bounded by
$C_0C_1k^{-1/2}s^{-1/2}G_{ct,k}(x-y)$.
Finally, $\int_0^t s^{-1/2}\,ds=2\sqrt t$.
This nonnegative integrable majorant proves both absolute convergence and validity of interchanging the integrals.
\end{proof}

In applications, the bounds for $H$ and $E$ must be proved from the operator under study, with constants independent of $k$.
The time power is decisive: $s^{-1/2}$ is integrable at zero, whereas $s^{-1}$ is not.
This distinction explains why the next chapter distributes derivatives between the two kernels before applying Duhamel's formula.

\section{Radial integration and Stieltjes integration by parts}
\label{heat:stieltjes-section}

Singularities at the origin, growth of mass in balls, and Gaussian decay often occur together.
The following method uses only a pushforward measure; it requires neither a density for the measure nor differentiability of its ball-mass function.
This is the form needed to extend ordinary polar integration to mass measures of positive currents.

\begin{theorem}[Integration by parts for mass in balls]
\label{heat:stieltjes}
Let $(Y,d)$ be a metric space, $x_0\in Y$, and $R>0$,
and let $\mu$ be a finite positive Borel measure on $\overline B(x_0,R)$.
Write
\[
F(r)=\mu\bigl(\{x:d(x,x_0)\leq r\}\bigr),\qquad0\leq r\leq R.
\]
For every $f\in C^1([0,R])$,
\begin{equation}
\int_{\overline B(x_0,R)}f(d(x,x_0))\,d\mu(x)
=f(R)F(R)-\int_0^R f'(r)F(r)\,dr.
\label{heat:stieltjes-closed}
\end{equation}
For integration over the open ball $B(x_0,R)$, replace $F(R)$ by
$F(R-)=\mu(B(x_0,R))$; the integral for $r<R$ is unchanged.
\end{theorem}
\begin{proof}
The pushforward under the distance function $\rho(x)=d(x,x_0)$,
$\nu=\rho_*\mu$, is a measure on $[0,R]$ satisfying
$\nu([0,r])=F(r)$.
By definition of pushforward, the integral is $\int_{[0,R]}f(s)\,d\nu(s)$.
The fundamental theorem of calculus gives
\[
f(s)=f(R)-\int_s^R f'(r)\,dr.
\]
Substitute this identity. Since $|f'|$ is bounded and $\nu$ finite, the absolute double integral is at most
$R\|f'\|_\infty\nu([0,R])$, so Fubini applies.
For fixed $r$, the mass of $s\leq r$ is $F(r)$.
Including or excluding equality changes only a $dr$-null set of radii with atoms;
equivalently, the diagonal has zero measure by integrating its singleton $dr$ sections.
This proves \eqref{heat:stieltjes-closed}.

For the open ball, repeat the argument with the pushforward measure on $[0,R)$.
Its total mass in the constant term is $F(R-)$, while the interior integral is unchanged.
Thus open and closed balls cannot be interchanged without accounting for mass on their boundary sphere.
\end{proof}

\begin{remark}[An atom at the origin]
Formula \eqref{heat:stieltjes-closed} already includes $f(0)\mu(\{x_0\})$.
For example, if $\mu=a\delta_{x_0}$, then $F(r)=a$ for every $r\geq0$,
and the right-hand side is $af(R)-a(f(R)-f(0))=af(0)$.
A Stieltjes notation integrating only over $(0,R]$ must add this atomic term separately.
The same proof applies to finite signed measures, with domination by total variation.
\end{remark}

\begin{corollary}[Volume growth gives a Gaussian integral bound]
\label{heat:radial-gaussian}
Let $\mu$ be a finite positive measure on a closed ball of finite radius centered at $x_0$.
If $\delta>0$, $C>0$ satisfy
$F(r)\leq Cr^\delta$ for $0<r\leq R$, then for every $\lambda>0$,
\begin{equation}
\int_{\overline B(x_0,R)}e^{-\lambda d(x,x_0)^2}\,d\mu(x)
\leq C\Gamma(1+\delta/2)\lambda^{-\delta/2}.
\label{heat:radial-gaussian-bound}
\end{equation}
\end{corollary}
\begin{proof}
Extend the measure by zero outside the ball and set $F(r)=F(R)$ for $r>R$.
Since $\delta>0$, the extension still satisfies $F(r)\leq Cr^\delta$, and $F(0)=0$.
Apply Theorem~\ref{heat:stieltjes} at radius $S>R$ with
$f(r)=e^{-\lambda r^2}$, and let $S\to\infty$.
The boundary term tends to zero, giving
\[
\int e^{-\lambda\rho^2}\,d\mu
=2\lambda\int_0^\infty r e^{-\lambda r^2}F(r)\,dr
\leq2\lambda C\int_0^\infty r^{\delta+1}e^{-\lambda r^2}\,dr.
\]
The substitution $s=\lambda r^2$ makes the right-hand side
$C\lambda^{-\delta/2}\int_0^\infty s^{\delta/2}e^{-s}\,ds$,
which is the claimed bound.
\end{proof}

The exponent $\delta$ is a mass-growth exponent and need not equal the ambient real dimension.
For the volume measure of a smooth $2p$-dimensional submanifold, for example, one may locally take $\delta=2p$.
For the mass of a positive current, the required growth bound must first follow from the appropriate monotonicity or local estimate.
Ambient dimension alone does not determine $\delta$.

\section{The exact model for a signed constant magnetic field}
\label{heat:magnetic-section}

For $x=(x_1,x_2)\in\R^2$, $k>0$, and $b\in\R$, consider
\[
D_1=\partial_{x_1},\qquad D_2=\partial_{x_2}-\ii kb x_1,
\qquad
P_{b,k}=-\frac1k(D_1^2+D_2^2).
\]
On compactly supported smooth functions,
$\langle P_{b,k}u,u\rangle=k^{-1}(\|D_1u\|^2+\|D_2u\|^2)\geq0$.
This is the constant-field model in a fixed gauge. The magnetic strength is the signed real number $b$.
A Schwartz function is a smooth function $f$ for which $\sup_x|x^\alpha\partial^\beta f(x)|<\infty$ for every pair of multi-indices; the set of such functions is denoted $\mathcal S(\R^2)$.

\begin{theorem}[The two-dimensional magnetic heat kernel]
\label{heat:magnetic-kernel}
For $b\ne0$, the heat kernel of $P_{b,k}$ is
\begin{align}
K_{b,k}(t;x,y)
={}&\frac{kb}{4\pi\sinh(bt)}
\exp\left(-\frac{kb}{4}\coth(bt)|x-y|^2\right)\notag\\
&\quad\cdot
\exp\left(\frac{\ii kb}{2}(x_1+y_1)(x_2-y_2)\right).
\label{heat:magnetic-formula}
\end{align}
At $b=0$ it extends continuously to the two-dimensional Gaussian kernel $G_{t,k}(x-y)$.
The formula holds for either sign of $b$, and
\begin{equation}
|K_{b,k}(t;x,y)|\leq G_{t,k}(x-y).
\label{heat:constant-diamagnetic}
\end{equation}
More precisely, this kernel gives the unique Schwartz-class strong solution for Schwartz initial data
and extends on $L^2$ to a strongly continuous self-adjoint contraction semigroup.
This semigroup specifies the realization of the heat operator in this section.
\end{theorem}
\begin{proof}
First verify the differential equation directly.
Write $r=x-y$ and
\[
a(t)=\frac{kb}{4}\coth(bt),\quad
c(t)=\frac{kb}{4\pi\sinh(bt)},\quad
\phi(x,y)=\frac{kb}{2}(x_1+y_1)(x_2-y_2).
\]
so $K=c(t)e^{-a(t)|r|^2}e^{\ii\phi}$.
Calculation gives
\[
\frac{D_1K}{K}=-2ar_1+\frac{\ii kb}{2}r_2,
\qquad
\frac{D_2K}{K}=-2ar_2-\frac{\ii kb}{2}r_1.
\]
Differentiate once more in each direction. The imaginary mixed terms cancel, giving
\[
\frac{(D_1^2+D_2^2)K}{K}
=-4a+\left(4a^2-\frac{k^2b^2}{4}\right)|r|^2.
\]
On the other hand,
\[
\frac{c'}c=-b\coth(bt)=-\frac{4a}{k},\qquad
 a'=\frac{kb^2}{4}-\frac{4a^2}{k}.
\]
Thus $\partial_tK+P_{b,k,x}K=0$.

For every real $v$, with continuous interpretation at zero,
$v/\sinh v$ is positive and even, and at most $1$;
$v\coth v$ is positive and even, and at least $1$.
Hence $c(t)\leq k/(4\pi t)$ and $a(t)\geq k/(4t)$,
proving \eqref{heat:constant-diamagnetic}.
As $b\to0$, we have
\[
\frac{b}{\sinh(bt)}\to\frac1t,
\qquad b\coth(bt)\to\frac1t,
\qquad e^{\ii\phi(x,y)}\to1,
\]
giving the zero-field limit.

Now verify the initial condition and the operator interpretation.
For $f\in C_c^\infty$, set $T_tf(x)=\int K(t;x,y)f(y)\,dy$.
After the change of variables $y=x+\sqrt{t/k}\,z$, the amplitude tends to the standard Gaussian,
and the phase tends to $1$ for fixed $x,z$.
The Gaussian bound and dominated convergence give
$T_tf\to f$ uniformly on compact sets.
Away from $\supp f$, the same Gaussian estimate supplies uniformly integrable $L^2$ tails,
so convergence also holds in $L^2$.
Moreover, $|T_tf|\leq G_{t,k}*|f|$, and the convolution estimate $\|a*b\|_2\leq\|a\|_1\|b\|_2$ gives
$\|T_tf\|_2\leq\|f\|_2$, since $\|G_{t,k}\|_1=1$ by Lemma~\ref{heat:gaussian-convolution}.
Density of $C_c^\infty$ extends $T_t$ to $L^2$,
preserving strong continuity and contractivity.

For Schwartz functions, the integral may be differentiated arbitrarily and remains Schwartz.
Indeed, each kernel derivative is a polynomial in $x,y,x-y$ times the same Gaussian-type kernel.
Using $|x|^N\leq C_N(|x-y|^N+|y|^N)$,
absorb powers of $y$ into the rapid decay of $f$ and powers of $x-y$ into Gaussian moments.
This bounds every Schwartz seminorm and justifies integration by parts for these solutions.
If two such solutions have the same initial value, their difference $w$ satisfies
\[
\frac d{dt}\|w\|_2^2
=-\frac2k(\|D_1w\|_2^2+\|D_2w\|_2^2)\leq0,
\]
so uniqueness follows. This also proves $T_tT_s=T_{t+s}$: as functions of $t$, both sides
satisfy the same equation with initial value $T_sf$. Density extends the identity to $L^2$.
The kernel identity $K(t;x,y)=\overline{K(t;y,x)}$ makes the semigroup self-adjoint.
Finally, kernel symmetry and integration by parts in $y$ give
$\partial_tT_tf=-T_t(P_{b,k}f)$ for $f\in C_c^\infty$.
Since $T_t(P_{b,k}f)\to P_{b,k}f$ in $L^2$,
the strong derivative at $t=0$ is $-P_{b,k}f$.
The constructed semigroup therefore has the stated differential operator as its negative generator on compactly supported smooth functions.
\end{proof}

\begin{remark}[Absolute values are allowed only in the amplitude]
Both $b/\sinh(bt)$ and $b\coth(bt)$ are even functions of $b$,
but $b$ in the phase cannot be replaced by $|b|$.
In fact, $K_{-b,k}(t;x,y)=\overline{K_{b,k}(t;x,y)}$.
When $b$ may pass through zero, the continuous extensions of
$bt/\sinh(bt)$ and $bt\coth(bt)$
are better suited to uniform estimates than a case distinction by the rank of the nonzero magnetic field.
\end{remark}

\begin{proposition}[Gauge transformations]
\label{heat:gauge}
On a Riemannian manifold with its volume measure, let $D_A=d-\ii kA$ on the trivial Hermitian line bundle and $P_A=k^{-1}D_A^*D_A$, where $A$ is a real one-form and $D_A^*$ is the formal $L^2$ adjoint;
take the self-adjoint domains to correspond under the unitary map below.
If $A'=A+d\psi$ with $\psi$ real and smooth, multiplication by
$Uu=e^{\ii k\psi}u$ is unitary and satisfies
\[
D_{A'}U=UD_A,\qquad P_{A'}=UP_AU^{-1}.
\]
When the heat kernels exist,
\[
K_{A'}(t;x,y)=e^{\ii k\psi(x)}K_A(t;x,y)e^{-\ii k\psi(y)}.
\]
Thus the diagonal kernel is gauge invariant for a line bundle; for vector bundles its fiber trace is also invariant under unitary gauge changes.
\end{proposition}
\begin{proof}
The Leibniz rule gives
\[
(d-\ii kA-\ii k\,d\psi)(e^{\ii k\psi}u)
=e^{\ii k\psi}(du-\ii kAu).
\]
Take adjoints and compose to obtain unitary conjugacy of the Laplace operators.
The exponential series for bounded operators, or Lemma~\ref{heat:borel-calculus} for the self-adjoint realizations, gives the same conjugacy of heat operators.
Writing this as an integral kernel proves the formula.
For $x=y$ the scalar factors cancel; in the matrix case use invariance of trace under conjugation.
\end{proof}

\begin{corollary}[Several magnetic planes and a constant matrix potential]
\label{heat:product-model}
Decompose $\R^m$ into $s$ two-dimensional coordinate planes and $m-2s$ free directions,
with magnetic parameters $b_1,\ldots,b_s\in\R$.
Let $V$ be a fixed $r\times r$ Hermitian matrix,
and let $P$ be the sum of the magnetic operators and the free operators $-k^{-1}\partial_j^2$, minus $V$.
Its heat kernel is the product of the scalar kernels, multiplied by $e^{tV}$.
On the diagonal,
\begin{equation}
K_P(t;x,x)
=\left(\frac{k}{4\pi t}\right)^{m/2}
 \prod_{j=1}^s\frac{b_jt}{\sinh(b_jt)}\ e^{tV},
\label{heat:model-diagonal}
\end{equation}
with each zero-parameter factor interpreted as $1$.
\end{corollary}
\begin{proof}
The scalar operators act on disjoint coordinate variables and therefore commute.
The time derivative of their product kernel gives their sum by the Leibniz rule;
the initial condition follows from the approximate-identity property in each coordinate and Fubini's theorem.
The constant matrix $V$ commutes with these scalar operators,
and $\partial_te^{tV}=Ve^{tV}$ corresponds precisely to the potential $-V$.
Uniqueness also follows by energy: the difference $w$ of two solutions satisfies
$\frac d{dt}\|w\|^2\leq2\|V\|\|w\|^2$,
so $e^{-2\|V\|t}\|w(t)\|^2$ is nonincreasing; zero initial value forces $w=0$.
The constructed kernel is therefore the required heat kernel.
Multiply the diagonal values from Theorem~\ref{heat:magnetic-kernel},
$kb_j/(4\pi\sinh(b_jt))$, by the free-direction factors
$(k/(4\pi t))^{1/2}$ to obtain \eqref{heat:model-diagonal}.
\end{proof}

If $V$ depends on $x$, this factorization generally fails.
Using $e^{tV(x_0)}$ in a local asymptotic means first fixing $x_0$ and freezing coefficients;
the freezing error must still be controlled separately.

\section{The flat complex torus}
\label{heat:torus-section}

Let
\[
X=\C^n/(\Z^n+\ii\Z^n),\qquad z_j=x_j+\ii y_j.
\]
Although the $z_j$ are not global functions, the forms $dz_j,d\bar z_j$ are invariant under lattice translations
and hence are global one-forms.
The following computation uses the Dolbeault complex of Chapter~\ref{sh:chapter} directly; the general elliptic Hodge theorem is not needed.

\begin{proposition}[Dolbeault cohomology of the complex torus]
\label{heat:torus-cohomology}
For the complex torus above, write $h^{p,q}(X)=\dim_{\C}H^{p,q}_{\dbar}(X)$ and, for increasing index sets $I,J$, write $dz_I$ and $d\bar z_J$ for the corresponding wedge products. Then
\[
H_{\dbar}^{p,q}(X)
=\operatorname{span}_{\C}
\{[dz_I\wedge d\bar z_J]:|I|=p,\ |J|=q\},
\qquad
h^{p,q}(X)=\binom np\binom nq.
\]
\end{proposition}
\begin{proof}
A smooth periodic function has a Fourier expansion
\[
f(x,y)=\sum_{(a,b)\in\Z^n\times\Z^n}
 f_{a,b}e_{a,b}(x,y),\qquad
 e_{a,b}=e^{2\pi\ii(a\cdot x+b\cdot y)}.
\]
Repeated integration by parts in $x_j,y_j$ shows that its coefficients decay faster than every power.
The series and all its derivatives therefore converge absolutely and uniformly.
To verify that the Fourier expansion equals the original function, use the multivariable Fej\'er kernels as an approximate identity.
Uniform continuity makes their convolutions converge uniformly to $f$. If all Fourier coefficients vanish,
these convolutions vanish identically, so $f=0$.
This also justifies the termwise operations below.

For each frequency $m=(a,b)$,
\[
\dbar(e_m\eta)=e_m\gamma_m\wedge\eta,
\qquad
\gamma_m=\pi\sum_{j=1}^n(\ii a_j-b_j)d\bar z_j,
\]
where $\eta$ is a constant-coefficient $(0,q)$-form.
For $m\ne0$, one has $\gamma_m\ne0$.
Let $\iota_{\gamma_m}$ be the contraction adjoint to exterior multiplication by $\gamma_m$.
In an orthonormal basis with first vector $\gamma_m/|\gamma_m|$, direct verification gives
\[
\gamma_m\wedge\iota_{\gamma_m}
+\iota_{\gamma_m}(\gamma_m\wedge\cdot)
=|\gamma_m|^2\operatorname{Id}.
\]
The complex at this frequency therefore has contracting homotopy
$h_m=|\gamma_m|^{-2}\iota_{\gamma_m}$.
If $u$ is $\dbar$-closed, every nonzero frequency satisfies
$u_m=\gamma_m\wedge h_mu_m$.
Since $\|h_m\|=|\gamma_m|^{-1}\leq C|m|^{-1}$,
the coefficients $(h_mu_m)_{m\ne0}$ still decay rapidly and define a smooth form $v$.
Termwise differentiation gives
$u=u_0+\dbar v$.

At zero frequency $\dbar$ vanishes, and every $\dbar v$ has zero constant Fourier coefficient.
No nonzero constant-coefficient $(0,q)$-form is therefore exact.
Thus $H^{0,q}_{\dbar}$ consists exactly of constant-coefficient forms.
A general $(p,q)$-form adds only fixed factors $dz_I$.
Apply the same computation to each $I$; since
$\dbar(dz_I\wedge v)=(-1)^p dz_I\wedge\dbar v$,
multiply the corresponding primitive by $(-1)^p$. This gives the stated basis and dimension.
\end{proof}

\begin{proposition}[The heat trace of a flat torus]
\label{heat:torus-trace}
On the real torus $\R^m/\Z^m$, the nonnegative operator $A=-\Delta$ has heat trace
\[
\tr(e^{-tA})=\sum_{a\in\Z^m}e^{-4\pi^2t|a|^2}<\infty,
\]
and heat kernel
\begin{equation}
K_t(x,y)=\sum_{\ell\in\Z^m}G_{t,1}(x-y+\ell).
\label{heat:periodized-kernel}
\end{equation}
\end{proposition}
\begin{proof}
The Fourier basis $e^{2\pi\ii a\cdot x}$ has eigenvalues $4\pi^2|a|^2$.
For fixed $t>0$, the multidimensional series factors into $m$ convergent one-dimensional Gaussian series.
The periodized Gaussian on the right converges, together with all derivatives, uniformly
on compact positive time intervals and compact sets in $x,y$, since polynomial factors are absorbed by slightly weaker Gaussian decay.
Its Fourier coefficient is
\[
\int_{[0,1]^m}\sum_{\ell\in\Z^m}G_{t,1}(z+\ell)
 e^{-2\pi\ii a\cdot z}\,dz
=\int_{\R^m}G_{t,1}(z)e^{-2\pi\ii a\cdot z}\,dz
=e^{-4\pi^2t|a|^2}.
\]
The last equality follows coordinatewise from the Gaussian Fourier transform, with normalization supplied by Lemma~\ref{heat:gaussian-convolution}.
It also follows by differentiating in the frequency parameter, integrating by parts, and solving
$F'(\xi)=-2t\xi F(\xi)$ with $F(0)=1$.
Thus \eqref{heat:periodized-kernel} is exactly the kernel of the spectral semigroup above.
\end{proof}

Here the dimension of the zero modes stays fixed, while positive frequencies are suppressed by $e^{-t\lambda}$.
This displays the basic mechanism of the heat method.
A different metric normalization for the $\dbar$-Laplacian changes only the constant in time,
not the cohomology dimensions in Proposition~\ref{heat:torus-cohomology}.


\section{Measure-zero spheres and Gaussian convolution on manifolds}

\begin{proposition}[Continuity of ball mass]
\label{heat:ball-continuity}
Write $d_g$ for the distance in the following metric-space statements. On a metric space with finite measure, $F_x(r)=\mu(\{y:d_g(x,y)\leq r\})$ is nondecreasing and right-continuous.
If $M$ is a compact Riemannian manifold of positive dimension and $\mu=dV_g$, each distance sphere
$\{y:d_g(x,y)=r\}$ has zero volume. Open and closed balls therefore have the same volume, and $F_x$ is continuous.
\end{proposition}
\begin{proof}
Monotonicity follows from inclusion. If $r_j\downarrow r$, the closed balls decrease to the closed ball of radius $r$;
continuity of a finite measure on decreasing sequences gives right continuity.

For volume measure, put $d=\dim_{\R}M\geq1$ and first show that every $y$ in the component of $x$ is reached by a minimizing geodesic from $x$.
Take piecewise smooth curves with lengths tending to $d_g(x,y)$, parametrized at constant speed on $[0,1]$.
Bounded lengths give equicontinuity; compactness and Arzel\`a--Ascoli yield a uniformly convergent subsequence.
Length is the supremum of the sums of distances over finite partitions, hence is lower semicontinuous under uniform convergence.
The limit realizes the minimum. In a sufficiently small normal coordinate ball, the short geodesic between two points minimizes length,
with uniqueness in the equality case, since the Gauss lemma separates radial length from a nonnegative tangential contribution.
Replacing short arcs of the minimizing curve by these geodesics shows that it is already locally geodesic.
It is therefore smooth, and for $r=d_g(x,y)$ there is $v\in T_xM$ with $|v|=r$ and $y=\exp_x(v)$.

The distance sphere of radius $r>0$ is consequently contained in
$\exp_x\{v:|v|=r\}$, the image of a compact smooth $(d-1)$-manifold under a smooth map.
In finitely many coordinate charts this map is Lipschitz. Subdivide the $(d-1)$-dimensional parameter domains into cubes of side $\delta$;
their images are covered by $O(\delta^{-(d-1)})$ balls of radius $C\delta$ in dimension $d$.
Their total volume is $O(\delta)$; let $\delta\downarrow0$ to obtain measure zero.
For $r=0$, the sphere is a single point and also has zero volume.
The union of increasing smaller closed balls is the open ball. This gives left continuity, using the zero boundary volume, and combines with right continuity already proved.
\end{proof}

\begin{lemma}[Convolution of two Gaussian kernels on a manifold]
\label{heat:manifold-convolution}
On a compact $d$-dimensional Riemannian manifold, for every $a,T>0$ there is $C$ such that, for $u,v>0$ with $u+v\leq T$,
\[
 \int_M u^{-d/2}e^{-d_g(x,z)^2/(au)}
             v^{-d/2}e^{-d_g(z,y)^2/(av)}\,dV(z)
 \leq C(u+v)^{-d/2}e^{-d_g(x,y)^2/(2a(u+v))}.
\]
The same estimate holds uniformly for metric families with uniform small-ball volume bounds and a uniform bound on total volume.
\end{lemma}
\begin{proof}
On finitely many coordinate charts, the metric is uniformly equivalent to the Euclidean metric, so small balls have volume at most $Cr^d$.
For larger radii, increase $C$, using the total-volume bound. Corollary~\ref{heat:radial-gaussian} then applies uniformly.
The triangle inequality and weighted Cauchy--Schwarz give
\[
 \frac{d_g(x,z)^2}{u}+\frac{d_g(z,y)^2}{v}
 \geq\frac{d_g(x,y)^2}{u+v}.
\]
Use half the exponent to extract the required off-diagonal factor.
From the remaining half retain only the first term to bound the spatial integral by $Cu^{d/2}$, or only the second to bound it by $Cv^{d/2}$.
Take the smaller bound and multiply by $u^{-d/2}v^{-d/2}$, obtaining
$C\max(u,v)^{-d/2}\leq C2^{d/2}(u+v)^{-d/2}$.
This proves the estimate.
\end{proof}

\begin{remark}
The lemma is suitable for one or finitely many Duhamel convolutions.
Repeated use in an infinite Volterra series would progressively weaken the Gaussian exponent.
The next chapter uses exponential weights to control the entire parametrix series with a Gaussian constant independent of the number of convolutions.
\end{remark}


\chapter{Heat-Kernel Localization, Frozen Coefficients, and Uniform Scaling}
\label{heat:local-chapter}

This chapter supplies the analytic passage from the exact magnetic model to a general manifold.
We prove a local scaling theorem with a uniform remainder before integrating the heat trace.
Three ingredients have distinct roles: a scalar parametrix supplies Gaussian bounds; Kato's inequality and local energy estimates give
bounds for the vector-bundle heat kernel and one covariant derivative; a quadratic-form Duhamel identity places one derivative on each side of the second-order error.
The time endpoints therefore produce the integrable factor $s^{-1/2}(1-s)^{-1/2}$ instead of $s^{-1}$.
Finite propagation then shows that local extension leaves the leading term unchanged.
This approach follows the coefficient-freezing and heat-kernel ideas of Demailly and related proofs
\cite{Demailly1985,DemaillyMorse1991,Bismut1987}; for higher-order expansions, see
\cite{MaMarinescuZelditch2015}.

The Gaussian factor $\exp(-k|x-y|^2/(4t))$ from Chapter~\ref{heat:models-chapter} concentrates on distances of order $\sqrt{t/k}$.
At that scale, a smooth coefficient changes little near the center.
The purpose of the estimates below is to make this observation uniform in the center and to control the contribution from outside the coordinate neighborhood.

\section{Scalar Gaussian bounds}

Below, uniform geometric bounds mean uniform equivalence to the Euclidean metric on coordinate balls of fixed radius,
uniform bounds for the volume density and the finitely many coefficient derivatives used, and a uniform positive normal-coordinate radius.
A compact manifold automatically has these properties. Multiplying its metric by any constant at least $1$
preserves such uniform bounds. Uniformly elliptic smooth metrics on Euclidean space with bounded coefficients and derivatives are also covered.
All time intervals are specified as $0<t\leq T$, and constants may depend on $T$.
Throughout this chapter $d=\dim_{\R}M$, $d_g(x,y)$ is Riemannian distance, and $dV_g$ is Riemannian volume.
The scalar Laplacian is $\Delta_g=d^*d\geq0$ on functions; on Euclidean space it is $-\sum_i\partial_i^2$.
A Laplace-type operator has scalar principal symbol $g^{ij}\xi_i\xi_j\operatorname{Id}$.
Its heat kernel acts from the fiber at $y$ to the fiber at $x$, as defined in paragraph~\ref{heat:kernel-notation}.
Covariant derivatives in the $y$ variable act on the dual source-bundle factor.
Repeated real coordinate indices are summed from $1$ to $d$; $O_T$ denotes a bound whose constant may depend on $T$ but not on the varying scale or center.

\begin{theorem}[Scalar Gaussian bounds with one derivative]
\label{heat:scalar-gaussian}
In the preceding settings, the heat kernel $h$ of the nonnegative scalar Laplacian $\Delta_g$ satisfies
\begin{align}
 0\leq h(t;x,y)&\leq C_Tt^{-d/2}e^{-d_g(x,y)^2/(C_Tt)},\label{heat:scalar-gaussian-formula}\\
 |d_xh(t;x,y)|&\leq C_Tt^{-(d+1)/2}e^{-d_g(x,y)^2/(C_Tt)}.
 \label{heat:scalar-gradient-formula}
\end{align}
The constants depend only on the stated uniform bounds and $T$.
For the Dirichlet realization on a domain, the scalar kernel is nonnegative and bounded above by that of the surrounding manifold without boundary.
\end{theorem}
\begin{proof}
\emph{Step 1: construct an initial parametrix.}
In a coordinate chart write
$\Delta_g=-a^{ij}(x)\partial_i\partial_j+b^j(x)\partial_j$,
with $dV_g=\rho(x)\,dx$. Freeze $a$ at $y$ and put
\[
 Z_y(t;x,y)=\frac{\exp[-\langle a(y)^{-1}(x-y),x-y\rangle/(4t)]}
 {(4\pi t)^{d/2}\sqrt{\det a(y)}\,\rho(y)}.
\]
Multiply by a partition of unity in $y$ and a cutoff in $x$, the latter equal to $1$ near the support of the former.
Sum these kernels to obtain $Z$. For a noncompact uniform atlas, use a locally finite partition with uniformly bounded overlap.
The frozen Gaussian has initial mass $\delta_x$ relative to $\rho(y)\,dy$, so $Z(t)f\to f$.

Set $R=(\partial_t+\Delta_{g,x})Z$.
The second-order residual contains $a(x)-a(y)$, whose norm is at most $C|x-y|$.
The Gaussian differentiation in Proposition~\ref{heat:gaussian-derivatives}, after a linear change of coordinates, shows that the frozen Gaussian's second derivatives have the form
$[O(t^{-1})+O(|x-y|^2t^{-2})]Z_y$.
The inequality $r^m e^{-r^2/(ct)}\leq C_m t^{m/2}e^{-r^2/(2ct)}$ gives
\begin{equation}
 |R(t;x,y)|\leq Ct^{-1/2}t^{-d/2}e^{-d_g(x,y)^2/(Ct)},\quad
 |d_xZ(t;x,y)|\leq Ct^{-1/2}t^{-d/2}e^{-d_g(x,y)^2/(Ct)}.
 \label{heat:parametrix-residual}
\end{equation}
The first-order coefficient terms have the same bound. Cutoff derivatives are supported away from the diagonal,
so their factors $t^{-m}e^{-c/t}$ are also absorbed into the right-hand side. The kernel $Z$ satisfies the analogous bound without $t^{-1/2}$.

\emph{Step 2: control the entire Volterra series.}
For two kernels, write
\[
 (K_1*K_2)(t;x,y)=\int_0^t\int_M K_1(t-s;x,z)K_2(s;z,y)\,dV_g(z)\,ds.
\]
Thus $*$ denotes time-and-space convolution, and $R^{*j}$ means the $j$-fold such convolution.
Formally set $\Phi=-R+R*R-R*R*R+\cdots$ and $H=Z+Z*\Phi$.
The following estimates prove convergence and the Gaussian bounds together.
For a $1$-Lipschitz function $f$ and $a\in\R$, weight a kernel by
$e^{af(x)}K(x,y)e^{-af(y)}$.
Complete the square using $a(f(x)-f(y))\leq |a|d_g(x,y)$.
Estimate \eqref{heat:parametrix-residual} bounds both row and column integrals of the weighted residual by
$Ct^{-1/2}e^{Ca^2t}$, and its supremum by
$Ct^{-(d+1)/2}e^{Ca^2t}$.
The kernels $Z$ and $d_xZ$ satisfy the same rules with time powers $0$ and $-1/2$, respectively.

For a convolution of $j$ residuals, fix a time partition $s_1+\cdots+s_j=t$.
At least one $s_i\geq t/j$.
Use the supremum bound for that kernel, row integrals to its left, and column integrals to its right, obtaining
\[
 C^j(j/t)^{d/2}e^{Ca^2t}\prod_{i=1}^j s_i^{-1/2}.
\]
Integrate over the time simplex and apply the Beta integral successively:
\begin{equation}
 |e^{af(x)}R^{*j}(t;x,y)e^{-af(y)}|
 \leq C^j j^{d/2}t^{-d/2}e^{Ca^2t}
       \frac{\Gamma(1/2)^j t^{j/2-1}}{\Gamma(j/2)}.
 \label{heat:volterra-majorant}
\end{equation}
The Beta formula itself follows by substituting $s=tu$ and applying Fubini to
\[\int_0^\infty\int_0^\infty x^{\alpha-1}y^{\beta-1}e^{-x-y}\,dxdy\]
with the change of variables $x+y=r$.
Here $\Gamma$ is the Gamma function defined in Proposition~\ref{heat:gaussian-derivatives}.
Its recurrence $\Gamma(u+1)=u\Gamma(u)$, proved by integration by parts, makes the ratio of terms with indices $j+2$ and $j$ tend to zero on every bounded time interval.
Applying this separately to the even and odd subseries proves convergence with the stated time weight on $0<t\leq T$.
The same argument for $Z*R^{*j}$ and $(d_xZ)*R^{*j}$ gives, respectively,
\[
 |e^{af(x)}H(t;x,y)e^{-af(y)}|
 \leq C_Tt^{-d/2}e^{Ca^2t},\qquad
 |e^{af(x)}d_xH(t;x,y)e^{-af(y)}|
 \leq C_Tt^{-(d+1)/2}e^{Ca^2t}.
\]
For the derivative, differentiate only the leftmost $Z$; its additional factor $s^{-1/2}$ remains integrable.
Choose $f(z)=d_g(z,y)$ and $a=d_g(x,y)/(2Ct)$ to obtain the required off-diagonal decay.
The weighted estimate avoids progressive deterioration of Gaussian constants under repeated two-kernel convolutions.

\emph{Step 3: verify the heat equation.}
Absolute convergence permits associating convolutions, and
$\Phi+R+R*\Phi=0$. Moreover,
$(\partial_t+\Delta_g)(Z*\Phi)=\Phi+R*\Phi$, so $H$ solves the heat equation distributionally.
The row and column integrals of $Z*\Phi$ are $O(\sqrt t)$, so this correction does not change the initial value.
Proposition~\ref{heat:smoothing-kernel}, applied to the nonnegative scalar realization, or local regularity on compact positive time intervals,
upgrades the distributional solution to a smooth one.
The difference of two $L^2$ solutions satisfies
$\frac d{dt}\|u\|_2^2=-2\|du\|_2^2\leq0$, giving uniqueness.
Thus $H$ is the kernel $h$ of the closed-form realization.

Finally, a solution with nonnegative initial value cannot first develop a negative interior minimum.
On a noncompact space, use a smooth equivalent of $\epsilon e^{At}(1+d_g(x,x_0)^2)$ as a barrier and pass to the limit,
or apply the energy estimate to the negative part $u_-\in H^1$.
Thus $h\geq0$. The Dirichlet solution has zero boundary values; subtract it from the surrounding nonnegative solution and apply the same comparison principle
to obtain kernel domination.
\end{proof}

\section{Kato's inequality and covariant gradient bounds}

\begin{lemma}[Kato's inequality for the heat equation]
\label{heat:kato}
Let $\nabla$ be a unitary connection on a Hermitian bundle and let $Q=Q^*\geq-M$, $M\geq0$.
For $P=\nabla^*\nabla+Q$,
\[
 |e^{-tP}u|(x)\leq e^{Mt}(e^{-t\Delta_g}|u|)(x),\qquad
 \|H_P(t;x,y)\|\leq e^{Mt}h(t;x,y).
\]
The conclusion also holds for Dirichlet realizations.
\end{lemma}
\begin{proof}
First calculate for a smooth solution $u$, setting $r_\eta=(|u|^2+\eta^2)^{1/2}$.
Since the connection is unitary,
\[
 (\partial_t+\Delta_g)|u|^2
 =-2|\nabla u|^2-2\operatorname{Re}\langle u,Qu\rangle,
 \qquad |d|u|^2|\leq2|u|\,|\nabla u|.
\]
With the nonnegative sign convention for $\Delta_g$, the chain rule gives
\[
 (\partial_t+\Delta_g)r_\eta
 =-\frac{|\nabla u|^2+\operatorname{Re}\langle u,Qu\rangle}{r_\eta}
   +\frac{|d|u|^2|^2}{4r_\eta^3}
 \leq M r_\eta.
\]
Apply the scalar comparison principle proved at the end of Theorem~\ref{heat:scalar-gaussian} and let $\eta\downarrow0$ to obtain the first inequality.
Approximate $L^2$ initial data by smooth data; alternatively pass first to the distributional limit in Kato's inequality.
Apply the first inequality to an approximate identity near $y$ times a fixed fiber vector $v$, then pass to the limit:
$|H_P(t;x,y)v|\leq e^{Mt}h(t;x,y)|v|$.
Taking the supremum over unit $v$ gives the kernel-norm bound.
Zero boundary conditions are preserved in the comparison, so the Dirichlet case is identical.
\end{proof}

\begin{lemma}[Local mean-value and energy estimates]
\label{heat:parabolic-energy}
In a region with uniform geometric bounds, set
$Q_a=(t_0-a^2,t_0]\times B(x_0,a)$, so $Q_{2r}=(t_0-4r^2,t_0]\times B(x_0,2r)$, with $r$ at most a fixed coordinate radius.
All cylinder integrals use $dt\,dV_g$, and the entire cylinder is assumed to lie in the region of the solution.
If a nonnegative smooth function $w$ satisfies $(\partial_t+\Delta_g)w\leq C_0w$, then
\begin{equation}
 w(t_0,x_0)\leq C r^{-d-2}\int_{Q_{2r}}w.
 \label{heat:mean-value}
\end{equation}
If $u$ solves $(\partial_t+\nabla^*\nabla+Q)u=0$ with $\|Q\|\leq C_0$, then
\begin{equation}
 \int_{Q_r}|\nabla u|^2\leq Cr^{-2}\int_{Q_{2r}}|u|^2.
 \label{heat:caccioppoli}
\end{equation}
\end{lemma}
\begin{proof}
Replace $w$ by $e^{-C_0(t-t_0)}w$ to remove the zeroth-order term; over the bounded time length under consideration this changes only the constant.
Take a space-time cutoff $\eta$ equal to $1$ near the final part of $Q_r$ and zero near the initial time and lateral boundary of $Q_{2r}$,
with $|\partial_t\eta|+|\Delta\eta|\leq Cr^{-2}$ and $|d\eta|\leq Cr^{-1}$.
Use the scalar heat-kernel representation for $\eta w$, obtained by applying Proposition~\ref{heat:inhomogeneous} to the compactly supported cutoff expression. In the Leibniz formula,
integrate $-2\langle d\eta,dw\rangle$ by parts spatially, moving derivatives onto the kernel and $\eta$.
Terms containing $\partial_t\eta$ lie at least $cr^2$ before the final time, where the kernel is at most $Cr^{-d}$.
Terms with spatial cutoff derivatives lie at least $cr$ from $x_0$; Theorem~\ref{heat:scalar-gaussian} gives the Gaussian kernel and gradient bounds
$Cr^{-d}$ and $Cr^{-d-1}$ there, even as time approaches the final time.
Every term is therefore bounded in absolute value by $Cr^{-d-2}\int_{Q_{2r}}w$.
The nonpositive heat-equation residual only decreases the representation, proving \eqref{heat:mean-value}.
The proof also shows that the mean-value constant is independent of $w$.

For the second estimate, pair the vector heat equation with $\eta^2u$ and take real parts.
Integration by parts in time gives a nonnegative final-time norm; integration by parts in space gives
$\int\eta^2|\nabla u|^2$ and the mixed term
$2\operatorname{Re}\int\eta\langle\nabla u,d\eta\otimes u\rangle$.
Use $2ab\leq\tfrac12a^2+2b^2$ to absorb half the gradient energy:
\[
 \frac12\int\eta^2|\nabla u|^2
 \leq C\int (|d\eta|^2+|\partial_t\eta|+C_0\eta^2)|u|^2.
\]
Choose $\eta=1$ on $Q_r$ and zero at the earliest time to obtain \eqref{heat:caccioppoli}.
\end{proof}

\begin{proposition}[One covariant derivative of the vector heat kernel]
\label{heat:covariant-gradient}
In addition to uniform geometric bounds, suppose the curvature of the unitary connection and its first covariant derivative, and $Q$ and its first covariant derivative,
are uniformly bounded. Then the kernel of $P=\nabla^*\nabla+Q$ satisfies
\begin{equation}
 \|\nabla_xH_P(t;x,y)\|+\|\nabla_yH_P(t;x,y)\|
 \leq C_Tt^{-(d+1)/2}e^{-d_g(x,y)^2/(C_Tt)},\qquad0<t\leq T.
 \label{heat:gradient-gaussian}
\end{equation}
No bounded connection one-form in a global gauge is required.
\end{proposition}
\begin{proof}
For a solution $u$ of the vector heat equation, commute covariant derivatives in a normal orthonormal frame at a point.
Commuting two derivatives produces the bundle-curvature term $F_{ij}u$; on the one-tensor $\nabla u$ it also produces Riemann curvature.
Comparing $\nabla_i\nabla_j\nabla_j u$ with
$\nabla_j\nabla_j\nabla_i u$ by successive commutations gives two kinds of terms:
two curvature terms multiplying $\nabla u$, and one term $\nabla F$ multiplying $u$.
Differentiating $Qu$ adds $Q\nabla u+(\nabla Q)u$.
Consequently, pointwise,
\[
 (\partial_t+\Delta_g)|\nabla u|^2
 \leq-2|\nabla^2u|^2+C(|\nabla u|^2+|u|^2).
\]
Here $2ab\leq a^2+b^2$ bounds each mixed term, with constants determined by the listed curvature bounds.
Add the calculation for $|u|^2$ in Lemma~\ref{heat:kato} to see that
$w=|\nabla u|^2+|u|^2$ satisfies Lemma~\ref{heat:parabolic-energy}.

Take $u(s,z)=H_P(s;z,y)v$, $|v|=1$. Near $(t,x)$ choose nested cylinders with
$r=\min(\sqrt t/8,r_0/8)$, all times lying in $[t/2,t]$.
First apply \eqref{heat:mean-value}, then \eqref{heat:caccioppoli}, to get
\[
 |\nabla u(t,x)|^2\leq Cr^{-d-4}\int_{Q_{4r}}|u|^2.
\]
Lemma~\ref{heat:kato} and Theorem~\ref{heat:scalar-gaussian} give, within this cylinder,
\[
 |u(s,z)|^2\leq C_Tt^{-d}
       \exp[-d_g(x,y)^2/(C_Tt)],
\]
since $d_g(z,y)\geq d_g(x,y)-4r$ and $r^2\leq Ct$.
The cylinder has space-time volume at most $Cr^{d+2}$.
Thus $|\nabla u(t,x)|^2\leq C_Tt^{-d-1}e^{-d_g(x,y)^2/(C_Tt)}$.
For $t$ above a fixed small constant, the fixed positive lower bound on $r$ and $t\leq T$ change only $C_T$.
Take square roots and the supremum over $v$ for the $x$-derivative estimate.
Self-adjointness gives $H(t;x,y)^*=H(t;y,x)$, yielding the same bound for the $y$ derivative.
\end{proof}

\begin{corollary}[Local and global bounds in the original scale]
\label{heat:global-gaussian}
Let $P_k=k^{-1}\nabla_k^*\nabla_k-V_k$, with $V_k=V_k^*$ and
$\sup_M\lambda_{\max}(V_k)\leq M$, where $\lambda_{\max}$ denotes the largest Hermitian eigenvalue.
On a fixed compact $d$-dimensional manifold, for $0<t\leq k$,
\[
 \|H_k(t;x,y)\|\leq C(k/t)^{d/2}
       \exp\left(Mt-\frac{k\,d_g(x,y)^2}{Ct}\right).
\]
For the Dirichlet kernel on an open set $\Omega$, replace $M$ by
$M_\Omega=\sup_\Omega\lambda_{\max}(V_k)$.
In the global kernel, $M$ cannot be replaced merely by a potential bound near $x,y$.
\end{corollary}
\begin{proof}
Apply Lemma~\ref{heat:kato} to the connection Laplacian semigroup at time $t/k$.
In that time variable the potential $-V_k$ becomes $-kV_k$, so its exponential factor is exactly $e^{Mt}$.
Then use Theorem~\ref{heat:scalar-gaussian} for $0<t/k\leq1$.
Dirichlet comparison takes place entirely within $\Omega$, so only the potential bound there is needed.
Proposition~\ref{heat:remote-well} will give a counterexample to the replacement excluded in the last sentence.
\end{proof}

\section{Putting second-order errors into quadratic forms}

\begin{lemma}[A quadratic-form Duhamel estimate with polynomial weights]
\label{heat:form-duhamel}
For $0<\epsilon\leq1$, work on the same Lebesgue space $L^2(\R^d,\C^r)$ and let
\[
 q_\epsilon(v,w)=\int\langle\nabla_\epsilon v,a_\epsilon\nabla_\epsilon w\rangle
                       +\langle v,Q_\epsilon w\rangle,
 \qquad q_0(v,w)=\int\langle\nabla_0v,\nabla_0w\rangle+\langle v,Q_0w\rangle.
\]
where $a_\epsilon$ is uniformly positive definite and bounded, and $\nabla_\epsilon,\nabla_0$ are unitary connections.
Let $P_\epsilon,P_0$ denote the lower-bounded self-adjoint operators associated with the closures of these forms, as in Proposition~\ref{heat:friedrichs}; write $H_\epsilon,H_0$ for their kernels.
Assume that the corresponding heat kernels and one covariant derivative satisfy the preceding uniform Gaussian bounds for $0<s\leq1$, and that
\begin{equation}
 |a_\epsilon-I|+|\nabla_\epsilon-\nabla_0|+|Q_\epsilon-Q_0|
 \leq C\epsilon(1+|Z|^N),\qquad Z\in\R^d.
 \label{heat:weighted-coefficients}
\end{equation}
Here $N\in\N_0$ is fixed and $|\nabla_\epsilon-\nabla_0|$ is the pointwise norm of the difference one-form. There is then a constant independent of $\epsilon$ such that
\[
 \|H_\epsilon(1;0,0)-H_0(1;0,0)\|\leq C'\epsilon.
\]
\end{lemma}
\begin{proof}
Set $C_\epsilon=\nabla_\epsilon-\nabla_0$, an endomorphism-valued one-form.
Adding and subtracting intermediate terms gives the exact identity
\begin{align*}
 q_\epsilon(v,w)-q_0(v,w)
 ={}&\int\langle\nabla_\epsilon v,(a_\epsilon-I)\nabla_0w\rangle
       +\langle\nabla_\epsilon v,a_\epsilon C_\epsilon w\rangle\\
    &+\langle C_\epsilon v,\nabla_0w\rangle
       +\langle v,(Q_\epsilon-Q_0)w\rangle.
\end{align*}
Each second-order error thus places one derivative on each of two different kernels.

First take $f,g\in C_c^\infty$ and set $v_s=e^{-(1-s)P_\epsilon}f$, $w_s=e^{-sP_0}g$.
Gaussian decay and the polynomial-weight bound put $C_\epsilon v_s,C_\epsilon w_s$ in $L^2$.
Switching connections therefore preserves finiteness of first-order energy, so both vectors lie in both form domains.
Spatial cutoffs followed by local smoothing approximate them in form norm by compactly supported smooth sections,
justifying the form-operator pairing identities.
Differentiate first for $\delta\leq s\leq1-\delta$. Lemma~\ref{heat:borel-calculus} places these vectors in the domains of every power of their respective operators, hence
\[
 \frac d{ds}\langle v_s,w_s\rangle
       =q_\epsilon(v_s,w_s)-q_0(v_s,w_s).
\]
Insert a compactly supported spatial cutoff first, then use Gaussian decay to remove the terms at infinity.
Integration gives the quadratic-form version of Duhamel's identity.

Now take $f,g$ to be approximate unit masses at $0$ times unit fiber vectors.
The common Gaussian factor of the two kernels at the intermediate point $Z$ is bounded by
\[
 [s(1-s)]^{-d/2}
 \exp\left(-\frac{|Z|^2}{C s(1-s)}\right)
\]
For fixed $N$, the substitution $Z=\sqrt{s(1-s)}\,Y$ gives
\[
 [s(1-s)]^{-d/2}\int(1+|Z|^N)
 e^{-|Z|^2/(Cs(1-s))}\,dZ\leq C_N.
\]
The four spatial integrals above are therefore bounded respectively by
\[
 C\epsilon[s(1-s)]^{-1/2},\qquad
 C\epsilon(1-s)^{-1/2},\qquad
 C\epsilon s^{-1/2},\qquad C\epsilon.
\]
All are integrable on $(0,1)$; the first integrates to $C\pi\epsilon$.
These bounds remain uniform when the endpoint masses are supported in the unit ball: completing the square as in \eqref{heat:gaussian-semigroup} centers the spatial Gaussian at a convex combination of the two endpoints, whose norm is at most one. Polynomial moments around that center have the same bound.
We may therefore let
$\delta\downarrow0$ first and then let the initial supports shrink to zero.
Strong continuity at the endpoints and smoothness of the positive-time kernels identify the heat-kernel difference; dominated convergence bounds its norm by $C'\epsilon$.
This also justifies every exchange of integrals and every endpoint limit.
\end{proof}

\begin{remark}[The role of quadratic forms]
Applying the second-order differential operator $P_\epsilon-P_0$ directly to the right-hand heat kernel
generally produces $s^{-1}$, whose time integral diverges at zero.
Moreover, $\|P_\epsilon-P_0\|_{L^2\to L^2}$ is usually infinite.
The four-term decomposition of Lemma~\ref{heat:form-duhamel} controls the difference through its energy form.
The Gaussian covariant-gradient bounds then make the time integral convergent.
\end{remark}

\section{Finite propagation separates the local problem from its surroundings}

\begin{lemma}[Finite propagation for the wave equation]
\label{heat:finite-propagation}
Let $P=\nabla^*\nabla+Q$, $Q\geq0$, act on a Hermitian bundle over a complete Riemannian manifold.
If $f$ is supported in $K$, then
\[
 \supp\cos(s\sqrt P)f\subset\{x:d_g(x,K)\leq |s|\}.
\]
In coordinates, if the principal-symbol matrix satisfies $a\leq c^2I$, the Euclidean propagation speed is at most $c$.
If two such operators agree on a ball, their wave solutions agree in interior cones that have not reached the boundary of that ball.
\end{lemma}
\begin{proof}
By Lemma~\ref{heat:borel-calculus}, for compactly supported smooth initial data $u(s)=\cos(s\sqrt P)f$ solves
$u_{ss}+Pu=0$, with $u(0)=f$ and $u_s(0)=0$.
Consider the moving region $\Omega_s=\{d_g(x,K)>s\}$ away from the initial support, with energy
\[
 E(s)=\frac12\int_{\Omega_s}
    (|u_s|^2+|\nabla u|^2+\langle u,Qu\rangle).
\]
Differentiate in time, use the wave equation, and integrate spatially by parts. The interior terms cancel.
The absolute boundary flux is at most
$\tfrac12(|u_s|^2+|\nabla_\nu u|^2)$; the energy lost as the boundary moves outward at unit speed
is at least this amount. Thus $E'(s)\leq0$.
Here $\nu$ denotes the outward unit normal to the moving boundary. The distance function is only Lipschitz, but smooth approximations give moving cutoffs
with $|d\rho|\leq1+o(1)$. Prove the cutoff integral inequality first, then pass to the limit by monotone convergence.
The initial exterior energy is zero, so $u_s$ and $\nabla u$ vanish outside the cone; the zero initial value then gives $u=0$ there.
Negative times follow from evenness of cosine.

Approximate general finite-energy initial data by smooth data.
In coordinates, the flux bound $a\leq c^2I$ gives the same argument with moving speed $c$.
For the difference of solutions of two operators, perform the same local energy calculation where the operators agree,
using a cone whose base is at the initial time and whose sides remain inside that region.
The difference has zero initial data and therefore vanishes in the cone. This proves the last assertion. More generally, if the potential is only locally bounded below, add $M|u|^2$ to the local energy so that it is positive. The additional interior terms are bounded by $C E(s)$, and Gronwall's inequality gives the same zero-energy conclusion. Thus finite propagation also applies to a nonnegative self-adjoint realization whose potential is not pointwise nonnegative, as needed below.
\end{proof}

\begin{proposition}[Exponential localization of the interior heat kernel]
\label{heat:wave-localization}
Let $P_1,P_2$ agree as Laplace-type operators on the coordinate ball $B(0,4R)$,
and suppose both are bounded below by the same $-C_0$. Assume uniform finite-order coefficient bounds and an ellipticity lower bound near the center,
and a fixed propagation-speed bound $c$ throughout the comparison ball.
For $R\geq1$,
\[
 \|H_{P_1}(1;0,0)-H_{P_2}(1;0,0)\|\leq C e^{-c_1R^2}.
\]
The constants are independent of coefficients outside the ball, but may depend on the common lower bound $C_0$.
The statement also applies when one operator is the Dirichlet realization inside the ball.
\end{proposition}
\begin{proof}
Add $C_0+1$ to both operators to make them positive; at the end multiply back by $e^{C_0+1}$.
For $P\geq0$, the one-dimensional Gaussian Fourier integral (computed in Proposition~\ref{heat:torus-trace}) and Lemma~\ref{heat:borel-calculus} give
\begin{equation}
 e^{-P}=\frac1{\sqrt{4\pi}}\int_\R e^{-s^2/4}\cos(s\sqrt P)\,ds.
 \label{heat:gaussian-functional-calculus}
\end{equation}
The scalar identity follows from the Gaussian Fourier transform; the operator integral is strongly absolutely convergent because the cosine norm is at most $1$.

Choose an even cutoff $\chi$, equal to $1$ on $[-1/2,1/2]$ and supported in $[-1,1]$.
Split the integral into the low-propagation part containing $\chi(cs/R)$ and the remaining tail.
Lemma~\ref{heat:finite-propagation}, including its locally lower-bounded-potential version, makes the local kernels of the low-propagation parts identical on $B(0,1)$.
The tail is $f_R(\sqrt P)$, where
\[
 f_R(\lambda)=\frac1{\sqrt{4\pi}}
 \int(1-\chi(cs/R))e^{-s^2/4}\cos(s\lambda)\,ds.
\]
Integrate by parts $2N$ times in $s$. Cutoff derivatives occur only for $|s|\geq R/(2c)$,
and every Gaussian derivative is a polynomial times the Gaussian. Hence
\begin{equation}
 \sup_\lambda(1+\lambda^2)^N|f_R(\lambda)|\leq C_Ne^{-c_2R^2}.
 \label{heat:fourier-tail}
\end{equation}
All polynomial factors can be absorbed by slightly decreasing $c_2$.

We must still pass from operator norm to point evaluation.
Choose an integer $m$ with $2m>d/2$.
Theorem~\ref{ana:interior-estimate} on nested interior balls and Lemma~\ref{ana:sobolev-embedding} give
\[
 |[(1+P)^{-m}v](0)|\leq C\|v\|_2.
\]
Thus point evaluation composed with $(1+P)^{-m}$ is represented by an $L^2$ vector of norm at most $C$.
Factor out $(1+P)^{-m}$ from each side of $f_R(\sqrt P)$. The intervening operator has norm at most
$\sup_\lambda(1+\lambda^2)^{2m}|f_R(\lambda)|$.
Estimate \eqref{heat:fourier-tail} gives the exponential bound for the tail kernel at $(0,0)$.
Apply this estimate to the two tails separately and subtract.
Before reaching the boundary, Dirichlet waves obey exactly the same local energy argument, and the interior elliptic estimates at the center are unchanged.
\end{proof}

\begin{remark}
The lower bound here is global. Making the potential very negative outside the ball can create large negative eigenvalues,
and then $C_0$ increases as well. Finite propagation does not remove this global information.
\end{remark}

\section{Radial gauge and uniform coefficients after scaling}

\begin{lemma}[The radial-gauge formula]
\label{heat:radial-gauge}
On a star-shaped coordinate ball centered at $0$, parallel-transport an orthonormal frame along radial lines.
Write $\nabla=d+A$ in this frame and $F=dA+A\wedge A$ for its curvature. The resulting connection form $A=A_i\,dx^i$ satisfies $x^iA_i=0$ and
\[
 A_i(x)=\int_0^1 s x^jF_{ji}(sx)\,ds
       =\tfrac12x^jF_{ji}(0)+O(|x|^2).
\]
Finite-order derivative bounds for the remainder are controlled by the corresponding curvature bounds.
\end{lemma}
\begin{proof}
Radial parallel transport is exactly the condition $x^iA_i=0$.
Using $F_{ji}=\partial_jA_i-\partial_iA_j+[A_j,A_i]$ and differentiating $x^jA_j=0$ gives
$x^jF_{ji}=x^j\partial_jA_i+A_i$; the commutator term vanishes because $x^jA_j=0$.
Therefore
$\frac d{ds}[sA_i(sx)]=s x^jF_{ji}(sx)$.
Integrate from $0$ to $1$; the endpoint term at $s=0$ is zero, giving the formula.
Next write
$F_{ji}(sx)-F_{ji}(0)=\int_0^s x^\ell\partial_\ell F_{ji}(ux)\,du$
to obtain the $O(|x|^2)$ remainder and its derivative bounds.
\end{proof}

\begin{proposition}[Scaling and extension to the whole space]
\label{heat:scaled-coefficients}
Let $M$ be compact, let $L$ be a Hermitian line bundle, and let $E$ be a fixed Hermitian vector bundle, both with smooth unitary connections.
Let $V,W$ be fixed smooth Hermitian endomorphisms of $E$, and set
\[
 P_k=k^{-1}(\nabla^{L^k\otimes E})^*\nabla^{L^k\otimes E}-V+k^{-1}W.
\]
Fix $T>0$. At $x_0$, take normal coordinates and radial orthonormal frames, and put
$\epsilon=\sqrt{t/k}$, $0<t\leq T$.
Pull back $tP_k$ by $x=\epsilon Z$ and unitarily conjugate the volume density to Lebesgue measure.
Its coefficients extend outside $|Z|<r/\epsilon$ to an operator $P_{\epsilon,t}$ on $\R^d$ with the following properties:
\begin{enumerate}
\item The extension agrees with the original operator inside that ball, and the principal symbol is uniformly elliptic.
\item The model is $P_{0,t}=\nabla_{0,t}^*\nabla_{0,t}-tV(x_0)$,
where $\nabla_{0,t}=d+\tfrac t2 Z^jF^L_{ji}(x_0)\,dZ^i$.
\item The coefficient differences satisfy \eqref{heat:weighted-coefficients}; one may take a fixed $N=4$.
\item The potentials, curvatures, and required covariant derivatives of both operators, together with the finite-order metric bounds,
are controlled by a constant depending only on $T$. Thus their heat kernels and one covariant derivative have uniform Gaussian bounds.
\end{enumerate}
All assertions are uniform for $x_0\in M$.
\end{proposition}
\begin{proof}
Choose $3r$ smaller than the uniform normal-coordinate radius.
Lemma~\ref{heat:radial-gauge} gives $A^L(x)=A_1^L(x)+O(|x|^2)$, where
$(A_1^L)_i(x)=\tfrac12x^jF^L_{ji}(x_0)$.
In the $Z$ variable the connection is
\[
 d+\epsilon k A^L(\epsilon Z)+\epsilon A^E(\epsilon Z)
 =d+tA_1^L(Z)+O_T(\epsilon|Z|^2)+\epsilon A^E(\epsilon Z),
\]
since $k\epsilon^2=t$.
The principal coefficient is $a(\epsilon Z)=I+O(\epsilon|Z|)$; normal coordinates actually give a stronger quadratic error,
but it is unnecessary here. The potential is $-tV(\epsilon Z)+\epsilon^2W(\epsilon Z)$.

We check that density conjugation introduces no uncontrolled first-order term.
For $\kappa(Z)=\sqrt{\det g(\epsilon Z)}$, multiplication by $\kappa^{1/2}$ is unitary onto the Lebesgue space.
The connection in the form becomes $\nabla-\beta$, where $\beta=\tfrac12d\log\kappa$ is real.
Expand and integrate the mixed terms by parts. The result is the divergence-form principal part of the original unitary connection, plus the real scalar potential
\[
 Q_\kappa=\partial_i(a^{ij}\beta_j)+a^{ij}\beta_i\beta_j.
\]
Its finite-order derivatives are uniformly bounded, and its difference from the model satisfies $C_T\epsilon(1+|Z|^4)$.
The density correction belongs in the potential; $\nabla-\beta$ itself is not a unitary connection.

Choose a fixed smooth cutoff $\chi$, equal to $1$ for $|x|\leq r$ and zero for $|x|\geq2r$.
Set $a_\epsilon=I+\chi(\epsilon Z)(a(\epsilon Z)-I)$.
For small $r$, this remains uniformly positive definite.
Define the connection as the model plus
\[
 \chi(\epsilon Z)
 [\epsilon kA^L(\epsilon Z)-tA_1^L(Z)+\epsilon A^E(\epsilon Z)]
\]
and define the potential by the corresponding cut-off difference from the model.
In the transition region, $|Z|\asymp\epsilon^{-1}$, meaning that their ratio is bounded above and below by positive constants independent of $\epsilon$.
Each derivative of $\chi(\epsilon Z)$ supplies a factor $\epsilon$.
Although the line-connection difference can reach $O_T(\epsilon^{-1})$, its first derivative and curvature remain $O_T(1)$.
The line connection is scalar, so its commutators with the connection, curvature, and potential of $E$ vanish.
The scaled connection of the fixed bundle is $O(\epsilon)$, with curvature $O(\epsilon^2)$;
cutoff terms have the same or smaller orders.
Thus curvature, its covariant derivatives, and the covariant derivatives of the potential are uniformly bounded.
The $Z$ derivatives of the metric are also controlled by the original finite-order bounds times the corresponding powers of $\epsilon$.

For $|Z|\leq2r/\epsilon$, Taylor remainders give the polynomial-weight bounds;
outside, all coefficients equal the model and their differences vanish.
This proves \eqref{heat:weighted-coefficients}.
The divergence-form operator is unitarily equivalent to the connection Laplacian for the metric $a_\epsilon^{-1}$ plus a bounded density potential,
so Lemma~\ref{heat:kato} and Proposition~\ref{heat:covariant-gradient} apply.
Compactness makes all coordinate, curvature, and potential bounds uniform in $x_0$.
\end{proof}

\section{Uniform diagonal asymptotics and heat traces}

\begin{theorem}[Local scaling for magnetic heat kernels]
\label{heat:scaled-diagonal}
Under Proposition~\ref{heat:scaled-coefficients}, let $H_k$ be the kernel of $e^{-tP_k}$.
In orthonormal real coordinates at $x$, write $\ii F^L_x$ as
$\sum_{j=1}^{\lfloor d/2\rfloor}b_j(x)\,dx_j\wedge dy_j$, allowing $b_j=0$.
For each fixed $T>0$, there is $C_T$ such that every $x\in M$, $k\geq1$, and $0<t\leq T$ satisfy
\begin{equation}
 \left\|(t/k)^{d/2}H_k(t;x,x)
 -(4\pi)^{-d/2}\prod_j\frac{tb_j(x)}{\sinh(tb_j(x))}\,e^{tV(x)}\right\|
 \leq C_T\sqrt{t/k}.
 \label{heat:uniform-diagonal}
\end{equation}
Zero-curvature factors have their continuous value $1$.
\end{theorem}
\begin{proof}
Fix $x_0,t,k$ and set $\epsilon=\sqrt{t/k}$.
In scaled coordinates, $e^{-tP_k}$ on the original manifold becomes a heat operator at time $1$.
It agrees with the whole-space extension $P_{\epsilon,t}$ on a ball of radius $r/\epsilon$.
Their lower bounds are uniform because their potentials have norm at most $C_T$; the finite-order coefficients near the center are uniformly bounded as well.
For sufficiently small $\epsilon$, Proposition~\ref{heat:wave-localization} bounds the central diagonal difference by
$C_Te^{-c/\epsilon^2}$.

Proposition~\ref{heat:scaled-coefficients} verifies every hypothesis of Lemma~\ref{heat:form-duhamel},
including the covariant-gradient bounds. The extension and model kernels thus differ on the time-$1$ diagonal by at most $C_T\epsilon$.
The constant magnetic model has curvature $t\ii F^L_{x_0}$ and potential contribution $-tV(x_0)$.
By \eqref{heat:model-diagonal}, its time-$1$ value at the origin is exactly
\[
 (4\pi)^{-d/2}\prod_j\frac{tb_j(x_0)}{\sinh(tb_j(x_0))}\,e^{tV(x_0)}.
\]

Scaling multiplies the kernel by $\epsilon^d$, since
$dV_g(\epsilon Z)=\epsilon^d\kappa(Z)dZ$.
At $Z=0$, density conjugation has value $\kappa(0)^{1/2}=1$, so there is no additional central factor.
The normalized original diagonal therefore differs from the model by at most
$C_T(\epsilon+e^{-c/\epsilon^2})\leq C'_T\epsilon$.
All coefficient bounds are uniform in $x_0$, proving the conclusion at every center.

The argument first covers $k\geq k_0(T)$.
For the remaining finitely many $k$, sufficiently small $t$ still satisfies the same scaling conditions.
On compact time intervals bounded away from zero, the smooth kernel is bounded on the compact manifold,
while $\sqrt{t/k}$ is bounded positively below. Increasing $C_T$ covers these remaining cases.
\end{proof}

\begin{corollary}[Heat-trace integration and the spatial variable]
\label{heat:semiclassical-trace}
Under the same hypotheses, for fixed $t>0$,
\begin{equation}
 \tr(e^{-tP_k})
 =\left(\frac{k}{4\pi t}\right)^{d/2}
   \int_M\tr(e^{tV(x)})\prod_j\frac{tb_j(x)}{\sinh(tb_j(x))}\,dV_g(x)
   +O_t(k^{(d-1)/2}).
 \label{heat:trace-main-term}
\end{equation}
For $0<t\leq T$, the remainder after multiplication by $(t/k)^{d/2}$ is $O_T(\sqrt{t/k})$.
\end{corollary}
\begin{proof}
The heat operator is positive and has a smooth kernel. Proposition~\ref{heat:kernel-trace} identifies its trace with the integral of the diagonal fiber trace.
Take the fiber trace in \eqref{heat:uniform-diagonal}. For $\operatorname{rank}E=r$, the absolute error is at most
$rC_T\sqrt{t/k}$. Integrate over the finite-volume manifold $M$.
Thus the leading coefficient of the global trace is the integral over $M$ of the pointwise curvature density.
\end{proof}

\begin{corollary}[Dirichlet boundaries contribute only lower-order terms]
\label{heat:dirichlet-trace}
If $\Omega$ is a fixed relatively compact smooth domain and $P_k$ has its Dirichlet realization,
then \eqref{heat:trace-main-term} remains valid with $M$ replaced by $\Omega$.
For fixed $t$, the remainder remains $O_t(k^{(d-1)/2})$.
\end{corollary}
\begin{proof}
Extend the coefficients to a compact neighborhood of $\overline\Omega$.
Write $\delta(x)=d_g(x,\partial\Omega)$ and $\epsilon=\sqrt{t/k}$.
For $\delta(x)\geq C\epsilon$, apply wave localization on a scaled interior ball whose radius is proportional to $\delta(x)/\epsilon$.
The normalized difference between the Dirichlet and surrounding kernels is at most
$C_Te^{-c\delta(x)^2/\epsilon^2}$.
For $\delta(x)<C\epsilon$, Corollary~\ref{heat:global-gaussian} controls each kernel by
$C_T\epsilon^{-d}$, so increasing the constant gives the same normalized exponential bound.
The normalized diagonal error is therefore at most
\[
 C_T\epsilon+C_Te^{-c\delta(x)^2/\epsilon^2}.
\]
The smooth boundary has finitely many normal collar charts with uniformly bounded volume density, so
$\operatorname{Vol}\{\delta(x)<s\}\leq Cs$ for small $s$.
Integrating $e^{-cu^2/\epsilon^2}$ directly in these coordinates gives $O(\epsilon)$;
outside the collar the contribution is exponentially smaller.
The integrated normalized error is thus $O_T(\epsilon)$, equivalent to the stated remainder.
A fixed regular boundary is required; the same remainder order is not automatic for arbitrary rough domains.
\end{proof}

\section{Comparison on shrinking coordinate balls}

\begin{proposition}[Quadratic-form comparison with a frozen connection]
\label{heat:local-form-comparison}
On a normal coordinate ball of radius $\rho$, freeze the line-bundle curvature and potential in radial gauge.
After conjugating the volume density, write $q_k$ for the actual form and $q_{k,0}$ for the constant magnetic model.
There is $C$, independent of the center, $k$, and sufficiently small $\rho$, such that for every $0<\eta\leq1$ and every $u$ in the Dirichlet form domain,
\[
 |q_k(u)-q_{k,0}(u)|
 \leq C(\rho+\eta)(q_{k,0}(u)+C\|u\|^2)
       +C(\rho+\eta^{-1}k\rho^4+k^{-1})\|u\|^2.
\]
In particular, for $\rho_k=k^{-1/2+a}$ with $0<a<1/4$, there is $\delta_k\to0$ such that
\[
 (1-\delta_k)q_{k,0}-\delta_k\|\cdot\|^2
 \leq q_k\leq
 (1+\delta_k)q_{k,0}+\delta_k\|\cdot\|^2.
\]
\end{proposition}
\begin{proof}
The principal coefficients differ by $O(\rho)$, and the potentials by $O(\rho+k^{-1})$.
The connection difference is $kO(\rho^2)+O(\rho)$, since the linear term of the line connection has already been frozen in radial gauge.
Expand the squared energy. Apply Young's inequality with parameter $\eta$ to the line-connection error and with parameter $\rho$ to the fixed-bundle error. Using
$2|\langle a,b\rangle|\leq\eta|a|^2+\eta^{-1}|b|^2$ gives
\[
 k^{-1}\bigl|\|\nabla_k u\|^2-\|\nabla_{k,0}u\|^2\bigr|
 \leq C(\eta+\rho) k^{-1}\|\nabla_{k,0}u\|^2
       +C\bigl((1+\eta^{-1})k\rho^4+k^{-1}(\rho+\rho^2)\bigr)\|u\|^2.
\]
Adding the metric difference, density potential, and potential variation gives the first estimate;
the model connection energy is $q_{k,0}$ plus a uniformly bounded potential term.
Choose $\eta=\sqrt{k}\rho_k^2=k^{-1/2+2a}$.
Since $a<1/4$, we have $\eta\to0$, while
$\eta^{-1}k\rho_k^4=\eta$, $\rho_k\to0$, and all remaining terms tend to zero.
Combine the finitely many error coefficients into $\delta_k$ to obtain the second estimate.
\end{proof}

\begin{remark}[Spectral comparison and kernel comparison]
Proposition~\ref{heat:local-form-comparison} and Theorem~\ref{heat:minmax}, through Corollary~\ref{heat:form-comparison}, give upper and lower bounds for local spectra and heat traces.
Proposition~\ref{heat:ims} combines these into global spectral estimates; the global operator on the compact manifold has compact resolvent by the elliptic energy estimate in Lemma~\ref{ana:energy} and Theorem~\ref{ana:rellich}.
For a local Dirichlet operator, extend the $H_0^1$ sections by zero into a fixed compact coordinate neighborhood and apply the same compact embedding.
This does not provide a pointwise order relation between matrix kernels.
The quadratic-form Duhamel identity uses the same Taylor remainders to obtain a uniform diagonal remainder.
This explains the relation between the two methods: min--max comparison controls eigenvalues, while Duhamel comparison controls kernels.
In the proof of holomorphic Morse inequalities, we first fix $t$ and let $k\to\infty$, then let $t\to\infty$.
The constant $C_T$ is allowed to depend on $T$, so these limits cannot be interchanged on the strength of the preceding estimates alone.
\end{remark}

\section{A fixed-scale parallel-transport leading term and smoothing positive currents}

\begin{proposition}[The first remainder for the parallel-transport parametrix]
\label{heat:parallel-transport-parametrix}
Let $M$ be a fixed compact $d$-dimensional Riemannian manifold, let $E$ have a smooth unitary connection,
and let $P=\nabla^*\nabla+Q$, with $Q=Q^*$ smooth.
Let $p_t(x,y)$ be the scalar heat kernel. For distances below the normal-coordinate radius, let
$\tau_{y\to x}$ denote parallel transport along the unique short geodesic.
Choose a nonnegative smooth cutoff $\chi(d_g(x,y))$, equal to $1$ near the diagonal and supported in that normal neighborhood.
For $0<t\leq1$,
\[
 H_P(t;x,y)=p_t(x,y)\chi(d_g(x,y))\tau_{y\to x}+R_t(x,y),\qquad
 \|R_t(x,y)\|\leq Ct\,t^{-d/2}e^{-d_g(x,y)^2/(Ct)}.
\]
In the interior of a fixed relatively compact region, any fixed smooth lower-bounded self-adjoint global extension agreeing on a larger neighborhood
has the same local expression. Changing such an extension adds an interior remainder $O(e^{-c/t})$, with constants allowed to depend on the extensions' lower bounds.
\end{proposition}
\begin{proof}
Take normal coordinates centered at $y$, and write $r=d_g(x,y)$ and
$dV_g=r^{d-1}j(r,\theta)\,drd\theta$, with $j(0,\theta)=1$.
Set
\[
 G=(4\pi t)^{-d/2}e^{-r^2/(4t)},\quad
 U(x,y)=j(x,y)^{-1/2}\tau_{y\to x},\quad Z=\chi(r)GU.
\]
The normal-coordinate volume density $j$ and the amplitude $U$ are smooth across the diagonal.
Where $\chi=1$, the radial Laplacian formula gives
\[
 (\partial_t+\Delta_g)G=\frac r{2t}(\partial_r\log j)G.
\]
The product rule also gives
\[
 (\partial_t+P)(GU)
 =G\left(PU+\frac rt\nabla_{\partial_r}U
                   +\frac r{2t}(\partial_r\log j)U\right).
\]
Radial parallel transport satisfies $\nabla_{\partial_r}\tau=0$, hence
$\nabla_{\partial_r}U=-\tfrac12(\partial_r\log j)U$.
The two terms containing $t^{-1}$ cancel exactly, leaving $GPU$.
Since $PU$ is bounded on the fixed normal neighborhood, the residual satisfies
$|R^{(0)}(t;x,y)|\leq Ct^{-d/2}e^{-d_g(x,y)^2/(Ct)}$.
Cutoff derivatives lie away from the diagonal, where every negative power of $t$ is absorbed by exponential decay.

Correct $Z$ by the Volterra construction of Theorem~\ref{heat:scalar-gaussian}.
The residual now has row and column integrals $O(1)$ rather than $O(t^{-1/2})$.
For $j$ residual factors, the time-simplex integral is $t^{j-1}/(j-1)!$;
for $Z*R^{(0)*j}$ it is $t^j/j!$.
The same exponential weights and largest-time-interval estimate give
\[
 |H_P-Z|\leq Ct\,t^{-d/2}e^{-d_g(x,y)^2/(Ct)}.
\]
Repeat this computation for the scalar operator with $Q=0$. Its parametrix is
$Z_{\rm sc}=\chi G j^{-1/2}$, and $|p_t-Z_{\rm sc}|$ has the same first-order remainder bound.
Subtracting the two proves the assertion; the difference between $\chi$ and $\chi^2$ is supported away from the diagonal and also belongs to the remainder.

For the local version, perform the same construction on a larger neighborhood.
The wave operators of the two fixed extensions agree before propagation reaches the exterior. The time-$t$ version of
\eqref{heat:gaussian-functional-calculus} and the interior elliptic estimates
give $O(t^{-N}e^{-c/t})$ by the proof of Proposition~\ref{heat:wave-localization}, with their fixed global lower bounds included in the constants.
Decrease $c$ to write this as $O(e^{-c'/t})$. Thus extension does not affect the stated interior leading term.
\end{proof}

\begin{corollary}[Heat-smoothing error for positive $(1,1)$ currents on a K\"ahler manifold]
\label{heat:positive-current-smoothing}
Let $X$ be a fixed compact K\"ahler manifold of real dimension $2n$, and $T$ a positive $(1,1)$ current.
Let $\Delta_d=dd^*+d^*d$ be the Hodge Laplacian of Proposition~\ref{hod:kahler-laplacians}. Set $T_t=e^{-t\Delta_d}T$, with the smooth kernel acting on the current, and let $\mu=T\wedge\omega^{n-1}/(n-1)!$ be its positive trace measure.
Then $T_t$ is a smooth real $(1,1)$-form, and there is a constant $C$, independent of $T$, such that
\[
 T_t(x)\geq-Ct\left[\int_X t^{-n}e^{-d_g(x,y)^2/(Ct)}\,d\mu(y)\right]\omega_x,
 \qquad0<t\leq1.
\]
If $T$ is closed, then $T_t$ is closed and represents the same de Rham class.
\end{corollary}
\begin{proof}
First verify the connection-Laplacian structure on two-forms. In a normal orthonormal frame at a point,
torsion freeness gives $d=\sum e^j\wedge\nabla_{e_j}$ and $d^*=-\sum\iota_{e_j}\nabla_{e_j}$.
On squaring $d+d^*$, anticommutation of exterior multiplication and contraction produces $-\sum\nabla_{e_j}^2$;
commuting derivatives in distinct directions produces only curvature endomorphisms.
Thus $\Delta_d=\nabla^*\nabla+Q$ with a fixed smooth curvature term $Q$,
so Proposition~\ref{heat:parallel-transport-parametrix} applies.
The K\"ahler condition $\nabla J=0$ makes Levi--Civita parallel transport complex-linear and unitary,
preserving type $(1,1)$, the real structure, and the cone of positive $(1,1)$ forms.
Since $p_t$ and $\chi$ are nonnegative, the leading term sends $T$ to a positive form.
Every matrix coefficient of a positive current is controlled by its trace measure (Theorem~\ref{pos:order-zero}),
so the operator norm of the remainder acting on $T$ is at most
$Ct\int t^{-n}e^{-d_g(x,y)^2/(Ct)}d\mu(y)$.
This gives the lower bound. Smoothness follows from the smooth kernel acting on distributions.

The identity $d\Delta_d=\Delta_dd$ follows directly from $\Delta_d=dd^*+d^*d$ and $d^2=0$,
so $dT=0$ implies $dT_t=0$.
For a closed current, $\partial_tT_t=-dd^*T_t$. Integrating distributionally from $0$ to $t$ gives
$T_t-T=-d\int_0^td^*T_s\,ds$.
The integral may first be taken from $\epsilon$ and then passed to a distributional limit, or defined by spectral calculus.
Thus the cohomology class is unchanged.
Proposition~\ref{hod:kahler-laplacians} says that on a K\"ahler manifold the Hodge Laplacian preserves type, so $T_t$ has type $(1,1)$.
\end{proof}

\section{A polynomial time window}

We now allow heat time to grow with $k$.
The fixed-time remainder does not control this regime because its constant depends on the upper time bound.
A separate argument uses the local model's spectral lower bound to compare the actual decay of the two semigroups.

\begin{lemma}[Long-time comparison retaining the spectral lower bound]
\label{heat:long-time-comparison}
Let two whole-space Laplace-type operators satisfy the short-time Gaussian and covariant-gradient hypotheses of Lemma~\ref{heat:form-duhamel},
with a global unweighted coefficient-difference bound $C\delta$, $0<\delta\leq1$.
If $P_0\geq\sigma$, $P_1\geq\sigma-C\delta$, and $|\sigma|\leq C$, then
\[
 \|H_1(t;0,0)-H_0(t;0,0)\|
 \leq C\delta(1+t)e^{-\sigma t+C\delta t},\qquad t\geq2.
\]
For $0<t\leq2$, the difference is at most $C\delta t^{-d/2}$.
All constants depend only on the stated uniform bounds.
\end{lemma}
\begin{proof}
The short-time conclusion follows from the proof of Lemma~\ref{heat:form-duhamel}, replacing total time $1$ by $t$.
After spatial convolution, the four time integrals have sizes $O(1)$, $O(\sqrt t)$, $O(\sqrt t)$, and $O(t)$,
with common spatial factor $t^{-d/2}$.

For $u\geq1$, the short-time rows $H(1/4;x,\cdot)$ have uniformly bounded $L^2$ norms.
Combining this with spectral decay of the semigroup on $L^2$ gives
\[
 \sup_{x,y}\bigl(\|H_i(u;x,y)\|+\|\nabla_yH_i(u;x,y)\|\bigr)
 \leq Ce^{-(\sigma-C\delta)u},
\]
For $i=0$, the term $C\delta$ can be omitted.
Explicitly, split the semigroup into times $1/4,u-1/2,1/4$.
The spectral lower bound controls the middle factor, while the short-time Gaussian bounds control both ends.
For a $y$ derivative, differentiate only the final factor and use its Gaussian gradient bound.

Likewise, kernel columns and their spatial covariant gradients satisfy
\[
 \|H_i(u;\cdot,0)\|_2+\|\nabla H_i(u;\cdot,0)\|_2
 \leq Ce^{-(\sigma-C\delta)u}.
\]
The gradient term needs an energy argument: uniform ellipticity and bounded potential give
$\|\nabla v\|_2^2\leq C\langle v,(P_i+C_1)v\rangle$.
Substitute $v=e^{-(u-1/2)P_i}H_i(1/2;\cdot,0)$.
In the spectral calculus, the multiplier
$(\lambda+C_1)^{1/2}e^{-(u-1/2)\lambda}$
for $\lambda\geq\sigma-C\delta$ is bounded by
$Ce^{-(\sigma-C\delta)u}$, since $u-1/2\geq1/2$.

Apply the four-term quadratic-form Duhamel identity at time $t$.
For $0<s<1$, the left kernel has time $t-s\geq1$, so use the supremum bound just proved.
The right kernel and its gradient have spatial $L^1$ norms at most $C$ and $Cs^{-1/2}$.
The integral on this interval is therefore at most $C\delta e^{-\sigma t+C\delta t}$.
Treat $t-1<s<t$ symmetrically.
On $1\leq s\leq t-1$, both kernels and gradients have the preceding $L^2$ spectral-decay bounds.
Apply Cauchy--Schwarz spatially and multiply by the time-interval length, at most $t$,
to obtain $C\delta t e^{-\sigma t+C\delta t}$.
Adding the three contributions proves the assertion; integrability of the short-time factor $s^{-1/2}$ has also justified both endpoints.
\end{proof}

\begin{theorem}[A relative leading term in a small-power time window]
\label{heat:polynomial-time-window}
Under Theorem~\ref{heat:scaled-diagonal}, write
\[
 E_0(t,x)=(4\pi t)^{-d/2}
       \prod_j\frac{tb_j(x)}{\sinh(tb_j(x))}\,e^{tV(x)},
 \qquad \varepsilon_*=\frac1{20(d+2)}.
\]
There exist $C,k_0$ such that, for $k\geq k_0$, $x\in M$, and $0<t\leq k^{\varepsilon_*}$,
\begin{equation}
 \|k^{-d/2}H_k(t;x,x)-E_0(t,x)\|
       \leq Ck^{-1/5}\|E_0(t,x)\|.
 \label{heat:long-relative}
\end{equation}
The fiber trace has a relative estimate of the same order.
For scalar kernels this is the usual relative error; for matrix kernels it is a norm comparison, not a multiplicative matrix identity.
\end{theorem}
\begin{proof}
Now fix the spatial scaling $x-x_0=k^{-1/2}Z$ and keep heat time $t$ variable.
In $Z$ space choose $R_k=k^{1/20}$, retaining the original coefficients inside radius $R_k$
and replacing them outside radius $2R_k$ by the constant-field model at the center.
The physical radius is $k^{-9/20}$.
Radial gauge and Taylor expansion give the global unweighted difference bound
\[
 |a_k-I|+|A_k-A_0|+|Q_k-Q_0|
 \leq C\bigl(k^{-1/2}R_k+k^{-1/2}R_k^2+k^{-1}\bigr)
 \leq Ck^{-2/5}=C\delta_k.
\]
Each cutoff derivative contributes $R_k^{-1}\leq1$, so the short-time geometric, curvature, potential, and covariant-derivative bounds remain uniform.

Set $B_j=|b_j(x_0)|$, let $v_{\max}$ be the largest eigenvalue of $V(x_0)$, and put
$\sigma=\sum_jB_j-v_{\max}$.
The model connection energy is bounded below by $\sum_jB_j$. On one magnetic plane,
$\| (D_1+\ii\operatorname{sgn}(b)D_2)u\|^2
=\|D_1u\|^2+\|D_2u\|^2-|b|\|u\|^2\geq0$.
Sum over the planes; free directions contribute nonnegative energy.
Thus $P_0\geq\sigma$.
Expand the difference of the two energy forms and use $2ab\leq\delta_k a^2+\delta_k^{-1}b^2$:
\[
 q_k^{\rm ext}(u)\geq(1-C\delta_k)\|\nabla_0u\|^2
                       -\langle u,V(x_0)u\rangle-C\delta_k\|u\|^2
 \geq(\sigma-C'\delta_k)\|u\|^2.
\]
The quantity $|\sigma|$ is uniformly bounded over all centers.
Lemma~\ref{heat:long-time-comparison} therefore bounds the extension-model kernel difference by
\[
 C\delta_k(1+t)e^{-\sigma t+C\delta_k t},\qquad t\geq2.
\]
On the other hand, $2u/(1-e^{-2u})\geq1$ for $u\geq0$, so the exact model formula gives
\[
 \|E_0(t,x_0)\|\geq(4\pi t)^{-d/2}e^{-\sigma t}.
\]
Division yields the relative error
$C\delta_k(1+t)t^{d/2}e^{C\delta_k t}$.
For $t\leq k^{\varepsilon_*}$, its power is at most
$k^{-2/5+\varepsilon_*(d/2+1)}=k^{-3/8}$, and $e^{C\delta_k t}$ stays bounded.
For $0<t\leq2$, the short-time estimate gives relative error $C\delta_k$.

We must also compare the original manifold kernel with the extension.
The time-$t$ wave-localization proof bounds the scaled diagonal difference by
\[
 C(1+t)^N\exp(Ct-cR_k^2/t)
\]
To obtain this estimate, replace the Gaussian in \eqref{heat:gaussian-functional-calculus} by $(4\pi t)^{-1/2}e^{-s^2/(4t)}$ and perform the same finite number of integrations by parts on the tail. At small times, any additional inverse powers of $t$ are absorbed by decreasing the Gaussian exponent constant.
The two global spectral lower bounds and the central interior elliptic constants are uniform.
Dividing by the model lower bound adds only $Ct^{d/2}e^{\sigma t}$.
Since $R_k^2=k^{1/10}$ and $\varepsilon_*<1/20$,
\[
 Ck^{\varepsilon_*}-c k^{1/10-\varepsilon_*}\longrightarrow-\infty
\]
has negative growth in a positive power of $k$, dominating every multiple of $\log k$. This relative error is smaller than any prescribed $k^{-N'}$.
Together with the $O(k^{-3/8})$ freezing error, it is certainly at most $Ck^{-1/5}$.

Finally, the model is a positive definite matrix, so $\tr E_0\geq\|E_0\|$,
while every matrix $A$ satisfies $|\tr A|\leq\operatorname{rank}(E)\|A\|$.
The trace therefore has a relative error of the same order.
\end{proof}

\begin{remark}
The stated exponent $\varepsilon_*$ leaves room in both error estimates.
The model's spectral lower bound controls the freezing error, while the propagation factor $e^{-cR_k^2/t}$ dominates the exterior growth $e^{Ct}$.
Holomorphic Morse inequalities require only the two successive limits described above; this growing-time estimate is a further refinement.
\end{remark}

\begin{corollary}[Direct diagonal comparison on arbitrary intermediate shrinking balls]
\label{heat:shrinking-ball-comparison}
For fixed smooth geometric data, take $\rho_k=k^{-1/2+a}$ with $0<a<1/2$.
Let $H^D_k,H^{D,0}_k$ be the actual and frozen-model Dirichlet kernels on the same coordinate ball of radius $\rho_k$.
For each fixed $T>0$,
\[
 \|H^D_k(t;0,0)-H^{D,0}_k(t;0,0)\|
 \leq C_{T,a}k^{d/2-1/2}t^{-d/2},\qquad0<t\leq T.
\]
\end{corollary}
\begin{proof}
By \eqref{heat:uniform-diagonal}, the original global kernel and the frozen whole-space model differ by at most
$C_T(k/t)^{d/2}\sqrt{t/k}$.
Compare each global kernel with its small-ball Dirichlet kernel. The scaled boundary distance is
$\rho_k\sqrt{k/t}=k^a/\sqrt t$.
Wave localization bounds each normalized boundary error by $C_Te^{-ck^{2a}/t}$.
For $0<t\leq T$, this is uniformly $O_{T,a}(k^{-1/2})$,
while $\sqrt{t/k}\leq\sqrt T k^{-1/2}$.
Adding the three terms proves the result.
The restriction $a<1/4$ in Proposition~\ref{heat:local-form-comparison}
is needed only when requiring the unweighted form error to vanish throughout the whole shrinking ball.
\end{proof}

\begin{proposition}[Interior first and second derivatives for constant-coefficient parabolic equations]
\label{heat:constant-parabolic-derivatives}
Suppose $u$ solves on $Q_R=(-R^2,0]\times B(0,R)$ the equation
\[
 u_t-a^{ij}u_{ij}-b^iu_i=cu,
\]
with constant coefficients $a,b,c$ and $\lambda I\leq a\leq\Lambda I$.
Then
\[
 \sup_{Q_{R/2}}|Du|\leq CR^{-1}\sup_{Q_R}|u|,\qquad
 \sup_{Q_{R/2}}|D^2u|\leq CR^{-2}\sup_{Q_R}|u|,
\]
where $C$ depends only on $d,\lambda,\Lambda,|b|R,|c|R^2$.
\end{proposition}
\begin{proof}
Set $v(s,z)=u(R^2s,Rz)$ to reduce to $R=1$, replacing $b,c$ by $Rb,R^2c$.
For $\mathcal P=\partial_t-a^{ij}\partial_i\partial_j-b^i\partial_i$, direct calculation gives
\[
 \mathcal P|u|^2\leq2|c||u|^2-2\lambda|Du|^2,\qquad
 \mathcal P|Du|^2\leq2|c||Du|^2-2\lambda|D^2u|^2.
\]
The second inequality uses constancy of the coefficients: each $u_j$ solves the same equation.
The scalar kernel is a Gaussian of covariance $2ta$ shifted by the drift $bt$.
Its bound and first-derivative bound for time at most one depend only on $\lambda,\Lambda,|b|$.
Thus the cutoff and mean-value proof of Lemma~\ref{heat:parabolic-energy} applies without change,
apart from the additional drift contribution in integration by parts,
$C|b|\,|\eta D\eta|\,|u|^2$, controlled by the fixed cutoff.
Apply energy estimates between $Q_1$ and an intermediate cylinder strictly containing $Q_{3/4}$, then the mean-value estimate to $|Du|^2+|u|^2$,
obtaining $\sup_{Q_{3/4}}|Du|\leq C\sup_{Q_1}|u|$.
Each $u_j$ still solves the same equation. Apply the first-derivative estimate again between $Q_{3/4}$ and $Q_{1/2}$:
$\sup_{Q_{1/2}}|D^2u|\leq C\sup_{Q_{3/4}}|Du|$.
Rescaling proves both estimates.
Constancy of the coefficients is used when differentiating the equation: each derivative again solves the same equation.
Variable coefficients create additional terms involving their derivatives.
\end{proof}

\begin{proposition}[Kernel derivatives of every order for fixed geometric data]
\label{heat:fixed-derivative-gaussian}
Let $P=\nabla^*\nabla+Q$ be a smooth self-adjoint Laplace-type operator on a fixed compact manifold.
For all nonnegative integers $a,b$, there is $C_{a,b}$ such that, for $0<t\leq1$,
\[
 \|\nabla_x^a\nabla_y^bH_P(t;x,y)\|
 \leq C_{a,b}t^{-(d+a+b)/2}e^{-d_g(x,y)^2/(C_{a,b}t)}.
\]
Replacing covariant derivatives by ordinary coordinate derivatives changes only constants and finitely many lower-order terms.
\end{proposition}
\begin{proof}
Fix an integer $m$.
Successive commutation of covariant derivatives writes the equation for $\nabla^j u$ as
\[
 (\partial_t+\nabla^*\nabla)(\nabla^ju)
 =\sum_{\ell=0}^j A_{j\ell}\nabla^\ell u,\qquad0\leq j\leq m.
\]
The $A_{j\ell}$ are combinations of the fixed curvature, potential, and finitely many covariant derivatives, so they are uniformly bounded.
This structure follows by induction: commuting the outermost $\nabla$ past the connection Laplacian
produces only curvature times a first derivative and differentiated curvature times the original section;
differentiate the existing lower-order coefficients by the Leibniz rule.
Thus the right-hand side contains no uncontrolled derivative above the order of $\nabla^ju$.

Pair these equations with cutoff multiples of $\nabla^ju$ and absorb the gradient mixed terms
as in the proof of \eqref{heat:caccioppoli}.
Proceed through $m+1$ nested parabolic cylinders to obtain
\[
 \int_{Q_r}|\nabla^j u|^2
 \leq C_mr^{-2j}\int_{Q_{2r}}|u|^2,\qquad0\leq j\leq m.
\]
At each step shrink the radius by a fixed ratio; all ratios may be chosen in advance depending on $m$.
A Bochner calculation for $w_m=\sum_{j=0}^m|\nabla^ju|^2$,
bounding every mixed term by $2ab\leq a^2+b^2$, gives
$(\partial_t+\Delta_g)w_m\leq C_mw_m$.
Lemma~\ref{heat:parabolic-energy} and the preceding energy estimates give
\[
 |\nabla^mu(t,x)|^2
 \leq C_mr^{-d-2-2m}\int_{Q_{2r}}|u|^2.
\]
Take $r=c_m\sqrt t$ and $u(s,z)=H_P(s;z,y)v$, with all cylinder times in $[t/2,t]$.
Kato's Gaussian bound controls the right-hand side and, as in the first-gradient proof, gives
$C_mt^{-d-m}e^{-d_g(x,y)^2/(C_mt)}$.
Take square roots and the supremum over unit $v$ for every $x$-derivative order;
self-adjointness gives every $y$-derivative order.
Finally, factor $H(t)=H(t/2)H(t/2)$, putting $a$ and $b$ endpoint derivatives on the two factors.
Lemma~\ref{heat:manifold-convolution}, on this fixed compact manifold, gives the asserted time power and exponential bound for mixed derivatives.
\end{proof}


\chapter{Holomorphic Morse Inequalities}
\label{morse:chapter}

Starting from the local scaling theorem proved in the preceding chapter, we establish the weak and strong holomorphic Morse inequalities on any compact complex Hermitian manifold,
for a smooth Hermitian holomorphic line bundle twisted by a fixed holomorphic vector bundle.
No K\"ahler assumption is needed: non-K\"ahler torsion remains in bounded lower-order terms.
The two points requiring particular care are curvature and heat-time normalization, and the order of the two limits.

\section{The long-time limit of the curvature model}
\label{heat:model-limit-section}

Use $F^L=\Theta^L$ for the Chern curvature of $L$ (Theorem~\ref{bun:chern-connection}); its real curvature is $\ii F^L$.
In a tangent space of complex dimension $n$, choose a real orthonormal coframe $dx_j,dy_j$ adapted to the complex structure, put $z_j=x_j+\ii y_j$ at that point, and write the curvature as
\[
\ii\Theta^L=\sum_{j=1}^n b_j\,dx_j\wedge dy_j,
\]
where $b_j\in\R$, and choose real volume form
$dV=dx_1\wedge dy_1\wedge\cdots\wedge dx_n\wedge dy_n$.
A basic expression for the constant-coefficient model in degree $(0,q)$ is
\begin{equation}
\Psi_q(t;b)=\frac1{(4\pi t)^n}
 \prod_{j=1}^n\frac{b_jt}{\sinh(b_jt)}
 \sum_{|J|=q}
 \exp\left(t\left(\sum_{j\notin J}b_j-\sum_{j\in J}b_j\right)\right).
\label{heat:psi-definition}
\end{equation}
Here $J$ runs through increasing $q$-element subsets of $\{1,\ldots,n\}$, and $d\bar z_J$ denotes their wedge product. Factors with zero parameter are interpreted by continuous extension.
This follows from Corollary~\ref{heat:product-model}:
the endomorphism $V$ in the model's potential $-V$ has value on $d\bar z_J$
$\sum_{j\notin J}b_j-\sum_{j\in J}b_j$.
The long-time limit identifies which curvature signs contribute in degree $q$.

\begin{proposition}[The model selects the negative curvature index]
\label{heat:model-limit}
If every $b_j\ne0$ and exactly $q_-$ are negative, then
\[
\lim_{t\to\infty}\Psi_q(t;b)
=\begin{cases}
 (2\pi)^{-n}\prod_j|b_j|,&q=q_-,\\
 0,&q\ne q_-.
\end{cases}
\]
If at least one $b_j=0$, the limit is zero for every $q$.
Moreover, for $|b_j|\leq M$ and $t\geq1$, there is the uniform bound
\begin{equation}
0\leq\Psi_q(t;b)
\leq\binom nq\frac{(2M+1)^n}{(4\pi)^n}.
\label{heat:model-dominating-bound}
\end{equation}
\end{proposition}
\begin{proof}
First suppose all $b_j\ne0$.
Each factor satisfies
\[
\frac{b_j}{\sinh(b_jt)}
=\frac{2|b_j|e^{-|b_j|t}}{1-e^{-2|b_j|t}}.
\]
For a fixed $J$, the total exponent is therefore
\[
t\left(\sum_{j\notin J}b_j-\sum_{j\in J}b_j-
\sum_j|b_j|\right)\leq0.
\]
Equality holds exactly when $J$ is the set of indices of negative eigenvalues.
In that case, the remaining denominators tend to $1$, and the constant is
$2^n(4\pi)^{-n}\prod_j|b_j|$. In all other cases the exponential decays strictly.
This proves the nondegenerate case.

If there are $\ell\geq1$ zero parameters, each contributes $t^{-1}$ to
\eqref{heat:psi-definition},
with zero contribution to the exponent.
The other factors either decay exponentially or tend to finite constants, as above.
The extra factor $t^{-\ell}$ therefore makes every term tend to zero.

Finally, for $u>0$, the inequality $e^{2u}-1\geq2u$ gives
\[
\frac{e^u}{\sinh u}
=\frac2{1-e^{-2u}}
=2+\frac2{e^{2u}-1}\leq2+\frac1u.
\]
For $b\ne0$, consequently,
$|b|e^{t|b|}/\sinh(t|b|)\leq2|b|+t^{-1}\leq2M+1$.
At $b=0$, the continuous value is $t^{-1}\leq1$.
Apply this bound to each factor of each $J$ term,
then sum over the $\binom nq$ subsets $J$ to obtain the uniform estimate.
\end{proof}

The number of negative eigenvalues need not be constant across the manifold.
Bound \eqref{heat:model-dominating-bound} remains valid when eigenvalues cross zero,
so dominated convergence applies to the model densities of a smooth curvature form on a compact manifold.
There is no need to split the proof into separate local rank cases.


\section{Why the same model applies on every Hermitian manifold}

\begin{proposition}[The leading curvature term of the Dolbeault operator]
\label{morse:weitzenbock}
Let $X$ be a compact Hermitian complex manifold, and $L,F$ Hermitian holomorphic bundles with $\operatorname{rank}L=1$.
Here $\Delta''=\dbar\dbar^*+\dbar^*\dbar$ is the Dolbeault Laplacian of Theorem~\ref{hod:hodge}, and $\Lambda^{0,\bullet}T^*X=\bigoplus_{q=0}^n\Lambda^{0,q}T^*X$. On $W=\Lambda^{0,\bullet}T^*X\otimes F$, there are a fixed unitary connection $\nabla^W$ and a fixed smooth Hermitian endomorphism $R$ such that
\begin{equation}
 2\Delta''_{L^k\otimes F}
 =(\nabla^{L^k}\otimes1+1\otimes\nabla^W)^*
  (\nabla^{L^k}\otimes1+1\otimes\nabla^W)+kQ_L+R.
 \label{morse:bochner-form}
\end{equation}
At a point choose a real orthonormal frame adapted to the complex structure, with
$\ii F^L=\sum_jb_jdx_j\wedge dy_j$. Let $N_j$ be the orthogonal projection onto forms containing the $j$th antiholomorphic basis form; on $d\bar z_J$ it equals one for $j\in J$ and zero otherwise. On the $(0,q)$ component indexed by $J$,
\[
 Q_L=\sum_{j=1}^nb_j(2N_j-1),\qquad
 Q_L|_{d\bar z_J\otimes F}
 =\left(\sum_{j\in J}b_j-\sum_{j\notin J}b_j\right)\operatorname{Id}_F.
\]
All operators preserve form degree.
\end{proposition}
\begin{proof}
Set $D_k=\sqrt2(\dbar+\dbar^*)$, so $D_k^2=2\Delta''$.
At the chosen point take the unit antiholomorphic coframe $\theta^j=d\bar z_j/\sqrt2$.
Write $e_j$ for exterior multiplication by $\theta^j$, and $\iota_j$ for its adjoint.
They satisfy
$e_j\iota_\ell+\iota_\ell e_j=\delta_{j\ell}$,
while exterior multiplications anticommute with each other, as do contractions.
The first-order part at the point is
\[
 D_k=2\sum_j(e_j\nabla_{\bar z_j}-\iota_j\nabla_{z_j})+B,
\]
where $B$ depends only on the fixed metric, frame, and bundle $F$.
In real directions the Clifford coefficients are
\[
 c_{x_j}=e_j-\iota_j,\qquad c_{y_j}=\ii(e_j+\iota_j),\qquad
 c_ac_b+c_bc_a=-2\delta_{ab}.
\]
Thus $D_k^2$ has scalar second-order principal part $-g^{ij}\partial_i\partial_j$.

We next recover the connection Laplacian uniquely from the first-order coefficients.
In a local orthonormal frame, with volume density $\rho$, write a self-adjoint Laplace-type operator as
$-g^{ij}\partial_i\partial_j+B^j\partial_j+C$.
For a connection $d+A$, the first-order coefficient of its connection Laplacian is
\[
 -\rho^{-1}\partial_i(\rho g^{ij})\operatorname{Id}-2g^{ij}A_i.
\]
This determines $A_i$ uniquely.
Formal self-adjointness imposes
$B^j+(B^j)^*=-2\rho^{-1}\partial_i(\rho g^{ij})\operatorname{Id}$,
so the recovered $A_i$ are skew-Hermitian and define a unitary connection.
Subtracting its connection Laplacian leaves a Hermitian zeroth-order term.

In the first-order expression for $D_k$, the parameter $k$ occurs only through the line connection $kA^L$.
In the first-order coefficients of the square, its contribution is exactly $-2g^{ij}kA_i^L$.
The recovered connection is therefore $kA^L+A^W$, with $A^W$ independent of $k$.
On subtraction, the quadratic terms $k^2A^LA^L$ and all undifferentiated mixed terms involving $kA^L$ cancel.
Choose a gauge with $A^L=0$ at the point. The remaining terms linear in $k$ come only from commuting covariant derivatives, and equal
\[
 k\sum_{a<b}c_ac_bF^L_{ab}.
\]
This expression is gauge invariant, so the pointwise computation identifies the global leading curvature term.
Curvature of the fixed bundle, metric derivatives, and Hermitian torsion contribute only to the $k$-independent endomorphism $R$.

Since $F^L$ has type $(1,1)$, diagonalize it unitarily at the point.
Then $F^L_{x_jy_j}=-\ii b_j$, and
\[
 c_{x_j}c_{y_j}
 =\ii(e_j\iota_j-\iota_je_j)=\ii(2N_j-1).
\]
Multiplication gives $b_j(2N_j-1)$; summing over $j$ proves the formula.
Finally, $D_k^2=2\Delta''$ preserves degree, so its first-order coefficients are block diagonal by degree,
and the recovered connection and remaining potential also preserve degree.
The argument applies to every Hermitian metric, with its torsion included in the fixed terms.
\end{proof}

\begin{theorem}[The uniform leading term of the Dolbeault heat kernel]
\label{morse:local-heat}
Let $K_{k,q}(t;x,y)$ be the kernel of
$\exp(-2t\Delta''_{L^k\otimes F,q}/k)$, and set $r=\operatorname{rank}F$.
For fixed $t>0$, uniformly in $x\in X$,
\begin{equation}
 k^{-n}\tr K_{k,q}(t;x,x)=r\Psi_q(t;b(x))+O_t(k^{-1/2}).
 \label{heat:needed-local-asymptotic}
\end{equation}
For $0<t\leq T$, the remainder after multiplication by $t^n$ is $O_T(\sqrt{t/k})$.
\end{theorem}
\begin{proof}
Divide \eqref{morse:bochner-form} by $k$ to obtain
$P_k=k^{-1}\nabla_k^*\nabla_k-V+k^{-1}R$, with fixed smooth endomorphism $V=-Q_L$.
These are exactly the hypotheses of Theorem~\ref{heat:scaled-diagonal}.
The model magnetic parameters are $b_j(x)$, and the model endomorphism $V$ on the $J$ component is
$\sum_{j\notin J}b_j-\sum_{j\in J}b_j$, acting scalarly on $F_x$.
Thus $\tr e^{tV(x)}$ equals
$r\sum_{|J|=q}e^{t(\sum_{j\notin J}b_j-\sum_{j\in J}b_j)}$.
Substitute in \eqref{heat:uniform-diagonal} and use real dimension $2n$ to obtain
\eqref{heat:needed-local-asymptotic} with the stated uniform remainder.
\end{proof}

\begin{remark}[A fixed zeroth-order term cannot be factored out arbitrarily]
In Proposition~\ref{morse:weitzenbock}, the endomorphism $R$ contains all fixed geometric terms.
For $P=P_0+k^{-1}R$ with $R=c\operatorname{Id}$ constant, one has
$e^{-tP}=e^{-ct/k}e^{-tP_0}$.
If $R$ does not commute with $P_0$, even this multiplicative factorization generally fails.
Proposition~\ref{heat:bounded-difference} supplies the appropriate operator perturbation formula; its kernel version is derived independently in Proposition~\ref{heat:kernel-duhamel}.
\end{remark}

\section{Heat-trace inequalities for the infinite-dimensional complex}

\begin{proposition}[Weak, strong, and adjacent-degree heat-trace estimates]
\label{morse:spectral-inequalities}
Abbreviate $\Delta''_{k,q}=\Delta''_{L^k\otimes F,q}$. Write $h_k^q=\dim_{\C}H^q(X,L^k\otimes F)$ and
$T_k^q(t)=\tr\exp(-2t\Delta''_{k,q}/k)$, setting both to zero outside the degree range.
For every $t>0$,
\begin{align*}
 h_k^q&\leq T_k^q(t),\\
 \sum_{j=0}^q(-1)^{q-j}h_k^j
 &\leq\sum_{j=0}^q(-1)^{q-j}T_k^j(t),\\
 h_k^q&\geq T_k^q(t)-T_k^{q-1}(t)-T_k^{q+1}(t).
\end{align*}
The strong inequality is an equality for $q=n$.
\end{proposition}
\begin{proof}
The Dolbeault comparison of Chapter~\ref{sh:chapter} and Theorem~\ref{hod:hodge} identify $h_k^q$ with the dimension of $\ker\Delta''_{k,q}$.
Since $\dbar\Delta''=\Delta''\dbar$, the eigenspaces of a fixed eigenvalue $\lambda$ in all degrees form a finite-dimensional complex.
For $\lambda>0$, a closed eigenform satisfies
$u=\lambda^{-1}\Delta''u=\dbar(\lambda^{-1}\dbar^*u)$, so this complex is exact.
This is the algebraic mechanism of Proposition~\ref{heat:finite-heat-morse}.
Apply the finite-dimensional strong inequality at each positive eigenvalue, multiply by $e^{-2t\lambda/k}$, and sum for the second assertion.
The first follows immediately from the nonnegative positive-eigenvalue contribution to the heat trace.
In top degree, every positive-eigenvalue subcomplex has Euler characteristic zero, giving equality.

For the third inequality, let an exact eigencomplex have dimension $c_j$ and boundary-space dimension $b_j$ in degree $j$.
Then $c_q=b_q+b_{q+1}$, with $b_q\leq c_{q-1}$ and $b_{q+1}\leq c_{q+1}$.
Hence $c_q-c_{q-1}-c_{q+1}\leq0$.
Multiply by the heat weight and sum over the positive spectrum. The zero-spectrum contribution is
$h_k^q-h_k^{q-1}-h_k^{q+1}\leq h_k^q$, proving the third inequality.

All sums converge absolutely: positive-time heat kernels are smooth, and on the compact manifold
$e^{-t\Delta''/k}$ is Hilbert--Schmidt,
so its square $e^{-2t\Delta''/k}$ is trace class.
Equivalently, Proposition~\ref{heat:kernel-trace} gives the trace as a finite diagonal integral.
Thus the finite alternating sum over degrees may be interchanged with the spectral sum.
\end{proof}

\section{Weak and strong Morse inequalities and asymptotic Riemann--Roch}
\label{heat:classical-morse-section}

Set
\[
 \alpha=c_1(L,h)=\frac{\ii}{2\pi}F^L,
 \qquad X(q)=\{x:\alpha_x\text{ has }q\text{ negative and }n-q\text{ positive eigenvalues}\}.
\]
The eigenvalues are taken relative to the chosen Hermitian metric; their positive and negative counts are independent of that choice. The sets $X(q)$ are open, and degenerate points lie in none of them.
We use $\chi(X,E)=\sum_{q=0}^n(-1)^q\dim_{\C}H^q(X,E)$ for the holomorphic Euler characteristic.
Write $X(\leq q)=\bigcup_{j=0}^qX(j)$, with out-of-range index sets understood as empty.
In an orthonormal frame,
\begin{equation}
 \frac{\alpha^n}{n!}=\frac{\prod_jb_j}{(2\pi)^n}\,dV_g.
 \label{morse:top-curvature}
\end{equation}
Thus the degenerate locus may have positive measure, but $\alpha^n$ vanishes there.

\begin{theorem}[Holomorphic Morse inequalities on compact complex manifolds]
\label{morse:main}
Let $X$ be a compact complex manifold of dimension $n$, let $L$ be a holomorphic line bundle with a smooth Hermitian metric,
and let $F$ be a holomorphic vector bundle of rank $r$. As $k\to\infty$,
\begin{align}
 h^q(X,L^k\otimes F)
 &\leq\frac{rk^n}{n!}\int_{X(q)}(-1)^q\alpha^n+o(k^n),
 \label{heat:classical-weak}\\
 \sum_{j=0}^q(-1)^{q-j}h^j(X,L^k\otimes F)
 &\leq\frac{rk^n}{n!}\int_{X(\leq q)}(-1)^q\alpha^n+o(k^n),
 \label{heat:classical-strong}\\
 h^q(X,L^k\otimes F)
 &\geq\frac{rk^n}{n!}\int_{X(q-1)\cup X(q)\cup X(q+1)}
                       (-1)^q\alpha^n+o(k^n).
 \label{morse:adjacent-lower}
\end{align}
Moreover,
\begin{equation}
 \chi(X,L^k\otimes F)
 =\frac{rk^n}{n!}\int_X\alpha^n+O(k^{n-1/2}).
 \label{heat:classical-rr}
\end{equation}
The last remainder is the order obtained directly by this proof; the full Riemann--Roch formula provides finer polynomial information.
\end{theorem}
\begin{proof}
\emph{First fix the heat time.}
For fixed $t>0$, divide the inequalities in Proposition~\ref{morse:spectral-inequalities} by $k^n$.
The diagonal trace formula and Theorem~\ref{morse:local-heat} make each normalized heat trace converge to
$r\int_X\Psi_q(t;b(x))\,dV_g$.
For example, the strong inequality gives
\[
 \limsup_{k\to\infty}k^{-n}
       \sum_{j=0}^q(-1)^{q-j}h_k^j
 \leq r\int_X\sum_{j=0}^q(-1)^{q-j}\Psi_j(t;b(x))\,dV_g.
\]
The left-hand side no longer depends on $t$.

\emph{Then let the heat time tend to infinity.}
Compactness bounds all $b_j(x)$ uniformly.
The majorant of Proposition~\ref{heat:model-limit} and finite volume permit dominated convergence term by term.
On $X(\ell)$, only the degree-$\ell$ model has a nonzero limit, namely
$(2\pi)^{-n}|\prod b_j|$; at degenerate points the limit is zero.
Thus on $X(\ell)$ with $\ell\leq q$, the strong alternating sum has limit
\[
 (-1)^{q-\ell}(2\pi)^{-n}|\prod b_j|
 =(-1)^q(2\pi)^{-n}\prod b_j,
\]
and it vanishes elsewhere. Using \eqref{morse:top-curvature} proves
\eqref{heat:classical-strong}.
Apply the same argument to one heat trace for \eqref{heat:classical-weak}.
For the adjacent-degree estimate, use $\liminf$. Its model limit is the positive degree-$q$ contribution minus
the degree-$(q-1)$ and degree-$(q+1)$ contributions, exactly the integral in \eqref{morse:adjacent-lower}.
Writing these $\limsup$ and $\liminf$ conclusions with an $o(k^n)$ remainder is a direct restatement of their definitions.

\emph{The full alternating sum can also be computed at finite heat time.}
At each $x$, the multiplicative identity in the exterior algebra gives
\[
 \sum_{q=0}^n(-1)^q\sum_{|J|=q}
       e^{t(\sum_{j\notin J}b_j-\sum_{j\in J}b_j)}
 =\prod_{j=1}^n(e^{tb_j}-e^{-tb_j}).
\]
Multiply by the common factor in \eqref{heat:psi-definition} to obtain the exact identity
\[
 \sum_{q=0}^n(-1)^q\Psi_q(t;b)=(2\pi)^{-n}\prod_jb_j.
\]
If some $b_j=0$, both sides vanish; the identity also follows by continuous extension.
Equality for the full alternating sum of the spectral complex identifies the Euler characteristic with the alternating heat trace.
Fix, for example, $t=1$, insert the local heat-kernel formula with uniform $O(k^{-1/2})$ remainder, and integrate
to obtain \eqref{heat:classical-rr}.
The full alternating sum therefore integrates the signed top-degree form over all of $X$.
\end{proof}

\begin{corollary}[Leading-order section growth for semipositive line bundles]
\label{morse:semipositive-growth}
If $\alpha\geq0$, then $h^q(X,L^k\otimes F)=o(k^n)$ for $q\geq1$, and
\[
 h^0(X,L^k\otimes F)=\frac{rk^n}{n!}\int_X\alpha^n+o(k^n).
\]
In particular, this applies to ample line bundles.
\end{corollary}
\begin{proof}
Semipositive curvature has no negative eigenvalues, so $X(q)=\varnothing$ for $q\geq1$.
Inequality~\eqref{heat:classical-weak} and nonnegativity of dimensions give $o(k^n)$ in higher degrees.
Move the finitely many higher-degree terms to the other side of asymptotic Riemann--Roch to get the formula for $h^0$.
If $L$ is ample, some $L^m$ defines a projective embedding.
Pull back the positive Fubini--Study metric of Proposition~\ref{bun:fs} to $L^m$, then take its positive $m$th root in each local frame.
This defines a smooth positive metric on $L$; the squared-modulus transition rule makes the local roots compatible.
The ample case is therefore a special case of the semipositive one.
\end{proof}

\begin{remark}[Leading asymptotics and actual vanishing differ]
An $o(k^n)$ term need not be zero. The corollary establishes leading-order growth, which already suffices for the section asymptotics of ample line bundles.
Vanishing of higher cohomology for all sufficiently large $k$ requires quantitative solving with positive curvature
(Corollary~\ref{pos:high-power-solving} and its vector-bundle version); a small error cannot simply be set equal to zero.
\end{remark}

\section{Coherent coefficients: reduction by support dimension}

Passing from vector bundles to arbitrary coherent sheaves cannot assume global vector-bundle resolutions on every compact complex manifold.
We instead use modifications and induction on support dimension, accounting also for the layers of a nonreduced support.
The extension uses two geometric theorems described in Appendix~\ref{dep:chapter}:

\begin{enumerate}
\item \emph{Analytic resolution of singularities and local freeness after modification.}
For a reduced compact complex analytic space and a coherent sheaf on it, choose on each top-dimensional irreducible component a proper modification from a smooth space
such that pullback modulo torsion is locally free. The modification may remain an isomorphism on a dense open set where the original space is smooth and the sheaf locally free.
This uses Hironaka's resolution and flattening theory, recorded in paragraph~\ref{dep:resolution}.
\item \emph{Grauert's proper mapping theorem.}
All higher direct images of a coherent sheaf under a proper holomorphic map are coherent; see paragraph~\ref{dep:proper-image}.
For analytic resolution, see \cite{ArocaHironakaVicente2018};
for flattening, see \cite{Hironaka1975}. The direct-image theorem must be read together with its original paper \cite{Grauert1960}
and correction \cite{Grauert1963}.
\end{enumerate}

These theorems reduce the problem to vector bundles on smooth compact manifolds.
The remaining sheaves have smaller-dimensional support; induction will show that their cohomology contributes a lower-order error.

For a sheaf $E$ on $Y$, $R^q\pi_*E$ denotes the $q$th higher direct image: choose a flabby resolution $I^\bullet$ and take the cohomology sheaf $\mathcal H^q(\pi_*I^\bullet)$.
The resolution comparison in Chapter~\ref{sh:chapter} makes this independent of the chosen resolution.

\begin{lemma}[Dimension control in direct-image comparison]
\label{morse:leray-comparison}
Let $\pi:Y\to X$ be a proper holomorphic map, $E$ a coherent sheaf, and $L$ a line bundle on $X$.
For the fixed degree $m$ under consideration, assume $H^m(Y,\pi^*L^k\otimes E)$ is finite-dimensional for every $k$. Suppose, for every relevant $p$ and $q>0$,
\[
 h^p(X,L^k\otimes R^q\pi_*E)=O(k^e),
\]
Then for this degree $m$,
\[
 h^m(Y,\pi^*L^k\otimes E)
 =h^m(X,L^k\otimes\pi_*E)+O(k^e).
\]
\end{lemma}
\begin{proof}
First explain the cohomological comparison being used.
Take a Godement flabby resolution $I^\bullet$ of $E$ and put $C^q=\pi_*I^q$.
Direct image preserves flabbiness because restrictions on $X$ are restrictions between inverse-image open sets.
The cohomology sheaves of $C^\bullet$ are $R^q\pi_*E$, and
$H^m\Gamma(X,C^\bullet)=H^m(Y,E)$.
Resolve $C^\bullet$ termwise by Godement resolutions to obtain the first-quadrant double complex
$D^{p,q}=\Gamma(X,G^p(C^q))$.
The Godement functors and their global-section functors here are exact: they are built from products of stalks and successive cokernels,
with exactness proved step by step in Proposition~\ref{sh:godement-exact}.

Taking cohomology first in the $p$ direction leaves only $p=0$, since each $C^q$ is flabby; total cohomology is $H^m(Y,E)$.
Taking first the $q$ direction and then the $p$ direction gives second page
\[
 E_2^{p,q}=H^p(X,R^q\pi_*E).
\]
Concretely, filter the total complex by the first index. The lowest component of a closed element must first be closed in the
$q$ direction; if it is a boundary, subtract a total boundary and move to the next column.
Representatives allowing $r-1$ successive eliminations form the $r$th page.
The next obstruction is a linear map $d_r$ from $(p,q)$ to $(p+r,q-r+1)$.
The total differential squares to zero, so $d_r^2=0$; the next page is its kernel modulo its image.
At fixed total degree there are only finitely many first-quadrant positions, so elimination eventually stops.
The resulting terms are successive quotients of a finite filtration on total cohomology.
This constructs the Leray spectral sequence used here and allows dimensions to be checked page by page.

The projection formula for finite-rank locally free coefficients, \ref{sh:projection-formula}, applies to these resolutions locally
on each trivializing open set for $L$. Hence
$R^q\pi_*(\pi^*L^k\otimes E)=L^k\otimes R^q\pi_*E$.
Every second-page term outside the row $q=0$ is therefore $O(k^e)$.
For $E_r^{m,0}$, outgoing differentials with $r\geq2$ have negative second target index and vanish.
Incoming differentials come from $(m-r,r-1)$, of total degree $m-1$ and positive second index, so their ranks are $O(k^e)$.
There are only finitely many such $r$.
The total cohomology is finite-dimensional by hypothesis, so $E_\infty^{m,0}$ is finite-dimensional. Only the finitely many incoming spaces just bounded are removed from $E_2^{m,0}$; hence this term is finite-dimensional as well, and
$\dim E_\infty^{m,0}=h^m(X,L^k\otimes\pi_*E)+O(k^e)$.
All other final-page terms in total degree $m$ are also $O(k^e)$.
Adding the dimensions of the finite filtration proves the assertion.
\end{proof}

\begin{theorem}[Support dimension controls cohomology growth]
\label{morse:support-growth}
Let $X$ be a compact complex manifold, $L$ any holomorphic line bundle, and $\mathcal G$ a coherent sheaf with support of dimension at most $d$.
With the two stated foundational inputs, for every fixed $q\geq0$,
\[
 h^q(X,L^k\otimes\mathcal G)=O(k^d).
\]
The constant may depend on $X,L,\mathcal G,q$.
\end{theorem}
\begin{proof}
Induct on $d$. For $d=0$, the support is finite, and the sheaf is a finite direct sum of finite-dimensional sheaves supported at points.
It is flabby, so its positive-degree cohomology vanishes. Tensoring with a line bundle does not change its vector-space dimension at any point.
Corollary~\ref{sh:zero-dimensional-coherent} gives the result.

Suppose the assertion is known for support dimension at most $d-1$.
First handle nonreduced support. Set
$\mathcal I=\sqrt{\operatorname{Ann}\mathcal G}$, the coherent ideal sheaf of the reduced support $Z$. Here $\operatorname{Ann}\mathcal G$ consists of local holomorphic functions annihilating every local section, and the radical consists of functions some power of which lies in this annihilator.
Its coherence follows from the analytic Nullstellensatz, Proposition~\ref{alg:analytic-nullstellensatz}.
In each stalk, $\mathcal I_x$ has finitely many generators $a_1,\ldots,a_s$,
each with a power $a_i^{N_i}\in\operatorname{Ann}\mathcal G_x$.
Every monomial of degree at least $1+\sum_i(N_i-1)$ contains some $a_i^{N_i}$,
so $\mathcal I_x^{N(x)}\mathcal G_x=0$.
This product sheaf is coherent, and a zero stalk vanishes on a neighborhood. Compactness supplies finitely many such neighborhoods; taking the maximum exponent gives
$\mathcal I^N\mathcal G=0$.

Filter $\mathcal G$ by $\mathcal I^j\mathcal G$.
The finitely many successive quotients $\mathcal G_j=\mathcal I^j\mathcal G/\mathcal I^{j+1}\mathcal G$
are annihilated by $\mathcal I$ and hence are coherent sheaves on the reduced analytic space $Z$.
Proposition~\ref{sh:long-exact}, applied to the successive short exact sequences, gives
\[
 h^q(X,L^k\otimes\mathcal G)
 \leq\sum_{j=0}^{N-1}h^q(X,L^k\otimes\mathcal G_j).
\]
It therefore suffices to treat each $\mathcal G_j$ on reduced support.
This filtration is essential: for example, the kernel of $\cO_{2D}\to\cO_D$ is still supported on the entire divisor $D$.

Let $Z_1,\ldots,Z_s$ be the $d$-dimensional irreducible components of $Z$.
There are finitely many because irreducible components of analytic sets are locally finite and $Z$ is compact.
On disjoint smooth dense open subsets of these components, the constant-rank condition on a matrix presentation makes
$\mathcal G_j$ locally free of rank $r_i$.
Resolve and make the sheaf locally free on each component, obtaining
$\pi_i:Y_i\to Z_i\hookrightarrow X$ and vector bundles $E_i$.
Set $Y=\coprod_iY_i$, $E|_{Y_i}=E_i$, and $\pi|_{Y_i}=\pi_i$.
The natural map
\[
 \mathcal G_j\longrightarrow\pi_*E
\]
is an isomorphism on these dense open subsets.
Its kernel, cokernel, and all $R^a\pi_*E$ for $a>0$ are coherent
and supported in dimension at most $d-1$. Higher direct images vanish on the isomorphism locus,
as do the kernel and cokernel; all lower-dimensional irreducible components already lie in the part of dimension at most $d-1$.

The induction hypothesis bounds the cohomology of these error sheaves by $O(k^{d-1})$ in each degree.
Split the natural map into two short exact sequences using its kernel, image, and cokernel. Then
\[
 h^q(X,L^k\otimes\mathcal G_j)
 =h^q(X,L^k\otimes\pi_*E)+O(k^{d-1}).
\]
Explicitly, the error is bounded by the dimensions of kernel cohomology in degrees $q,q+1$ and cokernel cohomology in degrees $q-1,q$.
The spaces on the smooth compact $Y_i$ are finite-dimensional by Theorem~\ref{hod:hodge}; thus Lemma~\ref{morse:leray-comparison} applies and gives
\[
 h^q(X,L^k\otimes\mathcal G_j)
 =\sum_i h^q(Y_i,\pi_i^*L^k\otimes E_i)+O(k^{d-1}).
\]
Each $Y_i$ is a smooth compact complex manifold of dimension $d$.
The vector-bundle weak Morse inequality \eqref{heat:classical-weak} bounds each term on the right by $O(k^d)$.
The curvature integral is fixed and finite; even when it is zero, $o(k^d)$ is still $O(k^d)$.
Sum over the finitely many filtration quotients to complete the induction.
\end{proof}

\begin{theorem}[Holomorphic Morse inequalities for arbitrary coherent sheaves]
\label{morse:coherent-reduction}
Let $X$ be a compact complex manifold of dimension $n$, let the holomorphic line bundle $L$ have a smooth Hermitian metric,
and let $\mathcal F$ be a coherent sheaf of generic rank $r$.
With the two stated foundational inputs, all weak, strong, and adjacent-degree inequalities of Theorem~\ref{morse:main}
hold for $L^k\otimes\mathcal F$, and the asymptotic Riemann--Roch integral is still over all of $X$.
\end{theorem}
\begin{proof}
Choose a proper modification $\pi:Y\to X$ from a smooth manifold such that
$E=\pi^*\mathcal F/\text{torsion}$ is a rank-$r$ vector bundle.
On a dense open subset of $X$, the sheaf $\mathcal F$ is locally free and $\pi$ is an isomorphism.
The kernel and cokernel of $\mathcal F\to\pi_*E$, and all higher direct images,
are therefore supported on analytic sets of dimension at most $n-1$.
Theorem~\ref{morse:support-growth} and Lemma~\ref{morse:leray-comparison}, whose finite-dimensionality hypothesis again follows from Theorem~\ref{hod:hodge} on $Y$, give
\begin{equation}
 h^q(X,L^k\otimes\mathcal F)
 =h^q(Y,\pi^*L^k\otimes E)+O(k^{n-1}).
 \label{morse:coherent-dimension-comparison}
\end{equation}
This comparison includes the torsion of the original sheaf; it cannot be discarded before the comparison is proved.

Equip $\pi^*L$ with the pullback metric, whose first Chern form is $\pi^*\alpha$.
On the isomorphism locus, its positive and negative indices agree with those of $\alpha$, and holomorphic change of variables preserves the integrals.
The exceptional set and its image are lower-dimensional analytic sets of zero Riemannian volume.
This follows on their smooth strata from Lipschitz coverings of lower-dimensional parameter domains, while the singular strata decrease further in dimension.
The top-degree curvature integrals over corresponding index sets on $Y$ and $X$ are therefore equal.

Apply Theorem~\ref{morse:main} to $Y,E,\pi^*L$,
then substitute \eqref{morse:coherent-dimension-comparison} degree by degree.
The finitely many $O(k^{n-1})$ errors are absorbed into $o(k^n)$,
giving the three inequalities and the leading Euler-characteristic term.
If $r=0$, the original support has dimension below $n$; the support-growth theorem directly gives $o(k^n)$ in every degree, again agreeing with the formulas.
\end{proof}

\begin{corollary}[Leading section growth for ample line bundles with coherent coefficients]
\label{morse:ample-coherent-growth}
If $L$ is ample and $\mathcal F$ has generic rank $r$, then
\[
 h^0(X,L^k\otimes\mathcal F)
 =\frac{rk^n}{n!}\int_Xc_1(L)^n+o(k^n).
\]
\end{corollary}
\begin{proof}
An ample line bundle has a positive-curvature metric, constructed in Corollary~\ref{morse:semipositive-growth}.
The coherent weak Morse inequalities give $o(k^n)$ for every degree $q\geq1$.
Absorb these finitely many higher-degree terms into the coherent asymptotic Riemann--Roch remainder to obtain the result.
\end{proof}

The noncompact $q$-convex case requires separate analysis: one must show that higher-degree weighted $L^2$ harmonic forms compute the original cohomology,
that the confining potential at infinity gives compact resolvent, and that cutoffs concentrate low-energy spectral spaces in a fixed relatively compact region.
The compact heat-kernel theorem alone contains none of these steps.
For the precise dependencies and noncompact version, see paragraph~\ref{dep:noncompact-morse}, \cite[Theorem 0.1]{Bouche1989},
and \cite[Section 6]{DemaillyMorse1991}.
For versions with boundary, also consult \cite{Berman2005,Berman2008}.

\section{A basic counterexample for effectiveness questions}
\label{heat:effective-warning}

\begin{example}[Very ampleness does not make the critical adjoint bundle ample]
\label{heat:fujita-counterexample}
Take $n\geq1$, $X=\mathbf P^n$, and $L=\mathcal O_{\mathbf P^n}(1)$.
The bundle $L$ is very ample, but
\[
K_X\otimes L^{n+1}\simeq\mathcal O_X
\]
is not ample.
The formula $K_{\mathbf P^n}=\mathcal O(-n-1)$ can be checked directly in coordinates:
on $Z_0\ne0$, use $z_j=Z_j/Z_0$; on $Z_i\ne0$, use
$w_0=1/z_i$ and $w_j=z_j/z_i$.
The coordinate-change Jacobian is a sign times $z_i^{-(n+1)}$.
The transitions for top holomorphic forms are therefore precisely those of $\mathcal O(-n-1)$.
Every positive tensor power of the trivial bundle is still trivial.
By Theorem~\ref{bas:compact-holomorphic}, holomorphic functions on a compact connected complex manifold are constant,
so its global sections do not separate distinct points and it cannot be very ample.
By the definition of ampleness, the trivial bundle is not ample either.
Thus semipositivity, global generation, and ampleness of $K_X+(n+1)L$ are distinct assertions.
\end{example}

Fujita's conjecture asks for bounds at specified tensor powers.
An assertion valid for all sufficiently large $m$ does not determine what happens at $m=n+1$.
For an effective linear global-generation bound, see Han's preprint~\cite{Han2026}; this concerns global generation and does not identify it with ampleness at the critical coefficient.

\paragraph{Exercises.}
First use finite complexes of length two or three to check
the index of the remainder in \eqref{heat:boundary-remainder}.
Next check \eqref{heat:trace-integral} on the flat torus.
Finally, let $b\to0$, $t\to0$, and $k\to\infty$ separately to check
the three distinct limits in \eqref{heat:magnetic-formula}.
The three calculations illustrate, respectively, the algebra of complexes, the diagonal trace formula, and the role of scaling.



\chapter{Potentials, Tauberian Methods, and Spectral Distributions}
\label{heat:extensions-chapter}

A heat kernel records both diffusion and the effect of a potential.
The Feynman--Kac formula expresses this through an average over paths, and a remote-well example shows how a distant potential can influence the kernel.
We then turn from kernels to spectra: Tauberian methods recover eigenvalue counts from heat traces.
We write $d=\dim_{\R}M$, $d_g$ for Riemannian distance, and $dV_g$ (or $dV$) for volume.
For a random path, $\mathbb E_x$ denotes expectation for paths started at $x$ and $\mathbb E_t^{x,y}$ expectation for the bridge from $x$ to $y$ over time $t$; the measures defining these expectations are constructed next.
A Dirac measure $\delta_a$ has mass one at $a$. For a lower-bounded self-adjoint operator $P$ with discrete spectrum $(\lambda_j)$, its counting measure is $\mu_P=\sum_j\delta_{\lambda_j}$ and its counting function is $N(\lambda)=\#\{j:\lambda_j\leq\lambda\}$, including multiplicity.
At times for which the sum is finite, the identity
\[
 \operatorname{tr}(e^{-tP})=\sum_j e^{-t\lambda_j}
 =\int e^{-t\lambda}\,d\mu_P(\lambda)
\]
explains the link: the heat trace is the Laplace transform of the counting measure.
The later sections ask how much of that measure can be recovered from its transform.

\section{Constructing path measures from the scalar heat kernel}

\begin{lemma}[Continuous heat-diffusion paths]
\label{heat:brownian-construction}
Let $M$ be a connected compact Riemannian manifold, and $p_t$ the kernel of $e^{-t\Delta_g}$.
There is a stochastic process $X_s$ starting at $x$, with continuous sample paths and finite-dimensional distributions
\[
 p_{s_1}(x,x_1)p_{s_2-s_1}(x_1,x_2)\cdots
 p_{s_m-s_{m-1}}(x_{m-1},x_m)\,dV(x_1)\cdots dV(x_m).
\]
Here $0<s_1<\cdots<s_m$ are observation times, and the process is a family of random variables with these joint laws. Its generator is $-\Delta_g$; in Euclidean coordinates each real component has variance $2s$, whereas the convention with generator $\tfrac12\sum_i\partial_i^2$ has variance $s$.
For $t>0$ and $x,y\in M$, there is also a continuous Brownian bridge defined by multiplying these densities by
$p_{t-s_m}(x_m,y)/p_t(x,y)$.
\end{lemma}
\begin{proof}
The scalar kernel is nonnegative and $e^{-t\Delta_g}1=1$, so every displayed density has total mass $1$.
The semigroup identity ensures consistency: integrating out an intermediate time leaves the corresponding lower-dimensional density.
First work on $[0,1]$ and use only the dyadic times $\mathbb D=\{j2^{-m}:m\in\N_0,\ 0\leq j\leq2^m\}$. Arbitrary finite time intervals follow by rescaling time in the estimates.
On the compact countable product $M^{\mathbb D}$, the consistent distributions define a positive functional on continuous cylinder functions.
Every continuous function is uniformly approximated by functions depending on finitely many coordinates: equip the product with metric
$\sum 2^{-j}\min(1,d_g(x_j,y_j))$, truncate the small tail, and use uniform continuity.
The functional therefore extends to all continuous functions. The measure-representation construction in the proof of Lemma~\ref{heat:borel-calculus} uses only compactness and continuous cutoffs and applies to this compact metric product as well, giving a probability measure.
This constructs the required countable-time extension without assuming a path measure in advance.

Theorem~\ref{heat:scalar-gaussian} and the radial integration argument of Theorem~\ref{heat:stieltjes} give, for every $p>0$,
\[
 \mathbb E\,d_g(X_s,X_{s+h})^p\leq C_p h^{p/2},\qquad0<h\leq1.
\]
Fix $0<\gamma<1/2$, and choose $p$ with $p(1/2-\gamma)>1$.
The level-$m$ dyadic grid has $2^m$ adjacent increments.
Markov's inequality and the union bound make the probability that any increment exceeds $2^{-m\gamma}$ at most
\[
 C_p2^m2^{-mp/2}2^{mp\gamma}
   =C_p2^{-m[p(1/2-\gamma)-1]}.
\]
This series converges. Indeed, the probability of a bad grid at some level $m\geq M$ is bounded by the corresponding series tail, which tends to zero as $M\to\infty$. Thus outside a null event all sufficiently fine grids satisfy the bound on every adjacent increment.
Dyadic expansions decompose the interval between two dyadic times into at most two grid intervals at each level.
Summing distances by a geometric series gives a uniform continuity estimate $C|s-u|^\gamma$.
The paths therefore extend uniquely and continuously to all real times in the interval.
Continuity of the heat kernel preserves the finite-dimensional distributions after extension.
The generator follows from
$\lim_{h\downarrow0}h^{-1}(e^{-h\Delta_g}f-f)=-\Delta_gf$.

We must also show that the bridge denominator is nonzero.
For small $t$ and $d_g(x,y)\leq c\sqrt t$, Proposition~\ref{heat:parallel-transport-parametrix} in the scalar case makes the local parametrix strictly positive and its error smaller than the leading term.
Connectedness gives a finite chain of small balls from $x$ to $y$. Repeated use of the semigroup identity,
integrating a positive kernel over positive-volume balls at each step, gives $p_t(x,y)>0$.
Bridge consistency again follows from the semigroup identity.
On $[0,t-\delta]$, its law relative to the process starting at $x$ has density
$p_{t-s}(X_s,y)/p_t(x,y)$, uniformly bounded there. The same increment argument gives interior continuity.
Near the terminal time, reverse time using symmetry of the kernel and repeat the argument starting at $y$.
The two continuous versions agree on the dense set of common times and patch to a continuous bridge from $x$ to $y$.
\end{proof}

\begin{theorem}[The scalar Feynman--Kac formula]
\label{heat:feynman-kac}
Let $W\in C^\infty(M,\R)$ on a connected compact manifold $M$. Then
\begin{align}
 (e^{-t(\Delta_g-W)}f)(x)
 &=\mathbb E_x\left[\exp\left(\int_0^tW(X_s)\,ds\right)f(X_t)\right],
 \label{heat:fk-semigroup}\\
 H_W(t;x,y)
 &=p_t(x,y)\mathbb E_t^{x,y}
       \left[\exp\left(\int_0^tW(X_s)\,ds\right)\right].
 \label{heat:fk-bridge}
\end{align}
\end{theorem}
\begin{proof}
Write $T_0(t)=e^{-t\Delta_g}$.
For the bounded multiplication operator $W$, iterate Proposition~\ref{heat:bounded-difference}:
\[
 e^{-t(\Delta_g-W)}
 =\sum_{m=0}^\infty\int_{0<s_1<\cdots<s_m<t}
 T_0(s_1)W T_0(s_2-s_1)W\cdots W T_0(t-s_m)\,ds_1\cdots ds_m.
\]
The $m=0$ term means $T_0(t)$.
The $m$th term has operator norm at most $\|W\|_\infty^m t^m/m!$, so the series converges absolutely.
It satisfies the bounded-perturbation Volterra equation. Iterating the norm bound for the difference of two solutions gives
$C\|W\|^m t^m/m!$ for every $m$; letting $m\to\infty$ proves uniqueness.
The sum is therefore the stated heat semigroup.

On the other hand, path continuity makes $\int_0^tW(X_s)ds$ an ordinary Riemann integral, of absolute value at most $t\|W\|_\infty$.
Expand the exponential and partition the $m$-dimensional cube into its $m!$ time orderings:
\[
 \exp\left(\int_0^tW(X_s)ds\right)
 =\sum_{m=0}^\infty\int_{0<s_1<\cdots<s_m<t}\prod_{j=1}^mW(X_{s_j})\,ds.
\]
The integrable bound $e^{t\|W\|_\infty}\|f\|_\infty$ permits interchanging expectation, integration, and summation.
Using the finite-dimensional distributions from Lemma~\ref{heat:brownian-construction}, each term is exactly the corresponding Dyson term applied to $f$.
This proves \eqref{heat:fk-semigroup}, first for bounded $f$ and then by density in $L^2$.

Repeat with the bridge distributions. Their final scalar-kernel factor supplies the last factor of the Dyson kernel,
with denominator $p_t(x,y)$.
Multiplying back gives the kernel series for $H_W$ term by term, proving \eqref{heat:fk-bridge}.
This is a kernel identity; it does not treat the probability-zero event $X_t=y$ as elementary conditional probability.
\end{proof}

\section{Remote potential wells and localization}

\begin{proposition}[A remote-well counterexample]
\label{heat:remote-well}
On $S^1=\R/\Z$, take $\Delta=-\partial_x^2$, $0<r<1/10$,
and a smooth $V\geq0$ with $V=0$ on $B_r=(-r,r)$ and
$V\geq v_1>0$ on some $U\Subset S^1\setminus\overline B_r$.
Let $P_k=k^{-1}\Delta-V$, with global heat kernel $H_k$.
There are no constants $C,\epsilon_1>0$, independent of $k$, for which
\[
 H_k(t;0,0)\leq C(k/t)^{1/2}e^{t\sup_{B_r}V}
 \quad\text{for every }0<t\leq k\epsilon_1.
\]
An exponential localization estimate for the difference between the global kernel and the Dirichlet kernel of $B_r$, using only the same local potential bound, also fails.
\end{proposition}
\begin{proof}
A spectral argument makes the influence of the distant well explicit.
Choose a smooth compactly supported function $f$ in $U$ with $L^2$ norm $1$.
Theorem~\ref{heat:minmax} bounds the first eigenvalue by
\[
 -\|V\|_\infty\leq\lambda_{1,k}
 \leq k^{-1}\|f'\|_2^2-v_1\leq-v_1/2
\]
for all sufficiently large $k$.
A first eigenfunction can be chosen real and nonnegative: replacing $u$ by $|u|$ does not increase its Rayleigh quotient,
and minimization with compact embedding gives a nonnegative minimizer. Its Euler equation and one-dimensional elliptic regularity make it smooth.
If it vanished at a point, its first derivative there would also vanish, and uniqueness for the second-order ODE would force it to vanish identically.
The normalized first eigenfunction $\phi_k$ is therefore strictly positive.

Its equation is
$\phi_k''=-k(V+\lambda_{1,k})\phi_k$.
Let $x_*$ be a minimum point, with $\phi_k'(x_*)=0$, and put
$Y=(\phi_k,k^{-1/2}\phi_k')$.
The coefficient matrix of this first-order system has norm at most $C\sqrt k$. Along a circle arc of length at most $1$,
Gronwall's inequality gives
$|Y(x)|\leq e^{C\sqrt k}\phi_k(x_*)$.
Since $\|\phi_k\|_2=1$ and the circle has length $1$, some point has $\phi_k\geq1$.
Hence $\phi_k(0)\geq\phi_k(x_*)\geq e^{-C\sqrt k}$.
The nonnegative diagonal spectral expansion gives
\[
 H_k(t;0,0)\geq e^{-t\lambda_{1,k}}\phi_k(0)^2
 \geq\exp(v_1t/2-2C\sqrt k).
\]
Fix any $0<\tau\leq\epsilon_1$ and take $t=k\tau$. The left-hand side grows exponentially,
whereas the asserted upper bound is $C\tau^{-1/2}$ because the local potential bound is zero. This is a contradiction.

On $B_r$, the Dirichlet operator is simply $k^{-1}\Delta^D$,
whose kernel at time $k\tau$ equals the fixed kernel of $\Delta^D$ at time $\tau$, independent of $k$.
Their difference therefore also grows exponentially and cannot satisfy a localization bound involving only the local potential.
\end{proof}

\begin{remark}[A semigroup interpretation of the same counterexample]
Alternatively, split time into quarters, a half, and a quarter.
Since $V\geq0$, every Dyson term is nonnegative, so the global kernels for the outer two intervals dominate the free kernels.
Restrict the middle interval to $U$. Dirichlet comparison and $V\geq v_1$ give the lower bound
$e^{v_1k\tau/2}p_U^D(\tau/2;z,w)$.
Integration over $U\times U$ yields a diagonal lower bound $C_\tau e^{v_1k\tau/2}$ with $C_\tau>0$.
This is the contribution of diffusion traveling to the well, remaining there, and returning.
It agrees with the bridge formula and shows why locating both endpoints in a small ball does not confine the entire heat path.
\end{remark}

\begin{proposition}[Localization with a global potential bound]
\label{heat:potential-localization}
Let $P_k=k^{-1}\nabla_k^*\nabla_k-V$ on a fixed compact manifold $M$, with unitary connection and Hermitian potential $V$.
Write $V_0=\sup_M\|V\|$ and $V_r=\sup_{B(x_0,r)}\|V\|$.
There are geometric constants $c,C>0$ such that, for normal coordinate balls and
$0<t\leq\min(k,ckr^2)$,
\[
 \|H_k(t;x_0,x_0)-H^D_{k,B_r}(t;x_0,x_0)\|
 \leq C(k/t)^{d/2}
       \exp\left(-\frac{kr^2}{Ct}+(V_0+V_r)t\right).
\]
In particular, the potential factor may be enlarged to $e^{2V_0t}$.
\end{proposition}
\begin{proof}
First use a smooth approximate identity supported near the center, multiplied by a unit fiber vector, as initial data.
Inside the ball, the difference $w$ of the global and Dirichlet solutions satisfies the same homogeneous heat equation with zero initial data
and boundary values equal to the global solution.
Lemma~\ref{heat:kato} makes $e^{-V_rt}|w|$ a scalar heat subsolution.
The comparison principle proved in Theorem~\ref{heat:scalar-gaussian} bounds it by the supremum of the boundary data over preceding times.
Corollary~\ref{heat:global-gaussian} bounds the boundary kernel data by
\[
 C(k/s)^{d/2}\exp\left(V_0s-\frac{kr^2}{C_1s}\right),\qquad0<s\leq t,
\]
For sufficiently small initial support, use $r/2$ first and absorb the change into $C_1$.
If $t\leq ckr^2$ with $c$ sufficiently small, then
$s^{-d/2}\exp[-kr^2/(C_1s)]$ increases on $(0,t]$,
since its logarithmic derivative is $-d/(2s)+kr^2/(C_1s^2)\geq0$.
The boundary supremum is therefore at most this Gaussian factor at $s=t$, multiplied by $e^{V_0t}$.
Multiply back by $e^{V_rt}$ to obtain the estimate.
Finally, let the initial data tend to a point mass at the center. Positive-time kernel smoothness and the uniform boundary bound justify the limit.
Take the supremum over unit fiber vectors to obtain the kernel-norm bound.
\end{proof}

\section{Kernel formulas for small zeroth-order perturbations}

\begin{proposition}[From operator Duhamel to kernel Duhamel]
\label{heat:kernel-duhamel}
On a fixed compact manifold, let $P=P_0+A$ on the common domain $D(P)=D(P_0)$, where $A$ is a bounded smooth endomorphism and both operators are lower-bounded and self-adjoint.
If their kernels satisfy the Gaussian bounds established here, then for positive time,
\begin{equation}
 H_P(t;x,y)-H_{P_0}(t;x,y)
 =-\int_0^t\int_MH_P(t-s;x,z)A(z)H_{P_0}(s;z,y)\,dV(z)ds.
 \label{heat:kernel-duhamel-formula}
\end{equation}
If $P_k=P_{0,k}+k^{-1}A$ and both kernels, for $0<t\leq T$, have bounds of the form
$C_T(k/t)^{d/2}e^{-kd_g(x,y)^2/(C_Tt)}$, then
\[
 \|H_k(t;x,y)-H_{0,k}(t;x,y)\|
 \leq C_T k^{d/2-1}t^{1-d/2}e^{-kd_g(x,y)^2/(C_Tt)}.
\]
If only a Gaussian bound containing $e^{Mt}$ is known for the full range $0<t\leq k$, the right-hand side must retain $e^{Mt}$ as well.
\end{proposition}
\begin{proof}
A bounded perturbation preserves the operator domain.
Differentiate on that domain the product
$e^{-(t-s)P}e^{-sP_0}$, then integrate to obtain the operator Duhamel identity.
The integrand has operator norm at most $e^{Ct}\|A\|$, so the integral defines a bounded operator,
and density extends the formula to all of $L^2$.
This is Proposition~\ref{heat:bounded-difference} applied to the multiplication operator $A$.

For $0<s<t$, composition of the two smooth kernels gives
\[\int H_P(t-s;x,z)A(z)H_{P_0}(s;z,y)dV(z).\]
Write the triple integral first for smooth initial data; absolute integrability permits Fubini.
Lemma~\ref{heat:manifold-convolution}, on the fixed compact manifold, gives integrability at the time endpoints.
Retain both independent factors
$s^{-d/2}(t-s)^{-d/2}$; spatial convolution combines them into $Ct^{-d/2}$.
The time integral is therefore $Ct^{1-d/2}$, retaining a slightly weaker off-diagonal Gaussian exponent.
This proves \eqref{heat:kernel-duhamel-formula} pointwise.

At the semiclassical scale, the two kernels supply $k^{d/2}$ each, spatial Gaussian integration supplies $k^{-d/2}$,
the perturbation supplies $k^{-1}$, and time integration supplies $t$. Together these give
$k^{d/2-1}t^{1-d/2}$.
Factors $e^{M(t-s)}$ and $e^{Ms}$ multiply to $e^{Mt}$ and cannot be dropped.
Finally, smoothness of the difference comes from the two heat operators themselves;
arbitrarily repeated termwise differentiation of the endpoint integral is neither needed nor justified without additional estimates.
\end{proof}

\section{How Laplace transforms determine counting measures}

Lemma~\ref{heat:bernstein} supplies the polynomial approximation and weak-measure comparison used below.

\begin{theorem}[Continuity of Laplace transforms]
\label{heat:laplace-continuity}
Let $\mu_k,\mu$ be locally finite positive measures on $[0,\infty)$, with finite Laplace integrals for every positive parameter, and suppose
\[
 \int e^{-t\lambda}\,d\mu_k(\lambda)
 \longrightarrow\int e^{-t\lambda}\,d\mu(\lambda)\qquad(t>0).
\]
Here the Laplace transform is the displayed integral, and $\mu(\{a\})=0$ means that $a$ is not an atom. Then for every $a\geq0$ with $\mu(\{a\})=0$,
$\mu_k([0,a])\to\mu([0,a])$.
The Laplace transform determines the measure uniquely in this class.
The same conclusion holds after translation when the common lower support bound is any fixed $-C$.
\end{theorem}
\begin{proof}
Push the finite measures $e^{-\lambda}\mu_k$ to $(0,1]$ under $z=e^{-\lambda}$,
assign zero mass at $z=0$, and call the resulting measures $\nu_k$. Define $\nu$ from $\mu$ similarly.
Their total masses are the Laplace integrals at $t=1$, hence uniformly bounded.
Their $m$th moments are
\[
 \int_0^1z^m\,d\nu_k(z)=\int_0^\infty e^{-(m+1)\lambda}\,d\mu_k(\lambda),
 \qquad m=0,1,2,\ldots.
\]
Lemma~\ref{heat:bernstein} therefore gives $\nu_k\rightharpoonup\nu$.

For $a>0$, one has
$\mu_k([0,a])=\int z^{-1}\mathbf1_{[e^{-a},1]}(z)\,d\nu_k(z)$.
The integrand is bounded, with only one interior discontinuity, at $e^{-a}$.
Approximate it above and below by continuous functions with linear transitions on an interval of length $\epsilon$.
Since $\nu(\{e^{-a}\})=e^{-a}\mu(\{a\})=0$, the difference between their limiting integrals tends to zero.
Weak convergence and squeezing prove the assertion.
For $a=0$ with $\mu(\{0\})=0$, use arbitrarily small positive continuity points $a'$ to bound
$0\leq\mu_k(\{0\})\leq\mu_k([0,a'])$, then let $a'\downarrow0$.

If two measures have the same Laplace transform, their corresponding measures $\nu$ have the same moments.
Polynomial approximation makes their integrals against all continuous functions equal, hence the measures coincide.
Multiplying by $z^{-1}$ on $(0,1]$ and reversing the change of variables recovers $\mu$.
Translating $\lambda$ only multiplies the Laplace transform by a fixed exponential factor, leaving the argument unchanged.
\end{proof}

\begin{theorem}[Karamata's Tauberian theorem]
\label{heat:karamata}
Use $f\sim g$ to mean $f/g\to1$ in the indicated limit. Let $\mu$ be a positive measure on $[0,\infty)$, with $a>0$, $\alpha>0$, and
\[
 \int_0^\infty e^{-t\lambda}\,d\mu(\lambda)\sim a t^{-\alpha}
 \qquad(t\downarrow0).
\]
Then
\[
 \mu([0,R])\sim\frac{a}{\Gamma(\alpha+1)}R^\alpha
 \qquad(R\to\infty).
\]
\end{theorem}
\begin{proof}
Define the scaled measure $\mu_R(B)=R^{-\alpha}\mu(RB)$.
For every fixed $t>0$,
\[
 \int e^{-t\lambda}\,d\mu_R(\lambda)
 =R^{-\alpha}\int e^{-(t/R)s}\,d\mu(s)\longrightarrow a t^{-\alpha}.
\]
Take $\mu_\infty$ to have density
$a\lambda^{\alpha-1}/\Gamma(\alpha)$.
The substitution $u=t\lambda$ gives its Laplace transform as $a t^{-\alpha}$.
It has no atom at $1$, so Theorem~\ref{heat:laplace-continuity} yields
\[
 R^{-\alpha}\mu([0,R])=\mu_R([0,1])
 \longrightarrow\frac{a}{\Gamma(\alpha)}\int_0^1\lambda^{\alpha-1}d\lambda
 =\frac{a}{\Gamma(\alpha+1)}.
\]
The last step uses $\Gamma(\alpha+1)=\alpha\Gamma(\alpha)$, obtained by one integration by parts.
\end{proof}

\begin{corollary}[Weyl's leading term]
\label{heat:weyl}
On a compact Riemannian manifold of dimension $d\geq1$, scalar Laplace eigenvalue counting satisfies
\[
 N(\Lambda)\sim
 \frac{\operatorname{Vol}(M)}{(4\pi)^{d/2}\Gamma(1+d/2)}\Lambda^{d/2}.
\]
The Dirichlet Laplacian on a smooth bounded domain has the same leading term, with the domain's volume.
\end{corollary}
\begin{proof}
The scalar diagonal construction in Proposition~\ref{heat:parallel-transport-parametrix} gives
$p_t(x,x)=(4\pi t)^{-d/2}(1+O(t))$, uniformly.
After integration, the heat trace is $(4\pi t)^{-d/2}\operatorname{Vol}(M)(1+O(t))$.
Apply Theorem~\ref{heat:karamata} to the discrete counting measure $\sum_j\delta_{\lambda_j}$.
For the Dirichlet problem, at a point of boundary distance $\delta(x)$ the difference from the surrounding kernel is bounded by
$Ct^{-d/2}e^{-c\delta(x)^2/t}$.
The smooth boundary-collar volume bound makes the integrated error $O(t^{(1-d)/2})$,
or $O(\sqrt t)$ relative to the leading term. The same Tauberian argument applies.
\end{proof}

\section{Magnetic spectral measures: the case with no free directions}

\begin{proposition}[The counting density of a constant magnetic field]
\label{heat:magnetic-density}
Let the real dimension be $d$, let the nonzero magnetic strengths be $B_1,\ldots,B_s>0$, and set $\nu=d/2-s\geq0$.
For $\nu>0$, $[u]_+^\nu=(\max\{u,0\})^\nu$. For $\nu=0$, use $[u]_+^0=\mathbf1_{\{u\geq0\}}$, where $\mathbf1$ is the indicator function. Also $\N_0=\{0,1,2,\ldots\}$.
Define the right-continuous nondecreasing function
\begin{equation}
 \mathcal N_B(\lambda)=
 \frac{2^{s-d}\pi^{-d/2}B_1\cdots B_s}{\Gamma(\nu+1)}
 \sum_{p\in\N_0^s}
 \left[\lambda-\sum_{j=1}^s(2p_j+1)B_j\right]_+^\nu.
 \label{heat:magnetic-density-formula}
\end{equation}
For $s=0$, the empty product is $1$ and the empty-index sum has one term.
Its Stieltjes measure is defined by $d\mathcal N_B((a,b])=\mathcal N_B(b)-\mathcal N_B(a)$ and satisfies
\[
 \int e^{-t\lambda}\,d\mathcal N_B(\lambda)
 =(4\pi t)^{-d/2}\prod_{j=1}^s\frac{B_jt}{\sinh(B_jt)}.
\]
\end{proposition}
\begin{proof}
The diagonal heat density of each magnetic plane has geometric-series expansion
\[
 \frac{B_j}{4\pi\sinh(B_jt)}
 =\frac{B_j}{2\pi}\sum_{p_j=0}^\infty e^{-t(2p_j+1)B_j}.
\]
It therefore corresponds to the discrete positive measure
$\frac{B_j}{2\pi}\sum_{p_j\geq0}\delta_{(2p_j+1)B_j}$.
Products over different planes correspond to convolutions of these measures.
The remaining $d-2s$ free real directions give $(4\pi t)^{-\nu}$.
For $\nu>0$, this corresponds to density
$(4\pi)^{-\nu}\lambda_+^{\nu-1}/\Gamma(\nu)$,
with cumulative distribution $(4\pi)^{-\nu}\lambda_+^\nu/\Gamma(\nu+1)$.
Convolve by translation at each discrete energy level to obtain \eqref{heat:magnetic-density-formula}.
The constant simplifies to
$(2\pi)^{-s}(4\pi)^{-\nu}=2^{s-d}\pi^{-d/2}$.

For $\nu=0$, there are no free directions; the remaining factor $1$ corresponds to $\delta_0$.
The cumulative distribution is therefore a step function at the Landau levels, rather than an ordinary zeroth power to be differentiated as a function.
With the stated right-continuous convention, the formula remains valid.
All terms are nonnegative, so Tonelli permits interchanging the series and Laplace integral.
\end{proof}

\begin{theorem}[Limiting eigenvalue counts for variable magnetic operators]
\label{heat:magnetic-counting}
Let $P_k$ satisfy Theorem~\ref{heat:scaled-diagonal}, or take its Dirichlet realization on a fixed smooth relatively compact domain.
Let $v_1(x),\ldots,v_r(x)$ be the eigenvalues of $V(x)$ with multiplicity,
and let $B(x)$ list the absolute values of the nonzero magnetic strengths of $\ii F^L_x$.
At every continuity point $\lambda$ of the following distribution function,
\begin{equation}
 \lim_{k\to\infty}k^{-d/2}N_k(\lambda)
 =\int_M\sum_{\ell=1}^r\mathcal N_{B(x)}(\lambda+v_\ell(x))\,dV_g(x).
 \label{heat:magnetic-counting-formula}
\end{equation}
For Dirichlet realizations, integrate over the domain instead.
In particular, the formula holds for all real $\lambda$ except at most countably many.
\end{theorem}
\begin{proof}
Set $\mu_k=k^{-d/2}\sum_j\delta_{\lambda_{j,k}}$.
Nonnegative connection energy and uniformly bounded potentials give a common finite lower support bound for these measures.
Corollary~\ref{heat:semiclassical-trace} (or Corollary~\ref{heat:dirichlet-trace} for a domain) gives the Laplace-transform limit
\[
 (4\pi t)^{-d/2}\int_M\sum_{\ell=1}^r e^{tv_\ell(x)}
                    \prod_j\frac{tB_j(x)}{\sinh(tB_j(x))}\,dV_g(x).
\]
By Proposition~\ref{heat:magnetic-density}, this is the Laplace transform of the positive measure $\mu$ obtained, for each $x,\ell$,
by translating $d\mathcal N_{B(x)}$ left by $v_\ell(x)$ and then integrating.

We verify that this integrated measure is locally finite.
For fixed $A$, positivity gives
$\mathcal N_{B(x)}(A+v_\ell(x))
\leq e^{A+v_\ell(x)}\int e^{-s}d\mathcal N_{B(x)}(s)$.
The last heat density is at most $(4\pi)^{-d/2}$ because $B/\sinh B\leq1$.
The $v_\ell$ are uniformly bounded and $M$ has finite volume, giving a uniform finite total bound on bounded intervals.
When magnetic rank varies, partition into finitely many measurable rank strata.
On each stratum the expression is a countable sum of nonnegative measurable terms formed from continuous ordered eigenvalues, so it is measurable.
Tonelli's theorem therefore applies.

Apply Theorem~\ref{heat:laplace-continuity} at points where $\mu$ has no atom to obtain counting convergence,
with cumulative distribution exactly \eqref{heat:magnetic-counting-formula}.
A locally finite measure has at most countably many positive-mass atoms on each compact interval:
for each $m$, only finitely many atoms have mass at least $1/m$; take their countable union.
Then take the union over integer intervals to prove the last assertion.
\end{proof}

\begin{example}[A matrix potential must be shifted eigenvalue by eigenvalue]
\label{heat:matrix-potential-warning}
With two free real directions, $\mathcal N_0(\lambda)=\lambda_+/(4\pi)$.
For $V=\operatorname{diag}(0,1)$, the correct counting density is
\[
 \frac{\lambda_++(\lambda+1)_+}{4\pi},
\]
with heat density $(1+e^t)/(4\pi t)$.
Replacing the potential by $\tr V=1$ gives only $(\lambda+1)_+/(4\pi)$, or incorrectly twice that if one then multiplies by rank.
Neither is correct.
\end{example}
\begin{proof}
The two eigendirections of the matrix potential are uncoupled, so the operator is a direct sum of two scalar operators.
Counting functions and heat traces add under direct sums, while the constant potential $-1$ shifts the spectrum left by $1$.
This gives both displayed formulas and verifies the need to sum over the individual $v_\ell$.
\end{proof}

\begin{remark}[The continuity-point condition is essential]
Take $\mu_k=\delta_{1/k}$ and $\mu=\delta_0$.
Their Laplace transforms converge for every positive parameter, but
$\mu_k(({-\infty},0])=0$ and $\mu(({-\infty},0])=1$.
Heat-trace convergence therefore does not by itself imply convergence of kernel dimensions at threshold zero.
The Morse argument obtains its inequalities from positive-eigenvalue subcomplexes before passing to a limit of heat traces.
\end{remark}


\chapter{Heat Regularization and Nadel Vanishing}
\label{nad:chapter}

Local convolutions do not automatically patch into a global line bundle metric: second-order errors from coordinate changes may concentrate near the singular set.
This chapter controls these errors with local heat kernels and proves that they disappear in the weighted Hilbert spaces.
The result is Nadel vanishing for general singular line bundle metrics.
The multiplier ideal measures the local integrability imposed by the metric; it disappears from an application only when that integrability condition holds for every holomorphic germ.
As in \eqref{bun:line-curvature}, $\Theta_h=F_h$ and $\ii\Theta_h$ denotes the real curvature.
A quasi-psh function is locally the sum of a plurisubharmonic function and a smooth real function.
For $h=e^{-\varphi}$ in a holomorphic frame, $\mathcal I(h)=\mathcal I(\varphi)$ is the multiplier ideal of Definition~\ref{pos:multiplier-definition}; its frame independence is Proposition~\ref{pos:multiplier-invariance}.

\section{From Positive Currents to Controlled Smooth Metrics}

\begin{proposition}[A global quasi-psh potential]\label{nad:global-potential}
Let $T$ be a closed real $(1,1)$ current on a complex manifold $X$, with a smooth lower bound locally.
There exist a smooth closed real $(1,1)$ form $\alpha$ and a global quasi-psh function $\varphi$
such that $T=\alpha+\ii\partial\dbar\varphi$.
On a compact K\"ahler manifold, if the de Rham class of $T$ is $[T]=2\pi c_1(L)$, then $T$ is the real curvature $\ii\Theta_h$ of
a singular metric with locally integrable weights.
\end{proposition}
\begin{proof}
Locally add $C\ii\partial\dbar|z|^2$ to make $T$ positive.
Theorem \ref{cur:local-potential}, followed by subtraction of $C|z|^2$, gives real local potentials $\varphi_j$ for $T$.
Each is a psh function minus a smooth function, hence quasi-psh.
On overlaps, $\partial\dbar(\varphi_j-\varphi_k)=0$,
so the difference is smooth and pluriharmonic.
Choose a subordinate smooth partition of unity $\{\eta_j\}$ by Theorem~\ref{bas:partition} and set $\varphi=\sum_j\eta_j\varphi_j$.
On $U_k$, write
$\varphi=\varphi_k+\sum_j\eta_j(\varphi_j-\varphi_k)$;
the final sum is smooth, so $\varphi$ is locally quasi-psh.
The form $\alpha=T-\ii\partial\dbar\varphi$ is smooth on each chart and is closed.

In the compact K\"ahler case, choose a smooth metric $h_0$ on $L$.
Then $T-\ii\Theta_{h_0}$ is a $d$-exact real $(1,1)$ current.
The current version of the $\partial\dbar$ lemma, \ref{hod:ddbar-lemma}, gives a real distribution $u$
whose complex Hessian $\ii\partial\dbar u$ is this difference.
Compare $u$ locally with the $L^1$ potential already obtained for $T$, minus a smooth potential for $\ii\Theta_{h_0}$. Their difference has zero complex Hessian,
hence is smooth. Thus $u\in L^1_{\rm loc}$ and has a quasi-psh representative.
The metric $h=h_0e^{-u}$ has curvature $T$.
\end{proof}

\begin{example}[Limits of smooth regularization]\label{nad:regularization-counterexamples}
First, a real $(1,1)$ current representing $2\pi c_1(L)$ need not come from a metric with locally integrable weights.
On a compact K\"ahler manifold, fix a line bundle $L$ with smooth metric $h_0$, take the scalar Dirac distribution $\delta_x$ (evaluation of a test density at $x$), and set
$T=\ii\Theta_{h_0}+\ii\partial\dbar\delta_x$.
This current is closed, real, and cohomologous to the original curvature. If $T=\ii\Theta_{h_0e^{-u}}$ with $u\in L^1$,
then $\partial\dbar(u-\delta_x)=0$. Taking the trace and applying elliptic regularity would make
$u-\delta_x$ smooth, forcing the Dirac distribution to be an $L^1$ function, a contradiction.
The local lower bound in the proposition is what guarantees locally integrable potentials.

Second, even a global psh function cannot always be approximated from above by smooth functions
whose losses relative to an arbitrary specified curvature lower bound tend uniformly to zero on fixed compact sets.
Take $X=\operatorname{Tot}(\cO_{\mathbf P^1}(-1))$, the total space of the tautological line bundle, with projection $\pi:X\to\mathbf P^1$ and zero section $E$.
Its natural squared fiber norm is locally
$\tau=|w|^2(1+|z|^2)$, and defines a global function.
Put $\varphi=\log\tau+\tau$. Theorem~\ref{cur:poincare-lelong} gives
\[
 \ii\partial\dbar\varphi=2\pi[E]+\theta,
 \qquad \theta=2\pi\,\pi^*\omega_{\rm FS}+\ii\partial\dbar\tau.
\]
Here $\omega_{\rm FS}$ has the normalization of Proposition~\ref{bun:fs}, so the factor $2\pi$ is required. The form $\theta$ is a smooth K\"ahler form: the first term is positive horizontally,
and the second is the pullback of the Euclidean squared-norm Hessian and is positive in the fiber direction along the zero section.
If smooth functions $\varphi_j$ on a relatively compact neighborhood of $E$ satisfied
$\ii\partial\dbar\varphi_j\ge(1-\delta)\theta$,$\delta<1$,
restriction to the compact curve $E$ and integration would give
$0\ge(1-\delta)\int_E\theta>0$, a contradiction.
The next lemma gives a different type of control: the curvature errors stay bounded and tend to zero almost everywhere, but need not tend to zero uniformly near the singular set.
\end{example}

\begin{lemma}[Heat regularization on relatively compact sets]\label{nad:heat-regularization}
Let $(X,\omega)$ be K\"ahler, and let $h=h_0e^{-\varphi}$ be a line bundle metric,
with $h_0$ smooth, $\varphi$ quasi-psh, and $T=\ii\Theta_h\ge\theta$,
where $\theta$ is a continuous semipositive $(1,1)$ form.
Here $D\Subset X$ means that $D$ is open with compact closure contained in $X$. For every such $D$, there are weights $\varphi_j$ smooth near $\overline D$ such that
\begin{enumerate}
\item $\varphi_j\ge\varphi$ and $\varphi_j\to\varphi$ almost everywhere;
\item the $\varphi_j$ have a common upper bound on $\overline D$;
\item $\ii\Theta_{h_0e^{-\varphi_j}}\ge\theta-\lambda_j\omega$, where
$0\le\lambda_j\le C_D$ and $\lambda_j\to0$ almost everywhere.
\end{enumerate}
The errors need not tend to zero on the singular locus.
\end{lemma}
\begin{proof}
Choose $\overline D\subset U_0\Subset U_1\Subset U_2\Subset X$.
Choose a relatively compact smooth real manifold with boundary containing $\overline U_2$, and double it across its boundary
to obtain a compact real manifold $M$. Extend the metric using collar coordinates and a partition of unity,
so that the Riemannian metric on $M$ agrees exactly with the original K\"ahler metric on $U_2$.
No complex structure is needed outside this neighborhood.
Choose a compactly supported cutoff $\chi$ equal to one on a neighborhood of $\overline U_1$ within $U_2$, and regard $\chi\varphi$ as an $L^1$ function on $M$,
denoted $\widetilde\varphi$. It has a finite upper bound on $M$.
Let $\Delta_\ell=dd^*+d^*d$ on real degree-$\ell$ forms, with $\Delta_0=d^*d$ on functions. Let $P_t=e^{-t\Delta_0}$ be the scalar heat semigroup and $H_t=e^{-t\Delta_2}$ the Hodge heat semigroup on two-forms. Write $d_g$ for the Riemannian distance.
We use the kernel expansion in Proposition~\ref{heat:parallel-transport-parametrix} and
the Gaussian derivative bounds in Proposition~\ref{heat:fixed-derivative-gaussian}.

\emph{1. The localized commutator produces only exponentially small errors.}
Let $J$ be the complex structure and define $J^*\eta(v)=\eta(Jv)$ on one-forms. On $U_2$, put $A=-J^*/2$ and extend it to a smooth endomorphism of the one-form bundle on $M$.
The second-order operator $Q=dAd$ equals $\ii\partial\dbar=\pi dd^c$ on $U_2$, with the normalization in \eqref{pos:dc-convention}.
Proposition~\ref{hod:kahler-laplacians} and the commutation identities in Proposition~\ref{pos:kahler-identities} give $\Delta_2Q-Q\Delta_0=0$ there.
The support of the nonzero commutator coefficients is therefore at positive distance from $\overline U_1$.
Differentiate $H_{t-s}QP_s$ first on smooth functions and integrate in $s$, as in the proof of Theorem~\ref{heat:duhamel}; this yields the following identity before the distributional limit is taken:
\[
 QP_t\widetilde\varphi-H_tQ\widetilde\varphi
 =\int_0^tH_{t-s}(\Delta_2Q-Q\Delta_0)P_s\widetilde\varphi\,ds.
\]
For $x\in\overline U_1$, integrate by parts in the integration variable, transferring all commutator derivatives to
$H_{t-s}(x,\cdot)$ and the smooth coefficients.
Their support is at distance at least $a>0$ from $x$, so the integrand is bounded in absolute value by
$C(t-s)^{-N}e^{-a^2/[C(t-s)]}\|P_s\widetilde\varphi\|_{L^1}$.
The scalar heat semigroup preserves positivity and contracts $L^1$, so the last norm is uniformly bounded.
Integration gives an $O(e^{-c/t})$ error, with the same conclusion for every fixed derivative order.
For a rigorous distributional justification, first prove the formula for smooth $L^1$ approximations of $\widetilde\varphi$.
The integration-by-parts bound uses only the $L^1$ norm, uniformly in the approximation, so one may then pass to the $L^1$ limit.
Only the metric on the K\"ahler neighborhood enters the local commutation identity.

\emph{2. The regularized weights dominate the singular weight.}
Put $\alpha=\ii\Theta_{h_0}$.
The inequality $\alpha+\ii\partial\dbar\varphi\ge0$ implies, on $U_2$,
$\Delta_0\varphi\le C$ distributionally for the nonnegative-sign Laplacian.
Split $\Delta_0\widetilde\varphi$ into a part near $\overline U_1$ and a part supported away from it.
The first has a fixed upper bound; pairing the distant part with the heat kernel gives another $O(e^{-c/t})$ term.
The identity
$P_t\widetilde\varphi-\widetilde\varphi=-\int_0^tP_s\Delta_0\widetilde\varphi\,ds$
therefore gives $P_t\widetilde\varphi\ge\varphi-Ct-O(e^{-c/t})$ on $\overline U_1$.
First obtain this distributionally, then pass to the canonical quasi-psh representative by shrinking ball averages to get a pointwise inequality.
Thus we may take
$\varphi_t=P_t\widetilde\varphi+Ct+C'e^{-c'/t}\ge\varphi$.
The approximate-identity property of the heat kernel gives almost-everywhere convergence. Positivity of the scalar heat semigroup and
the upper bound for $\widetilde\varphi$ give a uniform upper bound.

\emph{3. Positive-current mass controls the curvature error.}
Extend $\alpha$ smoothly to $M$. Step 1 and the smooth heat expansion give
\[
 \alpha+\ii\partial\dbar\varphi_t
   =H_t(\widetilde\alpha+Q\widetilde\varphi)+O(t)
       \quad\text{on }\overline D.
\]
The input current equals $T$ on $U_1$. Its exterior part is at positive distance from $D$ and contributes only an exponentially small error.
The leading heat-kernel term is a scalar Gaussian kernel times Levi--Civita parallel transport, with a smooth positive scalar density factor.
Short geodesics stay in the K\"ahler neighborhood, where parallel transport preserves the complex structure and the positive $(1,1)$ cone.
The leading average of $T\ge\theta$ is therefore at least $\theta-o(1)$, uniformly on
$\overline D$, by continuity of $\theta$.
The kernel remainder is controlled by
\[
 C t\int t^{-n}\exp\!\left(-\frac{d_g(x,y)^2}{Ct}\right)d\mu(y)
\]
where $\mu$ is the trace measure of $T$ on a fixed neighborhood plus a smooth volume measure.
Proposition \ref{cur:mass-growth} gives
$\mu(B(x,r))\le C r^{2n-2}$.
Split the integral into $d_g(x,y)\le\sqrt t$ and the annuli
$2^k\sqrt t<d_g(x,y)\le2^{k+1}\sqrt t$. This gives the uniform bound
$C\sum_{k\ge0}e^{-c4^k}2^{k(2n-2)}<\infty$.
On the other hand, at volume-almost-every differentiation point of the measure $\mu$,
$\mu(B(x,r))\le C_x r^{2n}$ for sufficiently small $r$.
The same annular estimate bounds the nearby contribution by $O_x(t)$; the contribution beyond a fixed distance decays exponentially.
Thus this error tends to zero almost everywhere.
Adding the preceding uniform continuous $o(1)$ term gives the required continuous nonnegative $\lambda_t$.
Any sequence $t_j\downarrow0$ gives all three assertions.
\end{proof}

\section{Weighted Projections and Curvature Errors}

The curvature error introduces an extra term into the $L^2$ estimate.
To solve the equation, we must show that this term disappears after projection onto closed forms.
The projections themselves vary with the metric, so we compare them in one fixed Hilbert space.

\begin{lemma}[Closed-subspace projections with varying metrics]\label{nad:projection-error}
Let $h_j\le h$ be smooth positive line bundle metrics converging almost everywhere to $h$, with a uniform positive lower bound relative to a fixed smooth metric.
Fix the base metric, and let $P_j$ be the orthogonal projection in $L^2(h_j)$ onto the distributionally $\dbar$-closed forms.
If $0\le\lambda_j\le C$, $\lambda_j\to0$ almost everywhere, and
$\|v_j\|_{h_j}\le C'$, then
every local weak limit of $a_j=P_j(\sqrt{\lambda_j}v_j)$ is zero.
\end{lemma}
\begin{proof}
Rescale the fixed metric to arrange $h_j\ge h_*$.
Put $H_j=(h_j/h_*)^{1/2}$ and $H=(h/h_*)^{1/2}$.
Then $1\le H_j\le H$ and $H_j\to H$ almost everywhere.
Contractivity of the projections and $\lambda_j\le C$ give
$\|H_ja_j\|_{h_*}\le\sqrt C C'$.
Pass to a weakly convergent subsequence $H_ja_j\rightharpoonup b$.
Since $H_j^{-1}\le1$ and converges almost everywhere, dominated convergence on each fixed test vector shows that
$a_j\rightharpoonup a=H^{-1}b$ in the fixed $L^2$ space.
Hence $\dbar a=0$ and $\|a\|_h=\|b\|_{h_*}<\infty$.

Because $h_j\le h$, $a$ lies in the range of every $P_j$.
Moreover, $H_ja\to Ha$ strongly in the fixed $L^2$ space, so
\[
 \|a\|_h^2
   =\lim_j(a,a_j)_{h_j}
   =\lim_j(a,\sqrt{\lambda_j}v_j)_{h_j}.
\]
The first equality combines the strong convergence $H_ja\to Ha$ and the weak convergence $H_ja_j\to b=Ha$;
the second uses self-adjointness of the orthogonal projection.
The absolute value of the last pairing is at most
$C'(\int\lambda_j|a|_h^2)^{1/2}$, which tends to zero by dominated convergence.
Thus $\|a\|_h=0$, and $a=0$.
The projections depend on $j$; convergence of $\sqrt{\lambda_j}$ alone
would not justify omitting this argument.
\end{proof}

\begin{theorem}[Global $L^2$ solvability for general singular line bundle metrics]\label{nad:global-l2}
Let $(X,\omega)$ be a K\"ahler manifold admitting a smooth psh exhaustion,
and let $h$ be a singular metric on a line bundle $L$ with $\ii\Theta_h\ge\theta$,
where $\theta$ is a continuous semipositive $(1,1)$ form.
Let $n=\dim_{\C}X$, let $\Lambda_\omega$ be the adjoint of wedge multiplication by $\omega$, and let $A=[\theta\wedge,\Lambda_\omega]$ be the induced operator on $(n,q)$ forms.
The inverse energy uses the extended-real convention of Proposition~\ref{pos:inverse-curvature}: it is infinite on a nonzero component in $\ker A$.
If $q\ge1$, $\dbar f=0$, and
$I=\int_X\langle A^{-1}f,f\rangle_h dV_\omega<\infty$,
then there is $u$ satisfying $\dbar u=f$ and $\int_X|u|_h^2dV_\omega\le I$.
\end{theorem}
\begin{proof}
First work on a relatively compact regular sublevel set $D\Subset X$ of the psh exhaustion. Proposition~\ref{l2:complete-metric} gives a complete K\"ahler metric $\widehat\omega_D$ there; fix $\delta>0$ and put
$\omega_\delta=\omega+\delta\widehat\omega_D$.
Lemma \ref{nad:heat-regularization} gives $h_j\le h$ with curvature at least
$\theta-\lambda_j\omega$, where the $\lambda_j$ are uniformly bounded and tend to zero almost everywhere.
In degree $(n,q)$, since $\omega\le\omega_\delta$, the negative curvature-error operator is bounded below by
$-q\lambda_j\operatorname{Id}$.

Let $T_j=\dbar:L^2_{n,q-1}(h_j)\to L^2_{n,q}(h_j)$ and
$S_j=\dbar:L^2_{n,q}(h_j)\to L^2_{n,q+1}(h_j)$,
let $P_j$ project onto $\ker S_j$.
Put $A_\delta=[\theta\wedge,\Lambda_{\omega_\delta}]$; all norms in this step use $\omega_\delta$. For $v\in\operatorname{Dom}T_j^*\cap\ker S_j$, Theorem~\ref{pos:bkn} and Lemma~\ref{l2:complete-core} give
\[
 \int\langle A_\delta v,v\rangle_{h_j}
 \le\|T_j^*v\|_{h_j}^2+q\|\sqrt{\lambda_j}v\|_{h_j}^2.
\]
Weighted Cauchy--Schwarz then yields
\[
 |(f,v)_{h_j}|^2\le I_{j,\delta}
       \bigl(\|T_j^*v\|_{h_j}^2+q\|\sqrt{\lambda_j}v\|_{h_j}^2\bigr),
\]
where $I_{j,\delta}$ is the inverse-curvature energy computed with $\theta,\omega_\delta,h_j$.
For every $v\in\operatorname{Dom}T_j^*$,
$T_j^*P_jv=T_j^*v$, $(f,P_jv)=(f,v)$. Indeed, $\operatorname{im}T_j\subset\ker S_j$ implies $(\ker S_j)^\perp\subset\ker T_j^*$; this is the projection step in Lemma~\ref{pos:hilbert-existence}.
Apply the estimate to $P_jv$, then use Hahn--Banach and Riesz on the range of
$v\mapsto(T_j^*v,\sqrt{q\lambda_j}P_jv)$
to obtain $u_j,b_j$ satisfying
\[
 f=T_ju_j+P_j(\sqrt{q\lambda_j}b_j),\qquad
 \|u_j\|_{h_j}^2+\|b_j\|_{h_j}^2\le I_{j,\delta}.
\]
Testing against every $v\in\operatorname{Dom}T_j^*$ proves $u_j\in\operatorname{Dom}T_j$
and the stated Hilbert-space identity, by the definition of the adjoint of the closed operator.

Lemma~\ref{l2:metric-monotonicity}, together with $h_j\leq h$, gives $I_{j,\delta}\le I$.
Lemma \ref{nad:projection-error} makes the error term tend weakly to zero locally.
Take a weak subsequence of $H_ju_j$, as in that lemma, to obtain $u$
with $\dbar u=f$ and $\|u\|_{h,\omega_\delta}^2\le\liminf I_{j,\delta}$.
Since $h_j\le h$ and converges almost everywhere, dominated convergence gives
$I_{j,\delta}\to I_{h,\delta}$.
Let $\delta\downarrow0$ and use the compact-set comparison and weak-limit argument of Theorem \ref{l2:incomplete-metric}
to obtain a solution on $D$ in the original metric $\omega$, with constant one.

Finally exhaust $X$ by increasing relatively compact psh sublevel sets.
Every solution has squared norm bounded by the global $I$, so diagonal weak limits on compact sets
give a global solution. Lower semicontinuity and monotone exhaustion recover the global estimate.
Finite energy also puts the datum in the required weighted $L^2$ spaces on each $D$,
because the continuous form $\theta$ has a finite upper bound on $\overline D$.
\end{proof}

\begin{remark}[A smooth bundle multiplied by a scalar singular weight]
The proof also applies to a smooth Hermitian vector bundle multiplied by a scalar quasi-psh weight,
provided the curvature lower bound is understood in the Nakano sense, or in the $m$-positive sense appropriate to the degree $(n,q)$.
During regularization, parallel transport acts on both the tangent and smooth bundle factors. Freezing the smooth bundle curvature gives an $o(1)$ error,
while the scalar Hessian remainder is still controlled by the same trace measure.
These positivity cones are preserved by unitary parallel transport and nonnegative averaging, so the projection argument applies unchanged.
This still does not assert a curvature theorem for arbitrary measurable vector bundle metrics.
\end{remark}

\section{Nadel Vanishing and Jets}

\begin{theorem}[Nadel vanishing]\label{nad:vanishing}
Here weakly pseudoconvex means that $X$ admits a smooth plurisubharmonic exhaustion. Let $(X,\omega)$ be weakly pseudoconvex K\"ahler, and suppose a singular metric $h$ on $L$ satisfies
$\ii\Theta_h\ge\epsilon(x)\omega$, with $\epsilon$ a positive continuous function.
Then, for every $q\ge1$,
\[
 H^q\bigl(X,\cO(K_X\otimes L)\otimes\cI(h)\bigr)=0.
\]
\end{theorem}
\begin{proof}
By Theorem \ref{l2:multiplier-coherent} and Proposition \ref{l2:local-resolution},
this cohomology is computed by global sections of the locally weighted $L^2$ Dolbeault resolution.
Take a closed $(n,q)$ form $f$ in that complex.
Let $\rho$ be a nonnegative smooth psh exhaustion. Choose a smooth convex increasing $\chi$ such that
\[
 \int_X |f|_h^2(1+\epsilon^{-1})e^{-\chi(\rho)}dV_\omega<\infty.
\]
Choose it successively on compact annuli: the original integral over $j\le\rho\le j+1$ is finite,
so it suffices to make $\chi(j)$ exceed the logarithm of one plus that integral, plus $j\log2$. Choose successively increasing slopes of a polygonal function,
then smooth it convexly while retaining these lower bounds.

The new metric $h'=he^{-\chi(\rho)}$ still has curvature at least $\epsilon\omega$.
The new and old weights differ by a bounded function on every compact set, so $\cI(h')=\cI(h)$,
and the local weighted $L^2$ resolution is unchanged.
Theorem \ref{nad:global-l2} with $\theta=\epsilon\omega$ gives
$\dbar u=f$ and
$\int|u|_{h'}^2\le\int|f|_{h'}^2/(q\epsilon)<\infty$.
Thus $u$ is locally square integrable for the original weight and belongs to the original resolution as a preimage of $f$.
Every closed class in positive degree vanishes, as required.
\end{proof}

\begin{corollary}[Finite jets at isolated points]\label{nad:jets}
Under the preceding theorem's assumptions, let $x_1,\ldots,x_N$ be
isolated points of $V(\cI(h))=\{x:\cI(h)_x\ne\cO_{X,x}\}$. Put $\mathcal M=\cO(K_X\otimes L)$, the sheaf of holomorphic sections. Then
\[
 H^0(X,\mathcal M)\longrightarrow
 \bigoplus_i\mathcal M_{x_i}\otimes_{\cO_{X,x_i}}\bigl(\cO_{X,x_i}/\cI(h)_{x_i}\bigr)
\]
is surjective. These are sheaf stalks, as in Definition~\ref{sh:def-stalk}. For a local metric weight $\varphi$, use the Lelong normalization in Definition~\ref{int:lelong-definition}. If $\nu(\varphi,x_i)\ge n+s_i$ with $s_i\in\N_0$, arbitrary jets of order $s_i$ can in particular be prescribed at these points.
\end{corollary}
\begin{proof}
Tensor the sequence $0\to\cI(h)\to\cO_X\to\cO_X/\cI(h)\to0$
with $K_X\otimes L$. Proposition~\ref{sh:long-exact} and Theorem~\ref{nad:vanishing} in degree one give surjectivity onto global sections of the quotient sheaf.
A local section at each isolated support point can be extended by zero on the rest of the support:
choose a small neighborhood containing no other support point, then glue to zero outside it.
Consequently, global sections of the quotient surject onto the displayed finite direct sum.
Finally use Proposition \ref{int:jet-order}, which gives
$\cI(h)_{x_i}\subset\mathfrak m_{x_i}^{s_i+1}$, and pass to the jet quotient.
\end{proof}



\chapter{Singular-Weight Estimates on General Complete Manifolds}
\label{csg:chapter}

We now extend the singular-weight estimate to manifolds that admit a complete K\"ahler metric but need not admit a plurisubharmonic exhaustion.
The curvature lower bound will be its measurable absolutely continuous part.
Following Demailly's method \cite[\S5]{Demailly1982}, we replace exhaustion by plurisubharmonic sublevel sets with compactly supported cutoffs and descending induction on the form degree.
The regularization from Chapter~\ref{nad:chapter} supplies the local estimates needed in that induction.
Throughout, $n=\dim_{\C}X$ and $r=\operatorname{rank}E$. The metric $\omega$ determines $d_g$, $dV_\omega$, and the adjoint $\Lambda_\omega$ of wedge multiplication by $\omega$.
The symbol $\Theta_h=F_h$ is the Chern curvature; capital $\Theta$ without a subscript below denotes a real Hermitian curvature tensor, including the factor $\ii$.
We use the rank-tested $m$-semipositivity of Definition~\ref{l2:m-positive}, the curvature action of Proposition~\ref{pos:nakano-operator}, and the generalized inverse energy of Proposition~\ref{pos:inverse-curvature}.
For a measure-coefficient current, the subscript $\mathrm{ac}$ means its coefficientwise absolutely continuous part with respect to $dV_\omega$ (Radon--Nikodym decomposition, Lemma~\ref{cur:radon-nikodym}).

\section{Measurable Curvature and Compactly Supported Cutoffs}

\begin{lemma}[Wedge multiplication and inverse curvature energy]\label{csg:wedge-energy}
At a point of an $n$-dimensional Hermitian complex manifold with a rank-$r$ Hermitian vector bundle, let $\Theta$ be an $m$-semipositive Hermitian curvature tensor,
where $m\ge\min(n-q+1,r)$, and let $A_q=[\Theta,\Lambda]$ act in bidegree $(n,q)$, with $1\le q\le n$.
For every $(0,1)$ covector $a$ and every $(n,q)$ form $f$,
\[
 \langle A_{q+1}^{-1}(a\wedge f),a\wedge f\rangle
 \le |a|^2\langle A_q^{-1}f,f\rangle.
\]
The inequality is interpreted in the extended real sense when its right-hand side is infinite.
\end{lemma}
\begin{proof}
The cases $q=n$ and $a=0$ are immediate. Otherwise, homogeneity allows us to assume $|a|=1$ and choose unitary coordinates with $a=d\bar z_1$.
Let $a^*$ denote the contraction adjoint to exterior multiplication by $a$.
Write $c_{j\bar k\alpha\bar\beta}$ for the curvature coefficients in unitary tangent and bundle frames, and $v_{jK,\alpha}$ for the antisymmetric coefficients of $v$; repeated indices in $jK$ give zero. For an $(n,q+1)$ form $v$, Proposition~\ref{pos:nakano-operator} expresses the curvature quadratic form as
\[
 Q_{q+1}(v)=\sum_{|K|=q}\sum_{j,k,\alpha,\beta}
 c_{j\bar k\alpha\bar\beta}v_{jK,\alpha}
                 \overline{v_{kK,\beta}}.
\]
The terms indexed by $K$ containing $1$ sum to precisely $Q_q(a^*v)$:
substituting $K=\{1\}\cup H$ into the alternating coefficients produces the same pair of signs on both sides, and hence the same product.
All remaining terms are nonnegative, since each coefficient matrix has at most $n-q$ nonzero rows
and therefore rank at most the permitted $m$. Thus
$Q_q(a^*v)\le Q_{q+1}(v)$.
Proposition~\ref{pos:inverse-curvature} now gives
\begin{align*}
 \langle A_{q+1}^{-1}(a\wedge f),a\wedge f\rangle
 &=\sup_v\{2\operatorname{Re}\langle f,a^*v\rangle-Q_{q+1}(v)\}\\
 &\le\sup_w\{2\operatorname{Re}\langle f,w\rangle-Q_q(w)\}
 =\langle A_q^{-1}f,f\rangle.
\end{align*}
The same variational formula handles zero eigenvalues, so no invertibility assumption is needed.
\end{proof}

\begin{lemma}[Global approximation of general quasi-psh weights]\label{csg:global-regularization}
Let $(X,\omega)$ be K\"ahler, let $h_0$ be a $C^2$ Hermitian metric on $E$,
and let $\varphi$ be quasi-psh, with $h=h_0e^{-\varphi}$. Assume that
\[
 \Theta=\ii\Theta_{h_0}
       +(\ii\partial\dbar\varphi)_{\rm ac}\otimes\operatorname{Id}_E
\]
is $m$-semipositive almost everywhere. There exist $C^2$ metrics $h_j=h_0e^{-\varphi_j}$,
continuous $m$-semipositive tensors $\Theta_j$, and continuous functions $\lambda_j\ge0$ such that
\[
 h_*\le h_j\le h,\quad h_j\to h\ \hbox{almost everywhere},\qquad
 \ii\Theta_{h_j}\ge_m\Theta_j-\lambda_j\omega\otimes\operatorname{Id}_E.
\]
Here $h_*$ is a fixed positive $C^2$ metric,
$\Theta_j\to\Theta$ and $\lambda_j\to0$ almost everywhere,
and the functions $\lambda_j$ are uniformly bounded on each compact set.
The notation $\ge_m$ means that the difference is nonnegative on matrices of rank at most $m$.
\end{lemma}
\begin{proof}
We first construct approximations on an arbitrary relatively compact set $D$.
The local heat-flow construction in Lemma \ref{nad:heat-regularization} gives
$\varphi_t\ge\varphi$, together with a parallel-transport main term and an error term for its complex Hessian.
We retain the average of the absolutely continuous curvature, instead of fixing a continuous form as the lower bound.
More precisely, use parallel transport along the short geodesic from $y$ to $x$,
with the K\"ahler connection on the tangent bundle and the $h_0$-unitary connection on $E$, and denote the resulting action by $\tau_{y,x}$.
For a nonnegative cutoff $\zeta$ supported at short distances and the scalar heat kernel $p_t$, set
\[
 \Theta_t(x)=\int\zeta(d_g(x,y))p_t(x,y)
                   \tau_{y,x}\Theta(y)\,dV_y.
\]
The integration takes place in a fixed neighborhood of $D$. Errors arising from the extension outside that neighborhood are
$O(e^{-c/t})$; see \ref{heat:fixed-derivative-gaussian} and
the local commutator argument in Lemma \ref{nad:heat-regularization}.
Definition~\ref{l2:m-positive} shows directly that the cone of $m$-semipositive tensors is closed under unitary transformations and nonnegative averaging, so $\Theta_t$ is continuous and $m$-semipositive.
Approximate identities converge at Lebesgue points of locally $L^1$ functions,
and consequently $\Theta_t\to\Theta$ almost everywhere.

Theorem~\ref{pos:order-zero} applied after adding a local smooth lower-bound correction makes the Hessian of a quasi-psh weight a current of order zero. Its singular part is a positive scalar $(1,1)$ current: the smooth correction has no singular part, and the positive-current density criterion, Theorem~\ref{cur:density}, applies to its coefficient measures.
Its average in the parallel-transport main term remains semipositive and may therefore be discarded from a curvature lower bound.
The continuous bundle curvature differs from its parallel-transport average by a quantity tending uniformly to zero on compact sets.
The remaining error is bounded by
\begin{equation}
 b_t(x)+Ct\int t^{-n}e^{-d_g(x,y)^2/(Ct)}\,d\mu(y)
 \label{csg:regularization-mass-error}
\end{equation}
where $b_t\to0$ uniformly on compact sets and
$\mu$ is the trace measure of the local positive closed current
$\ii\partial\dbar\varphi+C_D\omega$, plus a smooth volume measure.
Proposition \ref{cur:mass-growth} gives $\mu(B(x,r))\le Cr^{2n-2}$.
Summing over annuli with $r\asymp2^\ell\sqrt t$ shows that \eqref{csg:regularization-mass-error} is locally uniformly bounded.
At almost every point where the Lebesgue differentiation property for $\mu$ holds,
we have $\mu(B(x,r))\le C_xr^{2n}$ for sufficiently small $r$.
The same annular summation shows that the integral term, multiplied by $t$, tends to zero; the Gaussian factor suppresses the more distant annuli.
Thus a continuous nonnegative error $\lambda_t$ can be chosen that is locally uniformly bounded, tends to zero almost everywhere, and gives the required curvature lower bound.
The heat-kernel remainder for positive currents used here was established in Proposition \ref{heat:parallel-transport-parametrix}.
The error can remain nonzero on the singular locus; almost-everywhere convergence is the property used in the subsequent $L^2$ limit.

We now diagonalize the local constructions.
Choose a compact exhaustion $K_1\subset\operatorname{int}K_2\subset\cdots$
and a smooth function $M>\varphi+2$; local upper bounds and a partition of unity provide the latter.
Perform the preceding heat smoothing on a neighborhood of $K_{j+2}$.
Since $\varphi\le M-2$, the scalar heat flow preserves order, and the heat smoothing of the smooth function $M$ converges uniformly to $M$.
By taking $t_j$ sufficiently small, we can therefore arrange throughout a neighborhood of $K_{j+1}$ that
$\varphi\le\varphi_{t_j}\le M-1$.
Using a smooth cutoff equal to one near $K_j$ and supported in $\operatorname{int}K_{j+1}$,
take a convex combination of $\varphi_{t_j}$ and $M-1$. The upper bound applies throughout the transition region,
so this gives a global smooth function $\varphi_j$ with $\varphi\le\varphi_j\le M$.
Consequently $h_*=h_0e^{-M}\le h_j\le h$.

Set $\Theta_j=\Theta_{t_j}$ on $K_j$, multiply by a nonnegative cutoff outside, and extend by zero.
This preserves $m$-semipositivity. The continuous curvature loss in the transition region can be absorbed by increasing $\lambda_j$.
Each fixed compact set is eventually disjoint from these transition regions.
We quantify the comparison between different exterior extensions to ensure that their constants do not invalidate the conclusion as $j$ increases.
For each $K_i$, fix in advance a slightly larger relatively compact neighborhood $U_i$ and a fixed compact double,
and choose a fixed $C_i'$ on $U_i$ such that
$\ii\partial\dbar\varphi+C_i'\omega\ge0$.
Denote the error from this fixed construction by $\lambda_t^{(i)}$.
The local mass calculation above gives $\lambda_t^{(i)}\le C_i$ and convergence to zero almost everywhere.
The construction on the larger domain at stage $j$ and the fixed construction at stage $i$ use the same differential operator and input distribution on $U_i$.
Heat-kernel localization and its derivative bounds
(\ref{heat:wave-localization}, \ref{heat:fixed-derivative-gaussian})
show that the differences between their Hessians and between their main terms $\Theta_t$ on $K_i$ are bounded by
$C_{i,j}e^{-c_{i,j}/t}$, with $C_{i,j}<\infty$ and $c_{i,j}>0$.
The short-distance part uses only the fixed neighborhood $U_i$; all exterior mass and extension constants are included in $C_{i,j}$.

Choose $t_j<1/j$ so that the sum of these errors for the finitely many indices $i\le j$ is less than $1/j$,
and so that the constant added to ensure $\varphi_{t_j}\ge\varphi$ is less than $1/j$.
Finally, define $\lambda_j(x)$ to be the larger of zero and the maximum of the continuous tensor
$\Theta_j-\ii\Theta_{h_j}$ on unit-norm matrices of rank at most $m$.
This set of matrices is compact, so $\lambda_j$ is continuous and gives precisely the required curvature lower bound.
On each fixed $K_i$, for all sufficiently large $j$,
$\lambda_j\le\lambda_{t_j}^{(i)}+1/j$.
Thus the errors have a bound independent of $j$ and converge to zero on a common set of full measure.
The same localization comparison, together with the vanishing added constants, gives $\varphi_j\to\varphi$ and
$\Theta_j\to\Theta$ almost everywhere.
The finitely many earlier values of $j$ are covered by enlarging the local bound.
The construction uses an ordinary compact exhaustion and therefore requires no plurisubharmonic exhaustion.
\end{proof}

\section{Descending Induction for the General Complete \texorpdfstring{$L^2$}{L2} Theorem}

\begin{theorem}[Demailly's singular-weight estimate on a general complete manifold]\label{csg:general-complete}
Suppose $X$ admits a complete K\"ahler metric, and let $\omega$ be any K\"ahler metric.
Let $h=h_0e^{-\varphi}$, where $h_0$ is a $C^2$ Hermitian vector-bundle metric,
$\varphi$ is quasi-psh, and the absolutely continuous curvature $\Theta$ from the preceding lemma is $m$-semipositive,
with $1\le q\le n$ and $m\ge\min(n-q+1,r)$.
If $f$ is a closed $(n,q)$ form that is locally square integrable with respect to $h$, and
\[
 I=\int_X\langle[\Theta,\Lambda_\omega]^{-1}f,f\rangle_h\,dV_\omega<\infty,
\]
then there is a solution of $\dbar u=f$ satisfying
\[
 \int_X|u|_h^2dV_\omega\le I.
\]
Every norm uses the singular metric $h$; weighted and unweighted norms must be distinguished throughout these formulas.
\end{theorem}
\begin{proof}
First suppose the metric $\omega$ being used is complete. We argue by descending induction on $q=n,n-1,\ldots,1$.
Choose compactly supported smooth cutoffs $0\le\chi_\nu\le1$ such that
$\chi_\nu\to1$ uniformly on each compact set and
$|\dbar\chi_\nu|\le1/\nu$.
Lemma \ref{l2:complete-core} supplies such cutoffs; rescaling their radii gives the stated constant.
If $q<n$, the induction hypothesis and Lemma \ref{csg:wedge-energy} give $g_\nu$ satisfying
\[
 \dbar g_\nu=\dbar\chi_\nu\wedge f,\qquad
 \|g_\nu\|_h^2\le I/\nu^2.
\]
The $m$-semipositivity assumption on $\Theta$ also satisfies the weaker rank condition required in the next degree.
If $q=n$, the right-hand side vanishes by bidegree, so we take $g_\nu=0$; this starts the induction.

Fix $\nu$ and use the metrics $h_j$ from Lemma \ref{csg:global-regularization}.
Write $T_j=\dbar$ from $(n,q-1)$ to $(n,q)$ and $S_j=\dbar$ from $(n,q)$ to $(n,q+1)$,
with their maximal distributional domains, and let $P_j$ be the orthogonal projection in $L^2(h_j)$ onto $\ker S_j$. Here $\|\cdot\|_j$ and $(\cdot,\cdot)_j$ mean the $L^2$ norm and inner product for $h_j,\omega$.
Set
\[
 B_{j,\nu}=[\Theta_j,\Lambda_\omega]+\nu^{-1}\operatorname{Id},\qquad
 J_{j,\nu}=\int_{\supp\chi_\nu}
                 \langle B_{j,\nu}^{-1}f,f\rangle_{h_j}\,dV_\omega.
\]
Since $B_{j,\nu}\ge\nu^{-1}$ and $h_j\le h$,
$J_{j,\nu}\le\nu\int_{\supp\chi_\nu}|f|_h^2<\infty$.
This is where the hypothesis of local weighted square integrability is used.

For $v\in\operatorname{Dom}T_j^*\cap\operatorname{Dom}S_j$,
first apply Theorem~\ref{pos:bkn} to $\chi_\nu v$, then expand the product rules for both first-order operators.
The sum of the squared norms of exterior multiplication by $\dbar\chi_\nu$ and its adjoint contraction, applied to $v$, is
$|\dbar\chi_\nu|^2|v|^2$.
Using $|a+b|^2\le(1+1/\nu)|a|^2+(1+\nu)|b|^2$, we obtain
\begin{align}
 \int\langle B_{j,\nu}\chi_\nu v,\chi_\nu v\rangle_{h_j}
 &\le(1+1/\nu)(\|T_j^*v\|_j^2+\|S_jv\|_j^2)\notag\\
 &\quad+q\|\sqrt{\lambda_j}\chi_\nu v\|_j^2
             +3\nu^{-1}\|v\|_j^2.\label{csg:cutoff-estimate}
\end{align}
The coefficient $3/\nu$ combines the derivative error, at most $2/\nu$, with the additional $1/\nu$ in $B_{j,\nu}$.
Lemma~\ref{l2:complete-core} extends the inequality to the stated domains. It requires only $C^1$ coefficients in the first-order operators, so the $C^2$ metrics here suffice; the compactly supported BKN identity is valid at this regularity as well.

For arbitrary $v\in\operatorname{Dom}T_j^*$, write $v=P_jv+(1-P_j)v=v_0+v_1$.
Since $\chi_\nu f-g_\nu$ is closed and $g_\nu\in L^2(h_j)$,
\[
 (\chi_\nu f,v)_j=(\chi_\nu f,v_0)_j+(g_\nu,v_1)_j.
\]
Moreover, $T_j^*v_0=T_j^*v$ and $S_jv_0=0$ by the Hilbert-complex projection argument in Lemma~\ref{pos:hilbert-existence}.
Apply Cauchy--Schwarz with the weight $B_{j,\nu}$ and
\eqref{csg:cutoff-estimate} to the first term, and use
$\|g_\nu\|_j^2\le I/\nu^2$ for the second. Applying the same square inequality to these two scalars yields
\[
 |(\chi_\nu f,v)_j|^2\le
 C_{j,\nu}\bigl(\|T_j^*v\|_j^2
       +q\|\sqrt{\lambda_j}\chi_\nu P_jv\|_j^2\bigr)
       +\frac{D_{j,\nu}}\nu\|v\|_j^2,
\]
where we may take
$C_{j,\nu}=(1+1/\nu)^2J_{j,\nu}$ and
$D_{j,\nu}=6J_{j,\nu}+2I$.
Hahn--Banach and Lemma~\ref{heat:riesz}, as in the proof of Lemma~\ref{pos:hilbert-existence}, applied to the map
\[
 v\longmapsto
 \bigl(\sqrt{C_{j,\nu}}T_j^*v,
       \sqrt{qC_{j,\nu}}\sqrt{\lambda_j}\chi_\nu P_jv,
       \sqrt{D_{j,\nu}/\nu}\,v\bigr),
\]
therefore give $u_j,b_j,w_j$ such that
\[
 \chi_\nu f=T_ju_j+\sqrt q\,P_j(\sqrt{\lambda_j}\chi_\nu b_j)+w_j,
 \quad
 \|u_j\|_j^2+\|b_j\|_j^2\le C_{j,\nu},\quad
 \|w_j\|_j^2\le D_{j,\nu}/\nu.
\]
These identities are obtained by testing against the entire adjoint domain and hence also imply
$u_j\in\operatorname{Dom}T_j$.

Keep $\nu$ fixed and let $j\to\infty$.
The functions $\lambda_j\chi_\nu^2$ are uniformly bounded on all of $X$ and tend to zero almost everywhere, since their supports lie in a fixed compact set.
Lemma \ref{nad:projection-error} therefore makes the projection error tend weakly to zero locally.
That lemma also applies to the present vector bundle: $h_j/h_*$ is a positive scalar multiple of the identity,
so its use of $H_j$ and dominated convergence is unchanged.
For the inverse energy, $B_{j,\nu}^{-1}\le\nu$,
and domination by $\nu|f|_h^2\mathbf1_{\supp\chi_\nu}$ gives
\[
 J_{j,\nu}\longrightarrow
 J_\nu=\int_{\supp\chi_\nu}
  \langle([\Theta,\Lambda]+\nu^{-1})^{-1}f,f\rangle_h\le I.
\]
This passage uses the fixed regularizing term $\nu^{-1}$ to obtain a dominating function before taking any limit of inverse matrices.
Multiply the solutions by the metric square roots $H_j=(h_j/h_*)^{1/2}$, take a weakly convergent Hilbert-space subsequence,
and recover the untransformed limits as in Lemma \ref{nad:projection-error}. We obtain
\[
 \chi_\nu f=\dbar u_\nu+w_\nu,\qquad
 \|u_\nu\|_h^2\le(1+1/\nu)^2I,
 \quad\|w_\nu\|_h^2\le8I/\nu.
\]
Finally let $\nu\to\infty$. The sequence $u_\nu$ has a globally weakly convergent subsequence,
while $w_\nu\to0$ strongly and $\chi_\nu f\to f$ locally in $L^2$.
The limit therefore satisfies $\dbar u=f$, and weak lower semicontinuity gives $\|u\|_h^2\le I$.
This completes the induction through all degrees from $q=n$ down to $q=1$.

If the original $\omega$ is incomplete, choose a complete metric $\widehat\omega$ and set
$\omega_\delta=\omega+\delta\widehat\omega$.
The pointwise linear algebra of Lemma \ref{l2:metric-monotonicity} applies equally to measurable $m$-semipositive curvature,
so the inverse energy is at most the original $I$.
Apply the complete case just proved, and then use on each compact set the comparison
$\omega_\delta\le(1+C_K\delta)\omega$.
The weak-limit argument of Theorem \ref{l2:incomplete-metric} eliminates $\delta$ and proves the general statement.
\end{proof}

\begin{corollary}[The estimate under a uniform positive curvature bound]\label{csg:uniform-positive}
Under the assumptions of the preceding theorem, suppose the total curvature satisfies the Nakano lower bound
$\ii\Theta_h\ge\epsilon\omega\otimes\operatorname{Id}_E$, where $\epsilon>0$.
Then every closed weighted-square-integrable $(n,q)$ form has a solution satisfying
\[
 \|u\|_h^2\le\frac1{q\epsilon}\|f\|_h^2.
\]
In particular, this applies to singular line bundles satisfying the displayed positive lower bound.
\end{corollary}
\begin{proof}
Passing to the absolutely continuous part preserves the lower bound. Proposition~\ref{pos:curvature-eigenvalues} gives, in bidegree $(n,q)$,
$[\epsilon\omega\wedge,\Lambda]=q\epsilon\operatorname{Id}$.
Proposition~\ref{pos:inverse-curvature} therefore bounds the inverse energy by $(q\epsilon)^{-1}|f|_h^2$,
and the preceding theorem applies.
\end{proof}

\begin{corollary}[Scalar eigenvalue estimate for a line bundle]\label{csg:scalar-eigenvalues}
Under the assumptions of Theorem \ref{csg:general-complete}, let $E=L$ be a line bundle whose absolutely continuous curvature has eigenvalues
$0\le\lambda_1\le\cdots\le\lambda_n$ relative to $\omega$. Then
\[
 \|u\|_h^2\le\int_X\frac{|f|_h^2}{\lambda_1+\cdots+\lambda_q}\,dV_\omega
\]
provided the right-hand side is finite and $f$ is closed and locally weighted square integrable. At a zero denominator, the quotient is zero when $f=0$ and is infinite otherwise.
\end{corollary}
\begin{proof}
By Proposition~\ref{pos:curvature-eigenvalues}, in a unitary eigenbasis for the curvature the eigenvalues of the curvature operator in bidegree $(n,q)$ are sums of $q$ of the
$\lambda_j$, and the smallest is $\lambda_1+\cdots+\lambda_q$.
Compare the inverse operators componentwise and apply Theorem \ref{csg:general-complete}.
\end{proof}

\begin{remark}[Scalar singular weights and regularity]
The curvature of a general vector bundle is a tensor acting on $T^{1,0}X\otimes E$.
It cannot be described by only $n$ scalar eigenvalues; the preceding $m$-positivity condition and the corresponding curvature operator are required.
Moreover, an arbitrary measurable vector-bundle metric need not have a well-defined distributional curvature.
The singular metrics treated in this chapter are $C^2$ metrics multiplied by scalar quasi-psh weights, a class covering every singular-metric application in this book.

With a singular weight, the theorem initially produces a locally square-integrable distributional solution.
Smoothness of $f$ alone does not imply smoothness of the minimal-norm solution, because the weighted adjoint may have singular coefficients. If $q=1$, then $u$ has bidegree $(n,0)$;
for smooth $f$, local $\dbar$-elliptic regularity using a smooth reference metric still makes $u$ smooth.
For general $q$, smoothness of the minimal solution requires the full metric and weight entering its adjoint to be smooth as well.
\end{remark}


\part{Semester II: Positivity, Growth, and Applications}


\chapter{Ample Vector Bundles: Projectivization and Direct Images}\label{amp:chapter}
Projectivization turns a vector bundle into a line bundle on a larger manifold.
This lets us bring the positivity and embedding theorems of the first semester to bear on higher-rank bundles.
Fiberwise holomorphic sections then connect the two spaces: they give symmetric powers of the original bundle and the direct-image metrics used below.
The distinctions between Griffiths positivity, Nakano positivity, and ampleness remain essential throughout.

Unless a local statement says otherwise, $X$ is a compact complex manifold of dimension $n$, and $E$ has positive rank $r$.
For a holomorphic submersion $\pi:Y\to X$, put $T_{Y/X}=\ker(d\pi)$, $\Omega^j_{Y/X}=\Lambda^jT_{Y/X}^*$, and $K_{Y/X}=\det T_{Y/X}^*$.
A bundle is \emph{globally generated} if its global holomorphic sections span every fiber; for a coherent sheaf the evaluation map must be onto every stalk.
Ample, very ample, and positive line bundles have the meanings fixed in Chapter~\ref{emb:chapter}.
\section{Projective bundles and the relative Euler sequence}
Let $Y=\mathbb P(E^*)$ be the bundle of one-dimensional subspaces of $E^*$, with projection $\pi:Y\to X$ and tautological subbundle $\mathcal O_Y(-1)\subset\pi^*E^*$.
Set $\mathcal O_Y(1)=\mathcal O_Y(-1)^*$. With this convention,
$\pi_*\mathcal O_Y(m)=\operatorname{Sym}^mE$ for $m\ge0$, by Theorem~\ref{sh:projective-direct-image}.
The bundle $E$ is \emph{ample} if $\mathcal O_Y(1)$ is ample.
For rank one, each projective fiber is a point, so $Y=X$ and $\mathcal O_Y(1)=E$.
Thus this definition agrees with line-bundle ampleness.

\begin{proposition}[The Euler sequence and relative canonical bundle]\label{amp:euler}
If $\operatorname{rank}E=r$, then
\[
 0\to\mathcal O_Y\to\pi^*E^*\otimes\mathcal O_Y(1)
       \to T_{Y/X}\to0,
 \qquad K_{Y/X}=\mathcal O_Y(-r)\otimes\pi^*\det E.
\]
\end{proposition}
\begin{proof}
At $[\ell]\in\mathbb P(E_x^*)$, an infinitesimal displacement of the line is a linear map $\ell\to E_x^*/\ell$.
The natural map from $\Hom(\ell,E_x^*)$ to this space has kernel $\Hom(\ell,\ell)\simeq\C$.
These are the fibers of the stated sequence. In homogeneous local coordinates the maps are holomorphic, so the fiberwise sequences glue to a short exact sequence of bundles.
Taking determinants gives
$\det T_{Y/X}=\pi^*(\det E^*)\otimes\mathcal O_Y(r)$.
Dualizing gives the canonical-bundle formula; in particular, $K_{\mathbb P^{r-1}}=\mathcal O(-r)$.
\end{proof}

\begin{lemma}[A strictly plurisubharmonic fiber norm]\label{amp:dual-psh}
If $f>0$ is smooth and $\log f$ is strictly plurisubharmonic, then $\rho(z,\xi)=|\xi|^2f(z)$ is strictly plurisubharmonic where $\xi\ne0$.
In particular, the squared metric norm on the dual of a positive line bundle is strictly plurisubharmonic off the zero section, and its small disk bundles have strictly pseudoconvex boundary.
\end{lemma}
\begin{proof}
The vertical entry of the Levi matrix is $f$, its horizontal block is $|\xi|^2(f_{i\bar j})$, and its mixed entries are $\bar\xi f_i$ and their conjugates.
Completing the square, or taking the Schur complement of the vertical block, gives
\[
 \mathcal L_\rho(v,\eta)
 =f\left|\eta+\xi\,\partial\log f(v)\right|^2
       +|\xi|^2f\,\mathcal L_{\log f}(v).
\]
If $v\ne0$, the second term is positive; if $v=0$ and $\eta\ne0$, the first is positive.
When $|e|_L^2=e^{-\varphi}$, the dual frame has squared norm $e^{\varphi}$, so take $f=e^\varphi$. The dual metric, rather than $e^{-\varphi}$ again, is essential for the sign.
The level set $\rho=\varepsilon$ misses the zero section and has nonzero radial derivative. Its Levi form is positive on complex tangential vectors, proving strict pseudoconvexity.
\end{proof}

\begin{proposition}[Griffiths positivity implies ampleness]\label{amp:griffiths}
A strictly Griffiths-positive Hermitian metric on $E$ induces a strictly positive metric on $\mathcal O_{\mathbb P(E^*)}(1)$. Consequently $E$ is ample.
\end{proposition}
\begin{proof}
Choose a normal holomorphic frame at a base point, with metric matrix $H=I$ and $dH=0$, by Lemma~\ref{pos:normal-frames}. Griffiths positivity is Definition~\ref{pos:nakano-definition}.
In homogeneous fiber coordinates $\xi$, the weight of the dual tautological metric is $\log(\xi^{\mathsf T}H^{-1}\bar\xi)$.
Here $\xi$ is the column of holomorphic dual coordinates: the dual metric matrix is $(H^{-1})^{\mathsf T}$ under our conjugate-linear-first convention.
Its horizontal Levi term is the Griffiths curvature of $E$ evaluated in the direction determined by $\xi$, divided by $|\xi|^2$. The dual-connection sign thus gives a positive horizontal term for $\mathcal O(1)$.
The vertical term is
\[
 \frac{|\xi|^2|\eta|^2-|\langle\xi,\eta\rangle|^2}{|\xi|^4},
\]
which is positive after quotienting out the homogeneous radial direction: it is the Fubini--Study form.
The mixed terms vanish at the chosen point because $dH=0$.
Thus every nonzero tangent vector has positive curvature. The calculation can be made at every base point, giving a global positive metric. Theorem~\ref{emb:kodaira} proves ampleness.
\end{proof}

\section{The direct-image curvature needed for projective bundles}
\begin{theorem}[Positive direct images for a fixed local fiber]\label{amp:direct-image-curvature}
Suppose $\pi:Y\to X$ is locally holomorphically a product $U\times Z\to U$, where $Z$ is compact K\"ahler.
Let $A$ have strictly positive curvature on $Y$ and be holomorphically the pullback of a fixed fiber line bundle on each such product chart.
Assume the dimension of $V_x=H^0(Y_x,K_{Y_x}\otimes A|_{Y_x})$ is constant.
The product identifications give a holomorphic vector bundle $V=\pi_*(K_{Y/X}\otimes A)$ with fibers $V_x$. Its fiber-integral metric has strictly Nakano-positive curvature, in the sense of Definition~\ref{pos:nakano-definition}.
This includes $Z=\mathbb P^{r-1}$ and $A=\mathcal O_Y(r+m)$ used below.
\end{theorem}
\begin{proof}
Work on one small base chart with coordinates $t=(t_1,\ldots,t_n)$, put $d=\dim_\C Z$, and write the metric locally as $e^{-\varphi(t,z)}$.
Latin derivative indices refer to $t$ and Greek indices to fiber coordinates $z$. Changes of fiber line frame are independent of $t$, by the fixed-fiber hypothesis, so the base derivatives below are global on each fiber.
For smooth $A$-valued $(d,0)$-forms, define
\[
 \langle u,v\rangle_t=\int_Z\ii^{d^2}v\wedge\overline u\,e^{-\varphi}.
\]
This is conjugate-linear in $u$ and linear in $v$. Let $\mathcal H_t$ be the resulting $L^2$ completion, and let $P_t$ be its orthogonal projection onto $V_t$. Write $\Theta_{i\bar j}$ for the base components of the Chern curvature of $V$.
In integrals, $u_i\overline{u_j}e^{-\varphi}$ denotes $\ii^{d^2}u_i\wedge\overline{u_j}e^{-\varphi}$ in a line frame. This is the density for $\langle u_j,u_i\rangle_t$ and is positive when $i=j$.
For a holomorphic family $u$, the $(1,0)$ connection is
$D_i u=P_t(\partial_{t_i}u-\varphi_i u)$.
Indeed, differentiate the fiber integral defining the inner product and project onto $V_t$. The use of top-degree forms means that no separate varying volume density is needed.

At a fixed base point choose a holomorphic frame of $V$ whose Chern connection matrix vanishes there.
Differentiate once more in $\bar t$, using $P_t^2=P_t$. For $u_1,\ldots,u_n\in V_t$, the subspace-curvature formula is
\begin{equation}\label{amp:curvature-formula}
 \sum_{i,j}\langle u_j,\Theta_{i\bar j}u_i\rangle
 =\int_Z\sum_{i,j}\varphi_{i\bar j}u_i\overline{u_j}e^{-\varphi}
       -\left\|(I-P_t)\sum_i\varphi_i u_i\right\|^2.
\end{equation}
The sign can be checked term by term: the full Hilbert-bundle curvature gives the first integral, while the normal component removed by projection contributes its squared norm with a minus sign.
Equivalently, differentiate $P^2=P$ to write
$\partial P=P(\partial P)(I-P)+(I-P)(\partial P)P$ and expand the projected connection. This accounts for all terms.

Put $w=(I-P_t)\sum_i\varphi_i u_i$.
It is orthogonal to the holomorphic sections and satisfies
$\dbar_z w=\sum_i\dbar_z\varphi_i\wedge u_i$,
so it is the minimal-norm solution of this equation.
The fiberwise top-degree $L^2$ estimate gives
\[
 \|w\|^2\le\int_Z
 \sum_{i,j,\alpha,\beta}\varphi^{\bar\beta\alpha}
       \varphi_{i\bar\beta}\varphi_{\alpha\bar j}
       u_i\overline{u_j}e^{-\varphi},
\]
where $(\varphi^{\bar\beta\alpha})$ is the inverse vertical curvature matrix.
This is the inverse-curvature operator on $(\dim Z,1)$-forms from Proposition~\ref{pos:inverse-curvature} and Theorem~\ref{pos:compact-solving}; strict vertical positivity supplies their curvature hypothesis.
Substitution into \eqref{amp:curvature-formula} leaves the matrix
\[
 C_{i\bar j}=\varphi_{i\bar j}
  -\sum_{\alpha,\beta}\varphi_{i\bar\beta}
                       \varphi^{\bar\beta\alpha}\varphi_{\alpha\bar j}.
\]
This is the Schur complement of the vertical block of the full curvature matrix.
It is positive definite: minimize the positive quadratic form over vertical vectors by choosing the vertical vector to be minus the inverse vertical block times the mixed block.
The remaining integrand is therefore nonnegative everywhere and positive on a nonempty open set whenever some $u_i$ is nonzero. Its integral is positive, which is exactly strict Nakano positivity.
\end{proof}

\begin{corollary}[The determinant correction for an ample bundle]\label{amp:berndtsson}
If $E$ is ample, then $\operatorname{Sym}^mE\otimes\det E$ has a strictly Nakano-positive Hermitian metric for every $m\ge0$.
In particular, $\det E$ is positive and $X$ is projective.
\end{corollary}
\begin{proof}
The ample bundle $\mathcal O_Y(1)$ has a positive metric. Take $A=\mathcal O_Y(r+m)$.
Proposition~\ref{amp:euler} gives
$K_{Y/X}\otimes A=\mathcal O_Y(m)\otimes\pi^*\det E$.
The fiber section dimension is the constant $\binom{m+r-1}{r-1}$, and the projective-bundle formula gives
$\pi_*(K_{Y/X}\otimes A)=\operatorname{Sym}^mE\otimes\det E$.
Apply Theorem~\ref{amp:direct-image-curvature}. At $m=0$ this proves positivity of $\det E$, and Theorem~\ref{emb:kodaira} makes $X$ projective.
Only the locally trivial projective-bundle calculation is needed here. Direct-image positivity for varying fiber complex structures has additional terms; see \cite{Berndtsson2009}.
\end{proof}

\section{Serre vanishing and ampleness criteria}
\paragraph{Analytic and algebraic coherent sheaves.}
Serre's GAGA theorem identifies analytic coherent sheaves on a complex projective manifold with algebraic coherent sheaves and identifies their cohomology groups.
We use this theorem, stated in Appendix~\ref{dep:gaga}, to pass to homogeneous polynomial modules.
The local resolutions of Chapter~\ref{alg:chapter} alone do not provide this global algebraization.

\begin{theorem}[Serre vanishing and generation]\label{amp:serre}
If $L$ is ample and $\mathcal F$ is coherent, then for all sufficiently large $m$,
$H^q(X,\mathcal F\otimes L^m)=0$ for $q>0$, and $\mathcal F\otimes L^m$ is globally generated.
\end{theorem}
\begin{proof}
First suppose $L$ is very ample and embeds $i:X\hookrightarrow\mathbb P^N$.
The coherent sheaf $i_*\mathcal F$ has the same cohomology as $\mathcal F$: the direct image under a closed embedding is exact, its sections are defined on the same inverse-image open sets, and resolutions compute the assertion directly.
By GAGA regard $i_*\mathcal F$ as algebraic. We explain its construction from a finite graded $\C[Z_0,\ldots,Z_N]$-module without assuming Serre vanishing.
Take a local module generator $s$ on the standard affine chart $U_i=D_+(Z_i)=\{Z_i\ne0\}$. On $U_i\cap U_j$, localization expresses $s$ with a denominator that is a power of $Z_i/Z_j$. Multiplying by a sufficiently high common power extends it to each $U_j$, as a section of a twist. Differences of these extensions vanish after inverting $Z_i$, so another common power kills them. There are only finitely many charts and overlaps; one exponent works for all of them. The resulting global twisted section still generates the required germs on $U_i$.
Doing this for finitely many generators on all charts gives a surjection
$\bigoplus\mathcal O(-a)\twoheadrightarrow i_*\mathcal F$.
Repeat for its coherent kernel. The resulting map between finite sums has homogeneous polynomial entries; the cokernel of the corresponding graded-module map is the required finitely generated module $M$.
Its graded Hilbert free resolution sheafifies to
\[
 0\to\bigoplus\mathcal O(-a_{\ell j})\to\cdots\to
       \bigoplus\mathcal O(-a_{0j})\to i_*\mathcal F\to0.
\]
For completeness, this resolution is finite because a minimal graded resolution has zero differentials after reduction by the homogeneous maximal ideal, whereas the Koszul complex on the $N+1$ variables gives $\operatorname{Tor}_q(M,\C)=0$ for $q>N+1$.
Thus the minimal resolution terminates there. This is the argument of Theorem~\ref{alg:regular-resolution}, now applied to a global graded module.

Take $m$ larger than every $a_{ij}$.
Each $\mathcal O(m-a_{ij})$ has no higher cohomology: write it as
$K_{\mathbb P^N}\otimes\mathcal O(m-a_{ij}+N+1)$ and apply Corollary~\ref{pos:kodaira} to the positive second factor.
Split the finite resolution into short exact sequences. Starting at its left end, their long exact sequences give higher-cohomology vanishing for $i_*\mathcal F(m)$.
The same argument gives vanishing of $H^1$ of the sufficiently twisted first kernel, so
$H^0(\bigoplus\mathcal O(m-a_{0j}))\to H^0(i_*\mathcal F(m))$ is onto.
Each $\mathcal O(t)$ with $t\ge0$ is generated by homogeneous monomials; the quotient is consequently globally generated.
If only $L^a$ is very ample, write $m=ak+r$, $0\le r<a$, and apply what was proved to the finitely many sheaves $\mathcal F\otimes L^r$.
The maximum of their thresholds works for every sufficiently large $m$, not just for multiples of $a$.
\end{proof}

\begingroup
\tcbset{cgstatement/.append style={unbreakable}}
\begin{proposition}[Line-bundle ampleness criteria]\label{amp:line-criteria}
For a line bundle $L$ on a compact complex manifold, the following are equivalent: ampleness; existence of a positive-curvature metric; very ampleness of every sufficiently large power; higher-cohomology vanishing after sufficiently high twists of every coherent sheaf; and global generation after sufficiently high twists of every coherent sheaf.
Requiring only $H^1(L^m)=H^1(\mathcal I_x\otimes L^m)=0$ does not imply ampleness.
\end{proposition}
\endgroup
\begin{proof}
Ampleness and positive curvature are equivalent by Theorem~\ref{emb:kodaira} and Proposition~\ref{bun:fs}; Theorem~\ref{amp:serre} gives the cohomological and generation conditions. Here $\mathcal I_x$ is the ideal of a point and $\mathcal I_{x,y}=\mathcal I_x\cap\mathcal I_y$ is the ideal of two distinct reduced points.
Conversely, vanishing for all coherent sheaves applies to $\mathcal I_x^2$ and $\mathcal I_{x,y}$ for distinct points. Their long exact sequences make first-jet and two-point evaluations surjective.
Local separation persists on neighborhoods of points and of pairs away from the diagonal. The finite-cover and common-power argument in Theorem~\ref{emb:kodaira}, spelled out again in Proposition~\ref{amp:point-tests} below, then produces an embedding by one power of $L$.
Under the generation condition, first use $\mathcal F=\mathcal O$ to obtain base-point-free powers. Generation of $\mathcal I_x\otimes L^m$ at $y\ne x$ gives a section vanishing at $x$ and nonzero at $y$. At $x$, its images generate $\mathfrak m_x/\mathfrak m_x^2$, so tangent vectors are separated. The same compactness argument gives an embedding.
Once $L$ is positive, fix a very ample power $L^a$. Serre generation makes $L^{m-a}$ globally generated for all large $m$. Tensoring its sections with those of $L^a$ preserves point and tangent separation, so all sufficiently large powers are very ample.
Finally, take $X=\mathbb P^1$ and $L=\mathcal O$. We have $H^1(\mathcal O)=0$ and $\mathcal I_x\simeq\mathcal O(-1)$ also has $H^1=0$, but the degree-zero bundle $L$ is not ample.
\end{proof}

\begin{proposition}[Testing first cohomology and point ideals]\label{amp:point-tests}
In Proposition~\ref{amp:line-criteria}, the vanishing condition can be replaced by
$H^1(X,\mathcal F\otimes L^m)=0$ for every coherent $\mathcal F$ and all sufficiently large $m$, with threshold depending on $\mathcal F$.
The generation condition can be replaced by global generation of both $L^m$ and $\mathcal I_x\otimes L^m$ for each point $x$ and all sufficiently large $m$, with threshold depending on $x$.
\end{proposition}
\begin{proof}
Theorem~\ref{amp:serre} proves necessity. For sufficiency, either hypothesis gives two tasks: near each point, a finite set of sections of some power defines an embedding; and each pair of distinct points is separated by sections of some power. At first the powers may depend on the point or pair.
Under the $H^1$ hypothesis, use
\[
 0\to\mathcal I_x^2\otimes L^m\to L^m
   \to L^m\otimes(\mathcal O_X/\mathcal I_x^2)\to0.
\]
The long exact sequence gives surjective first-jet evaluation. In centered coordinates and a local frame choose sections $s_0,s_1,\ldots,s_n$ with first Taylor terms $1,z_1,\ldots,z_n$.
Near $x$, $s_0\ne0$ and the differentials of $s_j/s_0$ form a coordinate basis. Theorem~\ref{bas:holomorphic-inverse} gives a local embedding on a smaller neighborhood.
Applying the hypothesis to $\mathcal I_{x,y}$ similarly prescribes independent values at two points.

Under the pointwise generation hypothesis, choose a common $m$ for $L^m$ and $\mathcal I_x\otimes L^m$, and a section $s_0(x)\ne0$.
Reducing the generators of the latter sheaf modulo $\mathfrak m_x^2$ gives a spanning set of $\mathfrak m_x/\mathfrak m_x^2$.
Choose $n$ independent linear combinations and divide by $s_0$ to obtain local coordinates.
At $y\ne x$, the ideal stalk is $\mathcal O_y$, so some section $t$ vanishes at $x$ and is nonzero at $y$.
The evaluation vectors of $s_0,t$ at the two points are not proportional, proving separation.

Here is the common-power step, which uses compactness without assuming projectivity.
Choose finitely many neighborhoods $U_i$ covering $X$, with local-embedding section groups of powers $m_i\ge1$.
Let $M$ be a common multiple. All degree-$M/m_i$ monomials in the $i$th group are sections of $L^M$; on $U_i$ they remain base-point free and preserve its local embedding. For instance, where $s_0\ne0$,
\[
 \frac{s_j}{s_0}=\frac{s_js_0^{M/m_i-1}}{s_0^{M/m_i}}.
\]
The combined groups thus define a globally base-point-free system, injective and immersive on every $U_i$.
The set $W=\bigcup_iU_i\times U_i$ is an open neighborhood of the diagonal, and $C=(X\times X)\setminus W$ is compact and disjoint from it.
For each pair in $C$, choose separating sections in some power. In local fiber frames, a two-by-two evaluation minor is nonzero. It remains nonzero on a neighborhood of that pair. Finitely many such neighborhoods cover $C$, involving only finitely many additional powers $m'_j$.
Take $K$ divisible by $M$ and every $m'_j$, and raise all groups to the corresponding homogeneous monomial systems.
These are Veronese maps and preserve distinct nonzero evaluation lines.
The group from $L^M$ supplies global base-point freedom and separates distinct points in $W$; the other groups separate $C$. The original local ratios preserve injectivity of the differential.
Thus the combined sections define an injective holomorphic immersion. Compactness makes it an embedding. Enlarging to the full linear system preserves these properties, so $L$ is ample.
The hypotheses concern all coherent sheaves in the $H^1$ test, or global generation of point ideals. Testing only the two $H^1$ groups of Proposition~\ref{amp:line-criteria} is insufficient.
\end{proof}

\begin{lemma}[A regularity criterion for coherent sheaves]\label{amp:regularity}
Let $A$ be very ample on a projective manifold $X$, and let $V$ be coherent.
If $H^q(X,V\otimes A^{-q})=0$ for every $q>0$, then $V$ is globally generated.
\end{lemma}
\begin{proof}
Embed $i:X\hookrightarrow\mathbb P^N$ by $A$.
Closed-embedding direct image preserves sections, cohomology, and generation, so it suffices to prove the assertion for $F=i_*V$ on $\mathbb P^N$.
By GAGA we may work over Noetherian affine coordinate rings. Write $B(k)=B\otimes\mathcal O(k)$ for twists of a sheaf on $\mathbb P^N$; restriction to a hyperplane means tensoring with its structure sheaf.
Induct on $N$, the case $N=0$ being a single point.

Let $T$ be the largest zero-dimensional coherent subsheaf of $F$.
It exists because finite sums of such subsheaves remain zero-dimensional; if no finite sum were maximal, they would form a strictly increasing coherent-submodule chain, contradicting Corollary~\ref{alg:strong-noetherian}.
Its support is finite, so Proposition~\ref{sh:discrete-support} gives independent specification of its stalk sections, global generation, and vanishing of higher cohomology.
Set $\bar F=F/T$. It has no nonzero zero-dimensional subsheaf, since its inverse image would enlarge $T$.
For every integer $k$, the long exact sequence gives
\[
 H^q(\mathbb P^N,F(k))\simeq H^q(\mathbb P^N,\bar F(k))\quad(q>0),
 \qquad H^0(F(k))\twoheadrightarrow H^0(\bar F(k)).
\]
Thus $\bar F$ satisfies the same vanishings. Once it is generated, lift its generators and add those of $T$ to generate $F$.

Fix $x$. We need a linear equation $\ell$ vanishing at $x$ whose multiplication on $\bar F$ is injective.
The associated primes of a finite Noetherian module are prime ideals of the form $\operatorname{Ann}(v)$ for $v\ne0$. There are finitely many: use the finite prime-factor filtration from Lemma~\ref{alg:uniform-germ}.
If $Rv\simeq R/\mathfrak p$ first meets its $j$th step nontrivially, that intersection injects into the prime factor $R/\mathfrak p_j$.
A nonzero element in either domain quotient has annihilator respectively $\mathfrak p$ and $\mathfrak p_j$, so these primes coincide.
A finite affine cover therefore yields finitely many associated closed subsets.
For $\bar F$ they all have positive dimension: a zero-dimensional associated point would give a cyclic subsheaf supported there, extending by zero to a forbidden global zero-dimensional subsheaf.
Among hyperplanes through $x$, those containing any fixed positive-dimensional subset form a proper linear subspace. Finitely many proper subspaces cannot cover a complex vector space, so choose $\ell(x)=0$ outside all associated primes.

Multiplication $\ell:\bar F(-1)\to\bar F$ is injective.
Otherwise, in a nonzero local kernel choose an element with maximal annihilator. The argument of Lemma~\ref{alg:uniform-germ} makes this annihilator an associated prime containing $\ell$, a contradiction.
Put $H=(\ell=0)\simeq\mathbb P^{N-1}$ and $G=\bar F|_H$. Then
\[
 0\to\bar F(-q-1)\xrightarrow{\ell}\bar F(-q)
       \to i_{H*}G(-q)\to0.
\]
For $q>0$, the adjacent groups $H^q(\bar F(-q))$ and $H^{q+1}(\bar F(-q-1))$ vanish, hence $H^q(H,G(-q))=0$.
Induction makes $G$ globally generated.
Also $H^1(\bar F(-1))=0$, so $H^0(\bar F)\to H^0(G)$ is onto.
Lift finitely many generators at $x$; they generate $\bar F_x/\ell\bar F_x$.
The quotient $Q$ of $\bar F_x$ by their span satisfies $Q=\ell Q$, with $\ell\in\mathfrak m_x$. Lemma~\ref{alg:nakayama} gives $Q=0$.
Thus $\bar F$ is generated at every point. Compactness and local finite generation select finitely many global generators. Together with $T$ they generate $F$.
This argument includes sheaves with zero-dimensional torsion; compare \cite[Tags 02OI, 08A8]{Stacks}.
\end{proof}

\begin{theorem}[Uniform symmetric-power generation]\label{amp:symmetric-criterion}
Fix a very ample $A$ on a projective manifold.
A vector bundle $E$ is ample if and only if $\operatorname{Sym}^mE\otimes A^{-1}$ is globally generated for every sufficiently large $m$.
For the reverse implication, a single $m\ge1$ suffices.
\end{theorem}
\begin{proof}
If $E$ is ample, apply Serre vanishing on $Y=\mathbb P(E^*)$ to the finitely many sheaves $\pi^*A^{-1-q}$, $q=1,\ldots,\dim X$.
For one common threshold,
$H^q(Y,\mathcal O_Y(m)\otimes\pi^*A^{-1-q})=0$.
The projective-bundle direct-image formulas, Theorem~\ref{sh:projective-direct-image}, identify these with
$H^q(X,\operatorname{Sym}^mE\otimes A^{-1-q})$.
Lemma~\ref{amp:regularity} gives generation.
Conversely, generators of $\operatorname{Sym}^mE\otimes A^{-1}$ give every degree-$m$ homogeneous polynomial on each fiber. They define a base-point-free system for $\mathcal O_Y(m)\otimes\pi^*A^{-1}$, whose Fubini--Study metric is semipositive and strictly positive on nonzero vertical vectors, since the fiber map is Veronese.
Tensor with the positive metric on $\pi^*A$. Every vector with nonzero base projection now has positive curvature as well.
Thus $\mathcal O_Y(m)$ is positive; taking an $m$th root of its metric makes $\mathcal O_Y(1)$ positive and proves ampleness of $E$.
\end{proof}

\begin{proposition}[Cohomological criteria for ample vector bundles]\label{amp:vector-criteria}
A holomorphic vector bundle $E$ on a compact complex manifold is ample if and only if, for every coherent $\mathcal F$, the sheaves $\mathcal F\otimes\operatorname{Sym}^mE$ are globally generated for all sufficiently large $m$.
Equivalently, all their positive-degree cohomology groups vanish for sufficiently large $m$.
The thresholds may depend on $\mathcal F$.
In that case every sufficiently large $\operatorname{Sym}^mE$ is very ample, meaning that the tautological $\mathcal O(1)$ on its projective bundle is very ample.
\end{proposition}
\begin{proof}
Write $Y=\mathbb P(E^*)$, $L=\mathcal O_Y(1)$, and $E_m=\operatorname{Sym}^mE$. Let $\mathcal E_m$ be its sheaf of holomorphic sections; $(\mathcal E_m)_x$ is an $\mathcal O_{X,x}$-module stalk, not a vector-space fiber.
If $E$ is ample, then $L$ is ample and Corollary~\ref{amp:berndtsson} makes $X$ projective.
Serre vanishing for $\pi^*\mathcal F\otimes L^m$, followed by the coherent-coefficient version of Theorem~\ref{sh:projective-direct-image}, gives
\[
 H^q(X,\mathcal F\otimes E_m)
   \simeq H^q(Y,\pi^*\mathcal F\otimes L^m)=0\quad(q>0)
\]
for large $m$.
For generation, fix very ample $A$ and choose $b\ge1$ so that $\mathcal F\otimes A^b$ is generated.
Theorem~\ref{amp:symmetric-criterion}, applied to $A^b$, makes $E_m\otimes A^{-b}$ generated for all large $m$.
Tensor their generating sections to generate $\mathcal F\otimes E_m$; this also works when $\mathcal F$ is not locally free.

For the converse, construct separating sections on $Y$ before assuming that $X$ is projective.
Near $x\in X$, choose coordinates $z_1,\ldots,z_n$, a frame of $E$, and homogeneous fiber coordinates $\xi_0,\ldots,\xi_{r-1}$.
The direct-image formula identifies global sections of $E_m$ with those of $L^m$.
Under the vanishing hypothesis apply it to $\mathcal I_x^2$. For a sufficiently large $m\ge1$, the exact sequence gives
\[
 H^0(X,E_m)\twoheadrightarrow
 (\mathcal E_m)_x/\mathfrak m_x^2(\mathcal E_m)_x.
\]
Tensoring preserves injectivity here because $E_m$ is locally free.
On the fiber chart $\xi_0\ne0$, prescribe the first-order base expansions
\[
 P_0=\xi_0^m,\qquad P_a=\xi_a\xi_0^{m-1}\ (1\le a<r),
 \qquad z_jP_0\ (1\le j\le n).
\]
Dividing by the first section, their differentials above $x$ give $d(\xi_a/\xi_0)$ and $dz_j$.
Thus this finite group is base-point free near the specified fiber point and gives a local embedding. Use another nonzero homogeneous coordinate to cover the fiber.

Under the generation hypothesis, choose a common large $m$ so that $E_m$ and $\mathcal I_x\otimes E_m$ are both generated.
The first condition prescribes the fiber values $P_0,P_a$.
The second, modulo $\mathfrak m_x^2$, spans
$(\mathfrak m_x/\mathfrak m_x^2)\otimes(\mathcal E_m)_x/\mathfrak m_x(\mathcal E_m)_x$,
so first-order terms $z_jP_0$ can also be prescribed.
These latter sections vanish on the whole central fiber, hence have zero vertical derivatives; their ratios have horizontal derivatives $dz_j$.
The other ratios still have the vertical derivatives $d(\xi_a/\xi_0)$. Any extra horizontal terms leave the full differential matrix block triangular and invertible.
This again gives local embeddings.

Now separate two distinct points of $Y$.
If they lie in one fiber, full fiber evaluation gives the degree-$m$ Veronese system, which separates them.
If their base points are $x\ne y$, the vanishing hypothesis for $\mathcal I_{x,y}$ gives independent fiber evaluation at $x,y$, allowing the required polynomial values.
Under the generation hypothesis, use $\mathcal I_x\otimes E_m$ to choose a section zero on the entire $x$ fiber and nonzero at the specified point over $y$, and use $E_m$ to choose another nonzero at the specified point over $x$. The two evaluation vectors are not proportional.
Since $Y$ is compact, the finite-group, common-power, and Veronese argument in Proposition~\ref{amp:point-tests} applies to $L$ on $Y$.
It uses only the local embeddings and pair separation just proved, and gives an embedding by $L^K$.
Thus $E$ is ample. Only now do we use Corollary~\ref{amp:berndtsson} to conclude that $X$ is projective.

Finally fix very ample $A$ on $X$. For all large $m$, $V=E_m$ has $V\otimes A^{-1}$ globally generated.
On $\rho:Y_m=\mathbb P(V^*)\to X$, the generators define a base-point-free map for
$B=\mathcal O_{Y_m}(1)\otimes\rho^*A^{-1}$ that is a full linear embedding on each fiber.
Combine it with the embedding of the base defined by $A$.
Distinct base points are separated by the second factor, and distinct points in one fiber by the first.
A tangent vector killed by the product map first has zero base projection and is then zero by the fiber immersion.
The product is therefore an injective immersion, hence an embedding by compactness.
Its Segre composition is defined by products of sections of
$B\otimes\rho^*A=\mathcal O_{Y_m}(1)$.
The complete linear system still embeds, proving very ampleness of $V$.
\end{proof}

\section{Operations on ample bundles}
\begin{theorem}[Quotients, direct sums, and tensor products]\label{amp:operations}
A nonzero locally free quotient of an ample vector bundle is ample.
The direct sum $E\oplus F$ is ample if and only if both factors are ample.
The tensor product of an ample bundle with a nonzero globally generated bundle is ample, and the tensor product of two ample bundles is ample.
\end{theorem}
\begin{proof}
A quotient $E\twoheadrightarrow Q$ gives a closed embedding
$\mathbb P(Q^*)\hookrightarrow\mathbb P(E^*)$ preserving tautological bundles.
Restrict the positive curvature metric to prove ampleness of $Q$.
If a direct sum is ample, its factors are quotients and hence ample.
Conversely, Corollary~\ref{amp:berndtsson} makes $X$ projective; fix a very ample line bundle $A$.
There is $m_0$ such that $\operatorname{Sym}^pE\otimes A^{-1}$ and $\operatorname{Sym}^pF\otimes A^{-1}$ are generated for $p\ge m_0$.
For $p,q\ge m_0$, the bundle
$\operatorname{Sym}^pE\otimes\operatorname{Sym}^qF\otimes A^{-1}$
is a tensor product of these two generated bundles with $A$, hence is generated.
For the finitely many $q<m_0$, Serre vanishing on $\mathbb P(E^*)$ with fixed twists $\operatorname{Sym}^qF\otimes A^{-1-j}$, followed by regularity, gives the same conclusion for all sufficiently large $p$. Interchange the factors for $p<m_0$.
Thus a single sufficiently large bound on $p+q$ works for every splitting.
Use
$\operatorname{Sym}^k(E\oplus F)=\bigoplus_{p+q=k}\operatorname{Sym}^pE\otimes\operatorname{Sym}^qF$
and the generation criterion to prove ampleness of $E\oplus F$.
If $F$ is globally generated, it is a quotient of a finite trivial bundle, so $E\otimes F$ is a quotient of a finite direct sum of copies of $E$. The preceding results apply.

For two ample factors we need an additional argument; ampleness alone is not global generation.
In characteristic zero, every representation of a finite group is completely reducible: average any projection over the group to obtain an equivariant one.
Decompose the permutation action of $S_k$ as
$V^{\otimes k}=\bigoplus_\lambda M_\lambda\otimes W_\lambda(V)$,
where $M_\lambda$ are irreducible and
$W_\lambda(V)=\Hom_{S_k}(M_\lambda,V^{\otimes k})$.
Here $\lambda$ runs through irreducible representations of the permutation group $S_k$; write $\lambda^*$ for the label of $M_\lambda^*$. Schur's lemma says that an equivariant map between irreducibles is zero or an isomorphism, and each equivariant endomorphism is scalar. Indeed, kernel and image are invariant; subtracting an eigenvalue from an endomorphism then proves the scalar assertion.
Taking diagonal invariants and using this lemma gives
\[
 \operatorname{Sym}^k(V\otimes W)
   =\bigoplus_\lambda W_\lambda(V)\otimes W_{\lambda^*}(W).
\]
The equivariant projectors make these decompositions valid for bundles.
If $r=\operatorname{rank}E$, every nonzero $W_\lambda(E)$ occurs as a direct summand of
$\operatorname{Sym}^k(E\otimes\C^r)$, up to a positive finite multiplicity.
Indeed, $M_{\lambda^*}$ occurs in $(\C^r)^{\otimes k}$ whenever $M_\lambda$ does: this permutation representation has an invariant nondegenerate bilinear form and is self-dual.
The bundle $E\otimes\C^r=E^{\oplus r}$ is already known to be ample.
For all sufficiently large $k$, generation of its twisted symmetric power therefore makes all summands $W_\lambda(E)\otimes A^{-1}$ generated at once. The threshold is independent of $\lambda$ because they come from one direct-sum decomposition.
The same applies to $F$.
Each summand in the displayed formula, twisted by $A^{-1}$, is consequently a product of two bundles twisted by $A^{-1}$ and the generated line bundle $A$.
Thus $\operatorname{Sym}^k(E\otimes F)\otimes A^{-1}$ is generated for large $k$. The criterion proves the result.
\end{proof}

\begin{corollary}[Symmetric powers, exterior powers, and determinants]\label{amp:powers}
If $E$ is ample, then $\operatorname{Sym}^kE$ for $k\ge1$ and $\Lambda^kE$ for $1\le k\le r$ are ample.
If some $\operatorname{Sym}^kE$ with $k\ge1$ is ample, then $E$ is ample.
\end{corollary}
\begin{proof}
Theorem~\ref{amp:operations} makes $E^{\otimes k}$ ample.
Symmetrization and antisymmetrization in characteristic zero give locally free quotient bundles, proving the first assertion.
For the converse, the fiberwise Veronese embedding
$\nu:\mathbb P(E^*)\hookrightarrow\mathbb P((\operatorname{Sym}^kE)^*)$
satisfies $\nu^*\mathcal O(1)=\mathcal O(k)$.
Restrict the positive metric of the ambient tautological bundle and take its $k$th root. This makes $\mathcal O(1)$ positive on $\mathbb P(E^*)$, proving ampleness of $E$.
\end{proof}

\begin{proposition}[Extensions of ample bundles]\label{amp:extension}
In a short exact sequence $0\to E_1\to E_2\to E_3\to0$, ampleness of $E_1,E_3$ implies ampleness of $E_2$.
\end{proposition}
\begin{proof}
Corollary~\ref{amp:berndtsson} makes $X$ projective. Fix a very ample $A$ and put $n=\dim X$. Filter $\operatorname{Sym}^mE_2$ by the number of factors from $E_1$.
A local holomorphic splitting identifies its successive quotients as
$\operatorname{Sym}^pE_1\otimes\operatorname{Sym}^{m-p}E_3$, $0\le p\le m$.
Different local splittings differ by upper triangular transformations preserving this filtration, so it is global.
The ample bundle $E_1\oplus E_3$ has higher-cohomology vanishing for its high symmetric powers twisted by each of the finitely many $A^{-1-j}$, $j=1,\ldots,n$.
All the displayed quotients are direct summands of that symmetric power, so the vanishings hold simultaneously for every $p$.
Apply long exact sequences successively along the filtration to obtain
$H^j(X,\operatorname{Sym}^mE_2\otimes A^{-1-j})=0$
for every $j>0$ and large $m$.
Lemma~\ref{amp:regularity} makes $\operatorname{Sym}^mE_2\otimes A^{-1}$ generated, and Theorem~\ref{amp:symmetric-criterion} makes $E_2$ ample.
The filtration quotients and twists are locally free, so tensoring preserves the short exact sequences used in this argument.
\end{proof}

\begin{proposition}[Fixed twists and Euler characteristic growth]\label{amp:twist-growth}
If $E$ is ample of rank $r$ and $F$ is a fixed nonzero vector bundle, then $\operatorname{Sym}^mE\otimes F$ is ample for large $m$.
For $\dim X=n$, its Euler characteristic has degree $n+r-1$ in $m$, which cannot in general be replaced by $n$.
\end{proposition}
\begin{proof}
Choose very ample $A$, then choose $b$ with $F\otimes A^b$ globally generated.
Serre vanishing on $\mathbb P(E^*)$ and regularity give generation of
$\operatorname{Sym}^mE\otimes A^{-b-1}$ for large $m$.
Thus $\operatorname{Sym}^mE\otimes A^{-b}$ is the tensor product of the ample line bundle $A$ with a generated bundle and is ample.
Tensor it with the generated $F\otimes A^b$ and apply the operations theorem.
For the degree, $Y=\mathbb P(E^*)$ has dimension $n+r-1$ and
$\chi(X,\operatorname{Sym}^mE\otimes F)=\chi(Y,\mathcal O_Y(m)\otimes\pi^*F)$.
The Hilbert polynomial of the ample tautological bundle has degree $n+r-1$, with positive leading coefficient
$\operatorname{rank}(F)\int_Yc_1(\mathcal O(1))^{n+r-1}/(n+r-1)!$.
The independent algebraic proof of this polynomial formula appears in Theorem~\ref{idx:hilbert-polynomial}; the ampleness argument above uses only Serre vanishing.
For instance, if $X$ is a point, the dimension is
$\operatorname{rank}(F)\binom{m+r-1}{r-1}$, a polynomial of degree $r-1$.
\end{proof}

\section{Le Potier comparison and vanishing}
\begin{lemma}[A family of acyclic bundles on projective space]\label{amp:relative-forms}
On $\mathbb P^{r-1}$, for $1\le j\le r-1$,
$H^q(\mathbb P^{r-1},\Omega^j(1))=0$ for every $q\ge0$.
\end{lemma}
\begin{proof}
Dualizing Proposition~\ref{amp:euler} and taking exterior powers gives the exact Koszul-type sequence
\[
 0\to\Omega^j(1)\to\Lambda^j\C^{r*}\otimes\mathcal O(1-j)
 \to\cdots\to\C^{r*}\otimes\mathcal O\to\mathcal O(1)\to0.
\]
Where a homogeneous coordinate is nonzero, contraction with its normalized coordinate vector gives an explicit contracting homotopy and verifies exactness.
The line-bundle degrees range from $1-j$ to $1$, all strictly greater than $-r$.
Corollary~\ref{pos:kodaira} therefore kills their positive-degree cohomology. A negative-degree bundle has no sections: restrict to any projective line through a point and use Proposition~\ref{bun:sections-p1}.
The map on sections of the last two terms is the natural isomorphism $\C^{r*}\to H^0(\mathcal O(1))$.
Thus the global-section complex is exact as well.
Elimination through the finite exact sequence, as in Lemma~\ref{sh:double-complex}, gives the asserted vanishing in every degree.
\end{proof}

\begin{theorem}[Le Potier's comparison isomorphism]\label{amp:le-potier}
For a coherent sheaf $\mathcal F$ and a vector bundle $E$ on $X$, there is a natural isomorphism
\[
 H^q(X,\Omega_X^p\otimes E\otimes\mathcal F)
 \simeq H^q(Y,\Omega_Y^p\otimes\mathcal O_Y(1)\otimes\pi^*\mathcal F).
\]
If $E$ is ample, $\dim X=n$, and $\operatorname{rank}E=r$, then
$H^q(X,\Omega_X^p\otimes E)=0$ whenever $p+q>n+r-1$.
\end{theorem}
\begin{proof}
The sequence
$0\to\pi^*\Omega_X^1\to\Omega_Y^1\to\Omega_{Y/X}^1\to0$
filters $\Omega_Y^p$ with quotients
$\pi^*\Omega_X^{p-j}\otimes\Omega_{Y/X}^j$.
The locally trivial projection $\pi$ is flat, and the filtration quotients are locally free. Tensoring by $\pi^*\mathcal F$ therefore preserves the short exact sequences.
For $j>0$, Lemma~\ref{amp:relative-forms} and the relative computation in Theorem~\ref{sh:projective-direct-image} make every direct image, including degree zero, vanish after the $\mathcal O_Y(1)$ twist.
For general coherent $\mathcal F$, use a finite free resolution locally on the base and calculate each term; its finite length justifies elimination.
For $j=0$, the direct-image formula gives $\Omega_X^p\otimes E\otimes\mathcal F$ in degree zero and zero in higher degrees.
Apply the direct-image long exact sequences along the filtration. The Leray double complex then has just its degree-zero row, proving the natural cohomology isomorphism.
If $E$ is ample, $Y$ has positive $\mathcal O_Y(1)$ and dimension $n+r-1$.
Corollary~\ref{pos:akn} on $Y$ gives precisely the stated range.
\end{proof}

\begin{corollary}[Truncation of the Euler characteristic]\label{amp:euler-truncation}
If $K_X^{-1}\otimes E$ is ample of rank $r$, then $H^q(X,E)=0$ for $q\ge r$ and
$\chi(X,E)=\sum_{q=0}^{r-1}(-1)^qh^q(X,E)$.
\end{corollary}
\begin{proof}
Take $p=n$ in Theorem~\ref{amp:le-potier} and replace the coefficient bundle by $K_X^{-1}\otimes E$.
The inequality $n+q>n+r-1$ is equivalent to $q\ge r$, and
$\Omega_X^n\otimes K_X^{-1}\otimes E=E$.
All remaining terms of the Euler characteristic are zero.
\end{proof}


\chapter{Euler Characteristics, Hilbert Polynomials, and Integral Chern Classes}\label{idx:chapter}
The alternating sum of cohomology dimensions is often easier to compute than its individual terms.
Riemann--Roch expresses the holomorphic Euler characteristic through Chern classes; Chern--Gauss--Bonnet computes the topological Euler characteristic.
Their different inputs explain why the two numbers need not agree.
We derive higher-dimensional Riemann--Roch from the Atiyah--Singer index theorem and also give a separate algebraic proof of the Hilbert polynomial.
The final section concerns integral, rather than only real, Chern classes and uses the additional algebraic topology stated there.

\section{Higher-dimensional Riemann--Roch}
\paragraph{Notation for the index formula.}
For a first-order Dirac-type operator $D^+$ from even to odd sections, viewed as a map from $H^1$ to $L^2$, its Fredholm index is $\operatorname{ind}D^+=\dim\ker D^+-\dim\ker(D^+)^*$. For the operators used here on compact $X$, the compactness and regularity arguments of Lemma~\ref{hod:compact-resolvent} give finite-dimensional kernels and closed range.
The Chern character $\operatorname{ch}$ is defined in Theorem~\ref{bun:chern-weil}. If $x_1,\ldots,x_n$ are the formal degree-two Chern roots of $T^{1,0}X$, set
\[
 \operatorname{Td}(T^{1,0}X)=\prod_j\frac{x_j}{1-e^{-x_j}},
 \qquad \widehat A(TX_\R)=\prod_j\frac{x_j/2}{\sinh(x_j/2)}.
\]
These symmetric power series define cohomology classes through the Chern classes, with terms above real degree $2n$ discarded; $TX_\R$ means the underlying real tangent bundle.
A $\operatorname{Spin}^c$ structure is a lift of the oriented orthonormal frame bundle to $(\operatorname{Spin}(2n)\times U(1))/\{\pm(1,1)\}$. Its determinant line is associated to $[g,z]\mapsto z^2$, and its spinor bundle carries the Clifford action $c(\xi)^2=-|\xi|^2$ for real covectors $\xi$. The associated first-order Dirac operator has principal symbol $\ii c(\xi)$ in the Fourier convention $\partial_{x_j}\mapsto\ii\xi_j$. We use this construction only as part of the following external index input.

\paragraph{The $\operatorname{Spin}^c$ index theorem.}
We use the Atiyah--Singer $\operatorname{Spin}^c$ index formula: on a compact even-dimensional manifold, a $\operatorname{Spin}^c$ Dirac operator with determinant line bundle $B$, twisted by a complex bundle $E$, has index
$\int\operatorname{ch}(E)e^{c_1(B)/2}\widehat A(TX_\R)$.
We use this index theorem without proof; references and its precise hypotheses are given in Appendix~\ref{dep:index}.
Its full characteristic-class formula requires more than the leading heat-kernel asymptotic of Chapter~\ref{morse:chapter}.

\begingroup
\tcbset{cgstatement/.append style={unbreakable}}
\begin{theorem}[Hirzebruch--Riemann--Roch]\label{idx:hrr}
Let $X$ be a compact complex manifold of dimension $n$ and $E$ a holomorphic vector bundle. Then
\[
 \chi(X,E)=\sum_{q=0}^n(-1)^q\dim H^q(X,E)
 =\int_X\operatorname{ch}(E)\operatorname{Td}(T^{1,0}X).
\]
The Todd class uses the normalization fixed above.
\end{theorem}
\endgroup
\begin{proof}
On $E$-valued $(0,\bullet)$-forms consider
$D_E=\sqrt2(\dbar_E+\dbar_E^*)$, which interchanges even and odd degrees. Since $D_E^2=2\Delta''_E$, Theorem~\ref{hod:hodge} gives
\[
 \operatorname{ind}(D_E^+)
 =\dim\bigoplus_{q\text{ even}}\mathcal H^{0,q}(E)
  -\dim\bigoplus_{q\text{ odd}}\mathcal H^{0,q}(E)=\chi(X,E).
\]
The complex structure determines a natural $\operatorname{Spin}^c$ structure whose spinor bundle is $\Lambda^{0,\bullet}T^*X$ and whose determinant line bundle is $K_X^{-1}$. Concretely, let $\iota_{\xi^{0,1}}$ denote the Hermitian adjoint of exterior multiplication by $\xi^{0,1}$. A real covector $\xi$ acts by $\sqrt2(\xi^{0,1}\wedge-\iota_{\xi^{0,1}})$, whose square is $-|\xi|^2$.
The natural Dirac operator and $D_E$ have the same principal symbol. In the non-K\"ahler case their difference consists of zero-order torsion terms. They have equal indices, because the path between them consists of elliptic Fredholm operators and the index is locally constant along a continuous Fredholm family. To see the latter assertion, separate the finite-dimensional small spectral space near zero from its orthogonal complement. The operator on the complement remains invertible; on the finite-dimensional part, the difference between kernel and cokernel dimensions equals the fixed difference between domain and target dimensions.

The index formula therefore yields
$\chi(X,E)=\int_X\operatorname{ch}(E)e^{c_1(T^{1,0}X)/2}\widehat A(TX_\R)$.
The roots of the complexification of the real tangent bundle occur in pairs $\pm x_j$, so
$\widehat A(TX_\R)=\prod_j(x_j/2)/\sinh(x_j/2)$. Multiplying one factor at a time gives
\[
 e^{x/2}\frac{x/2}{\sinh(x/2)}=\frac{x}{1-e^{-x}}.
\]
This is the Todd factor, proving the formula. Only the component of total degree $2n$ is integrated, so every formal series automatically truncates to a finite sum.
The holomorphic structure on $E$ is needed to define the sheaf cohomology on the left; a smooth complex vector bundle alone does not specify it.
\end{proof}

\begin{corollary}[Asymptotic Riemann--Roch and the full polynomial]\label{idx:asymptotic-rr}
For a holomorphic line bundle $L$ and a rank-$r$ holomorphic bundle $E$ on a compact complex $n$-fold, write $[\eta]_j$ for the component of real cohomological degree $j$ of an inhomogeneous class $\eta$. Then
\[
 \chi(X,L^k\otimes E)
 =\sum_{j=0}^n\frac{k^j}{j!}
     \int_Xc_1(L)^j\wedge[\operatorname{ch}(E)\operatorname{Td}(T^{1,0}X)]_{2n-2j}.
\]
In particular,
$\chi(X,L^k\otimes E)=\frac{rk^n}{n!}\int_Xc_1(L)^n+O(k^{n-1})$.
\end{corollary}
\begin{proof}
Theorem~\ref{idx:hrr} and the multiplicativity in Proposition~\ref{bun:chern-functors} give
$\operatorname{ch}(L^k\otimes E)=e^{kc_1(L)}\operatorname{ch}(E)$.
Expand the exponential through degree $n$; higher terms vanish on a real $2n$-manifold. For each $j$, retain the complementary-degree component of the other factor. This gives the displayed identity.
For $j=n$, that complementary degree is zero, and $\operatorname{ch}_0(E)=r$, $\operatorname{Td}_0=1$. All remaining powers of $k$ have degree at most $n-1$. This proves the leading term and the remainder. The assertion concerns the alternating sum, not the asymptotic behavior of an individual $h^q$.
\end{proof}

\begin{corollary}[Chern--Gauss--Bonnet]\label{idx:cgb}
For a compact complex manifold, let $b_k(X)=\dim_\C H^k_{\mathrm{dR}}(X,\C)$ be its $k$th Betti number. Then
\[
 \chi_{\mathrm{top}}(X)=\sum_{k=0}^{2n}(-1)^kb_k(X)
                      =\int_Xc_n(T^{1,0}X).
\]
In general $\chi_{\mathrm{top}}(X)$ is not $\chi(X,T^{1,0}X)$.
\end{corollary}
\begin{proof}
If $x_i$ are the formal Chern roots of the tangent bundle, the cotangent roots are $-x_i$. Consequently
\[
 \sum_p(-1)^p\operatorname{ch}(\Omega_X^p)
 =\sum_I(-1)^{|I|}e^{-\sum_{i\in I}x_i}
 =\prod_i(1-e^{-x_i}).
\]
Multiplication by the Todd class cancels each denominator and leaves $\prod_i x_i=c_n(T^{1,0}X)$. Applying Theorem~\ref{idx:hrr} to each $\Omega_X^p$ and taking the alternating sum gives
$\int c_n=\sum_{p,q}(-1)^{p+q}h^{p,q}$.
For K\"ahler $X$, Theorem~\ref{hod:hodge} and Proposition~\ref{hod:kahler-laplacians} identify this with the topological Euler characteristic.
For arbitrary compact complex $X$, let $C^k=\bigoplus_{a+b=k}\mathcal A^{a,b}(X)$ with differential $d=\partial+\dbar$, and set
\[
 F^pC^k=\bigoplus_{\substack{a\ge p\\a+b=k}}\mathcal A^{a,b}(X).
\]
Then $F^0C=C$, $F^{n+1}C=0$, and $F^pC/F^{p+1}C$ is the Dolbeault complex of $(p,\bullet)$-forms, placed in total degree $p+q$. The short exact sequences of complexes
\[
 0\longrightarrow F^{p+1}C\longrightarrow F^pC
 \longrightarrow F^pC/F^{p+1}C\longrightarrow0
\]
give long exact cohomology sequences by the cochain-lifting argument in the proof of Proposition~\ref{sh:long-exact}. The quotient cohomology is finite-dimensional by Theorem~\ref{hod:hodge}. Descending induction on $p$ therefore makes every $H^k(F^pC)$ finite-dimensional. Write $\chi(F^pC)=\sum_k(-1)^k\dim H^k(F^pC)$. Additivity of alternating dimensions in the finite long exact sequences gives
\[
 \chi(F^pC)=\chi(F^{p+1}C)+\sum_q(-1)^{p+q}h^{p,q}.
\]
Summing over $p$ yields $\sum_k(-1)^kb_k=\sum_{p,q}(-1)^{p+q}h^{p,q}$. Thus the same formula holds without the K\"ahler assumption, and no degeneration of a spectral sequence is needed.
Finally, on $\mathbb P^1$ we have $T\mathbb P^1=\mathcal O(2)$ and hence $\chi(\mathbb P^1,T)=3$. But $b_0=b_2=1$ and $b_1=0$, so the topological Euler characteristic is $2$.
\end{proof}

\section{Hilbert polynomials without the index theorem}
\begin{lemma}[The Euler characteristic of projective space]\label{idx:projective-polynomial}
For every integer $t$,
\[
 \chi(\mathbb P^N,\mathcal O(t))=\binom{t+N}{N},
\]
where the right-hand side is interpreted as a degree-$N$ polynomial in $t$.
\end{lemma}
\begin{proof}
For $N=0$ both sides are one. A hyperplane $H\simeq\mathbb P^{N-1}$ gives
$0\to\mathcal O(t-1)\to\mathcal O(t)\to\mathcal O_H(t)\to0$, whence
$\chi_N(t)-\chi_N(t-1)=\chi_{N-1}(t)$.
By induction the right-hand side is $\binom{t+N-1}{N-1}$. Pascal's identity shows that the proposed polynomial satisfies the same recurrence.
At $t=0$, holomorphic functions are constant and Kodaira vanishing gives $H^q(\mathcal O)=0$ for $q>0$, since $\mathcal O=K\otimes\mathcal O(N+1)$ by Proposition~\ref{amp:euler}; use Corollary~\ref{pos:kodaira}. Thus $\chi_N(0)=1$. Running the recurrence in both directions proves the identity at every integer.
\end{proof}

\begin{theorem}[Support dimension and the degree of the Hilbert polynomial]\label{idx:hilbert-polynomial}
Let $X$ be projective, $A$ ample, and $\mathcal F\ne0$ coherent with support of dimension $d$. Then
$P(m)=\chi(X,\mathcal F\otimes A^m)$ is an integer-valued polynomial of degree $d$. Let $Y$ run over the $d$-dimensional irreducible components of the support, and let $\ell_Y(\mathcal F)$ be the generic length along $Y$: after algebraization, it is the composition length of the stalk $\mathcal F_{\eta_Y}$ at the generic point $\eta_Y$. For $\mathcal F=\mathcal O_X/\mathcal I$ it records the multiplicity of the possibly nonreduced structure. The leading coefficient is
\[
 \frac1{d!}\sum_{\dim Y=d}\ell_Y(\mathcal F)\int_Yc_1(A)^d>0.
\]
\end{theorem}
\begin{proof}
First suppose $A$ is very ample. The GAGA input in Appendix~\ref{dep:gaga} and the finite graded free resolution in Theorem~\ref{amp:serre} express the Euler characteristic of $i_*\mathcal F$ as a finite alternating sum of terms $\chi(\mathcal O(m-a))$. Lemma~\ref{idx:projective-polynomial} makes this a polynomial, integer-valued at every integer $m$.
We determine its degree by induction on the support dimension. After removing the maximal zero-dimensional subsheaf, choose a general hyperplane whose equation avoids the finitely many associated primes of the remaining sheaf, as justified in Lemma~\ref{amp:regularity}. Multiplication by this equation is injective. The quotient has support dimension one less, and
$P(m)-P(m-1)=P_{\mathcal F|H}(m)$ for the sheaf with no zero-dimensional subsheaf.
In dimension zero, higher cohomology vanishes and $P$ is the positive total length. After $d$ general hyperplane sections, the resulting zero-dimensional length is positive and equals the sum of the intersection numbers of the top-dimensional components, weighted by their generic lengths $\ell_Y$. Thus the $d$th difference is this positive integer and the next difference is zero. The polynomial has degree exactly $d$, with leading coefficient that integer divided by $d!$.
The intersection number equals $\int_Yc_1(A)^d$: apply Theorem~\ref{cur:poincare-lelong} successively to the hyperplane sections, replacing one Chern factor at a time by the corresponding divisor current, until only zero-dimensional intersection multiplicities remain.
A zero-dimensional subsheaf contributes a constant to the Euler characteristic and does not change the leading coefficient when $d>0$. This justifies its removal in the induction.

For general ample $A$, choose $a$ with $A^a$ very ample. The preceding argument gives a polynomial in $k$ on each residue class $m=ak+r$. To show that these are restrictions of a single polynomial in $m$, one may use HRR on a finite vector-bundle resolution. There is also a direct algebraic argument: multiplication by $[A]-1$ in the Grothendieck group lowers the support filtration. This group is generated by coherent-sheaf symbols $[\mathcal F]$, with the relation $[\mathcal F]=[\mathcal F']+[\mathcal F'']$ for every short exact sequence $0\to\mathcal F'\to\mathcal F\to\mathcal F''\to0$; its support filtration is generated by sheaves of support dimension at most a specified integer. Here $1=[\mathcal O_X]$, and $[A]$ acts by tensoring with the line bundle $A$. Indeed, $A$ is trivial on a nonempty open subset of each irreducible component. Therefore its difference from the trivial bundle, acting on the class of a coherent sheaf, is represented by sheaves of smaller support dimension. Define the forward difference by $\Delta P(m)=P(m+1)-P(m)$. It follows that $\Delta^{d+1}P(m)=0$ for every integer $m$, so $P$ is a single polynomial. Its leading coefficient is already determined on the subsequence of multiples of $a$, and rescaling yields the stated expression.
For completeness, the support-filtration assertion can be seen using a rational trivialization and its zero and pole divisors. Embed both $A$ and the trivial bundle into a common sufficiently large divisor line bundle. Tensoring these maps with a coherent sheaf need not preserve injectivity, but their kernels and cokernels are all supported on a divisor avoiding the generic points of the original top-dimensional support. The difference of the two classes therefore has smaller support dimension. Repeating this argument proves the required finite-difference identity without the index theorem.
\end{proof}

\section{Integral Chern classes and Peterson's criterion}
\paragraph{Topological notation.}
A CW complex $K$ is built by attaching real disks along their boundary spheres; $K^{(j)}$ is the union of its cells of dimension at most $j$. The group $\pi_j(Y)$ consists of based homotopy classes of maps $S^j\to Y$. Write $U(r)$ for the unitary group and $BU(r)$ for its classifying space, modeled by the infinite complex Grassmannian, whose tautological rank-$r$ bundle classifies rank-$r$ complex bundles. The cohomology comparisons and relative-cohomology operations used below are listed in Appendix~\ref{dep:good-cover}.

\begin{proposition}[Compatibility of integral and differential-form Chern classes]\label{idx:integral-chern}
A smooth complex vector bundle has natural integral Chern classes
$c_j(E)\in H^{2j}(X,\Z)$ whose images in real cohomology are the Chern--Weil classes defined earlier. Vanishing of the real classes alone does not exclude nonzero integral torsion classes.
\end{proposition}
\begin{proof}
The line-bundle case was proved in Lemma~\ref{prj:integral-gluing}. For a rank-$r$ bundle, let $\pi:\mathbb P(E^*)\to X$ be its projective bundle. Its integral cohomology is a free module over the cohomology of the base, with basis
$1,h,\ldots,h^{r-1}$, where $h=c_1(\mathcal O(1))$.
On a trivializing set this is the cellular calculation for $\mathbb P^{r-1}$. The same global classes give the basis on overlaps, so Mayer--Vietoris and the five lemma glue the local isomorphisms. More generally, for a finite-dimensional CW base, argue over the skeleta using relative cellular groups. The bundle is trivial over each cell; cup product with the global classes gives the relative isomorphism, and the long exact sequences of the pairs give the induction. Even if a skeleton has infinitely many cells, cellular cochains use the corresponding product and the same finite fiber basis still gives the isomorphism. There are only finitely many skeletal steps, so no interchange with an inverse limit is involved.
There are consequently unique integral classes $c_j(E)$ satisfying
$h^r-\pi^*c_1(E)h^{r-1}+\cdots+(-1)^r\pi^*c_r(E)=0$.
Uniqueness gives naturality.
Iterated projectivization produces a flag bundle $p:F\to X$; pullback is injective at each step because the displayed basis contains $1$. Hermitian orthogonal complements split the pulled-back smooth bundle as $p^*E=\bigoplus_{j=1}^r L_j$. Set $x_j=c_1(L_j)$ and let $\rho:M=\mathbb P((p^*E)^*)\to F$. On $M$, write $h=c_1(\mathcal O_M(1))$. The $j$th homogeneous coordinate is a smooth section $s_j$ of $B_j=\mathcal O_M(1)\otimes\rho^*L_j^{-1}$, and these sections have no common zero. On the open set $U_j=\{s_j\ne0\}$, the section trivializes $B_j$. Naturality of the line Chern class gives $(h-\rho^*x_j)|_{U_j}=0$.
The long exact sequence of the pair $(M,U_j)$ therefore lifts $h-\rho^*x_j$ to $H^2(M,U_j;\Z)$. Their relative cup product lies in
\[
 H^{2r}\!\left(M,\bigcup_jU_j;\Z\right)=H^{2r}(M,M;\Z)=0,
\]
since the $U_j$ cover $M$. Its image in absolute cohomology is $\prod_j(h-\rho^*x_j)$, so this product vanishes. These relative-cohomology facts are part of the topological input in Appendix~\ref{dep:good-cover}. Comparing with the unique projective-bundle relation identifies $p^*c_j(E)$ with the $j$th elementary symmetric function of the $x_j$.
The direct-sum formula for Chern--Weil classes gives the same expression in real cohomology. Injectivity of pullback then proves equality on the base.
\end{proof}

\paragraph{Bott's topological input.}
Bott periodicity and its integrality calculation say that, in the stable range $0\le j<2r$, $\pi_j(U(r))$ is $\Z$ for odd $j$ and zero for even $j$. For $1\le s\le r$, on $\pi_{2s}(BU(r))\simeq\Z$, the class $c_s$ evaluates as $\pm(s-1)!$ times the generator. This factorial is the source of the torsion hypothesis below and cannot be omitted; the exact external input and sources are in Appendix~\ref{dep:bott}.

\begin{theorem}[Peterson's triviality criterion]\label{idx:peterson}
Let $K$ be a CW complex of dimension at most $2r$, and let $E$ have complex rank $r$. Suppose that, for every $s\in\{1,\ldots,r\}$, $H^{2s}(K,\Z)$ has no nonzero element annihilated by $(s-1)!$. Then $E$ is trivial if and only if all its integral Chern classes vanish.
\end{theorem}
\begin{proof}
A trivial bundle has a constant classifying map and all positive-degree Chern classes vanish.
Conversely, let $f:K\to BU(r)$ classify $E$. The classifying construction uses local trivializations and a partition of unity to embed the bundle in a universal trivial bundle and sends each fiber to its image in a Grassmannian. Enlarging the ambient space and rotating embeddings shows that the resulting homotopy class is independent of the choices.
We deform $f$ to a constant map skeleton by skeleton. In the relevant range, odd homotopy groups of $BU(r)$ vanish and its even group in degree $2s$ is $\Z$. Once the map on the $(2s-1)$-skeleton has been made constant, the obstruction to extending the null-homotopy across the $2s$-skeleton is
$o_s\in H^{2s}(K,\Z)$. On a $2s$-cell, its value is the homotopy class of the sphere map obtained by collapsing the cell boundary.
The adjustment on a $(2s-1)$-cell is an adjustment of a homotopy, which has one extra parameter. The group involved is therefore $\pi_{2s}(BU(r))$, not the vanishing group $\pi_{2s-1}$. Changing the preceding homotopy changes the obstruction integers by a cellular coboundary. If the obstruction class vanishes, first make the opposite integer adjustments on the $(2s-1)$-cells and then fill each $2s$-cell.
Bott integrality gives, at this stage,
$c_s(E)=\pm(s-1)!\,o_s$: the map is constant on the lower skeleton, so the Chern class is computed on these sphere pieces.
Since $c_s(E)=0$ and multiplication by $(s-1)!$ is injective on the stated cohomology group, $o_s=0$. Continuing over the skeleta stays within the stable range because $\dim K\le2r$. The final map is null-homotopic and its classified bundle is trivial.
This completes the skeletal induction.
\end{proof}


\chapter{Numerical Positivity, Section Growth, and Big Line Bundles}\label{gp:chapter}
Let $X$ be a connected compact complex manifold of dimension $n$, let $\mathcal M(X)$ be its field of meromorphic functions, and set $a(X)=\operatorname{trdeg}_{\C}\mathcal M(X)$. Transcendence degree is the maximum size of an algebraically independent family. For a holomorphic line bundle $L$ on $X$, set $\kappa(L)=-\infty$ if $H^0(X,L^k)=0$ for every $k\ge1$; otherwise $\kappa(L)$ is the transcendence degree of the field generated by ratios of two nonzero sections of the same positive tensor power. The Kodaira dimension of $X$ is $\kappa(X)=\kappa(K_X)$.
\begin{learning}
Positive curvature produces sections, but the number of sections also records how much of a manifold a linear system can see. We compare the different positivity conditions and estimate section spaces by counting their possible Taylor coefficients. The distinction between powers of a line bundle and powers with a fixed twist will be important.
\end{learning}

\section{Eigenvalues, local frames, and positivity}
\begin{proposition}[The scope of simultaneous diagonalization]\label{gp:eigenvalues}
Given a positive Hermitian form $\omega$ and a Hermitian $(1,1)$-form $\theta$, at each point there is an $\omega$-unitary basis in which $\theta$ is diagonal. Its eigenvalues are independent of this basis.
Smooth unitary frames exist locally, but in general no smooth frame diagonalizes $\theta$ throughout a neighborhood. A superscript $*$ below is conjugate transpose; $e_j$ in the local form expression denotes a $(1,0)$ coframe.
\end{proposition}
\begin{proof}
In any basis let the matrices be $G>0$ and $A=A^*$. Apply the finite-dimensional spectral theorem to the Hermitian matrix $G^{-1/2}AG^{-1/2}$, then map its orthonormal eigenbasis back by $G^{-1/2}$. The resulting basis is $G$-unitary and diagonalizes $A$.
A change of unitary basis conjugates the matrix unitarily, preserving its characteristic polynomial and eigenvalues.
To construct a smooth unitary frame, apply Gram--Schmidt to an arbitrary smooth local frame. Every denominator is the square root of a positive smooth function, so the result is smooth.
Pointwise eigenvectors need not vary smoothly: the real-parameter family
$\left(\begin{smallmatrix}x&y\\y&-x\end{smallmatrix}\right)$ has eigenvalues $\pm\sqrt{x^2+y^2}$, which cannot even be labeled smoothly near the origin. The valid local expression is therefore $\theta=\ii\sum a_{j\bar k}e_j\wedge\bar e_k$, without assuming that the coefficient matrix is smoothly diagonal.
\end{proof}

Throughout the chapter $X$ is compact and connected, with a fixed Hermitian form $\omega$.
A line bundle $L$ is \emph{nef} if for every $\varepsilon>0$ it has a smooth metric with $\ii F_h\ge-\varepsilon\omega$; it is \emph{pseudoeffective} if it has a singular metric with nonnegative curvature current; and it is \emph{analytically big} if such a metric satisfies $\ii F_h\ge\varepsilon\omega$ for some $\varepsilon>0$.
It is \emph{semiample} if a positive tensor power is globally generated.
These definitions do not depend on the background metric, since any two Hermitian metrics on a compact manifold are uniformly comparable.

A \emph{quasi-plurisubharmonic} (quasi-psh) function is locally the sum of a psh function and a smooth real function, with its upper-semicontinuous representative understood; it is not identically $-\infty$ on a component.
\begin{lemma}[Compactness of normalized quasi-plurisubharmonic functions]\label{gp:qpsh-compactness}
On a compact connected $X$, the family of quasi-psh functions $u_j$ satisfying
$\ii\partial\dbar u_j\ge-C\omega$ and $\sup_Xu_j=0$, for a fixed $C>0$, is relatively compact in $L^1(X)$. Every limit satisfies the same current inequality.
\end{lemma}
\begin{proof}
On finitely many coordinate balls, add a sufficiently large fixed multiple $C'|z|^2$ so that $u_j+C'|z|^2$ is plurisubharmonic. These functions have a common local upper bound because $u_j\le0$.
We must also prevent convergence to $-\infty$. Choose $p_j$ with $u_j(p_j)\ge-1$ and pass to a subsequence with $p_j\to p$. In a fixed chart the submean inequality gives a common lower bound for the integrals over moving balls $B(p_j,r)$ of a fixed radius. For large $j$ they contain the fixed ball $B(p,r/2)$. Subtract the common upper bound for the integral over the complement to obtain a lower bound on that fixed ball. Here and below we use
$\int_U v\ge\int_Bv-C'\operatorname{vol}(B\setminus U)$ whenever $U\subset B$ and $v\le C'$.

Choose a finite cover by small coordinate balls with sufficiently large outer charts. A lower integral bound on one ball gives a lower bound on a fixed positive-volume overlap $W$ with a neighboring ball, by the same complement inequality. Hence there is a point of $W$ where the corresponding function is bounded below uniformly in $j$. Apply the submean inequality at that point, with a radius large enough to contain the neighboring small ball and small enough to remain in its outer chart. Subtracting the complementary integral gives a lower bound on that neighboring ball. The quadratic corrections and volume densities have fixed bounds on overlaps, so the estimates pass between charts. Finite connected chains propagate the bounds throughout the cover. Thus $\|u_j\|_{L^1}$ is uniformly bounded. This also rules out unbounded $L^1$ norms in the whole family: apply the same argument to any subsequence whose norms tend to infinity.

We spell out local $L^1$ compactness. On a slightly larger ball write $v_j=u_j+C'|z|^2$. In this paragraph
$\Delta_{\rm eucl}=\sum_{a=1}^{2n}\partial_{x_a}^2$, with the sign opposite to the positive spectral Laplacian. Plurisubharmonicity makes $\Delta_{\rm eucl}v_j$ a positive measure. For a nonnegative cutoff $\chi$ equal to one on a smaller ball, distributional integration by parts gives
$\int\chi\,\Delta_{\rm eucl}v_j=\int v_j\Delta_{\rm eucl}\chi$.
The $L^1$ bound therefore bounds the masses on the smaller ball.
Cut off these measures and extend them by zero. Their convolutions with a fundamental solution of the Laplacian are local potentials, and the remainders are harmonic.
Multiply the fundamental solution by a fixed cutoff that is one on all differences of points needed here. This leaves the local potentials unchanged and makes the kernel compactly supported and integrable.
Translations of this kernel are continuous in $L^1$. The convolutions consequently have a common modulus of $L^1$ translation continuity and a common compact support. Approximation by averages on a fine finite grid proves that they form a totally bounded subset of $L^1$.
The harmonic remainders are smooth by Lemma~\ref{cur:weyl} and have bounded $L^1$ norms; the mean-value formula then bounds all their derivatives on smaller balls, giving compactness there as well.
A diagonal subsequence over finitely many charts gives a global $L^1$ limit. Pairing with positive test forms and passing to the limit preserves the curvature lower bound. On each chart, add the same smooth quadratic correction to the limit. Theorem~\ref{cur:distribution-psh} gives its unique psh representative. Subtracting the correction yields a quasi-psh representative; uniqueness for distributions makes these representatives agree on overlaps.
\end{proof}

\begin{proposition}[Direct implications between positivity conditions]\label{gp:positivity-implications}
An ample line bundle is semiample and big. Semiampleness implies smooth semipositivity, which implies nefness, which implies pseudoeffectivity. Bigness also implies pseudoeffectivity.
If $A$ is ample and $N$ is nef, then $A\otimes N$ is ample.
\end{proposition}
\begin{proof}
A positive power of an ample bundle is very ample, giving semiampleness; a positive-curvature metric gives bigness.
If $L^m$ is generated by sections $s_j=f_je^{\otimes m}$, then
$\varphi=m^{-1}\log\sum_j|f_j|^2$ is a smooth weight with semipositive Levi form. On overlaps it transforms as a line-bundle weight. Thus $L$ has a smooth semipositive metric and is nef.
For nef $L$, fix a smooth metric $h_0$ and choose $h_j=h_0e^{-u_j}$ with
$\ii F_{h_0}+\ii\partial\dbar u_j\ge-j^{-1}\omega$.
Adding constants, arrange $\sup u_j=0$. Lemma~\ref{gp:qpsh-compactness} gives a subsequence converging in $L^1$ to $u$, and passage to currents gives
$\ii F_{h_0}+\ii\partial\dbar u\ge0$. Use the compatible quasi-psh representative supplied by that lemma and set $h=h_0e^{-u}$. Its local weights have nonnegative distributional Levi forms, so Theorem~\ref{cur:distribution-psh} makes them psh. Thus $h$ is a pseudoeffective metric.
The use of a global supremum is essential: normalization on just the first compact set of an exhaustion would not control the later sets.
On a disconnected compact manifold, normalize separately on each of its finitely many connected components before taking subsequences.
Bigness implies pseudoeffectivity directly from the curvature inequality.
Finally choose $\ii F_A\ge c\omega$ and a metric on $N$ with $\ii F_N\ge-c\omega/2$. Curvatures add under tensor product, so $\ii F_{A\otimes N}\ge c\omega/2$. Theorem~\ref{emb:kodaira} makes $A\otimes N$ ample.
\end{proof}

\section{Vanishing orders and Siegel's dimension bound}
\begin{lemma}[Three-circle estimates and propagation along finite chains]\label{gp:three-circle}
For a holomorphic function on a ball, $\log M(r)$, where $M(r)=\sup_{|z|\le r}|f(z)|$, is convex as a function of $\log r$.
Fix smooth metrics on $L$ and on a twisting bundle $F$, and write $\|s\|_{\infty,X}=\sup_X|s|$ and $\operatorname{ord}_p(s)$ for the least degree of a nonzero Taylor component in a local frame. For a fixed $L$ on a compact connected complex manifold, there are $a,b>0$ such that a section $s\in H^0(X,L^k\otimes F)$ vanishing to order $m$ at a fixed center satisfies
\[
 \|s\|_{\infty,X}\le e^{ak+b_F}e^{-bm}\|s\|_{\infty,X}.
\]
Here $F$ is any fixed vector bundle; $a,b$ are independent of $F$, whereas $b_F$ may depend on $F$.
\end{lemma}
\begin{proof}
Restrict $f$ to the disk in each complex direction and apply the one-variable Hadamard three-circle theorem. Its proof subtracts from $\log|f|$ the function $A\log|z|+B$ interpolating the two boundary bounds, then applies the subharmonic maximum principle on the annulus. Taking the supremum over directions gives the same estimate for balls in several variables.
If the Taylor expansion begins in degree $m$, the Cauchy coefficient estimate on concentric balls of radii $r<R$ gives
$M(r)\le C\theta^m M(R)$. Fixing $r/R$ sufficiently small absorbs the polynomial number of multi-indices into a slightly larger exponential factor $\theta<1$.

Choose finitely many outer coordinate balls, whose inner balls cover $X$, and holomorphic frames of $L$ on the outer balls. This cover must be independent of $F$.
On a fixed outer ball $B$, holomorphic sections of $F^*$ generate every fiber at an interior point $p$. To verify this, resolve the ideal of $p$ by the Koszul complex of the coordinate functions $z_i-p_i$ (Lemma~\ref{alg:koszul}) and tensor with $F^*$. Every term is a finite sum of copies of $F^*$. The vector-bundle Cartan B theorem, Theorem~\ref{stv:cartan-b}, makes their higher cohomology vanish. Splitting this finite resolution into short exact sequences yields $H^1(B,\mathcal I_pF^*)=0$, so evaluation $H^0(B,F^*)\to F_p^*$ is onto.
This forward reference is not circular: that theorem is proved from the earlier weighted estimates and does not use the section-growth results of this chapter.
On a fixed intermediate closed ball, compactness now gives finitely many holomorphic sections $\tau_1,\ldots,\tau_N$ generating every fiber. Consequently
\[
 c_F|v|\le\max_\ell|\tau_\ell(v)|\le C_F|v|,
 \qquad v\in F,
\]
with $0<c_F\le C_F<\infty$. We apply the three-circle estimate only to the scalar holomorphic functions $\tau_\ell(s/e_L^{\otimes k})$; no holomorphic trivialization of $F$ on a smaller cover is needed.

Join the inner balls by finite overlapping chains, each overlap containing a fixed smaller ball.
Changes of the fixed frames of $L$, and comparison with their smooth norms, cost at most $M^k$. The preceding norm comparisons cost constants depending on $F$ but not on $k$.
Transfer the initial exponential smallness to an overlap and then use three-circle interpolation to reach the next ball. At each step the bound has the form
$M_{j+1}\le e^{Ck+C_F}M_j^\gamma M_{\mathrm{global}}^{1-\gamma}$,
where $0<\gamma<1$ depends only on the fixed radii.
After finitely many steps, the factor $\theta^m$ becomes $\theta^{\gamma_1\cdots\gamma_Nm}$ and the comparison factors become $e^{C'k+C'_F}$. Chain lengths are bounded, so the products of interpolation exponents have a positive lower bound. Reach a ball containing a point of maximum global norm and include the remaining fixed metric comparisons. This proves the stated inequality.
The maximum of finitely many $|\tau_\ell(s)|$ obeys the same three-circle estimate; their number and norms affect only $b_F$.
\end{proof}

\begin{theorem}[A uniform vanishing-order bound]\label{gp:siegel-bound}
Fix $X,L$ and $p\in X$. There is $A>0$ such that for every fixed rank-$r$ vector bundle $F$ there is $B_F$ with
$\operatorname{ord}_p(s)\le Ak+B_F$ for each nonzero $s\in H^0(X,L^k\otimes F)$.
In particular,
\[
 h^0(X,L^k\otimes F)\le r\binom{n+\lfloor Ak+B_F\rfloor}{n}=O_F(k^n).
\]
\end{theorem}
\begin{proof}
Divide the inequality in Lemma~\ref{gp:three-circle} by the positive global norm of a nonzero section. Then
$1\le e^{ak+b_F-bm}$, hence $m\le(a/b)k+b_F/b$. Take $A=a/b$ and increase $B_F$ if necessary to make it nonnegative.
A section with zero jet through order $\lfloor Ak+B_F\rfloor$ vanishes to a larger order and must be zero. The jet map is therefore injective.
A local jet has $r$ components, each a polynomial of total degree at most $\lfloor Ak+B_F\rfloor$. Counting its monomials gives the displayed binomial coefficient. The argument identifies precisely which constants may depend on the twisting bundle.
\end{proof}

\begin{theorem}[Siegel's bound on algebraic dimension]\label{gp:algebraic-dimension}
The meromorphic function field $\mathcal M(X)$ of a compact connected complex $n$-fold has transcendence degree $a(X)\le n$.
For every line bundle, $\kappa(L)\le a(X)$.
\end{theorem}
\begin{proof}
Suppose $f_1,\ldots,f_d$ are algebraically independent meromorphic functions. Choose an effective divisor $D$ containing all their poles, let $L=\mathcal O(D)$, and let $s_0$ be its canonical section. Then $s_i=f_is_0$ are holomorphic.
The polar divisor is well defined because local rings of smooth complex manifolds are unique factorization domains: cancel common irreducible factors in a local quotient $g/h$ and record the multiplicities of the remaining denominator. These orders agree on overlaps and define a Cartier divisor. Here is the factoriality argument. Induct on $n$, starting with the field $R_0=\C$. For a nonzero nonunit $f\in R_n$, choose regular coordinates and write $f=uW$ by Weierstrass preparation. By induction and Gauss's lemma, $R_{n-1}[w]$ is a unique factorization domain. Since $W$ is monic, its irreducible factors can be normalized to be monic: their leading coefficients multiply to one and hence are units. Specializing $z'=0$, their product is $w^d$, so each monic factor specializes to a power of $w$ and is itself distinguished.
An irreducible polynomial factor $P$ remains irreducible in $R_n$. Indeed, if $P=gh$ with both germs nonunits, then $P(0,w)=w^e$ implies that $g(0,w)$ and $h(0,w)$ are nonzero and have positive orders. Preparation gives $g=u_gG$ and $h=u_hH$ with distinguished monic $G,H$. Uniqueness of preparation forces $P=GH$, contradicting polynomial irreducibility. Conversely, any analytic factorization of $f$ prepares in these same coordinates to a polynomial factorization of $W$, so polynomial uniqueness gives analytic uniqueness up to units. Thus $R_n$ is factorial.
Compactness leaves only finitely many polar components.
All monomials in the $f_i$ of total degree at most $k$, multiplied by $s_0^k$, are linearly independent sections of $L^k$: a linear dependence, divided by $s_0^k$, would be a polynomial relation among the $f_i$. Thus
$\binom{k+d}{d}\le h^0(X,L^k)\le Ck^n$.
Letting $k\to\infty$ gives $d\le n$.

Recall the definition of $\kappa(L)$ at the start of the chapter.
This agrees with the maximum dimension of the images of Kodaira maps: finitely many ratios can be placed in a common degree, giving one meromorphic projective map. Resolve its graph; the proper mapping theorem makes the image a compact analytic set, and Chow's theorem makes it projective algebraic. Its function field has transcendence degree equal to its dimension. These geometric comparison inputs are recorded in Appendix~\ref{dep:chapter}.
The ratio field is a subfield of $\mathcal M(X)$, so $\kappa(L)\le a(X)$.
\end{proof}

\begin{theorem}[Iitaka growth for an untwisted line bundle]\label{gp:iitaka-growth}
If $\kappa(L)=d\ge0$, then $h^0(X,L^k)=O(k^d)$. There are a positive integer $m$ and $c>0$ such that $h^0(X,L^{m\ell})\ge c\ell^d$ for all sufficiently large $\ell$.
\end{theorem}
\begin{proof}
Fix local coordinates and a frame at one point. Order Taylor monomials first by total degree and then lexicographically, and assign to a nonzero section the first exponent $\alpha\in\N_0^n$ with nonzero coefficient. The leading term of a product is the product of its leading terms, so this assignment is additive under multiplication. Two sections with the same leading term can be combined to cancel that term.
Theorem~\ref{gp:siegel-bound} gives $|\alpha|\le Ak+B$. Gaussian elimination in the finite-dimensional section space therefore produces a basis with distinct leading exponents. Its dimension equals the number of exponents that occur.

Record the values $(k,\alpha)\in\Z^{n+1}$ of the section ring $R(L)=\bigoplus_{k\ge0}H^0(X,L^k)$ and let $\Gamma$ be the subgroup they generate. Products of nonzero sections are nonzero on connected $X$, so $R(L)$ is an integral domain; $\operatorname{Frac}$ below denotes its fraction field. If $s$ values in this group are linearly independent, choose corresponding homogeneous rational sections, using products and quotients of sections as necessary. Distinct monomials in them have distinct values, hence cannot satisfy a polynomial relation. Therefore
$s\le\operatorname{trdeg}\operatorname{Frac}(\bigoplus_kH^0(L^k))=d+1$.
The equality does not require finite generation of the section ring. Indeed, let $K_0$ be the field of equal-degree ratios and fix a nonzero homogeneous element $s_0$ of degree $m>0$. If $s$ has degree $k$, then $s^m/s_0^k\in K_0$, so $s$ is algebraic over $K_0(s_0)$. On the other hand, $s_0$ is transcendental over $K_0$: clear homogeneous denominators in a putative finite polynomial relation; its terms have different degrees and hence must vanish separately in the graded direct sum.
The fraction field is thus algebraic, not necessarily finite, over $K_0(s_0)$ and has transcendence degree $d+1$.
For a fixed degree $k$, the nonempty slice $\Gamma\cap(\{k\}\times\Z^n)$ is a coset of a lattice of rank at most $d$. In the box $|\alpha|\le Ak+B$ it has $O(k^d)$ points: choose a lattice basis and linear coordinates on its real span; each coordinate ranges over an interval of length $O(k)$, uniformly in the translate. This proves the upper bound.

Choose $d$ algebraically independent ratios and clear denominators and degrees to write them as $s_i/s_0$ with all sections in one degree $m$.
The degree-$\ell$ homogeneous monomials in $s_0,s_1,\ldots,s_d$ are linearly independent, since a relation would give a polynomial relation among the ratios. Their number is $\binom{\ell+d}{d}$, proving the lower bound.
\end{proof}

\begin{proposition}[Kodaira dimension of a product]\label{gp:kodaira-product}
For compact connected complex manifolds $X,Y$ with nonnegative Kodaira dimensions,
$\kappa(X\times Y)=\kappa(X)+\kappa(Y)$.
If either dimension is $-\infty$, so is that of the product.
\end{proposition}
\begin{proof}
We have $K_{X\times Y}=\pi_X^*K_X\otimes\pi_Y^*K_Y$ and
\[
 H^0(X\times Y,K_{X\times Y}^m)
 =H^0(X,K_X^m)\otimes H^0(Y,K_Y^m).
\]
For the second equality, choose a basis of sections on $X$ and finitely many evaluation points giving a full-rank evaluation matrix, after choosing local fiber frames. Such points can be chosen successively: otherwise a nonzero linear combination would vanish at every point. Restrict a section on the product to these points and solve the fixed linear system. Its coefficients are holomorphic sections on $Y$. Substitution gives the desired tensor expansion on every $X$ fiber.
Dimensions consequently multiply. The upper bounds of Theorem~\ref{gp:iitaka-growth} give an exponent no larger than the sum, and the lower bounds on a common multiple of the two degrees give the reverse inequality. If all positive powers on one factor have zero section space, the tensor products are all zero.
\end{proof}

\begin{example}[A fixed twist can change the Iitaka growth exponent]\label{gp:twist-counterexample}
The untwisted growth estimate need not persist after tensoring by a fixed vector bundle. In particular, an estimate of the form
\[
 h^0(L^k\otimes F)\le C(X,L)\operatorname{rank}(F)k^{\kappa(L)}
                         +o(k^{\kappa(L)})
\]
cannot hold for all fixed twists with a constant independent of $F$. To see this, take $X=\mathbb P^1$, $L=\mathcal O$, and $F=\mathcal O(d)$. Then $\kappa(L)=0$, but the rank-one bundle $F$ has $d+1$ independent sections.
There is also a counterexample to the exponent itself. Let $C=\C/(\Z+\tau\Z)$ with $\operatorname{Im}\tau>0$, and let $V$ be the rank-two flat bundle whose monodromy under $1$ is the identity and under $\tau$ is
$\left(\begin{smallmatrix}1&1\\0&1\end{smallmatrix}\right)$.
Put $X=\mathbb P(V^*)$ and $L=\mathcal O_X(1)$.
The holomorphic sections of $\operatorname{Sym}^kV$ form a one-dimensional space. Indeed, expand a section in Fourier modes for its $1$-periodicity. At every nonzero frequency, the factor $e^{2\pi\ii m\tau}$ differs from the sole monodromy eigenvalue $1$, so that Fourier coefficient is zero. The invariant subspace at frequency zero is one-dimensional.
Explicitly, write the monodromy matrix as $e^N$ with $Nx=0,Ny=x$. On degree-$k$ polynomials its induced nilpotent operator satisfies
$N_k(x^{k-j}y^j)=j x^{k-j+1}y^{j-1}$.
This is one Jordan chain, so $\ker(e^{N_k}-1)=\ker N_k=\C x^k$.
Theorem~\ref{sh:projective-direct-image} gives $h^0(X,L^k)=1$, hence $\kappa(L)=0$.
If $B$ has positive degree $d$ on $C$, then $\operatorname{Sym}^kV\otimes B$ has a filtration with $k+1$ quotients isomorphic to $B$. Since $K_C$ is trivial and negative-degree line bundles have no sections, Theorem~\ref{hod:serre-duality} gives $H^1(C,B)=0$, and Theorem~\ref{cur:rr} gives $h^0(C,B)=d$.
The successive long exact sequences therefore give
$h^0(X,L^k\otimes\pi^*B)=(k+1)d$.
Thus even an $O(k^0)$ bound with a constant depending on $F$ is false.
Theorem~\ref{gp:siegel-bound} gives the uniform degree-$n$ bound for fixed twists, whereas Theorem~\ref{gp:iitaka-growth} describes untwisted powers.
\end{example}

\section{Bigness and Kodaira's criterion}
Call $L$ \emph{section-big} if $\limsup_{k\to\infty}k^{-n}h^0(X,L^k)>0$.
Theorem~\ref{gp:iitaka-growth} makes this equivalent to $\kappa(L)=n$.

\begin{theorem}[Kodaira's bigness criterion on a projective manifold]\label{gp:kodaira-big}
On a projective manifold the following are equivalent: $L$ is section-big; $mL=A+D$ for some $m>0$, ample $A$, and effective divisor $D$; and $L$ is analytically big.
\end{theorem}
\begin{proof}
The zero-dimensional case is immediate, so assume $n\ge1$.
Suppose first that $L$ is section-big. Take a very ample bundle $A$ and its embedding $X\hookrightarrow\mathbb P^N$.
A smooth hyperplane divisor can be chosen by the following incidence argument. A hyperplane section is singular at $x$ exactly when the hyperplane contains the projective tangent $n$-plane to $X$ at $x$. If $N>n$, the hyperplanes satisfying this condition at a fixed $x$ form $\mathbb P^{N-n-1}$. Their incidence space over $X$ is a compact smooth space of complex dimension $n+(N-n-1)=N-1$. Its projection to the dual projective space cannot cover that space: on finitely many compact coordinate patches the projection is Lipschitz, so its image has Hausdorff dimension at most $2N-2$, strictly less than $2N$. This uses the elementary Lipschitz dimension bound established in the discussion preceding Proposition~\ref{bas:analytic-zeros}. If $N=n$, no hyperplane contains a projective tangent $n$-plane, so the bad incidence is empty. Choose a hyperplane outside the bad image and call its smooth divisor $H$.
For $n=1$, the transverse section $H$ is compact and discrete, hence a finite set of points; therefore $h^0(H,L^k|_H)$ is its cardinality and is $O(1)$.
For $n\ge2$, apply the $(n-1)$-dimensional version of Theorem~\ref{gp:siegel-bound} to the finitely many connected components of $H$.
In either case,
$h^0(H,L^k|_H)=O(k^{n-1})$.
The exact sequence
$0\to L^k\otimes A^{-1}\to L^k\to L^k|_H\to0$
gives
$h^0(L^k\otimes A^{-1})\ge h^0(L^k)-h^0(H,L^k|_H)$.
Along the sequence witnessing bigness, the right side is positive. A nonzero section for some $m$ defines an effective divisor $D$ with $mL=A+D$.

Given this decomposition, multiply a positive metric on $A$ by the canonical singular metric of $D$. Theorem~\ref{cur:poincare-lelong} gives
$\ii F_{mL}=\ii F_A+2\pi[D]\ge\varepsilon\omega$.
Taking an $m$th root proves analytic bigness.
Multiplication by the $k$th power of the divisor section also gives an injection $H^0(A^k)\hookrightarrow H^0(L^{mk})$.
A very ample power of $A$ has $n$ locally independent coordinate ratios; their homogeneous monomials give at least $ck^n$ independent sections along a fixed sequence of powers. Hence $L$ is section-big.

It remains to deduce the decomposition from analytic bigness. Let $h=h_0e^{-\varphi}$ with $\ii F_h\ge\varepsilon\omega$ and fix a very ample $A$.
First fix a finite collection of inner and outer coordinate balls. For each $p$ in an inner ball, a cutoff on the outer ball, at a \emph{fixed radius}, produces a global quasi-plurisubharmonic function $\psi_p$ equal to $n\log|z-z(p)|^2$ near $p$ and satisfying
$\ii\partial\dbar\psi_p\ge-C\omega$, uniformly in $p$.
Only the fixed annulus away from the pole contributes a negative part; differentiating the cutoff formula gives the uniform bound.
Now choose $m$ so large that $m\varepsilon\omega$ absorbs this loss and the fixed negative curvature contributions of $A^{-1}$ and $K_X^{-1}$, leaving a strictly positive lower bound.

The Lelong number of $\varphi$ is zero almost everywhere. Indeed, at almost every point the trace measure of the local positive current has finite volume density, so its mass on a ball is $O(r^{2n})$; division by $r^{2n-2}$ tends to zero.
Skoda integrability, Theorem~\ref{int:skoda-integrability}, gives local integrability of $e^{-m\varphi}$ near almost every such point $p$.
Choose a sufficiently small neighborhood, a local frame of $L^m\otimes A^{-1}$, and a cutoff $\chi$ equal to one near $p$. Extend the resulting smooth local section by zero to a global section $\sigma$.
Its support may shrink with $m,p$, while the already chosen $\psi_p$ and its curvature lower bound remain fixed. Thus no new, uncontrolled curvature loss appears after $m$ has been chosen.
Add $\psi_p$ to the weight of $L^m\otimes A^{-1}$. The form $\dbar\sigma$ is supported away from $p$ and has finite weighted norm.
Rewrite the $(0,1)$ equation as an $(n,1)$ equation with the additional $K_X^{-1}$ twist. The large choice of $m$ has already absorbed its smooth curvature. Theorem~\ref{nad:global-l2}, applied to this strictly positive singular metric, gives $\dbar u=\dbar\sigma$ with
$\int|u|^2e^{-m\varphi}|z-p|^{-2n}<\infty$.
Near $p$, $u$ is holomorphic. Since $\varphi$ is locally bounded above, the integral forces $u(p)=0$: otherwise its radial part contains the divergent integral $\int_0^\delta dr/r$.
Thus $s=\sigma-u$ is a nonzero holomorphic section of $L^m\otimes A^{-1}$ with $s(p)=1$. Its divisor gives $mL=A+D$.
This direction uses both the singular-weight $L^2$ theorem and Skoda integrability, not just their smooth-weight models.
\end{proof}

\begin{proposition}[Moishezon and Morse criteria for bigness]\label{gp:morse-big}
A compact connected complex $n$-fold has a line bundle with $\kappa(L)=n$ if and only if $a(X)=n$.
Write $X(L,\le1)$ for the locus where the curvature is nondegenerate and has at most one negative eigenvalue, as in Theorem~\ref{morse:main}. If a smooth metric satisfies
$\int_{X(L,\le1)}c_1(L,h)^n>0$, then $\kappa(L)=n$ and $a(X)=n$.
\end{proposition}
\begin{proof}
The inequalities in Theorem~\ref{gp:algebraic-dimension}, $\kappa(L)\le a(X)\le n$, prove the first direction.
Conversely, if $a(X)=n$, use a common polar divisor of $n$ algebraically independent meromorphic functions to construct $L$, as in Theorem~\ref{gp:algebraic-dimension}. Their monomials give
$h^0(L^k)\ge\binom{k+n}{n}$, so $\kappa(L)=n$.
For the final assertion apply the strong holomorphic Morse inequality, already proved in Theorem~\ref{morse:main}, with $q=1$:
\[
 h^0(X,L^k)-h^1(X,L^k)
 \ge\frac{k^n}{n!}\int_{X(L,\le1)}c_1(L,h)^n+o(k^n).
\]
Since $h^1\ge0$, the positive leading coefficient gives $h^0\ge ck^n$. Theorem~\ref{gp:iitaka-growth} then forces $\kappa(L)=n$.
This is the Grauert--Riemenschneider bigness criterion obtained from the Morse inequalities.
\end{proof}

\begin{proposition}[The numerical dimension of a nef line bundle]\label{gp:numerical-dimension}
On a compact K\"ahler $n$-fold with K\"ahler form $\omega$, the numerical dimension of a nef line bundle has the equivalent descriptions
\begin{align*}
 \nu(L)&=\max\{j\in\{0,\ldots,n\}:c_1(L)^j\ne0\text{ in }H^{2j}(X,\R)\}\\
       &=\max\left\{j\in\{0,\ldots,n\}:\int_Xc_1(L)^j\wedge\omega^{n-j}>0\right\}.
\end{align*}
Here the Chern class is taken in real cohomology, so integral torsion does not contribute to numerical dimension.
Moreover $\kappa(L)\le\nu(L)$.
\end{proposition}
\begin{proof}
Choose representatives $\alpha_\varepsilon\ge-\varepsilon\omega$ of $c_1(L)$ and put $\beta_\varepsilon=\alpha_\varepsilon+\varepsilon\omega\ge0$.
The form $\beta_\varepsilon^j$ is positive and closed, with mass $\int\beta_\varepsilon^j\wedge\omega^{n-j}$.
If its limiting mass is zero, the coefficient trace bound of Theorem~\ref{pos:order-zero} gives $\beta_\varepsilon^j\to0$ as currents. Pairing with closed test forms then gives $c_1(L)^j=0$.
Conversely that class being zero makes its intersection with $\omega^{n-j}$ zero. Together with nonnegativity, these statements prove equality of the two definitions.

If $\kappa(L)=-\infty$ or $\kappa(L)=0$, the inequality follows from $\nu(L)\ge0$.
If $\kappa(L)=d\ge1$, choose a linear system whose image has dimension $d$. Resolution of its base ideal (Appendix~\ref{dep:resolution}) gives a sequence of blowups in smooth centers
$\mu:\widetilde X\to X$ with $\mu^*(mL)=M+D$, where $M$ is globally generated and $D$ is effective.
The blowup remains K\"ahler: adding a sufficiently small closed curvature form, positive in the vertical directions of the exceptional divisor, to the pulled-back background form fills its vertical degeneracy while preserving horizontal positivity. Local positivity and compactness give a global K\"ahler form $\widetilde\omega$.
The semipositive curvature of $M$ has rank $d$ on an open set, so
$\int c_1(M)^d\widetilde\omega^{n-d}>0$.
Set $\alpha=c_1(\mu^*(mL))$ and $\beta=c_1(M)$; both classes are nef.
Pair the identity
\[
 \alpha^d-\beta^d=[D]\sum_{j=0}^{d-1}\alpha^{d-1-j}\beta^j
\]
with $\widetilde\omega^{n-d}$.
Every term on the right is nonnegative: approximate the nef classes by forms bounded below by $-\varepsilon\widetilde\omega$, add $\varepsilon\widetilde\omega$, integrate on the regular locus of the effective divisor, and pass to the cohomological limit.
It follows that $\alpha^d\ne0$, hence $c_1(L)^d\ne0$, and therefore $d\le\nu(L)$.
The case $\kappa=-\infty$ is automatic.
The use of Hironaka's resolution theorem here is substantive; local Noetherianity alone does not supply it.
\end{proof}


\chapter{Global Applications of Nadel Vanishing}\label{nad:applications}
\begin{learning}
Nadel vanishing turns a curvature estimate into a statement about global sections. We begin with a fixed multiplier ideal, then relate numerical positivity to vanishing and generation. The examples show why the metric, the ideal, and the tensor power must each be specified.
\end{learning}
\section{High powers and a fixed multiplier ideal}
\begin{corollary}[High powers of a smooth positive line bundle with a fixed multiplier ideal]\label{nad:high-power}
Let $L$ be a smooth positive line bundle on a compact complex manifold, and let $\varphi$ be a global quasi-psh function.
Then $H^q(X,L^k\otimes\cI(\varphi))=0$ for every $q\ge1$ and all sufficiently large $k$.
\end{corollary}
\begin{proof}
The positive curvature of $L$ itself supplies a K\"ahler metric.
On the compact manifold, $\ii\partial\dbar\varphi\ge-C\omega$,
and the curvature of a smooth metric on $K_X^{-1}$ has a lower bound $-C'\omega$.
If $\ii\Theta_L\ge a\omega$, the curvature of
$L^k\otimes K_X^{-1}$ with the additional scalar weight $e^{-\varphi}$ is at least
$(ka-C-C')\omega$. Choose $k>(C+C')/a$ and apply Theorem~\ref{nad:vanishing}.
Smooth factors do not change the multiplier ideal, while $K_X$ and $K_X^{-1}$ cancel, proving the result.
\end{proof}

\begin{example}[Failure of high-power vanishing with a fixed ideal for a singular positive line bundle]
\label{nad:fixed-ideal-counterexample}
Let $X=\operatorname{Bl}_p\mathbf P^2$, let $H$ be the pullback of the line class, and let $E$ be the exceptional curve.
Put $L=H+E$. Since $2L=(2H-E)+3E$ and $2H-E$ is ample,
$L$ admits a strictly positive singular metric with logarithmic coefficient $3/2$ along $E$.
Ampleness can be checked directly. Take $p=[1:0:0]$; conics through $p$ give the map
$[x_0x_1:x_0x_2:x_1^2:x_1x_2:x_2^2]$.
In blowup coordinates $x_1=t,x_2=tu$, cancellation of the common factor gives
$[1:u:t:tu:tu^2]$, from which $t,u$ are recovered. The other blowup chart is analogous.
Away from $p$, the coordinate ratios likewise recover points and tangent directions, so the map is a closed embedding.
Thus $2H-E$ is in fact very ample.
Theorem \ref{pos:snc-ideal} gives $\cI(h)=\cO(-E)$.
Hence $L^k\otimes\cI(h)=\cO(kH+(k-1)E)$.

For $k\ge2$, a section restricts to $\cO_E(-(k-1))$, so it vanishes on $E$.
Removing successive positive multiples of $E$ gives
$h^0(kH+(k-1)E)=h^0(kH)=\binom{k+2}{2}$.
The final equality also follows by descending sections to $\mathbf P^2\setminus\{p\}$ and extending by Hartogs; the formula holds directly for $k=1$ as well.
Since $K_X=-3H+E$, Serre duality gives $h^2=0$:
the intersection of $K_X-kH-(k-1)E$ with the nef class $H$ is $-(k+3)$, precluding an effective divisor in that class.
Lines through $p$ exhibit $X\to\mathbf P^1$ as a $\mathbf P^1$ bundle,
identified in local coordinates with $\mathbf P(\cO\oplus\cO(-1))$.
The projective-bundle direct-image calculation \ref{sh:projective-direct-image} and
$H^{>0}(\mathbf P^1,\cO)=0$ therefore give $\chi(\cO_X)=1$.
Surface Riemann--Roch, \ref{idx:hrr}, together with $H^2=1$, $E^2=-1$, and $H\cdot E=0$, gives
\[
 \chi(kH+(k-1)E)
 =1+\tfrac12\bigl(k(k+3)-(k-1)(k-2)\bigr)=3k.
\]
Consequently,
\[
 h^1(kH+(k-1)E)=\binom{k+2}{2}-3k
      =\frac{(k-1)(k-2)}2>0\quad(k\ge3).
\]
Thus strict positivity of a singular metric does not give high-power vanishing with this fixed multiplier ideal; the preceding corollary requires a smooth positive metric on $L$.
\end{example}

\section{Vanishing in the Numerical-Dimension Range for Nef Line Bundles}

\begin{theorem}[Kawamata--Viehweg vanishing for nef line bundles]\label{nad:nef-numerical}
Recall from Proposition~\ref{gp:positivity-implications} that a nef line bundle admits smooth metrics with arbitrarily small negative curvature lower bounds. Let $L$ be nef on a smooth projective $n$-fold, and define
$\nu(L)=\max\{j:\int_Xc_1(L)^j\wedge c_1(A)^{n-j}>0\}$ for a fixed ample line bundle $A$.
This is the numerical dimension of Proposition~\ref{gp:numerical-dimension}; the products in this formula are cohomological intersection numbers.
Then
$H^q(X,K_X\otimes L)=0$ for $q>n-\nu(L)$.
\end{theorem}
\begin{proof}
\emph{First consider $\nu(L)=n$.}
Choose a smooth ample divisor $B$ sufficiently positive that $B-K_X$ is ample.
A nef bundle plus an ample bundle is ample, by Proposition \ref{gp:positivity-implications},
so Corollary~\ref{pos:kodaira}, applied to the ample bundle $kL+B-K_X$, gives $H^q(X,kL+B)=0$ for $q>0$.
Asymptotic Riemann--Roch, \ref{idx:asymptotic-rr}, gives
$h^0(kL+B)=c_1(L)^n k^n/n!+O(k^{n-1})$.
The dimension of the restricted section space on $B$ is $O(k^{n-1})$: choose a sufficiently positive very ample $C$ so that $L+C$ is ample.
Multiplication by the $k$th power of a section of $C$ that is not identically zero on any component of $B$
injects sections of $(kL+B)|_B$ into those of $(k(L+C)+B)|_B$.
The ample Hilbert polynomial gives the asserted bound for the latter.
The restriction exact sequence yields $h^0(kL)\ge c k^n$, proving that $L$ is big.

By Kodaira's bigness criterion, \ref{gp:kodaira-big}, write $mL=A_0+D$,
$A_0$ ample, $D$ effective.
This gives a singular metric $h_b$ on $L$ with curvature at least $c\omega$,
whose singular weight is $m^{-1}\log|s_D|^2$ plus a smooth function.
Nefness supplies smooth metrics $h_\epsilon$ with curvature at least $-\epsilon\omega$.
Choose a fixed sufficiently small $0<t<1$ such that the multiplier ideal of the singular weight of $h_b^t$ is $\cO_X$.
Theorem \ref{int:skoda-integrability} permits this choice:
the vanishing orders of local equations for $D$ have a uniform finite bound on the compact manifold, and multiplying them by $t/m$ makes them all less than one.
The mixed metric $h_t=h_b^t h_\epsilon^{1-t}$ is well defined on $L$,
has curvature at least $(tc-(1-t)\epsilon)\omega$, and still has trivial multiplier ideal.
Choose $\epsilon<tc/2$ and apply Theorem~\ref{nad:vanishing} in every degree $q>0$.

\emph{Now induct on the dimension.}
If $\nu(L)<n$, take a smooth hyperplane section $H$ in a sufficiently high power of the ample bundle $A$ used in the definition.
Replacing $c_1(A)$ by a positive integral multiple rescales the intersection numbers by positive factors and does not change numerical dimension.
The restriction $L|_H$ is nef, and
$\nu(L|_H)=\nu(L)$: the defining intersection numbers agree up to these positive factors,
and $\nu(L)\le n-1$ keeps the maximal nonzero term within the allowed range.
Tensor the divisor sequence and use adjunction to obtain
\[
 0\to K_X\otimes L\to K_X\otimes L(H)
   \to K_H\otimes L|_H\to0.
\]
Since $L(H)$ is ample, the middle term has vanishing higher cohomology.
If $q>n-\nu(L)$, then $q\ge2$ and
$q-1>(n-1)-\nu(L|_H)$. Induction annihilates $H^{q-1}$ of the restriction term,
and the long exact sequence gives $H^q(X,K_X\otimes L)=0$.
This proves the entire numerical-dimension range.
\end{proof}

\section{Vanishing with Arbitrary Effective Divisors}

\begin{lemma}[Intrinsic numerical dimension for nef bundles]\label{nad:numerical-invariance}
Let $X$ be a smooth projective $n$-fold, $L$ a nef line bundle, and $A$ an ample line bundle. For $0\leq j\leq n$, $\int_Xc_1(L)^j\wedge c_1(A)^{n-j}=0$ if and only if $c_1(L)^j=0$ in real cohomology.
Consequently, numerical dimension is independent of the ample bundle $A$ and is preserved under pullback by birational modifications.
\end{lemma}
\begin{proof}
Choose a positive curvature form $\omega$ for $A$. Nefness supplies smooth representatives of $c_1(L)$ satisfying
$\theta_\epsilon\ge-\epsilon\omega$.
The positive closed forms $(\theta_\epsilon+\epsilon\omega)^j$ have uniformly bounded mass, controlled by their cohomological intersection numbers.
Weak compactness of the coefficient measures gives a closed positive limit current $S$ representing $c_1(L)^j$.
If $\int S\wedge\omega^{n-j}=0$, its trace measure vanishes.
Proposition \ref{cur:coefficient-bound} forces $S=0$, and hence the cohomology class is zero.
The converse is immediate. For the weak compactness used here, apply weak-star compactness to bounded finite measures on finitely many coordinate charts,
then take a diagonal subsequence on a countable dense set of test functions.

For a birational modification $\mu:Y\to X$ between smooth projective manifolds, $\mu^*$ is injective on real cohomology:
for closed forms $a,b$ of complementary degrees, the integration formula
$\int_Y\mu^*a\wedge\mu^*b=\int_Xa\wedge b$
and Poincar\'e duality prove injectivity.
Thus $(\mu^*c_1(L))^j=0$ exactly when $c_1(L)^j=0$, proving invariance of numerical dimension.
\end{proof}

\subsection*{Reduction to normal crossings}

We use Hironaka's log resolution theorem and coherence of direct images under projective proper morphisms; these theorems are stated in Appendix~\ref{dep:chapter}.
Resolution puts the divisor and exceptional set into normal crossings.
Coherence then permits Serre vanishing to be applied to the higher direct images.
The argument below proves the required local vanishing and descends the resulting cohomological statement to $X$.

\begin{lemma}[Vanishing for fractional normal crossing divisors]\label{nad:snc-fractional-kv}
Let $F$ be a line bundle on a smooth projective manifold, with an equality of rational line bundles
$F=N+\sum_i c_iD_i$, where $N$ is nef and rational (some positive tensor power is an ordinary line bundle), the $D_i$ have normal crossings, and $0\le c_i<1$.
Then $H^q(X,K_X+F)=0$ for $q>n-\nu(N)$.
\end{lemma}
\begin{proof}
First suppose $N$ is big. Choose a strictly positive singular metric $h_b$ on $N$ with logarithmic effective-divisor singularities,
and smooth approximate metrics $h_\epsilon$ with curvature at least $-\epsilon\omega$.
Equip $F$ with the metric
$h_b^t h_\epsilon^{1-t}\prod_i|s_{D_i}|^{-2c_i}$.
For rational bundle powers, first clear all denominators, construct the metric on the resulting integral bundle, and take the corresponding root.
Because $F$ is integral, the resulting frame transformation law is consistent.
Choose the smooth reference-metric factors compatibly. Its real curvature is
$t\ii\Theta_b+(1-t)\ii\Theta_\epsilon+2\pi\sum_i c_i[D_i]$.
Once $t>0$ is fixed, sufficiently small $\epsilon$ makes the curvature strictly positive.

We must also prove that its multiplier ideal is the structure sheaf.
Choose $p>1$ with $pc_i<1$ for every $i$, and write $p^{-1}+(p')^{-1}=1$.
Theorem~\ref{pos:snc-ideal} shows that
$\prod_i|s_{D_i}|^{-2pc_i}$ is locally integrable.
Choose $t>0$ small enough that Theorem~\ref{int:skoda-integrability} makes the exponential of the local singular weight of $h_b^{p't}$ integrable.
H\"older's inequality makes the product of the two weights integrable, so the multiplier ideal is trivial.
Nadel vanishing gives vanishing in every positive degree.

If $\nu(N)<n$, choose a general smooth very ample hyperplane $H$ transverse to every
intersection stratum of the $D_i$. Such a hyperplane exists by applying the tangent-incidence argument used in Theorem~\ref{gp:kodaira-big} to each of the finitely many smooth divisor-intersection strata; the union of their bad hyperplanes is a proper subset of the dual projective space. Zero-dimensional strata can be avoided. Their restrictions still form a fractional normal crossing divisor with coefficients below one.
The nef part $N+H$ of $F+H$ is big, so the case just proved gives
$H^{>0}(X,K_X+F+H)=0$.
Use the adjunction exact sequence, Lemma \ref{nad:numerical-invariance}, and the same dimension induction as in
Theorem \ref{nad:nef-numerical} to obtain vanishing in every degree
$q>n-\nu(N)$.
\end{proof}

\begin{theorem}[Kawamata--Viehweg vanishing with an arbitrary effective divisor]\label{nad:kv-divisor}
Let $X$ be smooth projective, let $E$ be a line bundle, and suppose $mE=L+D$, where $m\ge1$ is an integer, $L$ is nef, and $D$ is effective.
Then
\[
 H^q\bigl(X,\cO(K_X+E)\otimes\cI(m^{-1}D)\bigr)=0
       \qquad(q>n-\nu(L)).
\]
\end{theorem}
\begin{proof}
Choose a projective log resolution $\mu:Y\to X$, and write
$\mu^*D=\sum_i a_iD_i$, with these divisors and the exceptional divisors having normal crossings.
Set
$F=\mu^*E-\lfloor\mu^*D/m\rfloor$.
This is an integral line bundle, and
\[
 F=\frac1m\mu^*L+\sum_i\left\{\frac{a_i}m\right\}D_i,
\]
Here $\lfloor a\rfloor$ is the greatest integer at most $a$ and $\{a\}=a-\lfloor a\rfloor$; floors of divisors are taken coefficientwise. All fractional parts lie in $[0,1)$. Lemmas~\ref{nad:snc-fractional-kv} and \ref{nad:numerical-invariance} give
$H^q(Y,K_Y+F)=0$,$q>n-\nu(L)$.

\emph{First identify the zeroth direct image.}
Write the Jacobian divisor as $K_{Y/X}=K_Y-\mu^*K_X=\sum b_iD_i$, including all relevant exceptional components. If $s_D$ is a local equation of $D$, write $\varphi_D=\log|s_D|^2$; then $\cI(m^{-1}D)$ means $\cI(\varphi_D/m)$.
In normal crossing coordinates, change of variables transforms the integral of $|f|^2e^{-\varphi_D/m}$ into the integral of
\[
 |f\circ\mu|^2\prod_i|w_i|^{2b_i-2a_i/m}
\]
times a smooth positive factor bounded above and below.
Coordinatewise integration makes integrability equivalent to
$\operatorname{ord}_{D_i}(f\circ\mu)+b_i\ge\lfloor a_i/m\rfloor$.
Consequently,
\[
 \cI(m^{-1}D)
 =\mu_*\cO_Y(K_{Y/X}-\lfloor\mu^*D/m\rfloor),
 \quad
 \mu_*\cO_Y(K_Y+F)=\cO_X(K_X+E)\otimes\cI(m^{-1}D).
\]
For the converse identification, a direct-image section defines a holomorphic function on the isomorphism locus,
which extends across the modification center of codimension at least two by Hartogs. The displayed integrability condition puts it in the multiplier ideal.
Properness ensures that finitely many coordinate neighborhoods covering a fiber cover the entire preimage after the base neighborhood is shrunk,
so the local integrability comparison holds in both directions.

\emph{Next prove that the higher direct images vanish.}
Choose a very ample line bundle $A$ on $X$. For every positive integer $k$,
the nef part of $F+\mu^*(kA)$ is $\mu^*(L/m+kA)$, which is nef and big.
The big case of the preceding lemma gives
$H^i(Y,K_Y+F+\mu^*(kA))=0$,$i>0$.
The coherent direct-image theorem makes $R^i\mu_*\cO_Y(K_Y+F)$ coherent.
For sufficiently large $k$, Theorem~\ref{amp:serre} kills the positive-degree cohomology on the base of all these sheaves tensored with $A^k$.
The Leray spectral sequence constructed in Lemma~\ref{morse:leray-comparison} is then concentrated in base degree zero and gives
\[
 H^0\bigl(X,R^i\mu_*\cO_Y(K_Y+F)\otimes A^k\bigr)
       \simeq H^i(Y,K_Y+F+\mu^*(kA))=0.
\]
For sufficiently large $k$, the generation assertion of Theorem~\ref{amp:serre} also makes the sheaf on the left globally generated.
Since it has no global sections, it is the zero sheaf.
Thus $R^i\mu_*\cO_Y(K_Y+F)=0$ for $i>0$.

Finally apply Leray to the untwisted sheaf to obtain
$H^q(X,\mu_*\cO_Y(K_Y+F))\simeq H^q(Y,K_Y+F)$.
Substitute the zeroth direct-image formula and the vanishing from the first paragraph to obtain the conclusion.
\end{proof}

\section{Siu's Uniform Generation Theorem}

\begin{theorem}[Siu's uniform global generation]\label{nad:siu-generation}
Let $X$ be a smooth projective $n$-fold, let $A$ be very ample, and put $G=K_X+(n+2)A$.
Then $G$ is ample, and for every line bundle $L$ with a semipositive-curvature singular metric $h$,
the sheaf $\cO(G+L)\otimes\cI(h)$ is globally generated.
\end{theorem}
\begin{proof}
Give $A$ the positive Fubini--Study metric induced by its embedding, with curvature $\omega_A$.
Fix $x\in X$. Choose a basis $\sigma_1,\ldots,\sigma_N$ of the sections of $A$ vanishing at $x$.
Very ampleness makes their common zero set exactly $\{x\}$, and in a local frame at $x$,
their first differentials span the cotangent space.
Therefore
$\psi_x=\log\sum_j|\sigma_j|_A^2$
satisfies
$\psi_x=\log|z-x|^2+O(1)$ and
$\ii\partial\dbar\psi_x\ge-\omega_A$.

Take a local representative of any $f\in\cI(h)_x$, viewed also as a local section of $G+L$ in a frame.
Choose a cutoff supported in a small coordinate ball, with $\chi=1$ near $x$.
Regard $\dbar(\chi f)$ as a form in top holomorphic degree; the remaining line bundle factor is
$(n+2)A+L$. Multiply its metric by $e^{-(n+1)\psi_x}$.
The total curvature is at least
$(n+2)\omega_A+\ii\Theta_h-(n+1)\omega_A\ge\omega_A$.
The datum is supported away from $x$, where the additional pole weight is bounded, so it has finite weighted squared norm.
Write $h=e^{-\varphi}$ in the chosen local holomorphic frame. Theorem~\ref{nad:global-l2} gives $u$ with
$\dbar u=\dbar(\chi f)$ and, locally,
\[
 \int |u|^2e^{-\varphi}|z-x|^{-2(n+1)}<\infty.
\]
The correction $u$ is holomorphic near $x$, though generally not on all of $X$.
The global section is $s=\chi f-u$; it is holomorphic and belongs to the sheaf defined by $\cI(h)$.

It remains to prove that the germs of $s$ and $f$ agree in $\cI(h)_x/\mathfrak m_x\cI(h)_x$.
On a small coordinate ball, apply general Skoda division, Theorem \ref{div:skoda}, to $u$,
with generators $g=(z_1-x_1,\ldots,z_n-x_n)$, $m=n-1$, and $\epsilon=1$.
The preceding integral is exactly the required input, so
$u=\sum_j(z_j-x_j)u_j$, with every $u_j$ locally square integrable for $e^{-\varphi}$.
Hence $u\in\mathfrak m_x\cI(h)_x$.
Coherence of the multiplier ideal makes this quotient finite dimensional. Choose a basis and perform the construction for each basis vector.
Nakayama's lemma, \ref{alg:nakayama}, shows that the germs of the resulting finitely many global sections generate
$\cI(h)_x$. Since $x$ was arbitrary, this proves global generation.

Finally we prove ampleness of $G$.
Take $L$ trivial and $h=1$. Repeat the cutoff construction for $K_X+(n+1)A$,
using pole coefficient $n$ instead of $n+1$. The curvature is still at least $\omega_A$.
The resulting $u$ is holomorphic near $x$ and satisfies $\int|u|^2|z-x|^{-2n}<\infty$,
so radial integration forces $u(x)=0$. A local $f$ with $f(x)\ne0$ therefore gives a global section not vanishing at $x$.
Thus $K_X+(n+1)A$ is globally generated.
The tensor product of a globally generated line bundle and a very ample line bundle is very ample:
the product of their section maps, followed by the Segre embedding, preserves the separation of points and tangent directions furnished by the very ample map.
Consequently, $G$ is even very ample, and in particular ample.
\end{proof}


\chapter{Stein Manifolds, Exhaustion Functions, and Noncompact Vanishing}
\label{stv:chapter}

On a compact manifold, integration by parts involves a single metric. In the noncompact case, three further points need justification:
completeness of the metric, integrability of the datum, and the passage from solutions on sublevel sets to a global solution.
Growing weights and complete metrics allow the $L^2$ estimates to address these questions. The model exhaustion is $\rho(z)=|z|^2$ on $\C^n$: its sublevel sets are relatively compact balls, and its Levi form is positive. Replacing $\rho$ by a rapidly increasing convex function will let us control the growth of data at infinity.

Throughout, $n=\dim_\C X$ and $r=\operatorname{rank}E$. An \emph{exhaustion} is a real function with relatively compact sublevel sets; we shift smooth exhaustions to be nonnegative. Here \emph{weakly pseudoconvex} means admitting a smooth psh exhaustion, and psh means plurisubharmonic (Definition~\ref{pos:psh-definition}). Curvature is written $\Theta=F$, with $\ii\Theta$ a Hermitian $\operatorname{End}(E)$-valued $(1,1)$-form; for a line bundle it is a real $(1,1)$-form.

\section{Growing Weights and \texorpdfstring{$m$}{m}-Positivity}

\begin{lemma}[Growing weights in one variable]\label{stv:growing-weight}
Let $\rho\ge0$ be a smooth exhaustion, let $a$ be continuous and nonnegative,
and let $\mu$ be a positive measure finite on compact sets. There is a smooth increasing convex function $\chi$ such that
\[
 \chi'(\rho(x))\ge a(x),\qquad \int_Xe^{-\chi(\rho)}\,d\mu<\infty.
\]
Finitely many functions $a$ and finitely many such measures can be handled simultaneously.
\end{lemma}
\begin{proof}
Put $a_j=\sup_{\rho\le j+2}a$ and
$b_j=\log\bigl(1+2^j\mu\{j\le\rho\le j+1\}\bigr)$.
Choose increasing slopes recursively so that a continuous convex polygonal function has, on each $[j,j+1]$,
slope greater than $a_j+1$ and value greater than $b_j+2$ at the integer $j$.
This is possible: when choosing the preceding slope, one can already meet the finitely many requirements at the next integer; choose the initial value large enough for $j=0$.
Extend affinely to the left on the real line, convolve with an even nonnegative smooth kernel of sufficiently small support,
and add a constant. Convolution preserves increasing slopes, and the available margins retain the required inequalities.
The integral on the $j$th shell is then at most $2^{-j}$, so summing proves the assertion.
For finitely many conditions, first take the maximum of the functions and the sum of the measures.
\end{proof}

\begin{proposition}[Griffiths positivity supplies a K\"ahler form]\label{stv:griffiths-kahler}
If a complex manifold admits a Griffiths positive Hermitian holomorphic vector bundle $E$,
then $X$ is K\"ahler.
\end{proposition}
\begin{proof}
In the curvature notation of Definition~\ref{pos:nakano-definition}, the Hermitian coefficient form of $\ii\operatorname{tr}F_E$ evaluated on a nonzero $(1,0)$ vector $v=(v_j)$ is
\[
 \sum_{j,k,a}R_{j\bar k a\bar a}v_j\bar v_k>0.
\]
Each summand over a fixed unitary fiber vector $e_a$ is positive by Griffiths positivity. This coefficient evaluation distinguishes the real positive form $\ii\operatorname{tr}F_E$ from its purely imaginary evaluation on $(v,\bar v)$.
The trace curvature is the curvature of $\det E$, by Proposition~\ref{bun:first-chern}. Its local expression
$-\ii\partial\dbar\log\det H$ shows that it is a real closed $(1,1)$ form.
It is therefore a K\"ahler metric.
\end{proof}

\begin{theorem}[$m$-positive vanishing on a weakly pseudoconvex manifold]\label{stv:m-positive-vanishing}
Suppose $X$ has a smooth psh exhaustion and $E$ is an $m$-positive vector bundle of rank $r$, in the sense of Definition~\ref{l2:m-positive}.
If $q\ge1$ and $m\ge\min(n-q+1,r)$, then
\[
 H^q(X,K_X\otimes E)=0.
\]
\end{theorem}
\begin{proof}
Since $m$-positivity implies Griffiths positivity, Proposition~\ref{stv:griffiths-kahler} gives a K\"ahler metric $\omega$;
Proposition \ref{l2:complete-metric} also supplies a complete K\"ahler metric.
Represent the cohomology class by a smooth closed $(n,q)$ form $f$.
In the coefficient calculation of Lemma \ref{l2:metric-monotonicity},
each nonzero matrix $(u_{jK,a})$ has rank at most $\min(n-q+1,r)$.
Strict $m$-positivity therefore makes
$A=[\ii\Theta_E,\Lambda_\omega]$ positive definite in degree $(n,q)$.
Its least eigenvalue $\epsilon(x)$ is continuous and positive.

Apply Lemma~\ref{stv:growing-weight} with the positive measure $(1+\epsilon^{-1})|f|^2dV_\omega$ to choose $\chi$ so that
$\int_X(1+\epsilon^{-1})|f|^2e^{-\chi(\rho)}dV_\omega<\infty$.
Replace the bundle metric by $he^{-\chi(\rho)}$. The additional curvature
$\ii\partial\dbar\chi(\rho)\otimes\operatorname{Id}_E$ is Nakano semipositive.
The curvature operator in degree $(n,q)$ is still bounded below by $\epsilon(x)$.
Theorem \ref{l2:incomplete-metric} gives $\dbar u=f$.
The weight and metric are smooth. The minimal-norm solution satisfies $\dbar^*u=0$ distributionally, so $\Delta''u=\dbar^*f$; Theorem~\ref{ana:elliptic-regularity} makes $u$ smooth. The Dolbeault isomorphism, Theorem~\ref{sh:dolbeault}, gives the required vanishing.
\end{proof}

\begin{theorem}[Cartan B for vector bundles]\label{stv:cartan-b}
If $X$ has a smooth strictly psh exhaustion and $E$ is any holomorphic vector bundle, then
\[
 H^q(X,\Omega_X^p\otimes E)=0\qquad(q\ge1).
\]
\end{theorem}
\begin{proof}
Set $\omega=\ii\partial\dbar\rho>0$ and
$F=\Omega_X^p\otimes E\otimes K_X^{-1}$.
Choose any smooth Hermitian metric on $F$. Its Nakano curvature has a finite lower bound on each compact set.
The term $\chi'(\rho)\omega$ can be made arbitrarily large on successive compact shells,
so Lemma \ref{stv:growing-weight} allows a choice of $\chi$ making
the metric $he^{-\chi(\rho)}$ on $F$ Nakano positive.
At the same time, arrange integrability of the chosen cohomology representative and its inverse-curvature energy.
Identify $E$-valued $(p,q)$ forms with $F$-valued $(n,q)$ forms,
and apply Theorem~\ref{l2:incomplete-metric} and Theorem~\ref{ana:elliptic-regularity}, exactly as in Theorem~\ref{stv:m-positive-vanishing}, to obtain a smooth preimage.
No positivity of the original $E$ is assumed; it is supplied by the freely increasing strictly psh weight.
\end{proof}

\section{From \texorpdfstring{$L^2$}{L2} Solvability to the Levi Problem}

\begin{proposition}[Arbitrary finite jets on a closed discrete set]\label{stv:discrete-jets}
Let $X$ admit a smooth strictly psh exhaustion, and let $E$ be holomorphic.
Let $\mathcal E$ be the sheaf of holomorphic sections of $E$; an order-$s$ jet at $x$ is an element of $\mathcal E_x/\mathfrak m_x^{s+1}\mathcal E_x$. Given a closed discrete set $\{x_j\}$, nonnegative integers $s_j$, and a holomorphic $E$-valued jet of order $s_j$ at each point, there is $S\in H^0(X,E)$ realizing them all.
\end{proposition}
\begin{proof}
Choose pairwise disjoint, locally finite coordinate balls $B_j$ on which $E$ is trivial.
Represent the prescribed jets by coordinate polynomials $P_j$, and choose cutoffs $\eta_j$
supported in $B_j$ and equal to one on smaller balls around $x_j$.
The locally finite sum $s=\sum_j\eta_jP_j$ is a global smooth section,
and $v=\dbar s$ is supported away from all the smaller balls.
Choose further cutoffs $0\le\zeta_j\le1$, equal to one near $\operatorname{supp}\eta_j$
and supported in $B_j$. Shrink the coordinate radii below one, and set
\[
 \psi=\sum_j(n+s_j)\zeta_j\log|z-x_j|^2,
 \quad
 \psi_\delta=\sum_j(n+s_j)\zeta_j
            \log(|z-x_j|^2+\delta^2).
\]
The sums are locally finite. Each $\psi_\delta$ is smooth and decreases to $\psi$ as $\delta\downarrow0$.
On the annuli where the cutoffs vary, $|z-x_j|$ is bounded away from zero,
so the negative parts of $\ii\partial\dbar\psi_\delta$ have a locally continuous upper bound independent of small $\delta$.
On the central balls, $\log(|z-x_j|^2+\delta^2)$ is psh and contributes no negative part.

Fix a complete K\"ahler metric $\widehat\omega$, and regard $v$ as an $(n,1)$ form with values in
$F=E\otimes K_X^{-1}$.
Using the strictly positive Levi form of $\rho$ and Lemma \ref{stv:growing-weight},
choose $\chi$ so that the Nakano curvature of $F$ with metric
$h_Fe^{-\chi(\rho)-\psi_\delta}$ is at least
$\widehat\omega\otimes\operatorname{Id}_F$, uniformly in small $\delta$, and
\[
 C:=\int_X|v|^2_{\widehat\omega,h_F}
              e^{-\chi(\rho)-\psi}dV_{\widehat\omega}<\infty.
\]
This integral is finite on each compact set because $v$ avoids all logarithmic poles.
Theorem~\ref{l2:complete-solving} gives $\dbar u_\delta=v$ with squared norm at most $C$.
For fixed $\delta_0$, every solution with $\delta<\delta_0$ also satisfies that bound
with weight $\chi(\rho)+\psi_{\delta_0}$.
Take successive locally weakly convergent subsequences, then let $\delta_0\downarrow0$.
Weak lower semicontinuity and monotone convergence give
\[
 \dbar u=v,\qquad
 \int_X|u|^2e^{-\chi(\rho)-\psi}dV_{\widehat\omega}\le C.
\]
Under the identification, $u$ is an $E$-valued $(0,0)$ form. On every central ball, $\dbar u=0$,
so the distributional Cauchy--Riemann equations make it holomorphic.
The smooth metric factors are bounded above and below there, hence
$\int|u|^2|z-x_j|^{-2(n+s_j)}dV<\infty$.
Expand the Taylor series and use orthogonality on spheres.
Every nonzero homogeneous term of degree at most $s_j$ would yield a divergent radial integral.
Thus the germ $u_{x_j}$ lies in $\mathfrak m_{x_j}^{s_j+1}\mathcal E_{x_j}$, as also follows from Proposition~\ref{pos:radial-ideal}.
The section $S=s-u$ is globally holomorphic and has all prescribed jets.
\end{proof}

\begin{corollary}\label{stv:nonzero-section}
On a manifold with a strictly psh exhaustion, every positive-rank holomorphic vector bundle has a nonzero global section.
In fact, evaluation $H^0(X,E)\to E_x$ is surjective at every point.
\end{corollary}
\begin{proof}
Apply Proposition \ref{stv:discrete-jets} to one point and a zero-order jet.
\end{proof}

\begin{definition}
Here a second-countable complex manifold is called Stein if it is holomorphically convex,
global holomorphic functions separate any two points, and at every point
$n$ global holomorphic functions can be chosen as local coordinates.
Holomorphic convexity means that the holomorphic hull of every compact set $K$,
$\widehat K=\{x:|f(x)|\le\sup_K|f|\text{ for every }f\in\cO(X)\}$, is compact.
\end{definition}

\begin{theorem}[The Levi problem]\label{stv:levi-problem}
A complex manifold $X$ is Stein if and only if it admits a smooth strictly psh exhaustion.
\end{theorem}
\begin{proof}
\emph{From an exhaustion to the Stein property.}
Zero-order interpolation at two points separates them. First-order interpolation at one point produces
$n$ global functions with independent differentials; Theorem~\ref{bas:holomorphic-inverse} makes them local coordinates.
If a compact set $K$ had noncompact holomorphic hull, choose in that hull a closed discrete sequence $x_j$ escaping every compact set.
Proposition \ref{stv:discrete-jets} gives a holomorphic function with $f(x_j)=j$.
But $x_j\in\widehat K$ would imply $j\le\sup_K|f|$, a contradiction.
The hull is closed, so relative compactness makes it compact. This proves the first implication.

\emph{From the Stein property to an exhaustion.}
Start with a compact exhaustion, repeatedly take holomorphic hulls, and place each term in the interior of the next.
This yields holomorphically convex compact sets $K_j$ with
$K_j\subset\operatorname{int}K_{j+1}$ and $\bigcup_jK_j=X$.
For every point $x$ in the compact shell $A_j=K_{j+2}\setminus\operatorname{int}K_{j+1}$,
we have $x\notin K_j$. Holomorphic convexity of $K_j$ supplies $g_x$ with
$|g_x(x)|>\sup_{K_j}|g_x|$.
Retain the strict inequality on a small neighborhood of $x$, and take a finite subcover of the shell.
Raise the resulting finitely many functions to sufficiently high powers and rescale them to obtain $f_{j,1},\ldots,f_{j,N_j}$ satisfying
\[
 \sum_{\ell=1}^{N_j}|f_{j,\ell}|^2\le4^{-j}\quad\hbox{on }K_j,
 \qquad
 \max_\ell|f_{j,\ell}|\ge j\quad\hbox{on }A_j.
\]
Both conditions can be imposed because the lower bound on each chosen neighborhood
strictly exceeds the upper bound on $K_j$, and the ratio tends to infinity under high powers.
More precisely, arrange $|f_{j,\ell}|\le2^{-j}/\sqrt{N_j}$ on $K_j$,
which gives the stated squared-sum bound.

The series $u=\sum_{j,\ell}|f_{j,\ell}|^2$ converges on compact sets with all derivatives:
on coordinate balls inside a larger compact set, the estimates in Theorem~\ref{bas:cauchy} bound derivatives of every fixed order by
the function bounds on a slightly larger compact set, and the geometric series is summable.
Apply vector-valued Cauchy estimates to each finite tuple $(f_{j,1},\ldots,f_{j,N_j})$;
the derivative bounds are controlled by $4^{-j}$ independently of $N_j$, preserving summability.
Thus $u$ is smooth and psh. The second inequality makes it tend to infinity along every sequence escaping compact sets.
Finally, the local-coordinate property and second countability give countably many global holomorphic functions $g_\ell$
whose differentials span the cotangent space at every point.
Choose $\epsilon_\ell>0$ successively small enough that
$\sum\epsilon_\ell|g_\ell|^2$ converges on every compact set with all derivatives.
Its Levi form is positive definite at each point, since some coordinate subfamily has independent differentials there.
Consequently,
$\rho=u+\sum\epsilon_\ell|g_\ell|^2$
is a smooth strictly psh exhaustion.
\end{proof}

\begin{remark}
A strictly psh exhaustion is also called strong $1$-completeness in some references.
An exhaustion that is strictly psh only outside a compact set is usually called strong pseudoconvexity; it permits compact exceptional analytic sets
and does not automatically imply the Stein property. The location of strict positivity is therefore part of the hypothesis.
\end{remark}

\section{Noncompact Akizuki--Kodaira--Nakano Vanishing}

\begin{lemma}[Controlling growth and derivatives simultaneously]\label{stv:derivative-weight}
Given a continuous $b:[0,\infty)\to(0,\infty)$ and an integer $N\ge1$,
there is a smooth increasing convex $\chi$ with $\chi''\ge1$ and
\[
 e^{-\chi(s)}\bigl(1+\chi'(s)+\chi''(s)\bigr)^N\le b(s)
 \quad(s\ge0).
\]
\end{lemma}
\begin{proof}
Write $y=\chi'$, choose $y(0)$ sufficiently large, and maintain
$1\le y'\le y^2/(4N)$. Then
\[
 \frac{d}{ds}\log(e^{-\chi}y^{2N})
 =-y+2N\frac{y'}y\le-\frac y2,
 \qquad 1+y+y'\le C_Ny^2.
\]
Thus $F=e^{-\chi}y^{2N}$ is decreasing.
To decrease $F$ by any further prescribed factor on a very short interval, use
$y'=y^2/(8N)$ to grow from $a$ to a chosen finite $B>a$.
The required time is $8N(1/a-1/B)<8N/a$, while $\chi$ increases by
$8N\log(B/a)$, so $F$ is multiplied by $(B/a)^{-6N}$.
Choosing $y(0)>128N$ allows each such growth interval to have length less than $1/8$.

Put $b_j=\min_{[j,j+1]}b>0$.
Choose $\chi(0)$ large enough that $C_N^NF(0)\le b_0$.
In the last half of $[j-1,j]$, for $j\ge1$, insert one such growth stage so that
$C_N^NF(j)\le b_j$; use $y'=1$ elsewhere.
Smooth the transitions between the two differential equations by choosing a smooth control
$a(s)\in[0,1]$ and setting
$y'=(1-a(s))+a(s)y^2/(8N)$, with $a$ flat at the transition endpoints.
Continue the growth until $F$ lies slightly below its target, then switch off smoothly.
Since $B$ can be arbitrarily large, the smooth transitions can be chosen short enough to retain the target bound,
while preserving $1\le y'\le y^2/(4N)$.
Only finitely many growth stages occur on a bounded interval, each ending at a finite $B$, so there is no finite-time blowup.
On $[j,j+1]$, $F(s)\le F(j)$, which proves the required inequality.
\end{proof}

\begin{theorem}[AKN vanishing on weakly pseudoconvex manifolds]\label{stv:noncompact-akn}
Suppose $X$ has a smooth psh exhaustion and $L$ has a smooth positive-curvature metric.
Then
\[
 H^q(X,\Omega_X^p\otimes L)=0\qquad(p+q>n).
\]
\end{theorem}
\begin{proof}
Put $\theta=\ii\Theta_L>0$, let $\rho\ge0$ be the given exhaustion,
and let $f$ be a smooth closed $(p,q)$ representative of the cohomology class.
For an increasing convex $\chi$, change both the bundle and base metrics:
\[
 h_\chi=he^{-\chi(\rho)},\qquad
 \omega_\chi=\theta+\ii\partial\dbar\chi(\rho).
\]
If $\chi''\ge1$, then
$\omega_\chi\ge\ii\partial\rho\wedge\dbar\rho$.
Thus $|d\rho|_{\omega_\chi}$ has a fixed upper bound. A curve escaping every compact set has
unbounded $\rho$ and hence infinite length, proving completeness of $\omega_\chi$.

We must also make $f$ integrable for this metric, which itself depends on $\chi$.
Fix a smooth reference metric $\omega_0$. On each compact shell $j\le\rho\le j+1$,
there are constants $c_j,C_j>0$ such that
\[
 c_j\omega_0\le\omega_\chi
 \le C_j(1+\chi'+\chi'')\omega_0.
\]
The lower bound controls the inverse matrices in the form norm; the upper bound controls volume. Therefore
\[
 |f|_{\omega_\chi,h}^2dV_{\omega_\chi}
 \le C'_j(1+\chi'+\chi'')^n|f|_{\omega_0,h}^2dV_{\omega_0}.
\]
After removing the derivative factor, the right-hand side has finite integral on each shell.
Apply Lemma \ref{stv:derivative-weight} with a continuous positive $b$ sufficiently small on each shell
to obtain $\int|f|_{\omega_\chi,h_\chi}^2<\infty$.

Since $\ii\Theta_{h_\chi}=\omega_\chi$,
Proposition \ref{pos:curvature-eigenvalues} gives
$[\ii\Theta_{h_\chi},\Lambda_{\omega_\chi}]=(p+q-n)\operatorname{Id}$.
Theorem~\ref{l2:complete-solving}, followed by Theorem~\ref{ana:elliptic-regularity} for the minimal solution, produces a smooth $u$ with $\dbar u=f$. Theorem~\ref{sh:dolbeault} proves vanishing.
Merely asking for a large weight would miss the growth of volume caused by its second derivatives;
the preceding lemma controls that additional growth.
\end{proof}

\begin{corollary}\label{stv:euler-positive}
On a compact complex manifold, if $L\otimes K_X^{-1}$ has a smooth positive-curvature metric, then
$\chi(X,L)=h^0(X,L)$.
\end{corollary}
\begin{proof}
Positive curvature itself gives a K\"ahler metric. Corollary~\ref{pos:kodaira} applied to
$K_X\otimes(L\otimes K_X^{-1})=L$ annihilates all higher cohomology.
The conclusion follows from the definition of the Euler characteristic.
\end{proof}

\section{Partial Curvature Positivity: The Compact Case and Noncompact Gluing}

\begin{lemma}[Counting eigenvalues]\label{stv:eigenvalue-count}
Let $0\le k\le n$ be an integer, $0\le\delta\le1$, and $0\le\lambda_j\le1$, with at least $n-k$ of the $\lambda_j$ at least $1-\delta$.
If $|I|=p$ and $|J|=q$, then
\[
 \sum_{j\in I}\lambda_j+\sum_{j\in J}\lambda_j-\sum_{j=1}^n\lambda_j
 \ge p+q-n-k-n\delta.
\]
\end{lemma}
\begin{proof}
The left-hand side equals
$\sum_{I\cap J}\lambda_j-\sum_{(I\cup J)^c}\lambda_j$.
Write $a=|I\cap J|$, $b=|(I\cup J)^c|$, so $a-b=p+q-n$.
At most $k$ terms in the intersection are below $1-\delta$, and each term in the negative sum is at most one.
The left-hand side is therefore at least $(a-k)(1-\delta)-b$;
if $a<k$, using zero for the positive contribution only improves the bound.
Now use $a\le n$ to obtain the asserted estimate.
\end{proof}

\begin{theorem}[Gigante--Girbau vanishing on compact K\"ahler manifolds]
\label{stv:compact-girbau}
Suppose the real curvature $\theta=\ii F_h$ of a smooth metric on $L$ is semipositive and has at least $n-k$ positive eigenvalues at every point, where $0\le k<n$.
If $X$ is compact K\"ahler, then
\[
 H^q(X,\Omega_X^p\otimes L)=0\qquad(p+q\ge n+k+1).
\]
\end{theorem}
\begin{proof}
Fix a K\"ahler metric $\omega$. Relative to it, the $(k+1)$st curvature eigenvalue in increasing order
has a positive lower bound $a>0$ on compact $X$.
Use $\omega_\varepsilon=\theta+\varepsilon\omega$.
The new curvature eigenvalues are $\lambda_j/(\lambda_j+\varepsilon)$.
They lie in $[0,1]$, and at least $n-k$ are at least $1-\varepsilon/a$.
Choose $\varepsilon<a/(2n)$. The preceding lemma gives
$[\theta,\Lambda_{\omega_\varepsilon}]\ge\frac12\operatorname{Id}$
in the bidegrees under consideration.
Theorem~\ref{pos:bkn} and Theorem~\ref{hod:hodge} force every harmonic representative to vanish. The omitted case $k=n$ has an empty degree range and is vacuous.
\end{proof}

\begin{proposition}[Solutions on sublevel sets]\label{stv:sublevel-girbau}
Let $X$ be weakly pseudoconvex K\"ahler, with $L$ satisfying the preceding theorem's curvature-rank condition.
Write $X_c=\{\rho<c\}\Subset X$.
If $f$ is smooth and closed of type $(p,q)$ near $\overline{X_c}$ and
$p+q\ge n+k+1$, then $\dbar u=f$ has a smooth solution on $X_c$.
\end{proposition}
\begin{proof}
The $(k+1)$st curvature eigenvalue has a positive lower bound $a$ on $\overline{X_c}$.
Choose $0<\varepsilon<a/(2n)$ and
$\lambda(t)=-A\log(c-t)$, with $A>2n$.
Set
\[
 \theta_\lambda=\theta+\ii\partial\dbar\lambda(\rho),\quad
 \omega_\lambda=\theta_\lambda+\varepsilon\omega,\quad
 h_\lambda=he^{-\lambda(\rho)}.
\]
Since $\lambda''=A/(c-t)^2$, the metric $\omega_\lambda$ is complete:
approaching the boundary forces divergence of $\int d\rho/(c-\rho)$.
As $\theta_\lambda\ge\theta$, at least $n-k$ eigenvalues relative to $\omega$ remain at least $a$.
The same count as in the compact case gives a curvature-operator lower bound of $1/2$.
Meanwhile, $\omega_\lambda$ grows at most like $(c-\rho)^{-2}$,
whereas $e^{-\lambda(\rho)}=(c-\rho)^A$.
Control the form norm by the fixed metric $\varepsilon\omega$ to obtain
$|f|_{\omega_\lambda,h_\lambda}^2dV_{\omega_\lambda}
\le C(c-\rho)^{A-2n}dV_\omega$.
The right-hand side is integrable; Theorem~\ref{l2:complete-solving} and local regularity (Theorem~\ref{ana:elliptic-regularity}) give the smooth solution.
\end{proof}

\begin{remark}[The global approximation input used here]\label{stv:runge-input}
The passage from sublevel sets to the whole manifold uses an approximation theorem for closed forms.
Under the preceding proposition's hypotheses, let $d<c$ be regular exhaustion values,
and let $\alpha$ be a locally square integrable closed $(p,s)$ form near $\overline{X_d}$, with
$p+s\ge n+k$. For every $K\Subset X_d$ and $\epsilon>0$,
there is a locally square integrable closed $(p,s)$ form $\beta$ on $X_c$ such that
$\|\alpha-\beta\|_{L^2(K)}<\epsilon$.
If $\alpha$ is defined only on $X_d$, first pass to a slightly smaller sublevel set.
This is a local-norm consequence of the Takegoshi--Ohsawa Runge approximation theorem for closed forms,
\cite[\S3]{TakegoshiOhsawa1981}.
The theorem of Takegoshi and Ohsawa uses the weight $-2n\log(d-\rho)$ and a complete metric.
A form locally square integrable near $\overline{X_d}$ is integrable for this weighted norm:
after combining the weight, changed form norm, and volume, the density relative to a fixed metric is bounded.
On each fixed $K$, the norm is equivalent to the ordinary $L^2$ norm, giving the formulation above.
Its proof requires boundary weights tending to infinity and uniform adjoint estimates.
We use this approximation theorem in the next proof. It allows the solutions on successive sublevel sets to be chosen compatibly.
\end{remark}

\begin{theorem}[Partial-positivity vanishing on weakly pseudoconvex K\"ahler manifolds]
\label{stv:noncompact-girbau}
Let $X$ be K\"ahler with a smooth psh exhaustion.
If $L$ has smooth semipositive curvature with at least $n-k$ positive eigenvalues everywhere, then
\[
 H^q(X,\Omega_X^p\otimes L)=0\qquad(p+q\ge n+k+1).
\]
This statement uses the approximation input explicitly recorded in Remark \ref{stv:runge-input}.
\end{theorem}
\begin{proof}
Choose strictly increasing regular values $c_j\to\infty$, and write $D_j=X_{c_j}$, starting with an auxiliary $D_0$.
For a global smooth closed $f$, Proposition \ref{stv:sublevel-girbau}
gives a smooth preimage $v_j$ on every $D_j$.
We construct preimages $u_j$ on $D_j$ satisfying
\[
 \dbar u_j=f,\qquad
 \|u_{j+1}-u_j\|_{L^2(D_{j-1})}<2^{-j}.
\]
Choose the first arbitrarily. Once $u_j$ has been constructed,
$v_{j+2}-u_j$ is closed on $D_j$, of type $(p,q-1)$,
with $p+q-1\ge n+k$.
Choose a smaller regular sublevel set inside $D_j$ containing $\overline{D_{j-1}}$,
then apply Remark \ref{stv:runge-input} to find a closed $b_j$ on $D_{j+1}$ with
\[
 \|b_j-(v_{j+2}-u_j)\|_{L^2(D_{j-1})}<2^{-j}.
\]
Set $u_{j+1}=v_{j+2}-b_j$ on $D_{j+1}$. The required equation and approximation bound follow immediately.

For fixed $D_m$, these difference norms are summable once $j\ge m+1$,
so $u_j|_{D_m}$ converges in $L^2(D_m)$.
The limits agree on overlaps and define a global $u\in L^2_{\rm loc}$.
Continuity of distributional differentiation gives $\dbar u=f$.
Proposition~\ref{l2:local-resolution}, Proposition~\ref{sh:dolbeault-resolution}, and Theorem~\ref{sh:resolution-theorem} show that the local $L^2$ and smooth Dolbeault resolutions compute the same sheaf cohomology,
so the class vanishes. A distributional solution suffices here; no smoothness assertion for the glued solution is needed.
\end{proof}

\section{Analytic Sets near the Boundary and Distinguishing Examples}

\begin{proposition}[Compact analytic sets stay away from a strictly pseudoconvex boundary]\label{stv:analytic-trapping}
Suppose $U\Subset X$ has smooth strictly pseudoconvex boundary.
There is a fixed compact $K\Subset U$ containing every compact analytic subset of $U$ with no zero-dimensional component.
\end{proposition}
\begin{proof}
Choose a strictly psh defining function $r$ near the boundary, with $U=\{r<0\}$ there.
Starting from any strictly Levi-positive boundary defining function, the modification $r+Cr^2$ corrects the normal eigenvalue:
tangential Levi positivity is retained, and the added $2C\partial r\wedge\dbar r$ controls the normal direction.
After shrinking the neighborhood, the full Hessian is positive definite.
Choose a smooth convex nondecreasing $\chi$, constant below an interior threshold
and with strictly positive derivative near zero. Extend $\chi(r)$ constantly to the interior to obtain
a smooth psh function $\psi$ on $U$, strictly psh on a fixed boundary collar.

For a compact irreducible analytic set $A$ of dimension $d\ge1$,
Theorem \ref{cur:analytic-integration} shows that $[A]$ is closed.
Stokes' theorem therefore gives
\[
 \int_{A_{\rm reg}}(\ii\partial\dbar\psi)^d=0.
\]
Noncompactness of $U$ causes no problem: insert a compactly supported cutoff equal to one near compact $A$.
The integrand is nonnegative everywhere. If $A$ met the boundary collar,
density of regular points in each local component would give a regular open piece with strictly positive integral, a contradiction.
Every positive-dimensional component is thus contained in the same compact set outside the collar.
\end{proof}

\begin{example}[Complete K\"ahler manifolds need not be Stein]\label{stv:punctured-examples}
The Poincar\'e-type metric on the punctured disk is
\[
 \frac{\ii\,dw\wedge d\bar w}{|w|^2(\log|w|)^2}
 =-4\ii\partial\dbar\log(-\log|w|^2).
\]
Direct differentiation of the radial function proves the identity; the minus sign is essential.
Radial length contains $\int dr/(r|\log r|)$, which diverges at both $r=0$ and $r=1$.

For $n\ge2$, the punctured ball $B^n\setminus\{0\}$ also has a complete K\"ahler metric:
start with a complete potential on the ball and, as in Lemma \ref{div:complement-complete}, add the potential
$-\sqrt{-\log|z|^2+C}$ to make every curve approaching the origin have infinite length.
It is nevertheless not Stein. Theorem~\ref{bas:hartogs} extends every holomorphic function to the whole ball,
so the holomorphic hull of a small sphere contains the small closed ball minus the origin, which is not relatively compact in the punctured ball.
Theorem \ref{stv:levi-problem} also rules out a strictly psh exhaustion.
\end{example}

\paragraph{Hartogs extension across an analytic set.}\label{stv:hartogs-analytic}
We will also use extension across an analytic set $A$ of complex codimension at least two in a complex manifold. Here is the deduction from Theorem~\ref{bas:hartogs}. Near a smooth point of $A$, choose coordinates with $A=\{z_1=\cdots=z_k=0\}$, $k\ge2$. Restrict the given function to the complement of $\{z_1=z_2=0\}$ and apply the theorem. The extension agrees with the given function on the full complement of $A$ by the identity theorem. These local extensions agree on overlaps for the same reason. Proposition~\ref{alg:analytic-stratification} leaves an analytic exceptional set of smaller dimension after the smooth strata of largest dimension have been removed. Repeat the construction for that set. Dimension drops at each step, so after finitely many local steps there is no exceptional set left. This proves holomorphic extension across all of $A$, with uniqueness again given by the identity theorem.

\begin{example}[Weak pseudoconvexity versus a strictly psh exhaustion]\label{stv:weak-not-stein}
Let $Y$ be a connected positive-dimensional compact complex manifold. The product $Y\times\C$ has psh exhaustion $|z|^2$,
but is not Stein: every holomorphic function is constant on each compact fiber $Y\times\{z\}$
and cannot separate two points in that fiber.
Likewise, if an analytic set $A$ of codimension at least two is removed from a connected compact complex manifold $X$,
all holomorphic functions extend to $X$ by the \hyperref[stv:hartogs-analytic]{preceding Hartogs argument} and are constant.
If $X\setminus A$ has at least two points, it is not Stein.
These examples distinguish weak pseudoconvexity, complete K\"ahler geometry, and the Stein property.
\end{example}

\section{Semipositive Curvature That Is Strictly Positive on a Dense Open Set}

\begin{proposition}[A direct proof in the K\"ahler case]\label{stv:dense-positive-kahler}
Let $X$ be compact K\"ahler, and let $L$ have smooth semipositive curvature, strictly positive on a dense open set.
Then $H^q(X,K_X\otimes L)=0$ for every $q\ge1$.
\end{proposition}
\begin{proof}
Take a harmonic $(n,q)$ representative $u$ by Theorem~\ref{hod:hodge}. Theorem~\ref{pos:bkn} gives
\[
 0\ge\int_X\langle[\ii\Theta_L,\Lambda]u,u\rangle dV\ge0.
\]
In bidegree $(n,q)$, the curvature-operator eigenvalues are sums of $q$ curvature eigenvalues,
so the operator is nonnegative everywhere and positive definite on the strictly positive open set.
The continuous nonnegative integrand has zero integral, hence $u$ vanishes on that dense open set.
Smoothness makes $u$ vanish everywhere. Theorem~\ref{hod:hodge} completes the proof.
\end{proof}

\begin{theorem}[The Grauert--Riemenschneider form on a general compact complex manifold]
\label{stv:dense-positive-general}
Let $X$ be connected and compact complex, and let $L$ have smooth semipositive curvature, strictly positive on a dense open set.
Then $H^q(X,K_X\otimes L)=0$ for every $q\ge1$.
The general non-K\"ahler case uses the two structural inputs stated below.
\end{theorem}
\begin{proof}
Semipositivity and strict positivity somewhere imply $\int_Xc_1(L)^n>0$, while the locus with exactly one negative curvature eigenvalue is empty.
The holomorphic Morse inequalities proved in this book therefore give
$h^0(X,L^m)\ge c m^n+o(m^n)$,$c>0$.
Proposition \ref{gp:morse-big} shows that $X$ is Moishezon.

The first structural input is: \emph{every smooth compact Moishezon manifold admits a modification
$\mu:Y\to X$ with $Y$ smooth projective.} This is the external input in Appendix~\ref{dep:moishezon}, together with Appendix~\ref{dep:resolution}.
The pullback $\mu^*L$ has semipositive curvature, strictly positive on the dense open set where $\mu$ is an isomorphism and the original curvature is strictly positive.
Apply Proposition~\ref{stv:dense-positive-kahler} on the projective manifold $Y$ to obtain
$H^q(Y,K_Y\otimes\mu^*L)=0$,$q\ge1$.

The second structural input is relative Grauert--Riemenschneider vanishing:
for such a proper modification, $R^j\mu_*K_Y=0$ for $j>0$ (Appendix~\ref{dep:gr}). Here $R^j\mu_*$ denotes the $j$th higher direct image.
Also $\mu_*K_Y=K_X$: top forms descend on the isomorphism locus
and their local coefficients extend across the codimension-at-least-two center in smooth $X$ by the \hyperref[stv:hartogs-analytic]{Hartogs argument for analytic sets}; holomorphic pullback gives the reverse inclusion.
The projection formula gives
$R^j\mu_*(K_Y\otimes\mu^*L)=L\otimes R^j\mu_*K_Y$.
The Leray spectral sequence has only its zeroth row, so
\[
 H^q(X,K_X\otimes L)
 \simeq H^q(Y,K_Y\otimes\mu^*L)=0.
\]
Thus the projective modification transfers the vanishing theorem to $X$, even when $X$ itself is not K\"ahler.
\end{proof}

\section{Higher Convexity}

\begin{definition}
If the Levi form of a smooth exhaustion has at least $n-q+1$ positive eigenvalues outside a compact set,
$X$ is called strongly $q$-convex; if this holds everywhere, it is strongly $q$-complete.
The remaining eigenvalues may be negative, so for $q>1$ this differs from a psh exhaustion.
\end{definition}

\begin{remark}[The theorems of Greene--Wu and Andreotti--Grauert]
\label{stv:q-convex-inputs}
Two theorems describe the wider role of partial Levi positivity.
\begin{enumerate}
\item Every connected noncompact complex $n$-manifold is strongly $n$-complete.
This is the theorem of Greene and Wu. See
\cite{GreeneWu1975} and Ohsawa's extension to analytic spaces, \cite{Ohsawa1984}.
It requires only one positive Levi eigenvalue and does not imply the Stein property.
\item If $X$ is strongly $q$-convex and $\mathcal F$ is a coherent analytic sheaf, then
$H^p(X,\mathcal F)$ is finite dimensional for $p\ge q$. If $X$ is strongly $q$-complete,
these groups vanish for $p\ge q$. This is the Andreotti--Grauert theorem,
\cite{AndreottiGrauert1962}.
\end{enumerate}
The second theorem applies to arbitrary coherent sheaves and is strictly stronger than the vector bundle form of Cartan B proved in full here.
Its proof uses local bumping for the relevant convexity index, cohomological approximation, and resolutions of coherent sheaves. These are further topics for which the cited papers and Appendix~\ref{dep:qconvex} provide the background.
A direct consequence is vanishing of the $n$th cohomology of every coherent sheaf on a noncompact complex $n$-manifold:
the first theorem gives strong $n$-completeness, and the second applies with $q=n$.
\end{remark}


\chapter{Adjoint Linear Systems and Effective Jet Separation}\label{eff:chapter}
\begin{learning}
Kodaira's theorem gives an embedding by a sufficiently high tensor power. Here we seek a numerical bound for the power, including prescribed values and derivatives at several points. The main construction combines Nadel vanishing with a dimension count: singular weights force the required vanishing orders, while a sequence of new sections reduces the multiplier-ideal zero locus to finitely many points. On surfaces, Reider's argument gives much sharper bounds by detecting the curves that obstruct separation.
\end{learning}

Throughout the global arguments, $X$ is a smooth projective complex manifold of positive dimension $n$. We may work on its connected components when needed.
We write $\mathcal O_{X,x}$ for the holomorphic local ring and $\mathfrak m_x$ for its maximal ideal (Lemma~\ref{alg:local-ring}). An order-$s$ jet of a line-bundle section means its germ modulo $\mathfrak m_x^{s+1}$ in the section-sheaf stalk; choosing a frame identifies this quotient with $\mathcal O_{X,x}/\mathfrak m_x^{s+1}$. A numerical polynomial is a polynomial with rational coefficients taking integer values at integers. Nefness and additive line-bundle notation are as in Chapter~\ref{gp:chapter}.

\section{Counting Jets and Values of Numerical Polynomials}
A first-order jet records a value and all first derivatives. More generally, an order-$s$ jet records every Taylor coefficient through degree $s$. For example, in two variables an order-two jet has the six coefficients of $1,z_1,z_2,z_1^2,z_1z_2,z_2^2$. Counting all these coefficients is the linear-algebra input to the interpolation construction.

\begin{lemma}[The dimension of a jet space]\label{eff:jets}
At a smooth point $x$ of complex dimension $n$,
$\dim_\C\mathcal O_{X,x}/\mathfrak m_x^{s+1}=\binom{n+s}{n}$.
Moreover, for $n\ge1,s\ge0$,
$\sum_{j=0}^s\binom{n+j-1}{n-1}=\binom{n+s}{n}$.
\end{lemma}
\begin{proof}
Reducing a convergent local power series modulo terms of total degree at least $s+1$ leaves the basis $z^\alpha$, $|\alpha|\le s$.
The number of nonnegative solutions of $\alpha_1+\cdots+\alpha_n=j$ is $\binom{n+j-1}{n-1}$, as seen by inserting $n-1$ separators among $j$ identical units.
Introduce the slack variable $\alpha_{n+1}=s-|\alpha|$.
Counting monomials of degree at most $s$ is then equivalent to counting $n+1$ nonnegative integers with sum $s$, giving $\binom{n+s}{n}$ and proving the summation identity.
The quotient $\mathfrak m^s/\mathfrak m^{s+1}$ counts only the homogeneous terms of degree $s$; an order-$s$ jet includes the lower-degree terms as well.
\end{proof}

\begin{lemma}[A large value in a short interval]\label{eff:numerical-polynomial}
Let $P$ be an integer-valued numerical polynomial of degree $d\ge1$ with leading coefficient $a_d/d!$, where $a_d$ is a positive integer.
Suppose $P(m)\ge0$ for integers $m\ge m_0$, with $m_0$ an integer.
For every integer $N\ge0$, some integer $m\in[m_0,m_0+dN]$ satisfies $P(m)\ge N$.
More precisely, for every integer $k\ge1$, some $m\in[m_0,m_0+kd]$ satisfies
$P(m)\ge a_dk^d/2^{d-1}$.
If $N\ge2d^2$, the first interval can be shortened to $[m_0,m_0+N]$.
\end{lemma}
\begin{proof}
If every value in the first interval were less than $N$, each integer point would be a root of one of
$P(m)=0,1,\ldots,N-1$, since the values are nonnegative integers.
These $N$ nonzero degree-$d$ polynomials have at most $dN$ roots in total, but the interval contains $dN+1$ integer points, a contradiction.
The $d$th difference with step $k$ gives
\[
 \sum_{j=0}^d(-1)^j\binom dj P(m_0+(d-j)k)=a_dk^d.
\]
Discard the negatively signed nonnegative terms.
The binomial coefficients with positive sign sum to $2^{d-1}$, so at least one corresponding value of $P$ is at least $a_dk^d/2^{d-1}$.

For the shorter interval, $d=1$ is already covered by the first assertion.
For $d=2$ and $N=2k$, the quadratic identity
$P(m_0)+P(m_0+2k)-2P(m_0+k)=a_2k^2$
shows that at least one endpoint has value at least $a_2N^2/8$.
For $N=2k+1$, use
$P(m_0)+P(m_0+N)-P(m_0+k)-P(m_0+k+1)=a_2k(k+1)$.
The larger endpoint value is at least $a_2(N^2-1)/8$.
For $N\ge8$, both bounds are at least $N$.
If $d\ge3$, take $k=\lceil2(N/2)^{1/d}\rceil$; then $k^d/2^{d-1}\ge N$.
It remains to check $kd\le N$.
The function $N/d-1-2(N/2)^{1/d}$ is increasing for $N\ge2$.
At $N=2d^2$ its value is $2d-1-2d^{2/d}\ge0$:
check $d=3$ directly, and for $d\ge4$ use $d^{2/d}\le\sqrt d\le d-\tfrac12$.
Thus $k\le N/d$, and the second assertion applies inside the shorter interval.
\end{proof}

\section{Constructing Singular Metrics for Interpolation}
Proposition~\ref{pos:jet-consequence} reduces jet separation to a vanishing theorem for a sheaf of sections with prescribed zeros. We therefore seek a positive singular metric whose multiplier ideal has zero-dimensional quotient and lies in the desired powers of the point ideals. Nadel vanishing will then make global sections surject onto that finite quotient.

We use additive notation for tensor products of line bundles.
Fix an ample line bundle $L$, choose $\mu L$ very ample, and let $\lambda_0,\ldots,\lambda_N$ be a basis of sections.
Set $\varphi_L=\mu^{-1}\log\sum_j|\lambda_j|^2$.
Its curvature $\omega_L=\ii\partial\dbar\varphi_L$ is strictly positive.
Absolute values in local formulas refer to coefficients in the current holomorphic frame; the total degrees ensure that all expressions glue with the correct metric transformation law.

\begin{lemma}[The metric construction]\label{eff:metric-induction}
Fix an integer $a\ge1$, a finite set of distinct points $x_1,\ldots,x_p$, and nonnegative integers $s_j$, and put
\[
 \delta=\sum_j\binom{n+(a+1)(n+s_j)-1}{n}.
\]
Let $b$ be an integer such that $H^q(X,(a+1)K_X+mL)=0$ for $q>0,m\ge b$.
Suppose an integer $b'\le b$ and a number $\eta\ge0$ satisfy
$b+\eta<\frac{a+1}{a}b'$
and the following numerical condition:
for every coherent ideal $I$ whose quotient $\mathcal O_X/I$ has positive-dimensional support, if
$P_I(m)=\chi(X,((a+1)K_X+mL)\otimes\mathcal O_X/I)$
is nonnegative for all integers $m\ge b$, then it is at least $\delta+1$ at some integer in $[b,b+\eta]$.
Then $aK_X+b'L$ admits a singular metric $h$ with strictly positive curvature, whose multiplier-ideal zero locus has dimension at most zero, and whose local weights satisfy
$\varphi\le a(n+s_j)\log|z-x_j|^2+O(1)$.
\end{lemma}
\begin{proof}
Start with $I_0=0$. The first section is chosen using the assumed ordinary cohomology vanishing; the first metric is defined only after this section has been chosen.
Suppose we have chosen sections
$\sigma_i\in H^0(X,(a+1)K_X+m_iL)$,
each vanishing to order $(a+1)(n+s_j)$ at every $x_j$, with
$m_i<\frac{a+1}{a}b'$.
Set $e_i=(a+1)b'-am_i\ge1$ and define the metric weight
\begin{equation}\label{eff:algebraic-weight}
 \varphi_\nu=\frac1{(a+1)\mu}\log\left(
   \sum_{i=1}^{\nu}|\sigma_i|^{2a\mu}
                    \Bigl(\sum_{j=0}^N|\lambda_j|^2\Bigr)^{e_i}\right).
\end{equation}
Expanding each summand gives squared norms of sections of
$(a+1)\mu(aK_X+b'L)$, because
$a\mu((a+1)K_X+m_iL)+e_i\mu L=(a+1)\mu(aK_X+b'L)$.
The expansion must include all degree-$e_i$ monomials with their multinomial coefficients.
Keeping only the powers $\lambda_j^{e_i}$ would not in general give the curvature bound below.
Factor one copy of $\sum|\lambda_j|^2$ from the sum.
The remaining expression is still a sum of squared absolute values of holomorphic sections, so
\[
 \ii\partial\dbar\varphi_\nu\ge\frac1{a+1}\omega_L.
\]
The vanishing orders of the $\sigma_i$ also give the required logarithmic upper bounds.
Write $I_\nu=\mathcal I(\varphi_\nu)$ and let $Z_\nu$ be the possibly nonreduced zero space defined by its quotient sheaf.

If $Z_\nu$ is positive-dimensional, then for $m\ge b$, Theorem~\ref{nad:vanishing} gives
$H^q(X,((a+1)K_X+mL)\otimes I_\nu)=0$ for $q>0$,
since the metric is on $aK_X+b'L$ and $m-b'\ge0$.
Together with the assumed ordinary vanishing, the short exact sequence yields
\[
 H^q(Z_\nu,((a+1)K_X+mL)|_{Z_\nu})=0\quad(q>0),
\]
and surjectivity of restriction from $X$ to $Z_\nu$ on sections.
For $\nu=0$, these statements are the hypothesis and the identity map, respectively, so the induction starts.
By Theorem \ref{idx:hilbert-polynomial}, this $h^0$ equals the Hilbert polynomial.
The numerical condition selects $m\in[b,b+\eta]$ for which it is at least $\delta+1$.
Lift $\delta+1$ linearly independent sections and map their span to the jets of orders through $(a+1)(n+s_j)-1$ at all $x_j$.
By Lemma~\ref{eff:jets} the target has dimension $\delta$, so the map has a nonzero kernel.
A nonzero section $\sigma_{\nu+1}$ in that kernel has the specified vanishing orders and is nonzero modulo $I_\nu$, since the restricted sections were linearly independent.
Set $m_{\nu+1}=m$; the interval condition gives the required strict inequality.

Adding the new summand makes $\varphi_{\nu+1}\ge\varphi_\nu$, hence $I_{\nu+1}\supset I_\nu$.
Locally,
$e^{-\varphi_{\nu+1}}\le C|\sigma_{\nu+1}|^{-2a/(a+1)}$:
on a small chart choose a nonvanishing $\lambda_j$ to absorb the remaining factors.
Therefore
$|\sigma_{\nu+1}|^2e^{-\varphi_{\nu+1}}\le C|\sigma_{\nu+1}|^{2/(a+1)}$
is locally integrable, proving $\sigma_{\nu+1}\in I_{\nu+1}$ in local frames.
Its class modulo $I_\nu$ is nonzero, so the inclusion is strict.
Each multiplier ideal is coherent by Theorem~\ref{l2:multiplier-coherent}. Strong Noetherianity on compact $X$, Corollary~\ref{alg:strong-noetherian}, forbids infinitely many such strict inclusions.
The process must therefore reach a zero-dimensional or empty $Z_\nu$ after finitely many steps, giving the desired metric.
For nonreduced $Z_\nu$, the needed condition is nonvanishing modulo $I_\nu$; pointwise nonvanishing on its reduced support is unnecessary.
\end{proof}

\begin{theorem}[A nef bound for adjoint line bundles]\label{eff:fujita-nef}
If $L$ is ample and $\dim X=n$, then $K_X+(n+1)L$ is nef and $K_X+(n+2)L$ is ample.
\end{theorem}
\begin{proof}
First apply Lemma~\ref{eff:metric-induction} with no prescribed points, so $\delta=0$.
The first part of Lemma~\ref{eff:numerical-polynomial}, with $N=1$ and support dimension at most $n$, permits $\eta=n$.
Fix $a\ge1$ and choose $b_1$ so large that $aK_X+b_1L$ is positive.
Add $b_1\omega_L$ to the curvature of any smooth metric on $aK_X$ to obtain such a metric.
Its multiplier ideal is $\mathcal O$, with empty zero locus, providing the starting condition.
If $aK_X+b_kL$ already carries a positive-curvature metric with zero-dimensional multiplier-ideal zero locus, Theorem~\ref{nad:vanishing} and Proposition~\ref{sh:discrete-support} give
$H^q(X,(a+1)K_X+mL)=0$ for $q>0,m\ge b_k$.
This is the ordinary vanishing needed for the next step.
Set
\[
 b_{k+1}=\left\lfloor\frac a{a+1}(b_k+n)\right\rfloor+1.
\]
Then $b_k+n<\frac{a+1}{a}b_{k+1}$.
Starting from sufficiently large $b_1$, this integer sequence decreases strictly until
$b_*=a(n+1)+1$:
for $b>b_*$, the next value is smaller than $b$ and at least $b_*$, while $b=b_*$ is fixed by the recursion.
After finitely many steps we obtain a strictly positive singular metric with zero-dimensional multiplier-ideal zero locus on
$F_a=aK_X+[a(n+1)+1]L$.

Fix any point $x$.
The sections of the very ample bundle $\mu L$ vanishing at $x$ have common zero locus exactly $x$ and locally generate $\mathfrak m_x$.
They define a semipositive singular metric on $n\mu L$ with weight
$n\log|z-x|^2+O(1)$.
Tensor it with the metric on $F_a$.
The resulting metric on $F_a+n\mu L$ remains strictly positive, and its multiplier-ideal zero locus remains zero-dimensional and contains $x$.
Theorem~\ref{nad:vanishing} and Proposition~\ref{sh:long-exact} give a section of $K_X+F_a+n\mu L$ nonzero at $x$.
Since $x$ was arbitrary, this line bundle is globally generated and therefore semipositive.
Divide its first Chern class by $a+1$ to obtain
\[
 c_1(K_X+(n+1)L)+\frac{n(\mu-1)}{a+1}c_1(L).
\]
Taking the corresponding metric root and subtracting this last small multiple of the fixed curvature yields smooth metrics on $K_X+(n+1)L$ whose negative curvature bounds tend to zero.
It is therefore nef.
Adding the ample bundle $L$ makes $K_X+(n+2)L$ ample by Proposition \ref{gp:positivity-implications}.
\end{proof}

\begin{theorem}[An effective jet-separation bound]\label{eff:effective-jets}
Let $L$ be ample, $G$ nef, $x_1,\ldots,x_p$ distinct points, and $s_j\ge0$.
If
\[
 m\ge 2+\sum_{j=1}^p\binom{3n+2s_j-1}{n},
\]
then $2K_X+mL+G$ simultaneously generates all jets of orders through $s_j$ at these points.
\end{theorem}
\begin{proof}
First assume $n\ge2$ and $p\ge1$.
Put $\delta=\sum_j\binom{3n+2s_j-1}{n}$ and $b=\delta+2$.
We have $\delta+1\ge2n^2$, as follows from
$\binom{3n-1}{n}\ge\binom{3n-1}{2}\ge2n^2$ for $n\ge2$.
By Theorem~\ref{eff:fujita-nef}, $K_X+(n+1)L$ is nef, so $K_X+mL$ is ample for $m\ge b>n+1$.
Corollary~\ref{pos:kodaira} gives
$H^q(X,2K_X+mL)=0$ for $q>0,m\ge b$.

In Lemma~\ref{eff:metric-induction} take $a=1$, $b'=b$, and $\eta=\delta+1$.
Every positive-dimensional zero space has dimension $d\le n$.
Theorem \ref{idx:hilbert-polynomial} gives leading coefficient a positive integer divided by $d!$, while the higher-cohomology vanishing in the induction makes its polynomial nonnegative at all integers $m\ge b$.
Since $\delta+1\ge2n^2\ge2d^2$, the short-interval conclusion of the numerical-polynomial lemma gives a value at least $\delta+1$ in $[b,b+\eta]$.
Moreover, $b+\eta=2\delta+3<2b$, so the induction hypotheses hold.
We obtain a strictly positive singular metric on $K_X+bL$, whose multiplier-ideal quotient $\mathcal O_X/I$ has zero-dimensional support and whose weight at $x_j$ satisfies
$\varphi\le(n+s_j)\log|z-x_j|^2+O(1)$.
Radial integrability gives $I_{x_j}\subset\mathfrak m_{x_j}^{s_j+1}$.
Indeed, if the first nonzero homogeneous term of a holomorphic germ $f$ is $P_d$, then
$\int_{|\zeta|=1}|f(r\zeta)|^2d\zeta
=r^{2d}(\int_{|\zeta|=1}|P_d(\zeta)|^2d\zeta+O(r))$,
with strictly positive leading coefficient in parentheses.
If $d\le s_j$, the weighted radial integral has exponent
$2d-2(n+s_j)+2n-1\le-1$ and diverges.
An estimate along only one radial line would not suffice for this integrability conclusion.
For $m\ge b$, the bundle $(m-b)L+G$ is nef.
Choose a smooth metric on it whose negative curvature is less than half the existing positive lower bound.
The tensor metric on $K_X+mL+G$ is still strictly positive and has the same multiplier ideal $I$, since a smooth positive factor does not affect local integrability.
Theorem~\ref{nad:vanishing} gives
$H^1(X,(2K_X+mL+G)\otimes I)=0$.
Global sections therefore surject onto the sections of the finite quotient sheaf $\mathcal O/I$, with the indicated line-bundle twist.
Passing further to the quotients $\mathcal O_{x_j}/\mathfrak m_{x_j}^{s_j+1}$ remains surjective, because the intermediate kernel is zero-dimensional and has vanishing $H^1$.
This gives all required jets.

If $n=1$, put $\ell=\sum_j(s_j+1)$.
The hypothesis gives $m\ge2+2\ell$.
The bundle
$K_C^{-1}-mL-G+\sum_j(s_j+1)x_j$
has degree at most $2-2g-m+\ell<0$.
Serre duality therefore gives
$H^1(C,2K_C+mL+G-\sum_j(s_j+1)x_j)=0$.
The short exact sequence for the prescribed point subscheme gives surjectivity on jets, proving the curve case.
With no prescribed points the assertion is vacuous.
\end{proof}

\begin{example}[Pseudoeffective cannot replace nef]\label{eff:pseudo-counterexample}
Let $X$ be the blowup of $\mathbb P^2$ at one point, with exceptional curve $E$.
Its local model is $\{((x,y),[u:v]):xv=yu\}$.
The normal-coordinate transition along the exceptional curve is that of
$\mathcal O_{\mathbb P^1}(-1)$, so $E^2=-1$.
Fix any ample bundle $L$ and put $d=L\cdot E>0$.
The local Jacobian calculation gives $K_X=\pi^*K_{\mathbb P^2}+E$, hence $K_X\cdot E=-1$.
For the effective, therefore pseudoeffective, bundle $G=NE$,
\[
 (2K_X+mL+G)\cdot E=-2+md-N.
\]
For any fixed integer $m$, choose an integer $N>\max\{0,md-2\}$.
The resulting degree is negative.
A globally generated line bundle remains globally generated on a curve; a nonzero section has a nonnegative total number of zeros, so its degree cannot be negative.
Thus a bound on $m$ independent of a pseudoeffective twist $G$ cannot guarantee global generation.
Analytically, a singular metric on $G$ can introduce additional multiplier ideals with positive-dimensional zero loci, so the resulting quotient need not prescribe arbitrary jets at the chosen points.
\end{example}

\begin{corollary}[Very ampleness and explicit bounds for canonical powers]\label{eff:very-ample}
Let $L$ be ample and $G$ nef, as in Theorem~\ref{eff:effective-jets}, and set $C_n=\binom{3n+1}{n}$.
Then $2K_X+mL+G$ is very ample for $m\ge2+C_n$, and
$m(K_X+(n+2)L)+G$ is very ample for $m\ge C_n-2n$.
If $K_X$ is ample, then $mK_X$ is very ample for $m\ge C_n+4$.
If $-K_X$ is ample, then $-mK_X$ is very ample for $m\ge C_n$.
\end{corollary}
\begin{proof}
First jets at one point require $\delta=C_n$; values at two distinct points require
$\delta=2\binom{3n-1}{n}$.
The ratio
$\binom{3n+1}{n}/\binom{3n-1}{n}=(3n+1)(3n)/((2n+1)(2n))\ge2$
shows that the same bound gives absence of base points, separation of tangent vectors, and separation of pairs of points.
Compactness then gives an embedding.
For the second assertion, put $G'=(m-2)(K_X+(n+1)L)+G$.
It is nef, since the stated bound implies $m\ge2$.
The identity
\[
 2K_X+(m+2n+2)L+G'=m(K_X+(n+2)L)+G
\]
allows the first assertion to apply whenever $m+2n+2\ge C_n+2$, giving the claimed bound.
Finally take $L=K_X$ and $L=-K_X$ in turn and collect coefficients in $2K_X+uL$.
The first-jet and two-point conditions supply the separation needed for very ampleness, in addition to positivity of the underlying powers.
\end{proof}

\section{Reider's Theorem and Adjoint Bundles on Surfaces}
In this section $X$ is a smooth projective surface and $A$ is a line bundle. In intersection expressions we also write $A$ for $c_1(A)$, so $A^2=\int_Xc_1(A)^2$. The length of a zero-dimensional scheme is the complex dimension of its finite structure-ring space, including multiplicities.
The surface proof uses extensions of coherent sheaves as well as intersection numbers. The next two lemmas explain the extension construction; readers meeting it for the first time can first examine the obstruction list in Theorem~\ref{eff:reider} and the numerical applications in Corollary~\ref{eff:fujita-low-dimension}.

\paragraph{Stability on surfaces.}
For a bundle $V$ on a surface, $c_2(V)$ in a numerical formula means $\int_Xc_2(V)$, and products of first Chern classes mean their intersection numbers.
We use the rank-two Bogomolov inequality: on a smooth projective surface, a slope-semistable rank-two bundle satisfies $c_1^2\le4c_2$.
For a fixed ample class $H$, semistability means that every rank-one subsheaf $M$ satisfies
$c_1(M)\cdot H\le c_1(E)\cdot H/2$.
Thus if $c_1^2>4c_2$, there is a rank-one subsheaf violating this inequality; saturation can only increase its slope.
We explain how this produces a class in the positive cone, taking care that a line subsheaf need not be a smooth subbundle injective on every fiber.
The stability-theoretic input is recorded in Appendix~\ref{dep:bogomolov}; the argument from it to Reider's theorem is given in detail below.

\begin{lemma}[Duality for coherent sheaves on a surface]\label{eff:coherent-duality}
A coherent sheaf $\mathcal F$ on a smooth projective surface has a natural nondegenerate pairing
\[
 \operatorname{Ext}^j_X(\mathcal F,\mathcal O_X)
 \times H^{2-j}(X,\mathcal F\otimes K_X)\longrightarrow\C.
\]
Here $\operatorname{Ext}^1$ classifies short exact extensions with subobject $\mathcal O_X$ and quotient $\mathcal F$.
For $\mathcal F=\mathcal O_Z\otimes A$, with $Z$ zero-dimensional, the resulting
$\operatorname{Ext}^2$ is the dual of $H^0(Z,(K_X+A)|_Z)$.
\end{lemma}
\begin{proof}
We first explain the passage from vector-bundle duality to a sheaf that need not be locally free.
The generation assertion of Theorem \ref{amp:serre} gives successive surjections from finite-rank vector bundles
$E_0\to\mathcal F$ and $E_1\to\ker(E_0\to\mathcal F)$.
The second kernel $E_2$ is locally free.
Indeed, over a regular local ring of dimension two, Theorem \ref{alg:regular-resolution} gives a free resolution of length at most two.
Compare it with the two surjections just constructed, using the following elementary fact twice.
If free modules $P,P'$ surject onto $M$ with kernels $K,K'$, then the pullback module
$P\times_M P'$ gives split short exact sequences in both directions, hence
$K\oplus P'\simeq K'\oplus P$.
Every second kernel is therefore stably free.
A finite stably free module over a local ring is free, and a basis chosen at one point extends to a neighborhood by coherence and Nakayama's lemma.
We obtain a global finite resolution
\[
 0\longrightarrow E_2\longrightarrow E_1\longrightarrow E_0
   \longrightarrow\mathcal F\longrightarrow0.
\]

Place $E_j$ in degree $-j$.
Resolve each term of $E_\bullet\otimes K_X$ by its Dolbeault complex.
Take local Dolbeault cohomology first. Theorem~\ref{sh:dolbeault-lemma} leaves the holomorphic complex $E_\bullet\otimes K_X$, whose only cohomology sheaf is $\mathcal F\otimes K_X$ in degree zero. Shifting the finite horizontal index by two permits the elimination argument of Lemma~\ref{sh:double-complex}. Thus the total smooth complex $\mathcal C^\bullet$ is a bounded fine complex in degrees $-2,\ldots,2$, with this same sole cohomology sheaf.
To apply the nonnegative-resolution theorem precisely, put $\mathcal B^0=\operatorname{im}(\mathcal C^{-1}\to\mathcal C^0)$. The exact negative tail
\[
 0\longrightarrow\mathcal C^{-2}\longrightarrow\mathcal C^{-1}
 \longrightarrow\mathcal B^0\longrightarrow0
\]
and fine-sheaf acyclicity show that $\mathcal B^0$ is acyclic and that $\Gamma(X,\mathcal C^{-1})\to\Gamma(X,\mathcal B^0)$ is surjective. Next, the exact sequence with terms $\mathcal B^0$, $\mathcal C^0$, and $\mathcal C^0/\mathcal B^0$ makes the quotient acyclic and makes $\Gamma(X,\mathcal C^0)\to\Gamma(X,\mathcal C^0/\mathcal B^0)$ surjective. Hence
\[
 0\longrightarrow\mathcal F\otimes K_X
 \longrightarrow\mathcal C^0/\mathcal B^0
 \longrightarrow\mathcal C^1\longrightarrow\mathcal C^2\longrightarrow0
\]
is an acyclic resolution of the form required by Theorem~\ref{sh:resolution-theorem}. The two surjections on global sections show that truncation has not changed the cohomology of $\Gamma(X,\mathcal C^\bullet)$. Its cohomology is therefore $H^*(X,\mathcal F\otimes K_X)$. This order of calculation requires no assertion that tensoring with smooth functions preserves exactness.
The Dolbeault total complex of the dual complex
$\mathcal Hom(E_\bullet,\mathcal O_X)$ defines
$\operatorname{Ext}^*_X(\mathcal F,\mathcal O_X)$.
Comparison of locally free resolutions and chain homotopies shows independence of choices.
Locally, lift maps into successive kernels; the difference of two lifts lifts to the preceding term, producing a chain homotopy.
Comparison through a double resolution then induces the same map on cohomology.
In degree one this gives the usual description by extensions.
After locally lifting generators of the quotient sheaf, their relations and differences of lifts on overlaps form a total one-cocycle.
Changing lifts adds a coboundary.
Conversely, take the quotient prescribed by the relations and glue using the overlap differences to construct a short exact extension.
The constructions are inverse.
Applied to a short exact sequence, this computation gives the long exact Ext sequence.

Integration, evaluation, and the Koszul signs pair the two total complexes in complementary degrees $j$ and $2-j$. Integration by parts makes their differentials transpose to one another, so the pairing descends to cohomology.
We prove its nondegeneracy by induction on the number of nonzero terms in a bounded complex of vector bundles. For one term, including any degree shift, it is precisely Theorem~\ref{hod:serre-duality}. For a longer complex, its highest-degree term is a subcomplex; quotienting removes that term. This gives a short exact sequence of complexes that is split in each degree. Dualizing reverses this short exact sequence, and the Dolbeault resolutions give long exact sequences on total cohomology by the cochain-lifting argument of Proposition~\ref{sh:long-exact}.
By induction the cohomology spaces for the one-term subcomplex and the shorter quotient are finite-dimensional and paired perfectly with the corresponding dual complexes. The long exact sequence also makes the middle spaces finite-dimensional. Dualize one long exact sequence over $\C$ and use the integration pairings to map the other into it; transpose differentials give commutative squares, with the conventional signs on connecting maps. The five lemma now makes each middle pairing an isomorphism to the dual space. Induction proves nondegeneracy for the original three-term resolution, and hence the asserted duality.
Integration and evaluation commute with sheaf morphisms, giving naturality.
In particular, the map induced by the quotient $A\to A|_Z$,
$\operatorname{Ext}^2(A|_Z,\mathcal O_X)\to H^2(X,A^{-1})$,
is transpose to evaluation
$H^0(K_X+A)\to H^0(Z,(K_X+A)|_Z)$.
\end{proof}

\begin{lemma}[Local extensions for schemes of length at most two]\label{eff:serre-extension}
Let $A$ be a line bundle and $Z$ a zero-dimensional subscheme of length $\ell\le2$ on a smooth projective surface.
If $H^0(K_X+A)\to H^0(Z,(K_X+A)|_Z)$ is not surjective but restriction to every proper subscheme is surjective, then there is a rank-two vector bundle in an exact sequence
$0\to\mathcal O_X\to E\to\mathcal I_Z\otimes A\to0$,
with $c_1(E)=A$ and $c_2(E)=\ell$.
\end{lemma}
\begin{proof}
Choose a nonzero linear functional $\lambda$ annihilating the evaluation image.
By the preceding lemma it corresponds to an element of
$\operatorname{Ext}^2(\mathcal O_Z\otimes A,\mathcal O_X)$.
The long exact Ext sequence of
$0\to\mathcal I_ZA\to A\to A|_Z\to0$
lifts it to $\operatorname{Ext}^1(\mathcal I_ZA,\mathcal O_X)$:
its image in $\operatorname{Ext}^2(A,\mathcal O_X)$ is transpose to global-section evaluation and is zero.
Call the middle sheaf of the resulting extension $E$.
We must prove that it is locally free.
At a length-one point the ideal is $(z_1,z_2)$.
For a length-two scheme supported at one point, the maximal ideal of its local quotient is one-dimensional and has square zero; otherwise Nakayama would make that maximal ideal zero.
Choose $z_2$ whose image generates it.
An independent tangent direction supplies a defining equation with nonzero linear term, which the implicit function theorem allows us to take as $z_1$.
Then $I_Z$ contains $(z_1,z_2^2)$, and equality follows because both quotients have dimension two.
If $Z$ consists of two distinct points, use the length-one description at each.

At each point write $R=\C\{z_1,z_2\}$ and $I_Z=(f,g)$.
The Koszul resolution of this regular sequence is
\[
 0\longrightarrow R\xrightarrow{\,(-g,f)\,}R^2
 \xrightarrow{\,(f,g)\,}I_Z\longrightarrow0.
\]
Applying $\operatorname{Hom}_R(-,R)$ gives
$\operatorname{Ext}^1_R(I_Z,R)=R/(f,g)$.
If $a\in R$ represents the local extension class, its middle module is explicitly
\[
 E_a=R^3/R(a,g,-f).
\]
When $a$ is a unit, eliminate the first generator to get $E_a\simeq R^2$.
When $a$ is not a unit, all three components of the relation vector lie in the maximal ideal, so
$\dim_\C E_a/\mathfrak mE_a=3$.
For a free rank-two module this dimension would be two, so $E_a$ is then not free.
Thus local freeness is equivalent to the local extension class being a unit.

Relate this unit condition to $\lambda$.
For $(f,g)=(z_1,z_2)$ or $(z_1,z_2^2)$, the local duality pairing is
\[
 (a,b)\longmapsto\frac1{(2\pi\ii)^2}
 \int_{|z_1|=\varepsilon_1,\ |z_2|=\varepsilon_2}
       \frac{a(z)b(z)\,dz_1\wedge dz_2}{f(z)g(z)}.
\]
This formula follows from the two connecting maps of the Koszul resolution:
use cutoff versions of $1/f$ and $1/g$ for the lifts, then integrate the Dolbeault connecting forms by parts onto the two small circles.
Two applications of the one-variable Cauchy formula give $a(0)b(0)$ in length one.
In length two, for $a=a_0+a_1z_2$ and $b=b_0+b_1z_2$, they give $a_0b_1+a_1b_0$.
Changing frames modifies the perfect pairing by a unit, preserving the following criterion.
In $\C[z_2]/(z_2^2)$, the socle, the subspace annihilated by the maximal ideal, is $\C z_2$.
The formula gives $\lambda_a(z_2)=a_0$.
Thus $a$ is a unit exactly when $\lambda_a$ is nonzero on the socle.
For length one the socle is the whole one-dimensional local quotient, giving the same criterion.

If $\lambda$ vanished on the socle at some point, it would descend after quotienting by that length-one ideal to a proper subscheme.
The descended functional would remain nonzero and would still annihilate all evaluations of global sections, contradicting minimality.
Every local extension class is therefore a unit, so $E$ is locally free everywhere.
In the free presentation, the section $\mathcal O_X\to E$ has components
$(-a^{-1}g,a^{-1}f)$, so its zero scheme has ideal exactly $(f,g)=I_Z$.
Taking determinants gives $c_1(E)=A$.
We compute the second Chern number from the already proved index formula. For a rank-$r$ bundle $V$ on a surface, expansion of Theorem~\ref{idx:hrr} through degree four gives
\begin{equation}\label{eff:surface-rr}
 \chi(X,V)=r\chi(X,\mathcal O_X)
  +\frac12c_1(V)\cdot\bigl(c_1(V)-K_X\bigr)-c_2(V).
\end{equation}
Indeed, $\operatorname{ch}(V)=r+c_1(V)+\tfrac12(c_1(V)^2-2c_2(V))$ in these degrees, the degree-two Todd term is $-K_X/2$, and its degree-four integral is $\chi(X,\mathcal O_X)$ by the same theorem with $V=\mathcal O_X$.
Euler-characteristic additivity in the two exact sequences gives
\[
 \chi(X,E)=\chi(X,\mathcal O_X)+\chi(X,A)-\ell;
\]
the zero-dimensional quotient has Euler characteristic $\ell$ by Corollary~\ref{sh:zero-dimensional-coherent}. Apply \eqref{eff:surface-rr} to $E$ and $A$, using $c_1(E)=A$ and $c_2(A)=0$. The constant and first-Chern terms cancel, leaving $c_2(E)=\ell$. This derives the required number without a separate top-Chern zero-counting theorem.
\end{proof}

\begin{theorem}[Reider's theorem in the form needed for Fujita's surface bounds]\label{eff:reider}
Let $X$ be a smooth projective surface and $A$ a nef line bundle. If $A^2\ge5$ and $K_X+A$ is not generated at a point $x$, there is an effective divisor $D$ through $x$ with
$(A\cdot D,D^2)=(0,-1)$ or $(1,0)$.
If $A^2\ge10$ and $K_X+A$ fails to separate a length-two subscheme, there is an effective divisor $D$ containing that subscheme or its minimal failing subscheme, with
\[
 (A\cdot D,D^2)\in\{(0,-1),(0,-2),(1,0),(1,-1),(2,0)\}.
\]
\end{theorem}
\begin{proof}
Take a minimal failing subscheme $Z$ of length $\ell=1$ or $2$ and construct $E$ by Lemma~\ref{eff:serre-extension}.
Since $A^2>4\ell$, the Bogomolov input gives a saturated rank-one subsheaf $M$ with
$(2c_1(M)-A)\cdot H>0$ and torsion-free quotient $Q=E/M$.
The sheaf $M$ is indeed a line bundle; here is the local argument.
Over $R=\C\{x,y\}$, the Koszul complex gives
$\operatorname{Tor}_2^R(Q,\C)=\{u\in Q:xu=yu=0\}=0$,
because $Q$ is torsion-free.
In a minimal local free resolution the differentials vanish modulo the maximal ideal, so the final rank equals the dimension of this Tor group.
Together with Theorem~\ref{alg:regular-resolution}, this gives a free resolution of $Q$ of length at most one.
Comparing it with $0\to M\to E\to Q\to0$ by the kernel comparison formula proved above makes $M$ stably free, hence free of rank one locally.

Wedge product gives a map $E\to\det E\otimes M^{-1}$.
Its kernel agrees generically with $M$.
An additional local kernel section would represent a nonzero class in $Q$ annihilated by a nonzero function, contrary to torsion-freeness.
The kernel is therefore exactly $M$ everywhere.
The image has the form $\mathcal I_W\otimes\mathcal O(D)$, with $D=A-M$.
At the generic point of a divisor the local ring is a discrete valuation ring, and the torsion-free quotient is free, so the wedge map is surjective there.
Thus $W$ has no divisorial component and is zero-dimensional.
Locally write the map $M\to E$ as $(u,v)$; then the ideal of $W$ is $(u,v)$.
The exact sequences with quotient $\mathcal I_W\otimes\mathcal O(D)$ give
$\chi(X,E)=\chi(X,M)+\chi(X,\mathcal O(D))-\operatorname{length}W$.
Applying \eqref{eff:surface-rr} and using $A=M+D$, the quadratic difference is
$\tfrac12((M+D)^2-M^2-D^2)=M\cdot D$, while the terms involving $K_X$ cancel. Thus
$c_2(E)=M\cdot D+\operatorname{length}W$. Since $c_2(E)=\ell$, we obtain
\[
 \ell=M\cdot D+\operatorname{length}W,\qquad
 (2M-A)^2=A^2-4\ell+4\operatorname{length}W>0.
\]
Together with $(2M-A)\cdot H>0$, this places $2M-A=A-2D$ in the same component of the positive cone as $H$.

The original section $s:\mathcal O_X\to E$ has nonzero image in $Q$.
Otherwise it would factor through $M$, and the zero divisor of that line-bundle section would be contained in the zero scheme $Z$ of $s$.
A nonzero holomorphic function with a zero has a divisorial zero-set component, rather than only isolated zeros on a surface.
The section of $M$ would therefore be nowhere zero, giving $M\simeq\mathcal O_X$.
This contradicts $(2M-A)\cdot H=-A\cdot H<0$.
Thus wedging $s$ with $M$ gives a nonzero section of $\mathcal O(D)$, defining an effective divisor $D$.
Locally, if $s=(f,g)$ and $M\to E$ is $(u,v)$, its equation is $ug-vf$.
It belongs to $(f,g)=I_Z$, so $D$ contains $Z$ \emph{scheme-theoretically} and is nonempty.
The second Chern-class identity also gives
\[
 A\cdot D-\ell\le D^2.
\]
The Hodge index theorem yields $D^2\le(A\cdot D)^2/A^2$, while the positive-cone condition gives
$0\le A\cdot D<A^2/2$.
For completeness, take a K\"ahler form $\omega$ representing $H$.
For a real harmonic $(1,1)$-form $\alpha$, the function $\Lambda\alpha$ is harmonic because the K\"ahler identities make $\Lambda$ commute with the Laplacian.
On a compact connected manifold it is constant, so write $\alpha=c\omega+\eta$ with $\Lambda\eta=0$ pointwise.
Wedge-product linear algebra on a surface gives
$\eta\wedge\omega=0$ and $\eta\wedge\eta=-|\eta|^2dV$.
The intersection form therefore has one positive direction and is negative definite on the remaining directions.
In coordinates where it is $t^2-|v|^2$, vectors $(t,v),(s,w)$ in the same positive cone satisfy $ts-v\cdot w>0$.
Since $A$ is nef and $A^2>0$, we have $A\cdot H>0$, so $A$ and $H$ lie in the same positive cone.
Hence $A\cdot(A-2D)>0$, or $A\cdot D<A^2/2$.
Effectivity of $D$ and nefness of $A$ give $A\cdot D\ge0$.
Decompose $D$ into $(A\cdot D)A/A^2$ and a component orthogonal to $A$.
The latter has nonpositive square, proving $D^2\le(A\cdot D)^2/A^2$.
For $A\cdot D=0$, the inequality is strict because a nonempty effective divisor has $H\cdot D>0$ and therefore nonzero cohomology class.

Put $a=A^2$ and $x=A\cdot D$.
For $\ell=1,a\ge5$, the inequalities $x-1\le x^2/a$ and $x<a/2$ exclude every integer $x\ge2$.
For $x=0$, the bounds $-1\le D^2<0$ give $D^2=-1$.
For $x=1$, the bounds $0\le D^2\le1/a<1$ give $D^2=0$.
For $\ell=2,a\ge10$, integers $x\ge4$ are excluded by
$x^2/a<x/2\le x-2$, and $x=3$ by $9/a<1=x-2$.
The remaining values $x=0,1,2$ give, respectively,
$D^2=-1,-2$, $D^2=-1,0$, and $D^2=0$.
The length-one possibilities are already included in this list.
\end{proof}

\begin{corollary}[Fujita bounds in dimensions one and two]\label{eff:fujita-low-dimension}
Let $L$ be ample. On a curve, $K_C+2L$ is globally generated and $K_C+3L$ is very ample.
On a surface, $K_X+3L$ is globally generated and $K_X+4L$ is very ample.
If $L^2\ge2$ on the surface, these improve to $K_X+2L$ and $K_X+3L$, respectively.
\end{corollary}
\begin{proof}
On a curve, $\deg L\ge1$, so the two degrees are at least $2g$ and $2g+1$; apply Theorem \ref{cur:ample}.
For global generation on a surface, take $A=3L$.
Then $A^2\ge9$ and every nonzero effective $D$ satisfies $A\cdot D\ge3$, contradicting Reider's failure list.
For very ampleness take $A=4L$.
Then $A^2\ge16$ and $A\cdot D\ge4$, excluding every length-two failure.
If $L^2\ge2$, take $A=2L$ and $A=3L$ instead.
Their squares are at least $8$ and $18$, and their intersections with nonzero effective divisors are at least $2$ and $3$, again excluding the respective lists.
Separation of schemes of lengths one and two is precisely absence of base points and separation of points and tangent vectors, giving the conclusions.
\end{proof}

\begin{remark}[Semiampleness and Fujita's conjecture]\label{eff:basepoint-free}
Semiampleness of $N=K_X+(n+1)L$ follows from the base-point-free theorem:
$N$ is nef and $N-K_X=(n+1)L$ is ample, so $aN$ is globally generated for some $a>0$.
This uses the base-point-free theorem in Appendix~\ref{dep:bpf}, together with the nefness proved in Theorem~\ref{eff:fujita-nef}.
It does not prove that $N$ itself is globally generated; that assertion and very ampleness of $K_X+(n+2)L$ itself are the sharp Fujita statements.
The nef bundle $K_X+(n+1)L$ need not be ample: on $\mathbb P^n$ with $L=\mathcal O(1)$, it is the trivial bundle.
\end{remark}


\chapter{Further Topics in Birational Geometry}\label{bir:chapter}
\begin{learning}
Curvature and section growth relate a complex manifold to projective geometry. We study finite maps, contractions, big line bundles, and rational curves. This chapter is intended for a subsequent seminar: it combines the earlier estimates with background results about analytic spaces and modifications, stated in Appendix~\ref{dep:chapter}.
\end{learning}

Unless specified otherwise, manifolds in this chapter are connected. A \emph{modification} is a proper surjective holomorphic map that is biholomorphic outside proper analytic subsets. A divisor is \emph{effective} when all its multiplicities are nonnegative; a rational divisor has rational multiplicities. Tensor and intersection notation is as fixed at the start of the book. We use the definitions of $a(X)$, $\kappa(L)$, nefness, and analytic bigness from Chapter~\ref{gp:chapter}.

\section{Finite maps and positivity on curves}
\begin{proposition}[A basepoint-free finite map]\label{bir:finite-ample}
Let $X$ be a compact complex manifold and $a$ a positive integer. If $L^a$ is globally generated and its associated map $f:X\to\mathbb P^N$ has zero-dimensional fibers, then $L$ is ample.
If $X$ is projective, strict positivity of $\deg(L|_C)$ on every irreducible curve $C$ is enough to ensure that the fibers are zero-dimensional.
\end{proposition}
\begin{proof}
Compactness makes $f$ proper. A compact zero-dimensional analytic fiber is finite, so $f$ is a finite analytic map and $L^a=f^*\mathcal O(1)$.
For a coherent sheaf $\mathcal F$, the direct image $f_*\mathcal F$ is coherent. In this finite case, coherence follows from the finite-projection model (Lemma~\ref{alg:analytic-finiteness} and Theorem~\ref{alg:oka}); the general higher-direct-image theorem is not required.
Moreover, $f_*$ is exact: its stalk at a point is the direct sum of the stalks over the finite fiber. To justify this description, choose disjoint small neighborhoods of the fiber points. Properness allows the base neighborhood to be shrunk so that its full inverse image lies in their union. The germs can then be represented and glued independently.
Thus higher direct images vanish. Proposition~\ref{sh:projection-formula} and the acyclic-resolution comparison (Theorem~\ref{sh:resolution-theorem}) give
\[
 H^q(X,\mathcal F\otimes L^{ak+r})
 =H^q(\mathbb P^N,f_*(\mathcal F\otimes L^r)\otimes\mathcal O(k)).
\]
Apply Theorem~\ref{amp:serre} for each $0\le r<a$ and take the largest threshold. All groups with $q>0$ vanish, so Proposition~\ref{amp:line-criteria} applies. The twist on the right is $\mathcal O(k)$ because $L^{ak}=f^*\mathcal O(k)$.

If $X$ is projective, a positive-dimensional fiber is a positive-dimensional projective analytic set. Successive general hyperplane sections produce a curve $C$ inside it. Since $L^a$ is pulled back from the target, it has degree zero on this contracted curve, contradicting the hypothesis. Hence all fibers are zero-dimensional.
\end{proof}

\begin{example}[Why the curve criterion needs projectivity]\label{bir:curve-free-torus}
A globally generated line bundle can have positive degree on every curve and still fail to be ample on a nonprojective manifold.
There are two-dimensional complex tori with no curves. On $\R^4/\Z^4$, choose a general constant complex structure $J$ for which no nonzero integral degree-two cohomology class has type $(1,1)$. Such choices exist: for each nonzero integral alternating form, $J$-invariance is a proper closed real-analytic condition. For a nonzero alternating form $b$, choose a real basis with $b(e_1,e_2)\ne0$. Set $Je_1=e_3$, $Je_3=-e_1$, $Je_2=t e_4$, and $Je_4=-t^{-1}e_2$, where $t>0$. Invariance would require $b(e_1,e_2)=t b(e_3,e_4)$, which fails for a suitable $t$. Changing the orientation of the basis treats either component of the space of complex structures. Thus each invariance condition is a proper real-analytic subset and has empty interior. There are countably many integral forms, so Baire's theorem allows all these proper conditions to be avoided simultaneously.
A curve would give a nonzero integral $(1,1)$-class. Nonvanishing follows by pairing its integration current with a K\"ahler form: $\int_C\omega>0$. This is a contradiction.
The trivial line bundle on such a torus is globally generated. Positivity on every curve is vacuous, but the bundle is not ample. One must assume projectivity or directly impose the finite-fiber condition of the preceding proposition.
\end{example}

\begin{proposition}[The numerical nef criterion in the projective case]\label{bir:numerical-nef}
Fix an ample line bundle $A$ with positive normalized curvature $\omega_A=c_1(A,h_A)$. On a smooth projective manifold, analytic nefness of $L$ is equivalent to $L\cdot C\ge0$ for every irreducible curve $C$, and also to
$\int_Yc_1(L)^{\dim Y}\ge0$ for every positive-dimensional irreducible subvariety $Y$.
\end{proposition}
\begin{proof}
For an analytically nef bundle choose curvature representatives $\alpha_\varepsilon\ge-\varepsilon\omega_A$.
Integrate the semipositive form $(\alpha_\varepsilon+\varepsilon\omega_A)^d$ over the regular locus of $Y$ and let $\varepsilon\to0$. The integral is a cohomological intersection number, hence a polynomial in $\varepsilon$; finite volume justifies each integral. Its limit is nonnegative. Curves are the case $d=1$.
Conversely, use Kleiman's projective nef-cone criterion, stated in Appendix~\ref{dep:ample-criteria}, as an external input. A class pairing nonnegatively with every curve belongs to the closure of the ample cone, so $L+\varepsilon A$ is ample for every positive rational $\varepsilon$. Take a positive metric on this rational line bundle and subtract $\varepsilon$ times the fixed curvature of $A$. This gives smooth metrics on $L$ with lower bounds tending to zero, exactly analytic nefness.
Kleiman's criterion is the substantive algebraic input; it is not obtained simply by replacing strict inequalities in Nakai--Moishezon by non-strict inequalities.
\end{proof}

\section{Strongly pseudoconvex disc bundles and Fourier weights}
\paragraph{Analytic-space inputs.}
We use Grauert's finiteness theorem and Remmert reduction: positive-degree coherent cohomology of a relatively compact strongly pseudoconvex domain is finite-dimensional, and its proper Stein reduction contracts the maximal connected compact analytic subsets. The finiteness input is stated in Appendix~\ref{dep:qconvex} and the proper-image input in Appendix~\ref{dep:remmert-chow}; their precise roles are identified below.

\begin{proposition}[Fourier weights imply Serre vanishing]\label{bir:disk-serre}
Suppose the zero section of $L^*$ has a circle-invariant, relatively compact, strongly pseudoconvex disc neighborhood $\Omega$.
For every coherent sheaf $\mathcal F$, one then has
$H^q(X,\mathcal F\otimes L^k)=0$ for $q>0$ and all sufficiently large $k$. Consequently $L$ is ample.
\end{proposition}
\begin{proof}
Let $\pi:L^*\to X$ be the bundle projection and choose a Stein trivializing cover of the base. The local domains $\pi^{-1}(U)\cap\Omega$ are pseudoconvex disc domains. Here we use the coherent-sheaf form of Cartan B, the $q=1$ case of the external Andreotti--Grauert input in Remark~\ref{stv:q-convex-inputs}; Theorem~\ref{stv:cartan-b} alone treats only vector bundles. Combined with Theorem~\ref{sh:leray}, it computes the cohomology of $\pi^*\mathcal F$ using this cover.
In a fiber coordinate $\xi$, multiplying a \Cech{} cochain with values in $\mathcal F\otimes L^k$ by $\xi^k$ produces a $\pi^*\mathcal F$-valued cochain on $\Omega$ of circle weight $k$.
In the opposite direction, apply the Fourier projection
\[
 P_ku(\xi)=\frac1{2\pi}\int_0^{2\pi}e^{-\ii k\theta}u(e^{\ii\theta}\xi)\,d\theta.
\]
It commutes with the cochain differential and, by holomorphic Taylor expansion, extracts exactly the degree-$k$ term. If a cocycle from the base becomes a coboundary on $\Omega$, project its primitive to weight $k$ and extract its coefficient. This yields a primitive on the base. Hence
$H^q(X,\mathcal F\otimes L^k)\hookrightarrow H^q(\Omega,\pi^*\mathcal F)$,
with distinct $k$ landing in distinct weight subspaces.
Grauert finiteness makes the target finite-dimensional. Nonzero vectors of different integer weights are linearly independent: in a finite linear relation, rotate by $\theta$, multiply by $e^{-\ii k\theta}$, and integrate to retain only its weight-$k$ summand. There can therefore be only finitely many nonzero weight subspaces. This proves the claimed vanishing.
Only finite linear relations are integrated here; no additional topological representation theory of a cohomology quotient is being assumed. Proposition~\ref{amp:line-criteria} gives ampleness.
The weight projection supplies a left inverse on the \Cech{} complex of the whole cover; this is what gives the injection on global cohomology.
\end{proof}

\begin{proposition}[Comparing contraction of a bundle and its tautological line bundle]\label{bir:blowdown-comparison}
Let $E$ be a holomorphic vector bundle of positive rank on a compact complex manifold $X$, and let $Y=\mathbb P(E^*)$.
The natural map from $\mathcal O_Y(-1)$ to $E^*$ is biholomorphic off the zero sections and maps the zero section $Y$ onto the zero section $X$. The two zero sections are analytically contractible to a point simultaneously. A contraction here is a proper holomorphic map of a neighborhood onto an analytic space, with that zero section as the fiber of one point, and biholomorphic on the complement.
\end{proposition}
\begin{proof}
The map is
$b(x,[\ell],v)=(x,v)$ for $v\in\ell\subset E_x^*$.
For $v\ne0$, its holomorphic inverse is $(x,v)\mapsto(x,[v],v)$.
The inverse image of a zero vector is the full projective fiber. Locally this is the blowup of the origin in a vector space, so $b$ is proper.
If the zero section of $E^*$ can be contracted, composing the contraction with $b$ contracts the zero section of $\mathcal O_Y(-1)$.
Conversely, suppose a neighborhood $U$ of $Y$ has a proper contraction $c:U\to W$ with $c^{-1}(o)=Y$. Properness of $b$ gives a neighborhood $V$ of $X$ such that $b^{-1}V\subset U$. Properness of $c$ then gives a sufficiently small neighborhood $W'$ of $o$ with
$c^{-1}W'\subset b^{-1}V$.
The map $c$ is constant on every $b$-fiber: off the zero section the fiber is a point, and every zero fiber lies in $Y$.
Thus $c^{-1}W'$ is a saturated open set. A proper surjection is a closed quotient map, so its image $V'$ under $b$ is open.

The equality $b_*\mathcal O=\mathcal O$ makes the local coordinate functions of $c$ descend uniquely to $V'$.
In the local blowup model, this equality can be checked directly. Off the zero section a function descends; near a fiber origin it is locally bounded because the inverse image of a small compact neighborhood is compact. Hartogs extension, or the one-variable removable-singularity theorem, extends it uniquely.
Embed $W'$ locally in a polydisc. The defining equations hold on the dense open set where descent is already known, hence also after extension. The descended map $V'\to W'$ is therefore holomorphic.
For compact $K\subset W'$, its inverse image is the compact set $b(c^{-1}K)$, so the descended map is proper. It is biholomorphic away from the zero section and contracts $X$ to a point. This proves the converse.
The descended map is holomorphic and proper, so it is a contraction in the analytic sense specified in the statement.
\end{proof}

\section{Bigness on a general compact complex manifold}
\paragraph{Background theorems.}
We use regularization of a strictly positive closed $(1,1)$-current to one with analytic singularities while preserving a strict lower bound; resolution of an analytic ideal by finitely many smooth-center modifications; and the existence of a smooth projective modification of a compact complex manifold with $a(X)=\dim X$.
These are, respectively, Demailly regularization, Hironaka resolution, and the Moishezon modification theorem. Appendix~\ref{dep:chapter} lists them. Smooth heat regularization alone does not imply these results.

\begin{theorem}[Big line bundles and Moishezon manifolds]\label{bir:moishezon}
On a compact connected complex $n$-fold, a line bundle is analytically big if and only if $\kappa(L)=n$.
The following are equivalent: $X$ is Moishezon; $a(X)=n$; $X$ carries a K\"ahler current in an integral class; and $X$ carries a big line bundle. Here Moishezon means bimeromorphic to a projective manifold, and a K\"ahler current is a closed positive $(1,1)$-current $T$ satisfying $T\ge\varepsilon\omega$ for a smooth positive Hermitian form $\omega$ and $\varepsilon>0$. An integral class lies in the image of $H^2(X,\Z)$ in real cohomology.
\end{theorem}
\begin{proof}
Suppose first that $L$ is analytically big. By the regularization input in Appendix~\ref{dep:moishezon}, choose on the same line bundle a still strictly positive metric with analytic singularities. The logarithmic coefficients in its Bergman regularization can be chosen rational. Fix such a metric, still denoted by $h$.
Resolve its ideal by $\mu:Y\to X$. Then
\[
 \mu^*c_1(L,h)=\beta+[D],\qquad \beta\ge\varepsilon\mu^*\omega,
\]
where $D$ is an effective rational divisor and $\beta$ is smooth and closed.
This identity comes from an actual metric, not just a de Rham class: after resolution the local ideal is generated by a monomial. Subtract its logarithm from the pulled-back weight. The remaining smooth weights define a metric on the rational line bundle $\mu^*L-D$, with first Chern form $\beta$.

There is an effective rational combination $E$ of exceptional divisors and a smooth curvature form $\theta_E$ such that $\mu^*\omega-\eta\theta_E$ is positive definite for some $\eta>0$.
For a single blowup along a smooth center, the exceptional divisor is the projectivized normal bundle, and $\mathcal O(-E)$ restricts to $\mathcal O(1)$ on each projective fiber. A Hermitian metric on the normal bundle gives a Fubini--Study metric whose curvature $-\theta_E$ is strictly positive in vertical directions; extend it smoothly to a neighborhood of the exceptional divisor.
The pullback of $\omega$ is already positive in horizontal directions. Multiplying $-\theta_E$ by a sufficiently small positive number fills the vertical degeneracy. The Schur-complement estimate for a block Hermitian matrix controls the bounded mixed terms while retaining horizontal positivity.
Away from the exceptional divisor, compactness and the existing strict positivity allow the coefficient to be reduced once more if needed. Repeat this for the finite sequence of blowups, choosing each new coefficient smaller than the positive margin already obtained. Pull back and add the earlier exceptional divisors to obtain $E$. Rational coefficients may be chosen because strict positivity is open in these coefficients.

For $0<\delta\le\varepsilon\eta$,
\[
 \beta-\delta\theta_E
 \ge(\varepsilon-\delta/\eta)\mu^*\omega
       +(\delta/\eta)(\mu^*\omega-\eta\theta_E)>0.
\]
Choose rational $\delta$ and an integer $m$ clearing the denominators of $D$ and $\delta E$. Define
$A=\mu^*L^m\otimes\mathcal O_Y(-mD-m\delta E)$.
The smooth weights above, divided by the corresponding divisor metric, give first Chern form $m(\beta-\delta\theta_E)$. Thus $A$ is positive.
We have an actual line-bundle identity $m\mu^*L=A+D'$ with $D'=mD+m\delta E$ effective. On a non-K\"ahler manifold, equality of de Rham Chern classes alone would not establish this identity.
Theorem~\ref{emb:kodaira} makes $Y$ projective, and Theorem~\ref{gp:kodaira-big} gives $\kappa(\mu^*L)=n$.
A modification does not add line-bundle sections. Indeed, $\mu_*\mathcal O_Y=\mathcal O_X$: a function descends on the isomorphism locus, is locally bounded near the compact fibers, and extends uniquely by normality and removability. The projection formula then gives
$H^0(Y,\mu^*L^k)=H^0(X,L^k)$, so $\kappa(L)=n$.

Conversely, $\kappa(L)=n$ implies $a(X)=n$, and the Moishezon modification theorem supplies a smooth projective $\mu:Y\to X$. Invariance of sections makes $\mu^*L$ big by section growth. Theorem~\ref{gp:kodaira-big} supplies a strictly positive curvature current $S\in c_1(\mu^*L)$.
Choose a K\"ahler form $\omega_Y$ with $S\ge c\omega_Y$. Since $\mu^*\omega\le C\omega_Y$, pushing forward gives
$\mu_*S\ge(c/C)\mu_*\mu^*\omega=(c/C)\omega$.
The last equality holds on the isomorphism locus and extends as an identity of currents by change of variables; the exceptional sets have zero smooth volume.
We must also check that this is the curvature of a singular metric on the original $L$. Fix a smooth first Chern form $\alpha$ of $L$ and write
$S=\mu^*\alpha+(\ii/2\pi)\partial\dbar v$, where $v$ is a globally integrable potential on $Y$. Proper pushforward commutes with distributional differentiation, giving
\[
 \mu_*S=\alpha+(\ii/2\pi)\partial\dbar(\mu_*v).
\]
The degree-zero current $\mu_*v$ is represented by a locally integrable function. On the isomorphism locus it is $v\circ\mu^{-1}$, and change of variables bounds its integral by that of $|v|$ times a bounded smooth Jacobian. It therefore defines a global singular weight on $L$. The strict lower bound proves analytic bigness.

Proposition~\ref{gp:morse-big} relates $a(X)=n$ to the existence of a line bundle with $\kappa(L)=n$; the Moishezon modification theorem relates it to bimeromorphic projectivity.
A big line bundle supplies an integral-class K\"ahler current through its first Chern current. Conversely, let $T$ be an integral-class K\"ahler current. Smooth Bott--Chern comparison, Theorem~\ref{hod:bc-aeppli}, gives
$T=\alpha+\ii\partial\dbar u$, with $\alpha$ a smooth real closed $(1,1)$-form in the integral class $[T]$. Taking real parts makes $u$ a real distribution.
Lemma~\ref{prj:integral-gluing} constructs a holomorphic line bundle with first Chern form $\alpha$. The local-potential theorem, Theorem~\ref{cur:local-potential}, and smoothness of pluriharmonic distributions show that $u$ is locally psh minus smooth, hence locally integrable. Add $2\pi u$ to the smooth metric weights of this line bundle. Its first Chern current is $T$, so the line bundle is big.
For the background on regularization, resolution, and the Moishezon property used here, see \cite{JiShiffman1993,Demailly1992,ArocaHironakaVicente2018}.
\end{proof}

\section{Rational curves and the canonical class}
\begin{definition}
A compact complex $n$-fold $X$ is \emph{uniruled} if a family of rational curves covers a dense open subset. Equivalently, there is a parameter space $S$ of dimension $n-1$ and a dominant meromorphic map $\mathbb P^1\times S\dashrightarrow X$. The requirement concerns one dominant family, rather than an unrelated choice of a rational curve through each point.
\end{definition}

\begin{proposition}[Two consequences of a covering family of rational curves]\label{bir:uniruled}
If a compact complex manifold $X$ is uniruled, then $K_X$ is not pseudoeffective and $\kappa(X)=-\infty$.
More generally, if $K_X$ is not pseudoeffective, then $\kappa(X)=-\infty$.
\end{proposition}
\begin{proof}
Take a dominant meromorphic family $F:\mathbb P^1\times S\dashrightarrow X$ and work first over the smooth locus of $S$. Its indeterminacy set $I$, where the meromorphic map is not holomorphic, has codimension at least two. Projection to $S$ is proper because $\mathbb P^1$ is compact, so $\pi(I)$ is analytic of dimension at most $n-2$ and cannot fill $S$.
After removing this set, the map is holomorphic along each entire $\mathbb P^1$ and on a parameter neighborhood. Choose a general parameter where the evaluation differential has full rank, as dominance allows. Shrink to a coordinate neighborhood $U\subset S$ so that the determinant is not identically zero on a general member.
For such a member $f:\mathbb P^1\to X$, differentiation in the curve and parameter directions gives a generically invertible bundle map
$\mathcal O(2)\oplus\mathcal O^{n-1}\to f^*T_X$.
Its determinant is a nonzero section of $f^*K_X^{-1}\otimes\mathcal O(-2)$, so Proposition~\ref{cur:degree} gives $\deg f^*K_X\le-2$.
Every global pluricanonical section therefore restricts to zero on a general member. These members cover a dense open set, so the identity theorem makes the section zero everywhere. Thus $\kappa(X)=-\infty$.

If $K_X$ had a positive curvature current, compose its local psh weights with $F$. The map has full rank on a dense open set, so a nonempty open parameter neighborhood contains points where the image has interior. A psh weight cannot be identically $-\infty$ on that interior. Its pullback is therefore a locally integrable psh function.
Use Fubini on finitely many charts covering $\mathbb P^1\times U$, after shrinking $U$ relatively compactly. For almost every parameter the pulled-back weights are integrable on all corresponding curve charts, and in particular are not identically $-\infty$. Choose such a parameter also outside the locus where the determinant vanishes identically.
The weights obey the transition law of the pulled-back line bundle and define a semipositive curvature current on $f^*K_X$. Its first Chern integral is the degree, which must be nonnegative, contradicting $\deg f^*K_X\le-2$. Thus $K_X$ is not pseudoeffective.
Finally, a nonzero section of $K_X^m$ defines an effective zero divisor whose integration current, divided by $m$, is a semipositive singular curvature representative of $K_X$. Hence $K_X$ would be pseudoeffective. The contrapositive proves the last assertion.
\end{proof}

\paragraph{The converse for compact K\"ahler manifolds.}\label{bir:ou}
For compact K\"ahler manifolds, non-pseudoeffectivity of $K_X$ does imply uniruledness by Ou's result \cite{Ou2025}. The cited preprint uses meromorphic integrability of foliations, a relative Albanese reduction, and relative Bergman-kernel metrics. The preceding proposition proves the other direction; the converse is the theorem cited here.
This must not be restated as ``$\kappa=-\infty$ if and only if uniruled.'' Non-pseudoeffectivity implies $\kappa=-\infty$ in only one direction; nonvanishing for a pseudoeffective canonical class is a separate problem.

\paragraph{Numerical dimension and abundance.}\label{bir:abundance}
First distinguish the definitions of $\nu$. For nef classes, use Proposition~\ref{gp:numerical-dimension}. For a pseudoeffective integral divisor $D$ on a smooth projective $n$-fold, the convention in this paragraph is
\begin{equation}\label{bir:numerical-convention}
 \nu_X(D)=\max\left\{j\in\{0,\ldots,n\}:\text{ for some ample }A,
 \ \limsup_{m\to\infty}\frac{h^0(X,\mathcal O_X(mD+A))}{m^j}>0\right\}.
\end{equation}
The ample twist $A$ is fixed before taking the limit. This is Nakayama's $\kappa_\sigma$ convention, equivalent to the numerical dimension used in \cite[Definitions 0.1(7), 0.3]{Eckl2016}; ordinary nonvanishing of $c_1(D)^j$ cannot replace it outside the nef case.
A minimal model here is a projective birational model $S$ with $\mathbb Q$-factorial terminal singularities and nef canonical $\mathbb Q$-Cartier divisor. The $\mathbb Q$-Cartier condition means that a positive integer multiple is a line-bundle divisor, and $\mathbb Q$-factorial means this for every divisor. Terminal means that the exceptional-divisor coefficients in $K_T=\pi^*K_S+\sum_Ea_EE$ on a smooth resolution are positive. The existence and properties of such models are part of minimal model theory.
Eckl's theorem \cite{Eckl2016} states that, if $S$ is a minimal model of a non-uniruled smooth projective variety $X$, then
$\nu_X(K_X)=\kappa_X(K_X)$ if and only if $K_S$ is semiample.
The model $S$ may have terminal singularities, $K_S$ is $\mathbb Q$-Cartier, and existence of the minimal model under discussion is part of the hypothesis. This is neither an equality $\kappa=\nu$ for arbitrary line bundles nor a proof of the abundance conjecture.
The elementary direction can be checked directly. Choose a positive divisible $m$ for which $mK_S$ is a globally generated line bundle, and let its map $f$ have image dimension $d$. Then $mK_S=f^*A$ for the restricted hyperplane bundle $A$ on that image. If $S$ is singular, take a smooth projective resolution $\pi:T\to S$ before making a differential-form calculation. The pullback $\pi^*(mK_S)=(f\circ\pi)^*A$ has a semipositive representative of generic rank $d$. Its $(d+1)$st power vanishes, whereas its $d$th power has positive intersection with a complementary power of a K\"ahler form on $T$. Thus the numerical dimension, computed by pullback, is $\nu(K_S)=d$. The image dimension of the section-ratio map gives $\kappa(K_S)=d$. The comparison with $X$ is part of the birational numerical-dimension framework in the cited theorem. The reverse implication requires the theory of good minimal models and abundance and is supplied by the cited theorem.


\chapter{Positive Currents Supported on Null Sets}
\label{pos:thin-chapter}

This chapter gives a short construction of a positive representative of a K\"ahler class whose coefficient measures are entirely singular with respect to smooth volume. This does not mean that its potential must take the value $-\infty$.
The construction is due to Tosatti \cite[Proposition~2.1]{Tosatti2026}. It uses a finite maximum of local potentials.

A useful real-variable analogy is $u(x)=|x|$. For a smooth compactly supported test function $\psi$, integration by parts on the two half-lines gives $\int_{\R}|x|\psi''(x)\,dx=2\psi(0)$. Thus $u''=2\delta_0$, where $\delta_0$ is the unit point mass at zero. A continuous potential can therefore have a derivative that is a singular measure. The construction below produces the corresponding phenomenon for positive currents.

\begin{definition}\label{pos:omega-psh}
Let $\omega$ be a smooth closed real $(1,1)$-form. A continuous real function $u$ is \emph{$\omega$-psh} if $\varphi+u$ is psh on every coordinate domain where $\omega=\ii\partial\dbar\varphi$.
Two local potentials differ by a pluriharmonic function, so this condition is independent of the potential. Proposition~\ref{pos:regularization} gives
$\omega+\ii\partial\dbar u\ge0$ as currents. In the construction below, positivity will be checked by taking maxima of local psh potentials.
\end{definition}

\begin{lemma}[Normalizing a strictly psh potential at a real minimum]
\label{pos:strict-real-minimum}
Suppose $\varphi$ is a smooth real-valued function near $0\in\C^n$ and
$\bigl(\varphi_{j\bar k}(0)\bigr)>0$.
Subtracting the real part of a holomorphic polynomial of degree at most two gives a function $\widetilde\varphi$ such that
\[
 \widetilde\varphi(0)=0,\qquad
 \widetilde\varphi(z)=\sum_{j,k}\varphi_{j\bar k}(0)z_j\bar z_k+o(|z|^2).
\]
On some small ball, there are constants $a,b>0$ with
$a|z|^2\le\widetilde\varphi(z)\le b|z|^2$, and
$\ii\partial\dbar\widetilde\varphi=\ii\partial\dbar\varphi$.
\end{lemma}
\begin{proof}
Group the real Taylor expansion by complex type:
\[
 \varphi(z)=\varphi(0)
 +2\operatorname{Re}\sum_j\varphi_j(0)z_j
 +\operatorname{Re}\sum_{j,k}\varphi_{jk}(0)z_jz_k
 +\sum_{j,k}\varphi_{j\bar k}(0)z_j\bar z_k+o(|z|^2).
\]
The first three terms are the real part of the holomorphic polynomial
$\varphi(0)+2\sum_j\varphi_j(0)z_j+\sum_{j,k}\varphi_{jk}(0)z_jz_k$.
Subtract it. This leaves the stated expansion and does not change $\partial\dbar$.
Let $\lambda>0$ and $M$ be the smallest and largest eigenvalues of the Hermitian quadratic term. On a smaller ball the remainder has absolute value at most $(\lambda/2)|z|^2$. We may therefore take $a=\lambda/2$ and $b=M+\lambda/2$.
\end{proof}

\begin{theorem}[A positive representative with null support, after Tosatti]
\label{pos:thin-current}
Let $(X,\omega)$ be compact and K\"ahler. There exists a continuous, locally Lipschitz real function $\Phi$ such that
\[
 T=\omega+\ii\partial\dbar\Phi
\]
is a closed positive $(1,1)$-current and $\supp T$ is contained in a closed null set $\Sigma$. Nullity is measured with respect to any smooth positive volume form, and $[T]=[\omega]$.
\end{theorem}
\begin{proof}
\emph{Step 1: fix a cover before choosing the constants.}
At each $x\in X$, Proposition~\ref{pos:smooth-local-potential} provides a smooth potential $\omega=\ii\partial\dbar\varphi_x$.
Lemma~\ref{pos:strict-real-minimum} normalizes it on a small coordinate ball centered at $x$ so that $\varphi_x(x)=0$ and $x$ is a strict real minimum. Denote the ball by $U_x$ and choose $a_x>0$ with $\varphi_x\ge a_x$ on $\partial U_x$. Its closure lies in a larger potential domain, so the corresponding strict lower bound also holds near the boundary. Put
\[
 V_x=\{z\in U_x:\varphi_x(z)<a_x/4\},\qquad
 J_x=(a_x/4,a_x/2).
\]
Each $V_x$ is open, contains $x$, and is relatively compact in $U_x$.
Compactness gives a finite subcover $V_1,\ldots,V_N$, with corresponding $U_i,\varphi_i,a_i,J_i$. This cover is now fixed: any later choice $c_i\in J_i$ preserves it.

\emph{Step 2: create local regions where the current vanishes.}
For constants $c_i\in J_i$ yet to be selected, define
\[
 \Phi_i=
 \begin{cases}
   \max\{c_i-\varphi_i,0\},&\text{on }U_i,\\
   0,&\text{on }X\setminus U_i.
 \end{cases}
\]
Since $c_i<a_i/2$, this function vanishes near $\partial U_i$. Its zero extension is therefore continuous and locally Lipschitz.
On $U_i$,
$\varphi_i+\Phi_i=\max\{c_i,\varphi_i\}$,
a maximum of two smooth psh functions, which is psh by Proposition~\ref{pos:psh-stability}. Thus $\Phi_i$ is $\omega$-psh there and also near its zero extension. Moreover, $\Phi_i>0$ on $V_i$ and, on the open set $\{\Phi_i>0\}$,
\[
 \omega+\ii\partial\dbar\Phi_i
 =\omega-\ii\partial\dbar\varphi_i=0.
\]

\emph{Step 3: avoid branches that tie identically.}
For $i<j$ and a connected component $W$ of $U_i\cap U_j$, the function $f=\varphi_i-\varphi_j$ satisfies $\partial\dbar f=0$ and is real analytic.
Here is a direct verification. The coefficients of $\partial f$ are holomorphic and this one-form is closed. On a small ball centered at zero, define
$G(z)=\int_0^1\sum_\ell f_{z_\ell}(tz)z_\ell\,dt$ and put $a_j=f_{z_j}$.
Using $\partial_ja_\ell=\partial_\ell a_j$, differentiate to obtain
\[
 \partial_{z_j}G
 =\int_0^1\left(a_j(tz)+t\sum_\ell z_\ell\partial_\ell a_j(tz)\right)dt
 =a_j(z).
\]
Thus $G$ is holomorphic and $\partial G=\partial f$. The real function $f-2\operatorname{Re}G$ has zero $\partial$ derivative and, by reality, zero $\dbar$ derivative. It is constant. Hence $f$ is locally real analytic.

If $f$ is constant with value $A$ on $W$, require $c_i-c_j\ne A$.
The manifold is second countable and locally connected, so $U_i\cap U_j$ has at most countably many connected components. Each component is open and contains a distinct member of a countable basis. Each pair $i,j$ therefore imposes at most countably many forbidden differences.
Choose $c_1,c_2,\ldots,c_N$ successively. When choosing $c_j$, the finitely many preceding indices exclude only countably many values $c_i-A$. They cannot exhaust the nonempty open interval $J_j$, so all requirements can be met.

\emph{Step 4: take the global maximum.}
With these constants fixed, let $\Phi=\max_i\Phi_i$. It is continuous and locally Lipschitz. On a domain with local potential $\varphi$,
$\varphi+\Phi=\max_i(\varphi+\Phi_i)$ is psh.
Consequently $T=\omega+\ii\partial\dbar\Phi$ is positive. Distributional differentiation gives $dT=0$, and $T-\omega$ is exact, so $[T]=[\omega]$.
Since the $V_i$ cover $X$, at least one $\Phi_i$ is positive at every point; therefore $\Phi>0$ everywhere.

\emph{Step 5: the tie set is null.}
Set
\[
 \Sigma_{ij}=\{x\in X:\Phi_i(x)=\Phi_j(x)=\Phi(x)\},\qquad
 \Sigma=\bigcup_{i<j}\Sigma_{ij}.
\]
Each set is closed by continuity, and the union is finite. Because $\Phi>0$, a tie can occur only in $U_i\cap U_j$, where it implies
\[
 \varphi_i-\varphi_j=c_i-c_j.
\]
On a connected component where the difference is constant, Step 3 makes this equation impossible. On every other component, it is the zero set of a nonzero real-analytic function and is Lebesgue-null by Proposition~\ref{bas:analytic-zeros}. Taking the countable union over components and the finite union over pairs proves nullity of $\Sigma$.
Smooth positive volume forms are mutually comparable on relatively compact coordinate domains, so nullity does not depend on the chosen volume form.

\emph{Step 6: the current vanishes off the tie set.}
At $x\notin\Sigma$ there is a unique maximizing branch, say $j$. Strict inequalities between finitely many continuous branches persist on a small neighborhood, so $\Phi=\Phi_j$ there. Since $\Phi_j(x)>0$, shrinking once more gives $\Phi_j=c_j-\varphi_j$. Step 2 then gives $T=0$ on this neighborhood. Thus $\supp T\subset\Sigma$, as required.
\end{proof}

Write $T_{\mathrm{ac}}$ for the form with locally integrable coefficients obtained by taking the absolutely continuous parts of the coefficient measures relative to smooth volume. The remaining coefficients are singular measures. Smooth coordinate changes preserve this decomposition, so it is intrinsic.
\begin{corollary}\label{pos:thin-mass}
If $n=\dim_\C X\ge1$, all coefficient measures of $T$ are singular with respect to smooth volume, and
\[
 T_{\mathrm{ac}}=0,\qquad
 \int_XT\wedge\omega^{n-1}=\int_X\omega^n>0.
\]
\end{corollary}
\begin{proof}
Theorem~\ref{pos:order-zero} identifies the coefficients of a positive current as measures. They are all supported on the same null set, so their absolutely continuous parts vanish.
Compactness and $d\omega=0$ allow distributional integration by parts:
$\langle\ii\partial\dbar\Phi,\omega^{n-1}\rangle=0$.
This proves the mass identity, whose right-hand side is positive because $\omega$ is K\"ahler.
\end{proof}

\par\penalty0
\begin{samepage}
\begin{remark}\label{pos:class-volume}
For a real $(1,1)$-class $\alpha$ on compact K\"ahler $X$, its positive-current volume is
\[
 \operatorname{Vol}(\alpha)=\sup_{S\ge0,\ [S]=\alpha}\int_X(S_{\mathrm{ac}})^n,
\]
where $S$ is closed and the value is zero if there is no positive representative. The absolutely continuous part is taken coefficientwise as above. The supremum, rather than the mass of one chosen representative, is essential in this definition; see \cite{Demailly2012}.
The support in Theorem~\ref{pos:thin-current} is null; the argument does not establish pluripolarity. A pluripolar set must locally lie in the $-\infty$ locus of a psh function that is not identically $-\infty$, whereas nullity is a volume condition.
Here the tie sets come from real-analytic level equations and may contain pieces of real codimension one. No potential taking the value $-\infty$ on them has been constructed. Thus this argument does not show that every K\"ahler class has a positive representative with pluripolar support.
Also, $T_{\mathrm{ac}}=0$ does not force the volume of its class to vanish. That volume takes a supremum over suitable positive representatives of the class; $T$ is just one representative.
\end{remark}
\end{samepage}


\chapter{Worked Examples: Triviality, Canonical Bundles, and Dimension}\label{exg:chapter}
\begin{learning}
A bundle can be trivial in one category and nontrivial in another. The examples below compare real, complex, smooth, and holomorphic bundles, then use canonical bundles to compute the Kodaira dimensions of curves, hypersurfaces, and products. The latter calculations can be read alongside Chapter~\ref{prj:chapter} before the more advanced applications.
\end{learning}

\section{The formula for real projective space}
\begin{proposition}[Complexified real tangent bundles and complex projective space]\label{exg:projective-chern}
Let $\lambda$ be the real tautological line bundle on $\mathbb R\mathbb P^N$, and set $a=c_1(\lambda\otimes_\R\C)$. Then
\[
 (T\mathbb R\mathbb P^N\otimes_\R\C)\oplus\C
 \simeq(\lambda\otimes_\R\C)^{\oplus(N+1)},\qquad
 c(T\mathbb R\mathbb P^N\otimes_\R\C)=(1+a)^{N+1}.
\]
Here $\C$ in a bundle direct sum denotes the trivial complex line bundle. Since $2a=0$, the rational or real Chern character is
$\operatorname{ch}(T\mathbb R\mathbb P^N\otimes\C)=N$.
On complex projective space, with $H=c_1(\mathcal O(1))$, we have
$c(T\mathbb P^N)=(1+H)^{N+1}$ and
$\operatorname{ch}(T\mathbb P^N)=(N+1)e^H-1$.
\end{proposition}
\begin{proof}
Fix a real line $\ell\subset\R^{N+1}$. Every nearby line is uniquely the graph of a small linear map $A:\ell\to\ell^\perp$, so the tangent space is naturally $\operatorname{Hom}_\R(\ell,\ell^\perp)$. Varying $\ell$ yields
$T\mathbb R\mathbb P^N\simeq\lambda^*\otimes\lambda^\perp$. Since
$\lambda^*\otimes\lambda\simeq\R$ and
$\lambda\oplus\lambda^\perp\simeq\R^{N+1}$, we obtain
$T\mathbb R\mathbb P^N\oplus\R\simeq(\lambda^*)^{\oplus(N+1)}$.
A real metric identifies $\lambda^*$ with $\lambda$, and complexification gives the first isomorphism.
Proposition~\ref{idx:integral-chern} and the direct-sum formula in Proposition~\ref{bun:chern-functors} give the total Chern class. The real metric trivializes $\lambda\otimes\lambda$, so the square of its complexification is trivial and $2a=0$. Torsion vanishes in rational cohomology; hence the Chern character $(N+1)e^a-1$ reduces to $N$.
The corresponding graph coordinates on complex projective space give the complex Euler sequence
$0\to\mathcal O\to\mathcal O(1)^{\oplus(N+1)}\to T\mathbb P^N\to0$.
The total Chern class is multiplicative in this sequence, whereas the Chern character is additive. This gives the two different formulas.
The real tangent bundle must be complexified before taking its Chern classes. The product $(1+a)^{N+1}$ is its total Chern class, whereas the exponential formula computes its Chern character.
\end{proof}

\section{Three classifications of line bundles}
Write $\operatorname{Pic}(X)$ for the group of holomorphic line bundles up to isomorphism, with tensor product as the group law. The sheaf $\mathcal O_X^*$ consists of nowhere-zero holomorphic functions; its star indicates units, not the dual of $\mathcal O_X$.
\begin{proposition}[What the exponential sequence implies]\label{exg:picard}
For a complex manifold $X$, the sequence
\[
 H^1(X,\mathcal O_X)\longrightarrow\operatorname{Pic}(X)
 \xrightarrow{c_1}H^2(X,\Z)\longrightarrow H^2(X,\mathcal O_X)
\]
is exact. If $H^{0,1}(X)=H^{0,2}(X)=0$, then
$c_1:\operatorname{Pic}(X)\to H^2(X,\Z)$ is an isomorphism.
On any smooth manifold, smooth complex line bundles and continuous complex line bundles are both classified by $H^2(X,\Z)$. Every continuous complex line bundle admits a smooth structure, unique up to smooth bundle isomorphism.
\end{proposition}
\begin{proof}
The kernel of $f\mapsto e^{2\pi\ii f}$ is the sheaf of locally constant integer-valued functions. Every nowhere-zero holomorphic function admits a local holomorphic logarithm. Hence
$0\to\Z\to\mathcal O_X\to\mathcal O_X^*\to1$ is exact on stalks.
Transition functions of line bundles are $\mathcal O_X^*$-valued one-cocycles, and changes of frame change them by a coboundary. Thus $H^1(\mathcal O_X^*)=\operatorname{Pic}(X)$. Proposition~\ref{sh:long-exact} gives the four displayed terms. Its connecting map is constructed from the integer differences of local logarithms on triple overlaps, so it is the first Chern class.
Theorem~\ref{sh:dolbeault} identifies the two outer groups with $H^{0,1}$ and $H^{0,2}$. Their vanishing gives injectivity and surjectivity, respectively.
Replace $\mathcal O_X$ by the sheaf of smooth complex-valued functions, or by the sheaf of continuous complex-valued functions. Local logarithms still exist. Both additive sheaves are fine by Proposition~\ref{sh:smooth-fine} and have no positive-degree cohomology by Theorem~\ref{sh:acyclic}, so the isomorphism classes of line bundles in either category are naturally identified with $H^2(X,\Z)$.
Under these identifications, forgetting smoothness induces the identity on $H^2(X,\Z)$ and is therefore bijective. This proves both existence and uniqueness of smoothing. The additional Dolbeault vanishing is a sufficient condition for holomorphic classification to agree with these classifications; without it, the exponential exact sequence must still be used.
\end{proof}

\begin{example}[A smoothly trivial but holomorphically nontrivial line bundle]\label{exg:flat-line}
Let $C=\C/(\Z+\tau\Z)$, where $\operatorname{Im}\tau>0$, and choose $\theta\in\R\setminus\Z$. The identifications
\[
 (z,v)\sim(z+1,v),\qquad (z,v)\sim(z+\tau,e^{2\pi\ii\theta}v)
\]
on $\C\times\C$ define a holomorphic line bundle $L_\theta$. Its transition factors have modulus one, so the standard Hermitian metric descends and has zero curvature.
Write $z=s+t\tau$ with $s,t\in\R$. The function $u(z)=e^{2\pi\ii\theta t}$ obeys the same transformation law and descends to a nowhere-zero smooth section. This directly proves smooth triviality without a classification theorem.
A nonzero holomorphic section would lift to an entire function $f$ with
$f(z+1)=f(z)$ and $f(z+\tau)=e^{2\pi\ii\theta}f(z)$. Its modulus is lattice-periodic and therefore bounded on the plane. Liouville's theorem makes $f$ constant, and the second transformation law forces that constant to vanish. This contradiction shows that $L_\theta$ has no nonzero holomorphic section and in particular is not holomorphically trivial.
For comparison, a principal divisor always defines a holomorphically trivial line bundle: multiplication by $f$ identifies $\mathcal O(\operatorname{div}f)$ with $\mathcal O_X$. The unitary transition factors in this example instead give a nontrivial holomorphic structure.
\end{example}

\begin{proposition}[Zero Chern class and a nonzero section]\label{exg:zero-class}
Let $X$ be connected, compact, and K\"ahler. If $c_1(L)=0\in H^2(X,\R)$ and $L$ has a nonzero holomorphic section, then $L$ is holomorphically trivial. The K\"ahler hypothesis cannot be omitted for general compact complex manifolds.
\end{proposition}
\begin{proof}
In dimension zero, $X$ is a point and the result is immediate. Suppose $n=\dim_\C X\ge1$. Let $D$ be the zero divisor of a nonzero section $s$, and let $\omega$ be a K\"ahler form. Theorem~\ref{cur:poincare-lelong} and Proposition~\ref{pos:current-calculus} give
\[
 \int_D\omega^{n-1}=\int_Xc_1(L)\wedge\omega^{n-1}=0.
\]
If $D$ is nonempty, its regular locus contains a nonempty open patch whose volume is strictly positive. The integral on the left would be positive, a contradiction. Thus $s$ never vanishes and gives a holomorphic trivialization $\mathcal O_X\simeq L$.

For a counterexample without the K\"ahler assumption, take the Hopf surface
$X=(\C^2\setminus\{0\})/(z\sim2z)$.
The dilation action is free and properly discontinuous: a compact subset meets only finitely many of its dilates. The compact fundamental annulus $1\le|z|\le2$ proves compactness. Radial logarithmic coordinates give a diffeomorphism $X\simeq S^3\times S^1$.
Define a line bundle by $(z,v)\sim(2z,2v)$. The positive smooth function $|z|$ is equivariant and gives a smooth trivialization, so $c_1(L)=0$. The holomorphic equivariant function $z_1$ descends to a nonzero section with zeros along $z_1=0$. If $L$ were holomorphically trivial, this section would correspond to a holomorphic function on a connected compact complex manifold. Such a function is constant and cannot be both nonzero and have zeros. Thus $L$ is not holomorphically trivial.
\end{proof}

\section{Canonical bundles and all possible Kodaira dimensions}
\begin{proposition}[The three Kodaira dimensions of compact curves]\label{exg:curve-kodaira}
A connected compact Riemann surface has Kodaira dimension $-\infty$, $0$, or $1$ according as its genus is $0$, $1$, or at least $2$.
\end{proposition}
\begin{proof}
Theorem~\ref{cur:rr} gives $\deg K_C=2g-2$, with $g=h^1(C,\mathcal O_C)$ the genus. If $g=0$, every positive power $K_C^m$ has negative degree and no nonzero section. Hence $\kappa=-\infty$ by definition.
If $g=1$, Theorem~\ref{hod:serre-duality} gives $h^0(K_C)=1$. A nonzero one-form has a zero divisor of degree zero and therefore has no zeros. Thus $K_C$ is trivial, every power has a one-dimensional section space, and $\kappa=0$.
If $g\ge2$, then $K_C$ has positive degree and $K_C^{1-m}$ has negative degree for $m\ge2$. Duality and Riemann--Roch give
$h^0(K_C^m)=m(2g-2)+1-g=(2m-1)(g-1)$.
The growth is linear, so $\kappa=1$. This calculation uses no higher-dimensional numerical criterion.
\end{proof}

\begin{proposition}[Adjunction and the Kodaira dimension of a hypersurface]\label{exg:hypersurface}
Let $X\subset\mathbb P^N$ be a connected smooth hypersurface of degree $d$, with $N\ge2$. Then
$K_X\simeq\mathcal O_X(d-N-1)$.
Its Kodaira dimension is $-\infty$, $0$, or $N-1$ according as $d<N+1$, $d=N+1$, or $d>N+1$. Also $\kappa(\mathbb P^N)=-\infty$.
\end{proposition}
\begin{proof}
Let $N_{X/\mathbb P^N}=T_{\mathbb P^N}|_X/T_X$ be the normal bundle. Locally write $X=(f=0)$. Smoothness means that $df$ does not vanish on $X$, and differentiation in the normal direction gives
$0\to T_X\to T_{\mathbb P^N}|_X\to\mathcal O_X(d)\to0$.
Check the sign of the last bundle through the conormal bundle. Local defining equations generate $\mathcal I_X=\mathcal O(-d)$, and differentiation induces
$\mathcal I_X/\mathcal I_X^2\simeq N_{X/\mathbb P^N}^*$.
If $f_i=g_{ij}f_j$, then $df_i=g_{ij}df_j$ on $X$, since $f_j\,dg_{ij}$ vanishes there. Thus the differentials generate the conormal bundle $\mathcal O_X(-d)$; dualizing gives the normal bundle $\mathcal O_X(d)$.
Taking determinants and dualizing the tangent sequence yields
$K_X=K_{\mathbb P^N}|_X\otimes\mathcal O_X(d)$.
Proposition~\ref{amp:euler} gives $K_{\mathbb P^N}=\mathcal O(-N-1)$, proving adjunction.
If $d-N-1<0$ and $K_X^m$ had a nonzero section, its zero divisor would have nonnegative integral against $\omega_{\rm FS}^{N-2}$, whereas the corresponding Chern-class integral is strictly negative. This is impossible.
If $d-N-1=0$, the canonical bundle is trivial. Connected compactness makes every pluricanonical section constant, so $\kappa=0$.
If $d-N-1>0$, the canonical bundle is ample; section growth, or a projective embedding by a positive power, gives $\kappa=\dim X=N-1$.
The same negative-degree argument applied to $K_{\mathbb P^N}$ proves the assertion for projective space.
The restriction $N\ge2$ is intentional: a hypersurface in $\mathbb P^1$ is a zero-dimensional set, to which this trichotomy must not be applied mechanically.
\end{proof}

\begin{corollary}[Examples of every allowed Kodaira dimension]\label{exg:all-kodaira}
For every $n\ge1$ and $d\in\{-\infty,0,1,\ldots,n\}$ there exists a compact complex $n$-fold $X$ with $\kappa(X)=d$.
\end{corollary}
\begin{proof}
For $d=-\infty$ take $\mathbb P^n$. Let $E$ be an elliptic curve and $C$ a curve of genus at least two. One may take a smooth plane quartic: the adjunction formula just proved, together with the degree formula for curves, gives genus three.
For $d\ge0$ set $X=E^{n-d}\times C^d$.
The canonical bundle of a product is the exterior tensor product of the canonical bundles of its factors, and Proposition~\ref{gp:kodaira-product} identifies its section space with the tensor product of theirs. Each elliptic factor contributes dimension one, and each higher-genus factor contributes linear growth in $m$. Thus the pluricanonical section dimension grows as $m^d$, giving $\kappa=d$.
\end{proof}

\begin{corollary}[The very ample polarization case]\label{exg:fujita-very-ample}
If $X$ is a smooth projective $n$-fold and $L$ is very ample, then $K_X+(n+1)L$ is globally generated and $K_X+(n+2)L$ is very ample.
\end{corollary}
\begin{proof}
Put $V=K_X+(n+1)L$. For $1\le q\le n$,
$H^q(X,V-qL)=H^q(X,K_X+(n+1-q)L)=0$,
since the coefficient of $L$ is at least one and Corollary~\ref{pos:kodaira} applies. For $q>n$, the Dolbeault resolution has no terms in that degree, so the groups also vanish. Lemma~\ref{amp:regularity}, with $A=L$, makes $V$ globally generated.
The bundle $V$ defines a map to projective space, while $L$ defines an embedding. Their product is an embedding because the second component already separates points and tangent vectors. Composing with the Segre embedding pulls back the hyperplane bundle to
$V\otimes L=K_X+(n+2)L$, proving very ampleness.
The proof uses very ampleness of $L$; the same bounds for an arbitrary ample $L$ are the subject of Fujita's conjecture. On $\mathbb P^n$, the bundle $K+(n+1)\mathcal O(1)$ is trivial and globally generated, so ``globally generated'' must not be replaced by ``ample.''
\end{proof}


\chapter{Sixteen Tutorials with Complete Solutions}\label{tutorial:chapter}
\begin{learning}
These problems bring together the definitions and methods developed in the course. Try each problem before reading its solution, and return to the cited results when a hypothesis or a normalization is unclear. Problems 1, 2, and 4--6 accompany the first semester. The cohomology calculation in Problem 3 also belongs to that route; its comparison of nefness, pseudoeffectivity, and bigness uses Chapter~\ref{gp:chapter}. Problems 7--16 provide second-semester examples and further practice, with the monomial and curve calculations offering the most direct starting points.
\end{learning}

\section{First semester: local data and global obstructions}
\subsection{Problem 1: Laurent series on the projective line}\label{tut:line-cech}
\paragraph{Problem.} Compute $H^0(\mathbb P^1,\mathcal O(d))$ and $H^1(\mathbb P^1,\mathcal O(d))$ for $d\in\Z$, without Riemann--Roch.
\paragraph{Solution.} Use $U_0=\{Z_0\ne0\}$ and $U_1=\{Z_1\ne0\}$, with coordinates $z=Z_1/Z_0$ and $w=1/z$. The formal local frames $e_0=Z_0^d$, $e_1=Z_1^d$ satisfy $e_1=z^de_0$. A global section is a pair of entire functions with
$f_0(z)=z^df_1(1/z)$.
For $d\ge0$, this says exactly that $f_0$ is a polynomial of degree at most $d$; for $d<0$ it forces zero.
We first justify computing cohomology on this fixed cover. For $\Omega=\C$ or $\C^*$, every smooth $(0,1)$-form $f\,d\bar z$ has a global smooth primitive.
Take an exhaustion $K_j$ by discs for $\C$, or compact annuli for $\C^*$. Choose cutoffs $\chi_j$ equal to one near $K_j$, supported in $K_{j+1}$, and put $\chi_0=0$.
Then $f_j=(\chi_j-\chi_{j-1})f$ is compactly supported and $f=\sum_jf_j$ locally finitely. Extend $f_j$ by zero to the plane and apply Lemma~\ref{sh:cauchy-operator} to obtain $u_j=Tf_j$ with $\dbar u_j=f_j\,d\bar z$.
For large $j$, $u_j$ is holomorphic near $K_{j-1}$. A finite Taylor or Laurent truncation $p_j$ satisfies
$\|u_j-p_j\|_{C^j(K_{j-2})}<2^{-j}$.
Each truncation is holomorphic on all of $\Omega$. Thus $u=\sum_j(u_j-p_j)$ converges with every derivative on compact sets and solves $\dbar u=f\,d\bar z$. For the finitely many initial terms take $p_j=0$.
Theorem~\ref{sh:dolbeault} now gives $H^q(\Omega,\mathcal O)=0$ for $q>0$. The bundle is trivial on both charts and their intersection, so Theorem~\ref{sh:leray} allows this \Cech{} calculation.

A one-cochain is a holomorphic function on $\C^*$, written $g(z)=\sum_{j\in\Z}a_jz^j$. A coboundary allows adding an entire function of $z$, or $z^d$ times an entire function of $w$. The first removes all exponents $j\ge0$, the second all $j\le d$.
The positive and negative Laurent tails converge on the entire $z$- and $w$-planes, respectively, because the original series converges normally on every compact annulus in $0<|z|<\infty$.
Only exponents $d<j<0$ remain. Therefore
\[
 h^0(\mathcal O(d))=\max(d+1,0),\qquad
 h^1(\mathcal O(d))=\max(-d-1,0).
\]
For $d\le-2$, the classes of $z^{d+1},\ldots,z^{-1}$ form a basis of $H^1$. There are no two-cochains, so higher cohomology vanishes. The remaining Laurent terms explicitly exhibit the gluing obstruction for negative line bundles.

\subsection{Problem 2: why the Euler sequence does not split holomorphically}\label{tut:problem-2}
\paragraph{Problem.} Show that
$0\to\mathcal O\to\mathcal O(1)^{\oplus2}\to\mathcal O(2)\to0$
on $\mathbb P^1$ does not split holomorphically, and explain why smooth splitting is still possible.
\paragraph{Solution.} A holomorphic splitting would identify the middle term with $\mathcal O\oplus\mathcal O(2)$. Tensoring by $\mathcal O(-2)$ would give
$\mathcal O(-1)^{\oplus2}\simeq\mathcal O(-2)\oplus\mathcal O$.
By \hyperref[tut:line-cech]{Problem 1}, the left side has no global sections and the right side has a one-dimensional section space, a contradiction.
For smooth splitting choose a Hermitian metric. The orthogonal complement of the subbundle $\mathcal O$ is a smooth line bundle, and the quotient map restricts to a smooth bundle isomorphism on this complement. In general the complement is not preserved by $\dbar$, so it need not be holomorphic. A Chern-class direct-sum identity does not imply holomorphic splitting.

\subsection{Problem 3: positivity and cohomology on a product of projective lines}\label{tut:problem-3}
\paragraph{Problem.} On $X=\mathbb P^1\times\mathbb P^1$, let $\pi_1,\pi_2$ be the projections and $\mathcal O(a,b)=\pi_1^*\mathcal O(a)\otimes\pi_2^*\mathcal O(b)$. For $L=\mathcal O(a,b)$, determine ampleness, nefness, pseudoeffectivity, and bigness. Compute the cohomology of $\mathcal O(2,-3)$.
\paragraph{Solution.} Normalize the Fubini--Study forms $\omega_1,\omega_2$ of the two factors to have integral one. The form $a\omega_1+b\omega_2$ represents $c_1(L)$. It is positive definite for $a,b>0$, giving ampleness, and semipositive for $a,b\ge0$, giving nefness and pseudoeffectivity.
Conversely, degrees on the two fiber classes are $a,b$. Ampleness forces them positive; nefness forces them nonnegative. If $T\in c_1(L)$ is a positive closed current, its wedges with the smooth semipositive forms $\omega_i$ are positive measures, so
\[
 a=\int_XT\wedge\omega_2\ge0,\qquad
 b=\int_XT\wedge\omega_1\ge0.
\]
These equalities are cohomological: an exact current pairs to zero with a smooth closed form by Stokes, while
$\omega_1^2=\omega_2^2=0$ and $\int_X\omega_1\wedge\omega_2=1$.
No restriction of a singular current to an individual fiber is needed.
Global sections are tensor products of sections on the factors, hence
\[
 h^0(X,L^k)=\max(ka+1,0)\max(kb+1,0)\quad(k\ge1).
\]
Quadratic growth occurs exactly when $a,b>0$, which is also the bigness condition.
For bundles on the two factors, write $A\boxtimes B=\pi_1^*A\otimes\pi_2^*B$. For the cohomological tensor decomposition, use product metrics and Theorem~\ref{hod:hodge}. The operator $\Delta''$ is the sum of the two factor operators. Its nonnegative energy vanishes exactly when both factor energies vanish, and tensor products of orthogonal eigenbases form a basis on the product. Thus
$H^q(X,A\boxtimes B)=\bigoplus_{i+j=q}H^i(A)\otimes H^j(B)$.
\hyperref[tut:line-cech]{Problem 1} gives $h^0(\mathcal O(2))=3$ and $h^1(\mathcal O(-3))=2$, with the other relevant groups zero. Therefore $(h^0,h^1,h^2)=(0,6,0)$ and the Euler characteristic is $-6$.

\subsection{Problem 4: local free resolutions do not imply global vanishing}\label{tut:problem-4}
\paragraph{Problem.} Find the error in: ``A coherent sheaf has a finite local free resolution, so its positive-degree cohomology vanishes.''
\paragraph{Solution.} Local freeness refers to small coordinate neighborhoods. Local resolutions and their contractions need not agree on overlaps. More directly, $\mathcal O(-2)$ is already locally free of rank one and resolves itself in length zero, but \hyperref[tut:line-cech]{Problem 1} gives
$H^1(\mathbb P^1,\mathcal O(-2))\simeq\C$.
Theorem~\ref{sh:resolution-theorem} computes cohomology by global sections when the resolution terms are acyclic for that functor; Theorem~\ref{alg:regular-resolution} asserts only local freeness. Local freeness does not imply acyclicity. Fine resolutions work because a separate partition-of-unity argument proves that their terms are acyclic.

\subsection{Problem 5: a minimal-norm Dolbeault solution}\label{tut:problem-5}
\paragraph{Problem.} On $\C$ with weight $e^{-a|z|^2}$, $a>0$, and the model normalization $|d\bar z|=1$, find the minimal-norm solution of $\dbar u=d\bar z$ and check the estimate's constant.
\paragraph{Solution.} Since $\partial\bar z/\partial\bar z=1$, the function $u_0=\bar z$ is a solution. Every other solution is $\bar z+f(z)$ with $f$ holomorphic. For $j\ge0$, angular integration gives
$\int_\C z^{j+1}e^{-a|z|^2}\,dx\,dy=0$,
so $\bar z$ is orthogonal to all holomorphic polynomials.
Taylor monomials of a holomorphic weighted-$L^2$ function are orthogonal in this radial inner product. Termwise angular integration and monotone convergence show that its Taylor partial sums converge in $L^2$. Thus $\bar z$ is orthogonal to the whole holomorphic kernel. Pythagoras proves it is the unique minimal-norm solution.
Direct integration gives
\[
 \|u_0\|^2=2\pi\int_0^\infty r^3e^{-ar^2}\,dr=\frac\pi{a^2},\qquad
 \|d\bar z\|^2=2\pi\int_0^\infty re^{-ar^2}\,dr=\frac\pi a.
\]
Hence $\|u_0\|^2=a^{-1}\|d\bar z\|^2$, attaining the constant $1/a$ in Theorem~\ref{pos:euclidean-solving}; compare Example~\ref{pos:gaussian-example}. Changing the geometric normalization of $d\bar z$ requires changing the adjoint and curvature-operator constants consistently.

\subsection{Problem 6: failure without closed range}\label{tut:problem-6}
\paragraph{Problem.} On the Hilbert space $\ell^2=\{x=(x_j)_{j\ge1}:\sum_j|x_j|^2<\infty\}$ let $T(x_1,x_2,\ldots)=(x_1,x_2/2,\ldots)$. Show that $T$ is injective with dense nonclosed range, and explain the role of a basic estimate.
\paragraph{Solution.} Coordinatewise, $\ker T=0$. Every finitely supported sequence is in the range, so the range is dense. The vector $y=(1,1/2,1/3,\ldots)$ belongs to $\ell^2$, but $Tx=y$ would force $x_j=1$ for every $j$, which is not square-summable. The finite truncations of $y$ are in the range and converge to $y$, proving nonclosedness.
For the standard basis, $\|e_j\|=1$ but $\|Te_j\|=1/j$. No uniform estimate $\|u\|\le C\|Tu\|$ can hold. In geometric arguments a curvature lower bound provides precisely such uniform control, ensuring solvability as in Corollary~\ref{pos:coercive-existence}; density of the range alone does not give it.

\Needspace{12\baselineskip}
\section{Second semester: singularities, heat kernels, and effective bounds}
\subsection{Problem 7: integral closure of two monomial ideals}\label{tut:problem-7}
\paragraph{Problem.} Determine the integral closure of $I=(x^2,y^3)\subset\C\{x,y\}$ and show why the one-power loss in dimension-two Brian\c{c}on--Skoda cannot be removed altogether.
\paragraph{Solution.} The necessary Newton condition in Corollary~\ref{bs:support} is $u/2+v/3\ge1$. Its minimal integral exponents under coordinatewise divisibility are $(2,0),(1,2),(0,3)$, so Lemma~\ref{bs:analytic-division} gives
$\overline I\subset(x^2,xy^2,y^3)$.
For the reverse inclusion, $xy^2$ is integral: it satisfies $T^2-x^2y^4=0$ and $x^2y^4\in I^2$. Integral closure is an ideal by the growth characterization in Theorem~\ref{div:integral-closure}, so all three generators and their analytic combinations lie in $\overline I$. Thus
$\overline I=(x^2,xy^2,y^3)$.
The new monomial is not in $I$, since neither $x^2$ nor $y^3$ divides it.
For $J=(x^2,y^2)$, one similarly has $xy\in\overline J\setminus J$. Hence the uniform inclusion $\overline{J^{k+1}}\subset J^k$ cannot be replaced by $\overline{J^k}\subset J^k$.
At $k=1$, every monomial in $\overline{J^2}$ has $u+v\ge4$. If both $u,v<2$, their sum is at most two, a contradiction. Thus one exponent is at least two and the monomial belongs to $J$. This verifies the inclusion and exhibits the reason for the exponent loss.

\subsection{Problem 8: equality at an integrability threshold}\label{tut:problem-8}
\paragraph{Problem.} Compute the multiplier ideals of
$\varphi=a\log|z_1|^2+b\log|z_2|^2$, $a,b\ge0$, and
$\psi=c\log(|z_1|^2+|z_2|^2)$, $c\ge0$.
\paragraph{Solution.} The weighted integral of $z_1^uz_2^v$ is a product of radial integrals. The first is finite exactly when $2u-2a+1>-1$, or $u>a-1$. The smallest permissible integer is $\lfloor a\rfloor$; similarly the second is $\lfloor b\rfloor$.
Angular orthogonality makes the integral of a general holomorphic function the sum of the nonnegative monomial integrals. Consequently
\[
 \mathcal I(\varphi)=(z_1^{\lfloor a\rfloor}z_2^{\lfloor b\rfloor}).
\]
For the radial weight, write $\mathfrak m=(z_1,z_2)$ for the local maximal ideal at the origin and suppose the first nonzero homogeneous term of $f$ has degree $m$. Its squared integral on the unit sphere is positive, so the spherical average is $Cr^{2m}+O(r^{2m+1})$ with $C>0$.
In real dimension four, integrability is therefore equivalent to that of
$\int_0^\epsilon r^{2m-2c+3}\,dr$, namely $m>c-2$. Thus
\[
 \mathcal I(\psi)=\mathfrak m^{\max(0,\lfloor c\rfloor-1)}.
\]
These are the special cases of Theorem~\ref{pos:snc-ideal} and Proposition~\ref{pos:radial-ideal}. At $c=2$, for example, the ideal is already $\mathfrak m$. Replacing the strict inequality by a non-strict one would incorrectly admit logarithmic divergence at the endpoint.

\subsection{Problem 9: where the mass goes when a logarithmic pole is smoothed}\label{tut:problem-9}
\paragraph{Problem.} For $\varphi_\epsilon(z)=\log(|z|^2+\epsilon^2)$, compute its curvature and prove its weak limit.
\paragraph{Solution.} Differentiation yields
\[
 \ii\partial\dbar\varphi_\epsilon
 =\frac{\ii\epsilon^2\,dz\wedge d\bar z}{(|z|^2+\epsilon^2)^2}
 =\frac{2\epsilon^2}{(|z|^2+\epsilon^2)^2}\,dx\,dy.
\]
This positive measure has mass
$4\pi\epsilon^2\int_0^Rr(r^2+\epsilon^2)^{-2}\,dr
=2\pi R^2/(R^2+\epsilon^2)$
on the disc of radius $R$, tending to $2\pi$.
For fixed $\delta>0$, its mass on $|z|\ge\delta$ is at most $2\pi\epsilon^2/\delta^2$, tending to zero.
For a compactly supported continuous test function $f$, split the integral into $|z|<\delta$ and its complement. On the first set, control $f-f(0)$ by its modulus of continuity; on the second, use the mass bound. Let $\epsilon\to0$ first and then $\delta\to0$. The limit is $2\pi f(0)$, proving weak convergence to $2\pi\delta_0$, where $\delta_0$ is the unit point mass at the origin. This agrees with Lemma~\ref{pos:one-variable-lelong}.
Each smoothed curvature is smooth, yet its maximum density is $2/\epsilon^2$. Smoothness for every fixed parameter does not provide a uniform curvature bound.

\subsection{Problem 10: Hodge numbers and heat traces of a flat complex torus}\label{tut:problem-10}
\paragraph{Problem.} For the standard flat metric on
$X=\C^n/(\Z^n+\ii\Z^n)$, compute $h^{p,q}$ and $\operatorname{Tr}(e^{-t\Delta''})$.
\paragraph{Solution.} Write $z_j=x_j+\ii y_j$. The Fourier basis is
$e^{2\pi\ii(a\cdot x+b\cdot y)}$, $(a,b)\in\Z^{2n}$.
For $\omega=\frac\ii2\sum dz_j\wedge d\bar z_j$, Proposition~\ref{hod:kahler-laplacians} gives $\Delta''=\frac12\Delta_d$ and eigenvalues $2\pi^2(|a|^2+|b|^2)$. There are $\binom np\binom nq$ form-basis elements $dz_I\wedge d\bar z_J$, and the operator acts coefficientwise.
Only constant coefficients give eigenvalue zero, so $h^{p,q}=\binom np\binom nq$. Summing all eigenvalues gives
\[
 \operatorname{Tr}_{p,q}(e^{-t\Delta''})
 =\binom np\binom nq\sum_{m\in\Z^{2n}}e^{-2\pi^2t|m|^2}.
\]
We can prove the required Gaussian summation identity directly.
The periodic function $G_t(x)=\sum_{j\in\Z}e^{-2\pi^2t(x+j)^2}$ converges uniformly with every derivative. Translating the integration intervals gives its $\ell$th Fourier coefficient:
\[
 \widehat G_t(\ell)=\int_\R e^{-2\pi^2tu^2}e^{-2\pi\ii\ell u}\,du
 =(2\pi t)^{-1/2}e^{-\ell^2/(2t)}.
\]
To evaluate this integral, put $F(\xi)=\int_\R e^{-au^2}e^{-2\pi\ii\xi u}\,du$.
Integration by parts gives $F'(\xi)=-(2\pi^2\xi/a)F(\xi)$, while polar coordinates applied to $F(0)^2$ give $F(0)=\sqrt{\pi/a}$. Solve this first-order equation and take $a=2\pi^2t$.
The Fourier series converges absolutely. At $x=0$ it gives
\[
 \sum_{j\in\Z}e^{-2\pi^2tj^2}
 =(2\pi t)^{-1/2}\sum_{j\in\Z}e^{-j^2/(2t)}.
\]
Taking $2n$ products gives the small-time leading term
$\binom np\binom nq(2\pi t)^{-n}$.
Its relative error is $O(e^{-1/(2t)})$: a nonzero integral vector has squared length at least one, and the remaining tail is dominated by a fixed summable Gaussian series.
As $t\to\infty$, every nonzero-frequency term vanishes. For $t\ge1$, the summable majorant $e^{-2\pi^2|m|^2}$ permits termwise passage to the limit. Only the constant term remains, recovering the Hodge dimension.

\subsection{Problem 11: how curvature signs select a cohomology degree}\label{tut:problem-11}
\paragraph{Problem.} Let $\lambda_1,\ldots,\lambda_n$ be real curvature eigenvalues and $t>0$. In the sum below, $J$ ranges over subsets of $\{1,\ldots,n\}$ of size $q$. Analyze, factor by factor,
\[
 B_q(t)=\frac1{(2\pi)^n}\sum_{|J|=q}
 e^{-2t\sum_{j\in J}\lambda_j}
 \prod_{j=1}^n\frac{\lambda_j}{1-e^{-2t\lambda_j}},
\]
interpreting a zero-eigenvalue factor as $1/(2t)$.
\paragraph{Solution.} For $\lambda=a>0$, an unselected factor tends to $a$; selecting it multiplies by $e^{-2ta}$ and gives zero. For $\lambda=-a<0$, the basic factor is $a/(e^{2ta}-1)$, tending to zero when unselected. Selecting it contributes $e^{2ta}$ and makes the limit $a$.
Without zero eigenvalues, the only possibly nonzero limiting term therefore selects exactly all negative eigenvalues. If their number is $q$, the limit is $(2\pi)^{-n}\prod_j|\lambda_j|$; otherwise it is zero.
A zero eigenvalue contributes $1/(2t)$, while every other factor with its selected or unselected exponential remains bounded as $t\to\infty$. Thus every term tends to zero in that case.
This explains the index domain $X(q)$ of Theorem~\ref{morse:main}: the nondegenerate curvature locus with exactly $q$ negative eigenvalues. Compare Proposition~\ref{heat:model-limit}. If $\lambda$ also tends to zero, the rate depends on $t|\lambda|$, so exchanging the long-time and degenerating-parameter limits requires additional estimates.

\subsection{Problem 12: the remainder in the strong Morse inequality}\label{tut:problem-12}
\paragraph{Problem.} Let $C^0\to C^1\to C^2$ be a complex with dimensions $(2,3,1)$ and both differential ranks equal to one. Compute its cohomology and identify the strong-inequality remainder.
\paragraph{Solution.} The first kernel has dimension $2-1=1$. The second kernel has dimension $3-1=2$, and quotienting by the one-dimensional incoming image leaves $h^1=1$. The final cokernel has dimension $1-1=0$. Hence $(h^0,h^1,h^2)=(1,1,0)$.
In general, put $c_j=\dim C^j$, $h_j=\dim H^j(C^\bullet)$, $r_j=\operatorname{rank}(d_j:C^j\to C^{j+1})$, and $r_{-1}=0$. As in \eqref{heat:boundary-remainder},
$c_j=h_j+r_{j-1}+r_j$.
Alternating summation cancels the internal ranks:
\[
 \sum_{j=0}^q(-1)^{q-j}(c_j-h_j)=r_q\ge0.
\]
Here at $q=1$ the left side is $(3-1)-(2-1)=1$, the uncancelled $r_1$. At $q=2$ it is zero, giving equality of Euler characteristics.
The heat-trace version repeats this identity on each positive-eigenvalue subcomplex. Its sign is supplied by finite-dimensional algebra, not by a special feature of the heat-kernel formula.

\subsection{Problem 13: multipoint jets and an exact threshold on the projective line}\label{tut:problem-13}
\paragraph{Problem.} At a nonempty finite set of distinct points $x_i\in\mathbb P^1$, prescribe jets of orders $s_i\ge0$. Show that $\mathcal O(d)$ realizes arbitrary such data exactly when $d\ge\sum_i(s_i+1)-1$.
\paragraph{Solution.} Choose an affine coordinate making all $x_i$ finite. Put $S=\sum_i(s_i+1)$ and
$g_i(z)=\prod_{j\ne i}(z-x_j)^{s_j+1}$.
For the prescribed Taylor polynomial $f_i$ at $x_i$, the quotient $f_i/g_i$ is holomorphic there because $g_i(x_i)\ne0$. Let $u_i$ be its Taylor truncation through degree $s_i$ and set $F=\sum_i g_i u_i$.
At $x_i$ its $i$th summand equals $f_i$ modulo $(z-x_i)^{s_i+1}$; every other summand has the required vanishing factor. Thus $F$ realizes all jets.
Each summand has degree at most $S-(s_i+1)+s_i=S-1$, proving sufficiency.
Conversely the target jet space has dimension $S$, whereas \hyperref[tut:line-cech]{Problem 1} gives $h^0(\mathcal O(d))=\max(d+1,0)$. For $d<S-1$, evaluation cannot be surjective. This model directly relates the number of jet conditions to the effective bound.

\subsection{Problem 14: why pseudoeffective twisting cannot replace nef twisting}\label{tut:problem-14}
\paragraph{Problem.} On $X=\operatorname{Bl}_p\mathbb P^2$, let $E$ be the exceptional curve and $L$ ample. For fixed $m$, construct a pseudoeffective $G$ such that $2K_X+mL+G$ is not globally generated.
\paragraph{Solution.} First check the intersection signs. The blowup charts are $(x,t)$ and $(s,y)$ with $y=tx$, $x=sy$, and $s=1/t$ on their overlap.
The exceptional curve is $x=0$ or $y=0$. In the normal quotient along it the frames satisfy $\partial_x=t\partial_y$, so its normal bundle is $\mathcal O_{\mathbb P^1}(-1)$ and $E^2=-1$.
Taking determinant duals in
$0\to T_E\to T_X|_E\to N_{E/X}\to0$
gives $K_X|_E=K_E\otimes N_{E/X}^{-1}=\mathcal O(-1)$ and hence $K_X\cdot E=-1$. These are the intersection numbers used in Example~\ref{eff:pseudo-counterexample}.
Put $d=L\cdot E>0$, choose an integer $N>\max\{0,md-2\}$, and let $G=NE$. It is effective and therefore pseudoeffective, but
\[
 (2K_X+mL+NE)\cdot E=-2+md-N<0.
\]
A globally generated line bundle restricts to a globally generated bundle on $E$. Choosing a section nonzero at one point, its effective zero divisor shows that this restricted bundle must have nonnegative degree. The negative intersection is a contradiction.
Since $G\cdot E=-N$, the twisting bundle is not nef. The example identifies a genuine role of the nef hypothesis in an effective jet theorem, rather than merely a limitation of one proof technique.

\subsection{Problem 15: a rank-zero coherent sheaf can have many sections}\label{tut:problem-15}
\paragraph{Problem.} For a line $H\subset\mathbb P^2$, compare the section growth of $\mathcal O_H(k)$ and of the double line $\mathcal O_{2H}(k)$.
\paragraph{Solution.} The sheaf $\mathcal O_H$ has rank zero on $\mathbb P^2$, but $h^0(\mathcal O_H(k))=k+1$ for $k\ge0$. Rank zero in a dimension-two Morse leading term means only $o(k^2)$ growth, not absence of sections.
The double line is defined by $\mathcal I_H^2$, and there is an exact sequence
\[
 0\longrightarrow\mathcal O_H(k-1)\longrightarrow
 \mathcal O_{2H}(k)\longrightarrow\mathcal O_H(k)\longrightarrow0.
\]
It comes from $\mathcal I_H/\mathcal I_H^2\simeq\mathcal O_H(-1)$ after twisting by $k$. For $k\ge1$, the outer $H^1$ groups vanish, so
$h^0(\mathcal O_{2H}(k))=k+(k+1)=2k+1$.
The leading coefficient two records the generic length of the nonreduced structure, exactly as in Theorem~\ref{idx:hilbert-polynomial}. Coherent-sheaf Morse estimates require a filtration by powers of the radical ideal: support on a line does not make the sheaf a vector bundle on the reduced line.

\par\penalty0
\begin{samepage}
\subsection{Problem 16: symmetric-power growth includes the rank}\label{tut:problem-16}
\paragraph{Problem.} On $\mathbb P^1$ let $E=\mathcal O(1)\oplus\mathcal O(2)$. For every integer $m\ge0$, compute the rank, first Chern class, and section dimension of $\operatorname{Sym}^mE$, and compare the growth degree with the general formula.
\paragraph{Solution.} A symmetric monomial contains $m-j$ copies of the first factor and $j$ of the second, giving
\[
 \operatorname{Sym}^mE=\bigoplus_{j=0}^m\mathcal O(m+j).
\]
The rank is $m+1$ and the degree is $\sum_{j=0}^m(m+j)=3m(m+1)/2$. This agrees with Theorem~\ref{bun:symmetric-chern}:
$c_1(\operatorname{Sym}^mE)=\frac{m(m+1)}2c_1(E)$ because $\deg E=3$.
Higher cohomology of every summand vanishes, and $h^0(\mathcal O(m+j))=m+j+1$. Thus
\[
 h^0(\operatorname{Sym}^mE)=\frac{3m(m+1)}2+(m+1)
 =\frac32m^2+\frac52m+1.
\]
The base has dimension one, but growth is quadratic. The relevant geometric dimension is $\dim\mathbb P(E^*)=1+2-1=2$, as explained in Proposition~\ref{amp:twist-growth}. Together with \hyperref[tut:problem-15]{Problem 15}, this shows why growth must be attached to the actual sheaf and its support, rather than inferred only from base dimension or generic rank.
\par
\end{samepage}


\chapter{Further Questions and Student Projects}
\label{research:chapter}

Several arguments in the preceding chapters suggest further questions about uniform estimates, singularity formation, and the geometry behind effective bounds. The projects below begin with concrete models and then indicate broader questions. Their scope varies considerably: some first calculations need only the material of the first semester, whereas the general problems require specialist reading.
The accompanying references describe the known starting points. Before taking up a question, compare the proposed statement with the cited papers and their subsequent developments; an omitted calculation in a textbook need not be new.

\section{Choosing a Project and Further Reading}
\label{research:scoring}

Each score runs from $1$ to $5$.
A higher difficulty score $D$ means more prerequisites and uncertainty; a higher value score $V$ means greater potential consequences for other problems; a higher starting-feasibility score $S$ means a more accessible first result using the book's tools.
The last score concerns the first stage, not the probability of solving the final problem. The scores are subjective guides to project selection and should be adjusted to the reader's preparation.

\begin{center}
\small
\begin{tabular}{@{}c p{0.65\linewidth} c c c@{}}
\toprule
No. & Direction & $D$ & $V$ & $S$\\
\midrule
1 & Quantitative constants in \Cech{}--Dolbeault comparison & 3 & 3 & 5\\
2 & Local weighted heat traces under relative-form perturbations & 4 & 4 & 3\\
3 & Two-parameter heat asymptotics near curvature zeroes & 4 & 5 & 3\\
4 & Scaling limits for Hermitian metrics of low regularity & 4 & 4 & 3\\
5 & Multiplier-ideal jumps in explicit degenerating families & 3 & 4 & 4\\
6 & Geometry and smoothing costs of thinly supported positive currents & 4 & 5 & 3\\
7 & Pluripolar support in nonprojective K\"ahler classes & 5 & 5 & 2\\
8 & Geometric limits of non-K\"ahler mass boundedness & 5 & 5 & 2\\
9 & Controlled degeneration in boundary Morse estimates & 4 & 4 & 2\\
10 & Geometric improvements of effective adjoint-bundle bounds & 5 & 5 & 2\\
\bottomrule
\end{tabular}
\end{center}

\paragraph{Literature as of 10 November 2026.}
The status descriptions below refer to the cited versions available on this date.
For real $(1,1)$-classes on a compact K\"ahler $n$-fold, nef means that smooth representatives can be chosen above $-\varepsilon\omega$ for every $\varepsilon>0$. Use the volume defined in Remark~\ref{pos:class-volume}, and write $\alpha^n=\int_X\alpha^n$ and $\alpha^{n-1}\cdot\beta=\int_X\alpha^{n-1}\wedge\beta$. Tosatti's 2026 preprint gives, for nef $\alpha,\beta$ on arbitrary compact K\"ahler manifolds,
\[
\operatorname{Vol}(\alpha-\beta)
\geq\alpha^n-n\alpha^{n-1}\cdot\beta,
\]
and discusses consequences for the volume function and duality of positive cones \cite{Tosatti2026}.
Related projects concern quantitative consequences, support geometry, or extensions beyond the compact K\"ahler setting.

For a smooth projective $n$-fold and an ample line bundle $L$, Han's 2026 preprint proves global generation of $K_X+mL$ for
$m\geq\lceil C_0n\rceil$, with the explicit constant $C_0=1.77629\ldots$ \cite{Han2026}.
The gap between this bound and the Fujita threshold $n+1$ leaves room for geometric refinements. \hyperref[research:fujita]{Question 10} focuses on such refinements in specified families.

\section{Question 1: Can comparison maps retain quantitative information?}
\label{research:cech}

\paragraph{Question.}
For a family of open covers with uniform geometric control, can the comparison between \Cech{} and Dolbeault cohomology be implemented by chain maps and chain homotopies with explicit norm bounds?
Determine the dependence on overlap multiplicity, derivatives of partition functions, and norms of solution operators on intersections, and how these constants grow when a cover degenerates.

\paragraph{Background.}
The partition-of-unity construction proves acyclicity of fine sheaves, while local $\dbar$ solution operators eliminate vertically closed elements in the comparison double complex.
The explicit constructions here are Proposition~\ref{sh:cech-contraction}, Lemma~\ref{sh:double-complex}, and Theorem~\ref{sh:dolbeault}; see also \cite{Demailly2012,Wells2008}.
The proposed estimate would add geometric control of the comparison maps to the qualitative isomorphism.

\paragraph{A first task.}
Fix finitely many polydisks and a smaller cover. Write each partition multiplication and local solution operator explicitly, recording the constants in $C^r$ or Sobolev norms.
For comparison, compute the same construction on a complex torus: the contracting homotopy in Proposition~\ref{heat:torus-cohomology} introduces only a factor $|m|^{-1}$ at nonzero frequency.
Then narrow selected intersections and test whether the constants must diverge.
An attainable upper bound together with a matching lower bound would provide more information than a formal comparison formula.

\paragraph{Main obstacle.}
Local solvability alone does not give a uniform solution norm, and arbitrary refinement may enlarge partition derivatives.
Fix a geometrically controlled family before seeking cover-independent estimates.
An undergraduate project can complete the explicit-cover calculation; broader uniformity requires further geometric analysis.

\section{Question 2: Can relative-form errors control localized heat traces?}
\label{research:weighted-trace}

\paragraph{Question.}
For positive integers $k$, let $A_k,B_k$ be nonnegative self-adjoint elliptic operators on a compact manifold, with a common form domain $V$. Put $q_A(u)=\|A^{1/2}u\|^2$ for a nonnegative self-adjoint $A$, and suppose
\[
|q_{A_k}(u)-q_{B_k}(u)|
\leq\epsilon_k\bigl(q_{B_k}(u)+\|u\|^2\bigr),
\qquad\epsilon_k\to0.
\]
Under what geometric hypotheses can one control, uniformly as $k\to\infty$ after the spatial and temporal rescalings of Chapter~\ref{heat:local-chapter}, the difference
\[
\tr\bigl(\chi e^{-tA_k}\chi\bigr)
-\tr\bigl(\chi e^{-tB_k}\chi\bigr)
\]
for $t>0$ and a real smooth cutoff $\chi$, acting by multiplication?
Specify the time range and the derivatives of $\chi$ entering the constants.

\paragraph{Background.}
Without $\chi$, Corollary~\ref{heat:form-comparison} gives heat-trace bounds.
For an $L^2$-bounded difference $A_k-B_k$, Proposition~\ref{heat:trace-perturbation} gives a trace-norm estimate.
Spatial localization introduces commutators with $\chi$ in addition to the form perturbation. Proposition~\ref{heat:nonmonotone} explains why operator-exponential monotonicity cannot supply the desired comparison.

\paragraph{Approach.}
Test two-dimensional matrices and one-dimensional periodic operators to decide whether form smallness alone suffices. If it does not, identify the missing local commutator control.
Interpret Duhamel's integral between the form domain and its dual, placing positive-time smoothing on both sides.
Split time at $t/2$, as in \eqref{heat:trace-norm-estimate}, so that the singular factors do not accumulate at one endpoint.
Finally test the hypotheses on magnetic Laplacians and express them directly in geometric coefficients.

\paragraph{A useful outcome.}
Aim for a localized trace-comparison lemma with explicit scales that can be used in heat-kernel localization, or a counterexample ruling out a proposed uniform estimate.
Closeness of quadratic forms by itself is not a proof of pointwise closeness of heat kernels.

\section{Question 3: How do two limits interact near zero curvature?}
\label{research:degenerate-curvature}

\paragraph{Question.}
When curvature eigenvalues approach zero near a hypersurface, can one construct a uniform model for $k\to\infty$ and $t=t(k)\to\infty$ that describes the transition between positive, negative, and zero curvature?
Start with a specified finite-order degeneration, such as
$b(x)=c x_1^\ell+O(|x|^{\ell+1})$, where $x$ is a real local coordinate, $c\ne0$ is real, and $\ell\ge1$ is an integer; determine the appropriate spatial and temporal scales.

\paragraph{Background.}
Proposition~\ref{heat:model-limit} shows that the constant-coefficient amplitude has no singularity at a zero eigenvalue.
It does not supply joint asymptotics for variable coefficients.
Full smooth fixed-positive-time scaling expansions already exist \cite{MaMarinescuZelditch2015}, as do results on spectral kernels with degenerate curvature and varying line bundles \cite{Chiang2024}.
The question here concerns uniformity when the geometric parameter and the heat time vary together.

\paragraph{Approach.}
Construct a finite-order degenerate model potential in one complex dimension. Balance the magnetic and differential terms to identify the scaling.
Check matching away from the degeneracy using the constant-field formula \eqref{heat:magnetic-formula}; near it, a nonquadratic model is needed.
First seek a local heat-trace leading term with an error estimate. Then investigate low-energy spectra and consequences for cohomology remainders.

\paragraph{Main obstacle.}
Constants in a fixed-$t$ expansion may grow rapidly with $t$. Without controlling that growth, replacing $t$ by $t(k)$ is unjustified.
The aim is to identify the actual transition between scales, not to interchange two limits formally.

\section{Question 4: What coefficient control survives at low regularity?}
\label{research:low-regularity}

\paragraph{Question.}
A $C^{1,1}$ weight has locally Lipschitz first derivatives. A continuity modulus is a nondecreasing function $\varpi:[0,1]\to[0,\infty)$ with $\varpi(r)\to0$ as $r\downarrow0$, bounding the oscillation at distance $r$. It is Dini if $\int_0^1\varpi(r)\,dr/r<\infty$; Dini-continuous derivatives have oscillation bounded by such a modulus. Under either regularity assumption, which scaled heat-trace conclusions remain valid?
Begin with an integrated leading term or an error modulus, with precisely defined weak coefficients, before asking for pointwise smooth diagonal expansions.

\paragraph{Background.}
The smooth arguments are developed in \cite{MaMarinescu2007,MaMarinescuZelditch2015}.
A Gaussian second derivative contributes $s^{-1}$, while a coefficient modulus at distance $\sqrt s$ contributes $\varpi(\sqrt s)$.
The endpoint integral is
\[
\int_0^t\frac{\varpi(\sqrt s)}s\,ds
=2\int_0^{\sqrt t}\frac{\varpi(r)}r\,dr.
\]
This substitution explains the relevance of Dini continuity, but does not itself prove a low-regularity heat-kernel theorem.

\paragraph{Approach.}
Define the operator by a closed quadratic form and specify weak curvature, volume measure, and test functions.
Smooth the metric to $h_\epsilon$ and estimate the errors as functions of both $\epsilon$ and $k$ before choosing $\epsilon=\epsilon(k)$.
Show that changing the smoothing method leaves the limit unchanged.
A manageable first goal is a precise integrable error bound on a torus or disk under a Dini condition, together with a failure example outside it.

\paragraph{Scope.}
Low-regularity scalar parabolic theory is extensive.
A new result should address the interaction of high tensor powers, bundle connections, changing curvature signs, and uniform constants, rather than merely restating a scalar estimate in bundle notation.

\section{Question 5: Multiplier ideals and thresholds in degenerating families}
\label{research:multiplier}

\paragraph{Question.}
Let $\tau$ range over a small complex parameter disk, let $f_{j,\tau}(z)$ be holomorphic in space and parameter, and let $c>0$. For the analytic-singularity weights
\[
\varphi_\tau=\log\left(|f_{1,\tau}|^2+\cdots+|f_{N,\tau}|^2\right),
\]
describe exactly the jumps of $\mathcal I(c\varphi_\tau)$ and the growth of local integrability constants as $\tau$ approaches a degenerate parameter.
Start with an explicit fixed-degree polynomial family in two complex variables.

\paragraph{Background.}
Guan--Zhou's strong openness theorem supplies a slightly larger integrable exponent \cite{GuanZhou2015}. The problem below concerns how such exponents and their constants vary with the parameter.
A jumping number is an exponent at which $\mathcal I(c\varphi)$ becomes strictly smaller as $c$ increases. For monomial ideals or a given log resolution these numbers have established computational methods; the coordinate calculation is Theorem~\ref{pos:snc-ideal}.
The target should be exact parameter dependence in a particular family, or an effective uniform estimate tied to geometry.

\paragraph{Approach.}
In the normal-crossing coordinate case, use polar integration to turn integrability into strict inequalities between exponents.
For the chosen two-variable family, perform explicit blowups, recording Jacobian factors and vanishing orders along exceptional divisors.
Infer candidate critical exponents, then verify every candidate by integration.
Include an example whose multiplier ideal cannot be recovered from the Lelong number alone, separating the information in a single numerical invariant from that in an ideal.

\paragraph{A first outcome.}
Describe the parameter strata on which the multiplier ideal is constant, prove the description, and obtain two-sided estimates as a critical parameter is approached.
Explicit families are accessible to undergraduates; uniform statements require additional singularity theory and $L^2$ methods.

\section{Question 6: Support geometry and the cost of smoothing positive currents}
\label{research:thin-support}

\paragraph{Question.}
After constructing a positive representative of a K\"ahler class supported on a closed measure-zero set, can one control the local support geometry, the distribution of mass at interfaces, and the tradeoff between positivity loss and derivative growth under smoothing?

\paragraph{Background.}
Tosatti's 2026 preprint constructs such measure-zero-supported currents using local potentials and finite maxima \cite{Tosatti2026}.
Theorem~\ref{pos:thin-current} gives the construction in our conventions, using pluriharmonic differences between local potentials.
The proposed project asks for further control of the support and its regularization. Pluripolar support is a stronger requirement and leads to the separate question below.

\paragraph{Approach.}
Begin with the maximum of two smooth local potentials.
At a transverse interface, compute the distributional derivative and the interface-supported contribution explicitly and check positivity.
Treat triple interfaces by strata.
Then choose a smooth convex regularized maximum and compute how its curvature lower bound, support-neighborhood thickness, and higher derivatives depend on the smoothing scale.
Connecting these quantities with errors in Monge--Amp\`ere mass concentration could produce a useful quantitative tool.

\paragraph{Main obstacle.}
The local potentials need not be real analytic; their pluriharmonic differences on overlaps are real analytic.
This does not give arbitrary level sets of individual potentials the same analytic structure.
Moreover, smoothing distributes curvature over a thin neighborhood, so the smoothed form cannot be treated pointwise as the original support current.

\section{Question 7: Pluripolar support and nonprojective K\"ahler geometry}
\label{research:plurisupported}

\paragraph{Question.}
On a given nonprojective compact K\"ahler manifold, does some K\"ahler class contain a closed positive $(1,1)$-current with pluripolar support?
If not for every manifold, can one identify verifiable sufficient conditions or obstructions?
A set is pluripolar if locally it lies in the $-\infty$ locus of a psh function not identically $-\infty$; keep this distinct from measure-zero support, as in Remark~\ref{pos:class-volume}.

\paragraph{Background.}
McCleerey's systematic study appeared in the \emph{Indiana University Mathematics Journal} in 2025 \cite{McCleerey2025}.
Tosatti's 2026 preprint also discusses the difficulty of general pluripolar-support constructions and uses measure-zero support for a role in volume comparison \cite{Tosatti2026}.
Such currents are of interest through their singularities, their cohomology classes, and the restrictions their support imposes on the ambient geometry.

\paragraph{Approach.}
Start with nonprojective complex tori or examples with explicitly understood $(1,1)$ cones.
Compare divisor currents, dynamical constructions, and envelope methods.
Determine the cohomology class, singularity type, and full support in each example; local Lelong numbers give only part of this information.
Seek first an existence result for one geometric family or an obstruction tied to a numerical property of the class.

\paragraph{Main obstacle.}
A current may have zero Lelong number everywhere while retaining complicated singularities; measure zero does not imply pluripolarity.
The full problem is too broad for a first undergraduate project, but understanding one explicit construction and its limitations is a suitable longer-term entry to positive-current theory.

\section{Question 8: Geometric limits of non-K\"ahler mass boundedness}
\label{research:nonkahler}

\paragraph{Question.}
For a compact Hermitian manifold $(X,\omega)$, define
\[
M(\omega)=\sup\left\{
\int_X(\omega+\ii\partial\dbar\varphi)^n:
\varphi\in C^\infty(X,\R),\quad
\omega+\ii\partial\dbar\varphi>0\right\}.
\]
Which verifiable geometric conditions imply $M(\omega)<\infty$?
In a concrete fibration or deformation family, can torsion, vertical energy, or another structural quantity distinguish boundedness from divergence?
When boundedness fails, can mass growth be related quantitatively to the regularity cost of the potentials?

\paragraph{Background.}
In the K\"ahler case, Stokes's theorem makes the integral equal to $\int_X\omega^n$.
For Hermitian manifolds, Boucksom--Guedj--Lu develop a volume theory for nonclosed Bott--Chern classes under a bounded-mass hypothesis \cite{BoucksomGuedjLu2025}.
This hypothesis is not universal: the 2026 preprint of Xia--Zhang gives a counterexample on the Hopf threefold
$(\C^3\setminus\{0\})/\langle z\mapsto e^{-1}z\rangle$ \cite{XiaZhang2026}.
This counterexample makes the choice of geometric hypotheses central: the question asks which conditions ensure boundedness and how their failure permits mass growth.

A Hermitian background is pluriclosed when $\partial\dbar\omega=0$. The bigness results of Li--Wang--Zhou under specific assumptions on pluriclosed backgrounds, nef classes, and bounded quasi-psh potentials also emphasize the need to track each structural hypothesis \cite{LiWangZhou2025}.
Bigness, mass boundedness, and sharp volume lower bounds are separate objectives.
To obtain a non-K\"ahler result, the proposed hypotheses must allow manifolds that carry no K\"ahler structure.

\paragraph{Approach.}
Expand $(\omega+\ii\partial\dbar\varphi)^n$ on an explicit low-dimensional model, retaining every $\partial\omega$ and $\partial\dbar\omega$ term produced by integration by parts.
Relate the mass error to a concrete energy and determine whether positivity permits that energy to diverge.
For models fibered by elliptic curves, the Fourier kernels in this book can help construct and test smooth potentials, but mixed base-fiber derivatives require separate control.
First restrict to a specified family of potentials and prove complete two-sided bounds; only then enlarge the admissible family.

\paragraph{Main obstacle.}
A locally positive perturbation may fail to remain positive globally.
Likewise, a formal sequence of large-energy functions need not make total mass diverge.
A useful result would identify precise mass-control conditions or explain how a known counterexample crosses their boundary, thereby supplying valid analytic hypotheses for volume comparison.

\section{Question 9: Morse estimates with degenerating boundary geometry}
\label{research:boundary}

\paragraph{Question.}
For a family of complex manifolds with boundary and controlled Levi-form degeneration, can one obtain uniform spectral estimates including the boundary contribution, with explicit dependence on the degeneracy?
First fix the cohomology and the boundary condition, such as the $\dbar$-Neumann realization associated with the quadratic form $\|\dbar u\|^2+\|\dbar^*u\|^2$ on the intersection of the maximal $\dbar$ domain and its adjoint domain. Different choices give different spectra and cohomology groups.

\paragraph{Background.}
Holomorphic Morse inequalities with boundary were studied by Berman \cite{Berman2005}.
The 2008 erratum must be read alongside the paper because certain sharpness assertions require added hypotheses \cite{Berman2008}.
Adding a boundary integral to a closed-manifold formula without establishing the boundary theory is insufficient.

\paragraph{Approach.}
Choose a separable model domain and compare how boundary conditions alter zero modes and low-energy spectra.
Construct the interior Gaussian and half-space boundary models separately, then check their matching.
A manageable first target is a lower bound on the inevitable growth of error constants as quantitative strict Levi positivity tends to zero.
A later target is a stable replacement scale when the maximal normalized order of contact of holomorphic curve germs with the boundary is finite (a finite-type condition).

\paragraph{Main obstacle.}
Interior elliptic estimates do not establish a boundary problem.
Without an appropriate subelliptic or closed-range estimate, finite dimensionality and the required heat-trace properties may not even be available.
This direction suits readers with additional PDE and several-complex-variable boundary theory.

\section{Question 10: Can effective adjoint bounds reflect more geometry?}
\label{research:fujita}

\paragraph{Question.}
The general bounds discussed above depend only on the dimension $n$.
Can additional local positivity, projective-bundle structure, controlled singularities, or special curve geometry improve the global-generation and embedding bounds for $K_X+mL$ toward their critical values?
Separate global generation from very ampleness, which also requires point and tangent separation.

\paragraph{Background.}
Example~\ref{heat:fujita-counterexample} rules out the false strengthening that $K_X+(n+1)L$ must always be ample.
Compare general linear global-generation bounds with Han's 2026 preprint \cite{Han2026}; for the traditional singular-metric, multiplier-ideal, and jet-separation route, see \cite{Demailly2012}.
A useful refinement should express the bound in terms of a geometric parameter that can be estimated in the chosen family.

\paragraph{Approach.}
Compute the smallest possible integers on curves, projective spaces, and explicit projective bundles, and identify why the critical examples fail.
Convert jet separation into vanishing with a finite-length quotient sheaf, then compute local weights forcing the multiplier ideal into the desired ideal powers.
For an ample $L$, its Seshadri constant at $x$ is $\varepsilon(L;x)=\inf_{C\ni x}(L\cdot C)/\operatorname{mult}_xC$, over irreducible curves through $x$, with $\operatorname{mult}_xC$ the local multiplicity. A bound involving this constant becomes effective when one also estimates it from the geometry of the chosen family.

\paragraph{A first outcome.}
A bound sharp on an explicit family, or a simple lemma explaining several older bounds at once, may be more informative than a tiny numerical improvement of a universal constant.
Optimal conjectures in all dimensions are long-term goals. An undergraduate project can instead complete all boundary examples and jet calculations in one low-dimensional family.

\section{Choosing a First Project}
\label{research:workflow}

A manageable project begins with a precise statement and a small collection of examples. A model where the proposed estimate can be calculated, a borderline case, and a counterexample to a missing hypothesis often reveal which part of the question is essential. Compare these examples with the cited literature before deciding what the project adds.

The first goal may be a single lemma: an integrable remainder in a localized heat trace, a critical integral in a degenerating family, or an interface formula for a positive current. Its hypotheses should specify the operator domains, limiting parameters, and the data on which constants depend. A complete proof of such a lemma gives a basis for deciding which broader generalizations are plausible.

\hyperref[research:cech]{Question 1}, \hyperref[research:multiplier]{Question 5}, and the explicit-model part of \hyperref[research:degenerate-curvature]{Question 3} offer the most accessible starting points in this book.
\hyperref[research:weighted-trace]{Question 2}, \hyperref[research:low-regularity]{Question 4}, and \hyperref[research:thin-support]{Question 6} suit readers with a stronger analysis background.
The full goals of Questions 7--10 require sustained specialist reading and discussion with a supervisor. A useful first project can consist of one precise proposition, its proof, and a clear account of its relation to existing results.


\appendix
\renewcommand{\chaptermark}[1]{\markboth{Appendix \thechapter}{}}
\titleformat{\chapter}[display]
 {\sffamily\bfseries\color{Navy}\raggedright}
 {\large\color{Teal}Appendix \thechapter}
 {8pt}{\Huge}[\vspace{10pt}{\color{Gold}\titlerule[1.2pt]}]


\chapter{Background Theorems}\label{dep:chapter}
Some of the later applications use results from topology, analytic spaces, and algebraic geometry beyond the core course. This appendix states the forms of those results that are used in the text and points to their proofs in the literature. A reader can consult the relevant paragraph when a chapter first calls on one of these theorems.

\section{Projective Geometry and Topology}

\paragraph{Good covers and topological cohomology comparisons.}\label{dep:good-cover}
A good cover is an open cover whose every nonempty finite intersection is contractible. We use the standard topological facts that smooth manifolds admit good covers, that the constant integer sheaf has no positive-degree cohomology on these intersections, and that its resulting \Cech{} and sheaf cohomology agree with singular cohomology. Over $\R$ or $\C$, de Rham cohomology identifies with the same cohomology through integration; the coefficient map from $\Z$ describes integral periods. Smooth manifolds also have the homotopy type of finite-dimensional CW complexes. We use the long exact sequence of a pair and the relative cup product
\[
 H^a(M,U;\Z)\otimes H^b(M,V;\Z)
 \longrightarrow H^{a+b}(M,U\cup V;\Z)
\]
for open subsets $U,V\subset M$, together with its compatibility with the absolute cup product. These are topological inputs from \cite{BottTu1982}, used in Lemma~\ref{prj:integral-gluing} and Proposition~\ref{idx:integral-chern}. The analytic gluing, metric, and projective-bundle calculations are supplied in those proofs. Fineness of the smooth-function sheaf does not make the constant integer sheaf fine, so the fine-resolution argument alone cannot replace this comparison input.

\paragraph{GAGA.}\label{dep:gaga}
For a complex projective algebraic variety, analytification gives an equivalence between the categories of algebraic and analytic coherent sheaves and preserves all sheaf cohomology groups \cite{Serre1956}. After obtaining a projective embedding, Theorem~\ref{amp:serre} uses this comparison to pass to finite graded modules, then proves vanishing and generation by clearing denominators and using Hilbert resolutions. The projectivity hypothesis is essential; generalizations require appropriate properness hypotheses.

\paragraph{The index formula.}\label{dep:index}
For a $\operatorname{Spin}^c$ Dirac operator on a compact smooth manifold of even dimension, with determinant line bundle $B$ and twisting complex vector bundle $E$,
\[
 \operatorname{ind}D_E^+=\int_X\operatorname{ch}(E)e^{c_1(B)/2}\widehat A(TX).
\]
This is a specific form of the Atiyah--Singer index theorem \cite{AtiyahSinger1968}. Chapter \ref{idx:chapter} derives holomorphic Riemann--Roch, Chern--Gauss--Bonnet, and the complete Euler characteristic polynomial from this input. Riemann--Roch for curves in Chapter \ref{prj:chapter}, the algebraic proof of Theorem \ref{idx:hilbert-polynomial}, and the leading-order Riemann--Roch formula in the heat-kernel chapter do not depend on it.

\paragraph{Bott periodicity and integrality.}\label{dep:bott}
Write $U=\bigcup_{r\ge1}U(r)$ for the stable unitary group. In the stable range, $\pi_{2j}(U)=0$ and $\pi_{2j+1}(U)=\Z$, while $c_s$ takes the value $\pm(s-1)!$ on a generator of $\pi_{2s}(BU(r))$. Starting from this input, the text gives each step of obstruction elimination in Peterson's criterion for triviality; see Theorem \ref{idx:peterson}. Periodicity and factorial integrality themselves belong to topological $K$-theory; see \cite{Bott1959,Peterson1959}. Vanishing of real Chern classes cannot replace the integral-class hypothesis.

\paragraph{Numerical criteria for ampleness.}\label{dep:ample-criteria}
The Nakai--Moishezon criterion states that a line bundle on a projective variety is ample if and only if its top self-intersection on every positive-dimensional irreducible subvariety is strictly positive. Kleiman's criterion characterizes the projective nef cone as the cone of classes pairing nonnegatively with all effective curves. Both are inputs from algebraic geometry; see \cite{Kleiman1966,Lazarsfeld2004}. The latter is used in Proposition \ref{bir:numerical-nef} to pass from the curve criterion to approximately semipositive metrics.
Propositions~\ref{amp:line-criteria} and \ref{amp:point-tests} give analytic criteria for ampleness. Proposition~\ref{bir:finite-ample} treats the finite-map criterion, and Example~\ref{bir:curve-free-torus} explains why a curve criterion without projectivity is insufficient.

\paragraph{Analytic contraction of the zero section.}\label{dep:contraction}
Grauert's negativity and contraction criterion relates strongly pseudoconvex neighborhoods of a compact zero section to the existence of an analytic contraction \cite{Grauert1962}. The direction asserting that a positive line bundle produces such neighborhoods can be checked directly. On its dual, take the squared norm $\rho$: its vertical Levi form is positive, and at nonzero vectors its horizontal Levi form is positive by the positive curvature of the original line bundle. Thus $\{\rho<\epsilon\}$ has strictly pseudoconvex boundary. The converse gives ampleness through the Fourier-weight argument of Proposition \ref{bir:disk-serre} and Grauert's finiteness theorem.
For a vector bundle, a positive metric on $\mathcal O_{\mathbb P(E^*)}(1)$ gives a homogeneous strictly psh squared norm on $E^*\setminus X$, whose small sublevel sets again form strongly pseudoconvex neighborhoods. Proposition \ref{bir:blowdown-comparison} proves the analytic descent from the tautological line bundle to the vector bundle in full. The existence theorem constructing an actual contraction space from such a neighborhood remains an input from Grauert's theory; a set-theoretic quotient does not acquire an analytic structure automatically.

\section{Analytic Spaces, Modifications, and Coherent Direct Images}
\paragraph{Hironaka resolution and flattening.}\label{dep:resolution}
In the complex analytic category in characteristic zero, locally finite modifications with smooth centers resolve the singularities of analytic spaces or ideals; only finitely many modifications are needed for the compact spaces considered here. For coherent sheaves, we also use the flattening form asserting that, after a proper modification, the pullback modulo torsion is locally free. For algebraic resolution, analytic resolution, and flattening, respectively, see \cite{Hironaka1964,ArocaHironakaVicente2018,Hironaka1975}.
The text proves the more elementary ingredients, including Weierstrass preparation, Noetherianity, finite projections, local free resolutions, and dimensions of analytic supports. These do not automatically supply global resolution. Global resolution is used in the Morse inequalities for coherent sheaves, the transformation of multiplier ideals of general effective divisors, and the relation between the Moishezon property and bigness.

\paragraph{Grauert's proper direct image theorem.}\label{dep:proper-image}
If $f:Y\to X$ is proper and holomorphic and $\mathcal F$ is coherent, every $R^qf_*\mathcal F$ is coherent. We use \cite{Grauert1960,Grauert1963}, together with the errata. In the coherent-sheaf reduction of Chapter \ref{morse:chapter}, this ensures that induction on support dimension applies to the error sheaves. The filtration accounting for nonreduced structure, the page-by-page dimension comparison in the Leray spectral sequence, and the $O(k^{d-1})$ error are all proved in the text. The special direct image formulas for local finite maps and projective bundles are proved separately in the book and do not require repeated appeals to the general theorem.

\paragraph{Remmert and Chow.}\label{dep:remmert-chow}
The image of a proper holomorphic map is analytic, and a closed analytic subset of projective space is algebraic. These two inputs are used to interpret graphs and images of meromorphic maps in terms of algebraic dimension, and to obtain the Stein reduction of strongly pseudoconvex spaces. See the chapters on analytic spaces and proper maps in \cite{Demailly2012}. The text also gives direct proofs for finite maps, zero-dimensional fibers, and descent in local coordinates; the general image theorem for singular analytic spaces remains a foundational input.

\paragraph{Regularization with analytic singularities and Moishezon modifications.}\label{dep:moishezon}
A strictly positive closed $(1,1)$ current admits a regularization with analytic singularities that preserves strict positivity, and a compact Moishezon manifold admits a smooth projective modification. For the first input see \cite{Demailly1992}; for the second and its relation to positive currents in integral classes see \cite{JiShiffman1993,Demailly2012}. Theorem \ref{bir:moishezon} gives all subsequent steps, including exceptional-divisor corrections after resolution, passage to a projective manifold, descent of sections, and pushforward of positive currents.
Regularization with analytic singularities serves a different purpose from the smooth heat regularization in Lemma \ref{nad:heat-regularization}. The latter allows controlled curvature errors that need not converge uniformly to zero, which suffices for the general $L^2$ theorem and Nadel vanishing. It does not automatically produce analytic singularities to which resolution of ideals can be applied.

\paragraph{Relative Grauert--Riemenschneider vanishing.}\label{dep:gr}
For the proper projective birational modifications $\mu:Y\to X$ between smooth complex spaces used in this book, we invoke the relative form
$\mu_*K_Y=K_X$ and $R^q\mu_*K_Y=0$ for $q>0$; see \cite{GrauertRiemenschneider1970}. Its only role here is to descend canonical-bundle vanishing from a projective or K\"ahler modification to a general compact complex manifold. For pulled-back line bundles, the projection formula and the Leray comparison then give the conclusions term by term. Without this relative vanishing, higher cohomology on a modification cannot simply be identified with higher cohomology on the original space.

\section{Noncompact Analysis}
\paragraph{Approximation of closed forms in the noncompact Gigante--Girbau theorem.}\label{dep:runge}
The book first chooses complete metrics and solves the equation on relatively compact sublevel sets. It then invokes the Takegoshi--Ohsawa Runge approximation theorem for closed forms to correct differences between solutions on successive subdomains, obtaining a distributional primitive through convergence in $L^2_{\mathrm{loc}}$. The $L^2$ Dolbeault resolution then yields vanishing. The exact degrees and points of application of this approximation input are specified in the text; see Sections 3 and 4 of \cite{TakegoshiOhsawa1981}.
A positive lower bound for the positive curvature eigenvalues in the compact case cannot be transferred directly to a noncompact space. Vanishing of cohomology on every subdomain does not remove the need for the global gluing argument.

\paragraph{Strong $q$-convexity, strong $q$-completeness, and finiteness.}\label{dep:qconvex}
Strong $q$-completeness means the existence of a smooth exhaustion whose Levi form has at least $n-q+1$ positive eigenvalues everywhere; strong $q$-convexity requires this condition only outside a compact set. Andreotti--Grauert theory gives finite-dimensional coherent-sheaf cohomology $H^j$ for $j\ge q$ on strongly $q$-convex spaces and vanishing of these groups on strongly $q$-complete spaces \cite{AndreottiGrauert1962}. Greene--Wu's result guarantees that every connected noncompact complex manifold of dimension $n$ is strongly $n$-complete \cite{GreeneWu1975}. These theorems concern the geometry at infinity. Their proofs use methods beyond compact Hodge theory.
For manifolds admitting an everywhere strictly psh exhaustion, the text separately supplies $L^2$ proofs of discrete interpolation, holomorphic convexity, and the Levi problem. A strongly pseudoconvex domain and a manifold with an everywhere strictly psh exhaustion must not be identified by definition.

\paragraph{General noncompact holomorphic Morse inequalities.}\label{dep:noncompact-morse}
Bouche's noncompact Morse inequalities \cite{Bouche1989} apply to a line bundle tensored with a holomorphic vector bundle. Let $X$ have complex dimension $n$ and be strongly $q$-convex, assume the curvature of a line bundle has at least $n-p+1$ nonnegative eigenvalues outside a compact set, and let $F$ be a holomorphic vector bundle of rank $r$. For $m\ge p+q-1$,
\[
 \sum_{j=m}^n(-1)^{j-m}h^j(X,L^k\otimes F)
 \le \frac{rk^n}{n!}\int_{X(L,h,\ge m)}(-1)^m c_1(L,h)^n+o(k^n).
\]
Here $h^j$ is the finite cohomology dimension in the stated range, and $X(L,h,\ge m)$ is the union of the nondegenerate curvature loci with at least $m$ negative eigenvalues. The exact convexity assumptions, metric choices, and integrability requirements are those of condition $P(q,p)$ in that paper. In the stated degrees, the sets of the relevant negative-curvature indices are contained in a compact region. These hypotheses provide the compactness and integrability needed in the noncompact estimate.
Chapter~\ref{morse:chapter} proves the compact vector-bundle result and its extension to coherent sheaves using resolution and proper direct images. In the noncompact setting, weighted Hodge comparison and the behavior of solutions on subdomains require the additional analysis in Bouche's theorem.
For a coherent sheaf on a noncompact space, there need not be one integer $N$ such that the $N$th power of the radical of its annihilator kills the entire sheaf. Thus the compact reduction by a finite filtration cannot be carried over without further hypotheses. The theorem stated here has vector-bundle coefficients.

\section{Effective Results and Birational Geometry}
\paragraph{Bogomolov instability.}\label{dep:bogomolov}
A slope-semistable rank-two vector bundle on a smooth projective surface satisfies $c_1^2\le4c_2$; contraposition gives a saturated destabilizing subsheaf when the discriminant is positive. Starting from this input, Chapter \ref{eff:chapter} fully constructs the locally free extensions associated with obstructions on subschemes of lengths one and two and derives Reider's numerical cases and the Fujita bounds for surfaces. This input is not a direct consequence of elementary Riemann--Roch; see the Bogomolov theorem explicitly used in \cite{Reider1988,Lazarsfeld2004}.

\paragraph{The basepoint-free theorem.}\label{dep:bpf}
On a smooth projective manifold, if an integral line bundle $N$ is nef and $aN-K_X$ is nef and big for some positive integer $a$, then $mN$ is globally generated for a sufficiently divisible positive integer $m$. This is the basepoint-free theorem of minimal model theory; see the treatment of adjoint linear systems in \cite{Lazarsfeld2004}.
The text first proves independently that $N=K_X+(n+1)L$ is nef and then observes that $N-K_X=(n+1)L$ is ample. Thus this input gives semiampleness of $N$. It does not give basepoint-freeness of $N$ itself, which is the sharper Fujita problem.

\paragraph{An equivalent formulation of abundance and the uniruledness criterion.}\label{dep:birational}
Eckl's theorem states that if $S$ is a minimal model of a smooth projective variety $X$ that is not uniruled, then
\[
 \nu_X(K_X)=\kappa_X(K_X)\quad\Longleftrightarrow\quad K_S\text{ is semiample}.
\]
Here $S$ may have the terminal singularities prescribed by minimal model theory, and $K_S$ is $\mathbb Q$-Cartier. The notation $\nu_X$ has the asymptotic-section meaning specified in \eqref{bir:numerical-convention}; it cannot be replaced by ordinary self-intersection for a general pseudoeffective divisor. This theorem asserts neither the existence of minimal models in every dimension nor that the equality always holds; see \cite{Eckl2016}.
Ou's result states that a compact K\"ahler manifold is covered by rational curves if and only if $K_X$ is not pseudoeffective \cite{Ou2025}. Proposition~\ref{bir:uniruled} proves that being covered by rational curves implies that the canonical class is not pseudoeffective, and that non-pseudoeffectivity implies $\kappa=-\infty$. Ou's reverse implication involves separate theories of foliations and relative metrics, for which the original paper remains the reference. The two implications proved here do not imply that $\kappa=-\infty$ forces coverage by rational curves.


\backmatter
\begin{thebibliography}{99}
\raggedright
\addcontentsline{toc}{chapter}{Bibliography}
\markboth{Bibliography}{}

\bibitem{Demailly2012}
J.-P. Demailly, \emph{Complex Analytic and Differential Geometry},
Lecture notes, version of 21 June 2012.
\url{https://www-fourier.univ-grenoble-alpes.fr/~demailly/manuscripts/agbook.pdf}.

\bibitem{DemaillyAnalytic2012}
J.-P. Demailly, \emph{Analytic Methods in Algebraic Geometry},
Surveys of Modern Mathematics \textbf{1}, International Press / Higher Education Press, 2012.
Author's manuscript: \url{https://www-fourier.univ-grenoble-alpes.fr/~demailly/manuscripts/analmeth_book.pdf}.

\bibitem{Huybrechts2005}
D. Huybrechts, \emph{Complex Geometry: An Introduction}, Universitext, Springer, 2005.
\url{https://doi.org/10.1007/b137952}.

\bibitem{BottTu1982}
R. Bott and L. W. Tu, \emph{Differential Forms in Algebraic Topology},
Graduate Texts in Mathematics \textbf{82}, Springer, 1982.

\bibitem{Fulton1998}
W. Fulton, \emph{Intersection Theory}, 2nd ed.,
Ergebnisse der Mathematik und ihrer Grenzgebiete, 3. Folge, \textbf{2}, Springer, 1998.

\bibitem{Evans2010}
L. C. Evans, \emph{Partial Differential Equations}, 2nd ed.,
Graduate Studies in Mathematics \textbf{19}, American Mathematical Society, 2010.

\bibitem{Stacks}
The Stacks Project Authors, \emph{The Stacks Project},
Chapters ``Sheaves on Spaces'' and ``Cohomology of Sheaves,'' online version, accessed November 2026.
\url{https://stacks.math.columbia.edu/}.

\bibitem{Hormander1965}
L. H\"ormander, \emph{$L^2$ estimates and existence theorems for the $\bar\partial$ operator},
Acta Mathematica \textbf{113} (1965), 89--152.
\url{https://doi.org/10.1007/BF02391775}.

\bibitem{Skoda1972}
H. Skoda, \emph{Application des techniques $L^2$ \`a la th\'eorie des id\'eaux d'une alg\`ebre de fonctions holomorphes avec poids},
Annales scientifiques de l'\'Ecole Normale Sup\'erieure (4) \textbf{5} (1972), no.~4, 545--579.
\url{https://doi.org/10.24033/asens.1237}.

\bibitem{HunekeSwanson2006}
C. Huneke and I. Swanson, \emph{Integral Closure of Ideals, Rings, and Modules},
London Mathematical Society Lecture Note Series \textbf{336}, Cambridge University Press, 2006.
Author's publicly available version and errata: \url{https://www.math.purdue.edu/~iswanso/book/index.html}.

\bibitem{Nadel1990}
A. M. Nadel, \emph{Multiplier ideal sheaves and K\"ahler--Einstein metrics of positive scalar curvature},
Annals of Mathematics \textbf{132} (1990), no.~3, 549--596.
\url{https://annals.math.princeton.edu/1990/132-3/p04}.

\bibitem{Lazarsfeld2004}
R. Lazarsfeld, \emph{Positivity in Algebraic Geometry I--II},
Ergebnisse der Mathematik und ihrer Grenzgebiete, 3. Folge, \textbf{48--49}, Springer, 2004.


\bibitem{Demailly1985}
J.-P. Demailly,
\emph{Champs magn\'etiques et in\'egalit\'es de Morse pour la $d''$-cohomologie},
Annales de l'Institut Fourier \textbf{35} (1985), no.~4, 189--229.
DOI: \href{https://doi.org/10.5802/aif.1034}{10.5802/aif.1034}.
\url{https://www.numdam.org/articles/10.5802/aif.1034/}.

\bibitem{DemaillyMorse1991}
J.-P. Demailly,
\emph{Holomorphic Morse inequalities},
in \emph{Several Complex Variables and Complex Geometry, Part 2}
(Santa Cruz, CA, 1989),
Proceedings of Symposia in Pure Mathematics \textbf{52}, Part~2,
American Mathematical Society, 1991, 93--114.
\url{https://www-fourier.univ-grenoble-alpes.fr/~demailly/manuscripts/morse.pdf}.

\bibitem{Bismut1987}
J.-M. Bismut,
\emph{Demailly's asymptotic Morse inequalities: A heat equation proof},
Journal of Functional Analysis \textbf{72} (1987), no.~2, 263--278.
DOI: \href{https://doi.org/10.1016/0022-1236(87)90089-9}{10.1016/0022-1236(87)90089-9}.
\url{https://www.imo.universite-paris-saclay.fr/~jean-michel.bismut/Bismut/1987.pdf}.

\bibitem{MaMarinescu2007}
X. Ma and G. Marinescu,
\emph{Holomorphic Morse Inequalities and Bergman Kernels},
Progress in Mathematics \textbf{254}, Birkh\"auser, 2007.
DOI: \href{https://doi.org/10.1007/978-3-7643-8115-8}{10.1007/978-3-7643-8115-8}.
\url{https://www.mi.uni-koeln.de/~gmarines/mmbook.html}.

\bibitem{MaMarinescuZelditch2015}
X. Ma, G. Marinescu and S. Zelditch,
\emph{Scaling asymptotics of heat kernels of line bundles},
Contemporary Mathematics \textbf{644} (2015), 175--202.
DOI: \href{https://doi.org/10.1090/conm/644/12768}{10.1090/conm/644/12768}.
\url{https://arxiv.org/abs/1406.0201}.

\bibitem{Hironaka1975}
H. Hironaka,
\emph{Flattening theorem in complex-analytic geometry},
American Journal of Mathematics \textbf{97} (1975), no.~2, 503--547.
DOI: \href{https://doi.org/10.2307/2373721}{10.2307/2373721}.

\bibitem{ArocaHironakaVicente2018}
J. M. Aroca, H. Hironaka and J. L. Vicente,
\emph{Complex Analytic Desingularization}, Springer, Tokyo, 2018.
DOI: \href{https://doi.org/10.1007/978-4-431-49822-3}{10.1007/978-4-431-49822-3}.
\url{https://link.springer.com/book/10.1007/978-4-431-49822-3}.

\bibitem{Grauert1960}
H. Grauert,
\emph{Ein Theorem der analytischen Garbentheorie und die Modulr\"aume komplexer Strukturen},
Publications Math\'ematiques de l'IH\'ES \textbf{5} (1960), 5--64.
DOI: \href{https://doi.org/10.1007/BF02684746}{10.1007/BF02684746}.
\url{https://www.numdam.org/articles/10.1007/BF02684746/}.
Read together with the following correction.

\bibitem{Grauert1963}
H. Grauert,
\emph{Berichtigung: Ein Theorem der analytischen Garbentheorie und die Modulr\"aume komplexer Strukturen},
Publications Math\'ematiques de l'IH\'ES \textbf{16} (1963), 35--36.
DOI: \href{https://doi.org/10.1007/BF02684291}{10.1007/BF02684291}.
\url{https://www.numdam.org/articles/10.1007/BF02684291/}.

\bibitem{Bouche1989}
T. Bouche,
\emph{In\'egalit\'es de Morse pour la $d''$-cohomologie sur une vari\'et\'e holomorphe non compacte},
Annales scientifiques de l'\'Ecole Normale Sup\'erieure (4)
\textbf{22} (1989), no.~4, 501--513.
DOI: \href{https://doi.org/10.24033/asens.1590}{10.24033/asens.1590}.
\url{https://www.numdam.org/articles/10.24033/asens.1590/}.

\bibitem{Berman2005}
R. Berman,
\emph{Holomorphic Morse inequalities on manifolds with boundary},
Annales de l'Institut Fourier \textbf{55} (2005), no.~4, 1055--1103.
DOI: \href{https://doi.org/10.5802/aif.2121}{10.5802/aif.2121}.
\url{https://www.numdam.org/articles/10.5802/aif.2121/}.
Read together with the following 2008 corrigendum.

\bibitem{Berman2008}
R. Berman,
\emph{Corrigendum to: Holomorphic Morse inequalities on manifolds with boundary},
Annales de l'Institut Fourier \textbf{58} (2008), no.~1, 377--381.
DOI: \href{https://doi.org/10.5802/aif.2354}{10.5802/aif.2354}.
\url{https://www.numdam.org/articles/10.5802/aif.2354/}.

\bibitem{Wells2008}
R. O. Wells, Jr.,
\emph{Differential Analysis on Complex Manifolds}, 3rd ed.,
Graduate Texts in Mathematics \textbf{65}, Springer, 2008.
DOI: \href{https://doi.org/10.1007/978-0-387-73892-5}{10.1007/978-0-387-73892-5}.

\bibitem{GuanZhou2015}
Q. Guan and X. Zhou,
\emph{A proof of Demailly's strong openness conjecture},
Annals of Mathematics \textbf{182} (2015), no.~2, 605--616.
DOI: \href{https://doi.org/10.4007/annals.2015.182.2.5}{10.4007/annals.2015.182.2.5}.
\url{https://annals.math.princeton.edu/2015/182-2/p05}.

\bibitem{Chiang2024}
Y.-L. Chiang,
\emph{Spectral Kernels and Holomorphic Morse Inequalities for Sequence of Line Bundles},
Preprint, arXiv:2404.18079v1, 28 April 2024.
\url{https://arxiv.org/abs/2404.18079}.

\bibitem{McCleerey2025}
N. McCleerey,
\emph{Plurisupported currents on compact K\"ahler manifolds},
Indiana University Mathematics Journal \textbf{74} (2025), no.~3, 605--645.
\url{https://www.iumj.indiana.edu/IUMJ/FULLTEXT/2025/74/60270}.
Public preprint: \url{https://arxiv.org/abs/2106.12017}.

\bibitem{LiWangZhou2025}
Y. Li, Z. Wang and X. Zhou,
\emph{A Note on Demailly's Transcendental Morse Inequalities Conjecture},
Frontiers of Mathematics \textbf{20} (2025), no.~5, 1195--1199.
DOI: \href{https://doi.org/10.1007/s11464-024-0145-y}{10.1007/s11464-024-0145-y}.
\url{https://journal.hep.com.cn/math/EN/10.1007/s11464-024-0145-y}.

\bibitem{Tosatti2026}
V. Tosatti,
\emph{Transcendental Morse inequality on K\"ahler manifolds},
Preprint, arXiv:2608.30153v2, 8 September 2026.
\url{https://arxiv.org/abs/2608.30153v2}.

\bibitem{Han2026}
J. Han,
\emph{A linear bound for Fujita's freeness conjecture},
Preprint, arXiv:2609.01574v2, 11 September 2026.
\url{https://arxiv.org/abs/2609.01574v2}.

\bibitem{BoucksomGuedjLu2025}
S. Boucksom, V. Guedj and C. H. Lu,
\emph{Volumes of Bott--Chern Classes},
Peking Mathematical Journal, published online 16 June 2025.
DOI: \href{https://doi.org/10.1007/s42543-025-00105-2}{10.1007/s42543-025-00105-2}.
\url{https://link.springer.com/article/10.1007/s42543-025-00105-2}.

\bibitem{XiaZhang2026}
M. Xia and K. Zhang,
\emph{A counterexample to the bounded mass property},
Preprint, arXiv:2608.21053v1, 21 August 2026.
\url{https://arxiv.org/abs/2608.21053v1}.


\bibitem{Serre1956}
J.-P. Serre, \emph{G\'eom\'etrie alg\'ebrique et g\'eom\'etrie analytique},
Annales de l'Institut Fourier \textbf{6} (1956), 1--42.
\url{https://www.numdam.org/articles/10.5802/aif.59/}.

\bibitem{AtiyahSinger1968}
M. F. Atiyah and I. M. Singer, \emph{The index of elliptic operators. III},
Annals of Mathematics (2) \textbf{87} (1968), no.~3, 546--604.
\url{https://annals.math.princeton.edu/1968/87-3/p07}.

\bibitem{Bott1959}
R. Bott, \emph{The stable homotopy of the classical groups},
Annals of Mathematics (2) \textbf{70} (1959), no.~2, 313--337.
\url{https://doi.org/10.2307/1970106}.
See also R. Bott, \emph{Quelques remarques sur les th\'eor\`emes de p\'eriodicit\'e},
Bulletin de la Soci\'et\'e Math\'ematique de France \textbf{87} (1959), 293--310,
\url{https://www.numdam.org/articles/10.24033/bsmf.1526/}.

\bibitem{Peterson1959}
F. P. Peterson, \emph{Some remarks on Chern classes},
Annals of Mathematics (2) \textbf{69} (1959), no.~2, 414--420.
Original issue contents: \url{https://www.jstor.org/stable/i307302}.

\bibitem{Hartshorne1966}
R. Hartshorne, \emph{Ample vector bundles},
Publications Math\'ematiques de l'IH\'ES \textbf{29} (1966), 63--94.
\url{https://www.numdam.org/item/PMIHES_1966__29__63_0.pdf}.

\bibitem{Berndtsson2009}
B. Berndtsson, \emph{Curvature of vector bundles associated to holomorphic fibrations},
Annals of Mathematics \textbf{169} (2009), no.~2, 531--560.
\url{https://doi.org/10.4007/annals.2009.169.531}.

\bibitem{Kleiman1966}
S. L. Kleiman, \emph{Toward a numerical theory of ampleness},
Annals of Mathematics (2) \textbf{84} (1966), no.~3, 293--344.
\url{https://annals.math.princeton.edu/1966/84-3/p01}.

\bibitem{Grauert1962}
H. Grauert, \emph{\"Uber Modifikationen und exzeptionelle analytische Mengen},
Mathematische Annalen \textbf{146} (1962), 331--368.
\url{https://doi.org/10.1007/BF01441136}.

\bibitem{Hironaka1964}
H. Hironaka, \emph{Resolution of singularities of an algebraic variety over a field of characteristic zero. I, II},
Annals of Mathematics (2) \textbf{79} (1964), 109--203 and 205--326.
For analytic desingularization, see also \cite{ArocaHironakaVicente2018}.

\bibitem{Demailly1992}
J.-P. Demailly, \emph{Regularization of closed positive currents and intersection theory},
Journal of Algebraic Geometry \textbf{1} (1992), no.~3, 361--409.
\url{https://www-fourier.univ-grenoble-alpes.fr/~demailly/manuscripts/regularization.pdf}.

\bibitem{JiShiffman1993}
S. Ji and B. Shiffman, \emph{Properties of compact complex manifolds carrying closed positive currents},
Journal of Geometric Analysis \textbf{3} (1993), no.~1, 37--61.
Author's publicly available version: \url{https://math.jhu.edu/~shiffman/publications/34%20closed%20positive%20currents%201993.pdf}.

\bibitem{GrauertRiemenschneider1970}
H. Grauert and O. Riemenschneider,
\emph{Verschwindungss\"atze f\"ur analytische Kohomologiegruppen auf komplexen R\"aumen},
Inventiones Mathematicae \textbf{11} (1970), 263--292.
\url{https://doi.org/10.1007/BF01403182}.

\bibitem{Reider1988}
I. Reider, \emph{Vector bundles of rank 2 and linear systems on algebraic surfaces},
Annals of Mathematics \textbf{127} (1988), no.~2, 309--316.
\url{https://annals.math.princeton.edu/1988/127-2/p04}.

\bibitem{Eckl2016}
T. Eckl, \emph{Numerical analogues of the Kodaira dimension and the Abundance Conjecture},
manuscripta mathematica \textbf{150} (2016), 337--356.
\url{https://doi.org/10.1007/s00229-015-0815-x}.
Public preprint: \url{https://arxiv.org/abs/1505.01262}.
See especially Theorem 0.7 (also Theorem 4.5); its assumption that a minimal model exists and its definition of numerical dimension are essential.

\bibitem{Ou2025}
W. Ou, \emph{A characterization of uniruled compact K\"ahler manifolds},
Preprint, arXiv:2501.18088, public version dated January 30, 2025.
\url{https://arxiv.org/abs/2501.18088}.


\bibitem{Demailly1982}
J.-P. Demailly,
\emph{Estimations $L^2$ pour l'op\'erateur $\bar\partial$ d'un fibr\'e vectoriel
holomorphe semi-positif au-dessus d'une vari\'et\'e k\"ahl\'erienne compl\`ete},
Annales scientifiques de l'\'Ecole Normale Sup\'erieure (4)
\textbf{15} (1982), no.~3, 457--511.
DOI: \href{https://doi.org/10.24033/asens.1434}{10.24033/asens.1434}.
\url{https://www.numdam.org/item/ASENS_1982_4_15_3_457_0/}.
See \S5 for complete-manifold solvability and \S8--9 for regularization of singular weights.

\bibitem{Chen2011}
B.-Y. Chen,
\emph{A simple proof of the Ohsawa--Takegoshi extension theorem},
arXiv:1105.2430v1, 12 May 2011.
\url{https://arxiv.org/abs/1105.2430}.

\bibitem{TakegoshiOhsawa1981}
K. Takegoshi and T. Ohsawa,
\emph{A vanishing theorem for $H^p(X,\Omega^q(B))$ on weakly $1$-complete manifolds},
Publications of the Research Institute for Mathematical Sciences
\textbf{17} (1981), no.~2, 723--733.
\url{https://ems.press/journals/prims/articles/3027}.
See \S3 for approximation of closed forms and \S4 for global gluing by corrections.

\bibitem{GreeneWu1975}
R. E. Greene and H. Wu,
\emph{Embedding of open Riemannian manifolds by harmonic functions},
Annales de l'Institut Fourier \textbf{25} (1975), no.~1, 215--235.
DOI: \href{https://doi.org/10.5802/aif.549}{10.5802/aif.549}.
\url{https://www.numdam.org/item/AIF_1975__25_1_215_0/}.

\bibitem{Ohsawa1984}
T. Ohsawa,
\emph{Completeness of noncompact analytic spaces},
Publications of the Research Institute for Mathematical Sciences
\textbf{20} (1984), no.~3, 683--692.
\url{https://ems.press/journals/prims/articles/3175}.

\bibitem{AndreottiGrauert1962}
A. Andreotti and H. Grauert,
\emph{Th\'eor\`emes de finitude pour la cohomologie des espaces complexes},
Bulletin de la Soci\'et\'e Math\'ematique de France
\textbf{90} (1962), 193--259.
DOI: \href{https://doi.org/10.24033/bsmf.1581}{10.24033/bsmf.1581}.
\url{https://www.numdam.org/item/BSMF_1962__90__193_0/}.

\end{thebibliography}

\end{document}
